\documentclass[11pt]{article}
\usepackage[a4paper,margin=27mm]{geometry}
\usepackage[T1]{fontenc}
\usepackage{lmodern,microtype}
\usepackage{amsmath,amssymb,amsthm,mathtools}
\usepackage{enumitem,booktabs,array,longtable}
\usepackage[hidelinks,hypertexnames=false]{hyperref}
\hypersetup{pdftitle={The Contamination Triangle}}

\newtheorem{theorem}{Theorem}[section]
\newtheorem{lemma}[theorem]{Lemma}
\newtheorem{proposition}[theorem]{Proposition}
\newtheorem{corollary}[theorem]{Corollary}
\newtheorem{conjecture}[theorem]{Conjecture}
\theoremstyle{definition}
\newtheorem{definition}[theorem]{Definition}
\theoremstyle{remark}
\newtheorem{remark}[theorem]{Remark}

\newcommand{\R}{\mathbb R}
\newcommand{\one}{\mathbf 1}
\newcommand{\Di}{\Delta^{\circ}}
\newcommand{\M}{\mathcal M}
\newcommand{\C}{\mathcal C}
\newcommand{\cl}{\operatorname{cl}}
\newcommand{\co}{\operatorname{conv}}
\newcommand{\supp}{\operatorname{supp}}
\newcommand{\relint}{\operatorname{relint}}
\newcommand{\aff}{\operatorname{aff}}
\newcommand{\lin}{\operatorname{line}}
\newcommand{\sd}{\mathbin{\triangle}}
\newcommand{\lp}{\underline{P}}
\newcommand{\up}{\overline{P}}
\newcommand{\dH}{d_{\mathrm H}}
\newcommand{\src}[1]{\texttt{#1}}
\newcommand{\status}[1]{\par\smallskip\noindent\textbf{Status:} #1\par\smallskip}
\newcommand{\Ltwo}{\mathcal L_2^{\mathrm s}}
\newcommand{\Nb}{N_0^{\mathrm b}}

\title{The Contamination Triangle}
\author{Research note on coverage under decomposable domains}
\date{16 October 2026 (rounds 14--19; audit of 15 October)}

\begin{document}
\emergencystretch=2em
\maketitle

\section{Introduction}\label{sec:intro}

\paragraph{The coverage question.}
A dominance argument for coherence has two halves. The first shows that, for any ranking of models that satisfies certain axioms, every model is weakly dominated by an undominated coherent model. For finitely many worlds this is the KKM theorem of \src{Loss Dominance Generalized v2} (Theorem~\ref{thm:kkm} below). The second half, \emph{coverage}, asks that every coherent model be undominated for at least one ranking that satisfies the axioms. Without coverage the argument shows only that the permissible models lie among the coherent ones. It does not show that every coherent model is permissible.

In the setting adopted here, models are sets of desirable gambles. A ranking compares models by the advice they give on binary decision problems. At each world, good advice is better than silence, silence is better than bad advice, and bad advice is better than contradictory advice (guidance, (G)). Comparisons depend only on the problems on which two models differ (independence, (IB)). Rankings are preorders, which may be incomplete and need not be numerical. Each ranking evaluates the models in a \emph{decomposable domain}, its class of evaluable models. Restricting the domain restricts only which models are evaluated; the decision problems are never altered or removed. To stop a target being protected merely by excluding its borderline competitors, the domain must contain every model that differs from the target only near the boundary of its cone of desirable gambles (condition (LF)).

\paragraph{Why the contamination triangle.}
Coverage is known for precise models, for every coherent model when there are two worlds, for normalised boxes, and for a large class of interior polygons (Section~\ref{sec:positive}). In three worlds the simplest imprecise models not known to be covered are those whose credal set is a triangle in the interior of the simplex. The most natural of these is the $\varepsilon$-contamination of the uniform prior,
\[
T_1\coloneqq(1-\varepsilon)u+\varepsilon\Delta,\qquad u=(\tfrac13,\tfrac13,\tfrac13),\quad 0<\varepsilon<1 .
\]
Its cone of desirable gambles has three flat facets that meet along three corner rays. Additive rankings built by mixing gauges cover hexagons and 9-gons arbitrarily close to $T_1$, but provably no triangle. In earlier work on coherence-preserving maps the same target was shown to have no continuous realiser. So $T_1$ is where the structure of the axioms, and not the convenience of a construction, decides coverage.

\paragraph{What is shown.}
\begin{enumerate}[leftmargin=*]
\item \emph{Additive rankings} (Sections~\ref{sec:reduction}--\ref{sec:main}). If an additive ranking on the $U$-domain (the class of measurable subsets of a fixed open cone $U$ around $\cl D_{T_1}\setminus\{0\}$) covers $T_1$, the Bayes models of its interior weightings form a Hausdorff-continuous family of credal sets with half-plane charts that passes through $T_1$. In three worlds the converse holds too. The main theorem (Theorem~\ref{thm:main}, Corollary~\ref{cor:coverage}) is that on every decomposable domain whose pairs near $\cl D_{T_1}$ (those with an orientation in a cone $D_{K'}$, $K'\subseteq\relint T_1$) are non-adjacent, and in particular on the $U$-domain, no additive ranking covers $T_1$. This holds for every $\varepsilon\in(0,1)$, at every interior target weighting, and with no regularity assumption. The proof applies to every all-facing triangle in general position, hence to every interior $\varepsilon$-contamination.
\item \emph{General decomposable domains} (Sections~\ref{sec:general} and~\ref{sec:gm}). Without non-adjacency, in every additive cover of $T_1$ the advice range of almost every mixed pair contains both acceptances, no vertex is pinned near the target weighting, and the d-faces of the pairs that cut the box witnesses at a vertex converge into $\partial\Delta$, with apexes at that vertex. Additive coverage of $T_1$ on some (LF) decomposable domain is reduced to a purely geometric problem, CF$(T_1)$, about Hausdorff-continuous credal families. The main new result (Theorem~\ref{thm:gm}) is that for every $\varepsilon\in(0,1)$ and every interior target weighting, no strictly guiding additive ranking with integral penalties covers $T_1$ on any decomposable domain with (LF). Hence CF$(T_1)$ has no solution, at the status of Theorem~\ref{thm:ccibu} (Corollary~\ref{cor:CFnone}). Three independent proofs are given, each confirmed by one referee; the inputs they share have had two independent readings. For bounded charge penalties the same holds given interior Bayes minimisers, along one of the three proofs (Corollary~\ref{cor:gmcharges}).
\item \emph{Coherent almost-preference rankings} (Section~\ref{sec:cp}). If each world's ranking is a coherent almost-preference in Walley's sense, then on the $U$-domain no such ranking covers $T_1$, complete or not, and with arbitrary finitely additive members: the integral field of Theorem~\ref{thm:integral} falls under Theorem~\ref{thm:M0closed}(a) (Corollary~\ref{cor:CPaU}; proved, its glue refereed once). Equivalently, on $\M_U$ every (CP-a) world adds up, and systems in which each world is (CP-a) or adds up are excluded too. On other non-adjacent domains coverage is impossible when every world is complete (under a mild domain condition (UN)) or strict comparisons are detected by countably additive previsions, and in general it reduces to one property (U) of a single integral field.
\item \emph{Rankings that add up} (\S\S\ref{sec:M0}--\ref{sec:M0-lex}). If every world of an arbitrary ranking on the $U$-domain adds up on the relevant comparisons, a supporting additive assessment exists whose exceptional weightings form a nonempty countable union of chords of $\Delta$, each cutting off one vertex. At those weightings coherent and incoherent minimisers coexist and the assessment cannot tell them apart; a cover would require the original ranking to do so. A two-world construction shows that a ranking can resolve such a tie while keeping every axiom. \textbf{New in round 11} (Theorem~\ref{thm:M0closed}): in three worlds this never helps. No strictly signed field on the $U$-domain whose Bayes models are coherent at almost every interior weighting has $T_1$ as a Bayes model at an interior weighting; hence no ranking whose worlds all add up covers $T_1$ on the $U$-domain, with or without (E$_\psi$). The chain of Section~\ref{sec:main} survives atom chords away from the target weighting; at a target weighting on a chord, the lower prevision of every coherent model that the assessment cannot tell apart from $T_1$ vanishes at the six crossings of $\partial D_{T_1}$, and the atoms that carry the ties contradict this. The new parts have been read by two independent referees; the ranking form rests on imports read once.
\item \emph{Building ordinal rankings directly} (\S\ref{sec:open-ordinal}). An exact extension theorem with a limit rule replaces guessed sections; the extension lemma as previously posed is false; an additive witness gives a positive but nearly tautological lemma. Flexible worlds have no finite obstruction, two-level layering is locally single-level additive, and in two worlds with a nested selection every world with a nonzero leading charge, and every complete world, adds up (for leading charges, round 13 removes the nestedness hypothesis, in any number of worlds: \S\ref{sec:open-core}). What remains open is pinned down in \S\ref{sec:two-questions}.
\item \emph{Worlds that do not add up: the two-world core} (\S\ref{sec:open-core}, new in round 12, extended in round 13). No index-local argument can exclude $T_1$ (a locality barrier), but every cover must show non-additive behaviour at $T_1$ itself: either some world's comparisons out of $T_1$ admit no strictly positive endorsing weight, or unendorsed relevant comparisons occur against incoherent models arbitrarily close to $T_1$ (Theorem~\ref{thm:R3}). On the edges of $\Delta$ the selected models form two-world systems, so the two-world results transfer to three worlds under an edge-nestedness hypothesis (Nest$_e$). \textbf{New in round 13:} a nonzero leading charge is strictly positive, and its world adds up, in any number of worlds and with no nestedness hypothesis (Theorem~\ref{thm:S21}, read by two referees); hence no cover of $T_1$ on $\M_U$ has every world (CP-a), or with a nonzero leading charge, or complete and edge-nested (Corollary~\ref{cor:Q37}). In two worlds a nonzero leading charge or monotone continuity at either world forces the selection to be nested (Theorems~\ref{thm:S22}, \ref{thm:MC}), and without (C) it need not be (Proposition~\ref{prop:R71}). In two worlds with a nested selection, the sets that a failing world cannot charge are initial segments of the chain, and simple certificates come in mirrored pairs, one at each world. The rate once proposed as the missing lemma is a single joint property of the two worlds that holds next to any additive world, so it cannot be that lemma (Theorem~\ref{thm:RATE}, Corollary~\ref{cor:SR}), and its linear form fails in a system with every axiom (Proposition~\ref{prop:ExA}). Two-world systems satisfying every axiom except closedness of the sections, at one world (a world that does not add up) or at both (a mirrored simple pair), are constructed explicitly (Propositions~\ref{prop:Q12}, \ref{prop:S31}). Until round 16 no world failing to add up was known in a full system, for any target; round 16 constructs one in two worlds (see the end of this item). \textbf{New in rounds 14--15} (\S\ref{sec:open-core}(f)): in three worlds, leading densities on the two edge regions of a world glue through the link class $R_{ij}\cap R_{ik}$ and propagate to all its comparisons, so the world adds up (Theorem~\ref{thm:B13}, refereed once); hence every world of a cover that does not add up has no leading density on one of its edge regions (Corollary~\ref{cor:B14}). A total world adds up as soon as it is almost representable, with no nestedness (Corollary~\ref{cor:B32}), and order continuity on edges excludes non-additive top levels with dense null sets, including those of signature type (Theorem~\ref{thm:B35}, Corollary~\ref{cor:B36}), all refereed once. Glued credal worlds either violate strict (G) or add up (Theorems~\ref{thm:C11}, \ref{thm:C16}); a full two-world system can have a world with no leading density, although it adds up (Theorem~\ref{thm:C21}); and on an edge, closed choice forces density unanimities to be comonotone (Theorem~\ref{thm:C23}); these round-15 results were refereed once (\src{C12\_referee}), with small repairs. On every decomposable domain, every face model at a weighting where $D_{T_1}$ is Bayes for an additive field with ($\gamma$) is crossing-pinned and contains $T_1$, and there are no apex ties there (Theorem~\ref{thm:B21}, refereed once). \textbf{New in round 16} (\S\ref{sec:open-core}(e), (f)): a full two-world system on $\M_U$, with every axiom and closed sections for every model, in which one world does not add up, of type (M2), next to an additive world (Theorem~\ref{thm:D21}, confirmed by two independent referees). It refutes Conjecture N17.B$'$, already in two worlds with a nested selection, and \src{Q1}'s Conjecture Q1.C1, and it answers question (R2) negatively. The failure is exactly a failure of integrability at a protected end of the chain: the same world, built on the design constructed there, is endorsed by the $\sigma$-finite density $1/(1-x)$, which is integrable on every initial segment of the chain (Proposition~\ref{prop:D2sigma}). \textbf{Audit of 15 October:} that end consists of two rays, one of which, $\R_+(1,0)$, lies in the interior of $U$, and every endorsing measure of that world, under any reading of endorsement with a mild property (R), puts infinite mass on every neighbourhood of that ray (Lemma~\ref{lem:F11}, Theorem~\ref{thm:F13}; observed by an external auditor, proved in note \src{F1}, refereed once). So Conjecture~\ref{conj:sigmaB}, the form of N17.B$'$ with densities locally integrable on $U$, is false; its repair with local integrability on the mixed problems only (Conjecture~\ref{conj:F1Cp}) is open. Next to an additive world every $\sigma$-finite endorser is a nondecreasing function of the chain coordinate (Theorem~\ref{thm:F12}); no certificate has all its targets in one compact subset of $U$ on which the payoffs are bounded away from zero (Theorem~\ref{thm:F14}, Corollary~\ref{cor:F15}); and there are full two-world systems in which every generator and the protected set lie in a region where both payoffs are bounded away from zero, with the protection at $\partial U$ (Theorems~\ref{thm:F16}, \ref{thm:F18}; refereed once). For $T_1$, the argument of Theorem~\ref{thm:M0closed} needs only the dominator comparisons of one admissible shadow (Theorem~\ref{thm:E11}); protection at ends on which $T_1$ is eventually correct at the protected world is invisible to it (Lemma~\ref{lem:E14}); and a cover in which every world is an end-protected decoration with integrable base densities that are, jointly, heavy at the ends is excluded (Theorem~\ref{thm:E15}). These results are proved in note \src{E1} and audited in note \src{F2}, whose audit was refereed once; they exclude less than \src{E1} first stated (part~(f) of \S\ref{sec:open-core}). \textbf{New in round 17} (\S\ref{sec:open-core}(f$'$)): failures of adding up that can be confined to sets of arbitrarily small measure are invisible to the argument of Theorem~\ref{thm:M0closed}. Next to two worlds with integrable densities that endorse all their relevant comparisons, no world with an a.e.\ finite endorser of such localized comparisons can be the third world of a cover (Theorem~\ref{thm:H12}); the heaviness hypothesis of Theorem~\ref{thm:E15} can be replaced by eventual $T_1$-correctness of the ends (Theorem~\ref{thm:H13}); and with finitely many ends and integrable base endorsers a cover needs a directed cycle in a conflict digraph between the worlds (Theorem~\ref{thm:H23}, Proposition~\ref{prop:HR1}). A full two-world system in which neither world adds up exists (Theorem~\ref{thm:H24}). These results were refereed once, and Proposition~\ref{prop:HR1} is proved in the referee report. Note \src{J1} (refereed once in round 19) excludes every such cycle whose ends can be put in an aligned order, which leaves, for a $2$-cycle next to a world without ends, the doubly misaligned cycles of Question~J1.Q; its referee reduces the rest of the class to three cases (Theorem~\ref{thm:JR2}, Corollary~\ref{cor:JR3}, Question~J1.Q$'$). \textbf{New in round 19} (\S\ref{sec:open-core}(f$''$)): a world that does not add up either carries a finite certificate, or has a non-null set on which no localized bounded endorser has a positive density (its non-localizable blind set), or is localizable with an integrable density (Theorem~\ref{thm:K31}, by the Hahn--Banach and Yosida--Hewitt theorems). The second kind is consistent with every axiom except closed sections (Theorem~\ref{thm:K33}); whether it exists in a full system is open. Every finite certificate contains a rigid comparison that competitiveness does not force (Theorem~\ref{thm:K21}, Corollary~\ref{cor:K22p}), and over Brier-type additive bases the natural decoration scheme for building one is empty (Theorem~\ref{thm:K26}, Proposition~\ref{prop:KR62}). If every world is localizable, a cover needs sheltering or infinite minority mass (Theorem~\ref{thm:K35}, Corollary~\ref{cor:K35p}), and for end-protected worlds with countably many ends only cyclic configurations with no aligned order remain (Theorems~\ref{thm:K14}, \ref{thm:K15}). These results were refereed once each. On every decomposable domain with (LF) at $D_{T_1}$, the target weighting of a strictly guiding integral additive field with ($\gamma$) carries no non-null ties, so the full-atom residual is empty (Theorem~\ref{thm:D1}, refereed once).
\item \emph{Weakened guidance} (\S\ref{sec:open-ties}). The theorem of Section~\ref{sec:main} survives permitted ties on problems with small stakes, for every band width in the sup norm and for Euclidean band widths below an explicit threshold.
\item Section~\ref{sec:open} states two concrete questions (\S\ref{sec:two-questions}), both now answered negatively (the first on the $U$-domain, the second on every decomposable domain with (LF)), states the one remaining route on the $U$-domain (a world that does not add up), and lists what is excluded and what is open; Section~\ref{sec:positive} recalls the positive results.
\end{enumerate}
\emph{Summary for $T_1$ on the $U$-domain.} Excluded so far: additive rankings; rankings whose worlds all add up (including the lexicographic class $\Ltwo$); (CP-a) rankings, complete or not; and rankings all of whose worlds are (CP-a), have a nonzero leading charge, are total and almost representable, or (round 15) are glued credal worlds (Corollary~\ref{cor:Q37}, Corollary~\ref{cor:B32}, Theorem~\ref{thm:C16}). So a cover needs a world that does \emph{not add up}. \textbf{Round 19} (\S\ref{sec:open-core}(f$''$), (g)) says what such a world can be. By Theorem~\ref{thm:K31}, it either carries a finite certificate (mode (M1)); or it has no finite certificate and a non-null non-localizable blind set $\Lambda_\infty$, a set on which, whatever family of sets of vanishing measure the comparisons are localized at, every bounded endorser of the localized comparisons has density zero; or it has no finite certificate and is localizable with an integrable density. In the first case, every certificate contains a rigid comparison that is not forced by competitiveness (Theorem~\ref{thm:K21}, Corollary~\ref{cor:K22p}), and over additive bases of Brier type the natural construction scheme, decorating a Bayes base, is empty (Theorems~\ref{thm:K25}, \ref{thm:K26}, Proposition~\ref{prop:KR62}); no full system with a certificate is known. In the second case, whether a full system can have a non-null $\Lambda_\infty$ is open (Question~Q$_{\mathrm{K3}}$); no axiom other than closedness of the sections excludes it (Theorem~\ref{thm:K33}), and next to two integrable worlds it needs a positive answer to the additive Question~Q$_{\mathrm A}$ (Theorem~\ref{thm:K37}). In the third case, if every world is localizable, the patch theorem (Theorem~\ref{thm:K35}, Corollary~\ref{cor:K35p}) leaves a cover only two mechanisms: \emph{sheltering}, two disagreeing worlds that protect a common region at every depth, which for end-protected worlds is Question~J1.Q$'$ (finitely many ends) and Question~J1.Q$'_\infty$ (countably many ends); or \emph{infinite minority mass}, which needs a world that is not localizable with an integrable density and is excluded if Question~Q$_{\mathrm{K3}}$ has answer No (\src{K3\_referee}). (Round 19 results are proved in notes \src{K1}--\src{K3} and refereed once; Propositions~\ref{prop:KR62}, \ref{prop:KQA}, Theorem~\ref{thm:JR2} and Corollary~\ref{cor:JR3} are proved in referee reports, one reader.) Worlds that do not add up exist in full systems: round 16 gives a full two-world system on $\M_U$ in which one world does not add up, next to an additive world (Theorem~\ref{thm:D21}; confirmed by two independent referees). It refutes Conjecture N17.B$'$ in two worlds and Conjecture Q1.C1, and answers (R2) negatively. Its failure is exactly non-integrability at a protected end: a $\sigma$-finite density, non-integrable only at that end, endorses every relevant comparison of that world (Proposition~\ref{prop:D2sigma}). That end contains the ray $\R_+(1,0)$, which lies in the interior of $U$, and no endorsing measure of that world is finite on compact subsets of $U$; so Conjecture~\ref{conj:sigmaB} is false (Theorem~\ref{thm:F13}; audit of 15 October, observed by an external auditor, proved in note \src{F1}, refereed once). No world that does not add up is known in three worlds or in a cover. Every world of a cover that does not add up (i) is not (CP-a) and has no nonzero leading charge; (ii) has, on one of its two edge regions, \emph{no leading density} (Theorem~\ref{thm:B13}, Corollary~\ref{cor:B14}; for countably additive densities: the gluing of finitely additive leading charges is open); (iii) is \emph{not a glued credal world}, and if it leads a finitely generated credal pair on its edges it is not total and carries a finite obstruction inside the glued band (Theorem~\ref{thm:C16}, Corollary~\ref{cor:C18}, Proposition~\ref{prop:C19}; round 15, refereed once); (iv) if total, is not almost representable on $M$ and has an edge family at its vertex that is not nested (Corollary~\ref{cor:B32}; CHECKED); that it cannot be total at all would follow from Conjecture~\ref{conj:B3C}, which is open. Moreover (v) a cover is not one in which every world is an end-protected decoration whose base has an integrable endorsing density, with these densities jointly heavy at the ends (Theorem~\ref{thm:E15}); world~1 of Theorem~\ref{thm:D21} has this form in its two-world system (Proposition~\ref{prop:E16}), but no result shows that a three-world installation of it does (round 16; proved in note \src{E1}, audited in note \src{F2}, audit refereed once). (vi) A cover does not have two worlds with integrable densities that endorse all their relevant comparisons together with a third world that has an a.e.\ finite density endorsing, in the cancellation form, its relevant comparisons that add no errors on some member of a family of sets of vanishing measure (Theorem~\ref{thm:H12}). So such a third world is neither an end-protected world with countably many ends whose base has such an endorser, end-heavy or not, nor a world that adds up $\sigma$-finitely. (vii) A cover is not one in which every world has an integrable density that endorses its comparisons localized on sets where $T_1$ is correct at that world (Theorem~\ref{thm:H13}); so (h) in (v) may be replaced by eventual $T_1$-correctness of every end, which (h) implies. (viii) If every world is end-protected with finitely many ends and has an integrable base endorser, a cover needs a directed cycle in the conflict digraph (Theorem~\ref{thm:H23}, Proposition~\ref{prop:HR1}), hence at least two worlds with $T_1$-wrong ends and at least two worlds without an integrable full endorser (Corollary~\ref{cor:H23p}); and no order of the worlds may be aligned (Theorem~\ref{thm:J12}), which closes the template of \src{H2} (Corollary~\ref{cor:J13}) and leaves the cases (R-i)--(R-iii) of Theorem~\ref{thm:JR2} and Corollary~\ref{cor:JR3}: for a $2$-cycle whose third world is incident to no arrow, only doubly misaligned cycles (Question~J1.Q), and in general Question~J1.Q$'$. (ix) With countably many ends and integrable base endorsers the same holds with alignment end by end (Theorems~\ref{thm:K14}, \ref{thm:K15}; Question~J1.Q$'_\infty$). (Round 17: (vi)--(viii) and Corollary~\ref{cor:H23p} are proved in notes \src{H1}, \src{H2} and refereed once in \src{H12\_referee}; Proposition~\ref{prop:HR1} is proved in that report; Theorem~\ref{thm:J12} and Corollary~\ref{cor:J13} are proved in note \src{J1} and were refereed once in round 19 (\src{J1\_referee}). (ix) is proved in note \src{K1}, refereed once.) A full two-world system in which neither world adds up exists (Theorem~\ref{thm:H24}, refereed once). Theorem~\ref{thm:M0closed} needs, for one admissible shadow, only its dominator comparisons into incoherent Bayes models (Theorem~\ref{thm:E11}), which may be taken within any $\eta$ of $T_1$ (Theorem~\ref{thm:E12}). At an end of a protected world on which $T_1$ is eventually correct at that world, a heavier shadow makes every Bayes model correct deep in the end (Lemma~\ref{lem:E14}), so the protected comparisons there are invisible to that argument. Under the hypotheses of Lemma~\ref{lem:E13}(a) every end is eventually $T_1$-correct at its world; for ends in a collar of $\partial U$ inside $\{h_i<0\}$ this holds with no hypothesis on any world (Proposition~\ref{prop:F21}(a)); and under (Nest$_e$) on both edges, with the further hypotheses of Proposition~\ref{prop:F21}(b), the ends lie at the near-unanimous boundary or at $\partial U\cap\{h_i<0\}$. Nothing proved confines them to near-unanimous limits. In a cover, for every admissible shadow, some world has unendorsed dominator comparisons into incoherent Bayes models on a non-null set of weights; this follows from (H$_U$) (Theorem~\ref{thm:E11}). For the reweighted shadows of Lemma~\ref{lem:E14}, such a comparison at an end-protected world whose ends are eventually $T_1$-correct is a base comparison, eventually zero at every end; so the failures must be carried by base comparisons or by comparisons at another world. Round 17: the same holds for the dominator comparisons into every regular model within a fixed cost of the Bayes model (Lemma~\ref{lem:H10}), so the failures cannot be confined to sets of vanishing measure at a world whose density is integrable off them. The remaining problems (\S\ref{sec:open-core}(g)) are: the two-world questions (R1), (R4) and Conjecture~\ref{conj:NDL}, which act on the edge systems ((R2) is answered); Conjecture~\ref{conj:F1Cp} (the repair of the refuted Conjecture~\ref{conj:sigmaB}) and the a.e.-finite target of \S\ref{sec:open-core}(e); Conjecture~\ref{conj:B3C} for total worlds (first test Conjecture~\ref{conj:B3Cp}); Question~C1.Q for worlds whose edge orders lead a credal pair; Conjecture~\ref{conj:C2C} for density-unanimity edges in three worlds; edge orders of other kinds without a leading density, which no argument confined to one edge can exclude (Corollary~\ref{cor:C25}); the questions of round 19 that follow the three cases above (\S\ref{sec:open-core}(g)): finite certificates (Question~Q-K2, Conjecture~\ref{conj:K26p}); non-localizable blind sets (Question~Q$_{\mathrm{K3}}$, its two-world form Conjecture~\ref{conj:K3C}, and Question~Q$_{\mathrm A}$); and sheltering (Questions~J1.Q, J1.Q$'$ and J1.Q$'_\infty$ for end-protected worlds, and sheltering between localizing families in general); and case (i) of Theorem~\ref{thm:R3}, which after Theorem~\ref{thm:E12} is no longer needed to obtain unendorsed comparisons near $T_1$. \emph{Off the $U$-domain} only additive rankings are excluded (Theorem~\ref{thm:gm}). For rankings that add up on a general decomposable domain, the transfer of the supporting assessment (Claim~B2.0) is not yet written. Independently of it, under (LF) at $D_{T_1}$ the target weighting of a strictly guiding integral additive field with ($\gamma$) is never exceptional (Theorem~\ref{thm:D1}, round 16, refereed once), so the full-atom residual, which occurs without (LF) (Propositions~\ref{prop:B2ex}, \ref{prop:D10}), is empty. Conditional on Claim B2.0, apex ties at the target weighting are excluded on every decomposable domain (Corollary~\ref{cor:B22}), and under (LF), at the status of Claim B2.0 with a reduced list of imports, $T_1$ is never among tied Bayes models of the supporting assessment (Corollary~\ref{cor:D17}). What remains is the branch $p^*\notin E$: the case $E=\varnothing$ at the status of Theorem~\ref{thm:ccibu}, and the case $E\ne\varnothing$, where route~I1 is audited only conditionally (\S\ref{sec:open-core}(f)). Rankings with a world that does not add up have not been examined off the $U$-domain.

\paragraph{Status of the results.}
Every result below carries a status label right after its statement. ``Refereed'' or ``confirmed'' means that a reader who did not write the note re-derived each step and returned the verdict CORRECT, CONFIRMED, CORRECT WITH REPAIR or CONFIRMED WITH REPAIR. All requested repairs are incorporated in the statements given here; where a referee's repair conflicts with the note, the referee's version is used. A result first proved \emph{in} a referee report has been checked by nobody else unless a second reading is recorded. CHECKED means verified with a proof sketch; NUMERICS means supported only by computation. Notes are cited by their file names (for example \src{N15} for \src{N15\_notes.md}); the Sources section at the end lists them.

{\small
\begin{longtable}{>{\raggedright\arraybackslash}p{6.4cm}>{\raggedright\arraybackslash}p{2.1cm}>{\raggedright\arraybackslash}p{5.6cm}}
\toprule
Result & Where & Status\\
\midrule
\endhead
Additive cover $\Rightarrow$ unique coherent interior Bayes family through the target & Prop.~\ref{prop:necessary} & proved; refereed once (R6)\\
Reduction to a credal family with half-plane charts & Prop.~\ref{prop:reduction} & proved; refereed once (N14)\\
Closed choice from interior Bayes uniqueness ($n=3$) & Thm.~\ref{thm:ccibu} & author-checked; re-derived without error in N13; no formal referee report\\
\textbf{No additive cover of $T_1$ (all-facing triangles) on non-adjacent domains} & Thm.~\ref{thm:main}, Cor.~\ref{cor:coverage} & \textbf{proved; refereed independently twice (N18, N20)}; imported lemmas refereed once each (N14, R12, R18)\\
Same with bounded charge penalties, given interior minimisers & Cor.~\ref{cor:charges} & proved (N14 Thm C2, in a referee report; splice step refereed in N24); use refereed (N18, N20)\\
General decomposable domains: compact collar replaces the conic $N_0$; ray transfer of charts & \S\ref{sec:general} preamble, Lem.~\ref{lem:raytransfer} & proved (N27); confirmed by one referee (N27 referee)\\
General decomposable domains: (TS) $\Rightarrow$ structure; dichotomy & Thms.~\ref{thm:TS}, \ref{thm:dichotomy} & proved; refereed once (N23); branch (II) also uses Cor.~\ref{cor:R232}\\
Lemma R23.1, Cor.\ R23.2, Cor.\ R23.4 & Lem.~\ref{lem:R231}, Cors.~\ref{cor:R232}, \ref{cor:R234} & proved in a referee report (N23); second reading CORRECT (CHECK\_REPORT.md)\\
\textbf{(TS) holds in every cover; (E5a) is empty} & Thm.~\ref{thm:TSalways}, Cor.~\ref{cor:TSalways} & \textbf{proved (N28); confirmed by two independent referees (N28 referee, TS second referee)}, with its inputs\\
Cut apexes at the vertex (cut d-faces converge into $\partial\Delta$); vertex-free interiors; non-precise credal sets on crossing chords & Thm.~\ref{thm:apex}, Prop.~\ref{prop:vertexfree} & proved (N28); confirmed by two referees (N28 referee, TS second referee; wording repair in Thm.~\ref{thm:apex})\\
Global one-sidedness of chords & Lem.~\ref{lem:R1N28} & proved in a referee report (N28 referee); second reading CONFIRMED (TS second referee)\\
Reduction to CF$(T_1)$ & Prop.~\ref{prop:CF} & CHECKED (N28); ($\Rightarrow$) in its every-$q$ form confirmed with repair by two referees (N28 referee, TS second referee); converse at the status of Thm.~\ref{thm:ccibu}\\
Diagonal symmetries & Prop.~\ref{prop:diag} & proved (N28), confirmed; ranking-level form proved in a referee report (N28 referee), one reader\\
\textbf{No additive cover of $T_1$ on any (LF) decomposable domain (integral penalties)} & Thm.~\ref{thm:gm} & \textbf{proved; three independent routes (I1, I2, C2/C1), each confirmed by one referee (I1, I2, C1C2 referees)}; shared inputs read twice, with the exceptions in Table~\ref{tab:gmdeps}\\
Same with bounded charge penalties, given interior minimisers & Cor.~\ref{cor:gmcharges} & proved along route I2 only (I2 referee, via N14 Thm.~C2); I1's charge clause CHECKED; Thm.~\ref{thm:D} not checked for charges\\
CF$(T_1)$ has no solution & Cor.~\ref{cor:CFnone} & at the status of Thm.~\ref{thm:ccibu} (through Prop.~\ref{prop:CF}($\Leftarrow$))\\
No straddling on the positive side of a chord; forced precise collars; chord structure & Lem.~\ref{lem:nostraddle}, Thm.~\ref{thm:collar}, Cor.~\ref{cor:chord} & proved independently in I2 and C2; each proof confirmed by one referee\\
Precise regions are constant or diagonal-projective; one map for all precise priors & Thms.~\ref{thm:D}, \ref{thm:T} & proved (PropD notes), not yet refereed; the constant-or-injective dichotomy proved independently in two referee reports (Prop.~P, Prop.~D$^*$), one reader each\\
No injective precise open region (sector caps excluded) & Thm.~\ref{thm:A} & proved (C1); confirmed with repair (C1C2 referee)\\
No locally definable solution of CF$(T_1)$ & Rem.~\ref{rem:B} & proved (C1); confirmed with repair (C1C2 referee); superseded\\
(CP-a): full support of countably additive parts is automatic & Thm.~\ref{thm:fullsupport} & proved; refereed once (N24)\\
(CP-a): integral reduction under (CP-a) alone & Thm.~\ref{thm:integral} & proved; refereed once (N24)\\
\textbf{(CP-a) with complete worlds: no cover of $T_1$} & Thm.~\ref{thm:complete}, Cor.~\ref{cor:complete} & \textbf{proved; refereed once (N24)}\\
(CP-a) with countably additive detection: no cover & Cor.~\ref{cor:DCA} & proved; refereed once (N13, N24)\\
General (CP-a): cover forces a tie weight; (U) as continuity & Cor.~\ref{cor:residual}, Prop.~\ref{prop:Ucont} & residual proved, refereed once (N24); continuity form proved by the referee, one reader; on $\M_U$ superseded by Cor.~\ref{cor:CPaU}\\
(U) for every (CP-a) system & Conj.~\ref{conj:U} & open on non-adjacent domains other than $\M_U$; on $\M_U$ its hypothesis is never met (Cor.~\ref{cor:CPaU})\\
M0: supporting assessment with properties ($\alpha$)--($\delta$) & Prop.~\ref{prop:M0} & proved; inputs refereed once or proved in referee reports (N13, N19, N24)\\
Structure of the exceptional set: countably many atom chords & Thm.~\ref{thm:Estructure} & proved (N25), confirmed by one referee; ``no fan points'' proved in a referee report (N25 referee), one reader\\
$p^*\in E(\psi)$ blocks the chain of \S\ref{sec:main}; genericity & Thm.~\ref{thm:pstarE}, Prop.~\ref{prop:generic} & proved (N25); confirmed with repair (N25 referee)\\
Consequences of (E$_\psi$): ideal tie-breaking, transport along chords & Prop.~\ref{prop:Epsi} & proved (N25); confirmed by one referee\\
Residuals R-a$'$, R-b$'$ of M0 & \S\ref{sec:M0-residual} & \textbf{closed} (Thm.~\ref{thm:M0closed})\\
Quasi-linearity of tie-breaking families; ball lemma & Lems.~\ref{lem:quasilin}, \ref{lem:ball} & proved (P4); confirmed by two referees (M0, M0 second); Prop.~F steps 1--4 read three times\\
Crossing pinning of face models (any interior target) & Prop.~\ref{prop:crosspin} & proved (M0 referee, Prop.~R.1); strengthened and confirmed (M0 second referee)\\
Chain of \S\ref{sec:main} survives atom chords; $p^*\in E$ & Prop.~\ref{prop:chainchords} & proved (P4); confirmed by two referees\\
Carrying atoms at $p^*\in E$; short argument & Prop.~\ref{prop:carrying} & proved (P4); confirmed by two referees\\
\textbf{No additive field with ($\gamma$) has $T_1$ as an interior Bayes model on $\M_U$} & Thm.~\ref{thm:M0closed}(a) & \textbf{proved; new parts read by two independent referees (M0, M0 second)}\\
\textbf{No ranking that adds up covers $T_1$ on $\M_U$} & Thm.~\ref{thm:M0closed}(b), Cor.~\ref{cor:M0closed} & \textbf{proved}, at the status of Prop.~\ref{prop:M0} without (E$_\psi$) (imports N13 Prop.~B, N19 Prop.~R3(a), read once)\\
Same for every all-facing triangle in general position & Thm.~\ref{thm:M0closed} status & CHECKED (both M0 referees)\\
\textbf{No (CP-a) ranking covers $T_1$ on $\M_U$} & Cor.~\ref{cor:CPaU}(a) & \textbf{proved}; glue refereed once (CPaU referee); inputs Thm.~\ref{thm:integral} (N24) and Thm.~\ref{thm:M0closed}(a) (two referees)\\
On $\M_U$ every (CP-a) world adds up; mixed (CP-a)/adding-up systems excluded & Cor.~\ref{cor:CPaU}(b) & proved in a referee report (CPaU referee, Route B), one reader; at the status of Thm.~\ref{thm:M0closed}(b)\\
No face model positive on a middle arc & Rem.~\ref{rem:midarc} & proved in a referee report (M0 referee); confirmed (M0 second referee)\\
One-sided corner inequalities plus box bounds are infeasible; mirror box bounds & Lem.~\ref{lem:P1box} & proved (P1); confirmed (M0 referee); lower bounds upgraded to proved there\\
Split-facet configuration (P1) & \S\ref{sec:M0-closed} & algebra consistent (NUMERICS); configuration impossible\\
\textbf{Two-world ranking resolving an exceptional additive tie} & Thm.~\ref{thm:N26.1} & \textbf{proved (N26); confirmed by one referee (N26 referee)}\\
Two-level lexicographic class $\Ltwo$: no 2-dim.\ faces, rigid chords, $p^*\notin E$ & Thm.~\ref{thm:Ltwo} & proved (N26); confirmed with repairs (N26 referee); does not transfer to general M0 rankings\\
$\Ltwo$: the pivot theorem alone removes the chords & Prop.~\ref{prop:R1N26} & proved in a referee report (N26 referee), one reader; superseded\\
No $\Ltwo$ ranking covers $T_1$ on $\M_U$ (Conj.~N26.R superseded) & Cor.~\ref{cor:M0closed}(c) & proved, as Thm.~\ref{thm:M0closed}(b)\\
Strict-preference variants conflict with (S) & Prop.~\ref{prop:A} & proved in a referee report (N13); one reader\\
Exact extension theorem with the limit rule (R5) & Thm.~\ref{thm:R5} & proved (P3); confirmed both directions (P2P3 referee)\\
Extension lemma as posed in Vacuous Targets \S5 is false; unions of extendable stages need not extend & Prop.~\ref{prop:extcounter} & proved (P3); confirmed, with the referee's narrowing (P2P3 referee)\\
Witnessed extension theorem & Thm.~\ref{thm:witness} & proved (P3); confirmed, close to tautological (P2P3 referee)\\
Villegas-type rigidity under (R5) & Conj.~\ref{conj:villegas} & conjecture (P3); referee remarks incorporated\\
Vector form of the test; flexible worlds; rigid acceptances; nested criterion & Lems.~\ref{lem:vectorform}, \ref{lem:nested}, Props.~\ref{prop:flexible}, \ref{prop:rigid} & proved (P2), parts CHECKED; confirmed (P2P3 referee)\\
Two-level layering is locally single-level additive & Thm.~\ref{thm:LX} & proved (P2); confirmed with class repair (P2P3 referee); two levels only\\
Any cover's selection fails continuity or straight charts & Prop.~\ref{prop:P23} & proved in a referee report (P2P3 referee), one reader; localisation to one facet arc is a GAP\\
One kink makes the facet-map system consistent & \S\ref{sec:open-ordinal}(b) & NUMERICS (P2, reproduced); kinks excluded by coherence\\
Two worlds, nested: strict escape & Prop.~\ref{prop:strictescape} & proved (P4; $b=0$ case P4B referee); confirmed\\
Two worlds: leading charges; lexicographic worlds with charge levels (no (Nest) needed, round 13) & Thm.~\ref{thm:E2}, Cor.~\ref{cor:S21}(a) & under (Nest) proved in a referee report (P4B referee), one reader; without (Nest) via Thm.~\ref{thm:S21}, two readers (S2, S2 referee)\\
Two worlds, nested: complete worlds add up & Thm.~\ref{thm:E3} & CHECKED via Savage's theorem (P4B referee), one reader\\
Closed but non-SP purely finitely additive world, excluded by strict escape & Prop.~\ref{prop:B2} & proved (P4); confirmed with repair (P4B referee); chain-free form Cor.~\ref{cor:S21}(b) (S2 referee repair)\\
Locality barrier: no exclusion from indices in a compact subset of $\Di$ & Prop.~\ref{prop:Q31} & proved (Q3); confirmed, reading repaired (Q3 referee)\\
Shadow dichotomy; two-valued shadows; N17 Conj.~5.2 false in its index reading & Prop.~\ref{prop:Q32} & proved (Q3); confirmed with repair (``at least one'', integrable densities) (Q3 referee)\\
Non-additive behaviour at $T_1$: unendorsable out-comparisons, or unendorsed comparisons against incoherent models within $\eta$ of $T_1$ & Thm.~\ref{thm:R3} & proved in a referee report (Q3 referee), one reader; replaces Q3's ``not forced'' claim (WRONG)\\
Far dominators force a far tie partner or $\neg$(LC$_{T_1}$); tie toggles & Prop.~\ref{prop:Q38}, Lem.~\ref{lem:Q39} & proved (Q3); confirmed with wording repairs (Q3 referee)\\
Edge complementarity & Lem.~\ref{lem:Q34} & proved (Q3, Q1); confirmed (Q3 referee)\\
Edge reduction under (Nest$_e$) & Prop.~\ref{prop:Q35} & proved (Q3; Q1 CHECKED); confirmed (Q3 referee); E3 step CHECKED\\
Three-world transfer (any $n$): leading charges and lexicographic worlds without (Nest$_e$); complete worlds, strict escape and Lemma~T under (Nest$_e$) & Thm.~\ref{thm:Q36} & (a), (b) proved without (Nest$_e$) (S2; confirmed, S2 referee); under (Nest$_e$) proved earlier (Q3; Q1 CHECKED; confirmed, Q3 referee); (c) CHECKED (Savage); (d) proved; (a$'$) proved (B1 Thm.~B1.3; confirmed, B1 referee)\\
\textbf{Leading charges are 0 on $M$ or strictly positive, with no nestedness (any $n\ge2$)}; E2 without (Nest), B2 without chain coordinates, LX's class & Thm.~\ref{thm:S21}, Cor.~\ref{cor:S21} & \textbf{proved (S2); confirmed (S2 referee), two readers}; Cor.\ (b) with repair; (c) proved in a referee report, one reader\\
Parallelogram lemma & Lem.~\ref{lem:PL} & proved (S2); confirmed with repair (S2 referee)\\
\textbf{No cover of $T_1$ with every world (CP-a), or with a leading charge nonzero on $M$, or complete and edge-nested} & Cor.~\ref{cor:Q37} & \textbf{proved} at the status of Thm.~\ref{thm:M0closed}(b); CHECKED with a complete world; confirmed (Q3 referee; leading-charge clause without (Nest$_e$), S2 referee)\\
(Nest$_e$) on the edge models $\co(K'\cup[x,y])$ is comonotonicity & Prop.~\ref{prop:comonotone} & proved in a referee report (Q3 referee), one reader\\
Additive ties: Bayes error sets nested on an edge; null ties in two worlds; (Nest$_e$) from additive edge worlds and a complete third world & Props.~\ref{prop:S20}, \ref{prop:RS2A} & S20 proved (S2), confirmed, (a) with wording repair (S2 referee); R-S2.A proved in a referee report (S2 referee), one reader\\
Edge families: straightening suffices; the edge inputs do not force (Nest$_e$) & Rem.~\ref{rem:straighten}, Prop.~\ref{prop:plateau} & proved (S1); confirmed, the remark with repair (S1 referee)\\
Two worlds, nested: cut form; Lemma T$_2$; left-fineness & Lem.~\ref{lem:cutform}, Thm.~\ref{thm:LF} & proved (Q1, Q2); confirmed (Q12 referee)\\
Two worlds, nested: uncharged sets are initial segments & Cor.~\ref{cor:IS} & proved (Q2); confirmed (Q12 referee)\\
Two worlds, nested: mirror pairs; MP-loop; MP$'$; no additive scaffold & Prop.~\ref{prop:MP}, Cors.~\ref{cor:MPloop}, \ref{cor:MPprime}, Prop.~\ref{prop:NS} & proved (Q2; MP = Q1.4 repaired); confirmed, NS with repair (Q12 referee); MP$'$ complete clause CHECKED\\
Spread and star certificates impossible next to an SP world & Prop.~\ref{prop:Q15} & proved (Q1); confirmed (Q12 referee)\\
Simple certificates = Ces\`aro-thin families; mirrored pairs give no contradiction; tilt world; ratio certificates & Props.~\ref{prop:N1}, \ref{prop:N2}, \S\ref{sec:open-core}(c), (e) & proved in a referee report (Q12 referee), one reader\\
Two worlds without nestedness: lattice of error sets, 1-D KKM, G-SE & Prop.~\ref{prop:lattice}, Lem.~\ref{lem:GSE} & proved (Q2); confirmed (Q12 referee)\\
Two worlds: a nonzero leading charge at either world forces a chain & Thm.~\ref{thm:S22} & proved (S2); confirmed (S2 referee)\\
Two worlds: Level Lemma; $\C_L$ modulo double ties is a linear continuum; strict escape without nestedness; nested representatives & Lem.~\ref{lem:level}, Cor.~\ref{cor:continuum}, Prop.~\ref{prop:S13} & proved (S1); confirmed (S1 referee)\\
Crossings in $\C_L$ are rigid; double ties only at the corner pairs & Lem.~\ref{lem:rigid} & proved (S1; CHECKED by script); confirmed, strengthened (S1 referee)\\
Either half of (MC$_{\mathrm{sel}}$), or monotone continuity at either world, forces a chain; a non-nested system needs both worlds without (MC) and without a leading charge & Thm.~\ref{thm:MC}, Cor.~\ref{cor:MC} & proved (S1); confirmed with repair (S1 referee); combined statement from S1 and S2 referees\\
Full two-world ranking with a crossing selection in which only (C) fails & Prop.~\ref{prop:R71} & proved in a referee report (S1 referee), one reader; NUMERICS\\
(S) forces a nested selection in two worlds & Conj.~\ref{conj:Q2N} & conjecture (Q2); proved when some world has a nonzero leading charge or is monotonely continuous (Thms.~\ref{thm:S22}, \ref{thm:MC}); open only when both worlds lack both\\
(S) forces one half of (MC$_{\mathrm{sel}}$) & Conj.~\ref{conj:S1R} & conjecture (S1; reduced by S1 referee)\\
Edges: weak edge undominatedness; under (EU$_{ij}$), Level Lemma, continuum and patterns; pairwise straightening & Lem.~\ref{lem:weakedge}, Props.~\ref{prop:S111}, \ref{prop:S112} & proved (S1), the propositions conditional on (EU$_{ij}$) (and monotone continuity); confirmed (S1 referee)\\
Global edge straightening & Conj.~\ref{conj:S113} & conjecture (S1)\\
Every axiom except closed sections at one world, with a world that does not add up & Props.~\ref{prop:Q12}, \ref{prop:Q13} & proved (Q1); confirmed, Q1.2 with repair (Q12 referee); NUMERICS for the exact section\\
Mirrored simple pair with every axiom except closed sections at both worlds & Prop.~\ref{prop:S31} & proved (S3); confirmed (S34 referee)\\
Sections of band-and-translates worlds (closedness forces $\inf\bar\mu(F)>0$; necessary, not sufficient); packing lemma; the fibrewise route fails (K) & Props.~\ref{prop:S32}, \ref{prop:S33}, \ref{prop:S34} & proved (S3); S3.2 confirmed with repair, S3.3, S3.4 confirmed (S34 referee)\\
Rates: (R3) at either world $\iff$ uniform strict escape $\iff$ no doubly-light sequence; true next to an end-vanishing SP weight & Thm.~\ref{thm:RATE}, Cor.~\ref{cor:SR} & proved (S4); confirmed (S34 referee)\\
No doubly-light sequences (restated (R3)) & Conj.~\ref{conj:NDL} & conjecture (S4)\\
Exact threshold of MP; loop lemma; blocking effect; uniform escape & Prop.~\ref{prop:MPsharp}, Lems.~\ref{lem:loop}, \ref{lem:block}, \ref{lem:UE} & proved (S4); MP$\sharp$(d) and the loop lemma confirmed with repair; MP$\sharp$(e) and the blocking lemma proved in a referee report (S34 referee), one reader\\
(QC): four equivalent forms; (QC) excludes mirrored pairs & Prop.~\ref{prop:QCJ}, Conj.~\ref{conj:QC} & implication proved (S3), confirmed; equivalences proved in a referee report (S34 referee), one reader; (QC) conjecture\\
Every axiom, both worlds additive, Lemma-T radius $\le\bar\mu(P)^2/2$; a linear rate would exclude ratio certificates & Props.~\ref{prop:ExA}, \ref{prop:LIN} & proved (S4); confirmed, LIN with repair (S34 referee)\\
\multicolumn{3}{l}{\emph{Rounds 14--15} (\S\ref{sec:open-core}(f); round-15 results refereed once, C12 referee)}\\
Sign classes; leading edge densities glue through the link class; transfer chain; they propagate, so the world adds up & Lems.~\ref{lem:B10}--\ref{lem:B12}, Thm.~\ref{thm:B13}, Thm.~\ref{thm:Q36}(a$'$) & proved (B1); confirmed (B1 referee), one reader\\
Edge comparisons alone do not determine a world & Prop.~\ref{prop:B1cex} & proved in a referee report (B1 referee), one reader\\
\textbf{Every world of a cover that does not add up lacks a leading density on some edge region} & Cor.~\ref{cor:B14} & \textbf{proved} at the status of Thm.~\ref{thm:M0closed}(b); confirmed (B1 referee); nonnegative densities CHECKED; countably additive only\\
Additive base with interior changes; aggregates of additive scores & Cor.~\ref{cor:B15} & proved (B1), with the referee's corrections; (b) the referee's ordinal form\\
Total worlds: leading charge $\iff$ almost-agreeing charge; AR worlds add up with no nestedness & Lem.~\ref{lem:B31}, Cor.~\ref{cor:B32} & proved (B3); confirmed, wording repair (B3 referee), one reader\\
Order continuity on edges & Lem.~\ref{lem:B33} & proved (B3); confirmed (B3 referee)\\
Ultrafilter edge example: crossing and not fine, but AR & Prop.~\ref{prop:B34} & proved (B3; ultrafilter lemma); confirmed with repairs (B3 referee); two-world realisation without (C) proved in that report\\
Leading top levels with dense null sets are trivial; signature (IRWAS) levels; exchangeable signatures are additive & Thm.~\ref{thm:B35}, Cor.~\ref{cor:B36}, Prop.~\ref{prop:B37} & proved (B3); confirmed (B3 referee); extension to leading top levels of any world in that report, one reader\\
Total worlds are AR; class-chain form for $n=2$ & Conj.~\ref{conj:B3C}, Prop.~\ref{prop:B3R3}, Conj.~\ref{conj:B3Cp} & conjectures; Prop.~\ref{prop:B3R3} proved in a referee report (B3 referee), one reader\\
Glue dichotomy: collapse or unanimity over the glued credal set & Thms.~\ref{thm:C11}, \ref{thm:C14} & proved (C1); confirmed (C12 referee), one reader; NUMERICS\\
No cover with every world glued credal, (CP-a), leading-charge, or total and AR & Thm.~\ref{thm:C16} & proved (C1); confirmed at the status of Thm.~\ref{thm:M0closed}(b) (C12 referee), with a scope repair\\
Credal band; total band worlds add up; finite obstruction & Thm.~\ref{thm:C17}, Cor.~\ref{cor:C18}, Prop.~\ref{prop:C19} & proved (C1); confirmed (C12 referee), one reader; Cor.\ CHECKED (Savage), with repair; Question C1.Q open\\
Full two-world system with a world that has no leading density or charge (it adds up) & Thm.~\ref{thm:C21} & proved (C2); confirmed with small repairs (C12 referee), one reader; NUMERICS\\
Edge comonotonicity of density unanimities; edge-only arguments cannot prove (LD) & Thm.~\ref{thm:C23}, Cor.~\ref{cor:C25} & proved (C2); confirmed (C12 referee), Cor.\ with its last sentence narrowed\\
Link rigidity; linear and box neighbours & Prop.~\ref{prop:C26}, Cor.~\ref{cor:C27} & proved (C2); confirmed (C12 referee), one reader; Conj.~\ref{conj:C2C} conjecture\\
\textbf{Face models at a target weighting on every decomposable domain: crossing-pinned, contain $T$, no apex ties} & Thm.~\ref{thm:B21} & \textbf{proved (B2); confirmed with repairs (B2 referee), one reader}; every all-facing triangle proved there\\
Rankings that add up on general domains: no apex ties at $p^*$; ties at $p^*$ are full corner atoms & Cor.~\ref{cor:B22}, Prop.~\ref{prop:B213} & conditional on Claim B2.0 (unwritten); Prop.~\ref{prop:B213} proved (B2), confirmed with repairs\\
Full-atom residual realised without (LF) & Prop.~\ref{prop:B2ex} & proved in a referee report (B2 referee), one reader\\
Route I1 against an exceptional set away from $p^*$ & \S\ref{sec:open-core}(f) & CHECKED (B2), corrected by the referee: conditional on unaudited inputs\\
\multicolumn{3}{l}{\emph{Round 16} (\S\ref{sec:open-core}(e), (f))}\\
\textbf{Full two-world system on $\M_U$ in which a world does not add up (type (M2)), next to an additive world}; designs; sections & Thm.~\ref{thm:D21}, Lems.~\ref{lem:D22}--\ref{lem:D24} & \textbf{proved (D2); confirmed by two independent referees (D2 referee, D2 referee 2)}, with small repairs\\
N17.B$'$ and N17.B$'$(2) false; Q1.C1 false; (R2) answered negatively & Thm.~\ref{thm:D21} & proved, from Thm.~\ref{thm:D21} (both D2 referees)\\
That world is endorsed by a $\sigma$-finite density non-integrable only at the protected end; that end contains the ray $\R_+(1,0)$, interior to $U$ & Prop.~\ref{prop:D2sigma}, Lem.~\ref{lem:F11} & proved in a referee report (D2 referee 2), one reader; the geometry of the end proved (F1), confirmed (F12 referee); NUMERICS\\
\multicolumn{3}{l}{\emph{Audit of 15 October} (\S\ref{sec:open-core}(e), (f); notes F1, F2, refereed once in F12 referee)}\\
\textbf{N17.B$'$ with densities locally integrable on $U$ is false ($n=2$)}; next to an additive world, $\sigma$-finite endorsers are nondecreasing in the chain coordinate & Conj.~\ref{conj:sigmaB}, Thms.~\ref{thm:F12}, \ref{thm:F13} & \textbf{refuted}, Thm.~\ref{thm:D21} being a counterexample: observed by an external auditor; proved (F1; a second route in F2); confirmed (F12 referee), one reader\\
Repaired form: densities locally integrable on $M$ & Conj.~\ref{conj:F1Cp} & conjecture (F1); its use for $T_1$ depends on (Nest$_e$) (F12 referee)\\
No certificate with all targets in one compact subset of $U$ with payoffs bounded away from $0$, next to an additive world & Thm.~\ref{thm:F14}, Cor.~\ref{cor:F15} & proved (F1); confirmed (F12 referee), one reader; last sentence of Cor.~\ref{cor:F15} a necessary condition only\\
Full two-world systems whose generators and protected set lie where both payoffs are bounded away from $0$, with protection at $\partial U$; where their endorsers blow up & Thms.~\ref{thm:F16}, \ref{thm:F18} & proved (F1); confirmed with a wording repair (F12 referee), one reader; NUMERICS\\
(LF) at $D_{T_1}$: no non-null ties at the target weighting; the full-atom residual is empty & Thm.~\ref{thm:D1} & proved (D1); confirmed with repairs (D1 referee), one reader\\
Three-atom field with $T_1$ interior to the face, without (LF) & Prop.~\ref{prop:D10} & proved (D1); confirmed (D1 referee); NUMERICS\\
Rankings that add up under (LF): $T_1$ never among tied Bayes models & Cor.~\ref{cor:D17} & at the status of Claim B2.0, with a reduced list of imports (D1 referee)\\
What Thm.~\ref{thm:M0closed} uses: dominator comparisons of one admissible shadow; Thm.~\ref{thm:R3} without its case split & Thms.~\ref{thm:E11}, \ref{thm:E12} & proved in a note (E1); audited, with the hypotheses written out (F2); audit confirmed (F12 referee), one reader; NUMERICS\\
Ends of an end-protected world are eventually $T_1$-correct; where ends can sit; end-heavy shadows; covers by end-protected decorations with integrable, jointly end-heavy base densities excluded; the failure in Thm.~\ref{thm:D21} is invisible in its two-world system & Lems.~\ref{lem:E13}, \ref{lem:E14}, Prop.~\ref{prop:F21}, Thm.~\ref{thm:E15}, Prop.~\ref{prop:E16} & proved in notes (E1; Prop.~\ref{prop:F21} in F2); audited (F2); audit confirmed (F12 referee), one reader; NUMERICS; heaviness in Thm.~\ref{thm:E15} removed in round 17 (Thm.~\ref{thm:H13}, Prop.~\ref{prop:HR1})\\
\multicolumn{3}{l}{\emph{Round 17} (\S\ref{sec:open-core}(f$'$), (g); notes H1, H2, J1; referees H12, J1 (round 19))}\\
Stability: dominator comparisons into any regular model within $\Phi^*_t$ of the Bayes cost are not all endorsed; admissibility follows from $T_1$'s localized out-comparisons; base endorsing families with no admissible tuple are impossible & Lems.~\ref{lem:H10}, \ref{lem:H11} & proved (H1); confirmed (H12 referee), one reader; Lem.~\ref{lem:H10}(b) with the referee's added bound; NUMERICS\\
\textbf{One world with an a.e.\ finite endorser of its localized comparisons, next to two worlds with integrable full endorsers: no cover} & Thm.~\ref{thm:H12} & \textbf{proved (H1); confirmed (H12 referee), one reader}; integrability-free use of Thm.~\ref{thm:M0closed}(a) justified by the referee\\
Every world with an integrable endorser of its $T_1$-localized comparisons: no cover; (h) of Thm.~\ref{thm:E15} replaced by $T_1$-correctness of the ends & Thm.~\ref{thm:H13} & proved (H1); confirmed (H12 referee), one reader; the note's ``E1.5 holds without (h)'' narrowed by the referee\\
Refined ends; a $T_1$-wrong end forces a partner world that fails to add up, and is shadowed; joint reweighting with (h$'$) & Lems.~\ref{lem:H20}, \ref{lem:H21}, Prop.~\ref{prop:H22} & proved (H2); confirmed (H12 referee), one reader; Lem.~\ref{lem:H21}(b) with the cancellation reading (referee's repair); Prop.~\ref{prop:H22} new only in (h$'$)\\
\textbf{Finitely many ends, integrable base endorsers: no cover unless the conflict digraph has a directed cycle}; no heaviness needed & Thm.~\ref{thm:H23}, Prop.~\ref{prop:HR1} & \textbf{Thm.~\ref{thm:H23} proved (H2); confirmed (H12 referee), one reader}; Prop.~\ref{prop:HR1} proved in a referee report (H12 referee), one reader, and independently in J1 (refereed once in round 19)\\
Such a cover needs two worlds with $T_1$-wrong ends and two without an integrable full endorser & Cor.~\ref{cor:H23p} & proved (H2), corrected (H12 referee); ``two worlds fail to add up'' only when the sets $V(W)$ are singletons\\
\textbf{Full two-world system in which neither world adds up} & Thm.~\ref{thm:H24} & \textbf{proved (H2); every axiom re-derived (H12 referee), one reader}; wording repair at $b=1$; NUMERICS\\
Near-Bayes models; closure lemma; pattern switching: no cover if some order of the worlds is aligned; the template of H2 and other configurations closed; shell exchanges in the surviving case & Lems.~\ref{lem:J1C}, \ref{lem:J10}, Thm.~\ref{thm:J12}, Cor.~\ref{cor:J13}, Prop.~\ref{prop:J15} & proved (J1); confirmed (J1 referee, round 19), one reader; Cor.~\ref{cor:J13}(a) strengthened and two one-line gaps closed by the referee; Example J1.4 CHECKED, with the referee's caveat on the radii\\
\multicolumn{3}{l}{\emph{Round 19} (\S\ref{sec:open-core}(f$'$)--(g); notes K1--K3, each refereed once; J1 referee)}\\
\textbf{Order alignment from arrow alignment; the residual of class (EP) is (R-i)--(R-iii)} & Thm.~\ref{thm:JR2}, Cor.~\ref{cor:JR3} & \textbf{proved in a referee report (J1 referee), one reader}; (R-ii) shells CHECKED; NUMERICS\\
Doubly misaligned $2$-cycles with an isolated third world; no aligned order in class (EP); in class (EP$_\infty$) & Questions J1.Q, J1.Q$'$, J1.Q$'_\infty$ (\S\ref{sec:open-core}(f$'$), (f$''$)) & open; J1.Q reworded by the J1 referee; a negative answer to J1.Q closes only (R-i)\\
\textbf{Classification: a world that does not add up has a certificate, or a non-null non-localizable blind set $\Lambda_\infty$, or is (Loc$_{\mathrm{int}}$)}; end-protection adds nothing without a base density & Thm.~\ref{thm:K31}, Prop.~\ref{prop:K32} & \textbf{proved (K3); confirmed (K3 referee), one reader}; (NC) made explicit; $\hat\Theta$, $\rho_\infty$ not canonical\\
The tilt world: every axiom except closed sections at one world, which is in (N) with $\Lambda_\infty=M$ & Thm.~\ref{thm:K33} & proved (K3), completing Q12 referee N3; confirmed (K3 referee), one reader; NUMERICS\\
Pairwise-meeting protection is localizable & Lem.~\ref{lem:K34} & proved (K3); confirmed (K3 referee), one reader; readings drawn from it withdrawn\\
\textbf{Patch theorem}: Thms.~\ref{thm:H12}, \ref{thm:H13} for $\sigma$-finite densities at several worlds; residual: sheltering or infinite minority mass; blow-ups sheltered or shadowed & Thm.~\ref{thm:K35}, Cor.~\ref{cor:K35p}, Lem.~\ref{lem:K36} & \textbf{proved (K3); confirmed (K3 referee), one reader}; the minority-mass case excluded if Q$_{\mathrm{K3}}$ has answer No (referee)\\
One world without a certificate next to two integrable worlds: property ($\gamma$) of the weak field, rigid blind set, admissibility & Lem.~\ref{lem:K37p}, Thm.~\ref{thm:K37} & proved (K3); confirmed (K3 referee), one reader; NUMERICS\\
Necessary conditions for Question Q$_{\mathrm A}$; where Thm.~\ref{thm:M0closed}(a) breaks & Prop.~\ref{prop:KQA} & proved in a referee report (K3 referee), one reader; the location of the break a reading\\
Next to an additive world in two worlds, every blind set is an initial segment of the chain & Prop.~\ref{prop:K38}, Conj.~\ref{conj:K3C} & proposition proved (K3), confirmed (K3 referee); conjecture open, its evidence weaker than stated\\
Non-localizable blind sets in full systems; the additive question with one weakly signed component & Questions Q$_{\mathrm{K3}}$, Q$_{\mathrm A}$ & open\\
\textbf{No certificate consists of divisible comparisons; forced comparisons are divisible} & Thm.~\ref{thm:K21}, Prop.~\ref{prop:K22}, Cor.~\ref{cor:K22p} & \textbf{proved (K2); confirmed (K2 referee), one reader}; CHECKED (script)\\
Two worlds: a certificate refutes Conj.~\ref{conj:Q2N}; frontier models (crossing needs interleaving); edge sweep & Thms.~\ref{thm:K23}, \ref{thm:K24}, Lem.~\ref{lem:K27} & proved (K2); confirmed (K2 referee), one reader; NUMERICS\\
Decoration scheme: every axiom from finite data; empty under (SPREAD) or (ADJ) at both edges; Brier satisfies (SPREAD) & Thms.~\ref{thm:K25}, \ref{thm:K26}, Prop.~\ref{prop:K26a} & proved (K2); confirmed (K2 referee), one reader; (ADJ) for Brier NUMERICS only\\
\textbf{Over the Brier base every endorser of the edge trades is $c\rho_1$ on $M$: the whole decoration template is empty} & Prop.~\ref{prop:KR62} & \textbf{proved in a referee report (K2 referee), one reader}\\
Bayes-shaped selections force adding up; a full system with a certificate next to two additive worlds & Conj.~\ref{conj:K26p}, Question Q-K2 & conjecture (K2); open\\
Countably many ends: safe closure, shadowing, per-world safe tests, conflict-free depths & Lems.~\ref{lem:K10}, \ref{lem:K11}, Thm.~\ref{thm:K12}, Cor.~\ref{cor:K13} & proved (K1); confirmed (K1 referee), one reader\\
\textbf{Countably many ends, integrable base endorsers: no cover if the conflict digraph is acyclic or some order is aligned} & Thms.~\ref{thm:K14}, \ref{thm:K15} & \textbf{proved (K1); confirmed (K1 referee), one reader}; sufficient condition with the referee's uniform reading; persistent conflict without an arrow CHECKED (K1 referee)\\
Anti-safe selection indices: necessary condition and exclusion test & Thm.~\ref{thm:K16} & proved (K1); confirmed (K1 referee), one reader; the note's ``iff'' with K1.Q is tautological\\
N17.B$'$: two worlds outside the classes above; three worlds & \S\ref{sec:open-ordinal}(c), \S\ref{sec:open-core}(g) & false in two worlds (Thm.~\ref{thm:D21}, round 16), so (R2) is answered negatively; its form with densities locally integrable on $U$ (Conj.~\ref{conj:sigmaB}) is false too (Thm.~\ref{thm:F13}), and Conj.~\ref{conj:F1Cp} is open; in three worlds no counterexample is known; for covers of $T_1$ the remaining items are (R1), (R4), Conj.~\ref{conj:NDL}, Conj.~\ref{conj:F1Cp}, Conj.~\ref{conj:B3C}, Question C1.Q, Conj.~\ref{conj:C2C}, other edge orders without a leading density, the interior step, the residual of \S\ref{sec:open-core}(g) (round 19: Questions Q-K2, Q$_{\mathrm{K3}}$, Q$_{\mathrm A}$, Conjs.~\ref{conj:K3C}, \ref{conj:K26p}, and sheltering, with Questions J1.Q, J1.Q$'$, J1.Q$'_\infty$ for end-protected worlds), and case (i) of Thm.~\ref{thm:R3}\\

Walley-2000 variant with invisible desirability & \S\ref{sec:cp-b} & open\\
Forced small-stakes ties: no additive cover & \S\ref{sec:open-ties} & proved; refereed once (N4)\\
Permitted small-stakes ties: Thm.~\ref{thm:main} survives below $\theta_x(T)$ & Thm.~\ref{thm:permitted} & proved (N27); confirmed by one referee (N27 referee)\\
Permitted ties, non-adjacent (LF) domains & Cor.~\ref{cor:permitted-nonadj} & CHECKED (N27); confirmed at CHECKED level\\
Euclidean band $\theta\in[\theta_x(T_1),1)$ & \S\ref{sec:open-ties} & open\\
Other triangle classes, $n\ge4$, segments, boundary targets & \S\ref{sec:open-other} & open\\
\bottomrule
\end{longtable}}

\section{Setting}\label{sec:setting}

\subsection{Problems, models and advice}
Let $\Omega\coloneqq\{1,\dots,n\}$ be the worlds, $n\ge2$; from Section~\ref{sec:reduction} on, $n=3$. A gamble is $f\in\R^n$; it pays $f_i$ at world $i$. Write $\Delta$ for the probability simplex, $e_i$ for its vertices, $\one=(1,\dots,1)$, $\supp p\coloneqq\{i:p_i>0\}$, and $\Di$ for the interior of $\Delta$ relative to $\aff\Delta$. Fix the standard Gaussian measure $\mu$ on $\R^n$; it has the same null sets as Lebesgue measure.

\paragraph{What a gamble represents.} Any gamble $f$ represents the binary choice $\{f,0\}$, which is equivalent to the binary choice $\{g,h\}$ for any gambles $g,h$ with $f=g-h$. That is all that $f$ represents. The gambles $f$ and $-f$ represent the same binary choice, with the roles of the two options swapped; we call this choice the \emph{pair} $\{f,-f\}$. A \emph{model} is a Lebesgue-measurable set $D\subseteq\R^n$ of desirable (accepted) gambles, taken modulo $\mu$-null sets. On the choice represented by $f$ it gives one of four kinds of advice:
\begin{itemize}[leftmargin=*]
\item $f\in D$ and $-f\notin D$: you prefer $f=g-h$ to $0$ and do not prefer $-f=h-g$ to $0$, which is true iff you prefer $g$ to $h$ and do not prefer $h$ to $g$; so $D$ advises choosing $g$ in $\{g,h\}$, for every $g,h$ with $f=g-h$;
\item $f\notin D$ and $-f\in D$: $D$ advises choosing $h$ in $\{g,h\}$, for every $g,h$ with $f=g-h$;
\item $f\notin D$ and $-f\notin D$: no advice (also called \emph{silence});
\item $f\in D$ and $-f\in D$: contradictory advice.
\end{itemize}
As bookkeeping only, each pair is indexed once by its member in $H\coloneqq\{f:f_1>0\}$, and for $f\in H$ we write the advice $a_D(f)$ as $+$ (accept $f$ only), $-$ (accept $-f$ only), $0$ (no advice) or $\times$ (both); nothing depends on this convention. At world $i$, accepting $f$ (that is, choosing $g$ in $\{g,h\}$) is \emph{good} if $f_i>0$ and \emph{bad} if $f_i<0$, and the roles swap for $-f$. Coordinate hyperplanes are null and are ignored. The \emph{background order} at world $i$, better first, is
\[
\text{good}\;\prec\;\text{no advice}\;\prec\;\text{bad}\;\prec\;\text{contradictory}.
\]
From here on the letters $g,h,x,\dots$ denote gambles, each representing its binary choice in this way; the options of the choice no longer appear.
A gamble is \emph{mixed} if it has a positive and a negative coordinate.

For a credal set $K$ (a nonempty closed convex subset of $\Delta$) put
\[
P_K(g)\coloneqq\min_{m\in K}m\cdot g,\qquad D_K\coloneqq\{g:P_K(g)>0\}\cup\{g\ge0:g\ne0\}.
\]
Every coherent set of desirable gambles agrees a.e.\ with some $D_K$ (\src{Full Coverage by Homothetic Credal Families v3}, App.~A), and $D_K\approx D_{K'}$ modulo null sets only if $K=K'$. A model is \emph{coherent} if it agrees a.e.\ with some $D_K$.

\subsection{Decomposable domains, rankings and axioms}
The domain of a ranking is the class of models that it evaluates. Restricting the domain means only restricting which models are evaluated. Decision problems are never altered or removed, and they are never given restricted options: every gamble $f$ still represents the binary choice $\{f,0\}$, equivalently $\{g,h\}$ with $f=g-h$, and the world orders of a ranking still concern every binary problem. This is the only way to make sense of a domain restriction in the accuracy-ordering framework.

\begin{definition}\label{def:domain}
A \emph{decomposable domain} is a nonempty class $\M$ of models (the \emph{evaluable models}) that is closed under splicing, i.e.\ $(X\cap A)\cup(Y\setminus A)\in\M$ whenever $X,Y\in\M$ and $A=-A$ is measurable, and under convergence in measure, i.e.\ $X\in\M$ whenever $X_m\in\M$ and $\mu(X_m\sd X)\to0$.
\end{definition}

\begin{remark}\label{rem:advicerange}
The decomposable domains are exactly the classes of the form
\[
\M=\{D\text{ measurable}:\ a_D(f)\in A(f)\text{ for a.e.\ }f\in H\},
\]
where $f\mapsto A(f)\subseteq\{+,0,-,\times\}$ is measurable (each set $\{f:a\in A(f)\}$ is measurable) and nonempty for a.e.\ $f$; the map $A$ is determined by $\M$ up to a null set. We call $A(f)$ the \emph{advice range} of the domain at $f$: the kinds of advice that evaluable models give on the binary choice represented by $f$. So ``the advice range at $f$ is $\{+,0\}$'' says that every evaluable model either accepts $f$ (and not $-f$) or gives no advice on that choice, and that evaluable models of both kinds exist; and when the advice range at $f$ is a single kind of advice, every evaluable model gives that advice on the choice represented by $f$. That choice is still a decision problem; it simply does not distinguish evaluable models.
\end{remark}
\status{proved (\src{Domain Closure and Extendability}, Prop.~2, author-checked): take a countable dense family $(X_k)$ in $\M$ and put $A(f)\coloneqq\{a_{X_k}(f):k\ge1\}$; every evaluable model has its advice in $A$ a.e., and every measurable selection from $A$ is a limit in measure of finite splices of the $X_k$. Splicing alone is not enough (same source, first part of Prop.~2). This is a special case of the theorem of Hiai and Umegaki on closed decomposable subsets of $L^1$~\cite{HU1977}.}

\paragraph{Earlier terminology.} Versions of this note and of the working notes written before 14 October 2026 called decomposable domains ``menu domains'', called the advice range a ``menu'' of available actions, called the $U$-domain of \S\ref{sec:aos} the ``accept-or-silence domain'', and called problems on which every evaluable model gives no advice ``forced silent''. That language misdescribed a restriction of the class of evaluable models as a restriction of the decision problems, and it has been removed from all current files; files and folders whose names used it were renamed on the same day.

A \emph{ranking} $L$ consists of a decomposable domain $\M_L$, its class of evaluable models, and a preorder $\preceq_i$ on $\M_L$ at each world $i$ (smaller is better). The world orders compare evaluable models by the advice they give on all binary problems. Write $X\prec_iY$ for the strict part and $X\sim_iY$ for ties. $Y$ \emph{dominates} $X$ if $Y\preceq_iX$ at every world and $Y\prec_jX$ at some world. $U_L$ is the set of undominated members of $\M_L$, and $\C_L$ the set of its coherent members. A model $C$ is \emph{competitive on} $S\subseteq\Omega$ if every $X\in\M_L$ has some $i\in S$ with $C\preceq_iX$. The following requirements are imposed within $\M_L$ (\src{Full Coverage \dots v3}, \S1; \src{Loss Dominance Generalized v2}, \S\S1--2).
\begin{enumerate}[leftmargin=*,label=\textbf{(\arabic*)}]
\item[\textbf{(G)}] \emph{Guidance.} If $X$ gives advice at least as good as $Y$ at world $i$ for a.e.\ $g$, then $X\preceq_iY$; if moreover it is strictly better on a non-null set, then $X\prec_iY$.
\item[\textbf{(IB)}] \emph{Independence across binary problems.} For measurable $A=-A$ and measurable $F,G\subseteq A^c$ (with all four models in $\M_L$): $(D\cap A)\cup F\preceq_i(E\cap A)\cup F$ iff $(D\cap A)\cup G\preceq_i(E\cap A)\cup G$.
\item[\textbf{(LA)}] $\C_L\ne\varnothing$.
\item[\textbf{(S)}] \emph{Closed competitive choice.} There is a map $p\mapsto C_p\in\C_L$ on $\Delta$ such that $C_p$ is competitive on $\supp p$ and $\{p:C_p\preceq_iX\}$ is closed for every $i$ and every $X\in\M_L$.
\item[\textbf{(C)}] A model tied at every world with a member of $\C_L$ belongs to $\C_L$.
\end{enumerate}
Restricting the domain leaves (G) unchanged: it compares the advice that two evaluable models give on each problem. For a class of rankings one also asks (E): every model is evaluated by some ranking (supplied by an all-model Brier ranking), and (A): every coherent model is in some $\C_L$. Condition (A) is \emph{coverage}.

\begin{theorem}\label{thm:kkm}
Let $\Omega$ be finite and let $L$ satisfy (LA) and (S). Then every $X\in\M_L$ is weakly dominated by an undominated selected model $C_p$, and every undominated model is tied at every world with a selected model.
\end{theorem}
\status{proved (\src{Loss Dominance Generalized v2}, Thm.~thm:kkm, via the KKM lemma; valid for preorders). Cited as ``Fact 0'' in the notes.}

\begin{definition}\label{def:LF}
A ranking $L$ \emph{covers} a coherent model $D$ if $D\in\C_L$ and (LF) holds at $D$: there is an open $N\supseteq\partial D\setminus\{0\}$ such that every measurable $R$ with $R\sd D\subseteq N$ a.e.\ belongs to $\M_L$. For a decomposable domain this means that for a.e.\ pair with an orientation $h\in N$ the advice range contains both ``accept $h$'' and no advice.
\end{definition}
This is the open-set form of (LF): $N$ need not contain a conic neighbourhood of $\partial D\setminus\{0\}$ (for instance $N=\{g:\operatorname{dist}(g,\partial D)<e^{-\|g\|}\}$ contains none; \src{N27} \S1). On the $U$-domain (\S\ref{sec:aos}) the explicit open cone $U$ serves as $N$. For additive rankings whose interior Bayes models are coherent, charts are invariant along almost every ray (Lemma~\ref{lem:raytransfer}), so the conic form holds at the level of charts; no strengthening of (LF) is needed anywhere in this note (\src{N27} Repair 27.1, confirmed in \src{N27\_referee}).

\subsection{The target}
Put $c_0\coloneqq(1-\varepsilon)/3$ and
\[
\kappa\coloneqq\frac{1+2\varepsilon}{1-\varepsilon}=1+\frac{3\varepsilon}{1-\varepsilon}>1 .
\]
The vertices of $T_1$ are $v_k=c_0\one+\varepsilon e_k$, and
\[
D_{T_1}=\Bigl\{g:\ \tfrac{1-\varepsilon}3(g_1+g_2+g_3)+\varepsilon g_k>0\ \text{for }k=1,2,3\Bigr\}\cup\{g\ge0:g\ne0\}.
\]
The facets of $D_{T}$ for a triangle $T=\co\{v_1,v_2,v_3\}$ are $F_k\coloneqq\{h:h\cdot v_k=0,\ h\cdot v_j\ge0\ (j\ne k)\}$, two-dimensional closed cones in $v_k^\perp$. Two facets $F_k,F_j$ meet in a \emph{corner ray} $r_{kj}\propto v_k\times v_j$. Coordinate $i$ \emph{crosses} on $F_k$ if $h_i=0$ at some $h\in\relint F_k$. For $T_1$ the corner rays are
\[
r_{12}=(-1,-1,\kappa+1),\quad r_{23}=(\kappa+1,-1,-1),\quad r_{31}=(-1,\kappa+1,-1),
\]
and on $F_1$ the coordinates $2,3$ cross at $w_{21}=(-1,0,\kappa)$ and $w_{31}=(-1,\kappa,0)$ (cyclically on $F_2,F_3$). These follow from $c_0(\kappa-1)=\varepsilon$.

\begin{definition}\label{def:allfacing}
A triangle $T\subseteq\Di$ is in \emph{general position} if no line through two of its vertices contains a vertex of $\Delta$. It is \emph{all-facing} if $v_k\in\operatorname{int}\co\{e_k,v_i,v_j\}$ for each $k$, with $\{i,j,k\}=\{1,2,3\}$.
\end{definition}
Every interior $\varepsilon$-contamination $T=(1-\varepsilon)p_0+\varepsilon\Delta$ ($p_0\in\Di$) is all-facing and in general position: its edges are parallel to those of $\Delta$ at positive distance, and $v_k$ lies on the open segment from $p_0\in\operatorname{int}T$ to $e_k$ (\src{R15} \S6, \src{R18} \S6).

\subsection{Additive rankings and charts}
An \emph{additive ranking} has measurable penalties $\ell_i(g,a)$, $a\in A(g)$, with an integrable envelope, and ranks world $i$ by $L_i(X)\coloneqq\int_H\ell_i(g,a_X(g))\,d\mu(g)$. Only the advice that evaluable models actually give enters: the penalty is defined for the kinds of advice in the advice range $A(g)$. On a problem on which all evaluable models give the same advice, the penalty is the same for every evaluable model and drops out of every comparison; it is still part of the ranking. On decomposable domains, (G) holds iff $\ell_i(g,\cdot)$ is strictly increasing in world-$i$ rank on $A(g)$ for a.e.\ $g$ (\src{R1} Lemma 1.1); such penalties are \emph{strictly guiding}. A model is $p$-Bayes if it minimises $p\cdot L$ on $\M$; equivalently its advice minimises $p\cdot\ell(g,\cdot)$ over the advice range $A(g)$ for a.e.\ $g$ (\src{R1} Lemma 1.2).

For a pair with $A(g)\subseteq\{+,0,-\}$ put $c\coloneqq\ell(g,0)-\ell(g,+)$, $c'\coloneqq\ell(g,0)-\ell(g,-)$ and $d\coloneqq c-c'$ (vectors over worlds). By (G), $c$ and $d$ have the strict sign pattern of $g$ and $c'$ that of $-g$. The \emph{chart} $R(h)$ of an orientation $h$ is the set of priors at which accepting $h$ is the unique minimiser. For advice range $\{+,0\}$, $R(g)=\{p:c\cdot p>0\}$. For $\{+,0,-\}$, $R(g)=\{c\cdot p>0,\,d\cdot p>0\}$, $R(-g)=\{c'\cdot p>0,\,d\cdot p<0\}$, and no advice is the minimiser on $\{c\cdot p\le0,\,c'\cdot p\le0\}$ (\src{R1} Prop.~2.1). Such a pair is \emph{adjacent} if its d-face $F(g)\coloneqq\{p\in\Di:c\cdot p=c'\cdot p>0\}$ is nonempty; otherwise it is \emph{split}. In words: a pair is adjacent when, at some interior weighting, accepting $g$ and accepting $-g$ are tied as the best advice and each wins on one side, so that the two acceptances compete directly across a boundary between interior weightings. The advice range of a split pair may still contain both acceptances, provided the region where no advice is best separates their charts. Adjacency therefore depends on the penalties, not only on the advice range. Mixed pairs with advice range $\{+,-\}$ (on which every evaluable model accepts one orientation and not the other) are always adjacent (\src{R4} Thm.~A).

\subsection{The \texorpdfstring{$U$}{U}-domain}\label{sec:aos}
Fix a credal set $K'$ in the relative interior of the target $T$ (for $T_1$, $K'\coloneqq u+s(T_1-u)$ with $0<s<1$) and put $U\coloneqq D_{K'}$, an open convex cone containing $\cl D_T\setminus\{0\}$. The \emph{$U$-domain} $\M_U$ is the class of measurable subsets $D$ of $U$ (modulo null sets); it is closed under splicing and under convergence in measure. This is a restriction on the class of evaluable models, and only that; its consequences, for every evaluated model $D\in\M_U$, are the following.
\begin{itemize}[leftmargin=*]
\item $U$ never contains both $f$ and $-f$, because $P_{K'}(f)>0$ and $P_{K'}(-f)>0$ cannot both hold (and $-f\notin U$ for $f\ge0$, $f\ne0$, since $K'\subseteq\Di$). So no evaluated model gives contradictory advice.
\item For $f\in U$ and any $g,h$ with $f=g-h$, $D$ advises choosing $g$ in $\{g,h\}$ iff $f\in D$, and gives no advice iff $f\notin D$.
\item For $f$ with $-f\in U$ the same holds with the roles of $g$ and $h$ swapped: $D$ advises choosing $h$ iff $-f\in D$, and otherwise gives no advice.
\item For $f$ with $f\notin U$ and $-f\notin U$, every evaluated $D$ has $f\notin D$ and $-f\notin D$, so $D$ gives no advice in $\{g,h\}$.
\end{itemize}
Nothing more follows from the restriction. The problems of the last kind (pairs $\{f,-f\}$ with $f\notin U$ and $-f\notin U$) remain decision problems, and the world orders still concern them; every evaluable model gives no advice on them, so they do not distinguish evaluable models. It is not enough that $f\notin U$: if $-f\in U$, evaluable models differ on $f$. (LF) holds at $D_T$ with any open $N$ such that $\partial D_T\setminus0\subseteq N\subseteq U$. On $\M_U$ we index each pair that meets $U$ by its orientation $g\in U$. Only the vector $c(g)=\ell(g,0)-\ell(g,+)$, $g\in U$, enters comparisons between evaluable models (on every pair $\{f,-f\}$ with $f\notin U$ and $-f\notin U$ all evaluable models give no advice, so its penalty drops out), and the $q$-Bayes model is
\[
B_q=\{g\in U:\ q\cdot c(g)>0\}.
\]
An additive ranking on a general decomposable domain is \emph{non-adjacent near $T$} if, for some credal set $K'\subseteq\relint T$, a.e.\ pair with an orientation in $D_{K'}$ is non-adjacent. The advice range of such a pair may contain both acceptances; what is excluded is that they compete directly as best advice across a boundary between interior weightings. This holds on $\M_U$, whenever the advice ranges of these pairs are contained in $\{+,0\}$ (a sufficient special case that concerns the domain alone), and whenever every interior Bayes credal set is two-dimensional (\src{R4} Thm.~B).

\subsection{Coherent almost-preferences}\label{sec:CPa-def}
Let $L^\infty(H,\mu)$ be the bounded measurable functions on $H$ modulo $\mu$-null sets. A \emph{state} is a positive normalised linear functional on $L^\infty(H,\mu)$, that is, a finitely additive probability on the measurable sets that vanishes on $\mu$-null sets.

\begin{definition}\label{def:CPa}
World $i$ of a ranking on a decomposable domain $\M$ is a \emph{coherent almost-preference over losses}, condition (CP-a), if there are
\begin{enumerate}[leftmargin=*,label=(\roman*)]
\item a bounded measurable penalty $\lambda_i(g,a)$ with $\lambda_i(g,a)<\lambda_i(g,b)$ whenever $a$ is strictly better than $b$ in world-$i$ rank at $g$ (for a.e.\ $g$), and
\item a coherent lower prevision $\lp_i$ on $L^\infty(H,\mu)$ in the sense of Walley~\cite{W1991}, with ``$f\ge0$'' read as $f\ge0$ a.e.,
\end{enumerate}
such that, writing $\Lambda_iX\coloneqq\lambda_i(\cdot,a_X(\cdot))$,
\[
X\preceq_iY\iff\lp_i(\Lambda_iY-\Lambda_iX)\ge0 .
\]
\end{definition}
Equivalently (\cite[Thm.~3.8.1]{W1991}, together with the lower envelope theorem \cite[Thm.~3.3.3]{W1991} and weak-* compactness \cite[\S3.6]{W1991}): $\{f:\lp_i(f)\ge0\}$ is a coherent class of almost-desirable gambles (axioms D0--D4); or $\preceq_i$ is the pull-back along $X\mapsto-\Lambda_iX$ of a coherent almost-preference (axioms R0--R5); or there is a nonempty weak-* closed convex set $\M_i$ of states with
\[
X\preceq_iY\iff P(\Lambda_iY-\Lambda_iX)\ge0\ \text{for every }P\in\M_i .
\]
The null-set reading is a restriction in Walley's own sense: (CP-a) holds iff there is a Walley-coherent lower prevision on bounded measurable gambles that is \emph{null-invariant} ($\lp_i(f)=\lp_i(f')$ whenever $f=f'$ a.e.), equivalently every member of its credal set vanishes on null sets (\src{N24} \S5). Not assumed: countable additivity, separability or non-atomicity of $\M_i$; completeness; conglomerability. World $i$ is \emph{complete} if $\M_i$ is a singleton. Instances include additive rankings with integral penalties ($\M_i=\{\rho\mu\}$), additive rankings with bounded charge penalties ($\M_i=\{P\}$), Pareto unanimity of additive rankings that share the encoding, and credal-set rankings. On $\M_U$ the canonical encoding is the error indicator $\Lambda_iX=1_{E_i(X)}$ with $E_i(X)\coloneqq X\sd(U\cap\{h_i>0\})$ (a function on $U$, carried to the pairs by the measure-preserving bijection $g\mapsto\operatorname{sign}(g_1)\,g$ from $U$ onto the representatives of the pairs meeting $U$; it vanishes on the pairs with $f\notin U$ and $-f\notin U$).

\subsection{Adding up and supporting additive assessments}\label{sec:addup-def}
This subsection fixes the terms used for rankings that are not assumed to be additive (Glossary of \src{START HERE}; \src{N17}, \src{N19}). Work on $\M_U$ with a ranking $L$ satisfying strict (G), (IB), (LA), (S), (C). The \emph{error set} $E_i(X)$ of a model at world $i$ is the set of $h\in U$ with $h_i>0$ that $X$ does not accept, together with those with $h_i<0$ that it accepts, as above. A model is \emph{regular} if it gives the correct advice on unanimous problems, i.e.\ accepts a.e.\ nonnegative gamble of $U$.

\begin{definition}\label{def:addup}
\begin{enumerate}[leftmargin=*,label=(\alph*)]
\item The \emph{relevant comparisons} at world $i$ are the pairs $B\preceq_iX$ with $B\in\C_L$ (coherent and undominated) and $X$ regular.
\item World $i$ \emph{adds up} if some finitely additive weight $\nu_i$ on $U$, zero on null sets and positive on non-null sets, endorses every relevant comparison: $\nu_i(E_i(B))\le\nu_i(E_i(X))$.
\item If $\nu_i$ is such a weight, let $\rho_i$ be the density of its countably additive part; then $\rho_i>0$ a.e.\ on $U$ (\src{N19} Lemma R1). The integral loss $\psi_i(X)\coloneqq\int_{E_i(X)}\rho_i\,d\mu$ is a \emph{supporting additive assessment}. It is an additive ranking on the same domain. It endorses the relevant comparisons when $\nu_i$ is countably additive, but it is not claimed to be numerically or topologically close to $\preceq_i$, and it can lose distinctions that $\preceq_i$ makes.
\end{enumerate}
\end{definition}

The weights in (b) are real-valued, hence bounded by $\nu_i(U)$. This is the convention of \src{N19} (its Lemma~R1 is stated for bounded charges), and (c) presupposes it, since the Yosida--Hewitt decomposition is a decomposition of bounded charges. \emph{Round 16:} the convention matters. In the two-world system of Theorem~\ref{thm:D21} no bounded weight that is positive on non-null sets endorses the relevant comparisons of world~1. A countably additive, $\sigma$-finite measure of infinite total mass endorses all of them (Proposition~\ref{prop:D2sigma}), and so does a suitable extended-valued finitely additive weight (status of Theorem~\ref{thm:D21}).

Restricting to regular $X$ is essential. Ranking first by the mass of unanimous problems of $U$ on which the model gives no advice, and only then by $\preceq_i$, preserves every axiom, the selection, $\C_L$ and the covered target, yet it destroys adding up for unrestricted comparisons (\src{N19} \S3.3, Prop.~R3, proved in that referee report). Such a priority manufactures apparent non-additivity without helping coverage, so only the relevant comparisons are tested.

Exactly one of the following holds.
\begin{itemize}[leftmargin=*]
\item[\textbf{(M0)}] Every world adds up.
\item[\textbf{(M1)}] At some world a finite family of relevant comparisons is an \emph{obstruction} to strictly positive error weights: some positive combination of their error-difference vectors $1_{E_i(X)}-1_{E_i(B)}$ is $\le0$ a.e.\ and not a.e.\ zero. The comparisons need not form an executable sequence of trades.
\item[\textbf{(M2)}] At some world there is no finite obstruction, but still no strictly positive weight endorses all relevant comparisons.
\end{itemize}
The division is logical bookkeeping. The substantive result is the reduction behind (M0) (Proposition~\ref{prop:M0}): supporting weights, if they exist, yield an additive system that is coherent everywhere except for ties at exceptional weightings. The restrictions known for (M1) (at least four comparisons, overlapping gains, crossing error sets; \src{N17} N17.3, refereed in \src{N19} \S4) concern (M1) certificates in this formulation only; they do not constrain every failure of adding up. On $\M_U$ a cover of $T_1$ must fall under (M1) or (M2) (Corollary~\ref{cor:M0closed}).

\section{Additive coverage as a realisation problem}\label{sec:reduction}

From now on $n=3$. This section reduces additive coverage of an interior target to a statement about a family of credal sets.

\begin{proposition}\label{prop:necessary}
Let $L$ be a strictly guiding additive ranking on a decomposable domain $\M$ satisfying (LA), (S), (C). Then:
\begin{enumerate}[leftmargin=*,label=(\alph*)]
\item for every $p\in\Di$ the $p$-Bayes model is a.e.\ unique and equals $D_{K_p}$ a.e.\ for a credal set $K_p$;
\item at each vertex $e_i$ the Bayes model is the world-$i$-best evaluable model, and it is coherent;
\item at every $p\in\partial\Delta$ some $p$-Bayes model is coherent;
\item if $L$ covers $T=D_K$ with $K\subseteq\Di$, then $T$ is the $p^*$-Bayes model for some $p^*\in\Di$.
\end{enumerate}
\end{proposition}
\status{proved; refereed once (R6). Sources: \src{R1} Thm.~2.2(N), Lemmas 1.4--1.5, Prop.~2.6; \src{R2} Lemma U. Uniqueness here is uniqueness of the Bayes minimiser at each interior weight, not of the selection, which may be reindexed (Consolidated Statements \S0.3).}

\begin{definition}\label{def:realisation}
Let $T\subseteq\Di$ be a credal set and $U=D_{K'}$ as above. A \emph{realisation} of $T$ on $\M_U$ at $p^*\in\Di$ is a measurable $c:U\to\R^3$ such that $c_i(g)$ has the strict sign of $g_i$ for a.e.\ $g$ and:
\begin{enumerate}[leftmargin=*]
\item[\textbf{(BU)}] for \emph{every} $q\in\Di$ the tie set $\{g\in U:q\cdot c(g)=0\}$ is null and $B_q=D_{K_q}$ a.e.\ for a credal set $K_q$;
\item[\textbf{(TP)}] $K_{p^*}=T$.
\end{enumerate}
\end{definition}
(These are conditions (R1), (R2) of \src{N12}.) By Proposition~\ref{prop:necessary}, a strictly guiding additive ranking on $\M_U$ that covers $T$ yields a realisation; Theorem~\ref{thm:ccibu} gives the converse. Write $P(q,g)\coloneqq P_{K_q}(g)$ and, for $g\in U$, $H(g)\coloneqq\{q\in\aff\Delta:q\cdot c(g)>0\}$, an open half-plane.

\begin{proposition}\label{prop:reduction}
Let $c$ be a realisation of $T$ at $p^*$. Then:
\begin{enumerate}[leftmargin=*,label=(\alph*)]
\item $K_q\supseteq K'$ for every $q\in\Di$, and $K_q$ is unique;
\item $q\mapsto K_q$ is Hausdorff-continuous on $\Di$, and $P$ is jointly continuous with $|P(q,g)-P(q',g')|\le\|g-g'\|_\infty+\|g\|\,\dH(K_q,K_{q'})$;
\item the vertex models $U\cap\{g_i>0\}=D_{\co(K'\cup\{e_i\})}$ are coherent;
\item $\tilde c\coloneqq p^*\odot c$ is a realisation of $T$ at $u$;
\item \emph{(half-plane charts)} for a.e.\ $g\in U$, $\{q\in\Di:P(q,g)>0\}$ equals $H(g)\cap\Di$ up to a null set, and $e_i\in H(g)$ if $g_i>0$, $e_i\notin\cl H(g)$ if $g_i<0$.
\end{enumerate}
Conversely, a Hausdorff-continuous family $q\mapsto K_q\supseteq K'$ with $K_{p^*}=T$ and vertex-respecting half-plane charts as in (e) comes from a realisation, and its tie sets are automatically null.
\end{proposition}
\begin{proof}[Proof sketch]
(a) $D_{K_q}$ is open, convex and a.e.\ inside the open convex $U$, so $D_{K_q}\subseteq U$, i.e.\ $K_q\supseteq K'$; injectivity of $K\mapsto D_K$ gives uniqueness. (b) If $q_m\to q$ then $1\{q_m\cdot c>0\}\to1\{q\cdot c>0\}$ off the null tie set of $q$, so $B_{q_m}\to B_q$ in measure; for cones with interior this is Hausdorff convergence of credal sets (\src{R1} Prop.~3.1). (c) $D_{K_1}\cap D_{K_2}=D_{\co(K_1\cup K_2)}$. (d) $\tilde q\cdot\tilde c(g)=(\tilde q\odot p^*)\cdot c(g)$, and normalising $\tilde q\odot p^*$ is a projective bijection of $\Delta$ fixing the vertices. (e) Fubini on $\{(q,g):1[q\cdot c(g)>0]\ne1[P(q,g)>0]\}$. Converse: \src{N12} Prop.~1.4 with the simplification of \src{N14} \S2.1.
\end{proof}
\status{proved; refereed once (N14: CORRECT, minor repair incorporated). Source: \src{N12} \S1 (Lemma 1.1, Props 1.2--1.4). The irregularity of $c$ has no effect beyond the family $q\mapsto K_q$.}

\begin{remark}\label{rem:N20}
Condition (BU) must hold at \emph{every} interior $q$, or Hausdorff continuity must be assumed. If (BU) held only for a.e.\ $q$, then $K_q$ could be altered on a null set of priors that contains $p^*$, and (TP) would carry no information (\src{N20} \S8). For a covering ranking this is automatic, by Proposition~\ref{prop:necessary}(a).
\end{remark}

\begin{theorem}\label{thm:ccibu}
Let $n=3$ and let $L$ be an additive ranking on a decomposable domain with integrable envelopes and penalties strictly increasing in world rank a.e. Suppose that for every $w\in\Di$ the minimiser of $w\cdot L$ on $\M$ is unique modulo null sets and coherent. Then $L$ satisfies (LA), (S), (C), and every undominated model is coherent and the only model with its loss profile.
\end{theorem}
\begin{proof}[Proof sketch]
The attainable set $\mathcal R=\{L(X):X\in\M\}$ is compact and convex (Lyapunov). Every Pareto-undominated profile has a supporting weight $w$. If $w$ is interior, use the hypothesis. If $w=e_i$, strict guidance gives a unique minimiser, and it is the limit of interior Bayes models. If $w=(w_1,w_2,0)$, a finite subgradient of $t\mapsto\min\{y_3:y\in\operatorname{argmin}w\cdot y,\ y_1=t\}$ turns the profile into a directional minimiser on $\operatorname{argmin}_aw\cdot\ell(g,a)$; directional minimisers are unique (a tie would persist at nearby interior weights) and are limits of interior Bayes models, hence coherent. Then select $x(p)$ by minimising $Q_p(x)=\sum_ip_i(x_i-a_i)^2$ over $\mathcal R$ (with $a_i<\min_{\mathcal R}x_i$), followed at edge indices by a secondary minimisation of the omitted coordinate. Each $x(p)$ is undominated. Coordinates of $x(p)$ are lower semicontinuous in $p$, which makes the sections closed, and strict monotonicity of $Q_p$ gives competitiveness on $\supp p$.
\end{proof}
\status{author-checked (\src{CCIBU.tex}; also checked by the coordinator, \src{BRIEF2}); Lemmas 2--4 re-derived without error in \src{N13} \S2; no formal referee report. The argument is specific to three worlds.}

\begin{corollary}\label{cor:realisation}
Let $T=D_K$ with $K\subseteq\Di$, and let $\M$ be a decomposable domain with (LF) at $T$. Some strictly guiding additive ranking with integrable envelopes on $\M$ that satisfies (LA), (S), (C) covers $T$ iff some strictly guiding additive ranking with integrable envelopes on $\M$ has all interior Bayes models a.e.\ unique and coherent, with $T$ among them. On $\M_U$ this holds iff $T$ has a realisation (Definition~\ref{def:realisation}).
\end{corollary}
\status{proved, from Proposition~\ref{prop:necessary} and Theorem~\ref{thm:ccibu} (so its status is that of Theorem~\ref{thm:ccibu} in the ``if'' direction). For the realisation form, penalties are recovered from $c$ by \src{R1} Prop.~2.1 (converse) and a positive rescaling for integrability, which does not change charts.}

\section{No additive cover of the contamination triangle}\label{sec:main}

\begin{theorem}\label{thm:main}
Let $T\subseteq\Di$ be an all-facing triangle in general position, and $p^*\in\Di$. Then $T$ has no realisation at $p^*$ on any $U$-domain $\M_U$. In particular, for every $\varepsilon\in(0,1)$, $T_1$ has no realisation at any interior target weighting.
\end{theorem}
\status{proved; refereed independently twice (\src{N18}, \src{N20}: both CORRECT, with the local repairs incorporated below). Source: \src{N15} (main chain), assembled with \src{N12} Lemma 2.1, Lemma 4.3, Thm.~2.4, Thm.~3.1 (refereed once, \src{N14}), \src{R10b} Lemma UVN and Cor.~PL2 (refereed once, \src{R12}), and \src{R15} Thm.~6.2 algebra (refereed once, \src{R18}). No hypothesis of differentiability, nondegeneracy or regularity of $q\mapsto K_q$ is used beyond the Hausdorff continuity that (BU) provides (Remark~\ref{rem:N20}).}

\begin{corollary}\label{cor:coverage}
Let $L$ be a strictly guiding additive ranking with integrable envelopes on a decomposable domain $\M$, satisfying (G), (IB), (LA), (S), (C). If $L$ is non-adjacent near $T$ (for $T$ as in Theorem~\ref{thm:main}), then $L$ does not cover $T$. In particular $T_1$ is not covered by any additive ranking on the $U$-domain, on any decomposable domain on which the pairs with an orientation in some $D_{K'}$, $K'\subseteq\relint T_1$, have advice ranges contained in $\{+,0\}$, or whenever all interior Bayes credal sets are two-dimensional.
\end{corollary}
\status{proved; refereed twice (\src{N18} \S1.10, \src{N20} \S7: Cor.~7.2 of \src{N15} CORRECT). It rests on the reduction of \src{N14} \S3(iv), which was proved in that referee report: under non-adjacency, $\hat K_q\coloneqq\co(K_q\cup K')$ satisfies Proposition~\ref{prop:reduction} with $U=D_{K'}$, so by its converse it is a realisation on $\M_{D_{K'}}$. The sources also list domains whose advice ranges near $\partial D_T$ have two elements (``two-action'' domains there) as an instance; the hypothesis actually used, through \src{N14} \S3(iv), is non-adjacency of all pairs meeting $D_{K'}$, which is how it is stated here.}

\begin{corollary}\label{cor:charges}
The conclusion of Corollary~\ref{cor:coverage} holds when the penalties are bounded charges vanishing on $\mu$-null sets, provided the interior Bayes minimisers exist.
\end{corollary}
\status{proved; \src{N14} Thm.~C2 (proved in that referee report; its key splice step is re-proved and refereed as \src{N22} Lemma 3.0 / \src{N24}), used in \src{N15} Cor.~7.3, refereed CORRECT conditional on existence in \src{N18} and \src{N20}. The existence proviso is discharged for complete rankings in Theorem~\ref{thm:complete}.}

The rest of this section proves Theorem~\ref{thm:main}. Fix a realisation $c$ of an all-facing triangle $T$ in general position at $p^*$. By Proposition~\ref{prop:reduction}(d) we may and do assume $p^*=u$; nothing below uses this before \S\ref{sec:contradiction}, and the earlier steps hold verbatim at any interior $p^*$. Put $r_0\coloneqq\operatorname{dist}(p^*,\partial\Delta)$ in $\aff\Delta$, so $B(p^*,r_0)\subseteq\Di$.

\subsection{Chords and straight levels}

For each facet $F_k$ with corner rays $r,r'$ put $h(\tau)\coloneqq(1-\tau)r+\tau r'$ and $\Psi(q,\tau)\coloneqq P(q,h(\tau))$. For each $q$, $\tau\mapsto\Psi(q,\tau)$ is concave on all of $\R$, because $P(q,\cdot)$ is concave and $h$ is affine. Since $K_{p^*}=T$, $\Psi(p^*,\tau)=0$ for $\tau\in[0,1]$. A gamble $g\in U$ is \emph{good} if $\{q:P(q,g)=0\}$ is null and the half-plane chart of Proposition~\ref{prop:reduction}(e) holds for it; a.e.\ $g$ is good. A \emph{cluster line} of a mixed $g\in U$ is a limit, with orientation, of the lines $\partial H(g_m)$ for good $g_m\to g$ (normalise $c(g_m)$ to unit length; a limit unit vector defines an affine function on $\aff\Delta$ and its zero line). Cluster lines exist by compactness.

\begin{lemma}\label{lem:signs}
On each facet $F_k\setminus\{0\}$, $h_k<0$. After relabelling $i,j$ so that $r=r_{ki}$, the coordinates $h_i,h_j$ each cross zero exactly once on $(0,1)$, at $\tau_i<\tau_j$, with sign patterns $(-,+)$ on $[0,\tau_i)$, $(+,+)$ on $(\tau_i,\tau_j)$ and $(+,-)$ on $(\tau_j,1]$.
\end{lemma}
\begin{proof}
Write $v_k=ae_k+bv_i+cv_j$ with $a,b,c>0$ (all-facing). For $0\ne h\in F_k$, $0=h\cdot v_k=ah_k+b\,h\cdot v_i+c\,h\cdot v_j$, and $h\cdot v_i,h\cdot v_j\ge0$ are not both zero (the $v$'s span $\R^3$), so $h_k<0$. At $r_{ki}$, $r_k<0$ and $r_i<0$ (the same fact on $F_i$), and $r_j>0$, since $r\le0$ would give $r\cdot v_j<0$. Likewise at $r_{kj}$: $r_k,r_j<0<r_i$. Along the facet $h_i,h_j$ are affine, so each crosses once; they cannot both be $\le0$ at one point, since then $h\le0$, $h\ne0$ and $h\cdot v_k<0$. Hence at the crossing of $i$ we have $h_j>0$, and $\tau_i<\tau_j$.
\end{proof}
\status{proved; \src{N15} \S1 for $T_1$, general proof supplied independently by both referees of \src{N15} (\src{N18} (R-c), \src{N20} \S4.0).}

\begin{lemma}\label{lem:chords}
\begin{enumerate}[leftmargin=*,label=(\alph*)]
\item For good $g$: $P(\cdot,g)\ge0$ on $\cl H(g)\cap\Di$, $P(\cdot,g)\le0$ on $\Di\setminus H(g)$, and $P(\cdot,g)=0$ on $\partial H(g)\cap\Di$.
\item For every mixed $g\in U$ and every cluster line $\ell$ with closed sides $\Pi^\pm$: $P(\cdot,g)\ge0$ on $\Pi^+\cap\Di$, $P(\cdot,g)\le0$ on $\Pi^-\cap\Di$, $P(\cdot,g)=0$ on $\ell\cap\Di$; and $e_i\in\Pi^+$ if $g_i>0$, $e_i\in\Pi^-$ if $g_i<0$.
\end{enumerate}
\end{lemma}
\begin{proof}
(a) The open set $\{P(\cdot,g)>0\}$ is a.e.\ inside the open half-plane $H(g)$, hence inside it; similarly $\{P<0\}\subseteq\Di\setminus\cl H(g)$; continuity in $q$ gives the closed statements. (b) Pass to the limit using joint continuity (Proposition~\ref{prop:reduction}(b)); the vertex statements are the limits of the strict signs of $c$.
\end{proof}
\status{proved; refereed once (\src{N14}: CORRECT). Source: \src{N12} Lemma 2.1.}

The level shift $P(q,g-t\one)=P(q,g)-t$ (since $m\cdot\one=1$) turns the chord statement at level $0$ into one at every level. For $g\in\R^3$ let $I_g\coloneqq\{t:\ g-t\one\in U\text{ and is mixed}\}=(\min_ig_i,\,P_{K'}(g))$.

\begin{lemma}\label{lem:levels}
Let $f\coloneqq P(\cdot,g)$ on $\Di$ and $t\in I_g$. Then $\{f>t\}$ and $\{f<t\}$ are each the intersection of $\Di$ with an open half-plane, or $\varnothing$, or $\Di$. Consequently, when both half-planes are proper, $\{f=t\}$ is the closed region between two lines that do not meet in $\Di$.
\end{lemma}
\begin{proof}
For $t'\in I_g$, Lemma~\ref{lem:chords}(b) applied to $g-t'\one$ gives a line with closed sides $\Pi^\pm_{t'}$ such that $f\ge t'$ on $\Pi^+_{t'}\cap\Di$ and $f\le t'$ on $\Pi^-_{t'}\cap\Di$. \emph{$\{f>t\}$ is convex:} if $f(q_1),f(q_2)>t$, choose $t'\in(t,\min f(q_i))$, which lies in $I_g$ because $f\le P_{K'}(g)$; then $q_1,q_2\in\operatorname{int}\Pi^+_{t'}$, and the segment lies in $\operatorname{int}\Pi^+_{t'}\cap\Di\subseteq\{f\ge t'\}$. \emph{$\{f\le t\}$ is convex:} it equals $\bigcap_{t'\in(t,t+\eta)}\Pi^-_{t'}\cap\Di$ for small $\eta$. Two complementary convex subsets of the convex set $\Di$, one of them open, are separated by a line, which gives the half-plane form. The case $\{f<t\}$ is symmetric. If the two boundary lines crossed at $x\in\Di$, the two disjoint open half-planes would share a wedge at $x$.
\end{proof}
\status{proved; refereed twice (\src{N18} \S1.1, \src{N20} \S1, wording repair: ``region between two lines'' rather than ``strip''). Source: \src{N15} Lemma 2.1.}

\begin{lemma}\label{lem:sandwich}
Let $f=P(\cdot,g)$, $f(p^*)=0$, and suppose $\{f>0\}=\Di\cap\{n\cdot(q-p^*)>0\}$ for a unit $n\in\one^\perp$ and $f\le0$ on $\lambda\cap B(p^*,r_0)$, where $\lambda\coloneqq p^*+\R u_\perp$ with $u_\perp\perp n$ a unit vector. Put $D^+(s)\coloneqq f(p^*+sn)$ and $D^-(s)\coloneqq-f(p^*-sn)$. Then $D^\pm$ are nondecreasing near $0$, and for $q=p^*+yu_\perp+xn$ with $|y|\le r_0/2$ and $x>0$ small,
\[
D^+\Bigl(\frac{x\,r_0}{r_0+|y|}\Bigr)\ \le\ f(q)\ \le\ D^+\Bigl(\frac{x\,r_0}{r_0-|y|}\Bigr),
\]
with the mirror statement for $x<0$.
\end{lemma}
\begin{proof}
Monotonicity: for $t>0$ in $I_g$, $\{f>t\}$ is a half-plane inside $\{f>0\}$ missing $p^*$, so it meets the ray $p^*+\R_+n$ in a terminal sub-ray. Upper bound: if $D^+(s_1)<f(q)$ with $s_1>xr_0/(r_0-|y|)$, choose a level $t$ strictly between; the convex set $\{f<t\}$ contains the open chord $\{p^*+y'u_\perp:|y'|<r_0\}$ and $p^*+s_1n$, hence the open triangle they span, whose height over $y$ is $s_1(1-|y|/r_0)>x$; so $f(q)<t$, a contradiction. The lower bound is the same argument with the convex set $\{f\le t\}$ and apex $q$.
\end{proof}
\status{proved; refereed twice (\src{N18} \S1.3, \src{N20} \S1, with the convex-hull proof of the upper bound and the weakened hypothesis $f\le0$ on $\lambda\cap B(p^*,r_0)$). The constants $r_0/(r_0\pm|y|)$ are sharp. Source: \src{N15} Lemmas 2.2--2.3.}

\begin{lemma}\label{lem:caps}
Let $\gamma:\R\to\R$ be concave.
\begin{enumerate}[leftmargin=*,label=(\alph*)]
\item If $\{\gamma>0\}\ne\varnothing$, then $\{\gamma=0\}=\partial\{\gamma>0\}$ (at most two points) and $\gamma<0$ off $\cl\{\gamma>0\}$; if $\gamma$ vanishes on a nondegenerate interval, $\gamma\le0$ everywhere.
\item If $\gamma(\tau_0)=0$ and $\gamma>0$ on $(\tau_0,\tau_0+\eta)$, then $\gamma(\tau)/(\tau-\tau_0)$ is nonincreasing on $\R\setminus\{\tau_0\}$, and $0<\gamma'_+(\tau_0)\le\gamma'_-(\tau_0)<\infty$. The mirror statement holds on the left.
\item For a jointly continuous $\Gamma(q,\tau)$ concave in $\tau$, $\partial^-_\tau\Gamma=\inf_{s>0}(\Gamma(q,\tau)-\Gamma(q,\tau-s))/s$ is upper semicontinuous and $\partial^+_\tau\Gamma=\sup_{s>0}(\Gamma(q,\tau+s)-\Gamma(q,\tau))/s$ is lower semicontinuous.
\end{enumerate}
\end{lemma}
\status{proved (standard convex analysis); refereed twice (\src{N18} \S1.4, \src{N20} \S2, endpoint slip repaired). Source: \src{N15} Lemma 3.1.}

\subsection{Vertex neighbourhoods and single-vertex pinning}

\begin{lemma}\label{lem:UVN}
Let $x\in U$ be mixed and $l$ a world. If $x_l>0$ there are a ball $B_1\ni x$ and a relative neighbourhood $V_l$ of $e_l$ in $\Delta$ with $P(q,h)>0$ for all $q\in V_l\cap\Di$, $h\in B_1$, and $H(h)\supseteq V_l\cap\Di$ for a.e.\ $h\in B_1$. If $x_l<0$ the same holds with $P(q,h)<0$ and $H(h)\cap V_l=\varnothing$.
\end{lemma}
\status{proved; refereed once (\src{R12}: CORRECT); second reading CONFIRMED (\src{TS\_second\_referee} \S1.1). Source: \src{R10b} Lemma UVN. It uses the strict minimiser of the vertex objective and dominated convergence as $q\to e_l$, together with upward (downward) closure of coherent cones.}

Fix a vertex $v=v_k$ of $T$ with a crossing $w\in\relint F_k$, $w_c=0$, and a small ball $\mathcal U$ around $w$ inside $U$, meeting no other facet, and inside the balls of Lemma~\ref{lem:UVN} at $w$ and at points $w_\pm\in F_k$ near $w$ with $(w_\pm)_c\gtrless0$. Put $\mathcal U^+=\mathcal U\cap D_T$, $\mathcal U^-=\mathcal U\setminus\cl D_T$ and
\begin{gather*}
\Gamma_v\coloneqq\{q\in\Di:D_{K_q}\supseteq\mathcal U^+\},\qquad
\beta_v\coloneqq\Gamma_v\setminus\alpha_v,\\
\alpha_v\coloneqq\{q\in\Gamma_v:v\in K_q\}=\{q\in\Gamma_v:D_{K_q}\cap\mathcal U^-=\varnothing\}.
\end{gather*}
$\beta_v$ is relatively open in $\Gamma_v$ (\src{R10b} Lemma D, as repaired in \src{R12}). For $q\in\beta_v$, $D_{K_q}$ contains an exterior collar of $F_k\cap\mathcal U$; for $q\in\alpha_v$, $P(q,h)<0$ for all $h\in\mathcal U^-$.

\begin{lemma}\label{lem:PL2}
Let $G\subseteq\Gamma_v$ be open and connected, and suppose a.e.\ $h\in\mathcal U^-$ has no d-face meeting $G$ (automatic on $\M_U$, where no advice range contains both acceptances). Then $G\subseteq\alpha_v$ or $G\subseteq\beta_v$.
\end{lemma}
\begin{proof}[Proof sketch]
It suffices to show that no $Q\in G$ lies in $\cl\beta_v\setminus\beta_v$. Suppose $Q$ does; then $Q\in\alpha_v$, so $P(Q,h)<0$ on $\mathcal U^-$, and $Q\notin\cl H(h)$ for good $h\in\mathcal U^-$. Choose $s_m\in\beta_v$ with $s_m\to Q$, and good $h^\pm_m\in\mathcal U^-$ converging to $w_\pm$, so close to the facet that $P(\cdot,h^\pm_m)>0$ near $s_m$ (collar). Then the chart lines of $h^\pm_m$ pass within $2|s_m-Q|$ of $Q$; normalised, $c(h^\pm_m)\to a_\pm$ with $a_\pm\cdot Q=0$. Every point of $\{a_+\cdot q>0\}$ near $Q$ has a neighbourhood a.e.\ accepting some $h^+_m\in\mathcal U^-$, so it lies in $\cl\beta_v$; every point of $\beta_v$ near $Q$ has $P(\cdot,h^-_m)>0$ for large $m$, so $a_-\cdot q\ge0$ there. Hence $\{a_+>0\}\subseteq\cl\{a_->0\}$ near $Q$, and since both lines pass through $Q$, $a_+=\lambda a_-$ with $\lambda>0$. But $(a_-)_c\le0$, because $c_c(h^-_m)<0$, while Lemma~\ref{lem:UVN} at $w_+$ for world $c$ gives $\{a_+\cdot q>0\}\supseteq V_c\cap\Di$, hence $(a_+)_c>0$ (a line through $e_c$ and the interior point $Q$ would cut $V_c$). This is a contradiction. So $\beta_v$ is relatively clopen in the connected set $G$.
\end{proof}
\status{proved; refereed once (\src{R12}: CORRECT); second reading CONFIRMED (\src{TS\_second\_referee} \S1.1). Source: \src{R10b} Lemma M (Steps 2--4) and Cor.~PL2. Its validity for sets $G$ that do not contain $p^*$, which is how it is used below, was re-checked by both referees of \src{N15} (\src{N18} \S1.5, \src{N20} \S3).}

\begin{lemma}\label{lem:nopin}
Let $\tau_c\in(\sigma_1,\sigma_2)\subseteq[0,1]$ be a crossing parameter of $F_k$, and let $G\subseteq\Di$ be nonempty, open and connected, with $\Psi(\cdot,\sigma_1)\ge0$ and $\Psi(\cdot,\sigma_2)\ge0$ on $G$. Then $\Psi>0$ on $G\times(\sigma_1,\sigma_2)$.
\end{lemma}
\begin{proof}
By concavity in $\tau$, $\Psi\ge0$ on $G\times[\sigma_1,\sigma_2]$. Choose the ball $\mathcal U$ around $w=h(\tau_c)$ so small that $F_k\cap\cl\tilde{\mathcal U}\subseteq\R_+h([\sigma_1,\sigma_2])$ for a slightly larger concentric ball $\tilde{\mathcal U}$. By homogeneity $K_q$ lies in the wedge $W\coloneqq\{m\in\aff\Delta:m\cdot h(\sigma_1)\ge0,\ m\cdot h(\sigma_2)\ge0\}$, with apex $v_k$, and then $q\in\Gamma_{v_k}$ (\src{R10b} \S1: write $h\in\mathcal U^+$ as $w''+(v_k\cdot h)\one$ with $w''\in F_k\cap\tilde{\mathcal U}$). Pick $j$ with $w_j>0$; by Lemma~\ref{lem:UVN}, $V_j\cap\Di\subseteq\Gamma\cap\beta$. Every good $h\in\mathcal U^+$ has $H(h)\supseteq G\cup(V_j\cap\Di)$, so $\operatorname{int}\co(G\cup V_j)\subseteq\Gamma$ (a Fubini step, checked in \src{N20} \S3). By Lemma~\ref{lem:PL2} this connected open set lies in $\alpha$ or in $\beta$; it meets $\beta$ at $V_j$, so it lies in $\beta$, and $G\subseteq\beta$, i.e.\ $v_k\notin K_q$ for $q\in G$. Since $h(\tau)$ is a strictly positive combination of $h(\sigma_1),h(\sigma_2)$ for $\tau\in(\sigma_1,\sigma_2)$, and these are independent modulo $\one$, $m\cdot h(\tau)>0$ on $W\setminus\{v_k\}$; compactness of $K_q$ gives $\Psi(q,\tau)>0$.
\end{proof}
\status{proved; refereed twice (\src{N18} \S1.5, \src{N20} \S3). Source: \src{N15} Lemma 4.1. This is the step where non-adjacency enters (through Lemma~\ref{lem:PL2} on a set that leaves every ball around $p^*$).}

\subsection{Pivot chords and corners}

Write $A_\tau\coloneqq\{\Psi(\cdot,\tau)>0\}$, $B_\tau\coloneqq\{\Psi(\cdot,\tau)<0\}$, $Z_\tau\coloneqq\{\Psi(\cdot,\tau)=0\}$, $P_\tau\coloneqq\Di\setminus B_\tau$, all subsets of $\Di$.

\begin{theorem}\label{thm:pivot}
For every facet $F_k$ and $\tau\in(0,1)$ there is a line $\lambda_\tau\ni p^*$ with $Z_\tau=\lambda_\tau\cap\Di$; $A_\tau$ and $B_\tau$ are the two open sides of $\lambda_\tau$ in $\Di$. If $n_\tau$ denotes the unit normal of $\lambda_\tau$ in $\one^\perp$ pointing into $A_\tau$, then $\tau\mapsto n_\tau$ is continuous and strictly monotone in angle, $\lambda_{\tau_i}=\lin(p^*,e_i)$ at each crossing, and $n_\tau\cdot(e_l-p^*)>0$ whenever $h(\tau)_l>0$. Every cluster line of $h(\tau)$ equals $\lambda_\tau$, with the same orientation.
\end{theorem}
\begin{proof}[Proof (steps in the order required by \src{N18} (R-a) and \src{N20} \S4)]
\emph{Step 0.} By Lemma~\ref{lem:UVN}, $A_\tau\supseteq V_l\cap\Di$ when $h(\tau)_l>0$ and $B_\tau\supseteq V_k\cap\Di$; by Lemma~\ref{lem:levels} at level $0\in I_{h(\tau)}$, $A_\tau$ and $B_\tau$ are proper open half-planes intersected with $\Di$, with boundary lines $\lambda^\pm_\tau$, and $Z_\tau$ lies between them.

\emph{Step 1 ($p^*\in\lambda^+_\tau$).} If $\tau<\tau_j$ take $\sigma_1\in(0,\min(\tau,\tau_i))$, $\sigma_2\in(\max(\tau,\tau_i),\tau_j)$; if $\tau\ge\tau_j$ take $\sigma_1\in(\tau_i,\min(\tau,\tau_j))$, $\sigma_2\in(\max(\tau,\tau_j),1)$. In both cases $[\sigma_1,\sigma_2]$ contains a crossing, and $h(\sigma_1),h(\sigma_2)$ share a positive coordinate $l$ (Lemma~\ref{lem:signs}). Then $G\coloneqq\operatorname{int}(P_{\sigma_1}\cap P_{\sigma_2})\supseteq V_l(\sigma_1)\cap V_l(\sigma_2)\cap\Di$ is nonempty, open and convex, and $p^*\in\cl G$. Lemma~\ref{lem:nopin} gives $G\subseteq A_\tau$, so $p^*\in\cl A_\tau\setminus A_\tau=\lambda^+_\tau\cap\Di$. A line through $p^*$ and $e_l$ with $h(\tau)_l>0$ would cut $V_l\cap\Di\subseteq A_\tau$, so $n_\tau\cdot(e_l-p^*)>0$.

\emph{Step 2 (monotone pieces).} If $\sigma<\tau<\sigma'$ and $A_\sigma\cap A_{\sigma'}\ne\varnothing$, concavity in $\tau$ gives $A_\sigma\cap A_{\sigma'}\subseteq A_\tau$; near $p^*$ this is an inclusion of wedges, so $n_\tau\in\operatorname{cone}(n_\sigma,n_{\sigma'})$. On $(0,\tau_j)$ all $A_\tau$ contain points near $e_j$, so the $n_\tau$ lie in one open half-circle and betweenness makes the angle monotone there; likewise on $(\tau_i,1)$.

\emph{Step 3 (continuity).} If the one-sided limits $n(\tau_0\pm)$ differed, then on the open wedge $W\coloneqq\{n(\tau_0+)\cdot\delta>0>n(\tau_0-)\cdot\delta\}$ and on $-W$ we would have $\Psi(p^*+\delta,\tau_0)=0$ for small $\delta$. The zero set is convex (Lemma~\ref{lem:levels}), so it would contain a neighbourhood of $p^*$, contradicting Step 1.

\emph{Step 4 (no constancy; crossing values).} If $n_\tau$ were constant on some $[a,b]$, with line $\lambda$, then $\Psi(q,\cdot)=0$ on $[a,b]$ for $q\in\lambda\cap\Di$, so $\Psi(q,\cdot)\le0$ on $\R$ (Lemma~\ref{lem:caps}(a)). Hence $\lambda\cap\Di\subseteq\{\Psi(\cdot,\tau)\le0\}$ for every $\tau$, which forces $\lambda^+_\tau=\lambda$ for all $\tau\in(0,1)$; by continuity $n_\tau$ would then be constant on $(0,1)$. But $n_\tau\cdot(e_i-p^*)$ is $>0$ where $h_i>0$ (Step 1) and $<0$ where $h_i<0$ (there $V_i\cap\Di\subseteq B_\tau$, which lies in the open negative side of $\lambda^+_\tau$). So the angle is non-constant on every interval, the two monotone pieces glue to one strictly monotone continuous angle on $(0,1)$, and continuity of $n$ and the sign change give $\lambda_{\tau_i}=\lin(p^*,e_i)$.

\emph{Step 5 (thin zero set).} If $\lambda^-_{\tau_0}\ne\lambda^+_{\tau_0}$, then $\lambda^-_{\tau_0}\not\ni p^*$, and a half-disc $D^-$ at $p^*$ on the negative side of $\lambda^+_{\tau_0}$ lies in $Z_{\tau_0}$. Pick $\tau_1\ne\tau_0$ near $\tau_0$ and $q\in A_{\tau_1}\cap D^-$. Then $\Psi(q,\tau_1)>0=\Psi(q,\tau_0)$, so by Lemma~\ref{lem:caps}(a) $q\in A_\tau$ for $\tau$ tending to $\tau_0$ from one side, and continuity of $n$ gives $n_{\tau_0}\cdot(q-p^*)\ge0$, contradicting $q\in D^-$.

\emph{Cluster lines.} A cluster line $\ell$ of $h(\tau)$ must meet $\Di$, since otherwise $\Psi(\cdot,\tau)$ would be one-signed on $\Di$. By Lemma~\ref{lem:chords}(b) its closed positive side misses $B_\tau$ and its closed negative side misses $A_\tau$, so $\ell\cap\Di\subseteq Z_\tau$ and $\ell=\lambda_\tau$ with the same orientation.
\end{proof}
\status{proved; refereed twice (\src{N18} \S1.5, \src{N20} \S4; CORRECT WITH REPAIR, repairs (R-a), (R-b), (R-c) incorporated). Source: \src{N15} Thm.~4.2. It excludes pinning of a vertex on any arc, non-unique clusters, chord jumps and fat zero sets.}

\begin{theorem}\label{thm:nosliding}
Let $r$ be the common corner ray of $F_k$ and $F_j$, parametrised so that $r=h(0)$ on both facets. Then the one-sided limits $n^k(0+)$ and $n^j(0+)$ of the chord normals coincide.
\end{theorem}
\begin{proof}
If not, $P(\cdot,r)$ is $\ge0$ and $\le0$ on two opposite open wedges at $p^*$, and by convexity of its zero set (Lemma~\ref{lem:levels}; $r\in U$ is mixed) it vanishes on a ball $B$ around $p^*$. By strict monotonicity on $F_k$ there is $q\in B$ with $n^k(0+)\cdot(q-p^*)<0$ and $q\in A^k_\tau$ for some $\tau\in(0,1)$. By continuity at $0+$, $q\notin A^k_\sigma$ for $\sigma\in(0,\eta)$, some $\eta>0$. So $\{\tau:\Psi^k(q,\tau)>0\}$ is an interval that contains $\tau$ and misses $(0,\eta)$; Lemma~\ref{lem:caps}(a) then gives $P(q,r)=\Psi^k(q,0)<0$, a contradiction.
\end{proof}
\status{proved; refereed twice (\src{N18} \S1.6, \src{N20} \S4). Source: \src{N15} Thm.~4.3. This excludes ``sliding edges''.}

\subsection{Magnitudes from rotation}

Fix a facet and orient $\tau$ so that the angle $\theta_\tau$ of $n_\tau$ is continuous and strictly increasing. Put $u_\tau\coloneqq$ the rotation of $n_\tau$ by $+\pi/2$ in $\one^\perp$, $D^+_\tau(x)\coloneqq\Psi(p^*+xn_\tau,\tau)$, $D^-_\tau(x)\coloneqq-\Psi(p^*-xn_\tau,\tau)$, and let $L_*(\tau_1)$, $L^*(\tau_1)$ be the lower and upper left Dini derivatives of $\theta$ at $\tau_1$. For $0<\xi<r_0/2$ the \emph{anchor} $q_\tau(\xi)\coloneqq p^*+\xi u_\tau$ lies on $\lambda_\tau$, in $A_\sigma$ for $\sigma>\tau$ near $\tau$ and in $B_\sigma$ for $\sigma<\tau$.

\begin{lemma}\label{lem:rotation}
Let $\tau_1\in(0,1)$, $0<\xi<r_0/2$, $S^\pm\coloneqq\partial^\pm_\tau\Psi(q_{\tau_1}(\xi),\tau_1)$. Then $0<S^+\le S^-<\infty$ and
\[
\frac{S^+}{L^*}\cdot\frac{r_0-\xi}{r_0\xi}\ \le\ \liminf_{x\downarrow0}\frac{D^+_{\tau_1}(x)}x\ \le\ \limsup_{x\downarrow0}\frac{D^+_{\tau_1}(x)}x\ \le\ \frac{S^-}{L_*}\cdot\frac{r_0+\xi}{r_0\xi},
\]
with $1/0=\infty$, $1/\infty=0$. The same holds for $D^-_{\tau_1}$ with $|\partial^-_\tau\Psi(q',\tau_1)|$ on the left and $|\partial^+_\tau\Psi(q',\tau_1)|$ on the right, $q'=p^*-\xi u_{\tau_1}$.
\end{lemma}
\begin{proof}[Proof sketch]
For $\tau_0<\tau_1$ near $\tau_1$ let $\Delta\theta\coloneqq\theta_{\tau_1}-\theta_{\tau_0}>0$ and $q\coloneqq q_{\tau_0}(\xi)$. Relative to $\lambda_{\tau_1}$, $q$ has transverse coordinate $\xi\sin\Delta\theta$ and tangential coordinate $\xi\cos\Delta\theta$, so Lemma~\ref{lem:sandwich} gives $D^+(\xi\sin\Delta\theta\,\tfrac{r_0}{r_0+\xi})\le\Psi(q,\tau_1)\le D^+(\xi\sin\Delta\theta\,\tfrac{r_0}{r_0-\xi})$. Since $\Psi(q,\cdot)$ vanishes at $\tau_0$ and is positive just to its right, Lemma~\ref{lem:caps}(b) gives $(\tau_1-\tau_0)\Psi(q,\tau_2)/(\tau_2-\tau_0)\le\Psi(q,\tau_1)\le(\tau_1-\tau_0)\,\partial^+_\tau\Psi(q,\tau_0)$ for $\tau_2>\tau_1$. As $\tau_0\uparrow\tau_1$, upper semicontinuity of $\partial^-$ (Lemma~\ref{lem:caps}(c)) bounds $\limsup\partial^+\Psi(q_{\tau_0},\tau_0)$ by $S^-$, and joint continuity followed by a supremum over $\tau_2\downarrow\tau_1$ gives $S^+$ on the left; $(\tau_1-\tau_0)/\sin\Delta\theta$ has limsup $1/L_*$ and liminf $1/L^*$, and every small $x$ is of the form $\xi\sin\Delta\theta\,r_0/(r_0\pm\xi)$. Positivity and finiteness of $S^\pm$ come from Lemma~\ref{lem:caps}(b), since $\Psi(q_{\tau_1},\cdot)$ changes sign strictly at $\tau_1$ (Theorem~\ref{thm:pivot}).
\end{proof}
\status{proved; refereed twice (\src{N18} \S1.7, \src{N20} \S5; both verified the semicontinuity directions analytically, noting that gauge-mixture numerics cannot test them). Source: \src{N15} Lemma 5.1.}

\begin{lemma}\label{lem:risingsun}
If $\theta$ is continuous and strictly increasing on $[a,b]$, then $L_*(\tau)>0$ for some $\tau\in(a,b]$.
\end{lemma}
\status{proved (elementary); refereed twice. Source: \src{N15} Lemma 5.2.}

\begin{proposition}\label{prop:regularity}
\begin{enumerate}[leftmargin=*,label=(\alph*)]
\item For every compact $J\subset(0,1)$ there are $C,t_0$ with $|\Psi(p^*+\delta,\tau)|\le C|\delta|$ for $\tau\in J$, $|\delta|\le t_0$.
\item $\theta'(\tau)$ exists and lies in $(0,\infty)$ for a.e.\ $\tau$; call this full-measure set $E'$.
\item For $\tau\in E'$ there is $c_\tau>0$ with $c_\tau x\le D^\pm_\tau(x)\le Cx$ for small $x$.
\item If $\lim_{x\downarrow0}D^\pm_\tau(x)/x$ exist, then $\lim_{t\downarrow0}\Psi(p^*+t\delta,\tau)/t$ exists for every $\delta$ and equals $\alpha[n_\tau\cdot\delta]_+-\beta[n_\tau\cdot\delta]_-$.
\end{enumerate}
\end{proposition}
\begin{proof}[Proof sketch]
(a) Lemma~\ref{lem:risingsun} gives points $a<b$ left of $J$ and $c$ right of $J$ with $L_*>0$; Lemma~\ref{lem:rotation} bounds $D^\pm$ linearly there, Lemma~\ref{lem:sandwich} extends the bound to $|\Psi(p^*+\delta,\sigma)|\le2C'|\delta|$ for $\sigma\in\{a,b,c\}$, and concavity in $\tau$ (extrapolation from the left, minimum of endpoints) transfers it to $J$. (b) Lebesgue's theorem for monotone $\theta$; $\theta'(\tau)=0$ would make the lower bound of Lemma~\ref{lem:rotation} infinite, contradicting (a). (c) Lemma~\ref{lem:rotation} with $L_*=L^*=\theta'$. (d) Lemma~\ref{lem:sandwich} at scale $t$.
\end{proof}
\status{proved; refereed twice (\src{N18} \S1.8, \src{N20} \S5.3). Source: \src{N15} Prop.~5.3. Part (b) also excludes singular (Cantor-type) chord maps. Oscillation of $D^\pm_\tau(x)/x$ inside $[c_\tau,C]$ is \emph{not} excluded locally; \S\ref{sec:proj} shows it is irrelevant.}

\subsection{Projectivity}\label{sec:proj}

\begin{lemma}\label{lem:twoslope}
Let $I$ be an open interval, $n:I\to\one^\perp$ continuous with $|n|=1$ and not constant on $I$, and $\alpha,\beta:I\to(0,\infty)$ continuous. Suppose that for every $\delta\in\one^\perp$ the function $f_\delta(\tau)\coloneqq\alpha(\tau)[n(\tau)\cdot\delta]_+-\beta(\tau)[n(\tau)\cdot\delta]_-$ is concave on $I$. Then $a\coloneqq\alpha n$ is affine on $I$ and $\beta/\alpha$ is constant.
\end{lemma}
\begin{proof}[Proof sketch]
On a small interval $J$ with $n(J)$ in an arc of length $<\pi/2$, let $C_J\coloneqq\{\delta:n(\tau)\cdot\delta>0\ \forall\tau\in J\}$. For $\delta\in C_J$, $f_\delta=a\cdot\delta$ is concave; for $\delta\in-C_J$, $f_\delta=b\cdot\delta$ with $b\coloneqq\beta n$ is concave. So $a,b$ are locally Lipschitz, and their distributional second derivatives are vector measures with $a''=-n\,m_a$, $b''=n\,m_b$, $m_a,m_b\ge0$ (using continuity of $n$). With $\times$ the determinant on $\one^\perp$, $d(a\times a')=a\times a''=0$ because $a\parallel n$; so $a\times a'\equiv L_a$ and likewise $b\times b'\equiv L_b$ are constant on $I$. Now $a\times a'=\alpha^2(n\times n')$ and $b\times b'=\beta^2(n\times n')$. If $L_a=0$ then $n'=0$ a.e.\ and $n$ is constant, which is excluded. So $(\beta/\alpha)^2=L_b/L_a$ is constant, $b=\rho a$, and $(m_b+\rho m_a)n=b''-\rho a''=0$ forces $m_a=0$: $a$ is affine.
\end{proof}
\status{proved; refereed once (\src{N14} \S2.6: CORRECT; regularity is derived from concavity and atoms of $a''$ are handled); its use in \src{N15} checked by both \src{N18} and \src{N20}. Source: \src{N12} Lemma 4.3.}

\begin{proposition}\label{prop:blowup}
Fix a facet and a compact interval $J_1\subset(0,1)$. There are $t_m\downarrow0$ and continuous $\alpha,\beta:J_1\to(0,\infty)$ such that, with $\Psi_t(\delta,\tau)\coloneqq\Psi(p^*+t\delta,\tau)/t$, the limits
\[
G(\delta,\tau)\coloneqq\lim_m\Psi_{t_m}(\delta,\tau),\qquad G_2(\delta,\tau)\coloneqq\lim_{\epsilon\downarrow0}G(\epsilon\delta,\tau)/\epsilon
\]
exist for all $(\delta,\tau)\in\one^\perp\times J_1$, $G_2(\delta,\tau)=\alpha_\tau[n_\tau\cdot\delta]_+-\beta_\tau[n_\tau\cdot\delta]_-$, and $\tau\mapsto G_2(\delta,\tau)$ is concave for every $\delta$.
\end{proposition}
\begin{proof}[Proof sketch]
Take $J_1\Subset J\Subset J'\Subset(0,1)$. By Proposition~\ref{prop:regularity}(a) and concavity, $\Psi_t(\delta,\cdot)$ is $L|\delta|$-Lipschitz on $J$ uniformly in $t$; and (Lemma~\ref{lem:sandwich} at scale $t$) $\Psi_t(\delta,\tau)\le\Psi_t(\delta',\tau)$ when $n_\tau\cdot\delta<n_\tau\cdot\delta'$ and $t$ is small (property (M)).

\emph{First blow-up.} Arzel\`a--Ascoli on a countable dense set of $\delta$ and a diagonal argument give $t_m$ and a limit $G$, concave and Lipschitz in $\tau$, that inherits (M).

\emph{Sweep.} Let $\tau_0\in E'\cap\operatorname{int}J$, $\vartheta\coloneqq\theta'(\tau_0)\in(0,\infty)$. Suppose the profile $x\mapsto G(\cdot,\tau_0)$ along $n_{\tau_0}$ had a jump $\mathcal J>0$ at $x_0$. Take a far point $\delta\approx Yu_{\tau_0}+x_0n_{\tau_0}$. Rotating the chord through an angle $\pm3\epsilon/Y$, i.e.\ moving $\tau$ by about $\pm3\epsilon/(Y\vartheta)$, moves $n_\sigma\cdot\delta$ across the jump. The Lipschitz bound in $\tau$ then gives $\mathcal J-o(1)\le L|\delta|\cdot6\epsilon/(Y\vartheta)\approx6L\epsilon/\vartheta$, so $\mathcal J=0$. Hence $G(\delta,\tau_0)=\varphi_{\tau_0}(n_{\tau_0}\cdot\delta)$ with $\varphi_{\tau_0}$ continuous and nondecreasing; (M) and equi-Lipschitz convergence on the dense set $E'\cap J$ extend the convergence to all $(\delta,\tau)$.

\emph{Window.} With $\delta_Y\coloneqq Yu_{\tau_0}+x_0n_{\tau_0}$, the concave functions $s\mapsto G(\delta_Y,\tau_0+s/(Y\vartheta))$, $|s|\le r$, converge uniformly to $s\mapsto\varphi_{\tau_0}(x_0+s)$: a single far point sees every transverse position as the chord turns through an angle of order $1/Y$. So $\varphi_{\tau_0}$ is concave on $\R$.

\emph{Second zoom.} Concavity and $\varphi(0)=0$ give $\alpha_{\tau_0}=\varphi'(0+)$ and $\beta_{\tau_0}=\varphi'(0-)$, with $c_{\tau_0}\le\alpha\le\beta\le C$ (Proposition~\ref{prop:regularity}(c)); equi-Lipschitz convergence of $G(\epsilon\delta,\cdot)/\epsilon$ on $E'\cap J_1$ extends the two-slope form to all of $J_1$ with $\alpha,\beta$ continuous, and $\alpha>0$ everywhere because a concave function $\ge0$ vanishing at an interior point vanishes identically.
\end{proof}
\status{proved; refereed twice (\src{N18} \S1.9, \src{N20} \S6; cosmetic repair $\tau_0\in\operatorname{int}J$). Source: \src{N15} \S6 Steps 1--4.}

Say that \textbf{(PC)} holds if for each facet $F_k$ there is a linear map $\Gamma_k:v_k^\perp\to\one^\perp$ with $\Gamma_kh\ne0$ on $\relint F_k$, such that every cluster line of $h\in\relint F_k$ is $\{q:\Gamma_kh\cdot(q-p^*)=0\}$ with positive side $\{\Gamma_kh\cdot(q-p^*)>0\}$. Say that \textbf{(NS)} holds if $\Gamma_kr\propto_+\Gamma_jr$ at each corner ray $r=F_k\cap F_j$.

\begin{corollary}\label{cor:PCNS}
(PC) and (NS) hold, and each $\Gamma_k$ has rank $2$.
\end{corollary}
\begin{proof}
By Proposition~\ref{prop:blowup} and Lemma~\ref{lem:twoslope} ($n$ is nowhere locally constant by Theorem~\ref{thm:pivot}), $a=\alpha n$ is affine on each compact $J_1\subset(0,1)$. Positively proportional affine maps with non-constant direction differ by a constant factor (\src{N20} \S6.5), so the normalisations patch and $a$ extends affinely to $[0,1]$; define $\Gamma_kh(\tau)\coloneqq a(\tau)$ and extend linearly to $v_k^\perp=\operatorname{span}\{r,r'\}$. The direction field $n_\tau$ is the zero-chord field and does not depend on the subsequence. With Theorem~\ref{thm:pivot} this is (PC). If $a(0)$ and $a(1)$ were parallel the direction would be constant, and if antiparallel $a$ would vanish inside, so the rank is $2$ and $\Gamma_kr=a(0)\ne0$. Finally $\Gamma_kr\propto_+n^k(0+)=n^j(0+)\propto_+\Gamma_jr$ by Theorem~\ref{thm:nosliding}, which is (NS).
\end{proof}
\status{proved; refereed twice (\src{N18} \S1.9--1.10, \src{N20} \S\S6.5, 7). Source: \src{N15} \S6 Step 5 and Thm.~7.1.}

\subsection{The contradiction}\label{sec:contradiction}

\begin{lemma}\label{lem:edgebox}
Let $a\ne b$. For a.e.\ pair $(A,B)\in U^2$ with $A_a<0<A_b$, $B_b<0<B_a$ and $A_bB_a<|A_a|\,|B_b|$, the closed edge $[e_a,e_b]$ does not meet $H(A)\cap H(B)$.
\end{lemma}
\begin{proof}
The ratio condition gives an open interval of $\lambda>0$ on which $(A+\lambda B)_a<0$ and $(A+\lambda B)_b<0$; by Fubini, for a.e.\ $(A,B)$ some such $A+\lambda B$ is good. Suppose $q_0\in[e_a,e_b]\cap H(A)\cap H(B)$. The open set $O\coloneqq H(A)\cap H(B)\cap\Di$ accumulates at $q_0$, and on $O$, $P(\cdot,A+\lambda B)\ge P(\cdot,A)+\lambda P(\cdot,B)>0$ a.e. If $A+\lambda B\notin U$ this contradicts $P\le P_{K'}\le0$; otherwise $O\subseteq H(A+\lambda B)$ and $q_0\in\cl H(A+\lambda B)$. But $c(A+\lambda B)$ is negative in coordinates $a$ and $b$, hence negative on the whole closed edge.
\end{proof}
\status{proved; refereed once (\src{N14} \S2.3: CORRECT, with an explicit sanity test). Source: \src{N12} Thm.~2.4. It needs no regularity and no coherence at edge priors.}

Under (PC) put $c_k(h)\in\R^3$ for the vector with $c_k(h)\cdot p^*=0$ whose restriction to $\one^\perp$ is $\Gamma_kh$; then the cluster line of $h$ is $\{c_k(h)\cdot q=0\}$, a line in all of $\aff\Delta$.

\begin{proposition}\label{prop:boxbridge}
Assume (PC), (NS) and the conclusions of Theorem~\ref{thm:pivot}.
\begin{enumerate}[leftmargin=*,label=(\alph*)]
\item \emph{Crossing rows.} For each crossing coordinate $i$ of $F_k$ there is $\lambda_{ik}>0$ with $c_k(h)_i=\lambda_{ik}h_i$ for all $h\in v_k^\perp$.
\item \emph{Corners.} $c_k(r)\propto_+c_j(r)$ at each corner ray $r=F_k\cap F_j$.
\item \emph{Box.} Let $A\in\relint F_k$, $B\in\relint F_{k'}$ and $q_0\in[e_a,e_b]$ satisfy $A_a<0<A_b$, $B_b<0<B_a$, $A_bB_a<|A_a||B_b|$. Then $c_k(A)\cdot q_0\le0$ or $c_{k'}(B)\cdot q_0\le0$.
\end{enumerate}
\end{proposition}
\begin{proof}
(a) At the crossing $w_{ik}$ the unique cluster line is $\lin(p^*,e_i)$, so $c_k(w_{ik})\cdot e_i=0$; the functional $h\mapsto c_k(h)_i$ on the plane $v_k^\perp$ vanishes on $\R w_{ik}=v_k^\perp\cap e_i^\perp$, hence equals $\lambda_{ik}h_i$, and $\lambda_{ik}>0$ because $c_k(h)\cdot e_i=\Gamma_kh\cdot(e_i-p^*)>0$ when $h_i>0$ (Theorem~\ref{thm:pivot}). (b) is (NS). (c) Suppose $c_k(A)\cdot q_0>0$ and $c_{k'}(B)\cdot q_0>0$. By uniqueness of cluster lines, the normalised vectors $c(A')$ of good $A'\to A$ converge to the normalised $c_k(A)$ (an affine function on $\aff\Delta$ is the restriction of a unique linear functional on $\R^3$), so $q_0\in H(A')$ for a.e.\ $A'$ in a small ball around $A$; likewise for $B$. The sign and ratio conditions are open. This produces a set of pairs of positive measure contradicting Lemma~\ref{lem:edgebox}.
\end{proof}
\status{proved; refereed once (\src{N14} \S2.5: CORRECT; nondegeneracy at the $b$-crossings from \src{R18} \S5). Source: \src{N12} Thm.~3.1 items 3--6. Its use with the conclusions of \S4.3 was checked in \src{N18} and \src{N20}.}

\paragraph{The contamination triangle.}
Take $T=T_1$ and $p^*=u$. On $F_k$ the crossing coordinates are the two others, and $c_k(h)\cdot u=0$ gives $c_k(h)_k=-\sum_{i\ne k}\lambda_{ik}h_i$. Normalise $\lambda_{31}=\lambda_{12}=\lambda_{23}=1$ and put $x_1\coloneqq\lambda_{21}$, $x_2\coloneqq\lambda_{32}$, $x_3\coloneqq\lambda_{13}$. Then
\begin{gather*}
c_1(h)\propto(-(x_1h_2+h_3),\ x_1h_2,\ h_3),\qquad
c_2(h)\propto(h_1,\ -(h_1+x_2h_3),\ x_2h_3),\\
c_3(h)\propto(x_3h_1,\ h_2,\ -(x_3h_1+h_2)).
\end{gather*}

\begin{proposition}\label{prop:T1}
For $T_1$, (PC), (NS) and the box of Proposition~\ref{prop:boxbridge}(c) are inconsistent, for every $\varepsilon\in(0,1)$.
\end{proposition}
\begin{proof}
\emph{Corner equations.} At $r_{12}=(-1,-1,\kappa+1)$: $c_1(r)=(x_1-\kappa-1,\,-x_1,\,\kappa+1)$ and $c_2(r)=(-1,\,1-x_2(\kappa+1),\,x_2(\kappa+1))$. These are positively proportional iff $c_1(r)=x_2^{-1}c_2(r)$, i.e.\ iff $x_1+1/x_2=\kappa+1$. Cyclically,
\begin{equation}\label{eq:corner}
x_1+\frac1{x_2}=x_2+\frac1{x_3}=x_3+\frac1{x_1}=\kappa+1 .
\end{equation}
With $\phi(t)\coloneqq\kappa+1-1/t$ this says $x_1=\phi(x_2)$, $x_2=\phi(x_3)$, $x_3=\phi(x_1)$, so $x_1$ is a fixed point of $\phi^3$. The M\"obius map $\phi$ has matrix $\bigl(\begin{smallmatrix}\kappa+1&-1\\1&0\end{smallmatrix}\bigr)$ with determinant $1$ and trace $\kappa+1>2$, so it is hyperbolic and $\phi^3$ has the same two fixed points $x_\pm$, the roots of $x^2-(\kappa+1)x+1=0$. Hence
\[
x_1=x_2=x_3=x\in\{x_+,x_-\},\qquad x+\frac1x=\kappa+1,\qquad x_+>1,\ x_-=1/x_+ .
\]

\emph{Box.} Let $\beta\coloneqq1/\kappa$, $q_0\coloneqq(s,1-s,0)$, and for small $\eta,\zeta>0$
\[
A=(-\beta(1+\eta),\,\eta,\,1)\in\relint F_1,\qquad B=(1,\,-\zeta,\,-\beta(1-\zeta))\in\relint F_3 ,
\]
near the two crossings of coordinate $2$ (membership uses $c_0(1-\beta)=\varepsilon\beta$). Then $c_1(A)\cdot q_0=x\eta(1-2s)-s$ and $c_3(B)\cdot q_0=sx-(1-s)\zeta$, and the ratio condition for $(a,b)=(1,2)$ is $\eta<\beta\zeta(1+\eta)$. Take $\zeta\coloneqq\theta sx/(1-s)$ and $\eta\coloneqq s/(\theta x(1-2s))$ with $\theta<1$. Both chart values are positive, and the ratio condition reads $(1-s)/((1-2s)(1+\eta))<\theta^2\beta x^2$. If $x^2>\kappa$, this holds for $\theta$ close to $1$ and $s$ small, contradicting Proposition~\ref{prop:boxbridge}(c). So $x^2\le\kappa$.

The transposition $\sigma$ of coordinates $1,2$ maps $T_1$ to itself and fixes $u$; $c^\sigma(g)\coloneqq\sigma c(\sigma g)$ is again a realisation, so the whole chain applies to it, and its facet data are $x^\sigma=1/x$. Hence also $x^{-2}\le\kappa$, i.e.
\[
\frac1\kappa\le x^2\le\kappa .
\]

\emph{Contradiction.} Since $t\mapsto t+1/t$ is increasing on $[1,\infty)$ and invariant under $t\mapsto1/t$, the box gives $x+1/x\le\sqrt\kappa+1/\sqrt\kappa$. With $y=\sqrt\kappa>1$,
\[
(\kappa+1)-\Bigl(y+\frac1y\Bigr)=\frac{(y-1)(y^2+1)}{y}>0,
\]
which contradicts $x+1/x=\kappa+1$. Concretely, $x_+^2>\kappa$ violates the upper bound and $x_-^2<1/\kappa$ the lower one; at $\varepsilon=0.3$, $x_+^2/\kappa\approx3.80$.
\end{proof}
\status{proved. The corner system is \src{R15} Prop.~5.2 and \src{R4} Prop.~G1 (refereed in \src{R18} and \src{R8}); the box value $1/\kappa\le x^2\le\kappa$ and the certificate are \src{R15} \S6 (refereed in \src{R18} \S6, which recomputed them by hand). \src{N20} \S7 confirms that this contradiction holds for every $\kappa>1$, i.e.\ every $\varepsilon\in(0,1)$. The transposition step replaces the explicit mirror certificate near $e_1$ used in \src{R15}; this is a change of presentation only. The corner equations alone are consistent (two sign-admissible solutions $x_\pm$), so the box is indispensable. The same equations, with the same contradiction, appear for continuous maps in \src{CONTAMINATION\_NOGO\_SCOPE\_AND\_GEOMETRY.md} \S2.}

\paragraph{General all-facing triangles.}
For an all-facing triangle in general position the same three ingredients give a contradiction by the algebra of \src{R15} Thm.~6.2. Work at $p^*=u$ with $\mu_{ik}=\lambda_{ik}/3$, write $v_{kl}$ for the $l$-th coordinate of $v_k$ and $\rho_k(i{:}j)\coloneqq\mu_{ik}/\mu_{jk}$. At the corner $r=r_{ac}$ with third index $b$, Lemma~\ref{lem:signs} gives $r_a,r_c<0<r_b$; put $A_r\coloneqq-r_c/r_b$, $B_r\coloneqq-r_a/r_b$, $s\coloneqq\rho_a(b{:}c)$, $t\coloneqq\rho_c(a{:}b)$. Then (NS) reads $A_r/s+B_rt=1$, and the boxes for the pairs $(a,b)$ and $(c,b)$ read $v_{ca}/v_{cc}\le st\le v_{aa}/v_{ac}$. Together with $r\perp v_a,v_c$ these give
\[
\frac{v_{cb}}{v_{cc}}\le\rho_a(b{:}c)\le\frac{v_{ab}}{v_{ac}},\qquad\frac{v_{ca}}{v_{cb}}\le\rho_c(a{:}b)\le\frac{v_{aa}}{v_{ab}} .
\]
Applying the second inequality at the other corner of facet $a$ pinches $\rho_a(b{:}c)=v_{ab}/v_{ac}$ on every facet. The corner equation then gives $t=v_{aa}/v_{ab}$, while pinching on facet $c$ gives $t=v_{ca}/v_{cb}$. So $v_{aa}v_{cb}=v_{ab}v_{ca}$, i.e.\ $\det(v_a,v_c,e_c)=0$, and $e_c$ lies on the line through $v_a,v_c$, contradicting general position.
\status{proved; refereed once (\src{R18} \S6: every step recomputed). In \src{R15} the box came from a forbidden-support lemma under extra hypotheses; here it comes from Proposition~\ref{prop:boxbridge}(c), which uses only chord lines (\src{N12} Thm.~3.1, \src{N14} \S2.5).}

\begin{proof}[Proof of Theorem~\ref{thm:main}]
Reduce to $p^*=u$ (Proposition~\ref{prop:reduction}(d)). Theorem~\ref{thm:pivot}, Theorem~\ref{thm:nosliding} and Corollary~\ref{cor:PCNS} give unique pivot chords, (PC) and (NS). Proposition~\ref{prop:boxbridge} turns these into crossing rows, corner equations and the box. Proposition~\ref{prop:T1} (for $T_1$) and the general pinching argument give a contradiction.
\end{proof}

\section{General decomposable domains}\label{sec:general}

Call a decomposable domain \emph{two-sided} if, near the target, evaluable models may accept $f$, accept $-f$, or give no advice: a.e.\ pair with an orientation near $\partial D_T$ has an advice range containing $\{+,0,-\}$. Without non-adjacency, pairs near $\partial D_T$ whose advice range is $\{+,0,-\}$ may have wedge charts $\{c\cdot q>0,\,d\cdot q>0\}$. Non-adjacency is used in exactly two places in \S\ref{sec:main} (\src{N21} \S1, confirmed in \src{N23} \S1): (U-glob) global straight levels and Lemma~\ref{lem:PL2} on a set that leaves every ball around $p^*$, inside Lemma~\ref{lem:nopin}; and (U-box) the box step, where the witnesses' actual charts must contain the edge prior. Everything else is local at $p^*$ or adjacency-tolerant. Throughout this section the ranking is strictly guiding additive on a decomposable domain with (LF) at $T$, and (BU) holds at every interior prior with $K_{p^*}=T$ (R4 notation: for a pair with advice range $\{+,0,-\}$, $R(g)=\{\alpha>0,d>0\}$, $R(-g)=\{\beta>0,d<0\}$ with $\alpha\coloneqq c\cdot q$, $\beta\coloneqq c'\cdot q$).

\paragraph{The neighbourhood $N_0$ (compact collar).} (LF) supplies only an open $N$, which need not contain a conic neighbourhood of $\partial D_T\setminus0$. Every use of a neighbourhood in this section concerns gambles near one fixed compact set of boundary gambles at bounded norm (\src{N27} \S1.1, inventory re-derived independently in \src{N27\_referee} \S1). Let $\Sigma\subseteq\partial D_T\setminus0$ be the compact set formed by the three facet segments $h([0,1])$ between the normalised corner rays, the crossings in every normalisation used, and the limits of the box witnesses of Proposition~\ref{prop:boxbridge}(c) in the normalisation of Proposition~\ref{prop:T1} (for $T_1$: $(-\beta,0,1)$ and $(1,0,-\beta)$) and of the test gambles of Theorem~\ref{thm:cuts}. Fix a compact $K''\subseteq\operatorname{int}T$ and $\eta>0$ such that the compact collar
\[
\Nb\coloneqq\{g:\operatorname{dist}(g,\Sigma)\le\eta\}
\]
lies in $N$, has $P_{K''}>0$ on it, misses $-\cl D_T$ and contains no $h\le0$; this is possible because $\Sigma$ is compact, $N$ is open and $P_{K''}>0$ on $\Sigma$. Below, $N_0$ means $\Nb$. Two numerical conditions are imposed:
\begin{enumerate}[leftmargin=*,label=(N\arabic*)]
\item level shifts are restricted to $|c|\le\bar c$ with $\bar c<\eta/\sqrt3$, so that $h(\tau)-c\one\in\operatorname{int}\Nb$ for $\tau\in[0,1]$;
\item the radius $r_1$ of Lemma~\ref{lem:B1} is chosen after $\bar c$, so small that $|\Psi_k(q,\tau)|<\bar c$ on $B(p^*,r_1)\times[0,1]$ for each facet (possible by joint continuity, since $\Psi_k(p^*,\cdot)=0$ on $[0,1]$).
\end{enumerate}
With these choices every statement of this section holds with the same proof (\src{N27} Repair 27.1; \src{N27\_referee} \S1.3: CONFIRMED, constants re-derived; \src{N28} Lemma 0 gives the same repair and was CONFIRMED WITH REPAIR, the repair being exactly (N1)--(N2)). The balls around crossings, the witness balls and the test-gamble neighbourhoods are taken inside $\operatorname{int}\Nb$, which is possible because their centres lie in, or converge to, $\Sigma$. For a.e.\ gamble $h$ of the open cone $\R_+N$, accepting $h$ lies in the advice range (\src{N28} Lemma 0(b), confirmed). The $U$-domain and Corollary~\ref{cor:coverage} are unaffected: they use the explicit cones $U$ and $D_{K'}$, never $N$ (\src{N27} \S1.5, confirmed).

\begin{lemma}\label{lem:raytransfer}
Let an additive ranking (strictly guiding, or with permitted ties as in \S\ref{sec:open-ties}) on a decomposable domain have, for every $q\in\Di$, accepted set $B_q\coloneqq\{h:q\in R(h)\}$ equal a.e.\ to $D_{K_q}\cup\{h\ge0\}$. Then:
\begin{enumerate}[leftmargin=*,label=(\alph*)]
\item for a.e.\ $(g,t)\in\R^3\times(0,\infty)$, $R(tg)\cap\Di=R(g)\cap\Di$;
\item under strict guidance, for a.e.\ mixed $g$ and a.e.\ $t>0$, the advice ranges at $g$ and $tg$ contain the same acceptances; they can differ only in whether they contain $\times$, and whether they contain no advice ($0$) at pairs with advice range $\{+,0,-\}$ whose no-advice cell $\Di\setminus(R(g)\cup R(-g))$ has empty interior;
\item for a.e.\ $g$ in the cone over $\operatorname{int}\Nb$ and every open convex $C\subseteq\{q:K_q\cap K''\ne\varnothing\}$: $R(-g)\cap C=\varnothing$, no d-face of $g$ meets $C$, and $R(g)\cap C=\{c(g)\cdot q>0\}\cap C$.
\end{enumerate}
\end{lemma}
\begin{proof}[Proof sketch]
(a) For fixed $(q,t)$ the $g$-section of $\{(q,g,t):1[q\in R(g)]\ne1[q\in R(tg)]\}$ lies in $(B_q\sd D')\cup t^{-1}(B_q\sd D')$ with $D'\coloneqq D_{K_q}\cup\{h\ge0\}$ a cone, hence is null; Fubini, and two relatively open convex sets that agree a.e.\ are equal. (b) An acceptance in the advice range of a mixed pair has a nonempty chart, an acceptance outside it an empty chart. (c) Apply (a) to $g/t\in\Nb$ and Lemma~\ref{lem:B1}; a d-face point would have nearby points of $C$ in $R(-g)$.
\end{proof}
\status{proved (\src{N27} Lemma 27.2, Cor.~27.3, Cor.~27.3$'$); confirmed by one referee (\src{N27\_referee} \S2, which also supplied the full $c$-form in (c)). No guidance is used in (a). Only the charts transfer along rays, not the penalty vectors.}

\begin{lemma}\label{lem:quasiconcave}
For every $g\in\R^3$, $q\mapsto P(q,g)$ is quasi-concave on $\Di$: the sets $\{P(\cdot,g)>c\}$ and $\{P(\cdot,g)\ge c\}$ are convex for every $c\in\R$. Consequently, if $A_\tau\coloneqq\{\Psi(\cdot,\tau)>0\}$ satisfies $A_\tau\cap B=H^+\cap B$ for a ball $B$ at $p^*$ and an open half-plane $H^+$ whose boundary passes through $p^*$, then $A_\tau\subseteq H^+$.
\end{lemma}
\status{proved; refereed once (\src{N23} \S2); second reading CONFIRMED (\src{TS\_second\_referee} \S1.1). Source: \src{N21} Lemma N21.1 and Cor.~N21.1$'$. It uses only convexity of every chart in the table of charts for each advice range (\src{R1}) and the chord property of good gambles, and holds for every $g$, including nonnegative gambles and pairs on which every evaluable model gives the same advice (\src{N28\_referee} \S2). Non-adjacency is sufficient, not necessary, for quasi-convexity (\src{N23} counterexample: the precise identity system with adjacent data).}

\begin{lemma}\label{lem:B1}
There is $r_1>0$, satisfying (N2), with $K_q\supseteq K''$ for $q\in B_1\coloneqq B(p^*,r_1)$. For a.e.\ $g\in N_0$ and every $q$ with $K_q\cap K''\ne\varnothing$: $q\notin R(-g)$, $q\notin F(g)$, and $q\in R(g)$ iff $c(g)\cdot q>0$. So near-facet gambles have half-plane charts on $B_1$.
\end{lemma}
\status{proved; refereed once (\src{N23} \S3, which strengthened it to a pointwise statement); second reading of the pointwise form CONFIRMED (\src{TS\_second\_referee} \S1.1). Source: \src{N21} Lemma N21.2.}

Say that \textbf{(TS)} holds if for every facet $F_k$ and every $\tau\in(0,1)$, $p^*\in\cl A_\tau\cap\cl B_\tau$.

\begin{theorem}\label{thm:TS}
On every decomposable domain with (LF), under (BU), (TP) and (TS), the conclusions of Theorems~\ref{thm:pivot} and~\ref{thm:nosliding}, Lemma~\ref{lem:rotation}, Propositions~\ref{prop:regularity} and~\ref{prop:blowup} hold with $\Di$ replaced by $B_1$ in every local statement and $r_0$ by $r_1$. In particular (NS) holds, and (PC) holds for $h$ near $\Sigma$, extended to the cone over $\Sigma$ by linearity of $\Gamma_k$.
\end{theorem}
\status{proved; refereed once (\src{N23} \S4: CORRECT, two wording repairs); second reading (\src{TS\_second\_referee} \S1.1): items 1--5 of \src{N21} Thm.~N21.4 re-derived, including the cluster clause for every good sequence, items 6--8 spot-checked, CONFIRMED. Source: \src{N21} Thm.~N21.4. The one new argument is the negative vertex sign at crossings via clusters of $c$-vectors on $B_1$. The wording of (PC) follows \src{N27\_referee} \S1.4: the cluster clause is proved only near $\Sigma$, which is all that is used downstream; the full conic clause also follows from Lemma~\ref{lem:raytransfer}.}

Fix a pair $(a,b)$ with a failing box, i.e.\ with $\Lambda_{ab}>1$ for the facet data of \S\ref{sec:contradiction}; here $\Lambda_{ab}$ is \src{R15}'s vertex-limit box ratio, and witnesses near the two $b$-crossings contradict the box of Proposition~\ref{prop:boxbridge}(c) on $[e_a,e_b]$ near $e_b$ when $\Lambda_{ab}>1$. For $T_1$ every world $b$ has exactly one such $a$, since $\Lambda_{ab}\Lambda_{cb}=\kappa^{-2}<1$ and $\Lambda\in\{x_+^2/\kappa,\,1/(\kappa x_+^2)\}$ (exact; \src{N23} \S6).

\begin{theorem}\label{thm:cuts}
Assume (PC), (NS) and a realisation of $T_1$ on a decomposable domain with (LF). For every $b$ there are good gambles $g_m$ with advice range $\{+,0,-\}$ near one of the two $b$-crossings $w^*\in F_{k^*}$, with d-faces $F(g_m)\cap\Di=(a_m,z_m)$, such that $a_m\to a^*\in[e_b,p^*]\setminus B_1$, $z_m\to z^*\in[e_b,e_a)$, $K_q\subseteq\lin(e_b,v_{k^*})$ for $q\in[a^*,z^*]\cap\Di$, and some interior $q_m\to z^*$ have $K_{q_m}\to\{e_b\}$ in the Hausdorff metric.
\end{theorem}
\status{proved; refereed once (\src{N23} \S6); second reading CONFIRMED (\src{TS\_second\_referee} \S3.1), with the cosmetic repair that the hypotheses ``(PC), (NS)'' include the cluster clause of (PC) and Lemma~\ref{lem:B1} ($N_0$, $K''$, $r_1$), which hold wherever the theorem is used. Its Step~1 uses the chart-form edge box (Lemma~\ref{lem:edgebox}). Source: \src{N21} Thm.~N21.6. Theorem~\ref{thm:apex} below sharpens the apex limit to $a^*=e_b$.}

\begin{corollary}\label{cor:availability}
Under the hypotheses of Theorem~\ref{thm:cuts}: accepting $h$ lies in the advice range for a.e.\ $h$ with $h_b>0$ that is not $\ge0$; the advice range of a.e.\ mixed pair contains both acceptances; $K_{e_b}=\{e_b\}$ for every $b$; and for every edge prior $q_0$ with support $S$, every limit of $K_q$ as $q\to q_0$ through $\Di$ lies in $\Delta_S$.
\end{corollary}
\status{proved; refereed once (\src{N23} \S6, wording repair); second reading CONFIRMED (\src{TS\_second\_referee} \S2). Source: \src{N21} Cor.~N21.7.}

\begin{lemma}\label{lem:R231}
Let $F_k$ have a crossing $\tau_c\in(\sigma_1,\sigma_2)$ and $0<\rho\le r_1$. If $\Psi(\cdot,\sigma_1)\ge0$ and $\Psi(\cdot,\sigma_2)\ge0$ on $B(p^*,\rho)$, then $v_k\in K_q$, $\Psi(q,\cdot)\equiv0$ on $[\sigma_1,\sigma_2]$ and $\Psi(q,\cdot)\le0$ on $[0,1]$ for every $q\in B(p^*,\rho)$.
\end{lemma}
\begin{proof}[Proof sketch]
As in Lemma~\ref{lem:nopin}, $B(p^*,\rho)\subseteq\Gamma_{v_k}$; Lemma~\ref{lem:B1} rules out d-faces there; Lemma~\ref{lem:PL2} applies to the ball; and now the $\alpha$-point is $p^*$ itself, since $K_{p^*}=T\ni v_k$. So the whole ball lies in $\alpha_{v_k}$.
\end{proof}
\status{proved in a referee report (\src{N23} \S7.1, Lemma R23.1). Second reading: CORRECT (recorded in \src{CHECK\_REPORT.md}); every step re-derived, including the use of Lemma~\ref{lem:PL2} on a ball that contains $p^*$ and the choice $\mathcal U\subseteq\operatorname{int}N_0$ needed for Lemma~\ref{lem:B1}.}

\begin{corollary}\label{cor:R232}
Failure of the $B$-half of (TS) alone is impossible: on each facet, if $p^*\in\cl A_\tau$ for every $\tau\in(0,1)$, then $p^*\in\cl B_\tau$ for every $\tau$. Hence (TS) is equivalent to its $A$-half (TS$_A$).
\end{corollary}
\status{proved in a referee report (\src{N23} \S7.2, Cor.~R23.2). Second reading: CORRECT (recorded in \src{CHECK\_REPORT.md}); in particular items 1--4 of \src{N21} Thm.~N21.4 use only the $A$-half of (TS). It refutes the residual case (E5b) left open in \src{N21}.}

\begin{theorem}\label{thm:dichotomy}
Let a strictly guiding additive ranking on a decomposable domain with (LF) cover $T_1$ (integral penalties, or bounded charges given interior minimisers). Then either
\begin{enumerate}[leftmargin=*,label=(\Roman*)]
\item (TS) holds, and then (PC), (NS) hold, the advice range of a.e.\ mixed pair contains both acceptances, all vertex models are precise, interior limits at edge priors lie in the edge face, and every vertex carries the cut structure of Theorem~\ref{thm:cuts}; or
\item (TS) fails, and then some facet $F_k$ is pinned: there are an open convex $G$ with $p^*\in\cl G$, on which $v_k\in K_q$ and $\Psi(q,\cdot)\equiv0$ on an interval around a crossing; every segment from $G$ towards a vertex neighbourhood $V_l$ contains an interior $Q$ with $K_Q=\{v_k\}$; and accepting $h$ lies in the advice range for a.e.\ $h$ with $v_k\cdot h>0$.
\end{enumerate}
\end{theorem}
\status{proved; refereed once (\src{N23} \S8, assembling \src{N21} Thm.~N21.8 and Prop.~N21.9, both refereed). Branch (II) as stated also needs Corollary~\ref{cor:R232}, which excludes a failure of the $B$-half alone; that corollary has now had a second reading (CORRECT). The further refinement $A_{\tau_c}\cap B_1=\varnothing$ at the crossing (\src{N23} Prop.~R23.3) is not used in the statement and remains a one-reader result. Branch (II) never occurs (Theorem~\ref{thm:TSalways}).}

\begin{corollary}\label{cor:R234}
In every covering of $T_1$ as in Theorem~\ref{thm:dichotomy}, the advice range of a.e.\ mixed pair contains at least one acceptance. Hence no such covering exists on a decomposable domain with (LF) on which, for a non-null set of mixed pairs, the advice range contains neither acceptance. If (TS) holds and, for a non-null set of mixed pairs, the advice range lacks one acceptance, there is no covering either.
\end{corollary}
\status{the first two sentences: proved in a referee report (\src{N23} \S7.5, Cor.~R23.4); second reading (\src{CHECK\_REPORT.md}): CORRECT, from Corollary~\ref{cor:availability} if (TS) holds and from Corollary~\ref{cor:R232} with \src{N21} Prop.~N21.9(ii)--(iii) otherwise. The last sentence: \src{N21} Thm.~N21.8(b), refereed once (\src{N23}); by Corollary~\ref{cor:TSalways}(b) it no longer needs the hypothesis (TS).}

\begin{proposition}\label{prop:noWLOG}
Restricting the domain gives no reduction to the non-adjacent case: removing ``$-$'' from the advice range of a non-null set of pairs (that is, removing from the class every model that accepts $-g$ there) removes the vertex selections from the class of evaluable models, so \src{R11} Lemma RS does not apply, and removing it while keeping the penalties changes the Bayes family and can destroy coherence (an explicit two-world example exists).
\end{proposition}
\status{proved; refereed once (\src{N23} \S5, scope repair: what is shown is that the restriction route fails, not that no reduction exists). Source: \src{N21} Prop.~N21.5.}

\subsection{Competing acceptances: no pinning, and the cuts reach the vertices}\label{sec:general-N28}

In this subsection a \emph{cover} means a strictly guiding additive ranking (integral penalties, or bounded charges vanishing on null sets with interior minimisers) on a decomposable domain with (LF) at $D_T$, satisfying (LA), (S), (C) and covering $T$; $K_q$ are its interior Bayes credal sets and $K_{p^*}=T$. For a facet $F_k$ write $A^k_\tau,B^k_\tau$ for the sets $A_\tau,B_\tau$ of \S\ref{sec:main}. Define the \emph{wall} of a probability $m_0$ and the set of priors \emph{avoiding} the vertex $e_b$ by
\[
\mathcal W_{m_0}\coloneqq\{q\in\Di:K_q=\{m_0\}\},\qquad \mathcal A_b\coloneqq\{q\in\Di:e_b\notin K_q\}.
\]

\begin{lemma}\label{lem:walls}
In every cover:
\begin{enumerate}[leftmargin=*,label=(\alph*)]
\item each wall $\mathcal W_{m_0}$ is convex and relatively closed in $\Di$;
\item each $\mathcal A_b$ is open and convex;
\item for a.e.\ pair $g$ and every $m_0\in\Delta$ with $m_0\cdot g\ne0$, $\cl F(g)\cap\Di$ contains no $q$ with $m_0\in K_q$;
\item for a vertex $v_m$ of $T$, $i\ne m$ and $\tau\in(0,1)$: $\mathcal W_{v_m}\subseteq A^i_\tau$, which is convex and does not contain $p^*$.
\end{enumerate}
\end{lemma}
\begin{proof}[Proof sketch]
(a) $\mathcal W_{m_0}=\bigcap_{h:\,m_0\cdot h>0}\{P(\cdot,h)>0\}$, an intersection of convex sets (Lemma~\ref{lem:quasiconcave}); for $m\in K_q\setminus\{m_0\}$ the gamble $h=(m_0-m)-((m_0-m)\cdot m)\one$ separates. (b) $\mathcal A_b=\{P(\cdot,\one-e_b)>0\}$. (c) \src{R4} Thm.~B: $K_q\subseteq g^\perp$ on $\cl F(g)\cap\Di$ for a.e.\ $g$. (d) On $\mathcal W_{v_m}$, $P(q,h^i(\tau))=v_m\cdot h^i(\tau)>0$, while $P(p^*,h^i(\tau))=0$.
\end{proof}
\status{proved (\src{N28} Lemmas 2.1--2.4); confirmed by one referee (\src{N28\_referee} \S2); second reading CONFIRMED (\src{TS\_second\_referee} \S1.1).}

\begin{theorem}\label{thm:TSalways}
Let $T$ be an all-facing triangle in general position (in particular $T_1$, for every $\varepsilon\in(0,1)$). In every cover of $T$, (TS) holds. Hence branch (II) of Theorem~\ref{thm:dichotomy} never occurs: there is no single-vertex pinning near $p^*$, by a ball or by a thin set.
\end{theorem}
\begin{proof}[Proof sketch]
Suppose (TS) fails. By Corollary~\ref{cor:R232} the $A$-half fails on some facet $F_m$, and branch (II) gives a \emph{wall of $F_m$ toward $V_l$} for some $l\ne m$: an open set $O$ with $p^*\in\cl O$ and $v_m\in K_q$ on $O$, such that every segment from $O$ to $V_l\cap\Di$ meets $\mathcal W_{v_m}$. (i) Such a wall forbids (TS) at every $\tau$ on $F_l$: since $h^l(\tau)_l<0$, the $c$-cluster sign puts $e_l$ on the closed negative side of the chord, so some point $v\in V_l\cap\Di$ lies strictly beyond it, while the wall points lie in $A^l_\tau$ (Lemma~\ref{lem:walls}(d)) inside the positive half-plane (Lemma~\ref{lem:quasiconcave}); for $o$ close to $p^*$ the wall point on $[o,v]$ would be closer to $p^*$ than $\operatorname{dist}(p^*,\mathcal W_{v_m})>0$. (ii) By Lemma~\ref{lem:R231} and \src{R10b} Thm.~PL, $B$-failures on $F_l$ cannot straddle a crossing next to the pinned set, so $F_l$ has walls toward both other vertices. (iii) Two relatively closed walls that meet every segment from sets accumulating at $p^*$ toward all three vertex neighbourhoods have $p^*$ in the convex hull of their union. (iv) Both walls lie in the convex set $A^n_\tau$ ($n$ the third index), which misses $p^*$.
\end{proof}
\status{proved (\src{N28} Prop.~3.1, Lemmas 3.2--3.4, Thm.~3.5); confirmed by one referee (\src{N28\_referee} \S3, every step re-derived, including that the walls of branch (II) meet \emph{every} segment from the pinned set to $V_l\cap\Di$, as the surround step needs). Second, independent reading CONFIRMED (\src{TS\_second\_referee} \S1, which re-derived every step from the statements of the cited lemmas, not from the first reports). Its inputs \src{N21} Prop.~N21.9 and Thm.~N21.4 (refereed in \src{N23}), \src{R10b} Thm.~PL, Lemma UVN, Lemma D, Lemma M, Cor.~PL2 (refereed in \src{R12}), and Lemma~\ref{lem:R231} and Corollary~\ref{cor:R232} (\src{N23}, with a second reading) were re-derived in the same report, so each now has two readings (items 6--8 of N21.4 only spot-checked).}

\begin{corollary}\label{cor:TSalways}
\begin{enumerate}[leftmargin=*,label=(\alph*)]
\item Every cover of an all-facing $T$ in general position is in branch (I) of Theorem~\ref{thm:dichotomy}: (TS), (PC), (NS) hold, and Theorem~\ref{thm:cuts} and Corollary~\ref{cor:availability} hold at every world with a failing box.
\item $T_1$ has no additive cover on any decomposable domain with (LF) on which, for a non-null set of mixed pairs, the advice range lacks one of the two acceptances: for instance, advice ranges contained in $\{+,0\}$ or in $\{0,-\}$ there, or a non-null set of mixed pairs on which every evaluable model gives no advice.
\item In every cover of $T_1$, $K_q\to\{e_b\}$ as $q\to e_b$ for every $b$, and every limit of $K_q$ at an edge prior with support $S$ lies in $\Delta_S$.
\end{enumerate}
\end{corollary}
\status{proved (\src{N28} Cor.~3.6); confirmed by one referee; second reading of (b), (c) CONFIRMED (\src{TS\_second\_referee} \S2), with the dependency note that (b), (c) rest on \src{N21} Thm.~N21.6, whose Step~1 uses the chart-form edge box. (b), (c) are \src{N21} Cor.~N21.7 and Thm.~N21.8(b), whose only extra hypothesis was (TS). So in every cover the advice range of almost every mixed pair contains both acceptances, and competition between them is not merely permitted but required.}

For the rest of this subsection fix a world $b$ of $T_1$, the index $a$ with a failing box $\Lambda_{ab}>1$, and the third index $c$. The two $b$-crossings are $w'\in\relint F_a$ and $w\in\relint F_c$ (for $(a,b,c)=(1,2,3)$, $w'\propto(-1,0,\kappa)$ and $w\propto(\kappa,0,-1)$). Under (TS) both crossing chords are the line $\lambda_b\coloneqq\lin(p^*,e_b)$. The tangent lines from $e_b$ to $T_1$ are $\ell'\coloneqq\lin(e_b,v_a)$ and $\ell\coloneqq\lin(e_b,v_c)$. One has $\{m\in\Delta:m\cdot w'\le0,\ m\cdot w\le0\}=\{e_b\}$ (\src{N28} Fact 4.0, confirmed). With the data $x_+$ the failing pairs $(a,b)$ are $(1,2),(2,3),(3,1)$; with $x_-$ they are $(2,1),(3,2),(1,3)$ (\src{N28\_referee}, NUMERICS from Proposition~\ref{prop:boxbridge}(c)).

\begin{theorem}\label{thm:apex}
In Theorem~\ref{thm:cuts}, the apex limit is $a^*=e_b$. Hence the cut d-faces $(a_m,z_m)$ of Theorem~\ref{thm:cuts} converge into $\partial\Delta$, to the segment $[e_b,z^*]$ of the edge $[e_b,e_a]$, with apexes converging to $e_b$. The same holds for every all-facing triangle in general position at every world with a failing box.
\end{theorem}
\begin{proof}[Proof sketch]
If $a^*\in\Di$, let $w^{**}$ be the other $b$-crossing. Along $[a^*,z^*]$, $K_q\subseteq\ell_{w^*}$ gives $P(q,w^{**})\ge0$. If $z^*\ne e_b$, convexity of $\{P(\cdot,w')\ge0\}$ and $P(\cdot,w')\le0$ off the positive side give a triangle in $B_1$ on which $P(\cdot,w')=0$, contradicting thin zero sets. If $z^*=e_b$, then $P(a^*,w^{**})=0$ forces $e_b\in K_{a^*}$; separating $a^*$ from the open convex $\mathcal A_b\ni p^*$ puts a ball around a point of $(a^*,e_b)$ outside $\mathcal A_b$, and d-face points of $g_m$ in that ball contradict Lemma~\ref{lem:walls}(c).
\end{proof}
\status{proved (\src{N28} Thm.~4.1); confirmed by one referee (\src{N28\_referee} \S4.1); second reading CONFIRMED WITH REPAIR (\src{TS\_second\_referee} \S\S3.2--3.3). This settles the case $a^*\in\Di$ that \src{N23} \S7.6 singled out. \emph{Repair.} Earlier versions stated here, in the introduction and in \S\ref{sec:general-open} that competition between accepting $g$ and accepting $-g$ at a positive distance from $\partial\Delta$ is impossible. That is not proved: only the cut d-faces of Theorem~\ref{thm:cuts} are shown to converge into $\partial\Delta$. Nothing proved here excludes an adjacent pair, far from $\partial D_{T}$, whose d-face contains a fixed interior prior $q$ with $\dim K_q\le1$ and $K_q\subseteq g^\perp$ (\src{TS\_second\_referee} \S3.3). According to that report, no proof in \src{N28}, \src{I1}, \src{I2} or \src{C2} used the broader sentence.}

\begin{lemma}\label{lem:R1N28}
In every cover of an all-facing $T$ in general position, for every facet $F_k$ and $\tau\in(0,1)$, with $H^\pm_\tau$ the open sides of $\lambda_\tau$ ($H^+_\tau$ containing $A^k_\tau$): $P(\cdot,h(\tau))>0$ only on $H^+_\tau$, $P(\cdot,h(\tau))<0$ on all of $\Di\cap H^-_\tau$, and $P(\cdot,h(\tau))\le0$ on $\lambda_\tau\cap\Di$. So the chord line separates $\Di$ globally.
\end{lemma}
\begin{proof}
By Theorem~\ref{thm:TSalways}, $Z_\tau\cap B_1=\lambda_\tau\cap B_1$. If $q_1\in\Di\cap H^-_\tau$ had $P(q_1,h(\tau))\ge0$, the convex set $\{P(\cdot,h(\tau))\ge0\}$ (Lemma~\ref{lem:quasiconcave}) would contain the triangle spanned by $q_1$ and $\lambda_\tau\cap B(p^*,r_1/2)$, on which also $P\le0$; so $P=0$ on a triangle meeting $B_1\setminus\lambda_\tau$, a contradiction. The positive statement is Lemma~\ref{lem:quasiconcave}.
\end{proof}
\status{proved in a referee report (\src{N28\_referee} \S4.2, Lemma R1); second reading CONFIRMED (\src{TS\_second\_referee} \S4.1), for every all-facing $T$ in general position. It supersedes the CHECKED remark after \src{N28} Prop.~4.3.}

\begin{proposition}\label{prop:vertexfree}
In every cover of $T_1$:
\begin{enumerate}[leftmargin=*,label=(\roman*)]
\item no interior credal set $K_q$, $q\in\Di$, contains a vertex of $\Delta$;
\item for every facet, $\tau\in(0,1)$ and $q\in\Di\setminus H^+_\tau$, $K_q$ meets $\{m:m\cdot h(\tau)\le0\}$, and it meets $\{m:m\cdot h(\tau)<0\}$ if $q\in H^-_\tau$;
\item for $q\in\lambda_b\cap\Di$, $K_q$ contains a point with $m\cdot w'\le0$ and a different point with $m\cdot w\le0$, so $K_q$ is not a singleton; on the whole open $e_a$-side of $\lambda_b$, $K_q$ meets $\{m\cdot w'<0\}$, and on the whole open $e_c$-side it meets $\{m\cdot w<0\}$;
\item no precise prior lies on a crossing chord. In particular a precise identity cap $K_q=\{q\}$ on a full neighbourhood of $e_b$ (\src{R4} \S9), and every diagonal-projective precise region meeting $\lambda_b$, are incompatible with every cover.
\end{enumerate}
\end{proposition}
\begin{proof}[Proof sketch]
(i) If $e_b\in K_y$, separate $y$ from $\mathcal A_b$, which contains points near $e_a$ and $e_c$ (Corollary~\ref{cor:TSalways}(c)); the separating functional is negative near $e_b$, and by Theorem~\ref{thm:apex} some apex $a_m\in\cl F(g_m)$ lies there, against Lemma~\ref{lem:walls}(c). (ii) is Lemma~\ref{lem:R1N28}. (iii) is (ii) for $w'$ and $w$ at the crossing parameter; a common point would be $e_b$, excluded by (i). (iv) A singleton on $\lambda_b$ would be $\{e_b\}$.
\end{proof}
\status{proved (\src{N28} Prop.~4.3); confirmed by one referee (\src{N28\_referee} \S4.3); second reading CONFIRMED (\src{TS\_second\_referee} \S4.2). The strict forms in (ii) and (iii) use Lemma~\ref{lem:R1N28} (now read twice); (i) and (iv) use only the weak form of (ii), and hold on the whole chord line on both sides of $p^*$ (\src{C1C2\_referee} \S5). On the chord itself only the weak inequalities hold; Corollary~\ref{cor:chord} gives the exact chord structure. Theorem~\ref{thm:A} extends (iv) to every injective precise open region, wherever it lies. The arithmetic for the identity cap was checked numerically in \src{N28} \texttt{checks.py} (3).}

\subsection{Reduction to a problem about credal families}\label{sec:general-CF}

\begin{definition}\label{def:CF}
Let $T\subseteq\Di$ be a credal set. The \emph{concurrent-chart family problem} CF$(T)$ asks for a Hausdorff-continuous map $q\mapsto K_q$ from $\Di$ to credal sets, with $K_{p^*}=T$ for some $p^*\in\Di$, such that for a.e.\ mixed $g\in\R^3$ there are $c,c'\in\R^3$ with $\operatorname{sign}c_i=\operatorname{sign}g_i=-\operatorname{sign}c'_i$ ($i=1,2,3$) and, for \emph{every} $q\in\Di$,
\[
P_{K_q}(g)>0\iff c\cdot q>\max(0,c'\cdot q),\qquad P_{K_q}(-g)>0\iff c'\cdot q>\max(0,c\cdot q).
\]
\end{definition}
In words: for almost every pair, the partition of the prior triangle into ``accept $g$'', ``accept $-g$'' and ``neither'' is the cell decomposition of $\max\{0,c\cdot q,c'\cdot q\}$, three lines through one point, with the vertex signs of \src{R1}.

\begin{proposition}\label{prop:CF}
If an additive ranking with integral penalties covers $T_1$ on some decomposable domain with (LF), then CF$(T_1)$ has a solution. Conversely, every solution of CF$(T_1)$ yields such a cover, on the all-model domain, provided Theorem~\ref{thm:ccibu} holds. The prior $p^*$ may be chosen arbitrarily in $\Di$.
\end{proposition}
\begin{proof}[Proof sketch]
($\Rightarrow$) Proposition~\ref{prop:necessary} gives the family $K_q$; it is Hausdorff-continuous by \src{R1} Prop.~3.1 ($L^1$-continuity of $q\mapsto B_q$, then the homeomorphism $K\mapsto1_{D_K}$; Proposition~\ref{prop:reduction} is stated for $U$-domain realisations and is not the right citation here). By Corollary~\ref{cor:TSalways}(b) the advice range of a.e.\ mixed pair contains both acceptances, so advice ranges are $\{+,0,-\}$ or $\{+,-\}$; a $\{+,-\}$ chart is the cell decomposition of $(\lambda d+e_0\one,(\lambda-1)d+e_0\one)$ for $0<\lambda<1$ and small $e_0>0$. \emph{Every-$q$ form.} The map $(q,g)\mapsto1[q\in R(g)]$ is jointly measurable, because $c$, $c'$ and $A(\cdot)$ are measurable in $g$; for each interior $q$ its $g$-section differs from $\{g:P(q,g)>0\}$ by a null set (a.e.\ uniqueness of the $q$-Bayes model at every $q$). By Fubini there is a co-null set of $g$, which we take to be the intersection of the co-null sets for $g$ and for $-g$ (the \emph{good} set, symmetric under $g\mapsto-g$), for which $R(\pm g)\cap\Di$ and $\{P(\cdot,\pm g)>0\}$ are a.e.\ equal; both are open and convex (Lemma~\ref{lem:quasiconcave}), hence equal. So for good $g$ the chart condition of Definition~\ref{def:CF} holds at \emph{every} $q\in\Di$. ($\Leftarrow$) Take advice ranges $\{+,0,-\}$, penalties from \src{R1} Prop.~2.1 (converse), rescaled per pair for integrability; a measurable choice of $(c,c')$ exists because the admissible set is Borel (test $C\subseteq R(c,c')\Leftrightarrow\min_CP(\cdot,g)>0$ over closed rational polygons $C$) and the Jankov--von Neumann theorem applies. Tie sets lie in $\pm\partial D_{K_q}$, so Bayes models are unique and coherent at every interior $q$, and Theorem~\ref{thm:ccibu} gives (LA), (S), (C). The last sentence is Proposition~\ref{prop:diag}(a).
\end{proof}
\status{CHECKED in \src{N28} (Prop.~5.1); confirmed with repair by one referee (\src{N28\_referee} \S5.1); ($\Rightarrow$) in the every-$q$ form re-derived and CONFIRMED WITH REPAIR by a second referee (\src{TS\_second\_referee} \S5), whose repairs are incorporated: the citation of \src{R1} Prop.~3.1 for Hausdorff continuity, the symmetric good set, and joint measurability for Fubini. CF data are not unique (the apex may lie outside $\Di$); uses of a specific $c(g)$ downstream are legitimate because, for $g$ in the cone over $N_0$, the advice range contains silence (by (LF)) and $c$ is the penalty vector $\ell(g,0)-\ell(g,+)$ (same report). Earlier repairs incorporated: the measurable-selection test above replaces the note's ``test on a countable dense set of $q$'', which presumed convexity of $\{P(\cdot,g)>0\}$ for an arbitrary candidate; ($\Rightarrow$) needs both acceptances in the advice range a.e., so it is specific to $T_1$ (and to targets whose covers have both acceptances in the advice range a.e.); ($\Leftarrow$) has the status of Theorem~\ref{thm:ccibu} (author-checked, no formal referee report).}

\begin{proposition}\label{prop:diag}
Let $D,D'$ be positive diagonal matrices and $\nu$ the normalisation $x\mapsto x/(\one\cdot x)$.
\begin{enumerate}[leftmargin=*,label=(\alph*)]
\item If $K$ solves CF$(T)$ at $p^*$, then $\tilde q\mapsto K_{\nu(D'^{-1}\tilde q)}$ solves CF$(T)$ at $\nu(D'p^*)$, with chart data $D'^{-1}c,D'^{-1}c'$.
\item If $K$ solves CF$(T)$ at $p^*$, then $q\mapsto\nu(D^{-1}K_q)$ solves CF$(\nu(D^{-1}T))$ at $p^*$, with data $(\tilde c,\tilde c')(Dg)\coloneqq(c,c')(g)$.
\item The map $X\mapsto DX$ on models, with $A(g)\mapsto A(D^{-1}g)$ on advice ranges, preserves decomposable domains, splicing, convergence in measure, null sets, guidance ranks, (IB), coherence ($D\cdot D_K=D_{\nu(D^{-1}K)}$) and (LF). Hence a ranking covers $D_T$ iff its transport covers $D_{\nu(D^{-1}T)}$; for additive rankings, multiply the penalties by the Radon--Nikodym factor of the transported Gaussian measure.
\end{enumerate}
\end{proposition}
\status{(a), (b): proved (\src{N28} Prop.~5.2); confirmed by one referee. (a) is Proposition~\ref{prop:reduction}(d). (c): proved in a referee report (\src{N28\_referee} \S5.2, Prop.~R3), one reader only. By (c), for every $T$ and every class of rankings closed under this transport, coverability of $T$ depends only on the orbit of $T$ under $T\mapsto\nu(D^{-1}T)$; diagonal gamble rescaling is a symmetry of the whole axiom system, and no appeal to Proposition~\ref{prop:CF} or Theorem~\ref{thm:ccibu} is needed.}

By \S\ref{sec:general-N28}, every solution of CF$(T_1)$ coming from a cover satisfies (TS), (PC), (NS), $K_q\to\{e_b\}$ at the vertices with edge confinement, cut apexes at the vertices, vertex-free interiors, global one-sidedness of chords, non-precise credal sets on the crossing chords, and the wall and d-face constraints of Lemma~\ref{lem:walls}. By Proposition~\ref{prop:CF}($\Leftarrow$) every solution comes from a cover (at the status of Theorem~\ref{thm:ccibu}), so the residual called (E4$''$) in \src{N28} is CF$(T_1)$ itself, together with these necessary conditions; the list is not exhaustive (\src{N28\_referee} \S6). Section~\ref{sec:gm} shows that these conditions, together with three further consequences of the chart structure, are inconsistent: there is no cover, and CF$(T_1)$ has no solution at the status of Theorem~\ref{thm:ccibu}.

\subsection{What was open on general decomposable domains}\label{sec:general-open}
This subsection records the residuals as they stood before Section~\ref{sec:gm}, with the corrections required by the referees of round~10.
\begin{itemize}[leftmargin=*]
\item \textbf{(E5a)} \emph{Single-vertex pinning near $p^*$.} Closed: it never occurs (Theorem~\ref{thm:TSalways}).
\item \textbf{(E4$'$)} \emph{The near-all-orientation cut regime.} Reduced to CF$(T_1)$ (Proposition~\ref{prop:CF}), and now closed: no cover exists (Theorem~\ref{thm:gm}). What Theorem~\ref{thm:apex} proves is that the cut d-faces of Theorem~\ref{thm:cuts} converge into $\partial\Delta$, with apexes at $e_b$. An earlier version said here that competition between the two acceptances at a finite distance from $\partial\Delta$ is excluded; that is not proved (\src{TS\_second\_referee} \S3.3), and it is not used. A precise cap on a full neighbourhood of a vertex is excluded (Proposition~\ref{prop:vertexfree}(iv)).
\item \emph{The sector-cap example.} \src{N28} inferred that, with the identity cap of \src{R4} \S9 excluded, the claim that (E4$'$) is consistent at the scale of the vertices had lost its support. The referee showed that this inference was unjustified at the time (\src{N28\_referee} \S4.3). The cuts of \src{R4} \S9 run from the apex to the failing edge and never reach $\lambda_b$. The identity $K_q=\{q\}$ restricted to a \emph{sector} $S_b$ at $e_b$ adjacent to the failing edge passes every constraint derived in \S\ref{sec:general-N28}, Lemma~\ref{lem:R1N28} included, and carries the cuts (angular extent $\varphi_{\max}\approx0.367,0.196,0.127,0.026$ at $\varepsilon=0.05,0.2,0.3,0.6$, in an angular parameter measured from the failing edge with $\varphi=\tfrac12$ on $\lambda_b$; status NUMERICS, \src{N28\_referee} counterexample C1). \emph{It is not part of any cover, nor of any solution of CF$(T_1)$.} The identity is the diagonal-projective map with $w=\one$, which is injective, and Theorem~\ref{thm:A} excludes every injective precise open region (\src{C1} Cor.~5.1, confirmed in \src{C1C2\_referee} \S4). Along route C2, precise caps along the failing edge are in fact forced (Theorem~\ref{thm:collar}), but with $w$ equal to $\kappa$ off $b$ and $1$ at $b$, not the identity (\src{C2} \S6, confirmed in \src{C1C2\_referee} \S4); and Section~\ref{sec:gm} shows that no cover exists at all. A first-order blow-up suggested that precise priors near $e_b$ occur only strictly on the failing side of $\lambda_b$ (NUMERICS N1 of \src{N28\_referee}); this is consistent with Theorem~\ref{thm:collar} and with Corollary~\ref{cor:chord}.
\item \emph{The chord layer (correction).} The previous version described a non-precise layer along $\lambda_b$ that pokes ``strictly beyond $\ell$ on the open $e_c$-side and beyond both $\ell$ and $\ell'$ on the chord''. The second half is not proved: on $\lambda_b$, Proposition~\ref{prop:vertexfree} gives only $P(q,w')\le0$ and $P(q,w)\le0$ (\src{TS\_second\_referee} \S4.3). It is contradicted by Corollary~\ref{cor:chord} (\src{C2} Cor.~3.3, confirmed in \src{C1C2\_referee} \S1): on $\lambda_b\cap\Di$, $q\ne p^*$, the credal set $K_q$ lies inside the tangent wedge at $e_b$ and touches both $\ell'$ and $\ell$. What is proved off the chord is that $K_q$ meets $\{m\cdot w'<0\}$ (lies partly strictly beyond $\ell'$) on the whole open $e_a$-side of $\lambda_b$, and meets $\{m\cdot w<0\}$ on the whole open $e_c$-side (Proposition~\ref{prop:vertexfree}(iii)). (The unproved sentence came from a summary bullet at the end of \src{N28\_referee} \S4, not from a computation; \src{C1C2\_referee} \S1.)
\item \emph{Conjecture N21/N28} (CF$(T_1)$ has no solution; \src{N28} Conj.~N28, restated in \src{N28\_referee} \S6, replacing \src{N21} Conj.~N21) is now Corollary~\ref{cor:CFnone}, at the status of Theorem~\ref{thm:ccibu}. Its ranking form, which needs no appeal to Theorem~\ref{thm:ccibu}, is Theorem~\ref{thm:gm}.
\end{itemize}

\section{No additive cover on general decomposable domains}\label{sec:gm}

This section shows that the necessary conditions of Section~\ref{sec:general} are inconsistent. A \emph{cover} is as in \S\ref{sec:general-N28}, with integral penalties unless bounded charges are mentioned explicitly: a strictly guiding additive ranking on a decomposable domain with (LF) at $D_{T_1}$ that satisfies (LA), (S), (C) and covers $T_1$, with interior Bayes credal sets $K_q$ and $K_{p^*}=T_1$. Since no cover exists, every statement below that begins ``in every cover'' is vacuously true; such statements are the steps of the proofs, and they are stated so that each step can be checked separately.

\subsection{The theorem}\label{sec:gm-thm}

\begin{theorem}\label{thm:gm}
Let $\varepsilon\in(0,1)$. No strictly guiding additive ranking with integral penalties and integrable envelopes on a decomposable domain with (LF) at $D_{T_1}$ satisfies (LA), (S), (C) and covers $T_1$. This holds at every interior target weighting $p^*$.
\end{theorem}
\status{proved, by three proofs that share only the inputs listed in Table~\ref{tab:gmdeps}.
\emph{Route I1} (\S\ref{sec:gm-I1}): \src{I1} Thm.~I1, confirmed by one referee (\src{I1\_referee}: CONFIRMED; its repairs R1--R7 are expository and are incorporated below).
\emph{Route I2} (\S\ref{sec:gm-other}): \src{I2} Thm.~6, confirmed by one referee with repairs (\src{I2\_referee} \S8; Lemma~1 repaired, and the projective step rests on Prop.~D$^*$, proved in that report).
\emph{Route C2, and the short route} (\S\ref{sec:gm-other}): \src{C2} Thm.~5.1 and \src{C1} Thm.~A, confirmed with repairs by one referee (\src{C1C2\_referee} \S6, Theorem~R; Prop.~P proved in that report; the finite exceptional set of $\varepsilon$ left in \src{C2} is empty by the referee's hand elimination).
The imported inputs have had two independent readings, with the exceptions flagged in Table~\ref{tab:gmdeps}; each new step has had one reading, and Theorem~\ref{thm:D} (one ingredient of the short route) none beyond the overlapping Prop.~P and Prop.~D$^*$. No differentiability, regularity, continuity of chart data or linear scaling at the vertices is assumed anywhere. The box inequality $1/\kappa\le x^2\le\kappa$ of Proposition~\ref{prop:T1} does not appear in route I1's own steps; but box witnesses enter all three routes through the inputs (Theorem~\ref{thm:cuts}, hence Corollary~\ref{cor:TSalways}(b),(c), uses the chart-form edge box of Lemma~\ref{lem:edgebox}), and routes I2 and C2 also use the box certificate of Proposition~\ref{prop:T1} directly (\src{I1\_referee} R7, \src{TS\_second\_referee} \S2). It settles Conjecture~N21/N28 in its ranking form.}

\begin{corollary}\label{cor:CFnone}
For every $\varepsilon\in(0,1)$, CF$(T_1)$ (Definition~\ref{def:CF}) has no solution.
\end{corollary}
\status{at the status of Theorem~\ref{thm:ccibu} (author-checked; no formal referee report). Every step of every route uses properties proved for \emph{covers} (Corollary~\ref{cor:TSalways}, Proposition~\ref{prop:vertexfree}), so a solution of CF$(T_1)$ is excluded only once it is known to come from a cover. That is the converse direction of Proposition~\ref{prop:CF}: it builds a ranking from the solution and needs Theorem~\ref{thm:ccibu} to verify (LA), (S), (C). Hence the status (\src{I1} Cor.~(b), \src{I1\_referee} \S8, \src{I2\_referee} \S8, \src{C1C2\_referee} \S6). Theorem~\ref{thm:gm} itself does not depend on Theorem~\ref{thm:ccibu}. Source of the conjecture it settles: \src{N28} Conj.~N28 (formerly \src{N21} Conj.~N21).}

\begin{corollary}\label{cor:gmcharges}
The conclusion of Theorem~\ref{thm:gm} holds when the penalties are bounded charges vanishing on $\mu$-null sets, provided the interior Bayes minimisers exist.
\end{corollary}
\status{proved along route I2 only, at the status of \src{N14} Thm.~C2 (proved in that referee report, one reader; its splice step refereed in \src{N24}). \src{I2\_referee} \S8 confirms route I2 for charges: by \src{N14} Thm.~C2, at every interior $q$ the Bayes minimiser coincides a.e.\ with the Bayes model of the countably additive parts, an integral rule, so the interior Bayes family and its charts are those of integral penalties, and \S\ref{sec:general-N28} already covers charges. The other routes have not been read for charges. \src{I1} Cor.~(c) (route I1 for charges) is CHECKED only; its referee did not referee it (\src{I1\_referee} \S0). Theorem~\ref{thm:D} has not been checked for charges: \src{PropD} Lemma~1.1(c) notes that its chart hypothesis (DL) holds once a pointwise chart field reproduces the Bayes models a.e., which is what \src{N14} Thm.~C2 supplies, but no reader has checked that step. So the short route and route C2 are recorded here for integral penalties only. The existence proviso is discharged for complete rankings in Theorem~\ref{thm:complete}; in general, whether coverage forces interior minimisers to exist when the penalties are charges is open (\src{N14} \S5.1).}

\paragraph{The three routes.} All three start from the structure of \S\ref{sec:general-N28}: (TS), (PC), (NS) with facet data $x\in\{x_+,x_-\}$ satisfying the corner equation $x+1/x=\kappa+1$, both acceptances in the advice range of a.e.\ mixed pair, $K_q\to\{e_b\}$ at the vertices, edge confinement, vertex-free interiors, and charts in the exact cell form of CF$(T_1)$ at every interior prior.
\begin{itemize}[leftmargin=*]
\item \emph{Route I1: the edges pin the chart lines} (\S\ref{sec:gm-I1}). Edge confinement makes the limits of $K_q$ at edge priors one-dimensional, and the facet chords from $p^*$ determine them exactly. This pins the chart line of every gamble whose zero line crosses that edge, which forces explicit supporting lines of $K_q$ near $p^*$. Three of them bound a convex set on both sides of a certain line through $p^*$ only if the facet maps are diagonal, which happens iff $x=1$. The corner equation gives $x\ne1$.
\item \emph{Route I2: the box returns as a precise collar} (\S\ref{sec:gm-other}). On the positive side of a facet chord, a non-singleton $K_q$ never straddles the facet line. So near each vertex the facet gambles of the two crossing facets have exact half-plane behaviour at non-precise priors, and the failing box forces every prior near the failing edge to be precise. That precise collar is diagonal-projective with one weight vector for all precise priors; at the interface between the collar and the non-precise chord layer, a blow-up gives one relation per vertex, and the three relations contradict the corner equation.
\item \emph{Route C2, and a short route} (\S\ref{sec:gm-other}). The same forced collars, glued to the non-precise layer, pin the weight vector at each vertex to a value that depends on the vertex, against the transport theorem. Shorter still: forced collars are injective precise regions, and no injective precise open region exists in any cover (Theorem~\ref{thm:A}). The short route uses no conic algebra and no Bezout argument.
\end{itemize}

\subsection{Common inputs}\label{sec:gm-inputs}

\paragraph{Normalisations.} Multiplying the penalties of world $i$ by $p^*_i>0$ leaves every world preorder, the domain, (LA), (S), (C) and (LF) unchanged and re-indexes the Bayes family by $q\mapsto\nu(p^*\odot q)$, which sends $u$ to $p^*$; so we may take $p^*=u$ (Proposition~\ref{prop:diag}(a); \src{I1\_referee} \S5). The transposition of coordinates $1,2$ maps covers of $T_1$ to covers of $T_1$, fixes $u$ and sends the facet data $x$ to $1/x$ (proof of Proposition~\ref{prop:T1}); since $x_-=1/x_+$ we may take $x=x_+$, so that $x>1$ and $x^2>\kappa$. The facet vectors are then $c_1(h)=(-(xh_2+h_3),\,xh_2,\,h_3)$, $c_2(h)=(h_1,\,-(h_1+xh_3),\,xh_3)$, $c_3(h)=(xh_1,\,h_2,\,-(xh_1+h_2))$ (\S\ref{sec:contradiction}), and the failing pairs $(a,b)$ are $(1,2),(2,3),(3,1)$. For a facet gamble $h\in\relint F_k$, its \emph{chord} is the line $\{q\in\aff\Delta:c_k(h)\cdot q=0\}$ through $p^*$; by (PC) it is the cluster line of $h$.

\paragraph{Chart cells.} By Proposition~\ref{prop:CF}($\Rightarrow$), for a.e.\ mixed $g$ (the \emph{good} gambles; a symmetric co-null set) there are $c,c'$ with the sign patterns of $g$ and $-g$, $d\coloneqq c-c'$, such that, with the plane cells
\[
R(g)=\{c\cdot q>0,\ d\cdot q>0\},\quad R(-g)=\{c'\cdot q>0,\ d\cdot q<0\},\quad S(g)\coloneqq\{c\cdot q\le0,\ c'\cdot q\le0\},
\]
for \emph{every} $q\in\Di$: $P(q,g)>0\iff q\in R(g)$ and $P(q,-g)>0\iff q\in R(-g)$. The complement of $S(g)$ in the plane is $R(g)\cup R(-g)\cup\{d\cdot q=0<c\cdot q=c'\cdot q\}$; on the closed d-face, $K_q\subseteq g^\perp$ (Lemma~\ref{lem:walls}(c), i.e.\ \src{R4} Thm.~B). On $S(g)\cap\Di$, neither acceptance is strict, so $K_q$ meets the probability line $\ell_g\coloneqq\{m\in\aff\Delta:m\cdot g=0\}$. For $g$ in the cone over $\operatorname{int}N_0$ the advice range contains silence (by (LF)), so $c$ is the penalty vector, and on $B_1$ the chart of $g$ is $\{c\cdot q>0\}\cap B_1$ (Lemmas~\ref{lem:B1} and~\ref{lem:raytransfer}(c)). For pairs with advice range $\{+,-\}$ the CF data are those of the proof of Proposition~\ref{prop:CF}, with $S(g)=\varnothing$.

\begin{lemma}\label{lem:veccluster}
In every cover, let $h\in\relint F_k$ and let good gambles $g_m\to h$ lie in the cone over $\operatorname{int}N_0$. Then $c(g_m)/|c(g_m)|\to c_k(h)/|c_k(h)|$ in $\R^3$. Consequently, if $q_0\in\aff\Delta$ has $c_k(h)\cdot q_0\ne0$, there are $\delta,\rho>0$ such that $c(g)\cdot q$ has the sign of $c_k(h)\cdot q_0$ for every good $g$ in that cone with $|g-h|<\delta$ and every $q$ with $|q-q_0|<\rho$.
\end{lemma}
\begin{proof}[Proof sketch]
A limit $\hat c$ of the normalised $c(g_m)$ defines an affine function on $\aff\Delta$. On $B_1$ the charts are half-planes, so by the cluster clause of (PC) (Theorem~\ref{thm:TS}) its zero line is the chord of $h$, with the same positive side. An affine function on $\aff\Delta$ determines its coefficient vector in $\R^3$, because $\aff\Delta$ spans $\R^3$; and $\hat c\not\propto\one$, since $\hat c$ has the weak sign pattern of the mixed $h$. Hence $\hat c$ is the normalised $c_k(h)$. The uniform statement follows by contradiction along a subsequence.
\end{proof}
\status{proved: the cluster clause of Theorem~\ref{thm:TS} (\src{N21} Thm.~N21.4 item~5) with Lemma~\ref{lem:B1}. The vector form and its uniform version are written out in \src{I1\_referee} \S1 and were re-derived independently in \src{I2\_referee} \S1, \src{C1C2\_referee} \S1 and \src{TS\_second\_referee} \S\S1.1,~6. It is \emph{not} available at corner rays, where Theorem~\ref{thm:pivot} covers only $\tau\in(0,1)$ (\src{TS\_second\_referee} \S6). Positive homogeneity and Lemma~\ref{lem:raytransfer} extend it from the facet segments of $\Sigma$ to all of $\relint F_k$ (\src{C1C2\_referee} \S1).}

\subsection{First proof: the edges pin the chart lines}\label{sec:gm-I1}

Throughout this subsection fix a cover of $T_1$, normalised by $p^*=u$ and $x=x_+$. For an edge $E=[e_i,e_j]$ write $z(t)\coloneqq(1-t)e_i+te_j$ for its points read as probabilities and $y(s)\coloneqq(1-s)e_i+se_j$ for the same points read as priors. For $t\in(0,1)$ exactly one of the two tangent lines from $z(t)$ to $T_1$ has $e_j$ on the side of $T_1$; it is $\ell_h$ for a facet gamble $h=h_j(t)$ touching some vertex $v_k$. The \emph{edge map} $\mu_j$ sends $t$ to the parameter $s$ at which the chord of $h_j(t)$ meets $E$; $\mu_i$ is defined in the same way from the tangent line with $e_i$ on the side of $T_1$, with $s$ still measured from $e_i$. Put $Z_o\coloneqq\mu_o^{-1}$. A prior $y(s)\in\relint E$ is \emph{non-failing} if $Z_j(s)<Z_i(s)$; the non-failing priors form an open set $N_E$.

\begin{lemma}\label{lem:edgemaps}
In every cover of $T_1$, normalised as above:
\begin{enumerate}[leftmargin=*,label=(\alph*)]
\item for $h\in\relint F_k$ and every $l$ with $h_l\ne0$, $c_k(h)_l$ has the sign of $h_l$; so every facet chord separates the vertices of $\Delta$ strictly, as $\ell_h$ does;
\item $\mu_i$ and $\mu_j$ are increasing homeomorphisms of $[0,1]$ onto itself fixing $0$ and $1$;
\item each of $\mu_i,\mu_j$ is piecewise M\"obius with two pieces, the break being where the tangent line passes through two vertices of $T_1$ (at $t=c_0$ or $t=1-c_0$); on the piece where the tangent line touches $v_k$, $\mu_o$ is the restriction of a M\"obius map that fixes the endpoint $e_l$ of $E$ for $l\ne k$, and that fixes $e_k$ iff the chord of the probability line through $v_k$ and $e_k$ passes through $e_k$.
\end{enumerate}
\end{lemma}
\begin{proof}[Proof sketch]
(a) For $k=1$, coordinates $2,3$ follow from the crossing rows. Along the facet segment from $r_{31}$ to $r_{12}$, $xh_2+h_3$ is affine with endpoint values $(x(2+\varepsilon)+\varepsilon-1)/(1-\varepsilon)$ and $\kappa+1-x$, which are positive iff $1/(\kappa+1)<x<\kappa+1$; this holds for $x_\pm$ by the corner equation. The other facets follow by cyclic symmetry. (b) By (a) the chord meets $\relint E$, so $\mu_o$ is well defined; at the ends the tangent line tends to a crossing gamble, whose chord is $\lin(p^*,e_i)$ or $\lin(p^*,e_j)$. Injectivity: the chord normal $n_h$ is strictly monotone on each facet (Theorem~\ref{thm:pivot}) and continuous at the corners (NS). It cannot fold back at a corner. The signs of $n\cdot(e_l-u)$ cut the circle into six sectors met in the cyclic order $(+,-,-)$, $(+,+,-)$, $(-,+,-)$, $(-,+,+)$, $(-,-,+)$, $(+,-,+)$, and by (a) facet~1 carries $n$ through $(-,+,-)\to(-,+,+)\to(-,-,+)$, across a whole sector, which fixes the direction of rotation; facets~2 and~3 continue in the same direction. So $n$ runs once round the circle, monotonically. On the arc used by $\mu_j$, $h_j>0>h_i$ confines $n$ to an open arc shorter than $\pi$, so distinct $t$ give distinct chords, which meet $E$ in distinct points. (c) $\mu_o$ is the composite of the projective bijections $E\to$ pencil of $v_k\to$ pencil of $p^*\to E$, the middle one induced by $\Gamma_k$. The line through $v_k$ and $e_l$ is a crossing line iff $l\ne k$, and then its chord is $\lin(p^*,e_l)$.
\end{proof}
\status{proved (\src{I1} Lemmas~1--2); confirmed by one referee (\src{I1\_referee} \S3; (a) recomputed symbolically by the endpoint argument given here; (b) CONFIRMED WITH REPAIR R2, the no-fold-back argument written out above). Explicit maps on $[e_1,e_2]$: facet~3 gives $s=tx/(tx-t+1)$, which fixes both ends; facet~1 fixes $e_2$ only, with $s(0)=(x-1)/(2x-1)$ (sympy, \src{I1}, \src{I1\_referee}).}

\begin{proposition}\label{prop:edgelimits}
In every cover of $T_1$, normalised as above, let $y(s)\in\relint E$. Every Hausdorff limit $K^*$ of $K_q$ as $q\to y(s)$ through $\Di$ is a subset of $E$; identify it with a subset of $[0,1]$ through the parameter $t$. Then
\begin{enumerate}[leftmargin=*,label=(\roman*)]
\item $\min K^*\le Z_j(s)$ and $\max K^*\ge Z_i(s)$;
\item $K^*$ contains no two points $t'<t<t''$ with $t\in(0,Z_j(s))\cup(Z_i(s),1)$.
\end{enumerate}
Hence $K_q\to[Z_j(s),Z_i(s)]$ if $y(s)\in N_E$, and every limit is a single point of $[Z_i(s),Z_j(s)]$ if $Z_j(s)>Z_i(s)$.
\end{proposition}
\begin{proof}[Proof sketch]
$K^*\subseteq E$ is edge confinement (Corollary~\ref{cor:TSalways}(c)). Fix $t$ and let $\hat h$ be the normalised $h_j(t)$. By Lemma~\ref{lem:edgemaps}(a), $c_k(\hat h)\cdot y(s)$ has the sign of $s-\mu_j(t)$, and by Lemma~\ref{lem:veccluster} so has $c(g)\cdot q$ for good $g$ near $\hat h$ and $q\in\Di$ near $y(s)$. If $t<Z_j(s)$, then $c(g)\cdot q>0$, so $q\notin S(g)$, so $K_q$ does not strictly straddle $\ell_g$; in the limit $K^*$ does not strictly straddle $t$. If $t>Z_j(s)$, then $c(g)\cdot q<0$, so $q\notin R(g)$ and $K_q$ meets the closed negative side of $\ell_g$; in the limit $\min K^*\le t$. The case $o=i$ is the mirror image. In the non-failing case (ii) forbids $\min K^*<Z_j(s)$ and $\max K^*>Z_i(s)$; in the failing case $(0,Z_j)\cup(Z_i,1)=(0,1)$, so $K^*$ is a point. The two parameters $t=c_0,1-c_0$ at which $h_o(t)$ is a corner ray are not covered by Lemma~\ref{lem:veccluster}; they are excluded, which is harmless because (ii) concerns open sets of $t$ and (i) infima and suprema over $t$.
\end{proof}
\status{proved (\src{I1} Prop.~E); confirmed by one referee (\src{I1\_referee} \S3, with repair R3: the two corner-ray parameters, excluded as above; also \src{TS\_second\_referee} \S6). At the failing ends it reproduces the precise edge limits of \src{R4} \S9 and \src{N28\_referee}; at the midpoint of $[e_1,e_2]$ (the end of the crossing chord $\lin(u,e_3)$) the limit is the shadow of $T_1$ from $e_3$, with parameters $(1-\varepsilon)/(2+\varepsilon)$ and $(1+2\varepsilon)/(2+\varepsilon)$, non-failing for every $\varepsilon$ (NUMERICS at 12 values of $\varepsilon$, \src{I1\_referee}).}

\begin{proposition}\label{prop:clines}
In every cover of $T_1$, normalised as above, let $g$ be good, suppose $\ell_g$ meets $\relint E$ at $z(t)$ with $e_o$ on its positive side ($o\in\{i,j\}$), and suppose the prior $y\coloneqq y(\mu_o(t))$ is non-failing. Then $y\in\cl R(g)\cap S(g)$; in particular $c(g)\cdot y=0$ and $c'(g)\cdot y\le0$.
\end{proposition}
\begin{proof}[Proof sketch]
For $o=j$: at edge priors slightly beyond $y$ toward $e_j$, Proposition~\ref{prop:edgelimits} puts every limit of $K_q$ strictly on the positive side of $\ell_g$, so nearby interior $q$ lie in $R(g)$. At edge priors slightly before $y$ the limit is a non-degenerate interval straddling $t$ (non-failing, and $Z_j<t<Z_i$ by continuity), so nearby $K_q$ strictly straddle $\ell_g$ and $q\in S(g)$. $S(g)$ is closed, so $y\in\cl R(g)\cap S(g)$.
\end{proof}
\status{proved (\src{I1} Thm.~R); confirmed by one referee (\src{I1\_referee} \S4, which checks that the null sets of good $g$ do not depend on $q$ and records that for $\{+,-\}$ pairs the hypotheses cannot hold, a by-product rather than an inconsistency). It applies to every good gamble whose zero line crosses $E$, not only to near-facet gambles.}

\begin{lemma}\label{lem:forcedlines}
In every cover of $T_1$, normalised as above, let $q\in\Di$ and let $\Lambda\ni q$ be an oriented line that meets two edges $E_1,E_2$ at non-failing priors $y_1\in N_{E_1}$, $y_2\in N_{E_2}$, with $e_{o_m}$ the end of $E_m$ on the positive side of $\Lambda$. Let $z_m\in E_m$ be the probability with parameter $Z^{E_m}_{o_m}(y_m)$, and let $g_L$ be a gamble with $\ell_{g_L}=\lin(z_1,z_2)$ and $e_{o_1},e_{o_2}$ on its positive side. Then $\min_{m\in K_q}m\cdot g_L=0$.
\end{lemma}
We call $\lin(z_1,z_2)$ a \emph{forced supporting line} of $K_q$: it touches $K_q$, and $K_q$ lies on its closed positive side.
\begin{proof}[Proof sketch]
For good $g$ near $g_L$, $\ell_g$ crosses $E_1,E_2$ near $z_1,z_2$ with the same orientations, and their $\mu$-images lie in the open sets $N_{E_m}$. By Proposition~\ref{prop:clines} both images lie on the ray $\{c(g)=0,\ c'(g)\le0\}$, which is convex, so the whole segment of the line $\{c(g)\cdot q=0\}$ inside $\Delta$ lies in $\cl R(g)\cap S(g)$, and $P(\cdot,g)=0$ on it. For good $g_n\to g_L$ these segments converge to $\Lambda\cap\Delta$ (continuity of $\mu$), and joint continuity of $P$ gives $P(q,g_L)=0$.
\end{proof}
\status{proved (\src{I1} Lemma~S); confirmed by one referee (\src{I1\_referee} \S4; NUMERICS on a genuine CF solution, the normalised-box system: $|\min_{K_q}g_L\cdot m|\le1.1\cdot10^{-15}$ over 400 random $(q,\Lambda)$).}

Work now on the facet $F_1$, with the edges $E_1\coloneqq[e_1,e_2]$ and $E_2\coloneqq[e_3,e_1]$, which meet at $e_1$. Let $\mu^{(1)}$ be facet~1's M\"obius piece of the edge map of $E_1$ toward $e_2$, and $\mu^{(2)}$ facet~1's M\"obius piece of the edge map of $E_2$ toward $e_3$. Put $h^*\coloneqq(0,1,-1)$, so that $\ell_{h^*}=\lin(v_1,e_1)$ (not a facet line), and $C_1\coloneqq\{q:c_1(h^*)\cdot q=0\}$, a line through $u$; here $c_1(h^*)=(1-x,x,-1)$.

\begin{lemma}\label{lem:window}
In every cover of $T_1$, normalised as above:
\begin{enumerate}[leftmargin=*,label=(\alph*)]
\item There are three distinct priors $y^{(1)},y^{(2)},y^{(3)}\in\relint E_1$ on the $e_2$-side of the midpoint of $E_1$, and a neighbourhood $U$ of $u$, such that for every $q\in U$ and every $i$: $y^{(i)}\in N_{E_1}$, and $Z^{E_1}_2(y^{(i)})$ lies in facet~1's piece; and $\lin(q,y^{(i)})$ meets $\relint E_2$ at a prior of $N_{E_2}$ whose $Z^{E_2}_3$-image lies in facet~1's piece.
\item For $q\in U$ let $\rho_q:E_1\to E_2$ be the central projection from $q$ and $\pi_q\coloneqq(\mu^{(2)})^{-1}\circ\rho_q\circ\mu^{(1)}$, a projectivity from $E_1$ to $E_2$. Then $\pi_u$ is the perspectivity with centre $v_1$; $e_1\in C_1$ iff $x=1$; and if $x\ne1$, then $\pi_q(e_1)=e_1$ iff $q\in C_1$, and $q\mapsto\pi_q(e_1)$ has nonzero derivative transverse to $C_1$ at $u$.
\item In affine coordinates $(X,Y)$ on $\aff\Delta$ with origin $e_1$ and axes along $E_1$ and $E_2$, write $\pi(X)=(\alpha X+\beta)/(\gamma X+\delta)$, and let $\ell(X)\coloneqq(1/X,1/\pi(X),-1)$ be the coefficients of the line through $(X,0)$ and $(0,\pi(X))$. For distinct nonzero $X_1,X_2,X_3$ put $\mathcal D\coloneqq\det[\ell(X_1);\ell(X_2);\ell(X_3)]$. Then $\mathcal D=0$ whenever $\beta=0$, and $\partial\mathcal D/\partial\beta$ at $\beta=0$ equals $(\delta/\alpha^2)(w_1-w_2)(w_2-w_3)(w_3-w_1)$ with $w_i=1/X_i$.
\end{enumerate}
\end{lemma}
\begin{proof}[Proof sketch]
(a) Fix the three priors first, then shrink $U$. At the midpoint $m_{12}$ of $E_1$, Proposition~\ref{prop:edgelimits} gives $Z^{E_1}_2=(1-\varepsilon)/(2+\varepsilon)$, which exceeds $c_0$ and so lies in facet~1's piece $(c_0,1)$, and $m_{12}\in N_{E_1}$. The line $\lin(u,m_{12})$ ends at $e_3$; near $e_3$ the endpoint slopes of the two edge maps of $E_2$ are $1/x$ and $\kappa x$, so $N_{E_2}$ contains an initial segment $(e_3,e_3+\delta_0)$ iff $\kappa x^2>1$, which holds for $x_+$; and toward $e_3$ facet~1 supplies the piece on $[0,1-c_0]$. Choose $y^{(i)}$ on the $e_2$-side of $m_{12}$, so close to $m_{12}$ that $\lin(u,y^{(i)})$ meets $E_2$ within $\delta_0$ of $e_3$ with $Z^{E_2}_3$-image below $1-c_0$, and the conditions on $E_1$ hold; these are open conditions. Since $y^{(i)}\ne m_{12}$, $u\notin\lin(y^{(i)},e_3)$, so the conditions persist for $q$ in a small $U$. (The required closeness to $m_{12}$ depends on $\varepsilon$.) (b) $\rho_u$ follows chords through $u$, and $\mu^{(1)},\mu^{(2)}$ are induced by the same pencil map of facet~1, so $\pi_u(z)=\lin(v_1,z)\cap E_2$. By Lemma~\ref{lem:edgemaps}(c), $\mu^{(m)}(e_1)=a_m\coloneqq C_1\cap\lin E_m$; so $\pi_q(e_1)=e_1$ iff $\lin(q,a_1)$ meets $\lin E_2$ at $a_2$, iff $q\in\lin(a_1,a_2)=C_1$. Transversality holds because $\lin(q,a_1)$ pivots about $a_1\notin\lin E_2$ and $(\mu^{(2)})^{-1}$ is a diffeomorphism of the projective line; at $\varepsilon=1/7$ (where $x_+=2$) the point $a_2$ is at infinity and the statement is read in a projective chart of $\lin E_2$. Finally $c_1(h^*)_1=1-x$. (c) At $\beta=0$, $\ell(X_i)=w_iA+B$ for fixed $A,B$, so $\mathcal D=0$; differentiating, the $\gamma$-part cancels and the $\delta$-part is a Vandermonde determinant.
\end{proof}
\status{proved (\src{I1} Lemmas~W, P, D); confirmed by one referee (\src{I1\_referee} \S5): (a) CONFIRMED WITH REPAIR R5 (order of quantifiers: the window points are fixed before $U$ is shrunk, as above; the limit $y\to m_{12}$ itself is degenerate); (b) CONFIRMED with repair R6 (projective reading at $\varepsilon=1/7$; NUMERICS $\varrho/r=\pm2.31\cdot10^{-3}$ there); (c) CONFIRMED, exact in sympy. The window exists for every $\varepsilon\in(0,1)$ (NUMERICS at 12 values of $\varepsilon$, including $1/7$ and $0.995$).}

\begin{proof}[First proof of Theorem~\ref{thm:gm} (sketch)]
Normalise $p^*=u$ and $x=x_+$, and take the window of Lemma~\ref{lem:window}(a). Put $z_i\coloneqq(\mu^{(1)})^{-1}(y^{(i)})$, which does not depend on $q$, and $L_i(q)\coloneqq\lin(z_i,\pi_q(z_i))$ for $q\in U$. Since $\pi_q(z_i)$ is the $Z^{E_2}_3$-image of $\lin(q,y^{(i)})\cap E_2$, Lemma~\ref{lem:forcedlines} applies with $\Lambda=\lin(q,y^{(i)})$: each $L_i(q)$ is a forced supporting line of $K_q$, with $e_2,e_3$ on its positive side. Let $g_i(q)$ be the corresponding gambles, normalised continuously in $q$.

At $q=u$, $\pi_u$ is the perspectivity from $v_1$, so the $L_i(u)=\lin(v_1,z_i)$ are three facet lines through $v_1$ in increasing order, and $g_2(u)=a\,g_1(u)+b\,g_3(u)$ with $a,b>0$ and no $\one$-component. For $q$ near $u$ write $g_2(q)=a(q)g_1(q)+b(q)g_3(q)+\varrho(q)\one$ with $a(q),b(q)>0$. Then $\varrho(q)$ is the value of $g_2(q)$ at the point $L_1(q)\cap L_3(q)$, and $\varrho(q)=\sigma_0\,\mathcal D(\pi_q)$ with $\sigma_0$ a locally constant nonzero factor (orientations and normalisation).

In the coordinates of Lemma~\ref{lem:window}(c), $\pi_u(X)=Y_vX/(X-X_v)$ with $v_1=(X_v,Y_v)$ on neither edge line, so $\alpha\delta\ne0$ near $u$, and $\beta(q)=0$ iff $\pi_q(e_1)=e_1$. By Lemma~\ref{lem:window}(b), since $x_+\ne1$, $\beta(q)$ vanishes exactly on $C_1$ near $u$ and changes sign across it with nonzero derivative; by Lemma~\ref{lem:window}(c), so does $\varrho(q)$. Pick $q\in U$ on the side of $C_1$ where $\varrho(q)>0$. Lemma~\ref{lem:forcedlines} gives $K_q\subseteq\{m\cdot g_1\ge0\}\cap\{m\cdot g_3\ge0\}$, so $m\cdot g_2\ge\varrho(q)>0$ on $K_q$ (using $m\cdot\one=1$), whereas the same lemma gives $\min_{K_q}m\cdot g_2=0$. This is a contradiction.
\end{proof}
\status{proved (\src{I1} Thm.~I1); confirmed by one referee (\src{I1\_referee} \S5: the identity $l_2\cdot(l_1\times l_3)=-\det(l_1,l_2,l_3)$ was checked symbolically, and an independent certificate using the full piecewise edge maps, with no M\"obius fitting, gives $\varrho/r$ of opposite signs on the two sides of $C_1$, stable in $r$, at every $\varepsilon$ tested, e.g.\ $\pm6.452\cdot10^{-3}$ at $\varepsilon=0.3$). On the normalised-box system, a genuine CF solution, the edge limits, pinned chart lines and forced supporting lines are reproduced exactly and the certificate does not fire, as it must not: its facet maps are diagonal, so every $C_k$ passes through $e_k$ (\src{I1\_referee} \S7). The inputs are listed in Table~\ref{tab:gmdeps}.}

\paragraph{Where the contradiction lives.} The relation that replaces the box is
\[
\text{(Diag)}\qquad e_1\in C_1:\ \text{the chord of the probability line through $v_1$ and $e_1$ passes through $e_1$.}
\]
It says that facet~1's pencil map is the restriction of a projective map fixing $e_1$; with the crossing rows, which fix $e_2$ and $e_3$, this makes $c_1$ restricted to $v_1^\perp$ the restriction of a diagonal map. For the facet data of $T_1$, (Diag) holds iff $x=1$, and the corner equation $x+1/x=\kappa+1>2$ gives $x\ne1$. So the facet maps of a cover would have to be diagonal, and the corner equation forbids it. The contradiction appears at first order in $|q-p^*|$, on one side of $C_1$; the failing ends of the edges (where the sector caps of \S\ref{sec:general-open} lived) are never reached (\src{I1} \S6, interpretation).

\subsection{Precise collars and precise regions}\label{sec:gm-precise}

A prior $q\in\Di$ is \emph{precise} if $K_q$ is a singleton, written $\{\sigma(q)\}$; a \emph{precise region} is a connected open set of precise priors. For $w\in(0,\infty)^3$ the \emph{diagonal-projective} map $\sigma_w(q)\coloneqq w\odot q/(w\cdot q)$ is a bijection of $\Di$ fixing the vertices of $\Delta$.

\begin{lemma}\label{lem:nostraddle}
In every cover of $T_1$, let $h\in\relint F_k$ and $q\in\Di$ with $c_k(h)\cdot q>0$. Then $P(q,h)>0$, or $P(q,-h)>0$, or $K_q$ is a singleton contained in $h^\perp$. More precisely, there is $r>0$ such that no $g$ with $|g-h|<r$ has $\min_{m\in K_q}m\cdot g<0<\max_{m\in K_q}m\cdot g$.
\end{lemma}
\begin{proof}[Proof sketch]
By Lemma~\ref{lem:veccluster}, $c(g)\cdot q>0$ for every good $g$ near $h$ in the cone over $\operatorname{int}N_0$. So $q\notin S(g)$, and $q$ lies in $R(g)$, in $R(-g)$, or on the d-face, where $K_q\subseteq g^\perp$. In no case does $K_q$ strictly straddle $g^\perp$. Strict straddling is open in $g$ and good gambles are dense, which gives $r$. If $K_q$ contains $m_0\in h^\perp$ and some $m_1\ne m_0$, a gamble near $h$ separates them strictly: $h-\delta\one$ if $m_1\cdot h>0$; $h+\delta\one$ if $m_1\cdot h<0$; and $h+\delta\bigl((m_1-m_0)-((m_1-m_0)\cdot m_0)\one\bigr)-\delta'\one$ with $0<\delta'<\delta|m_1-m_0|^2$ if $m_1\cdot h=0$. This contradicts the first part. If $K_q$ misses $h^\perp$, it lies strictly on one side.
\end{proof}
\status{proved independently in \src{I2} (Lemma~1) and in \src{C2} (Lemma~3.1); each proof confirmed by one referee, different for the two notes (\src{I2\_referee} \S1, with repair R1: the case $m_1\cdot h<0$ was missing and the gamble for $m_1\cdot h=0$ needed the shift $-\delta'\one$, both as above; \src{C1C2\_referee} \S1). This is where the global form of the CF chart enters: $c(g)\cdot q>0$ excludes the silence cell at \emph{every} $q$, not only near $p^*$.}

Fix a world $b$ of $T_1$, with failing index $a$ and third index $c$, crossings $w'\in\relint F_a$ and $w\in\relint F_c$, and tangent lines $\ell'=\ell_{w'}$ and $\ell=\ell_w$ from $e_b$, as in \S\ref{sec:general-N28}. The \emph{tangent wedge} at $e_b$ is $\mathcal T_b\coloneqq\{m\in\Delta:m\cdot w'\ge0,\ m\cdot w\ge0\}$, the part of $\Delta$ between the two tangent lines on the side of $T_1$.

\begin{corollary}\label{cor:chord}
In every cover of $T_1$, for every $q\in\lambda_b\cap\Di$ with $q\ne p^*$: $K_q\subseteq\mathcal T_b$, $P(q,w')=P(q,w)=0$, and $K_q$ meets $\ell'$ and $\ell$ at points other than $e_b$.
\end{corollary}
\begin{proof}[Proof sketch]
On $\lambda_b$ the value $c_a(h^a(\tau))\cdot q$ vanishes at the crossing parameter with nonzero slope in $\tau$, so a strict straddle of $\ell'$ by $K_q$ would persist to nearby $\tau$ on the positive side, which Lemma~\ref{lem:nostraddle} forbids (the singleton case is excluded by Proposition~\ref{prop:vertexfree}(iv)). So $K_q$ lies on one closed side of $\ell'$, and likewise of $\ell$. Lemma~\ref{lem:quasiconcave} gives $P(q,w')\le0$ and $P(q,w)\le0$; every combination of sides other than the positive ones then puts a point of $\{m:m\cdot w'\le0,\ m\cdot w\le0\}=\{e_b\}$ into $K_q$, against Proposition~\ref{prop:vertexfree}(i).
\end{proof}
\status{proved (\src{C2} Cor.~3.3); confirmed by one referee (\src{C1C2\_referee} \S1), on both sides of $p^*$; NUMERICS agree (per-prior pencil MILP of \src{C2}, rerun by the referee). \src{TS\_second\_referee} \S4.3 found no proved statement in conflict with it. It corrects \S\ref{sec:general-open}.}

\begin{theorem}\label{thm:collar}
In every cover of $T_1$, normalised by $x=x_+$, fix $b$ with failing index $a$ and third index $c$.
\begin{enumerate}[leftmargin=*,label=(\alph*)]
\item There is a neighbourhood $V$ of $e_b$ such that for every non-precise $q\in V\cap\Di$ and $k\in\{a,c\}$: for $h\in\relint F_k$, $P(q,h)>0$ iff $c_k(h)\cdot q>0$; and $P(q,h)=0$ for the facet gamble $h\in\relint F_k$ whose chord passes through $q$.
\item There are $s_0>0$ and a connected open set $O_b\subseteq\Di$ of precise priors that contains a relative neighbourhood in $\Di$ of every edge prior $(1-s)e_b+se_a$, $0<s<s_0$.
\end{enumerate}
\end{theorem}
We call $O_b$ a \emph{precise collar} of the failing edge at $e_b$; it may be a horn, its width shrinking to $0$ at $e_b$.
\begin{proof}[Proof sketch]
(a) For non-precise $q$, the open sets $\{\tau:P(q,h^k(\tau))>0\}$ and $\{\tau:P(q,-h^k(\tau))>0\}$ cover the interval $J_k(q)\coloneqq\{\tau\in(0,1):c_k(h^k(\tau))\cdot q>0\}$ by Lemma~\ref{lem:nostraddle}, so one of them is all of it. Near $e_b$ it is the first: some $\tau^*$ has $h^k(\tau^*)_b>0$, $c_k(h^k(\tau^*))\cdot e_b>0$ and, since $K_q\to\{e_b\}$, $P(q,h^k(\tau^*))>0$. Off $J_k(q)$, $P\le0$ by Lemma~\ref{lem:quasiconcave}, and $P=0$ at the endpoint by continuity. $V$ is shrunk so that the chord parameters of both facets through every $q\in V$ lie in $(0,1)$. (b) The box certificate of Proposition~\ref{prop:T1} gives, for small $s$, $A\in\relint F_a$ and $B\in\relint F_c$ with $c_a(A)\cdot q_0>0$, $c_c(B)\cdot q_0>0$ at $q_0=(1-s)e_b+se_a$ and $A_bB_a<|A_a||B_b|$ (this uses $x_+^2>\kappa$). At non-precise $q$ near $q_0$, (a) gives $P(q,A)>0$ and $P(q,B)>0$, hence $P(q,A+\lambda B)>0$ for $\lambda$ with $(A+\lambda B)_a<0$ and $(A+\lambda B)_b<0$. Edge confinement puts every limit of $K_q$ at $q_0$ into $[e_a,e_b]$, where $m\cdot(A+\lambda B)<0$; so all priors near $q_0$ are precise. The union of these neighbourhoods is connected.
\end{proof}
\status{proved independently in \src{I2} (Thms.~2--3) and in \src{C2} (Cor.~3.2 and Thm.~4.1); each confirmed by one referee (\src{I2\_referee} \S\S2--3; \src{C1C2\_referee} \S2, which checks the certificate for $\varepsilon\in\{0.01,0.05,0.3,0.6,0.9,0.99\}$ and notes that $s_0$ shrinks as $\varepsilon\to1$). This is the box of Proposition~\ref{prop:T1} returning on general domains: cuts at non-precise priors cannot absorb its failure, and only precise priors can. The step that distinguishes $T_1$ from covered targets is the existence of the box certificate, which Lemma~\ref{lem:edgebox} forbids for half-plane charts (\src{I2\_referee} \S7).}

\paragraph{The d-line condition.} The next two theorems hold for any family $q\mapsto K_q$ of credal sets on $\Di$, under one condition on charts. Say that \textbf{(DL)} holds on $Q\subseteq\Di$ if there is a null set $\mathcal N\subseteq\R^3$ such that every mixed $g\notin\mathcal N$ satisfies one of: $P(\cdot,g)\le0$ on $Q$; $P(\cdot,-g)\le0$ on $Q$; or there is $d\in\R^3$ with $\operatorname{sign}d_i=\operatorname{sign}g_i$ ($i=1,2,3$) such that $\{q\in Q:P(q,g)>0\}\subseteq\{d\cdot q>0\}$ and $\{q\in Q:P(q,-g)>0\}\subseteq\{d\cdot q<0\}$. In words: for almost every pair, the priors at which accepting $g$ is best and those at which accepting $-g$ is best lie on opposite open sides of one line, whose normal has the strict sign pattern of $g$. No measurability of $g\mapsto d$ is required.

\begin{lemma}\label{lem:DLholds}
(DL) holds on $\Di$ for every solution of CF$(T)$, for every $T$ (take $d=c-c'$), and for the interior Bayes family of every strictly guiding additive ranking with integral penalties on a decomposable domain whose $p$-Bayes model is a.e.\ $D_{K_p}$ at every $p\in\Di$; in particular for every cover (Proposition~\ref{prop:necessary}(a)).
\end{lemma}
\status{proved (\src{PropD} Lemma~1.1(a),(b)); not yet refereed. For rankings, the pointwise penalty chart is open and convex and, by Fubini, a.e.\ equal to $\{P(\cdot,g)>0\}$; an open null difference is empty, which gives the exact inclusion, and strict guidance gives the signs of $d=\ell(g,-)-\ell(g,+)$. Not checked for bounded charge penalties (\src{PropD} Lemma~1.1(c)); see Corollary~\ref{cor:gmcharges}.}

\begin{theorem}\label{thm:D}
Let $q\mapsto K_q$ assign a credal set to each $q\in\Di$, let $O\subseteq\Di$ be open and connected with $K_q=\{\sigma(q)\}$ for $q\in O$ and $\sigma$ continuous on $O$, and suppose (DL) holds on $O$. Then either $\sigma$ is constant on $O$, or $\sigma=\sigma_w$ on $O$ for a single $w\in(0,\infty)^3$; in the second case $\sigma$ is injective on $O$.
\end{theorem}
\begin{proof}[Proof sketch]
\emph{Segment law.} If $q_1,q_2\in O$ are precise, $[q_1,q_2]\subseteq O$ and $q\in[q_1,q_2]$, then $K_q$ meets $[\sigma(q_1),\sigma(q_2)]$: otherwise some mixed $g\notin\mathcal N$ is positive at $\sigma(q_1),\sigma(q_2)$ and negative on $K_q$; then $P(q_1,g)>0$ and $P(q,-g)>0$, so (DL) would put $q_1,q_2$ on the positive open side of the d-line of $g$ and $q$ on the negative one. \emph{No segment images.} If $\sigma$ mapped a convex open $V\subseteq O$ onto a non-degenerate segment $I$, then for a.e.\ $m\in\relint I$ and a.e.\ $g\in m^\perp$ the fibre $\sigma^{-1}(m)\cap V$ would be the d-line of $g$ inside $V$; so all such $g$ would share one d-line and hence, by the strict signs of $d$, one pair-class $\pm\operatorname{sign}g$. But the pencil $m^\perp$ ($m$ not a vertex) meets at least two pair-classes in open cones. \emph{Constant or injective} on convex open sets follows from the segment law and the previous step by a triangle argument. \emph{Injective implies diagonal-projective}, without the fundamental theorem of projective geometry: by invariance of domain $\sigma$ is a homeomorphism onto an open set and maps lines into lines; the zero line $\Lambda(g)=\sigma^{-1}(\ell_g)$ depends continuously on $g$ and, being the d-line for a.e.\ $g$, has a normal with the strict sign pattern of $g$ for a.e.\ $g$, hence, by continuity, a zero coordinate where $g$ has one. So $\sigma$ preserves the pencils of lines through the vertices, $\sigma_1/\sigma_2=F_3(q_1/q_2)$ cyclically, and a Pexider equation gives $\sigma=\nu(w\odot q^{\alpha})$; collinearity of images forces $\alpha=1$. A connectedness argument makes $w$ global.
\end{proof}
\status{proved (\src{PropD} Thm.~D, Lemmas~SL, Z, 2--4); not yet refereed. It replaces \src{R4} Prop.~D, which \src{R4} recorded as CHECKED and \src{R8} found to have a GAP in the step ``non-constant $\Rightarrow$ injective'' (folds, plateaus and one-dimensional images were not excluded); that status is withdrawn, and R4 Prop.~D is cited in this note only through Theorem~\ref{thm:D} or the independent proofs below. Independently, the constant-or-injective dichotomy for covers and CF solutions was proved in two referee reports: Prop.~P (\src{C1C2\_referee} \S3) and Prop.~D$^*$ (\src{I2\_referee} \S4.2), each read by one reader only (its author). Each, combined with the injective case of \src{R4} Prop.~D (invariance of domain and the classical local fundamental theorem of projective geometry; CHECKED, written out as \src{C1} Lemma~5.0 and confirmed in \src{C1C2\_referee} \S5), gives the conclusion of Theorem~\ref{thm:D} for the regions used below. Theorem~\ref{thm:D} needs only (DL) and continuity of $\sigma$ on $O$: no (PC), (NS), target, regularity, or half-plane or three-ray form of the charts. NUMERICS (\src{PropD} \S8): images of segments under $\nu(w\odot q^\alpha)$ are collinear only at $\alpha=1$.}

\begin{theorem}\label{thm:T}
Suppose (DL) holds on $\Di$, and let $O$ be a precise region on which $\sigma=\sigma_w$. Then:
\begin{enumerate}[leftmargin=*,label=(\alph*)]
\item for a.e.\ mixed $g$ with $\ell_g\cap\sigma_w(O)\ne\varnothing$, the normal $d$ of (DL) can be taken positively proportional to $w\odot g$; hence on all of $\Di$, $\{P(\cdot,g)>0\}\subseteq\{q:(w\odot g)\cdot q>0\}$ and $\{P(\cdot,-g)>0\}\subseteq\{q:(w\odot g)\cdot q<0\}$;
\item every precise $q'\in\Di$, inside $O$ or not, has $K_{q'}=\{\sigma_w(q')\}$;
\item if $O'$ is another precise region on which $\sigma$ is not constant, then $\sigma=\sigma_{w'}$ on $O'$ with $w'$ positively proportional to $w$; and no open subset of $\Di$ carries a constant precise $K_q$.
\end{enumerate}
\end{theorem}
\begin{proof}[Proof sketch]
(a) $\sigma_w(O)$ is open, so it meets both sides of $\ell_g$, and on a convex ball around a point of $\sigma_w^{-1}(\ell_g)$ the d-line of $g$ is $\{(w\odot g)\cdot q=0\}$ with matching orientation. (b) For $q'$ precise with $K_{q'}=\{\sigma'\}$, (a) gives $\sigma'\cdot g>0\Rightarrow\sigma_w(q')\cdot g\ge0$ for every $g$ whose line meets $\sigma_w(O)$; applied to all $g\perp m_0$ for two points $m_0\in\sigma_w(O)$ not collinear with $\sigma'$, it forces $\sigma_w(q')=\sigma'$. (c) follows from (b) and injectivity of $\sigma_w$.
\end{proof}
\status{proved (\src{PropD} Thm.~T); not yet refereed. The same statements were proved as \src{I2} Prop.~4(b),(c) (one weight $w$ for all precise priors, at all three vertices), confirmed by one referee (\src{I2\_referee} \S4.3), and, for two regions with non-proportional weights, as \src{C2} Prop.~5.2, confirmed by one referee (\src{C1C2\_referee} \S4).}

\begin{theorem}\label{thm:A}
In every cover of $T_1$ there is no nonempty open $O\subseteq\Di$ on which $K_q=\{\sigma(q)\}$ with $\sigma$ injective. At the status of Theorem~\ref{thm:ccibu}, the same holds in every solution of CF$(T_1)$.
\end{theorem}
\begin{proof}[Proof sketch]
By Theorem~\ref{thm:D}, $\sigma=\sigma_w$ on each component of $O$. By Proposition~\ref{prop:vertexfree}(iv), $O$ avoids the three crossing chords, which cut $\Di$ into six open triangles; take $O$ connected, inside one of them, $C=\operatorname{int}\co\{p^*,e_b,f\}$, with two sides $s_1=(p^*,e_b)$, $s_2=(p^*,f)$ on chords. The lines that meet $O$ and cross both $s_1$ and $s_2$ form a nonempty open set. A.e.\ such line is $L_g\coloneqq\{(w\odot g)\cdot q=0\}$ for a good $g$, and since both sides of $L_g$ meet $O$, $L_g\cap O$ lies in the closed d-face of $g$; so the d-line of $g$ is $L_g$. The closed d-face is a segment of that line reaching $\partial\Delta$, so it contains the crossing point $y_1(g)\in s_1$ or $y_2(g)\in s_2$, where $K_{y}\subseteq\ell_g$. \emph{Measurable fibration (repaired):} the sets $G'_i\coloneqq\{g \text{ good}:P(y_i(g),g)=0=P(y_i(g),-g)\}$ are Borel, since $P$ is continuous, and one of them, say $G'_1$, is non-null; its image under $g\mapsto(y_1(g),\text{direction of }L_g)$, a submersion, is analytic and of positive measure, so some $y\in s_1$ carries two gambles $g,g'$ with $\ell_g\ne\ell_{g'}$ and $K_y\subseteq\ell_g\cap\ell_{g'}$. Then $y\in\lambda_b$ is precise, against Proposition~\ref{prop:vertexfree}(iv).
\end{proof}
\status{proved (\src{C1} Thm.~A); confirmed with repair by one referee (\src{C1C2\_referee} \S5): Step~5 needed measurability of the fibred sets, supplied by the Borel sets $G'_i$ defined through $P(y,\pm g)=0$, as above. \src{C1} relied on its Lemma~5.0 (the injective case of \src{R4} Prop.~D, CHECKED) for $\sigma=\sigma_w$; Theorem~\ref{thm:D} now supplies it without the fundamental theorem of projective geometry. The input Proposition~\ref{prop:vertexfree}(iv) is proved for covers, so for bare CF solutions the statement has the status of Theorem~\ref{thm:ccibu}. NUMERICS illustration at $\varepsilon=0.3$: 1283 lines through one chord point and admissible cap points leave the chamber through both chords (\src{C1} \texttt{thmA\_cap\_geometry.out}).}

\begin{corollary}\label{cor:nocap}
The referee's sector-cap example (\src{N28\_referee}, counterexample C1: the identity $K_q=\{q\}$ on a sector at $e_b$ next to the failing edge), the identity on any open set, and every diagonal-projective precise region anywhere in $\Di$, are part of no cover of $T_1$.
\end{corollary}
\status{proved (\src{C1} Cor.~5.1), from Theorem~\ref{thm:A}, since $\sigma_w$ is injective and the identity is $\sigma_{\one}$; confirmed (\src{C1C2\_referee} \S\S4--5). Proposition~\ref{prop:vertexfree}(iv) had excluded only regions meeting a chord. The facet gambles tested in \src{N28\_referee} do not detect the obstruction; generic gambles whose lines leave the chamber through both chords do (\src{C1} \S5).}

\begin{remark}\label{rem:B}
Before the routes above were completed, \src{C1} proved that CF$(T_1)$ has no \emph{locally definable} solution: none whose graph $\{(q,m):q\in Q,\ m\in K_q\}$ is definable in an o-minimal expansion of the real field for every compact $Q\subseteq\Di$ (for instance semialgebraic or piecewise-projective solutions, or solutions with finitely many pieces on each compact set). The proof combines Theorem~\ref{thm:A}, the absence of $C^1$ rank-one precise regions, and a dimension count on the d-faces of adjacent pairs, which are non-negligible in every solution because Theorem~\ref{thm:main} excludes solutions whose pairs are a.e.\ split.
\end{remark}
\status{proved (\src{C1} Thm.~B, at the status of Proposition~\ref{prop:vertexfree}(iv) for CF solutions and Theorem~\ref{thm:main} with \src{N14} \S3(iv)); confirmed with three repairs by one referee (\src{C1C2\_referee} \S5: definability of the CF chart on each compact, dimension counts on projective classes, and the finite-segment fact in case $k=1$). \emph{Superseded} by Theorem~\ref{thm:gm} and Corollary~\ref{cor:CFnone}, which need no definability. With Prop.~P or Theorem~\ref{thm:D}, \src{C1}'s set $P_{\mathrm{wild}}$ is empty and every open precise region of a CF$(T_1)$ solution is constant (\src{C1C2\_referee} \S3, \src{PropD} \S5); \src{C1} \S7 (wild solutions, Cantor fans) is moot.}

\subsection{Second and third proofs}\label{sec:gm-other}

\begin{proof}[Second proof of Theorem~\ref{thm:gm} (route I2, sketch)]
Normalise $p^*=u$, $x=x_+$, and fix $b=2$, so $(a,c)=(1,3)$. Theorem~\ref{thm:collar}(b) gives a precise collar $O_b$. By Theorem~\ref{thm:D} (or Prop.~D$^*$ of \src{I2\_referee}) and Lemma~\ref{lem:DLholds}, $\sigma$ is constant or $\sigma_w$ on $O_b$; a constant would be $e_b$, because $O_b$ accumulates at $e_b$ and $K_q\to\{e_b\}$, and this is excluded by Proposition~\ref{prop:vertexfree}(i). So $\sigma=\sigma_w$ on $O_b$, and by Theorem~\ref{thm:T} every precise prior anywhere has $K_q=\{\sigma_w(q)\}$, with the same $w$ at all three vertices.

\emph{Transition points.} For $q\in V$ (Theorem~\ref{thm:collar}(a)) let $h^a_\partial(q)\in\relint F_a$ and $h^c_\partial(q)\in\relint F_c$ be the facet gambles whose chords pass through $q$, and $\Phi_b(q)\coloneqq\ell_{h^a_\partial(q)}\cap\ell_{h^c_\partial(q)}$, the intersection of their zero lines; it is constant along lines through $p^*$. A \emph{transition point} is a precise prior in $V$ that is a limit of non-precise priors. At such a point, Theorem~\ref{thm:collar}(a) and joint continuity of $P$ give $K_q=\{\Phi_b(q)\}$. On every small arc around $e_b$ from $\lambda_b$ (non-precise, Proposition~\ref{prop:vertexfree}(iii)) to the collar, the first precise point is a transition point $q_r$. So $G_b\coloneqq\sigma_w-\Phi_b$, which is real-analytic near $e_b$ with $G_b(e_b)=0$, vanishes at points $q_r\to e_b$ inside the closed sector between $\lambda_b$ and the failing edge.

\emph{Blow-up.} With $q=e_2+t(y_1,\cdot,y_3)$, the linearisation of $\Phi_b$ at $e_b$ is $y\mapsto(y_1-y_3)v$ with
\[
v=\frac{\kappa}{1-\kappa^2}\Bigl(\frac1x-\kappa x,\ \frac\kappa x-x\Bigr),\qquad v_1,v_3>0,
\]
and that of $\sigma_w$ is $y\mapsto(\omega_ay_1,\omega_cy_3)$ with $\omega_a=w_1/w_2$, $\omega_c=w_3/w_2$. Zeros of $G_b$ accumulating at $e_b$ inside the closed sector force $\det DG_b=\omega_a\omega_c+\omega_av_3-\omega_cv_1=0$ with kernel in the open sector, i.e.\ $1/\omega_c=G(\omega_a)\coloneqq(v_1-\omega_a)/(v_3\omega_a)$ with $0<\omega_a<v_1$. By cyclic symmetry the three vertices give $r_2=G(r_1)$, $r_3=G(r_2)$, $r_1=G(r_3)$ with $r_1=w_1/w_2$, $r_2=w_2/w_3$, $r_3=w_3/w_1$, so $r_1r_2r_3=1$. The M\"obius map $G$ has matrix $\bigl(\begin{smallmatrix}-1&v_1\\v_3&0\end{smallmatrix}\bigr)$, with real eigenvalues of opposite sign and unequal modulus, so $G^3$ fixes only the fixed points of $G$; with $r_1r_2r_3=1$ this forces $r_1=r_2=r_3=1$, which needs $v_1-v_3=1$. But $v_1-v_3=\kappa(x+1/x)/(\kappa+1)=\kappa$ by the corner equation, and $\kappa>1$.
\end{proof}
\status{proved (\src{I2} Lemma~1, Thms.~2, 3, Prop.~4, Lemma~5, Thm.~6); confirmed by one referee (\src{I2\_referee} \S\S1--6, with repair R1 to Lemma~1 and Prop.~D$^*$ replacing the citation of \src{R4} Prop.~D, whose status \src{I2} had misstated). The blow-ups were recomputed exactly at all three vertices without using the cyclic symmetry (\src{I2\_referee} \S6, \texttt{blowup\_exact.out}). The global weight of Theorem~\ref{thm:T} is essential: any two of the three relations are solvable at a unique positive $w$, and the third then fails (NUMERICS, \texttt{solve\_system.out}, \texttt{third\_relation.out}). The referee's sector cap has $w=\one$ and $\det DG_b=1-\kappa\ne0$, so it cannot meet a non-precise prior at its boundary.}

\begin{proof}[Third proof of Theorem~\ref{thm:gm} (route C2, sketch)]
As in the second proof, each collar $O_b$ carries $\sigma=\sigma_{w^{(b)}}$, constant being excluded. \emph{Gluing.} Every boundary point $q$ of $O_b$ close to $e_b$ is a limit of non-precise priors, at which $K_q$ touches the zero lines of the two facet gambles whose chords pass through $q$ (Theorem~\ref{thm:collar}(a)); so $Q^k(q)\coloneqq(w^{(b)}\odot q)\cdot H_k(q)=0$ for $k\in\{a,c\}$, where $H_k(q)$ is that facet gamble of $F_k$, linear in $q$ by (PC). Such points exist at every small distance from $e_b$ (arc argument), so the two conics $\{Q^a=0\}$ and $\{Q^c=0\}$ have infinitely many common zeros accumulating at $e_b$. \emph{Pinning.} The two quadratic forms are never proportional, so a common component is a line, and since the common zeros accumulate at $e_b$ it passes through $e_b$; solving the common-line system by hand for every $x>1$, the only positive solution (with $w^{(b)}_b=1$) is $w^{(b)}$ equal to $\kappa$ off $b$ and $1$ at $b$, the common line having slope $1/x$. The other algebraic candidate, $w=(x,1,1/x)$ at $b=2$, makes the conics share the line $x^2q_1-(x^2+1)q_2+q_3=0$, which misses $e_b$; it is excluded by the accumulation at $e_b$ alone. \emph{Contradiction.} So $w^{(1)}\propto(1,\kappa,\kappa)$ and $w^{(2)}\propto(\kappa,1,\kappa)$, which are not proportional since $\kappa\ne1$; this contradicts Theorem~\ref{thm:T}(c) (\src{C2} Prop.~5.2: a gamble whose line meets both images would need its d-line to be both $\{(w^{(1)}\odot g)\cdot q=0\}$ and $\{(w^{(2)}\odot g)\cdot q=0\}$).
\end{proof}
\status{proved (\src{C2} Lemma~3.1, Cor.~3.2, Thm.~4.1, Lemma~4.2, Prop.~4.3, Prop.~5.2, Thm.~5.1); confirmed with repairs by one referee (\src{C1C2\_referee} \S\S1--4, route A). Repairs: Prop.~P replaces \src{C2}'s citation of the unrepaired \src{R4} Prop.~D; candidate (I) of \src{C2} Prop.~4.3 is a shared line not through $e_b$, not a coincidence of the conics, and needs no sign test; the referee's hand elimination for all $x>1$ makes the exceptional finite set of $\varepsilon$ in \src{C2} \emph{empty}, so the route holds for every $\varepsilon\in(0,1)$. \src{C2} \S2 (finite discretisations DCF, DCF$^+$) is NUMERICS: finite systems are feasible, and their solutions are artefacts of the finite grids.}

\begin{proof}[Short route (sketch)]
Theorem~\ref{thm:collar}(b) gives a precise open collar $O_b$. By Theorem~\ref{thm:D} with Lemma~\ref{lem:DLholds}, $\sigma$ is constant or $\sigma_w$ on $O_b$; constant is excluded by Proposition~\ref{prop:vertexfree}(i), as in the second proof. So $\sigma$ is injective on the open set $O_b$, which Theorem~\ref{thm:A} forbids.
\end{proof}
\status{proved, at the status of its three new inputs: Theorem~\ref{thm:collar} (two independent proofs, each refereed once), Theorem~\ref{thm:D} (not yet refereed; alternatively Prop.~P or Prop.~D$^*$ with the CHECKED injective case of \src{R4} Prop.~D, one reader each) and Theorem~\ref{thm:A} (one referee). Observed in \src{PropD} \S6 and, with Prop.~P, in \src{C1C2\_referee} \S6 (route B). It uses no conic algebra, no Bezout argument and no transition algebra, and does not raise the status above that of its inputs.}

\subsection{Dependencies}\label{sec:gm-deps}

Table~\ref{tab:gmdeps} lists what Theorem~\ref{thm:gm} rests on. ``Two readings'' means that two referees, independently, re-derived the statement and confirmed it; ``one reading'' means one referee; a result proved in a referee report counts as read only by its author unless a second reading is recorded.

{\small
\begin{longtable}{>{\raggedright\arraybackslash}p{5.2cm}>{\raggedright\arraybackslash}p{2.2cm}>{\raggedright\arraybackslash}p{6.7cm}}
\caption{Inputs and new steps of Theorem~\ref{thm:gm}.}\label{tab:gmdeps}\\
\toprule
Input or step & Routes & Readings and status\\
\midrule
\endfirsthead
\toprule
Input or step & Routes & Readings and status\\
\midrule
\endhead
\multicolumn{3}{l}{\emph{Imported inputs}}\\
Prop.~\ref{prop:necessary}(a); Hausdorff continuity (\src{R1} Prop.~3.1), joint continuity of $P$ & all & proved; refereed once (\src{R6}; \src{N14})\\
Prop.~\ref{prop:CF}($\Rightarrow$), every-$q$ chart form & all & two readings, with repairs (\src{N28\_referee} \S5.1, \src{TS\_second\_referee} \S5)\\
Thm.~\ref{thm:TSalways} and its inputs (\src{N21.9}, \src{N21.4} items 1--5, \src{R10b} UVN/D/M/PL2/PL, Lem.~\ref{lem:R231}, Cor.~\ref{cor:R232}) & all & two readings (\src{N28\_referee}/\src{N23}/\src{R12}; \src{TS\_second\_referee} \S1); N21.4 items 6--8 (localised N15 chain, itself refereed twice) only spot-checked in the second reading\\
Cor.~\ref{cor:TSalways}(b),(c): both acceptances; $K_q\to\{e_b\}$; edge confinement & all & two readings (\src{N28\_referee}, \src{N23}; \src{TS\_second\_referee} \S2)\\
Thm.~\ref{thm:cuts} & all (via Cor.~\ref{cor:TSalways}) & two readings (\src{N23} \S6; \src{TS\_second\_referee} \S3.1)\\
Thm.~\ref{thm:apex}, with the wording repair & all (via Prop.~\ref{prop:vertexfree}) & two readings (\src{N28\_referee} \S4.1; \src{TS\_second\_referee} \S3.2)\\
Lem.~\ref{lem:R1N28} & none directly & two readings (\src{N28\_referee}, author; \src{TS\_second\_referee} \S4.1); it gives the strict forms in Prop.~\ref{prop:vertexfree}(ii),(iii), while the routes use only the weak forms (Lem.~\ref{lem:quasiconcave})\\
Prop.~\ref{prop:vertexfree} (i)--(iv) & all & two readings (\src{N28\_referee} \S4.3; \src{TS\_second\_referee} \S4.2)\\
Lems.~\ref{lem:quasiconcave}, \ref{lem:B1}, \ref{lem:walls} & all & two readings (\src{N23}, \src{N28\_referee}; \src{TS\_second\_referee} \S1.1)\\
Lem.~\ref{lem:veccluster} (vector clusters) & all & from Thm.~\ref{thm:TS} item 5 (two readings); vector form written out in \src{I1\_referee} and re-derived in \src{I2\_referee}, \src{C1C2\_referee}, \src{TS\_second\_referee}\\
Crossing rows, corner equation, box certificate (Props.~\ref{prop:boxbridge}(a),(b), \ref{prop:T1}) & all & refereed (\src{N14}, \src{R18}; assembly \src{N23} \S6); main chain of \S\ref{sec:main} refereed twice\\
$K_q\subseteq g^\perp$ on closed d-faces (\src{R4} Thm.~B) & all & refereed once (\src{R8}); CF form re-derived in \src{I1\_referee}, \src{I2\_referee}, \src{C1} Lemma~4.1 (confirmed in \src{C1C2\_referee})\\
Normalisations ($p^*=u$, $x=x_+$) & all & elementary; re-checked in \src{I1\_referee} \S5\\
\midrule
\multicolumn{3}{l}{\emph{New steps}}\\
Lems.~\ref{lem:edgemaps}, \ref{lem:forcedlines}, \ref{lem:window}, Props.~\ref{prop:edgelimits}, \ref{prop:clines}; first proof & I1 & one reading (\src{I1\_referee}; repairs R1--R7, expository)\\
Lem.~\ref{lem:nostraddle}, Thm.~\ref{thm:collar} & I2, C2, short & proved independently in \src{I2} and \src{C2}; each proof read once, by different referees\\
Cor.~\ref{cor:chord} & (corrects \S\ref{sec:general-open}) & one reading (\src{C1C2\_referee})\\
Thm.~\ref{thm:D}, Lem.~\ref{lem:DLholds} & short (or I2, C2) & not yet refereed (\src{PropD}); the constant-or-injective dichotomy proved independently in Prop.~P (\src{C1C2\_referee}) and Prop.~D$^*$ (\src{I2\_referee}), one reader each, plus the injective case of \src{R4} Prop.~D (CHECKED; classical)\\
Thm.~\ref{thm:T} & I2, C2 & not yet refereed (\src{PropD}); same content as \src{I2} Prop.~4(b),(c) and \src{C2} Prop.~5.2, each read once\\
Transition points and blow-up (second proof) & I2 & one reading (\src{I2\_referee})\\
Gluing, pinning, two caps (third proof) & C2 & one reading (\src{C1C2\_referee}), with repairs\\
Thm.~\ref{thm:A} & short & one reading (\src{C1C2\_referee}), with repair\\
\midrule
\multicolumn{3}{l}{\emph{Only for the corollaries}}\\
Thm.~\ref{thm:ccibu} (Prop.~\ref{prop:CF}($\Leftarrow$)) & Cor.~\ref{cor:CFnone} & author-checked; no formal referee report\\
\src{N14} Thm.~C2 & Cor.~\ref{cor:gmcharges} & proved in a referee report, one reader; splice step refereed in \src{N24}\\
\bottomrule
\end{longtable}}

The once-refereed inputs on which all three routes depend are Proposition~\ref{prop:necessary} together with the continuity statement \src{R1} Prop.~3.1, \src{R4} Thm.~B, and the crossing-row and box-certificate assembly of \src{N23} \S6; everything else that the routes import has now had two independent readings. Before the theorem is promoted from one reader to two, a second reading of one complete route is needed; the short route is the shortest, and its only unrefereed link is Theorem~\ref{thm:D}, for which Prop.~P or Prop.~D$^*$ can substitute.

\section{Coherent almost-preferences and rankings that add up}\label{sec:cp}

Standing hypotheses \textbf{(H)}: $n=3$; a decomposable domain $\M$; the ranking satisfies strict (G), (IB), (LA), (S), (C), and covers $T_1$ with (LF). Let $N_i$ be the set of problems whose advice range has at least two elements, i.e.\ on which evaluable models can give different advice (on $\M_U$: the pairs $\{f,-f\}$ with $f\in U$ or $-f\in U$, not only those whose representative in $H$ lies in $U$). Under (CP-a) each state splits as $P=P^c+P^s$ (Yosida--Hewitt) with $P^c=\rho_P\mu$, $\rho_P\in L^1_+$, and $P^s$ purely finitely additive; any finitely many purely finitely additive parts have a common carrier of arbitrarily small $\mu$-measure (\src{N14} Lemma Y; \src{N11} Fact 1.2).

\subsection{Normal form and full support}

\begin{proposition}\label{prop:normalform}
Under (CP-a), $\preceq_i$ is a preorder, (IB) and weak (G) hold automatically, and every $P\in\M_i$ weakly respects the ranking: $X\preceq_iY\Rightarrow P(\Lambda_iX)\le P(\Lambda_iY)$ (property (W)). Strict (G) at $i$ holds iff $\up_i(1_A)\coloneqq\max_{P\in\M_i}P(1_A)>0$ for every non-null $A\subseteq N_i$.
\end{proposition}
\status{proved; refereed once each (\src{N13} for \src{N11} Prop.~1.1; \src{N24} for \src{N22} Lemma 1.0).}

\begin{theorem}\label{thm:fullsupport}
Under (CP-a) and strict (G) at world $i$, the set $\mathcal P^+_i\coloneqq\{P\in\M_i:\rho_P>0\text{ a.e.\ on }N_i\}$ is nonempty, and $(1-t)P+tQ\in\mathcal P^+_i$ whenever $P\in\mathcal P^+_i$, $Q\in\M_i$, $t\in[0,1)$.
\end{theorem}
\begin{proof}[Proof sketch]
Key lemma: if a weak-* compact convex set $\M$ of states has $P^c(A)=0$ for every $P\in\M$ and some non-null $A$, there is a non-null $N\subseteq A$ with $P(N)=0$ for all $P\in\M$. For one step, Kneser's minimax theorem applied to $(P,f)\mapsto P(f)$ on $\M\times\{0\le f\le1_{A'},\int f=(1-\eta^2)\mu(A')\}$ gives a large $B\subseteq A'$ of uniformly small $\M$-mass, because each single member has small carriers. Iteration and countable additivity of $\mu$ then give $N$. With the countable chain condition, the essential union of the supports of the $\rho_P$ must cover $N_i$, by Proposition~\ref{prop:normalform}; a countable convex combination of members (normalised partial sums, norm-closed) gives an element of $\mathcal P^+_i$.
\end{proof}
\status{proved; refereed once (\src{N24} \S1: CORRECT WITH REPAIR, cosmetic). Source: \src{N22} Lemma 1.1, Thm.~1.2. Of independent interest: a weak-* compact convex set of states on $L^\infty(\mu)$ that charges every non-null set contains a state whose normal part is equivalent to $\mu$.}

\subsection{The integral reduction and the residual (U)}

For $P=(P_1,P_2,P_3)\in\prod_i\mathcal P^+_i$ define the integral field $\psi^P_i(X)\coloneqq\int\lambda_i(g,a_X(g))\,\rho_{P_i}(g)\,d\mu(g)$, let $Z_p(P)$ be the set of $g$ at which $p\cdot\psi^P(g,\cdot)$ has two or more minimisers on $A(g)$, and put $E(P)\coloneqq\{p\in\Di:\mu(Z_p(P))>0\}$.

\begin{theorem}\label{thm:integral}
Assume (H) and (CP-a). For every $P\in\prod_i\mathcal P^+_i$, $\psi^P$ is a strictly guiding integral field on the same domain $\M$ (so (LF) at $T_1$ is kept), and:
\begin{enumerate}[leftmargin=*,label=($\alph*$)]
\item $T_1$ is $\psi^P$-Bayes at some $p_T\in\Di$;
\item for every $p\in\Delta$ some $\psi^P$-$p$-Bayes model is coherent;
\item $E(P)$ is Lebesgue-null, and at every $p\in\Di\setminus E(P)$ the $\psi^P$-$p$-Bayes model is a.e.\ unique and coherent;
\item at each vertex the Bayes model is the world-best evaluable model, which is coherent.
\end{enumerate}
\end{theorem}
\begin{proof}[Proof sketch]
($a$) $T_1$ is competitive on $\Omega$ (\src{R11} Lemma IC, refereed in \src{R13}). Splicing off a common small carrier of the three singular parts shows that no model is better than $T_1$ for all three integral parts; Lyapunov and separation give a weight, which the Lemma IC construction pushes into $\Di$. ($b$) Approximate the infimum of the full charge objective by models equal to a pointwise integral minimiser off the carriers; by Theorem~\ref{thm:kkm} select weak dominators; by (W) for the full members they are near-minimisers of the integral objective; compactness in measure of coherent models gives a coherent minimiser. ($c$) Tonelli on $H\times\Delta$; off $E(P)$ the minimiser is pointwise unique, so it is the coherent one. ($d$) Strict guidance and the vertex lemma.
\end{proof}
\status{proved; refereed once (\src{N24} \S2). Source: \src{N22} Thm.~2.1, building on \src{N11} Cor.~3.2 and \src{N13} Prop.~B (refereed in \src{N13}, \src{N24}). No detection, attainment or non-atomicity hypothesis is used; atoms enter only through finitely many small carriers.}

\begin{corollary}\label{cor:residual}
Assume (H) and (CP-a) on a decomposable domain non-adjacent near $T_1$. Then $E(P)\ne\varnothing$ for every $P\in\prod_i\mathcal P^+_i$. Equivalently, $T_1$ cannot be covered once the following holds:
\begin{quote}
\textbf{(U)} some $P\in\prod_i\mathcal P^+_i$ has no non-null pointwise Bayes tie at any interior weight.
\end{quote}
\end{corollary}
\begin{proof}
If $E(P)=\varnothing$, then (on $\M_U$; on other non-adjacent domains after the reduction used in Corollary~\ref{cor:coverage}) the vector $c_i(g)\coloneqq\rho_{P_i}(\bar g)\bigl(\lambda_i(\bar g,0)-\lambda_i(\bar g,a_g)\bigr)$, $g\in U$, where $\bar g$ is the representative in $H$ of the pair of $g$ and $a_g\in\{+,-\}$ is the advice ``accept $g$'', satisfies (BU) at \emph{every} interior weight and (TP) at $p_T$, which Corollary~\ref{cor:coverage} excludes. Uniqueness off a null set of weights would not suffice (Remark~\ref{rem:N20}).
\end{proof}
\status{proved; refereed once (\src{N24} \S2.2, which checked that the every-weight form of (BU) is what is supplied). Source: \src{N22} Cor.~2.3.}

\begin{theorem}\label{thm:complete}
Assume (H), and suppose every world is complete: $\M_i=\{P_i\}$ for states $P_i$ with arbitrary purely finitely additive and atomic parts. Assume \textbf{(UN)}: for each $i$ there is a non-null $U_+\subseteq N_i$ on which every coherent model of $\M$ gives the advice in the advice range that is best at world $i$ (true on $\M_U$ with $U_+=U\cap\{g\gg0\}$, and on the all-model domain). Write $\varphi_i\coloneqq P_i\circ\Lambda_i$. Then:
\begin{enumerate}[leftmargin=*,label=(\alph*)]
\item for every $p\in\Delta$, $\inf_\M p\cdot\varphi$ is attained by a selected model;
\item for every $p\in\Di$ the minimiser is a.e.\ unique, coherent and selected;
\item $P\in\prod_i\mathcal P^+_i$ and $E(P)=\varnothing$.
\end{enumerate}
\end{theorem}
\begin{proof}[Proof sketch]
(a) Take a minimising sequence, selected weak dominators $C_{q_m}$ (Theorem~\ref{thm:kkm}), and subsequences with $q_m\to q^*$, $C_{q_m}\to C_\infty$ coherent and $\varphi(C_{q_m})\to z$. Worsening $C_m$ off a small carrier shows that the attainable world-$i$ values contain intervals of length bounded below, which uses (UN); so there are $X$ with $\varphi_i(X)$ just above $z_i$. Completeness turns $\varphi_i(C_{q_m})<\varphi_i(X)$ into $C_{q_m}\preceq_iX$, and the closed sections of (S) give $C_{q^*}\preceq_iX$. Hence $\varphi(C_{q^*})\le z$ and $C_{q^*}$ attains. (b) A selected weak dominator of a minimiser has equal profile; completeness makes this a tie at every world, and (C) gives coherence. Uniqueness follows from complementary splices and a bounded non-null window. (c) Ties of $\varphi$ are ties of $L$, so detection is automatic.
\end{proof}
\status{proved; refereed once (\src{N24} \S3: CORRECT; window-radius repair incorporated). Source: \src{N22} Thm.~3.1.}

\begin{corollary}\label{cor:complete}
On the $U$-domain, and on every decomposable domain non-adjacent near $T_1$ that satisfies (UN), no ranking satisfying strict (G), (IB), (LA), (S), (C) whose worlds are complete coherent almost-preferences over losses covers $T_1$. Equivalently, no additive ranking with bounded charge penalties vanishing on null sets covers $T_1$ there; the existence proviso of Corollary~\ref{cor:charges} is discharged by (S).
\end{corollary}
\status{proved; refereed once (\src{N24} \S3). Source: \src{N22} Cor.~3.2, using Theorem~\ref{thm:main}.}

\begin{corollary}\label{cor:DCA}
Assume (H), (CP-a) and countably additive detection: whenever $X\prec_iY$, some countably additive $Q\in\M_i$ has $Q(\Lambda_iY-\Lambda_iX)>0$ (for example, every member of $\M_i$ is countably additive, or $\M_i$ is the weak-* closure of its countably additive members). Then the countable aggregate of a norm-dense sequence of countably additive members is an integral field with a.e.-unique coherent Bayes models at every interior weight and $T_1$ among them. Consequently $T_1$ is not covered on decomposable domains non-adjacent near $T_1$.
\end{corollary}
\status{proved; refereed once (\src{N13} for \src{N11} Thm.~2.1, which uses \src{R11} Thm.~W with \src{R13} Prop.~S; \src{N24} \S3.3 for the application of Corollary~\ref{cor:residual}).}

\begin{proposition}\label{prop:Ucont}
On $\M_U$, fix $P\in\prod_i\mathcal P^+_i$ and let $K_q$ be the credal sets of the unique coherent $\psi^P$-Bayes models for $q\in\Di\setminus E(P)$. Then $q\mapsto K_q$ is Hausdorff-continuous on $\Di\setminus E(P)$, and $p_0\notin E(P)$ iff $K$ has a limit at $p_0$ along $\Di\setminus E(P)$. So (U) holds iff some member choice has a coherent Bayes family that extends continuously to all of $\Di$.
\end{proposition}
\status{proved in a referee report (\src{N24} \S6.1, Prop.~R1), one reader only; the non-adjacent version is stated there at the level of an expectation and was not re-checked.}

\begin{remark}\label{rem:EPchords}
On $\M_U$ the field $\psi^P$ is strictly signed and integrable, and its Bayes models are coherent at a.e.\ interior weight (Theorem~\ref{thm:integral}($c$)). So Theorem~\ref{thm:Estructure} below applies to it: $E(P)$ is a countable union of open chords of $\Delta$, each the common chart line of a non-null set of gambles and each cutting off one vertex, with no isolated exceptional points. Hence on $\M_U$, (U) fails exactly when every member choice $P$ has such a chord. (This application is ours; the sources state Theorem~\ref{thm:Estructure} for the M0 setting, but its proof uses only the properties just listed.) In the same way Theorem~\ref{thm:M0closed}(a) applies to the strictly signed field $c^P$ underlying $\psi^P$, so on $\M_U$ no (CP-a) ranking covers $T_1$ (Corollary~\ref{cor:CPaU}). Because (H) includes coverage, the structural statement ``(U) fails exactly when every member choice $P$ has such a chord'', and Proposition~\ref{prop:Ucont}, are therefore vacuous on $\M_U$ as stated. Their content survives for (CP-a) systems that satisfy every axiom of (H) except coverage: by \src{N24} \S2.1 only Theorem~\ref{thm:integral}($a$) uses coverage, so ($b$)--($d$), $E(P)$ and the chord description make sense for such systems (\src{CPaU\_referee} \S4, repair R9; this reading has not been re-checked step by step).
\end{remark}

\begin{proposition}\label{prop:CD}
If every $\M_i$ has a countable subfamily detecting all strict comparisons (for example, $\M_i$ weak-* separable: finitely generated, contained in the countably additive states, or the closed convex hull of a norm-separable set), there is $\hat P\in\prod_i\mathcal P^+_i$ with $E(\hat P)$ contained in the set of interior weights at which $\inf_\M p\cdot(\hat P\circ\Lambda)$ is not attained. Hence countable detection together with attainment of one charge field at every interior weight excludes a covering on non-adjacent domains.
\end{proposition}
\status{proved in a referee report (\src{N24} \S6.2, Prop.~R2), one reader only. Attainment at every interior weight is open even for finitely generated incomplete $\M_i$ (\src{N22} Remarks 3.3).}

At a tie weight a hypothetical cover has a specific anatomy: if $P$ consists of countably additive members, the face of Bayes profiles at $p_0\in E(P)$ has dimension $\ge1$, is realised by selected models, and carries a strict comparison that $P$ does not see (\src{N22} Prop.~4.1, refereed with repair in \src{N24} \S4).

\begin{conjecture}\label{conj:U}
(U) holds for every (CP-a) system satisfying (H). On non-adjacent domains this is equivalent to: no (CP-a) ranking covers $T_1$.
\end{conjecture}
\status{open on non-adjacent domains other than $\M_U$ (\src{N22} Conj.~N22.U). It subsumes the earlier detection, atom and existence questions (\src{N11} Conj.~A, B; \src{N14} existence proviso). On the $U$-domain its hypothesis is never met: (H) includes coverage of $T_1$, and no (CP-a) ranking covers $T_1$ there (Corollary~\ref{cor:CPaU}, proved; glue refereed once in \src{CPaU\_referee}). So on $\M_U$ the conjecture, and Corollary~\ref{cor:residual}, hold vacuously, and its equivalent form ``no (CP-a) ranking covers $T_1$'' is what has been proved. This is a statement about the conditional, not about (U): (U) is a property of a system, and Corollary~\ref{cor:CPaU} does not show that (U) holds for any (CP-a) system on $\M_U$; its proof tolerates $E(P)\ne\varnothing$. On other non-adjacent domains (including (UN) domains) the conjecture is open (\src{CPaU\_referee} \S\S4, 5, repair R5).}

\subsection{Strict preference}
Walley's strict-preference model \cite[\S\S3.7.7--3.7.9]{W1991}, read in a preorder framework as $\prec\cup{=}$, gives condition (CP-s): $X\preceq_iY$ iff $X=Y$, or $\lp_i(\Lambda_iY-\Lambda_iX)>0$, or $\Lambda_iY-\Lambda_iX\ge0$ a.e.\ and $\ne0$. Walley 2000 \cite[\S6]{W2000} gives (CP-b): $X\preceq_iY$ iff $\Lambda_iY-\Lambda_iX\in\mathcal D_i\cup\{0\}$ for a coherent set $\mathcal D_i$ of desirable gambles (D1--D4); (CP-s) is the case $\mathcal D_i=\{\lp_i>0\}\cup\{f\ge0,f\ne0\}$.

\begin{proposition}\label{prop:A}
Let two worlds $\omega\ne j$ satisfy (CP-s) with weak-* compact sets of countably additive states with a.e.-positive densities, on $\M_U$ (or any decomposable domain whose advice ranges contain, for every selected model, advice worse at every world than the model's own on a non-null set of unanimous problems, and advice better at $\omega$ than the model's own wherever the model errs at $\omega$), and suppose the world-ideal models $W_\omega\ne W_j$ on a non-null set. Then no selection $u:\Delta\to\C$ satisfies (S), even without (LA), (C) or a target.
\end{proposition}
\begin{proof}[Proof sketch]
Under strict positivity, $Y\prec Z$ iff $\lp(\Lambda Z-\Lambda Y)>0$ at these worlds, so upper sets $\{Y:C\prec Y\}$ are open (Walley's openness axiom). For a selected $C^*=C_{p_0}$ build an incoherent $X$ with $\lp_\omega(\Lambda_\omega X-\Lambda_\omega C^*)=0$; then $p_0\notin\{p:C_p\preceq_\omega X\}$, which is closed, while every $p$ with $C_p\prec_\omega C^*$ lies in it. So near $p_0$ on the edge $[e_\omega,e_j]$ no selected model is strictly better than $C^*$ at $\omega$ or $j$; competitiveness forces $C_p=C^*$ there. The selection is locally constant on the connected edge, so $W_\omega=C_{e_\omega}=C_{e_j}=W_j$.
\end{proof}
\status{proved in a referee report (\src{N13} \S3, Prop.~A), one reader only. Consequence: additive rankings are \emph{not} instances of (CP-s) or of (CP-b) with $\mathcal D=\mathcal D^+$ that are compatible with (S), so for these variants ``impossibility contains the additive impossibility'' fails; (CP-a) is the only Walley variant for which the reduction is proved and known to be non-vacuous. For general $\M_i$ the inconsistency is \src{N13} Conj.~N13.A (open).}

\subsection{The Walley-2000 variant}\label{sec:cp-b}
Under (CP-b) put $\M_i\coloneqq\{P:P\ge0\text{ on }\mathcal D_i\}$; then (IB) and strict (G) are automatic and every $P\in\M_i$ satisfies (W) (\src{N11} Prop.~6.3, refereed in \src{N13}). If countably many countably additive members strictly separate $\mathcal D_i$, the conclusions of Corollary~\ref{cor:DCA} hold (\src{N11} Thm.~6.4, refereed in \src{N13}); with strongly non-atomic members instead, \src{N11} Thm.~6.4 gives the finitely additive reduction of \src{N11} Thm.~5.5. If $\mathcal D_i$ contains \emph{invisible} desirable gambles ($P(f)=0$ for all $P\in\M_i$), the ranking is genuinely non-Archimedean and no real additive field obeying (W) detects it. Under a lexicographic representation with countably additive levels the interior lexicographic Bayes models are unique, coherent and selected (\src{N11} Prop.~6.5), but whether $T_1$ is lexicographically Bayes at some non-standard interior prior is the open lexicographic separation problem (\src{R11} \S5, \src{N10} on \src{N9.9}).
\status{the variant with invisible desirability is open.}

\subsection{Rankings that add up: the supporting assessment and its exceptional set}\label{sec:M0}

This subsection and the next two concern an \emph{arbitrary} ranking that adds up, in the sense of Definition~\ref{def:addup}; the ranking need not be additive, complete or a coherent almost-preference. Standing hypotheses \textbf{(H$_{\mathrm{M0}}$)}: $n=3$; the $U$-domain $\M_U$ for $T_1$; a ranking with strict (G), (IB), (LA), (S), (C) and $T_1\in\C_L$ ((LF) is automatic on $\M_U$); a selection $p\mapsto C_p$ as in (S); and (M0). Let $\psi$ be a supporting additive assessment, put $c_i(g)\coloneqq\rho_i(g)\operatorname{sign}(g_i)$, and write
\[
B_q\coloneqq\{g\in U:q\cdot c(g)>0\},\quad Z_q\coloneqq\{g\in U:q\cdot c(g)=0\},\quad E(\psi)\coloneqq\{q\in\Di:\mu(Z_q)>0\}.
\]
$B_q$ is the strict $\psi$-$q$-Bayes set, $Z_q$ the tie set, and $E(\psi)$ the set of \emph{exceptional weightings}. The $\psi$-$q$-Bayes models are the models $B_q\cup S$ with $S\subseteq Z_q$ measurable. Let $\mathcal R\coloneqq\{\psi(X):X\in\M_U\}$ (compact and convex by Lyapunov), $V_\psi(t)\coloneqq\min_{y\in\mathcal R}t\cdot y$, and let $\mathfrak F_t\coloneqq\operatorname{argmin}_{\mathcal R}t\cdot y$ be the \emph{face} at $t$.

\begin{proposition}\label{prop:M0}
Under (H$_{\mathrm{M0}}$), $\psi$ is a strictly guiding additive ranking on the same domain $\M_U$, and:
\begin{enumerate}[leftmargin=*,label=($\alph*$)]
\item[($\alpha$)] $T_1$ is $\psi$-Bayes at some interior weighting $p^*$;
\item[($\beta$)] at \emph{every} weighting $p\in\Delta$ at least one $\psi$-$p$-Bayes model is coherent;
\item[($\gamma$)] $E(\psi)$ is Lebesgue-null, and at every $q\in\Di\setminus E(\psi)$ the $\psi$-$q$-Bayes model is a.e.\ unique and coherent, $B_q=D_{K_q}$;
\item[($\delta$)] the vertex Bayes models are coherent.
\end{enumerate}
If $E(\psi)=\varnothing$, then $c$ is a realisation of $T_1$ at $p^*$, which Theorem~\ref{thm:main} excludes. So a cover forces $E(\psi)\ne\varnothing$.
\end{proposition}
\status{proved. For countably additive weights: \src{N17} Prop.~2.1, refereed CORRECT in \src{N19}; for general weights: via \src{N13} Prop.~B (proved in that referee report; its ($\beta$) re-derived in \src{N24} \S2.1); the restriction to relevant comparisons is \src{N19} Prop.~R3(a) (proved in that referee report, one reader). The citation follows \src{N25\_referee} \S3, which notes that \src{N22} Thm.~2.1 is stated for (CP-a) and that the M0 form is \src{N19} \S2.2 with \src{N13} Prop.~B. The last sentence is the argument of \src{N22} Cor.~2.3 (refereed in \src{N24} \S2.2).}

So the unresolved issue in (M0) is not that coherence disappears at exceptional weightings. At $t\in E(\psi)$ coherent and incoherent minimisers \emph{coexist} (Theorem~\ref{thm:Estructure}(d)); $\psi$ assigns them equal profiles and cannot tell them apart; by Theorem~\ref{thm:kkm} and (C) every undominated model is coherent, so the original ranking must tell them apart, by strict comparisons that $\psi$ does not register.

\begin{theorem}\label{thm:Estructure}
Let $c$ be a strictly signed integrable field on $\M_U$ whose Bayes models are coherent at a.e.\ interior weighting (property ($\gamma$); for instance a supporting assessment under (H$_{\mathrm{M0}}$)). Then:
\begin{enumerate}[leftmargin=*,label=(\alph*)]
\item $V_\psi$ is concave on $\Delta$; at $t\in\Di$ it is differentiable iff $t\notin E(\psi)$, and its superdifferential at $t$ is $\mathfrak F_t$.
\item $E(\psi)=\bigcup_k(\ell_k\cap\Di)$, a countable (possibly empty) union of \emph{atom chords}. For each $k$ there is an \emph{oriented atom} $G_k$: a non-null set of mixed gambles $g\in U$ whose vectors $c(g)$ are all positive multiples of one vector $b_k$; then $\ell_k=\{q:q\cdot b_k=0\}$ is their common chart line. Each $G_k$ lies in one open orthant, $b_k$ has no zero coordinate, and $\ell_k$ misses the vertices of $\Delta$ and cuts off exactly one of them, that of the \emph{odd world} of $b_k$ (the world whose sign differs from the other two). There are no \emph{fan points}: no interior weighting at which a non-null set of gambles whose chart directions are spread without atoms has its chart lines through that weighting.
\item $E(\psi)$ meets every line other than the $\ell_k$ in a countable set; if it contains a segment, it contains the whole open chord.
\item For every $q\in\Di$, $B_q$ is coherent, and it is the coherent $q$-Bayes model with the largest credal set, $B_q=D_{K^\flat_q}$. The map $q\mapsto K^\flat_q$ is upper semicontinuous on $\Di$, equals $K_q$ off $E(\psi)$, and is continuous exactly at the points of $\Di\setminus E(\psi)$. Every extreme point of $\mathfrak F_q$ has a coherent realiser, while models $B_q\cup S$ with $S\subseteq Z_q$ bounded and non-null are incoherent.
\item For an oriented atom $G$ with line $\ell$: at every $q\notin E(\psi)$ on the positive side of $\ell$ the Bayes model contains $G$, and on the negative side it misses $G$. At all but countably many points $t$ of the chord, $B_t$ and $B_t\cup G$ are coherent, so $G$ is a coherent conic shell between two coherent cones.
\end{enumerate}
\end{theorem}
\begin{proof}[Proof sketch]
(a) $V_\psi(t)=\int_U\sum_{i:g_i>0}t_i\rho_i\,d\mu-\int_U[t\cdot c]_+\,d\mu$, and $\Phi'(t;\delta)+\Phi'(t;-\delta)=\int_{Z_t}|\delta\cdot c|\,d\mu$ for the convex part $\Phi$; $c\propto\one$ is impossible on $Z_t$. (b) With $\gamma(g)\coloneqq c(g)/|c(g)|$ and $m\coloneqq\gamma_*\mu$ on the sphere, $\mu(Z_q)=m(q^\perp\cap S^2)$; the atoms of $m$ give the chords and the non-atomic part gives the fan points, of which there are countably many since two great circles meet in two points. \emph{No fans:} along each line $g+\R\one$ the tied good gambles share one oriented chart line through the fan point, and $P(q,\cdot)$ is $1$-Lipschitz in the sup norm; two nearby fibres with different lines then give $|J(g)|+|J(g')|\le2\operatorname{dist}$ for the lengths of their tied intervals, which fails at a density point. (d) Directional limits $B_q\cup(Z_q\cap\{\delta\cdot c>0\})$ are limits in measure of coherent Bayes models, hence coherent; $B_q$ is a countable intersection of them; Straszewicz's theorem for extreme points. (e) For $g\in G$ the sign of $q\cdot c(g)$ is that of $q\cdot b$.
\end{proof}
\status{(a), (c), (d), (e) and the chord description in (b): proved (\src{N25} Thm.~2.1, Cor.~2.2, Thm.~3.3, Thm.~4.1, the claim referred to as N25 Claim 1); confirmed by one referee (\src{N25\_referee} \S\S1--2). ``No fan points'': proved in a referee report (\src{N25\_referee} \S2, Prop.~F); its steps 1--4 (one chart line per fibre and the estimate $|J(g)|+|J(g')|\le2d$) were re-derived in \src{M0\_referee} \S1.2 and \src{M0\_second\_referee} \S2.1, so they have three readings, while step 5 (no fans proper) has one. ``Cuts off one vertex'' is the referee's observation that it follows from \src{N25} Thm.~4.1(a), since $b_k$ has mixed signs and no zero coordinate. Incoherence of bounded non-null splices in (d) is the splice fact of \src{R2} Lemma U (two cones cannot differ by a bounded non-null set). Prop.~F uses only ($\gamma$): no (E$_\psi$), no (S), no lexicographic structure. The \src{N26\_referee} (\S3.3) records fan points as open for general M0 rankings and exhibits a strictly signed field with a fan point; that field does not satisfy ($\gamma$), so it does not conflict with Prop.~F, which the \src{N26} referee did not have. Nothing bounds the number of chords, and nothing prevents $E(\psi)$ from being dense. Crossings of two atom chords (with a two-dimensional face) are not excluded in general.}

\begin{theorem}\label{thm:pstarE}
Assume ($\alpha$) and ($\gamma$), and suppose $p^*\in E(\psi)$. Then there are a direction $\delta$, a sequence $s_m\downarrow0$ with $p^*+s_m\delta\notin E(\psi)$, and a relatively open set of gambles $h\in\partial D_{T_1}\setminus0$ with $P_{T_1}(h)=0$ but $P(p^*+s_m\delta,h)\to P_{K_\delta}(h)\ne0$, where $D_{K_\delta}$ is the directional limit $B_{p^*}\cup(Z_{p^*}\cap\{\delta\cdot c>0\})$. So $P(\cdot,h)$ does not converge to $0$ at $p^*$, and the chain of \S\ref{sec:main} cannot start at $p^*$: Lemma~\ref{lem:levels} at level $0$ through $p^*$ and Step~1 of Theorem~\ref{thm:pivot} both need that convergence. If $p^*$ is a generic point of an atom chord of sign pattern $\pm\sigma$, the affected $h$ have pattern $\sigma$ or $-\sigma$: they lie on one middle arc of a facet and near the opposite corner ray.
\end{theorem}
\status{proved (\src{N25} Thm.~9.1); confirmed with repair (\src{N25\_referee} \S4). Repair: the note's ``$q\mapsto P(q,h)$ has no limit at $p^*$'' is not proved; what is proved is non-convergence to $P_{T_1}(h)=0$, which is all the blocking conclusion needs. Addition (CHECKED, \src{N25\_referee} \S4, transferring \src{N26} Prop.~4.3): if the tie set at $p^*$ is a single oriented atom $G$ and no other chord passes through $p^*$, then $T_1=B_{p^*}\cup G$ is the accepting extreme of the face and $\sigma$ is the sign pattern of a corner ray.}

Let $\mathcal N$ be the convex cone of supporting assessments (the countably additive parts of endorsing weights), and let the \emph{atom set} $A(\psi)$ be the union of the oriented atoms of $\psi$.

\begin{proposition}\label{prop:generic}
Under (H$_{\mathrm{M0}}$):
\begin{enumerate}[leftmargin=*,label=(\alph*)]
\item some $\psi^*\in\mathcal N$ has an atom set of least measure, and $A(\psi^*)$ is contained, modulo null sets, in $A(\psi')$ for every $\psi'\in\mathcal N$;
\item for $\kappa\in(0,\infty)^3$, $\kappa\odot\psi\in\mathcal N$, and $E(\kappa\odot\psi)$ and the target weighting are the images of $E(\psi)$ and $p^*$ under one projective bijection of $\Delta$ fixing the vertices; so ``$p^*\in E$'' is invariant under rescaling;
\item if the ranking covers $T_1$, then every $\psi\in\mathcal N$ has $E(\psi)\ne\varnothing$, hence an atom chord, and $\mu(A(\psi))>0$. Genericity cannot remove the chords.
\end{enumerate}
\end{proposition}
\status{(a): proved (\src{N25} Thm.~6.2); confirmed with repair (\src{N25\_referee} \S1): normalise the weights, $\|\nu_k\|\le1$, rather than $\int|c_k|\le1$, so that the countable mixtures $\sum_kw_k\nu_k$ converge and have countably additive part $\psi^w$. (b): proved (\src{N25} Prop.~6.4(a)); confirmed. (c): proved in a referee report (\src{N25\_referee} \S3), one reader, from Proposition~\ref{prop:M0} and Theorem~\ref{thm:Estructure}; it needs no (E$_\psi$).}

\paragraph{The hypothesis (E$_\psi$).} Say that \textbf{(E$_\psi$)} holds if $\psi$ itself endorses the relevant comparisons: $B\in\C_L$, $X$ regular and $B\preceq_iX$ imply $\psi_i(B)\le\psi_i(X)$. It holds when the endorsing weights are countably additive (then $\psi_i=\nu_i\circ E_i$). It is not supplied by \src{N22} or \src{N24} in general; with purely finitely additive parts, endorsement gives only $\psi_i(C)+\nu^{\mathrm s}_i(E_iC)\le\psi_i(Y)+\nu^{\mathrm s}_i(E_iY)$ with uncontrolled singular parts $\nu^{\mathrm s}_i$ (\src{N25} N25.12, \src{N25\_referee} \S3). Call a selected $C_p$ \emph{$j$-improvable} if some regular $Z$ has $Z\prec_jC_p$ and $\psi_i(Z)\le\psi_i(C_p)$ for both $i\ne j$. On a chord $\ell$, let $\ell^\circ$ be its points on no other chord, so that $Z_t=G\cup G'$ for $t\in\ell^\circ$ ($G,G'$ the two oppositely oriented atoms on $\ell$); a \emph{split} is a measurable $S\subseteq G\cup G'$.

\begin{proposition}\label{prop:Epsi}
Assume (H$_{\mathrm{M0}}$) and (E$_\psi$).
\begin{enumerate}[leftmargin=*,label=(\alph*)]
\item For $t\in\Di$, every point of $\mathfrak F_t$ is the profile of a selected $t$-Bayes model; every $t$-Bayes $Y$ is weakly dominated by a selected $C$ with $\psi(C)=\psi(Y)$, strictly at some world if $Y$ is incoherent.
\item $q\mapsto\psi(C_q)$ is continuous at every $p$ whose selected credal set lies in $\Di$.
\item For each world $j$, the projected profiles $(\psi_i(C_p))_{i\ne j}$ of $j$-improvable selected models lie in countably many graphs of $1$-Lipschitz functions over the line $\one^\perp\subseteq\R^2$; in particular they form a null set.
\item Hence at a two-dimensional face (which by Theorem~\ref{thm:Estructure} occurs only where atom chords cross), for area-a.e.\ tied profile no selected realiser is $j$-improvable for any $j$ (``ideal tie-breaking''); on an atom face, all but countably many tied profiles have no selected realiser improvable at the odd world.
\item Along a chord, (IB) makes the comparisons among the models $B_t\cup S$ the same at every $t\in\ell^\circ$; the undominated splits do not depend on $t$, and for each undominated split $S$ the model $B_t\cup S$ is coherent at every $t\in\ell^\circ$.
\end{enumerate}
\end{proposition}
\status{proved (\src{N25} Lemma 1.2, Lemmas 7.1--7.4, Thm.~7.5, Cor.~7.6, Lemma 8.1, Thm.~8.2); confirmed by one referee (\src{N25\_referee} \S\S1, 4). (E$_\psi$) is genuinely needed: with charges, the tail estimate behind (b) fails (the obstruction of \src{N22} Remarks 3.3). (c) is the analogue, without lexicographic structure, of \src{R2} Thm.~7; it costs one dimension.}

\subsection{Residual configurations of (M0)}\label{sec:M0-residual}
Under (H$_{\mathrm{M0}}$), without (E$_\psi$), a cover of $T_1$ forces, for every $\psi\in\mathcal N$: $E(\psi)$ is a nonempty countable union of open atom chords, there are no fan points, and each chord cuts off the vertex of its odd world. With $\psi^*$ as in Proposition~\ref{prop:generic}(a), one of the following holds (\src{N25\_referee} \S7, superseding \src{N25} \S10). Both residuals are now \textbf{closed}: Theorem~\ref{thm:M0closed} below shows that neither configuration occurs. They are kept here because the closing argument follows them.
\begin{itemize}[leftmargin=*]
\item[\textbf{R-a$'$}] $p^*\in E(\psi^*)$: the target weighting lies on an atom chord. Sub-cases: (i) a single oriented atom at $p^*$; (ii) both oriented atoms of the chord non-null; (iii) $p^*$ at a crossing of two or more chords, with a two-dimensional face. The chain of \S\ref{sec:main} cannot start (Theorem~\ref{thm:pstarE}). In (i), $T_1$ is the accepting extreme and $\sigma$ a corner pattern. Under (E$_\psi$), Proposition~\ref{prop:Epsi} applies. \emph{Status:} \textbf{CLOSED} (round 11): excluded by Propositions~\ref{prop:crosspin} and~\ref{prop:carrying} (Theorem~\ref{thm:M0closed}).
\item[\textbf{R-b$'$}] $p^*\notin E(\psi^*)$, and $E(\psi^*)$ is a nonempty countable union of atom chords, possibly crossing and possibly dense. \emph{Status:} \textbf{CLOSED} (round 11): the chain of \S\ref{sec:main} tolerates chords away from $p^*$ (Proposition~\ref{prop:chainchords}), so $p^*\in E(\psi)$ for every $\psi\in\mathcal N$.
\end{itemize}
Closed earlier: every configuration with fan points (\src{N25}'s residual R-c and the fan half of R-a), by Prop.~F; and \src{N25} Conj.~N25.L, which is vacuous because a countable $E(\psi)$ is empty.

\emph{Localisation (superseded in round 11).} The proof of Theorem~\ref{thm:main} as written uses Hausdorff continuity of $K$ on sets that together cover $\Di$ (\src{N25} Prop.~9.3, CHECKED, confirmed as a reading of the proof). The edge box (Lemma~\ref{lem:edgebox}) holds for every strictly signed integral field whose Bayes models are coherent at a.e.\ weighting (\src{N25} Prop.~9.2, proved, confirmed; also \src{N26} \S4.3). Lemma~\ref{lem:chords} and Lemma~\ref{lem:levels} survive off a null $E$ (\src{N26\_referee} \S4.2, proved there for the class $\Ltwo$; on our reading the argument uses only continuity of $K$ on $\Di\setminus E$ and nullity of $E$, and \src{N25\_referee} \S5 makes the same observation for straight levels). What genuinely crosses an atom chord is the clopen argument of Lemma~\ref{lem:PL2} on $\operatorname{int}\co(G\cup V_l)$: every chord disconnects $\Di$ and separates $p^*$ from some vertex. The remaining failure point is evaluation at $p^*$ when $p^*\in E$. In round 11 both points were resolved: continuity of $K$ is never needed, because quasi-linearity of every tie-breaking family (Lemma~\ref{lem:quasilin}) gives straight levels at \emph{every} weighting and the upper semicontinuous family $K^\flat$ supplies the one semicontinuity needed; Lemma~\ref{lem:PL2} holds for $K^\flat$ across chords by the ball lemma (Lemma~\ref{lem:ball}); and evaluation at $p^*\in E$ is replaced by crossing pinning (Proposition~\ref{prop:crosspin}).

\subsection{Case (M0) is closed on the \texorpdfstring{$U$}{U}-domain}\label{sec:M0-closed}

This subsection shows that both residuals of \S\ref{sec:M0-residual} are empty: on $\M_U$ no ranking that adds up at every world covers $T_1$ (Theorem~\ref{thm:M0closed}). The argument is additive. It concerns any strictly signed field with property ($\gamma$), and it uses no (E$_\psi$), no lexicographic structure and no closed choice beyond what Proposition~\ref{prop:M0} already used. Sources: \src{P4} Part A; the first M0 referee report \src{M0\_referee}, which found the closing step; and the second, independent report \src{M0\_second\_referee}, which strengthened it.

\paragraph{Standing data.} Let $c$ be a strictly signed measurable field on $U=D_{K'}$, $K'=u+s(T_1-u)$, such that $B_q=D_{K_q}$ a.e.\ for a.e.\ $q\in\Di$ (property ($\gamma$)). Then $E\coloneqq\{q\in\Di:\mu(Z_q)>0\}$ is null (Tonelli), and Theorem~\ref{thm:Estructure} applies to $c$; integrability is not used below. $B_q$, $Z_q$ and $K^\flat_q$ are as in \S\ref{sec:M0}, and $P^\flat(q,g)\coloneqq P_{K^\flat_q}(g)$. A gamble $g\in U$ is \emph{good} if $\{q\in\Di\setminus E:P(q,g)>0\}$ agrees with $H(g)\cap\Di$ up to a null set; a.e.\ $g$ is good (Fubini).

\paragraph{Tie-breaking families.} For a measurable $A\subseteq U$ put $B^A_q\coloneqq B_q\cup(Z_q\cap A)$: the strict Bayes set, with the ties in $A$ broken towards acceptance. Call $A$ \emph{admissible} if $B^A_q$ is coherent for \emph{every} $q\in\Di$, say $B^A_q=D_{K^A_q}$, and put $P^A(q,g)\coloneqq P_{K^A_q}(g)$. The empty set is admissible, with $K^\varnothing=K^\flat$ (Theorem~\ref{thm:Estructure}(d)). So is $A_\delta\coloneqq\{g:\delta\cdot c(g)>0\}$ for every $\delta\in\one^\perp$ parallel to no atom chord (all but countably many directions); then $B^{A_\delta}_q$ is the directional limit of Theorem~\ref{thm:Estructure}(d) (\src{N25} Lemma~3.1). Write $K_\delta\coloneqq K^{A_\delta}_{p^*}$. Call such $\delta$ \emph{generic}.

\begin{lemma}\label{lem:quasilin}
Let $A$ be admissible.
\begin{enumerate}[leftmargin=*,label=(\alph*)]
\item For every $g\in\R^3$ the function $q\mapsto P^A(q,g)$ is quasi-concave and quasi-convex on $\Di$: all its level sets $\{>t\}$, $\{\ge t\}$, $\{<t\}$, $\{\le t\}$ are convex.
\item $K^A_q=K_q$ and $K^A$ is Hausdorff-continuous at every $q\in\Di\setminus E$; $K^A_q\to\co(K'\cup\{e_l\})$ as $q\to e_l$ through $\Di$, so Lemma~\ref{lem:UVN} holds for $P^A$.
\item For good $g$, the set $\{P^A(\cdot,g)>0\}$ contains $H(g)\cap\Di$ and lies in $\cl H(g)$, and its points on $\partial H(g)$ lie in $E$. For $A=\varnothing$, $\{P^\flat(\cdot,g)>0\}=H(g)\cap\Di$ exactly.
\end{enumerate}
\end{lemma}
\begin{proof}[Proof sketch]
(a) For fixed $g$, the set of $q$ with $g\in B^A_q$ is an open half-plane ($g\notin A$) or a closed one ($g\in A$), so $B^A_{q_1}\cap B^A_{q_2}\subseteq B^A_q\subseteq B^A_{q_1}\cup B^A_{q_2}$ for $q\in[q_1,q_2]$. If $P^A(q,g_0)<t<\min_iP^A(q_i,g_0)$, a ball around $g_0-t\one$ lies in $D_{K^A_{q_1}}\cap D_{K^A_{q_2}}$ but outside $D_{K^A_q}$, which contradicts the first inclusion; the bound by the maximum is symmetric. (b) Off $E$ the tie rule affects a null set of gambles, so $B^A_{q_m}\to B_q$ in measure, and cones converging in measure have Hausdorff-converging credal sets (\src{R1} Prop.~3.1). As $q\to e_l$, $q\cdot c(g)\to c_l(g)\ne0$ for every $g$. (c) A convex set that agrees a.e.\ with an open half-plane contains it and lies in its closure; boundary points off $E$ are excluded by continuity at $q$; $K^\flat$ is upper semicontinuous, so $\{P^\flat(\cdot,g)>0\}$ is open.
\end{proof}
\status{proved (\src{P4} Lemmas~2.1, 2.2); confirmed by two independent referees (\src{M0\_referee} \S1, \src{M0\_second\_referee} \S1). The case $A=\varnothing$ of (a) was proved independently as \src{P1} Lemma~P1.2 (straight levels for $K^\flat$ at every interior weighting), confirmed in \src{M0\_referee} \S2. So Lemma~\ref{lem:levels} holds for every admissible family, at every weighting and every level, with no continuity in $q$.}

\begin{lemma}\label{lem:ball}
Let $x\in\Di$ and let $O$ be an open ball of gambles with $O\subseteq Z_x$ modulo null sets. Then the vectors $c(g)$, $g\in O$, are a.e.\ positive multiples of one vector. In particular $O$ lies, modulo null sets, in one open orthant, and no ball centred on a coordinate plane $\{g_c=0\}$ lies in $Z_x$ modulo null sets.
\end{lemma}
\begin{proof}[Proof sketch]
Along a.e.\ fibre $g+\R\one$ meeting $O$, the good tied gambles share one oriented chart line through $x$, and two fibres at sup-distance $d$ with different lines satisfy $|J(g)|+|J(g')|\le2d$, where $J$ is the set of tied levels (steps 1--4 of Prop.~F; see the proof of Theorem~\ref{thm:Estructure}(b)). Fibres through a smaller concentric ball have $|J|$ bounded below, so the line is locally constant, hence constant. An oriented line through $x$ fixes $c$ up to a positive factor, because $\aff\Delta$ spans $\R^3$, and strict signs give the orthant.
\end{proof}
\status{proved (\src{P4} Lemma~2.3). Confirmed by two referees: \src{M0\_referee} \S1.2 and \src{M0\_second\_referee} \S2.1. It rests on steps 1--4 of Prop.~F (\src{N25\_referee} \S2), which have now been read three times (\src{N25\_referee}, \src{M0\_referee}, \src{M0\_second\_referee}). Step 5 of Prop.~F (no fan points) is not used here or anywhere in this subsection.}

\paragraph{Face models.} Let $T\subseteq\Di$ be a credal set and $p^*\in\Di$ with $D_T=B_{p^*}\cup S$ for a measurable $S\subseteq Z_{p^*}$. A \emph{$p^*$-face model} is a credal set $K$ with $D_K=B_{p^*}\cup R$ modulo null sets for some measurable $R\subseteq Z_{p^*}$: a coherent model that the field cannot tell apart from $T$ at $p^*$. Then $D_K\sd D_T\subseteq Z_{p^*}$ and $D_K\subseteq U$. Examples are $T$ itself, $K^\flat_{p^*}$, and every directional limit $K_\delta$. On a facet $F_k$ of $D_{T_1}$, with crossings at $\tau_i<\tau_j$ (Lemma~\ref{lem:signs}), call $h((\tau_i,\tau_j))$ the \emph{middle arc} and $h([0,\tau_i))$, $h((\tau_j,1])$ the \emph{corner portions}.

\begin{proposition}\label{prop:crosspin}
Let $T\subseteq\Di$ be a credal set with $D_T=B_{p^*}\cup S$, $S\subseteq Z_{p^*}$, $p^*\in\Di$. Then every $p^*$-face model $K$ has $P_K(w)=0$ at every $w\in\partial D_T\setminus\{0\}$ that has a zero coordinate. For $T=T_1$ (and for every all-facing triangle), on each facet $F_k$ the function $\tau\mapsto P_K(h(\tau))$ vanishes at both crossings, is $\ge0$ on the closed middle arc, and is $\le0$ on both corner portions.
\end{proposition}
\begin{proof}
Let $w_c=0$ and let $m_w\in T$ attain $P_T(w)=0$. If $P_K(w)>0$, put $g_1\coloneqq w-s(\one-e_c)$ with $0<s<P_K(w)$. Then $(g_1)_c=0$, $P_T(g_1)\le m_w\cdot g_1=-s(1-m_{w,c})<0$ and $P_K(g_1)\ge P_K(w)-s>0$, so a ball around $g_1$ lies in $D_K\setminus\cl D_T\subseteq Z_{p^*}$. If $P_K(w)<0$, put $g_2\coloneqq w+s(\one-e_c)$ with $0<s<|P_K(w)|$. Then $(g_2)_c=0$, $P_T(g_2)\ge s(1-\max_{m\in T}m_c)>0$ and $P_K(g_2)\le P_K(w)+s<0$, so a ball around $g_2$ lies in $D_T\setminus\cl D_K\subseteq S\subseteq Z_{p^*}$. Either ball is centred on $\{g_c=0\}$, against Lemma~\ref{lem:ball} at $x=p^*$. For $T_1$ the boundary gambles with a zero coordinate are the six crossings, two in the relative interior of each facet (Lemma~\ref{lem:signs}), and the signs follow by concavity of $\tau\mapsto P_K(h(\tau))$.
\end{proof}
\status{proved. Stated for $T_1$ with the hypothesis $p^*\in E$ in \src{M0\_referee} \S5.2 (Prop.~R.1); confirmed with a strengthening repair in \src{M0\_second\_referee} \S2.2, which drops ``$p^*\in E$'' (if $Z_{p^*}$ is null, no ball lies in it) and the triangle hypothesis. The form above is the repaired one; two readers. The only external input is Lemma~\ref{lem:ball}, applied to the tie set $Z_{p^*}$ of the field; $K$ enters only through $D_K\sd D_T\subseteq Z_{p^*}$, so it does not matter that a directional limit is only a limit in measure of Bayes models. What is specific to all-facing triangles is the \emph{location} of the zero-coordinate gambles. NUMERICS (\src{m0ref\_py/crossing\_balls.py}, \src{m0ref2\_py/geometry\_check.py}): 9600 and 450 constructions, no failure.}

\begin{proposition}\label{prop:chainchords}
Let $c$ be as above, $p^*\in\Di\setminus E$, and $K^\flat_{p^*}=T_1$. Then Lemma~\ref{lem:PL2} holds for $K^\flat$ on every open connected set, chords allowed; Lemma~\ref{lem:nopin} holds for every admissible family; and the conclusions of Theorems~\ref{thm:pivot} and~\ref{thm:nosliding}, Lemma~\ref{lem:rotation}, Propositions~\ref{prop:regularity}, \ref{prop:blowup} and~\ref{prop:boxbridge} and Corollary~\ref{cor:PCNS} hold for $P^\flat$. Proposition~\ref{prop:T1} then gives a contradiction. Consequently, if a strictly signed field with ($\gamma$) has $D_{T_1}$ as a Bayes model at an interior weighting $p^*$, then $p^*\in E$.
\end{proposition}
\begin{proof}[Proof sketch]
This is a step-by-step audit of \S\ref{sec:main} with $P=P^\flat$. Since $p^*\notin E$, no line through $p^*$ is an atom chord, so every line through $p^*$ meets $E$ in a countable set (Theorem~\ref{thm:Estructure}(c); without Prop.~F step 5 this still holds, since the remaining exceptional points are countable and each chord meets the line at most once). Lemmas~\ref{lem:chords}(a) and~\ref{lem:levels}: Lemma~\ref{lem:quasilin}, which needs no cluster lines and no joint continuity. Lemma~\ref{lem:UVN}: Lemma~\ref{lem:quasilin}(b). Lemma~\ref{lem:PL2} for $K^\flat$: $\beta$ is relatively open by upper semicontinuity of $K^\flat$; at a point $Q\in\alpha\cap\cl\beta$, every Hausdorff cluster value $K^*$ of $K^\flat_q$ contains $v$, since otherwise $D_{K^*}$ contains a collar of $F_k$ near the crossing that misses $D_{K^\flat_Q}$ and so lies in $Z_Q$ modulo null sets, and the collar contains a ball centred on $\{g_c=0\}$, against Lemma~\ref{lem:ball} at $Q$; the original argument then runs with the exact charts of $K^\flat$. Lemma~\ref{lem:nopin}: $\operatorname{int}\co(G\cup V_j)\subseteq\Gamma$ follows directly from the exact charts, and $K^A\subseteq K^\flat$ transfers $v_k\notin K$. Step~4 of Theorem~\ref{thm:pivot}: only $q\in\lambda\setminus E$ (a co-countable set) is used; continuity of $K^\flat$ at such $q$ excludes $\Psi(q,\tau)<0$, and an open half-plane through $p^*$ that misses a co-countable part of $\lambda$ has boundary $\lambda$. Lemma~\ref{lem:rotation} and Propositions~\ref{prop:regularity}, \ref{prop:blowup}: Lemma~\ref{lem:caps}(c) is needed only at the limit anchor $q_{\tau_1}(\xi)$, which is chosen off $E$ (a co-countable set of $\xi$ for each of the countably many $\tau$ used); the approaching anchors enter only pointwise. Cluster clause: a robust form from the exact charts and the $1$-Lipschitz dependence on $g$. Lemma~\ref{lem:edgebox}: it holds for every field with ($\gamma$) (\src{N25} Prop.~9.2). Proposition~\ref{prop:T1}: $p^*\odot c$ and the transposed field $c^\sigma(g)\coloneqq\sigma c(\sigma g)$ preserve ($\gamma$), move $E$ together with $p^*$, and fix $U$ and $D_{T_1}$. The other steps are pointwise in $q$. For the last sentence: if $p^*\notin E$, then $S\subseteq Z_{p^*}$ is null and $T_1=K^\flat_{p^*}$.
\end{proof}
\status{proved (\src{P4} Thm.~3.1 and Thm.~A1); confirmed by two referees: \src{M0\_referee} \S1.1 (both items the author asked to be rechecked, Step~4 and the anchor choice, re-derived; wording repair incorporated: joint continuity is used only at the limit anchor) and \src{M0\_second\_referee} \S4.1 (step-by-step table against \S\ref{sec:main}). This closes residual R-b$'$. For every all-facing triangle in general position, through the \src{R15} algebra in place of Proposition~\ref{prop:T1}: CHECKED (\src{P4}; \src{M0\_referee} \S7; \src{M0\_second\_referee} \S5.3).}

\paragraph{Carrying atoms.} Suppose now that $p^*\in E$. On a chord $\ell$ with points $\ell^\circ$ on no other chord (\S\ref{sec:M0}), and an oriented atom $G$ on $\ell$, the model $B_t\cup G$ is coherent for $t\in\ell^\circ$ (Theorem~\ref{thm:Estructure}(e)); write $D_{K^+(t)}\coloneqq B_t\cup G$.

\begin{proposition}\label{prop:carrying}
Let $c$ be as above, $p^*\in E$, and $D_{T_1}=B_{p^*}\cup S$ with $S\subseteq Z_{p^*}$.
\begin{enumerate}[leftmargin=*,label=(\alph*)]
\item $S$ is non-null; equivalently, $T_1\ne K^\flat_{p^*}$.
\item $K^\flat_{p^*}\supsetneq T_1$ and $S=D_{T_1}\setminus D_{K^\flat_{p^*}}$. The set of $h\in\partial D_{T_1}\setminus\{0\}$ with $P^\flat(p^*,h)<0$ is a union of at most three open arcs, each containing one corner ray, contained in the corner portions of the two adjacent facets, and containing no crossing. Over each such arc $\alpha_G$ the shell $\{h+s\one:h\in\alpha_G,\ 0<s<|P^\flat(p^*,h)|\}$ lies, modulo null sets, in one oriented atom $G$ (a \emph{carrying atom}), whose chord passes through $p^*$ and whose sign pattern is that of the corner ray. $S$ is the union of its parts in the carrying atoms.
\item For a carrying atom $G$ on $\ell$, the values $\pi^+(h)\coloneqq P_{K^+(t)}(h)$ and $\pi^\flat(h)\coloneqq P^\flat(t,h)$ do not depend on $t\in\ell^\circ$, and $\pi^+\ge0>\pi^\flat$ on $\alpha_G$.
\item If $\pi^+$ vanishes on an open subinterval of $\alpha_G\cap\relint F_k$ for some facet $F_k$, there is a contradiction.
\end{enumerate}
\end{proposition}
\begin{proof}[Proof sketch]
(a) If $S$ is null, then $D_{K_\delta}\supseteq D_{T_1}$, i.e.\ $K_\delta\subseteq T_1$, for every generic $\delta$. A proper subset of $T_1$ misses a vertex $v_k$ and is then positive on all of $\relint F_k$, crossings included, against Proposition~\ref{prop:crosspin}. So $K_\delta=K_{-\delta}=T_1$ for all generic $\delta$, and $Z_{p^*}$ is null, against $p^*\in E$. (b) $S=D_{T_1}\setminus D_{K^\flat_{p^*}}\subseteq Z_{p^*}$. For each $m\in K^\flat_{p^*}$ the set $\{h\in F_k:m\cdot h<0\}$ is a sub-cone of $F_k$ containing a corner ray, and no crossing lies in the negative set, by the ball construction of Proposition~\ref{prop:crosspin}. Over each arc the shell is open, connected and contained in $Z_{p^*}$, so Lemma~\ref{lem:ball} and connectedness give one atom $G$ with $c\propto b$; $G\subseteq Z_{p^*}$ gives $p^*\cdot b=0$, and the sign pattern is that of the gambles near the corner ray. (c) Frozen support (\src{N25} Thm.~4.1(e)): $G\cap(g+\R\one)=\{g-s\one:\pi^\flat(g)\le s<\pi^+(g)\}$, and the shell over $h\in\alpha_G$ lies in $G$, so $\pi^+(h)\ge0$ and $\pi^\flat(h)\le-|P^\flat(p^*,h)|<0$. (d) Let $r=r_{kj}$ be the corner ray of $\alpha_G$, with odd world $o$, and take a generic $\delta$ with $\delta\cdot b>0$. At $t\in\ell^\circ$, $K^{A_\delta}_t=K^+(t)$, so $\varphi_t(\tau)\coloneqq P^{A_\delta}(t,h(\tau))$ is concave and vanishes on the subinterval; by Lemma~\ref{lem:caps}(a), $\varphi_t\le0$ on $\R$. A gamble $h(\tau_m)$ on the middle arc of $F_k$ is positive at $j$ and at $o$, and $\ell$ separates $e_o$ from $e_j$ ($b_o>0>b_j$). By Lemma~\ref{lem:UVN} for $P^{A_\delta}$ and convexity (Lemma~\ref{lem:quasilin}), $\{P^{A_\delta}(\cdot,h(\tau_m))>0\}$ contains segments from near $e_o$ to near $e_j$; their crossings with $\ell$ fill an open interval of $\ell$, which meets $\ell^\circ$. There $\varphi_t(\tau_m)>0$, a contradiction.
\end{proof}
\status{proved (\src{P4} Prop.~4.1(i), (ii), Lemmas~2.4, 2.5, 4.2, Thm.~4.3); confirmed by two referees (\src{M0\_referee} \S1; \src{M0\_second\_referee} \S\S2.3, 4.3). In (a), the argument through Proposition~\ref{prop:crosspin} is the second referee's; \src{P4} uses its shell Lemma~2.4(a), and both are confirmed. Theorem~4.3 of \src{P4} needs $\varphi_t=0$ only on a subinterval, which is how (d) is stated. Inputs: \src{N25} Thm.~4.1(a), (b), (d), (e), (f) and Lemma~3.1, confirmed once in \src{N25\_referee} and re-read in both M0 reports.}

\begin{theorem}\label{thm:M0closed}
Fix $\varepsilon\in(0,1)$, $0<s<1$, $K'=u+s(T_1-u)$ and $U=D_{K'}$.
\begin{enumerate}[leftmargin=*,label=(\alph*)]
\item \emph{Additive form.} No strictly signed field $c$ on $U$ whose Bayes models are coherent at a.e.\ interior weighting has $D_{T_1}$ as a Bayes model at an interior weighting: there is no $p^*\in\Di$ with $D_{T_1}=B_{p^*}\cup S$ for a measurable $S\subseteq Z_{p^*}$.
\item \emph{Ranking form.} On $\M_U$ there is no ranking with strict (G), (IB), (LA), (S), (C) and $T_1\in\C_L$ in which every world adds up on its relevant comparisons (case (M0)). This holds with or without (E$_\psi$).
\end{enumerate}
\end{theorem}
\begin{proof}
(a) If $p^*\notin E$, Proposition~\ref{prop:chainchords} gives a contradiction. Let $p^*\in E$. By Proposition~\ref{prop:carrying}(a), (b), $S$ is non-null and is carried by a corner atom $G$ on a chord $\ell\ni p^*$, with normal $b$ and an arc $\alpha_G$ inside the corner portions. Since $\pi^+\ge0$ is continuous on $\alpha_G$, Proposition~\ref{prop:carrying}(c), (d) give some $h\in\alpha_G\cap\relint F_k$ with $\pi^+(h)>0$. Take a generic $\delta$ with $\delta\cdot b>0$; then $G\subseteq Z_{p^*}\cap\{\delta\cdot c>0\}\subseteq D_{K_\delta}$. For $0<s<\pi^+(h)$ the gamble $h-s\one$ has a neighbourhood contained in $G$ modulo null sets (\src{N25} Thm.~4.1(d)), hence in the open convex set $D_{K_\delta}$; so $P_{K_\delta}(h)>s>0$. But $K_\delta$ is a $p^*$-face model and $h$ lies in a corner portion, so Proposition~\ref{prop:crosspin} gives $P_{K_\delta}(h)\le0$, a contradiction. (b) Under (H$_{\mathrm{M0}}$), Proposition~\ref{prop:M0} gives a supporting assessment whose field satisfies ($\alpha$) at an interior weighting and ($\gamma$); apply (a).
\end{proof}
\status{(a) \textbf{proved}. Sources: \src{P4} Part A with \src{M0\_referee} \S5 (Prop.~R.1, Theorem~R). Every new part has been read independently by two referees (\src{M0\_referee}, \src{M0\_second\_referee}): CONFIRMED, and Prop.~R.1 CONFIRMED WITH REPAIR (a strengthening, incorporated in Proposition~\ref{prop:crosspin}). It strictly strengthens Theorem~\ref{thm:main}, which assumes (BU) at \emph{every} interior weighting; nullity of $E$ follows from ($\gamma$), and integrability is not used. (b) \textbf{proved at the status of Proposition~\ref{prop:M0}} ($\alpha$), ($\gamma$) without (E$_\psi$): for general finitely additive weights these rest on \src{N13} Prop.~B and \src{N19} Prop.~R3(a), each proved in a referee report and read once; for countably additive weights on \src{N17} Prop.~2.1, refereed in \src{N19}. Neither M0 referee re-derived these imports; a second reading of them would put (b) at the standard of (a). Not used: (E$_\psi$); Prop.~F step 5; the hypothesis (HC) of \src{P1}; Lemmas~P1.1, P1.3 of \src{P1}; the results of \src{N26}; Theorem~\ref{thm:ccibu}; lexicographic structure; \src{P4} Part B. Dependency list as completed by \src{M0\_second\_referee} \S6: \src{N25} Thm.~2.1, Lemmas~3.1, 3.2, Thms.~3.3, 4.1 (a), (b), (d), (e), (f); Prop.~F steps 1--4; Lemmas~\ref{lem:signs}, \ref{lem:quasilin}, \ref{lem:ball}; Propositions~\ref{prop:crosspin}, \ref{prop:chainchords}, \ref{prop:carrying}; the odd-vertex cut of Theorem~\ref{thm:Estructure}(b). \emph{Scope:} the same statement for every all-facing triangle in general position, and for every $K'\subseteq\relint T_1$, is CHECKED only (\src{M0\_referee} \S7, \src{M0\_second\_referee} \S\S5.3, 7; NUMERICS on 2000 random all-facing triangles). \emph{Domain:} proved on the $U$-domain only; on general decomposable domains only the additive Theorem~\ref{thm:gm} is proved, and rankings that add up there have not been examined.}

\begin{corollary}\label{cor:M0closed}
On $\M_U$:
\begin{enumerate}[leftmargin=*,label=(\alph*)]
\item every ranking with strict (G), (IB), (LA), (S), (C) that covers $T_1$ has a world that does not add up, i.e.\ a world with no weight, nonnegative on its relevant comparisons, whose countably additive part is positive a.e.\ on mixed problems ($\neg$SP$^{\mathrm{reg}}$; case (M1) or (M2));
\item residuals R-a$'$ and R-b$'$ are empty, and so are the triple-point configuration of \src{P4} and the split-facet configuration of \src{P1} (below);
\item no ranking in the class $\Ltwo$ of \S\ref{sec:M0-lex} covers $T_1$, since its primary endorses every comparison;
\item question (Q1) of \S\ref{sec:two-questions} is answered negatively.
\end{enumerate}
\end{corollary}
\status{proved, from Theorem~\ref{thm:M0closed} (\src{M0\_referee} \S7; confirmed in \src{M0\_second\_referee} \S7). In (a), ``not (M0)'' is equivalent to $\neg$SP$^{\mathrm{reg}}$ at some world by \src{N19} Lemma~R1 in one direction and trivially in the other. Item (c) supersedes Proposition~\ref{prop:R1N26} and Conjecture~\ref{conj:N26R} as routes to this conclusion. \src{P4} Thm.~A3 (no (M0) cover under (E$_\psi$); proved and confirmed) is a special case.}

Every (CP-a) ranking on $\M_U$ is excluded as well. For $P\in\prod_i\mathcal P^+_i$ write, as in the proof of Corollary~\ref{cor:residual}, $\bar g\coloneqq\operatorname{sign}(g_1)\,g$ for the representative in $H$ of the pair of $g\in U$, $a_g$ for the advice ``accept $g$'', and
\[
c^P_i(g)\coloneqq\rho_{P_i}(\bar g)\bigl(\lambda_i(\bar g,0)-\lambda_i(\bar g,a_g)\bigr),\qquad g\in U,\ i=1,2,3 .
\]

\begin{corollary}\label{cor:CPaU}
Fix $\varepsilon\in(0,1)$, $0<s<1$, $K'=u+s(T_1-u)$ and $U=D_{K'}$.
\begin{enumerate}[leftmargin=*,label=(\alph*)]
\item On $\M_U$ there is no ranking with strict (G), (LA), (S), (C) and $T_1\in\C_L$ in which every world is a coherent almost-preference over losses ((IB) and weak (G) are then automatic, Proposition~\ref{prop:normalform}). Here each $\lambda_i$ is any bounded measurable penalty that is strictly increasing in world-$i$ rank a.e., and each $\M_i$ is any nonempty weak-* closed convex set of null-invariant states: complete or not, with purely finitely additive and atomic parts allowed.
\item More generally, on $\M_U$ there is no ranking with strict (G), (IB), (LA), (S), (C) and $T_1\in\C_L$ in which every world is (CP-a) or adds up.
\end{enumerate}
\end{corollary}
\begin{proof}
(a) Suppose such a ranking exists; it satisfies (H) ((LF) is automatic on $\M_U$). Fix $P\in\prod_i\mathcal P^+_i$; this set is nonempty by Theorem~\ref{thm:fullsupport}, since every world has (CP-a) and strict (G). \emph{Strict signs.} $P_i\in\mathcal P^+_i$ gives $\rho_{P_i}>0$ a.e.\ on $N_i$, which on $\M_U$ contains every pair meeting $U$, including the nonnegative gambles of $U$. By Definition~\ref{def:CPa}(i) and the background order, $\lambda_i(\bar g,0)-\lambda_i(\bar g,a_g)$ is $>0$ where $g_i>0$ and $<0$ where $g_i<0$, for a.e.\ $g$ (coordinate planes are null). So $c^P$ is strictly signed on $U$. \emph{Translation.} The Bayes models of the integral field $\psi^P$ (integrable envelope $\|\lambda_i\|_\infty\rho_{P_i}$) are exactly its pointwise minimisers (\src{R1} Lemma~1.2). On a pair meeting $U$, $p\cdot\psi^P(\bar g,\cdot)$ has two minimisers on $\{a_g,0\}$ iff $p\cdot c^P(g)=0$; pairs with neither orientation in $U$ have advice range $\{0\}$ and no ties. Since $g\mapsto\bar g$ is measure-preserving ($\mu$ is symmetric), $Z_p(P)$ corresponds to $Z_p(c^P)$ and $E(P)=E(c^P)$. Hence Theorem~\ref{thm:integral}($a$) gives $p_T\in\Di$ and a measurable $S\subseteq Z_{p_T}(c^P)$ with $D_{T_1}=B_{p_T}\cup S$ (note $D_{T_1}\subseteq U$), which is the hypothesis of Theorem~\ref{thm:M0closed}(a) at $p^*=p_T$; and Theorem~\ref{thm:integral}($c$) gives $B_q=D_{K_q}$ a.e.\ for $q$ off the null set $E(c^P)$, which is property ($\gamma$). Theorem~\ref{thm:M0closed}(a) excludes this. Neither ($\beta$) nor ($\delta$) of Proposition~\ref{prop:M0} is needed, and the argument does not pass through (U): it tolerates $E(P)\ne\varnothing$.

(b) A (CP-a) world $i$ with strict (G) adds up on $\M_U$. Put $w_i(g)\coloneqq|\lambda_i(\bar g,a_g)-\lambda_i(\bar g,0)|$, which is $>0$ a.e., and $W_i\coloneqq U\cap\{g_i>0\}$. Then $\Lambda_iX=\Lambda_iW_i+w_i1_{E_i(X)}$ on the pairs meeting $U$, and $\Lambda_iX=\Lambda_iW_i$ on the others. For $P_i\in\mathcal P^+_i$, $\nu_i(A)\coloneqq P_i(w_i1_A)$ is finitely additive, vanishes on null sets, satisfies $\nu_i(A)\ge\int_Aw_i\rho_{P_i}\,d\mu>0$ for non-null $A$, and by (W) endorses \emph{every} comparison at world $i$. So every world adds up, which is (M0), and Theorem~\ref{thm:M0closed}(b) applies.
\end{proof}
\status{(a) \textbf{proved}. The glue (strict signs of $c^P$; the translation of Theorem~\ref{thm:integral}($a$), ($c$) into the hypotheses of Theorem~\ref{thm:M0closed}(a); $E(P)=E(c^P)$) has had one reading (\src{CPaU\_referee} \S\S1--3: CORRECT WITH REPAIR, repairs R1--R8 incorporated; discrete sanity checks \src{cpauref\_py/glue\_check.py}, NUMERICS, no verdict rests on them). Inputs: Theorem~\ref{thm:fullsupport} (\src{N22} Thm.~1.2; refereed once, \src{N24} \S1); Theorem~\ref{thm:integral}($a$), ($c$) (\src{N22} Thm.~2.1 with \src{N11} Cor.~3.2(d), \src{R11} Lemma~IC and \src{N13} Prop.~B; refereed in \src{N13} and \src{N24} \S2.1), the weakest links; Theorem~\ref{thm:M0closed}(a) (new parts read independently twice). Not used: (U), Corollary~\ref{cor:residual}, Theorem~\ref{thm:complete}, (UN), Proposition~\ref{prop:Ucont}, (E$_\psi$), Proposition~\ref{prop:M0}, \src{N19} Prop.~R3(a), Theorem~\ref{thm:ccibu}, Prop.~F step 5. On $\M_U$ it subsumes Corollaries~\ref{cor:complete} and~\ref{cor:DCA}. (b) proved in a referee report (\src{CPaU\_referee} \S6, ``Route B''), one reader, at the status of Theorem~\ref{thm:M0closed}(b); it places (CP-a) worlds inside (M0) and covers mixed systems. Routes (a) and (b) share their core (\src{N13} Prop.~B), so they are not independent; (a) has the better-documented status. \emph{Scope:} proved on the $U$-domain with $K'=u+s(T_1-u)$; for every $K'\subseteq\relint T_1$ and every all-facing triangle in general position CHECKED, inherited from Theorem~\ref{thm:M0closed}. Not established: other non-adjacent domains (for instance under (UN)), general decomposable domains, (CP-s), (CP-b) with invisible desirability, $n\ge4$. There Conjecture~\ref{conj:U} and Corollary~\ref{cor:residual} remain the live statements.}

\paragraph{How the residuals were closed.} \src{P4} first reduced case (M0) without (E$_\psi$) to a \emph{triple-point configuration}: $p^*$ a common point of at least three atom chords, carrying three corner atoms with positively spanning normals, each \emph{partial}, i.e.\ straddling $\partial D_{T_1}$ along its corner arc (\src{P4} Cor.~4.4, confirmed as a list of necessary conditions; the name is unrelated to condition (TP) of Definition~\ref{def:realisation}). Full carrying atoms ($G\subseteq D_{T_1}$) were already excluded by Proposition~\ref{prop:carrying}(d), and they are automatic when $p^*$ lies on one chord, when at most two atoms carry $S$, or when the carrying normals lie in an open half-plane (\src{P4} Prop.~4.1(iii)). \src{M0\_referee} observed that \src{P4} had applied the ball lemma to $K^\flat$ and to subsets of $T_1$ but not to the directional limits $K_\delta$, which are neither; doing so (Proposition~\ref{prop:crosspin}) empties the triple-point configuration.

\begin{remark}\label{rem:midarc}
If $c$ is as above, $D_{T_1}=B_{p^*}\cup S$ with $S\subseteq Z_{p^*}$ and $p^*\in E$, then every $p^*$-face model $K$ has $P_K\equiv0$ on every closed middle arc; hence $K\supseteq T_1$ and $S=Z_{p^*}$. This is not needed for Theorem~\ref{thm:M0closed}, whose hypotheses are in fact never met. It records the local reason why no tie gambles sit outside $\cl D_{T_1}$ over a middle arc: two atoms through $p^*$ (one over the middle arc, one over the adjacent corner portion) would force $P(\cdot,w)>0$ on one sector at $p^*$ cut out by their chords and $P(\cdot,w)=0$ on the two sectors adjacent to it (which are vertically opposite to each other), off $E$, against convexity of $\{P^\flat(\cdot,w)\le0\}$, where $w$ is the crossing between the two arcs.
\end{remark}
\status{proved in a referee report, as Prop.~R.2 of \src{M0\_referee} (\S5.4). Confirmed by the second referee, \src{M0\_second\_referee} \S2.4 (strictness at $w$ and the hull-of-opposite-sectors step checked; NUMERICS for the sector-hull step, 20000 random line pairs, no failure).}

\paragraph{The results of \src{P1}.} \src{P1} attacked case (M0) from the construction side. Its proved results are as follows.

\begin{lemma}\label{lem:P1box}
Let $\kappa>1$ and $x_1,x_2,x_3>0$ satisfy, for each $k$ (indices mod $3$), $x_k+1/x_{k+1}\ge\kappa+1$ and $1/\kappa\le x_kx_{k+1}\le\kappa$. This is impossible.
\end{lemma}
\begin{proof}
Put $a=x_k$, $b=x_{k+1}$. From $ab\le\kappa$: $\kappa/b+1/b\ge a+1/b\ge\kappa+1$, so $b\le1$. From $ab\ge1/\kappa$: $1/b\le\kappa a$, so $(1+\kappa)a\ge\kappa+1$ and $a\ge1$. Hence $x_k\ge1\ge x_{k+1}$ for every $k$, so $x_1=x_2=x_3=1$, and then $2\ge\kappa+1$, a contradiction.
\end{proof}
\status{proved (\src{P1} Lemma~P1.1); confirmed (\src{M0\_referee} \S2; sympy and random checks rerun, no feasible point). It strengthens the algebra of Proposition~\ref{prop:T1}: the box contradiction survives when the corner equations are weakened to the one-sided inequalities. The six box bounds are those of Proposition~\ref{prop:T1}: the upper ones are the $(1,2)$ certificate written with general facet data ($x_1x_3\le\kappa$) and its cyclic images; the lower ones come from the mirror, since the transposed field has facet data $(1/x_2,1/x_1,1/x_3)$. \src{P1} had the lower bounds at CHECKED/NUMERICS level; \src{M0\_referee} \S2 upgraded them to \textbf{proved} (symbolic check \src{m0ref\_py/transpose\_box.py}), one reader.}

The case $A=\varnothing$ of Lemma~\ref{lem:quasilin}(a) is \src{P1} Lemma~P1.2 (confirmed). \src{P1} Lemma~P1.3 and Thm.~P1.4 (an atom through $p^*$ at a corner turns the corner equation into the strict inequality $x_k+1/x_{k+1}>\kappa+1$, so by Lemma~\ref{lem:P1box} corner atoms cannot rescue a cover) were proved conditionally on a hypothesis (HC) that the chain runs sector by sector; they are confirmed with a scope repair in \src{M0\_referee} \S2 and superseded by Theorem~\ref{thm:M0closed}. \src{P1} \S4 found that \emph{split facets} (atom chords through $p^*$ whose arcs lie in the middle regions of the facets, giving each facet two linear pieces) make the facet-data algebra at $p^*$ consistent: NUMERICS, reproduced and extended to strict corner inequalities by \src{M0\_referee} (\src{m0ref\_py/tp\_split.py}). As a statement about the algebra this is correct; as an escape it is \textbf{wrong}: the configuration is impossible (Theorem~\ref{thm:M0closed}; locally, Remark~\ref{rem:midarc}). \src{P1} \S4.3(a), (b) are confirmed; its conclusion that nothing local closes the configuration is wrong.

\paragraph{What singles out $T_1$.} \src{M0\_second\_referee} \S5 checked that Theorem~\ref{thm:M0closed}(a) conflicts with no known positive result, and located the inputs specific to $T_1$, one for each branch of the proof.
\begin{itemize}[leftmargin=*]
\item \emph{Branch $p^*\notin E$} (Proposition~\ref{prop:chainchords}). Everything up to (PC), (NS), crossing rows and the box holds for any interior target whose facets carry crossings with a shared positive coordinate. The contradiction is the corner-equation and box algebra (Proposition~\ref{prop:T1}; the \src{R15} pinching for all-facing triangles). For covered targets (gauge mixtures, normalised boxes, the hexagon around $u$, the 9-gons) every interior Bayes model is unique ($E=\varnothing$), only this branch is relevant, and the algebra is consistent (\src{R10a} Thm.~M; \src{Conditional Shunning of Error.tex}).
\item \emph{Branch $p^*\in E$} (Propositions~\ref{prop:crosspin}, \ref{prop:carrying}). Proposition~\ref{prop:crosspin} itself holds for every interior target. The inputs specific to all-facing triangles are those of Lemma~\ref{lem:signs}: (i) each facet has \emph{two} crossings in its relative interior, so concavity gives ``$\le0$ on the corner portions''; (ii) each corner ray lies in an orthant with two negative coordinates, so carrying atoms have corner patterns; (iii) the middle arc of $F_k$ is positive at $j$ and $o$, which a corner-pattern chord separates. The boundary cycle of $D_K$ meets each plane $\{g_c=0\}$ exactly twice, so there are six crossings in all, and two relative-interior crossings per facet force exactly three facets. A normalised box has all six crossings at its corner rays; there Proposition~\ref{prop:crosspin} gives $K\subseteq$ box, the opposite inclusion, and no contradiction. In two worlds the boundary rays have no zero coordinate, so Proposition~\ref{prop:crosspin} is vacuous and Theorem~\ref{thm:N26.1} is consistent with it. The systems of \src{Conditional Shunning of Error.tex} with ties are not strictly signed.
\end{itemize}
NUMERICS (\src{m0ref2\_py}): all 2000 random all-facing triangles satisfy (i)--(iii); among 591 random interior polygons with at least four hull vertices, every one fails (i). \emph{Status:} the reading of the two branches is CONFIRMED (\src{M0\_second\_referee} \S5); the extension of Theorem~\ref{thm:M0closed}(a) to every all-facing triangle in general position is CHECKED only. Not established: general decomposable domains, other targets, $n\ge4$.

\subsection{Resolving exceptional ties: two worlds, and the lexicographic class}\label{sec:M0-lex}

Let $\Ltwo$ be the class of \emph{two-level lexicographic rankings with additive levels} on $\M_U$: world $i$ ranks $X$ by the lexicographic order of $(\psi_i(X),S_i(X))$, where the \emph{primary} $\psi_i(X)=\int_{E_i(X)}\rho_i\,d\mu$ has $\rho_i>0$ a.e., and the \emph{secondary} $S_i(X)=\int\sigma_i(g,a_X(g))\,d\mu(g)$ has an integrable envelope and is otherwise arbitrary. The primary endorses \emph{every} comparison at world $i$, so every world adds up, and (E$_\psi$) holds with $\psi$ the primary. A model is \emph{improvable at $i$} if some $Z$ has $\psi_i(Z)=\psi_i(X)$ and $S_i(Z)<S_i(X)$.

\begin{theorem}\label{thm:N26.1}
Let $n=2$ and let $\M_U$ be the $U$-domain with $K'=\{\tfrac12\}$. There is an explicit ranking in $\Ltwo$ that satisfies strict (G), (IB), (LA), (S) (with an explicit selection and closed sections for every model), (C), and (LF) at the target $T\coloneqq D_K$, where $K$ is the set of probabilities whose world-$1$ coordinate lies in $[0.25,0.625]$, such that:
\begin{enumerate}[leftmargin=*,label=(\alph*)]
\item the primary has exactly one exceptional interior weighting $t_0=0.3$, whose tie set is a non-null cone of mixed problems;
\item $T$ is a coherent primary-minimiser at $t_0$ that is neither one-sided limit of the Bayes models (it is intermediate);
\item every incoherent $t_0$-minimiser is strictly dominated, at both worlds, by the coherent model with the same primary profile;
\item the undominated models form a continuous monotone path of coherent models, and the selection index is not the primary's Bayes weighting.
\end{enumerate}
\end{theorem}
\begin{proof}[Proof sketch]
For a mixed pair put $r\coloneqq-g_2/(g_1-g_2)$, which is the same for $g$ and $-g$. Pairs with $r<\tfrac12$ have their $U$-orientation in $H$; pairs with $r>\tfrac12$ have $U$-orientation $-g$, with negative first coordinate, and evaluable models may accept $-g$ there. Take $\rho_1=1-\theta(r)$, $\rho_2=\theta(r)$ on $U$, with $\theta$ nondecreasing on $r<\tfrac12$ and equal to $t_0$ on $[0.2,0.3]$, and strictly increasing on $r>\tfrac12$. The secondaries are built from the path $u\mapsto C_u$ of undominated models so that each $C_u$ is the a.e.-unique pointwise minimiser of $S_1+u\psi_1$ and of $S_2+(3-u)\psi_2$. Then no selected model is improvable at either world, the sections of (S) are primary sublevel sets of a continuous monotone function of $u$, and at $t_0$ a bathtub comparison shows that every incoherent tied model loses strictly at both worlds to the coherent model of equal primary mass.
\end{proof}
\status{proved (\src{N26} \S1, N26.1); confirmed by one referee (\src{N26\_referee} \S1: every axiom re-derived by hand, with independent numerics). It is a genuine M0 mechanism: a ranking that adds up, satisfies every axiom, and resolves an exceptional additive tie that its supporting assessment cannot. In two worlds every target is already covered (\src{Full Coverage \dots v3}), so it adds no coverage. Referee repair to an unused remark (CHECKED): other assessments that endorse only the relevant comparisons can break the tie.}

\begin{theorem}\label{thm:Ltwo}
Let $n=3$, let a ranking in $\Ltwo$ on $\M_U$ satisfy (LA), (S), (C), and let $\psi$ be its primary.
\begin{enumerate}[leftmargin=*,label=(\alph*)]
\item \emph{Multiplier normal form.} At a world $i$ and a level $x$ interior to the range of $\psi_i$, put $g_i(x)\coloneqq\min\{S_i(Z):\psi_i(Z)=x\}$. Every undominated model at level $x$ that is not improvable at $i$ minimises $S_i+\lambda\psi_i$ pointwise for every $\lambda\in-\partial g_i(x)$. If two such models at one level differ on a non-null set, the multiplier at that level is unique; the set of multiplier \emph{values} arising at such levels is countable.
\item \emph{Splice lemma.} If two distinct undominated $t_0$-Bayes models ($t_0\in\Di$) share a multiplier at every world, a bounded splice of them is undominated and incoherent, so (LA), (S), (C) fail.
\item Every interior face of the primary has dimension at most $1$. Hence $E(\psi)$ is a countable union of non-crossing chords; the tie set of a chord $\ell$ is the union of two oppositely oriented atoms, of sign patterns $\sigma$ and $-\sigma$, with chart line $\ell$.
\item \emph{Rigidity.} No chart line of a gamble outside the tie set of $\ell$ crosses $\ell$ inside $\Delta$; the faces at the points of $\ell$ are translates of one fixed segment; $\ell$ cuts off a vertex, and the credal sets of the face models meet both edges at its endpoints.
\item If $T=D_K\in\C_L$ with $K\subseteq\Di$, every interior weighting $p$ at which $T$ is $\psi$-Bayes lies outside $E(\psi)$, and $\operatorname{dist}(p,E(\psi))>0$ whenever $K_q\to K$ as $q\to p$.
\end{enumerate}
Consequently, a ranking in $\Ltwo$ on $\M_U$ that covers $T_1$ has $E(\psi)$ a nonempty countable union of rigid, vertex-cutting, non-crossing chords at positive distance from $p^*$.
\end{theorem}
\status{proved (\src{N26} Lemmas 2.1--2.3, Thm.~3.1, Cor.~3.2, Thm.~3.3, Cor.~3.4, Thm.~4.1, Cor.~4.2, i.e.\ N26.2--N26.7); confirmed with repairs by one referee (\src{N26\_referee} \S\S2--3). Repairs incorporated: in (a), what is countable is the set of multiplier values, not the set of levels carrying a value; in (c), the tie set has the two patterns $\pm\sigma$, not one; in (d), the face is a translate of a fixed segment in direction $d$, not $[y^-,y^-+D]$. \emph{Scope.} These results concern the class $\Ltwo$ with $\psi$ its primary. They do \emph{not} transfer to general M0 rankings: the proofs use the lexicographic structure at four points (undominated implies weakly $\psi$-Pareto; \src{R2} Thm.~7; a multiplier from the convex value function of an additive secondary; the splice is undominated because it is a pointwise Lagrangian minimiser) (\src{N26\_referee} \S3.2, \src{N25\_referee} \S6). For general M0 rankings, ``no isolated exceptional points'' holds by Prop.~F, but ``no two-dimensional faces'' and ``$p^*\notin E$'' are not established. Theorem~\ref{thm:N26.1} shows that (e) is genuinely a three-world fact. Within $\Ltwo$ one may take $\psi$ to be the primary, which lies in $\mathcal N$ and satisfies (E$_\psi$); then R-a$'$ is excluded and only R-b$'$-type configurations remain (\src{N26\_referee} \S3.3). \emph{Round 11:} by Corollary~\ref{cor:M0closed}(c) no ranking in $\Ltwo$ on $\M_U$ covers $T_1$ at all, so the ``consequently'' clause is vacuous for $T_1$; (a)--(e) remain valid as structure theorems for $\Ltwo$ and for other targets.}

\begin{proposition}\label{prop:R1N26}
Let a ranking in $\Ltwo$ on $\M_U$ satisfy (LA), (S), (C), with $T_1\in\C_L$ primary-Bayes at $p^*$. Suppose that the cluster-line conclusion of Theorem~\ref{thm:pivot} holds for the primary's field on the middle range $[\tau_i,\tau_j]$ of each facet: every $h(\tau)$ has a unique oriented cluster line $\lambda_\tau\ni p^*$, $\tau\mapsto\lambda_\tau$ is continuous, and $\lambda_{\tau_i}=\lin(p^*,e_i)$, $\lambda_{\tau_j}=\lin(p^*,e_j)$. Then $E(\psi)=\varnothing$; with Theorem~\ref{thm:main}, no ranking in $\Ltwo$ on $\M_U$ covers $T_1$.
\end{proposition}
\begin{proof}[Proof sketch]
The lines $\lambda_\tau$ sweep every direction through $p^*$, so one of them crosses a given chord transversally; gambles near that $h(\tau)$ have chart lines crossing the chord, against Theorem~\ref{thm:Ltwo}(d).
\end{proof}
\status{proved in a referee report (\src{N26\_referee} \S4.3, Prop.~R1), one reader only. Superseded as a route: its conclusion (no $\Ltwo$ cover of $T_1$ on $\M_U$) is Corollary~\ref{cor:M0closed}(c), which needs neither its hypothesis nor Conjecture~\ref{conj:N26R}.}

\begin{conjecture}\label{conj:N26R}
In $\Ltwo$ on $\M_U$, the cluster-line conclusion of Theorem~\ref{thm:pivot} holds for the primary off a countable union of rigid chords at positive distance from $p^*$. Sharper: a cover of $T_1$ in $\Ltwo$ would need rigid chords whose tie sets accumulate in measure at every point of some crossing ray of $\partial D_{T_1}$ (every ball around such a point meets some chord's tie set in a non-null set).
\end{conjecture}
\status{open as stated (\src{N26} Conj.~N26.R, sharpened in \src{N26\_referee} \S4), but \textbf{superseded} for coverage: its purpose, and the whole lexicographic residual for $\Ltwo$, is settled by Corollary~\ref{cor:M0closed}(c) (no $\Ltwo$ ranking covers $T_1$ on $\M_U$). Original status note: By Proposition~\ref{prop:R1N26} the first sentence implies that no $\Ltwo$ ranking on $\M_U$ covers $T_1$; Theorem~\ref{thm:nosliding}, Corollary~\ref{cor:PCNS} and Proposition~\ref{prop:boxbridge} need no re-check. The second sentence is the referee's sharpened residual N26.R$'$: by Remark~R2 of that report (CHECKED, sketch), Lemma~\ref{lem:PL2} and Lemma~\ref{lem:nopin}, and with them Theorem~\ref{thm:pivot}, survive whenever no chord's tie set meets a small ball around a point of a crossing ray.}

\section{What is excluded and what is open}\label{sec:open}

\subsection{Two concrete questions, both now answered, and the remaining target}\label{sec:two-questions}
An ordinary additive ranking cannot vindicate $T_1$ on the $U$-domain (Corollary~\ref{cor:coverage}), nor, with integral penalties, on any decomposable domain with (LF) (Theorem~\ref{thm:gm}). A more general ranking on that domain could escape this result in either of two ways.
\begin{itemize}[leftmargin=*]
\item \emph{One} (M0): positive error weights endorse its relevant comparisons, but the supporting additive assessment loses distinctions at exceptional weightings, where it ties coherent models with incoherent ones. The original ranking would have to distinguish them in a way compatible with (IB) and closed choice.
\item \emph{Two} (M1 or M2): no positive weighting of errors endorses all the relevant comparisons. This may have a finite witness (M1), or emerge only from the whole collection of comparisons (M2).
\end{itemize}
On the $U$-domain the first escape is now ruled out (Theorem~\ref{thm:M0closed}); the second has been neither ruled out nor turned into a covering construction (\S\S\ref{sec:open-ordinal}, \ref{sec:open-core}). For additive rankings, general two-sided domains no longer offer an escape (Section~\ref{sec:gm}); non-additive rankings on two-sided domains, including rankings that add up, have not been examined. Two questions were concrete:
\begin{enumerate}[leftmargin=*,label=\textbf{(Q\arabic*)}]
\item \emph{Can exceptional additive ties be resolved ordinally?} \emph{Answered on $\M_U$: no.} On $\M_U$, can a ranking with strict (G), (IB), (LA), (S), (C) that adds up at every world cover $T_1$, by resolving the exceptional ties of its supporting assessments (the atom chords of Theorem~\ref{thm:Estructure}) in configuration R-a$'$ or R-b$'$ (\S\ref{sec:M0-residual})? Two things are \emph{not} to be demanded: that the supporting assessment itself satisfy (S) or (C), and that its Bayes weightings equal the selection indices (an auxiliary numerical representation can fail to supply the obvious selection without the original ranking failing closed choice; compare the conditional-probability example of \src{Loss Dominance Generalized v2}). A negative answer would prove that closed choice and independence prevent the needed distinctions, or permit supporting weights without exceptional ties. A positive answer would exhibit a concrete mechanism by which ordinal rankings outperform additive ones. Theorem~\ref{thm:N26.1} supplies such a mechanism in two worlds; in three worlds, Theorem~\ref{thm:Ltwo} and Conjecture~\ref{conj:N26R} describe what is left for lexicographic rankings. For coherent almost-preference rankings the same question is (U) (Conjecture~\ref{conj:U}, Remark~\ref{rem:EPchords}).
\emph{Answer.} No ranking on $\M_U$ that adds up at every world covers $T_1$, in either configuration and with or without (E$_\psi$) (Theorem~\ref{thm:M0closed}(b), Corollary~\ref{cor:M0closed}). The exceptional ties are not the obstacle: by the additive form (Theorem~\ref{thm:M0closed}(a)), no strictly signed field whose Bayes models are coherent at a.e.\ interior weighting has $T_1$ as a Bayes model at any interior weighting, exceptional or not, so no ordinal tie-breaking can help. The ranking form holds at the status of Proposition~\ref{prop:M0} ($\alpha$), ($\gamma$) without (E$_\psi$); the new parts of the additive form have had two independent readings. The two-world mechanism of Theorem~\ref{thm:N26.1} does not transfer: in two worlds the boundary rays of the target have no zero coordinate, so crossing pinning (Proposition~\ref{prop:crosspin}) has nothing to act on. For (CP-a) rankings on $\M_U$ the answer is the same (Corollary~\ref{cor:CPaU}, proved; glue refereed once). \emph{Domain caveat:} this is proved on the $U$-domain only; on general decomposable domains only the additive theorem (Theorem~\ref{thm:gm}) is proved.
\item \emph{Can additive rankings on two-sided domains cover $T_1$?} Equivalently (Proposition~\ref{prop:CF}), does CF$(T_1)$ have a solution? \emph{Answered: no.} No strictly guiding additive ranking with integral penalties covers $T_1$ on any decomposable domain with (LF), for any $\varepsilon\in(0,1)$ and any interior target weighting (Theorem~\ref{thm:gm}, three independent proofs, each refereed once), and CF$(T_1)$ has no solution (Corollary~\ref{cor:CFnone}). The case left open in the previous version, sector caps at the failing edges glued to a non-precise layer along the crossing chords, is excluded: the identity caps are part of no cover (Corollary~\ref{cor:nocap}), and the chord layer is not what was described (Corollary~\ref{cor:chord}). What remains:
\begin{itemize}[leftmargin=*]
\item \emph{The converse direction.} The ranking form of the answer needs no appeal to Theorem~\ref{thm:ccibu}. The CF form (Corollary~\ref{cor:CFnone}), for solutions not known to come from a ranking, holds at the status of Theorem~\ref{thm:ccibu}, through Proposition~\ref{prop:CF}($\Leftarrow$); Theorem~\ref{thm:ccibu} is author-checked and has no formal referee report.
\item \emph{Integral versus charge penalties.} For bounded charges vanishing on null sets the answer is proved only given interior Bayes minimisers, and has been read along one route only (Corollary~\ref{cor:gmcharges}); whether coverage forces such minimisers to exist when the penalties are charges is open (\src{N14} \S5.1), except for complete rankings (Theorem~\ref{thm:complete}).
\item \emph{Readings.} Each route has had one referee; Theorem~\ref{thm:D} has had none (its dichotomy has two independent one-reader proofs). A second reading of one complete route would promote the theorem (\S\ref{sec:gm-deps}).
\end{itemize}
\end{enumerate}
\paragraph{The remaining target.} On the $U$-domain the \emph{only} remaining route to an ordinal cover of $T_1$ is a ranking with strict (G), (IB), (LA), (S), (C) and $T_1\in\C_L$ in which some world does not add up on its relevant comparisons, case (M1) or (M2) (Corollary~\ref{cor:M0closed}(a)). What is known about such a world:
\begin{itemize}[leftmargin=*]
\item it is not a coherent almost-preference over losses (Theorem~\ref{thm:fullsupport}: a member with a.e.-positive countably additive part endorses its comparisons; any $n$; this is Route B of Corollary~\ref{cor:CPaU}(b));
\item if it is flexible, it is not cell-fine and has no finite obstruction, so its failure is lexicographic layering (Proposition~\ref{prop:flexible}); two-level layering with additive levels is single-level additive on every injective region free of exceptional weightings (Theorem~\ref{thm:LX}); and in the class $\Ltwo$ no cover exists at all (Corollary~\ref{cor:M0closed}(c));
\item it has no leading charge nonzero on $M$, so in particular it is not lexicographic with charge levels, in any number of worlds and with no nestedness hypothesis (Theorem~\ref{thm:S21}, round 13);
\item in the two-world analogue, a partner world with a nonzero leading charge or monotone continuity (for instance an additive one) makes the selection nested automatically (Theorems~\ref{thm:S22}, \ref{thm:MC}); with a nested selection the world must moreover be incomplete (Theorem~\ref{thm:E3}, CHECKED) and have spread-out certificates (\S\ref{sec:open-ordinal}(c)); its uncharged sets are initial segments, and a simple certificate forces a mirrored one at the other world (\S\ref{sec:open-core}(c));
\item in three worlds, when the selection is edge-nested (Nest$_e$), such a world of a cover of $T_1$ is moreover incomplete (Theorem~\ref{thm:Q36}(c), Corollary~\ref{cor:Q37}); the finer two-world restrictions (strict escape, Lemma~T, uncharged initial segments, mirrored simple certificates) are proved only for its comparisons between subsets of a single $R_{ij}$, i.e.\ for the edge system (Proposition~\ref{prop:Q35}), and say nothing about comparisons out of interior selected models, whose error sets cross; without (Nest$_e$), or a nested edge family replacing it (Remark~\ref{rem:straighten}), none of these finer restrictions transfers; and every cover shows non-additive behaviour at $T_1$ itself (Theorem~\ref{thm:R3}).
\end{itemize}
So the precise target is a three-world ranking on $\M_U$ covering $T_1$ with a world that does not add up, has no leading charge nonzero on $M$ and is not (CP-a) (both unconditionally), and is, in the sense of the two-world results and of \src{P2}, \emph{incomplete and non-lexicographic}, or whose edge selections are not nested; equivalently, a counterexample to Conjecture N17.B$'$ (\src{N19}) in three worlds. An obstruction would prove N17.B$'$ for covers of $T_1$. (Round 16: N17.B$'$ is false in two worlds, Theorem~\ref{thm:D21}; so a proof of it for covers of $T_1$ must use an input that the two-world system of that theorem lacks, such as the third world or the target.) The two-world restrictions do transfer to three worlds along the edges of $\Delta$, where complementarity is automatic (Lemma~\ref{lem:Q34}), under edge nestedness (Nest$_e$); what does not transfer automatically is the passage from edge comparisons to comparisons out of interior selected models, whose error sets cross. \emph{Summary.} For $T_1$ on $\M_U$, every ranking class examined so far is excluded except rankings with a world that does not add up; the decisive remaining problem is a two-world question, (R1), (R2), (R4) and Conjecture~\ref{conj:NDL} of \S\ref{sec:open-core}(g), together with edge nestedness (or straightening) for complete worlds, with the two further steps stated there (the passage to interior comparisons for general acceptance cones, and case (i) of Theorem~\ref{thm:R3}). Rounds 14 and 15 sharpen this target: such a world also lacks a leading density on one of its edge regions and is not a glued credal world; the updated list is \S\ref{sec:open-core}(g), summarised in the introduction. Round 16 shows that a world that does not add up can occur in a full system (two worlds, Theorem~\ref{thm:D21}), and that the protected part of the mechanism by which it does so is invisible to the argument of Theorem~\ref{thm:M0closed} at ends on which $T_1$ is eventually correct at the protected world; covers in which every world is an end-protected decoration with integrable, jointly end-heavy base densities are excluded, and nothing proved excludes the other installations (Theorems~\ref{thm:E11}, \ref{thm:E15}; audited in \src{F2}, audit refereed once; \S\ref{sec:open-core}(g)). Round 17 removes the heaviness hypothesis, excludes failures that can be localized next to two worlds with integrable full endorsers, and reduces the case of finitely many end-protected worlds with integrable base endorsers to conflict cycles with no aligned order of the worlds (\S\ref{sec:open-core}(f$'$), (g)). Round 19 classifies the worlds that do not add up into those with a finite certificate, those with a non-null non-localizable blind set, and those that are localizable, and says what a cover needs in each case (\S\ref{sec:open-core}(f$''$), (g)).

For rankings beyond the additive class, the present results do not force exact coverage to be replaced by approximation.

\subsection{Weakened guidance}\label{sec:open-ties}
Two target-independent weakenings of (G) were tested, both allowing good advice and silence to tie at world $i$ on mixed problems with small world-$i$ stake, $|g_i|\le\theta\|g\|$ (the band).
\begin{itemize}[leftmargin=*]
\item \emph{Forced ties} (\src{BRIEF3}). For additive rankings on decomposable domains, ``NSI$_\theta$ covers $T_1$'' fails for every $\theta>0$, for both the sup and the Euclidean norm, every $\varepsilon\in(0,1)$, and indeed for $\varepsilon$-contamination of every interior baseline in every $n\ge3$. \emph{Status:} proved, with two independent complete proofs (\src{N1}, \src{N2}); refereed once (\src{N4}). Scope: additive rankings only; forced-tie reading only.
\item \emph{Permitted ties} (\src{BRIEF4}; ties allowed, placed by the ranking). The conditional no-go of \src{R15} Thm.~6.2 (under the first-order hypotheses (DF), (ND) and non-adjacency) transfers verbatim for every $\theta>0$. \emph{Status:} proved; refereed once (\src{N7}, on \src{N5}, \src{N6}). The unconditional chain of \S\ref{sec:main} also transfers, below an explicit threshold (Theorem~\ref{thm:permitted}).
\end{itemize}
In both readings (C) forces (BU) at every interior weighting for every $\theta$ and both norms, because \src{R1} Lemma 1.4 uses no guidance; the caveat ``$\theta<1$ (sup), $\theta<1/\sqrt3$ (Euclidean)'' of \src{BRIEF4}, \src{N5} and \src{N7} is unnecessary (\src{N27} Prop.~27.4, confirmed in \src{N27\_referee} \S3.1).

A \emph{permitted-tie realisation} of $T$ at $p^*$ on $\M_U$ is a measurable $c:U\to\R^3$ with $c_i(g)\ge0$ if $g$ is mixed with $0<g_i\le\theta\|g\|$, $c_i(g)>0$ if $g_i>\theta\|g\|$ or $g$ is unanimous, $c_i(g)<0$ if $g_i<0$, together with (BU) at every $q\in\Di$ and (TP). For a facet $F_k$ and a crossing coordinate $i\ne k$ of $F_k$ consider
\begin{itemize}[leftmargin=*]
\item[\textbf{(X$_{ik}$)}] a non-null set of $x\in U$ has $v_k\cdot x<0$ and $c_i(x)>0$;
\item[\textbf{(VP)}] for every $x\in\relint F_k$ and $l\ne k$ with $x_l>0$, $\liminf_{q\to e_l,\,q\in\Di}P(q,x)>0$ (vertex positivity),
\end{itemize}
and put $\theta_x(T)\coloneqq\min_k\min_{i}\sup_{h\in F_k\setminus0}h_i/\|h\|$, the minimum over facets $F_k$ and crossing coordinates $i\ne k$.

\begin{theorem}\label{thm:permitted}
Let $T\subseteq\Di$ be an all-facing triangle in general position, $p^*\in\Di$, $\theta>0$, and the norm sup or Euclidean. If (X$_{ik}$) holds for every facet $F_k$ and every crossing coordinate $i$ of $F_k$, or if (VP) holds, then $T$ has no permitted-tie realisation at $p^*$ on any $U$-domain $\M_U$. (X$_{ik}$) holds automatically when $\theta<\theta_x(T)$. Consequently no additive ranking with permitted ties, integrable envelopes, (LA), (S), (C) on $\M_U$ covers $T$ in that range. For $T_1$, at every $\varepsilon\in(0,1)$ and every interior target weighting:
\begin{enumerate}[leftmargin=*,label=(\alph*)]
\item sup norm: $\theta_x=1$, so the conclusion holds for every $\theta\in(0,1)$, exactly the range in which the band is a proper subset of the mixed problems;
\item Euclidean norm: $\theta_x=\sqrt{(\kappa^2+1)/(\kappa^2+2)}$ ($\approx0.837$, $0.928$, $0.999$ at $\varepsilon=0.05$, $0.3$, $0.9$), so the conclusion holds for every $\theta<\sqrt{(\kappa^2+1)/(\kappa^2+2)}$, and in particular for every $\theta<\sqrt{2/3}$ uniformly in $\varepsilon$.
\end{enumerate}
\end{theorem}
\begin{proof}[Proof sketch]
A covering ranking yields a permitted-tie realisation ((BU) from \src{R1} Lemma 1.4, $p^*$ interior by \src{N1} Prop.~2.3, which uses only strict badness). In the chain of \S\ref{sec:main} the only input sensitive to ties is Lemma~\ref{lem:UVN}(a) at relative-interior facet points: positivity near $e_l$ of a facet gamble with $x_l>0$. It is equivalent to (VP) at $x$, and \src{N5} Thm.~3.3 (facet clause) supplies it from (X$_{lk}$). Every other sign use is strict badness ($c_i<0$ where $g_i<0$, never tied) or holds with weak signs (Lemma~\ref{lem:PL2}, Step 4 of Theorem~\ref{thm:pivot}, Proposition~\ref{prop:boxbridge}(a)). The corner-ray gap found in \src{N7} is never used. The transposition of Proposition~\ref{prop:T1} keeps the band, since both norms are permutation invariant. For $T_1$ in the Euclidean norm the maximiser of $h_l/|h|$ on $F_k$ is the normalised projection of $e_l$ onto $v_k^\perp$, which lies in $\relint F_k$, and $\theta_x^2-\tfrac23=(\kappa^2-1)/(3(\kappa^2+2))>0$.
\end{proof}
\status{proved (\src{N27} Thm.~27.6 with Prop.~27.4, Lemmas 27.5, 27.5$'$); confirmed by one referee (\src{N27\_referee} \S\S3, 6; thresholds recomputed symbolically, numerics reproduced). It rests on the facet clause of \src{N5} Thm.~3.3 (refereed once, \src{N7}). The referee adds that the Euclidean maximiser maximises $x_l/|x|$ over the whole half-space $\{v_k\cdot x\le0\}$, so enlarging $U$ does not raise the threshold for $T_1$.}

\begin{corollary}\label{cor:permitted-nonadj}
For $\theta<\theta_x(T)$, no additive ranking with permitted ties covers $T$ on any decomposable domain with (LF) (arbitrary open $N$) that is non-adjacent near $T$.
\end{corollary}
\status{CHECKED (\src{N27} Cor.~27.8, through \src{N27} Lemma 27.7, proved: the set of gambles accepted at some interior weighting is a.e.\ a countable union of the cones $D_{K_q}$, so the set of acceptances given by Bayes models is conic without vertex facts); confirmed at CHECKED level, with two wording repairs, by one referee (\src{N27\_referee} \S4, which also supplied the half-plane form of non-adjacent charts for advice range $\{+,0,-\}$ under weak signs).}

\emph{What remains open.}
\begin{itemize}[leftmargin=*]
\item \emph{The Euclidean window $\theta\in[\theta_x(T_1),1)$.} There no gamble at all beyond a facet plane is off-band at the crossing coordinate, so (X$_{lk}$) holds exactly when the ranking leaves untied a non-null set of $l$-positive points of $U\cap\{v_k\cdot x<0\}$; it is not automatic, but it does not fail automatically either. The open case is the \emph{all-tied configuration}: for some crossing pair $(l,k)$, $c_l=0$ a.e.\ on $U\cap\{x_l>0,\ v_k\cdot x<0\}$. The missing input is (VP); equivalently, every Hausdorff cluster set $W$ of $K_q$ as $q\to e_l$ satisfies $W\subseteq\{m:m\cdot w_{lk}\ge0,\ m\cdot r>0\}$, with $w_{lk}$ the $l$-crossing of $F_k$ and $r$ the corner ray at the $l$-positive end of $F_k$ (CHECKED). A directional vertex limit $W'=\co(K'\cup\{e_l,v_k\})$ meets the local vertex requirements and violates (VP) on the $l$-positive arc of $F_k$ (CHECKED, local; this repairs \src{N27}'s $W=[e_l,v_k]$, which violates $K_q\supseteq K'$). \emph{Conjecture} (Euclidean norm, unsupported either way): no permitted-tie realisation of $T_1$ exists for $\theta\in[\theta_x(T_1),1)$. In the sup norm nothing is open for $\theta\in(0,1)$. The window needs $\theta>\sqrt{2/3}\approx0.816$, far from negligible stakes.
\item Not examined: decomposable domains that are not non-adjacent near the target, under permitted ties, and the charge version (Corollary~\ref{cor:charges}) under permitted ties.
\end{itemize}

\subsection{Lexicographic and other non-additive rankings}
\begin{itemize}[leftmargin=*]
\item Two-level lexicographic rankings with additive levels: interior lexicographic Bayes models are a.e.\ unique and coherent (\src{N9} N9.5, \src{R11} Prop.~LX). A face-continuum theorem kills every design in which a secondary trades off against another world's primary at a continuum of rates (\src{N9} N9.6). The strongest plateau design (lexicographic Bayes map constant on $\Di$) and the free-collar design both fail, the latter with an explicit incoherent witness (\src{N9} N9.7, N9.8). \emph{Status:} proved; refereed once (\src{N10}). Coverage of $T_1$ by general lexicographic rankings is open; the precise residual is in \src{N10} \S4 (F4). For the class $\Ltwo$ on $\M_U$ the exceptional set of the primary consists of rigid, vertex-cutting, non-crossing chords away from the target weighting (Theorem~\ref{thm:Ltwo}); the question formerly left there (Conjecture~\ref{conj:N26R}) is superseded: no $\Ltwo$ ranking covers $T_1$ on $\M_U$ (Corollary~\ref{cor:M0closed}(c)). Two-level layering with a weakly signed primary is single-level additive on injective regions free of exceptional weightings (Theorem~\ref{thm:LX}).
\item Arbitrary (IB) preorders: a long list of necessary structure is proved (vertex anchoring, interior target index, closed tie fibres, edge repulsion; \src{N8}, \src{N9}, refereed in \src{N10}), but none of it excludes coverage. By \src{N9} N9.9 any ranking-level impossibility contains the additive one, which \S\ref{sec:main} now supplies on non-adjacent domains. \emph{Rankings that add up} (M0) are analysed in \S\S\ref{sec:M0}--\ref{sec:M0-closed}: on $\M_U$ they never cover $T_1$ (Theorem~\ref{thm:M0closed}), so question (Q1) is answered. \emph{Status:} for arbitrary (IB) preorders, open; what remains is the route of the next item.
\item \emph{Rankings that do not add up} ((M1) or (M2) of \S\ref{sec:addup-def}). On $\M_U$, a cover outside (M0) needs a world with no weight, nonnegative on the relevant comparisons, whose countably additive part is positive a.e.\ on mixed problems (condition $\neg$SP$^{\mathrm{reg}}$; \src{N17}, \src{N19}). This arises either through a finite obstruction (M1), whose certificates in the stated formulation use at least four comparisons with overlapping gains and crossing error sets, or with no finite obstruction at all (M2). An (M1) certificate need not be an executable sequence of losing trades, and the restrictions just listed do not constrain (M2). \src{N19} states the necessary condition as conditional on excluding the integral class of \src{N13} Prop.~B, i.e.\ integral fields with unique coherent Bayes models at \emph{almost every} weight. Theorem~\ref{thm:main} removes the regularity part of that condition; the ``almost every'' part was exactly the (M0) residual of \S\ref{sec:M0-residual} (as framed in \src{BRIEF8} and \src{N25} \S10), and for coherent almost-preferences it is (U) (an identification made here, not in the sources); on $\M_U$ both are now closed (Theorem~\ref{thm:M0closed}, Corollary~\ref{cor:CPaU}), so this route is the only one left there (\S\ref{sec:two-questions}, ``The remaining target''). Until round 16 no system was known, for any target, in which (M1) or (M2) occurs on the relevant comparisons together with (S). Round 16 gives one: a full two-world system on $\M_U$ in which (M2) occurs at one world (Theorem~\ref{thm:D21}). No such system is known in three worlds. The direct construction attempt of round 11, and the restrictions it produced, are in \S\ref{sec:open-ordinal}; the round-12 results on this route are in \S\ref{sec:open-core}. \emph{Status:} route open; \src{N19} Conj.~N17.B$'$ is false in two worlds (Theorem~\ref{thm:D21}, round 16); before that it was narrowed in two worlds (Theorems~\ref{thm:E2}, \ref{thm:E3}, \S\ref{sec:open-core}(c), (d)) and in three (Theorem~\ref{thm:Q36}: for leading charges unconditionally, for complete worlds under edge nestedness).
\end{itemize}

\subsection{Building ordinal rankings directly}\label{sec:open-ordinal}
By Corollary~\ref{cor:M0closed}(a), an ordinal cover of $T_1$ on $\M_U$ needs a world that does not add up. Round 11 (\src{BRIEF10}) also tried to build such a cover directly: preorders, one per world, with strict (G), (IB), (LA), (S) (an explicit coherent selection, closed sections for \emph{every} evaluable model, competitiveness on $\supp p$), (C) and $T_1\in\C_L$, or else a genuine obstruction. No cover was constructed. This subsection records the tools and the obstructions found (\src{P2}, \src{P3}, \src{P4} Part B; referee reports \src{P2P3\_referee}, \src{P4B\_referee}).

\paragraph{(a) The exact consistency test and the limit rule.} The exact consistency test of \src{Coverage after Neutral Abstention} (\S3, Thm.~closure) fixes a coherent selection containing the target and closed sets $A_i(X)$ satisfying the covering condition (for every $p$ and $X$ some $i\in\supp p$ has $p\in A_i(X)$). It closes the installed comparisons under four rules: reflexivity, transitivity, (IB)-transfer, and (R4) if $X\preceq_iC_p$ at every $i$, install $C_p\preceq_iX$ at every $i$. Its forbidden outcomes are (a) reversal of a strict guidance improvement, (b) $C_p\preceq_iX$ with $p\notin A_i(X)$, and (c) an incoherent model tied at every world with a selected one. No forbidden outcome occurs iff preorders with (G), (IB), (S), (C) exist with exactly those sections. For a construction its drawback is that the sections must be guessed in advance.

Fix a decomposable domain $\M$ and a coherent selection $C$. A \emph{successful system for $C$} is a family of preorders with (G), (IB), every $C_p$ undominated and competitive on $\supp p$, every section $\{p:C_p\preceq_iX\}$ closed, and (C). A \emph{choice} $\kappa$ assigns to each pair $(p,X)$ a world $\kappa(p,X)\in\supp p$. The \emph{limit rule} is (R5): if $p_k\to p$ and $C_{p_k}\preceq_iX$ for all $k$, install $C_p\preceq_iX$. Let $\mathrm{Cl}_5(\kappa)$ be the least family of relations that contains the background order and the chosen comparisons $C_p\preceq_{\kappa(p,X)}X$ and is closed under the five rules (it exists by Knaster--Tarski and is reached by stage $\omega_1$); call it \emph{clean} if it has no outcome (a) and no outcome (c).

\begin{theorem}\label{thm:R5}
\begin{enumerate}[leftmargin=*,label=(\arabic*)]
\item If $\kappa$ is total and $\mathrm{Cl}_5(\kappa)$ is clean, then $\mathrm{Cl}_5(\kappa)$ is a successful system for $C$.
\item If $(\preceq_i)$ is a successful system for $C$ and $\kappa(p,X)$ is any $i\in\supp p$ with $C_p\preceq_iX$, then $\mathrm{Cl}_5(\kappa)\subseteq(\preceq_i)$ and $\mathrm{Cl}_5(\kappa)$ is clean.
\item Hence a successful system for $C$ exists iff some total $\kappa$ has $\mathrm{Cl}_5(\kappa)$ clean.
\end{enumerate}
\end{theorem}
\begin{proof}[Proof sketch]
(1) The rules give preorders and (IB); the background order gives weak (G), and strict (G) because its reversal is outcome (a); (R4) gives undominatedness; totality gives competitiveness; (R5) gives closed sections, since $\Delta$ is metrizable; (C) follows from Theorem~\ref{thm:kkm} and the absence of outcome (c). (2) A successful system contains the generators and is closed under the five rules (under (R4), since otherwise $X$ would dominate $C_p$; under (R5) by closed sections), so it contains the least closure; outcomes (a) and (c) would be inherited.
\end{proof}
\status{proved (\src{P3} Thm.~3.1); confirmed in both directions (\src{P2P3\_referee} \S A2: (R5) interacts correctly with (IB)-transfer and (R4)). The theorem is exact but formal: the existence of a clean total $\kappa$ is as hard as the original problem. Its gain over Thm.~closure is that no sections are prescribed; closed sections are produced by an infinitary rule, and the only free data are the choices. Related, also proved and confirmed (\src{P3} Thm.~4.1, \src{P2P3\_referee} \S A3): on $\M_U$, all requirements except closed sections are of finite character (propositional compactness), and compactness fails exactly at (R5). Whether $T_1$ on $\M_U$ admits a system with every requirement except closed sections is therefore a finite-character question; it is open.}

On $\M_U$ a comparison $X\preceq_iY$ depends only on the \emph{exchange} $\sigma=1_X-1_Y$, a measurable function $U\to\{-1,0,1\}$. Installing $C_p\preceq_iX$ therefore installs $C_q\preceq_iC_q-\sigma$ at every $q\in\mathrm{Comp}(\sigma)\coloneqq\{q:\{\sigma=1\}\subseteq C_q,\ \{\sigma=-1\}\cap C_q=\varnothing\text{ a.e.}\}$, and every successful system satisfies \emph{exchange transversality}: the set of worlds at which $\sigma$ is accepted meets $\supp q$ for every $q\in\mathrm{Comp}(\sigma)$ (\src{P3} Lemma~2.1, Cor.~2.2, \S5; proved, confirmed on $\M_U$ in \src{P2P3\_referee} \S\S A1, A4, with a repair to the general-domain formula: $\mathrm{Comp}(e)=\{q:a_{C_q}=a_{C_p}\text{ on }Z\}$).

The extension lemma as posed in \src{Vacuous Targets and Transfinite Construction} \S5 reads, in the form formalised in \src{P3} \S1: a \emph{state} consists of relations that contain the background order, are closed under the four rules and produce no forbidden outcome, together with relatively open \emph{reserved sets} $R_i(X)$ on which $C_q\preceq_iX$ is permanently forbidden, satisfying the per-model covering condition (every $p$ has some $i\in\supp p$ with $p\notin R_i(X)$); the claim is that every unmet pair $(p,X)$ can be met at some unreserved $i\in\supp p$ without a forbidden outcome, unions being taken at limit stages.

\begin{proposition}\label{prop:extcounter}
Let $\rho>0$ be bounded measurable, let world $i$ rank by $\psi^\rho_i(X)\coloneqq\int_{E_i(X)}\rho|h_i|\,d\mu$ on $\M_U$ (for $T_1$, any $\varepsilon$), and select $q\mapsto B_q=U\cap\{q\cdot h>0\}=D_{\co(K'\cup\{q\})}$. This is a successful system for every such $\rho$, and:
\begin{enumerate}[leftmargin=*,label=(\alph*)]
\item there is a clean state whose reservations satisfy the per-model covering condition for every model, with an unmet pair that cannot be met; so the extension lemma as posed is false;
\item there is an increasing sequence of partial choices $\kappa_m$, each contained in a successful system for the same selection (so each $\mathrm{Cl}_5(\kappa_m)$ is clean), whose union has no successful extension: at the limit point, (R5) reverses a strict guidance comparison.
\end{enumerate}
\end{proposition}
\begin{proof}[Construction]
(a) Let $A$ be the union of two small balls around $a=(-0.1,1,1)$ and $a'=(1,-0.1,-0.1)$, inside $D_{T_1}\cap B_r$ with $r=(\tfrac12,\tfrac12,0)$, and $\sigma\coloneqq1_A$. Reserve world~1 for $X_r=B_r-\sigma$ near $r$, and worlds 2 and 3 for $X_u=U\setminus A$ near $u$. Every model satisfies the per-model covering condition, yet $(u,X_u)$ cannot be met: worlds 2 and 3 are reserved at $u$, and meeting it at world~1 would accept $\sigma$ at world~1, which by transport is the reserved comparison $B_r\preceq_1X_r$ at $r$. (b) Let $p_0\in\Di\setminus K'$, $r_k\to p_0$ along the segment towards $e_1$, $W'\coloneqq U\cap\{h_1<0\}\cap\{r_1\cdot h>0\}$ and $X\coloneqq B_{p_0}\setminus W'$, a strict world-1 improvement on $B_{p_0}$. Let $\kappa_m$ install $B_{r_k}\preceq_1X$ for $k\le m$; it is realised by $\psi^\rho$ with a weight of order $1/\delta_m$ on part of the shell $B_{p_0}\setminus B_{r_1}$ ($\delta_m$ the distance from $r_m$ to $p_0$). The union, through (R5), installs $B_{p_0}\preceq_1X$.
\end{proof}
\status{proved (\src{P3} \S6, Examples~1 and~2); confirmed, with independent Monte Carlo checks of the geometry (\src{P2P3\_referee} \S\S A5, A6). \emph{The referee's narrowing.} (i) The selection does not contain $T_1$, and must not; these are counterexamples to a general lemma, not obstructions for $T_1$. (ii) What (b) proves is that neither ``extendable to a successful system with the same selection'' nor ``$\mathrm{Cl}_5$ of the stage is clean'' is preserved by unions of $\omega$-chains of partial choices, so a transfinite construction needs an invariant that unions preserve (the reserved-neighbourhood invariant of \src{Vacuous Targets and Transfinite Construction}, or containment in one fixed successful system); the note's stronger reading, that a construction must commit with foresight that finite information does not supply, is not proved. The repaired lemma with \emph{exchange-compatible} open reservations (\src{P3} EL$^*$) is open in general; its simplest two-world dead ends are excluded (proved in \src{P3}).}

\begin{theorem}\label{thm:witness}
Let $\psi$ be an additive ranking with strictly guiding integral penalties and integrable envelope, $\Psi$ its world preorders, and $C$ a coherent selection with: (W1) for every $p$ and $X$ some $i\in\supp p$ has $\psi_i(C_p)\le\psi_i(X)$; (W2) no $X$ has $\psi(X)\le\psi(C_p)$ coordinatewise with $\psi(X)\ne\psi(C_p)$; (W3) $p\mapsto\psi_i(C_p)$ is lower semicontinuous for each $i$. Let $\kappa$ be $\psi$-admissible ($\psi_{\kappa(p,X)}(C_p)\le\psi_{\kappa(p,X)}(X)$). Then $\mathrm{Cl}_5(\kappa)\subseteq\Psi$ with no outcome (a), and $\mathrm{Cl}_5(\kappa)$ is clean iff no incoherent $X$ with $\psi(X)=\psi(C_p)$ is tied with $C_p$ at every world in $\mathrm{Cl}_5(\kappa)$. If moreover (W4) no incoherent $X$ has $\psi(X)=\psi(C_p)$ for any $p$, then every $\psi$-admissible $\kappa$ gives a successful $\mathrm{Cl}_5(\kappa)$, and $\Psi$ itself is a successful system.
\end{theorem}
\status{proved (\src{P3} Thm.~7.1); confirmed (\src{P2P3\_referee} \S A7), which remarks that the content is close to tautological: if a $\Psi$ with (W1)--(W4) exists, it is already a successful system, and the useful part is that the (C)-tie is the only failure mode. \emph{Scope.} It reproduces the precise and Brier covers, normalised boxes and gauge mixtures (confirmed), and the two-world intervals of \src{Full Coverage \dots v3} at the level of the cited statement. It cannot produce an additive cover of $T_1$ (Theorem~\ref{thm:gm}). It never reaches the vacuous model under strict guidance on the all-model domain with $n=3$ (\src{Three Complete Rankings Cannot Cover}), or $n\ge3$ (\src{Complete Rankings for Any Number of Worlds}); this is not shown for an arbitrary decomposable domain (repair R-P3c). A witnessed cover of $T_1$ on $\M_U$ would be an (M0) ranking with (E$_\psi$) (\src{P3} \S8.2, confirmed), so it is excluded by Theorem~\ref{thm:M0closed}. An additive \emph{pre-system} through $T_1$ ((W1)--(W3) with $T_1$ selected, without (C)) has ($\gamma$) by the KKM argument of \src{P3} \S8.3 and ($\alpha$) by the referee's derivation (repair R-P3d); so Theorem~\ref{thm:M0closed}(a) excludes it too (an observation made here, CHECKED, not refereed; \src{P2P3\_referee} \S C.3 had excluded pre-systems with $E(\psi)=\varnothing$).}

\begin{conjecture}\label{conj:villegas}
Let world $i$ be complete, and let the selection be continuous in measure, with world-$i$ error sets nested along monotone selection paths whose shell exhaustions $S_k\uparrow S$ generate the measure algebra of $U$. Then (R5) forces a countably additive representation of $\preceq_i$.
\end{conjecture}
\status{CONJECTURE (\src{P3} P3.C1; label confirmed in \src{P2P3\_referee} \S A8). Proved: (R5) gives Villegas's monotone continuity~\cite{V1964}, but only along the sequences $A'\cap S_k$. If true, covers of $T_1$ with complete worlds and such selections would be complete (CP-a) systems, excluded by Theorem~\ref{thm:complete}. \emph{Referee's remarks} (R-P3e): (i) nested error sets along paths are non-generic in three worlds, since by Lemma~\ref{lem:nested} they force the credal sets to share both tangent lines from $e_i$, a codimension-2 condition; the general content of (R5) is continuity of $\preceq_i$ along $Y_k\to Y$ without monotonicity. (ii) Countable additivity is more than needed: Theorem~\ref{thm:complete} already handles complete worlds represented exactly by a finitely additive state, so the nearby target is ``(R5) forces exact representability by a state''. In two worlds with a nested selection, complete worlds are handled by Theorem~\ref{thm:E3}.}

\paragraph{(b) Exchange rates and layering.}

\begin{lemma}\label{lem:vectorform}
On $\M_U$, with $\sigma_i(h)\coloneqq\operatorname{sign}h_i$, every $X,C\in\M_U$ satisfy $1_{E_i(X)}-1_{E_i(C)}=\sigma_i\cdot(1_C-1_X)$ at every world. Hence $C\preceq_iX$ iff $\sigma_i(1_C-1_X)$ lies in the acceptance set $\mathcal A_i$ of world $i$, and the exact test reads: choose sets $\mathcal A_i$ of measurable $\{-1,0,1\}$-valued functions that contain the nonnegative ones, contain no nonzero function $\le0$, and are closed under sums that stay $\{-1,0,1\}$-valued; put $A_i(X)\coloneqq\{p:\sigma_i(1_{C_p}-1_X)\in\mathcal A_i\}$; and check that these sets are closed, that the covering condition holds, that the (R4)-closure stays inside $\mathcal A_i$, and outcome (c).
\end{lemma}
\status{proved (\src{P2} P2.1, Lemma~1.1); confirmed (\src{P2P3\_referee} \S B1).}

Fix a finite partition $\mathcal Q=(Q_1,\dots,Q_N)$ of $U$ into non-null cells and a finite non-atomic $\nu$ equivalent to $\mu$. For $\delta$ as above put $m(\delta)\coloneqq(\nu(\{\delta=1\}\cap Q_c),\nu(\{\delta=-1\}\cap Q_c))_{c\le N}$. A world is \emph{$(\mathcal Q,\nu)$-flexible} if acceptance of $\delta$ depends only on $m(\delta)$ and survives scaling $m(\delta)$ by any $s\in(0,1]$; it is \emph{cell-fine} if, for all cells $c\ne c'$, gaining mass $s$ in $Q_c$ against losing mass $\epsilon s$ in $Q_{c'}$ is accepted for some $\epsilon>0$ and all small $s>0$.

\begin{proposition}\label{prop:flexible}
For a flexible world: (a) the accepted mass vectors are the trace of a convex cone; (b) there is $w\ge0$, $w\ne0$, with $\sum_cw_c\bigl(\nu(\{\delta=1\}\cap Q_c)-\nu(\{\delta=-1\}\cap Q_c)\bigr)\ge0$ for every accepted $\delta$; in particular the world has no finite obstruction, so (M1) is impossible at it; (c) if it is cell-fine, $w$ can be taken with every $w_c>0$, so the world adds up on all comparisons; (d) in general, every accepted $\delta$ is endorsed by a lexicographic hierarchy of additive levels with cellwise-constant densities, the first nonzero level being positive on each cell.
\end{proposition}
\status{(a)--(c) proved (\src{P2} Lemma~3.1(a)--(d)); confirmed (\src{P2P3\_referee} \S B4). (d) CHECKED (\src{P2} Lemma~3.1(e), via lexicographic separation~\cite{MLS1987}); confirmed at CHECKED level. \emph{Scope:} flexibility is invariance under a \emph{fixed finite} partition; additive worlds with non-cellwise-constant densities are not flexible, and ``flexible violations of adding up are lexicographic layering'' holds for this class only.}

Consider lexicographic worlds: world $i$ ranks by $(F_i,S_i)$ lexicographically, with $F_i(X)=\int\varphi_i(h,a_X(h))\,d\mu$ and $S_i(X)=\int\psi_i(h,a_X(h))\,d\mu$ (integrable envelopes); the primary chart vector $c^F_i\coloneqq\varphi_i(\cdot,0)-\varphi_i(\cdot,+)$ is \emph{weakly} signed (zeros allowed); and either the secondary vector $c^S_i\coloneqq\psi_i(\cdot,0)-\psi_i(\cdot,+)$ is strictly signed, or $c^F_i$ is strictly signed and $c^S_i$ arbitrary (the class $\Ltwo$). Assume strict (G), (IB), (LA), (S), (C) on $\M_U$. Let $B_t$ be the lexicographic $t$-Bayes model and $E$ the set of $t$ at which the primary or, where the primary vanishes, the secondary has a non-null tie set (contained in countably many lines and points).

\begin{theorem}\label{thm:LX}
Let $W\subseteq\Di$ be an open disk with $W\cap E=\varnothing$ on which $t\mapsto B_t$ is injective. There is an open disk $W''\subseteq W$ with the following property. Let $M''$ be the set of problems on which $t\mapsto a_{B_t}(h)$ is not constant on $W''$. Then either (i) for every world $i$ there is $\alpha_i\in\R$ with $c^S_i=\alpha_ic^F_i$ a.e.\ on $M''$ and $c^F_i$ strictly signed there ($\alpha_i>0$ when $c^S_i$ is strictly signed), or (ii) $c^F=0$ a.e.\ on $M''$ and $B_t$ restricted to $M''$ is the secondary $t$-Bayes model (possible only when $c^S$ is strictly signed). In either case, for models that differ only on $M''$, every world's order is the additive order of one strictly signed field.
\end{theorem}
\begin{proof}[Proof sketch]
Two facts: by \src{R2} Thm.~7 the primary values at improvable undominated models form a countable set (in test form, a non-local minimum produces $C_p\preceq_iZ\prec_iC_p$); and a non-improvable model at a non-extremal level is a pointwise Lagrangian minimiser of $\psi_i+\lambda\varphi_i$, $\lambda\in\R$ (Lyapunov). For one world, $W$ is the union of the Lagrangian set and countably many closed level sets of $F_i(B_t)$. If a level set has interior, the concave value function has constant $\partial_iV$ there, and the distributional identity $\partial_i^2V=-\int\delta(t\cdot c^F)(c^F_i)^2\,d\mu$ forces $c^F_i=0$ on the problems whose primary chart line meets that disk. Otherwise, by Baire category, $B_t$ is Lagrangian for all $t\in W$; since a point does not separate a disk, injectivity puts all moving problems in one atom of the Lagrangian chain, where $c^S_i$ is proportional to $c^F_i$. Combining the worlds, the mixed cases contradict injectivity.
\end{proof}
\status{proved (\src{P2} Thm.~LX); confirmed, proof checked step by step (\src{P2P3\_referee} \S B5). The class statement is the referee's repair R-P2c: for a strictly signed primary (the class $\Ltwo$) only the second case of the proof occurs, alternative (ii) is vacuous, and $\alpha_i\in\R$ (when $\alpha_i\le0$ the secondary never decides). \emph{Scope:} two levels only. ``Layering is neutralised by closed sections'' holds for two-level layering; multi-level flexible worlds (Proposition~\ref{prop:flexible}(d)) are not covered, and the multi-level version is open (\src{P2} \S8, nearby result~1). Corollary LX.1 of \src{P2} (CHECKED, confirmed): the second level can do non-additive work only on thin (non-injective) regions, on a dense exceptional set, or on problems whose lexicographic advice is the same at every interior weighting. Round 13: the class is narrower than stated, since a weakly signed primary vanishes either a.e.\ on $M$ or nowhere on $M$ up to null sets (Corollary~\ref{cor:S21}(c); \src{S2\_referee} \S2.5, one reader).}

\begin{proposition}\label{prop:rigid}
\begin{enumerate}[leftmargin=*,label=(\alph*)]
\item Suppose the acceptance set of world $j$ is $\mathcal A^{\mathrm{flex}}\cup\mathcal A^{\mathrm{rig}}$, where $\mathcal A^{\mathrm{flex}}$ is endorsed by a strictly positive weight $\nu_j$ and every $\delta\in\mathcal A^{\mathrm{rig}}$ has $\|\delta\|_{L^1(\mu)}\ge\eta>0$. Then every comparison $C\preceq_jX$ with $\mu(C\sd X)<\eta$ is endorsed by $\nu_j$. The exotic acceptances of \src{N17} N17.7 and N17.8 are of this kind.
\item Let incoherent $X_k\to T_1$ in measure have selected weak dominators $C_{p_k}$ that use rigid acceptances, with $p_k\to\bar p$. Then ($\alpha$) $C_{\bar p}=T_1$ up to ties, and the selection is discontinuous at the target fibre; or ($\beta$) some coherent $C\ne T_1$ is tied with $T_1$ at every world; or ($\gamma$) some world $j$ is non-Archimedean at $T_1$: some $C$ has $C\preceq_jY_m$ for a sequence $Y_m\to T_1$ in measure, but not $C\preceq_jT_1$.
\end{enumerate}
\end{proposition}
\status{(a) proved (\src{P2} Lemmas~5.1, 5.2); confirmed (\src{P2P3\_referee} \S B6). (b) CHECKED (\src{P2} Prop.~5.3); confirmed at CHECKED level, with the wording repair that in ($\beta$) the model $C$ need not be far from $T_1$. \src{P2}'s (M1) design with rigid corner exchanges fails exactly here: rigid acceptances supply only dominators at distance at least $\eta$, and none of ($\alpha$)--($\gamma$) was realised or refuted.}

\begin{lemma}\label{lem:nested}
Let $M,M'$ be credal sets containing $K'$, with $e_i\notin M$. Then $E_i(D_M)\subseteq E_i(D_{M'})$ modulo null sets iff $M\subseteq\co(M'\cup\{e_i\})$ and $M'\subseteq\Delta\cap(\operatorname{cone}M-\R_+e_i)$; equivalently, $M$ and $M'$ have the same two tangent lines from $e_i$, and on every ray from $e_i$ both boundary points of $M'$ lie beyond those of $M$.
\end{lemma}
\status{proved (\src{P2} Lemma~6.1); confirmed (\src{P2P3\_referee} \S B7, with a Monte Carlo spot check). So nested pairs of coherent models are of codimension~2 in three worlds (NUMERICS: none of 2352 ordered pairs of translates of $T_1$ is nested at world~1), and the two-world mechanism of part~(c) below has nothing to act on generically. Theorem~\ref{thm:LX} recovers its conclusion in three worlds only for two-level lexicographic worlds, so ``crossing error sets are not by themselves an escape'' is established for that class only (referee's repair).}

\begin{proposition}\label{prop:P23}
Let $L$ be any ranking on $\M_U$, additive or not, with strict (G), (IB), (LA), (S), (C) that covers $T_1$, with selection $p\mapsto C_p=D_{K_p}$. Then $T_1$ is selected at an interior index $p^*$, and the family $q\mapsto K_q$ on $\Di$ violates at least one of: (i) Hausdorff continuity on $\Di$; (ii) vertex-respecting half-plane charts, i.e.\ for a.e.\ $g\in U$ the set $\{q\in\Di:P_{K_q}(g)>0\}$ agrees a.e.\ with $H(g)\cap\Di$ for an open half-plane $H(g)$ with $e_i\in H(g)$ if $g_i>0$ and $e_i\notin\cl H(g)$ if $g_i<0$.
\end{proposition}
\begin{proof}
The target may be taken selected (\src{Coverage after Neutral Abstention}, end of \S3). Its index is interior by the edge-Pareto condition (\src{P3} \S5.2): a selected model at an index $q\in\partial\Delta$ accepts the gambles of $U$ that are positive at every world of $\supp q$, and $D_{T_1}$ rejects some of them. $K_q\supseteq K'$ is automatic on $\M_U$. If (i) and (ii) held, the converse of Proposition~\ref{prop:reduction} would give a realisation of $T_1$ at $p^*$, which Theorem~\ref{thm:main} excludes.
\end{proof}
\status{proved in a referee report (\src{P2P3\_referee} \S B3), one reader; it replaces \src{P2} P2.3 (CHECKED there as ``linearity of the facet maps is not an additivity input''). By the converse of Proposition~\ref{prop:reduction}, straight charts together with continuity \emph{are} additivity at the level of the selection. \textbf{GAP:} \src{P2}'s localisation, that a cover must avoid a straight-chart family near $T_1$ on at least one facet arc, is not proved. The chain of \S\ref{sec:main} uses charts at weightings across all of $\Di$ (Lemma~\ref{lem:UVN} near the vertices, Lemmas~\ref{lem:PL2} and~\ref{lem:nopin} on sets that leave every ball around $p^*$, the box at the edges); a version localised near $p^*$ would be a new theorem.}

\emph{Holonomy (NUMERICS).} The constraint system that \S\ref{sec:main} produces for the facet maps (strict chart signs, crossing rows, (NS) at the corners, monotone rotation, the edge box, and linearity of each facet map) is inconsistent (Proposition~\ref{prop:T1}). With one kink between the two crossings of one facet, and the other two facets linear, the remaining finite system is consistent at $\varepsilon=0.05,0.3,0.6$, and with one kink per facet at $\varepsilon=0.9$ (\src{P2} \S2.2; reproduced independently with exact box margins in \src{P2P3\_referee} \S B2, relative margins $0.077$, $0.478$, $1.385$). This shows only that the contradiction of Proposition~\ref{prop:T1} needs linearity on all three facets. It is not a construction route: kinked data are excluded by coherence, since Lemma~\ref{lem:twoslope} forces the facet maps to be affine (at the break of the $\varepsilon=0.3$ solution the required concavity fails, \src{p23ref\_py/rotation\_check.out}), and the kinked direction field is not projective. A frozen corner sector does not help (a search outcome); a frozen middle sector is chart-consistent but is an atom chord through the target weighting, i.e.\ configuration R-a$'$, now closed by Theorem~\ref{thm:M0closed}.

\paragraph{(c) Two worlds with a nested selection, and the step to three.} Hypothesis \textbf{(Nest)}: $n=2$ on $\M_U$, strict (G), (IB), (LA), (S), and the world-1 error sets of the members of $\C_L$ form a chain modulo null sets (\src{N19} \S4, with that report's repair). Then $\beta(p)\coloneqq\mu(E_1C_p)/\mu(M)$ ($M$ the mixed problems of $U$) is continuous from $\Delta$ onto $[0,1]$, $\C_L$ is the selection, its error sets form a chain $c_b$ ($b\in[0,1]$, $\mu(c_b)=b\,\mu(M)$) with $d_b\coloneqq M\setminus c_b$, sections are closed in $p$ iff they are closed in $b$, and competitiveness reads (K): for every interior $b$ and every world-1 error set $A$ of a regular model, $c_b\le_1A$ or $d_b\le_2A^c$ (\src{P4B\_referee} Lemma~N0, proved; read $[s,t)$ as $c_t\setminus c_s$). Write $\varphi_1(A)\coloneqq\max\{b:c_b\le_1A\}$, which is attained (\src{P4} Lemma~B0). \emph{N17.B$'$} at world~1 (\src{N19}) asks for a weight $\nu$, positive on non-null sets, with $\nu(A)\ge\nu(c_b)$ whenever $c_b\le_1A$ (SP$^{\mathrm{reg}}_1$).

\begin{proposition}\label{prop:strictescape}
Under (Nest), for every $b\in[0,1)$ and every $A\supseteq c_b$ with $A\setminus c_b$ non-null, $\varphi_1(A)>b$; dually for world~2.
\end{proposition}
\begin{proof}
For $b'\in(\varphi_1(A),1)$, (K) gives $d_{b'}\le_2A^c$, and closed world-2 sections give $d_{\varphi_1(A)}\le_2A^c$. If $\varphi_1(A)=b$, then $A^c\subsetneq d_b$, against strict (G) at world~2. If $b=0$ and $\varphi_1(A)=0$, the same argument gives $M=d_0\le_2A^c$, again against strict (G).
\end{proof}
\status{proved (\src{P4} Lemma~B1 for $b\in(0,1)$; the case $b=0$ is Cor.~N1 of \src{P4B\_referee} \S2); confirmed (\src{P4B\_referee} \S2). It uses (K) at interior $b$, closed world-2 sections at regular models and strict (G) at world~2 only; no (C), (LA) or world-1 closedness.}

A charge $\pi\ge0$ vanishing on null sets is a \emph{leading charge} of world~1 if $\pi(B\setminus A)>\pi(A\setminus B)$ implies $A\le_1B$, and $A\le_1B$ implies $\pi(B\setminus A)\ge\pi(A\setminus B)$.

\begin{theorem}\label{thm:E2}
Under (Nest), every leading charge of world~1 is either $0$ or has a countably additive density $>0$ a.e.; in the second case it is strictly positive and endorses every relevant comparison, so SP$^{\mathrm{reg}}_1$ holds. In particular, if world~1 is lexicographic with countably many levels that are bounded charges $\ge0$ vanishing on null sets (countably additive, purely finitely additive or mixed), its first nonzero level is strictly positive and N17.B$'$ holds at world~1.
\end{theorem}
\begin{proof}[Proof sketch]
Put $f(b)\coloneqq\pi(c_b)$ and split $\pi=\rho\lambda+\sigma$ (Yosida--Hewitt). A left jump of $f$ produces a trimmed initial segment whose section is $[0,\beta)$, not closed; Proposition~\ref{prop:strictescape} makes $f$ right-continuous. If $\{\rho=0\}$ is non-null, $f$ vanishes on an initial segment $[0,s]$, and either $\pi=0$ or a section $[0,1)$ fails to be closed.
\end{proof}
\status{proved in a referee report (\src{P4B\_referee} \S8.3, Thm.~E2 and Cor.~E2$'$), one reader. It subsumes \src{P4} Thm.~B3 (lexicographic worlds with countably many countably additive levels: closed sections force the levels to be entered in order along the chain, and the first positive level endorses; proved without competitiveness, confirmed, with the remarks that entering in order is necessary but not sufficient, NUMERICS, and that the mechanism is that of \src{R2} Thm.~7) and the repaired Proposition~\ref{prop:B2}(b). With (K), the multi-level in-order worlds of \src{P4} Thm.~B3 cannot occur in a full two-world system. \emph{Round 13:} (Nest) is not needed; strict (G), (IB) and (S) suffice (Theorem~\ref{thm:S21} with $n=2$, Corollary~\ref{cor:S21}(a); \src{S2} Thm.~S2.1, confirmed in \src{S2\_referee} \S\S1, 2.1). The statement under (Nest) is the special case.}

\begin{theorem}\label{thm:E3}
Under (Nest), if $\le_1$ is complete, then N17.B$'$ holds at world~1.
\end{theorem}
\begin{proof}[Proof sketch]
No non-null set is $\le_1$ every member of a sequence of non-null sets shrinking to finitely many points (\src{P4B\_referee} Lemma~T, from Proposition~\ref{prop:strictescape} and closed sections through two ``drop'' lemmas). With completeness, compactness gives a finite partition of $M$ into pieces each $\le_1$ a given non-null set, so $\le_1$ is a fine qualitative probability in Savage's sense, and has an almost agreeing finitely additive probability~\cite{Sav1954,Fis1986}; it is strictly positive, vanishes on null sets and endorses every relevant comparison.
\end{proof}
\status{CHECKED (\src{P4B\_referee} \S8.4, Thm.~E3; the only imported step is Savage's theorem, whose proof is Boolean and applies on the measure algebra); one reader. Lemma~T is proved in that report (\S8.2), one reader. This closes a GAP in \src{P4} \S7, where complete non-lexicographic worlds had not been excluded.}

\begin{proposition}\label{prop:B2}
Identify $M$ with $[0,1]$ so that $c_b=[0,b)$, with Lebesgue measure $\lambda$. Let $x_0\in(0,1)$ and let $\nu^+$ be a purely finitely additive charge vanishing on null sets, with $\nu^+([0,1])=1$ and $\nu^+((x_0,x_0+h))=1$ for every $h>0$ (a right density at $x_0$).
\begin{enumerate}[leftmargin=*,label=(\alph*)]
\item The lexicographic world $\mathrm{lex}(\nu^+,\lambda)$ satisfies strict (G) and (IB), and all its sections are closed. If $\nu^+$ has infinitely many disjoint sets of positive mass (for instance a Banach limit of $\lambda(\cdot\cap(x_0,x_0+h))/h$), it fails SP$_1$. If $\nu^+$ takes finitely many values (for instance $\{0,1\}$), SP$_1$ holds.
\item For every such $\nu^+$, no two-world nested system with closed sections, (K) and strict (G) has this world~1: with $D=(x_0+h,1)$, $\varphi_1(c_{x_0}\cup D)=x_0$, against Proposition~\ref{prop:strictescape}.
\end{enumerate}
\end{proposition}
\status{proved (\src{P4} Prop.~B2); confirmed with repair (\src{P4B\_referee} \S3). Repairs incorporated: \src{P4}'s ``fails SP'' is false for $\{0,1\}$-valued $\nu^+$ (then $\lambda+\nu^+$ endorses every comparison) and holds under the added hypothesis; \src{P4}'s witness with $0<\nu^+(D)<1$ need not exist and is replaced by $D=(x_0+h,1)$, which works for every $\nu^+$. So closed sections alone do not give N17.B$'$; competitiveness through the other world is needed. \emph{Round 13:} (b) holds for every ranking with strict (G), (IB) and (S), nested or not, in the chain-free form of Corollary~\ref{cor:S21}(b) (\src{S2\_referee} \S2.2).}

\emph{The precise open class in two worlds.} N17.B$'$ holds at world~1 under (Nest) whenever world~1 is (CP-a) (Theorem~\ref{thm:fullsupport}; no (S) is needed), has a nonzero leading charge (Theorem~\ref{thm:E2}; in particular every lexicographic world with charge levels; by Corollary~\ref{cor:S21}(a) this needs no (Nest)), or is complete (Theorem~\ref{thm:E3}, CHECKED). Round 13: if either world has a nonzero leading charge or is monotonely continuous on $M$ (for instance additive with integrable density), the selection is automatically nested (Theorem~\ref{thm:S22}, Corollary~\ref{cor:MC}), so in those classes the hypothesis (Nest) costs nothing. A counterexample must have a world~1 that is incomplete, not (CP-a), and without a nonzero leading charge (a non-linear acceptance cone, like the exotic cones of \src{N17} N17.7); that has strict escape at every $b\in[0,1)$ and closed sections; and whose certificates of $\neg$SP$_1$ (sums $1_B+\sum_k\lambda_k(1_{A_k}-1_{c_{a_k}})\le\epsilon$ a.e., for every $\epsilon>0$; \src{P4} \S5.5, confirmed) are \emph{spread out}: no non-null set is $\le_1$ sets shrinking to finitely many points, as in $c_a\le_1A_j$ with disjoint $A_j$ spread over a set of positive measure. A proof would need a uniform version of the right-drop lemma or a closedness statement for a countable family of cuts. \emph{Status:} open (\src{P4B\_referee} \S8.5). \emph{Round 16:} the class is inhabited. Theorem~\ref{thm:D21} gives a counterexample to N17.B$'$ under (Nest) and (C), with world~2 additive: world~1 is, as Theorems~\ref{thm:fullsupport}, \ref{thm:E2} and~\ref{thm:E3} (CHECKED) require, not (CP-a), without a nonzero leading charge and incomplete; it has closed sections, and its certificates are ratio certificates with thin sources whose targets meet near the end of the chain. So N17.B$'$ fails in two worlds, also in the nested form N17.B$'$(2) of \src{P4} \S5.5 and \src{P4B\_referee} \S8.5. It holds with a $\sigma$-finite endorsing measure (Proposition~\ref{prop:D2sigma}), but with no endorsing measure whose density is locally integrable on $U$, so Conjecture~\ref{conj:sigmaB} is false (Theorem~\ref{thm:F13}); the form with local integrability on $M$ is Conjecture~\ref{conj:F1Cp}.

\emph{Three worlds.} The removal lemma (\src{P4} Lemma~B4; \src{R2} Thm.~7 with a lexicographic primary of countably many levels) holds for every $n$: if world $i$ is lexicographic with countably many countably additive levels, a selected $C_{p_0}$ errs at $i$ on a non-null set on which the levels below $j$ vanish, and indices $p_m\to p_0$ have strictly smaller valuation below level $j$, then (S) fails (proved, confirmed in \src{P4B\_referee} \S6). Necessary conditions for a three-world counterexample to N17.B$'$ (proved; confirmed): at the failing world, either an (M1) certificate (at least four comparisons between members of $\C_L$ and regular models, overlapping gains, crossing error sets, cyclic) or an (M2) failure; if that world is lexicographic with countably additive levels, every selected model that errs on a deep-level set is a weak local minimum of the shallow-level valuation along the selection; and every section is closed for every model, including the trimmed models of the removal lemma. For comparisons out of \emph{interior} selected models nothing in the two-world results transfers as it stands: strict escape needs a complementary world along a chain ($E_2=M\setminus E_1$), and crossing error sets destroy the cut structure that \src{P4} Thm.~B3(ii) uses (\src{P4B\_referee} \S6, which also corrects \src{P4}'s attribution: Thm.~B3 does not use strict escape). \emph{Correction (round 12).} The earlier statement here, that nothing in the two-world results transfers to $n=3$, was too strong. On the \emph{edges} of $\Delta$ the complementary world is supplied automatically (Lemma~\ref{lem:Q34}), and under edge nestedness (Nest$_e$) strict escape, Lemma~T, Theorems~\ref{thm:E2} and~\ref{thm:E3} and the results of \S\ref{sec:open-core}(c) hold on each edge; leading charges and complete worlds then globalise (Proposition~\ref{prop:Q35}, Theorem~\ref{thm:Q36}; \src{Q1} \S5.2, \src{Q3} \S4, \src{Q3\_referee} \S4). \emph{Status:} N17.B$'$ in three worlds is open (in two worlds it is false, Theorem~\ref{thm:D21}); it is proved for leading-charge worlds with no nestedness hypothesis (Theorem~\ref{thm:S21}, round 13), and, under edge nestedness, CHECKED for complete worlds.

\subsection{Worlds that do not add up: the two-world core}\label{sec:open-core}
Round 12 (\src{BRIEF11}) asked whether a world that does not add up (condition $\neg$SP$^{\mathrm{reg}}$ of Corollary~\ref{cor:M0closed}(a)) can exist inside a full system satisfying every axiom, for any target and any number of worlds: either construct one, with closed sections for \emph{every} evaluable model, or prove that none exists. Neither was achieved in round 12. No such world was constructed (\src{Q1}), and no exclusion was proved (\src{Q2}, \src{Q3}). Round 16 constructs one in two worlds (Theorem~\ref{thm:D21}, part~(e)). What was proved narrows the question to a two-world problem, recorded below together with the precise remaining question (part (g)). Sources: \src{Q1}, \src{Q2}, \src{Q3}; referee reports \src{Q3\_referee} (on \src{Q3}, \src{Q1} \S5 and \src{Q2} \S6) and \src{Q12\_referee} (on the two-world parts of \src{Q1} and \src{Q2}). Where a referee's repair conflicts with a note, the repaired version is given. \emph{Round 13} (\src{BRIEF12}) attacked items (R4) (nestedness from (S)) and (R1) (mirrored simple pairs) of part~(g), from both sides. Neither was settled in general. What was proved: leading charges are strictly positive with no nestedness hypothesis, in any number of worlds (Theorem~\ref{thm:S21}), so Theorem~\ref{thm:Q36}(a), (b) and the leading-charge clause of Corollary~\ref{cor:Q37} no longer need (Nest$_e$); in two worlds a nonzero leading charge or monotone continuity at either world forces a nested selection (Theorems~\ref{thm:S22}, \ref{thm:MC}), and without (C) nestedness can fail in a full ranking (Proposition~\ref{prop:R71}); the rate of item (R3) is a single joint property that holds next to an additive world (Theorem~\ref{thm:RATE}, Corollary~\ref{cor:SR}); and mirrored simple pairs exist when closed sections are dropped at both worlds (Proposition~\ref{prop:S31}). Sources: \src{S1}, \src{S2}, \src{S3}, \src{S4}; referee reports \src{S1\_referee} (on \src{S1}), \src{S2\_referee} (on \src{S2}) and \src{S34\_referee} (on \src{S3} and \src{S4}). Round-13 results are marked as such; the referees' repaired versions are used throughout. \emph{Rounds 14 and 15} (\src{BRIEF13}, \src{BRIEF14}) are part~(f): edge leading densities glue through the link class and propagate (Theorem~\ref{thm:B13}), so every world of a cover that does not add up lacks a leading density on some edge region (Corollary~\ref{cor:B14}); total worlds, glued credal worlds and density unanimities on edges; and face models at a target weighting on general decomposable domains. Part~(g), formerly part~(f), states what a cover still needs and the remaining questions. \emph{Round 16} (notes \src{D1}, \src{D2}, \src{E1}; referee reports \src{D1\_referee}, \src{D2\_referee}, \src{D2\_referee2}): a full two-world system in which a world does not add up (end of part~(e)); (LF) excludes the full-atom residual on general decomposable domains, and what the argument of Theorem~\ref{thm:M0closed} needs, with the exclusion of covers in which every world is an end-protected decoration with integrable, jointly end-heavy base densities (end of part~(f)). \emph{Audit of 15 October} (notes \src{F1}, \src{F2}; referee report \src{F12\_referee}): Conjecture~\ref{conj:sigmaB} is false, and nearby two-world systems (end of part~(e)); the exact hypotheses of the results of \src{E1}, and where ends can sit (end of part~(f)); and what a cover on $\M_U$ still needs (part~(g)). \emph{Round 17} (notes \src{H1}, \src{H2}, \src{J1}; referee report \src{H12\_referee} on \src{H1} and \src{H2}): failures that can be localized are invisible, admissibility of base endorsers is automatic, the conflict digraph, a full two-world system in which neither world adds up, and pattern switching (part~(f$'$)); the residual (part~(g)). \emph{Round 19} (notes \src{K1}, \src{K2}, \src{K3}; referee reports \src{J1\_referee}, \src{K1\_referee}, \src{K2\_referee}, \src{K3\_referee}): the referee's corrections and additions to \src{J1} (end of the \src{J1} block of part~(f$'$)); the classification of failures of adding up, finite certificates, the patch theorem and countably many ends (part~(f$''$)); and part~(g) rewritten.

\paragraph{Notation.} On $\M_U$ write $\sigma_i(h)\coloneqq\operatorname{sign}h_i$, $P_i\coloneqq U\cap\{h_i>0\}$ (the world-$i$ ideal model, so $E_i(X)=X\sd P_i$), $M$ for the mixed problems of $U$ and $N_+$ for its nonnegative gambles; $U=M\cup N_+$ modulo null sets, and $U$ contains no nonzero $h\le0$. By (G) and (IB), world $i$ is the same thing as a strictly monotone order $\le_i$ on world-$i$ error sets, $E_i(X)\le_iE_i(Y)$ iff $X\preceq_iY$, satisfying de Finetti cancellation $A\le_iB\iff A\setminus B\le_iB\setminus A$ (\src{N9} N9.4). For a regular model $X$ and a coherent $B$, both error sets lie in $M$, so relevant comparisons are comparisons of subsets of $M$. A \emph{shadow} is a vector $\nu=(\rho_i\mu)_i$ of measures with $\rho_i\in L^1(\mu)$ and $\rho_i>0$ a.e.; it \emph{endorses} a comparison $B\preceq_iX$ if $\nu_i(E_iB)\le\nu_i(E_iX)$, and its $t$-Bayes model is $B_t\coloneqq\{g\in U:\sum_it_i\rho_i(g)\sigma_i(g)>0\}$. For $i\ne j$ let $e_{ij}\coloneqq[e_i,e_j]$ be an edge of $\Delta$ and $R_{ij}\coloneqq\{h\in U:h_ih_j<0\}$. Standing hypotheses \textbf{(H$_U$)}: $n=3$, the $U$-domain for $T_1$, and a ranking with strict (G), (IB), (LA), (S) (selection $p\mapsto C_p$), (C) and $T_1\in\C_L$; this is (H$_{\mathrm{M0}}$) of \S\ref{sec:M0} without (M0).

\paragraph{(a) What an exclusion of $T_1$ must use.} The first result says that no argument confined to selection indices in a compact subset of $\Di$ can exclude $T_1$. Fix $L>2$ and put, for $g\in U$, $w_i(g)\coloneqq L$ if either $g\in D_{T_1}$ and $g_i>0$, or $g\notin D_{T_1}$ and $g_i<0$, and $w_i(g)\coloneqq1$ otherwise; let $\nu^L_i\coloneqq w_i\mu$, $\psi^L_i(X)\coloneqq\nu^L_i(E_iX)$, and $Q_L\coloneqq\{t\in\Delta:\min_it_i>1/(L+1)\}$.

\begin{proposition}\label{prop:Q31}
\begin{enumerate}[leftmargin=*,label=(\alph*)]
\item For $t\in Q_L$ the tie set of $\nu^L$ is null and $B_t=D_{T_1}$ a.e. If $t_k=\min_it_i<1/(L+1)$, then $B_t$ misses $S_k\coloneqq D_{T_1}\cap\{g_k>0>g_i,g_j\}$ ($\{i,j,k\}=\{1,2,3\}$), so $\mu(B_t\sd D_{T_1})\ge\eta_0\coloneqq\min_k\mu(S_k)>0$, where $\eta_0$ does not depend on $L$.
\item For every compact $Q\subseteq\Di$ there is $L$ with $Q\subseteq Q_L$, and then the additive, strictly guiding ranking $\psi^L$ on $\M_U$ with $C_q\coloneqq T_1$ for $q\in Q$ satisfies global strict (G), (IB), (LA) and $T_1\in\C_L$, and, at every index $q\in Q$, competitiveness of $C_q$ on $\supp q$, relative closedness in $Q$ of the sections $\{q\in Q:C_q\preceq_iX\}$, and (C) for models tied with $T_1$.
\end{enumerate}
Consequently a derivation of a contradiction from these premises, with selection indices restricted to a fixed compact subset of $\Di$, does not exist: an exclusion of covers of $T_1$ must use selection indices contained in no compact subset of $\Di$. Global consequences of (S), such as Theorem~\ref{thm:kkm}, count as using all indices.
\end{proposition}
\status{proved (\src{Q3} Lemma~2.1, Prop.~Q3.1); confirmed (\src{Q3\_referee} \S1, Lemma~2.1 re-derived by hand; NUMERICS \src{q3ref\_py/shadow\_family.py} (1)), with the reading repaired: the conclusion holds in the precise form stated, for index-local arguments. \src{Q3}'s extension to arguments ``local in model space'' is heuristic and is not claimed; in model space the shadow argument does force behaviour arbitrarily close to $T_1$ (Theorem~\ref{thm:R3}). So the localisation GAP of Proposition~\ref{prop:P23} cannot be closed by index-local means.}

\begin{proposition}\label{prop:Q32}
Assume (H$_U$), and let $\nu$ be a shadow. At least one of the following holds.
\begin{enumerate}[leftmargin=*,label=(\alph*)]
\item Some comparison $T_1\preceq_iX$ with $X$ regular has $\nu_i(E_iT_1)>\nu_i(E_iX)$.
\item $T_1$ is $\nu$-Bayes at an interior weight; the set $I_\nu$ of $t\in\Di$ at which the tie set is null and $B_t$ is incoherent is not Lebesgue-null; and at every $t\in I_\nu$ some selected $C$ and some world $j$ have $C\preceq_jB_t$ and $\nu_j(E_jC)>\nu_j(E_jB_t)$, a relevant comparison that $\nu$ does not endorse.
\end{enumerate}
For $\nu=\nu^L$, case (b) holds whatever the cover is, $I_{\nu^L}\cap Q_L=\varnothing$, and every $B_t$ with $t\in I_{\nu^L}$ has $\mu(B_t\sd D_{T_1})\ge\eta_0$.
\end{proposition}
\begin{proof}[Proof sketch]
If (a) fails, \src{N17} Prop.~2.1($\alpha$) (which uses only regular $X$) makes $T_1$ $\nu$-Bayes at an interior weight. Theorem~\ref{thm:M0closed}(a), applied to the strictly signed field $\rho\sigma$, says that the strict Bayes sets are not coherent at a.e.\ interior weight; ties are null off a null set of weights (Tonelli). At $t\in I_\nu$, \src{N17} Lemma~5.1 (Theorem~\ref{thm:kkm} with null ties) gives the dominator. For $\nu^L$, $T_1=B_u$ and Proposition~\ref{prop:Q31}(a) gives the rest; at $\min t=1/(L+1)$ the tie set contains a set $S_k$, so those weights are not in $I_{\nu^L}$.
\end{proof}
\status{proved (\src{Q3} Props.~Q3.2, Q3.3); confirmed with repair (\src{Q3\_referee} \S\S2.1--2.2): ``exactly one of (a), (b)'' is corrected to ``at least one'' ((b) concerns $\nu$, (a) the cover; both can hold), and the densities must be integrable. It closes \src{N19}'s GAP on ``continuum''. NUMERICS \src{q3py/propL\_shadow.py} ($\varepsilon=0.3$, $s=0.9$, $L=21$). \emph{Readings.} (i) \src{N17}'s Conjecture~5.2 is false in its index reading: for $\nu^L$ the forced incoherent $B_t$ cannot be taken with $t$ tending to the weight at which $T_1$ is Bayes. Its model-space half ($B_t\to T_1$) is in fact forced (Theorem~\ref{thm:R3}). (ii) The $t$ here is a shadow weight, not a selection index; the dominator $C$ is selected at an index about which nothing is said. Rescaling the worlds of $\nu^L$ moves the collar $\Delta\setminus Q_L$ projectively, so the invariant content is that $I_\nu$ avoids a neighbourhood of the weight at which $T_1$ is $\nu$-Bayes. Propositions~\ref{prop:Q31} (selection indices) and~\ref{prop:Q32} (shadow weights) do not combine into one ``boundary'' statement. (iii) \src{Q3}'s claims that the premise of Proposition~\ref{prop:rigid}(b) ``is not forced by any shadow argument'' and that ``load-bearing reversals near $T_1$ are not forced'' are \textbf{wrong} (\src{Q3\_referee} \S2.3): one shadow is not the shadow argument, and a cover must meet the disjunction for every shadow at once. The correct statement is Theorem~\ref{thm:R3}.}

For a cover, let $\mathcal N_i$ be the set of world-$i$ shadows $\nu_i$ that endorse every comparison $T_1\preceq_iX$ with $X$ regular, and $\mathcal N\coloneqq\prod_i\mathcal N_i$; so $\nu\in\mathcal N$ iff Proposition~\ref{prop:Q32}(a) fails for $\nu$.

\begin{theorem}\label{thm:R3}
Assume (H$_U$). Exactly one of the following holds.
\begin{enumerate}[leftmargin=*,label=(\roman*)]
\item Some world $i$ has $\mathcal N_i=\varnothing$: for every shadow, some comparison $T_1\preceq_iX$ out of $T_1$ itself, with $X$ regular, is not endorsed.
\item $\mathcal N\ne\varnothing$, and for every $\eta>0$ there are a selected $C\ne T_1$, an incoherent model $X$ with $\mu(X\sd D_{T_1})<\eta$, and a world $j$ such that $C\preceq_kX$ at every world $k$, and $C\preceq_jX$ is a relevant comparison that some $\nu'\in\mathcal N$ does not endorse.
\end{enumerate}
\end{theorem}
\begin{proof}[Proof sketch]
\emph{Bathtub closure.} If $\nu_i\in\mathcal N_i$, $\lambda>0$, and $f\ge0$ is measurable with $f\le\lambda$ on $E_iT_1$ and $f\ge\lambda$ off it, and $f\nu_i$ is a shadow (integrable density, $>0$ a.e.), then $f\nu_i\in\mathcal N_i$, and the same holds for comparisons $T_1\preceq_iX$ with $X$ arbitrary: with $A=E_iT_1$ and $E=E_iX$, $\int_{E\setminus A}f\,d\nu_i-\int_{A\setminus E}f\,d\nu_i\ge\lambda(\nu_i(E\setminus A)-\nu_i(A\setminus E))\ge0$. \emph{Case (ii).} Take $\nu\in\mathcal N$ and normalise by bathtub closure so that errors of $T_1$ weigh $\le1$ and correct points $\ge1$. Fix a mixed ray $w\in D_{T_1}$ and a closed circular cone $C_\eta$ around it with $\cl C_\eta\subseteq D_{T_1}\setminus0$, $\mu(C_\eta)<\eta$, every point of $C_\eta$ a midpoint of two points of $D_{T_1}\setminus C_\eta$. Pick $j$ with $w_j<0$ and an interior $t$ near $e_j$ (off a null set of weights) at which at least half of $C_\eta$ has $t\cdot c<0$; call that part $C'$. Multiply the density of world $i$ by $M\coloneqq2/\min_it_i+1$ wherever world $i$ is correct for $T_1$, except on $C_\eta$. By bathtub closure the result $\nu'$ is in $\mathcal N$, and its $t$-Bayes model is $T_1\setminus C'$ with null ties, which is incoherent (a coherent cone containing $D_{T_1}\setminus C_\eta$ contains $D_{T_1}$) and within $\eta$ of $T_1$. \src{N17} Lemma~5.1 gives the selected dominator; $C\ne T_1$ because $\nu'$ endorses the comparisons out of $T_1$.
\end{proof}
\status{proved in a referee report (\src{Q3\_referee} \S2.4: Lemma~R1, Prop.~R2, Thm.~R3), one reader. Correction (15 October): \src{Q3\_referee} states bathtub closure with the further hypothesis that $f$ is bounded on bounded sets. The proof does not use it, and the construction does not always satisfy it: the normalisation step uses $f=\max(1,1/\rho_i)$ off $E_iT_1$, which is unbounded on bounded sets wherever $\rho_i$ has no positive lower bound there (shadows only require $\rho_i>0$ a.e.). The hypothesis is therefore omitted here; what the construction needs, and has, is that $f\nu_i$ is a shadow: its densities are bounded by $M\max(\rho_i,1)$, integrable because $\mu$ is finite; NUMERICS \src{q3ref\_py/shadow\_family.py} (2)--(4), including an explicit thin-cone version for $\nu^L$ with $\mu(B'_t\sd T_1)/\mu(U)=1.4\cdot10^{-3},9\cdot10^{-5},10^{-5},\dots$. It replaces item~3 of \src{Q3} \S6. So the shadow argument always forces non-additive behaviour at $T_1$: in $T_1$'s own out-comparisons, or against incoherent models converging to $T_1$. What it does not settle is the \emph{position of the dominators}: near $T_1$, or far. In case (ii) there are incoherent $X_k\to T_1$ with selected weak dominators and unendorsed comparisons, as Proposition~\ref{prop:rigid}(b) needs, with two differences: the endorsing shadow varies with $k$, and the dominators may approach $T_1$. If they stay at distance $\ge\eta$ along a subsequence, Proposition~\ref{prop:Q38} applies. \emph{Round 16:} Theorem~\ref{thm:E12} (proved in a note, audited in \src{F2}, audit refereed once) gives the conclusion of case~(ii) without the case split, with admissible shadows in place of shadows in $\mathcal N$.}

Say that \textbf{(LC$_{T_1}$)} (lower closedness at $T_1$) holds at world $j$ if $C_k\preceq_jX_k$, $C_k\to C$ and $X_k\to T_1$ in measure imply $C\preceq_jT_1$.

\begin{proposition}\label{prop:Q38}
Assume (H$_U$). Let $X_k\to T_1$ in measure, let the coherent models $C_k$ weakly dominate $X_k$ at every world (they need not be selected), and let $C_k\to C^\infty$ in measure (true along a subsequence, by compactness). If (LC$_{T_1}$) holds at every world along $(C_k,X_k)$, then $C^\infty\sim_jT_1$ at every world, $C^\infty\in\C_L$, and $\mu(C^\infty\sd T_1)=\lim_k\mu(C_k\sd X_k)$. Hence dominators at distance $\ge\eta$ force a tie partner of $T_1$ in $\C_L$ at distance $\ge\eta$; and in Proposition~\ref{prop:rigid}(b), alternative ($\alpha$) without ($\beta$) forces the failure of (LC$_{T_1}$) at some world.
\end{proposition}
\status{proved (\src{Q3} Prop.~Q3.8); confirmed with a wording repair (\src{Q3\_referee} \S6): alternative ($\gamma$) of Proposition~\ref{prop:rigid}(b) is the failure of (LC$_{T_1}$) with $C_k\equiv C$ fixed, so $\neg$(LC$_{T_1}$) is implied by ($\gamma$) and is weaker, a generalised (moving-dominator) form of ($\gamma$), not a sharpened one. Under (LC$_{T_1}$) the partner in ($\beta$) is at distance $\ge\eta$, which answers the wording objection of \src{P2P3\_referee} \S B6. Neither closed sections nor the indices are used.}

Call a measurable $A\subseteq U$ a \emph{tie toggle} of $T_1$ if $T_1\sd A$ is tied with $T_1$ at every world, and let $\mathcal Z$ be the class of tie toggles.

\begin{lemma}\label{lem:Q39}
Assume (H$_U$). (a) $\mathcal Z$ contains $\varnothing$ and is closed under finite disjoint unions and under relative differences $B\setminus A$ ($A\subseteq B$, both in $\mathcal Z$). (b) For $A\in\mathcal Z$, $T_1\sd A\in\C_L$ by (C); in particular it is coherent. (c) If every $\preceq_i$ is closed in measure, $\mathcal Z$ is closed in measure, hence also under countable increasing unions. If every world is closed in measure (for instance, of the type of the fine (NC)-violating world of \src{N17} N17.7, which is closed in measure), far dominators of near-$T_1$ models force a far tie partner.
\end{lemma}
\status{proved (\src{Q3} Lemma~Q3.9, via sum closure of acceptance sets, \src{N17} N17.1(a)); confirmed with repair (\src{Q3\_referee} \S6): (a) gives only finite unions and relative differences, so ``Dynkin-type'' needs (c); and the application to far dominators needs (LC$_{T_1}$) at \emph{every} world, not only at the N17.7 world. Alternatives ($\beta$) and ($\gamma$) of Proposition~\ref{prop:rigid}(b) remain open: for general orders $\mathcal Z$ may be as small as $\{\varnothing,A\}$ (a rigid tie), and whether (S) forces lower closedness is open (\src{Q3} \S5; the Villegas remark there is CHECKED, one import).}

\paragraph{(b) Edges: three worlds reduce to two.} On an edge of $\Delta$ only two worlds compete, and the selected models form a two-world system on $R_{ij}$.

\begin{lemma}\label{lem:Q34}
Let $n\ge2$ and let a ranking on $\M_U$ satisfy strict (G). If $C\in\M_U$ is competitive on $\{i,j\}$ (for instance $C=C_p$ with $p\in e_{ij}$), then a.e.\ $C$ is correct at world $i$ or at world $j$. Hence $E_i(C)\cup E_j(C)\subseteq R_{ij}$, $C$ agrees with $P_i$ (equivalently $P_j$) off $R_{ij}$, and $E_j(C)=R_{ij}\setminus E_i(C)$. For $A\subseteq R_{ij}$ the \emph{edge test model} $X^{ij}_A$, equal to $P_i$ off $R_{ij}$ with $E_i(X^{ij}_A)=A$, is regular and has $E_j(X^{ij}_A)=R_{ij}\setminus A$.
\end{lemma}
\begin{proof}
If $A\coloneqq\{h:\text{$C$ errs at $i$ and at $j$}\}$ were non-null, $X\coloneqq C\sd A$ would have $E_i(X)=E_i(C)\setminus A$ and $E_j(X)=E_j(C)\setminus A$, so $X\prec_iC$ and $X\prec_jC$ by strict (G), against competitiveness. Off $R_{ij}$ the signs of $h_i$ and $h_j$ agree, so the two worlds agree on correctness; on $R_{ij}$ exactly one is correct.
\end{proof}
\status{proved (\src{Q3} Lemma~Q3.4; \src{Q1} Lemma~Q1.7); confirmed (\src{Q3\_referee} \S3, re-derived; NUMERICS \src{shadow\_family.py} (5)). It uses only strict (G) and competitiveness on $\supp p\subseteq\{i,j\}$, not (C) or coherence. Vertex indices are included.}

Say that \textbf{(Nest$_e$)} holds at world $i$ on $e_{ij}$, for the given selection, if $\{E_i(C_p):p\in e_{ij}\}$ is a chain modulo null sets; by Lemma~\ref{lem:Q34} this holds at world $i$ iff it holds at world $j$. It is the edge analogue of (Nest) of part~(c) of \S\ref{sec:open-ordinal}.

\begin{proposition}\label{prop:Q35}
Let $n\ge2$, let a ranking on $\M_U$ satisfy strict (G), (IB), (LA) and (S), and assume (Nest$_e$) at world $i$ on $e_{ij}$. Put $\beta(p)\coloneqq\mu(E_iC_p)/\mu(R_{ij})$ for $p\in e_{ij}$. Then:
\begin{enumerate}[leftmargin=*,label=(\arabic*)]
\item $\beta$ is continuous from $e_{ij}$ onto $[0,1]$, $\beta(e_i)=0$, $\beta(e_j)=1$, and models with equal $\beta$ agree a.e.; write $c_b\coloneqq E_i(C_p)$ and $d_b\coloneqq R_{ij}\setminus c_b=E_j(C_p)$ for $\beta(p)=b$;
\item for $A\subseteq R_{ij}$ the sections $\{b:c_b\le_iA\}$ and $\{b:d_b\le_jR_{ij}\setminus A\}$ are closed;
\item \textbf{(K)}: for every $b\in(0,1)$ and every $A\subseteq R_{ij}$, $c_b\le_iA$ or $d_b\le_jR_{ij}\setminus A$;
\item consequently strict escape (Proposition~\ref{prop:strictescape} with \src{P4B\_referee} Cor.~N1), the drop lemmas DL and DR, Lemma~T, Theorem~\ref{thm:E2}, and, from part~(c) below, Lemma~T$_2$, Theorem~\ref{thm:LF}, Corollary~\ref{cor:IS} and Proposition~\ref{prop:MP} hold verbatim for $\le_i,\le_j$ restricted to measurable subsets of $R_{ij}$, with $M$ replaced by $R_{ij}$ and worlds $(1,2)$ by $(i,j)$;
\item if $\le_i$ is complete, its restriction to subsets of $R_{ij}$ is fine in Savage's sense (the first step of Theorem~\ref{thm:E3}).
\end{enumerate}
\end{proposition}
\status{proved (\src{Q3} Prop.~Q3.5; \src{Q1} Prop.~Q1.8, CHECKED there); confirmed (\src{Q3\_referee} \S4.1, which checks every hypothesis of \src{P4B\_referee} \S\S1--2, 8 on the edge system and upgrades \src{Q1} Prop.~Q1.8 to PROVED, at the standard of \src{P4B\_referee} \S8, one reader, plus this check). Item~5 inherits the CHECKED status of Theorem~\ref{thm:E3}. The extension of item~4 to part~(c) (LF, IS, MP and \src{Q1} Thm.~Q1.4; the referee's later results of part~(c) are not claimed on edges) is \src{Q3\_referee} \S\S4.3, 8, which corrects \src{Q2}'s CONJECTURE Q2.E (its restriction set was vacuous and it omitted nestedness) to this form and proves it from Lemma~\ref{lem:Q34}, item~4 and \src{Q2} \S\S2--3; \src{Q1} Cor.~Q1.10 is the case of Proposition~\ref{prop:Q15}: if world $j$ has SP$^{\mathrm{reg}}_j$, world $i$ has no spread certificate inside one $R_{ij}$ (CONFIRMED, conditional on \src{Q1} Thm.~Q1.4, which is refereed in \src{Q12\_referee}). Lemma~T is proved in \src{P4B\_referee} \S8.2: under (Nest), no non-null set is $\le_1$ every member of a sequence of non-null sets contained in slabs shrinking to finitely many points of the chain.}

A charge $\pi\ge0$ on $U$, bounded and vanishing on null sets, is a \emph{leading charge of world $i$ on $M$} if, for all measurable $A,B\subseteq M$, $\pi(B\setminus A)>\pi(A\setminus B)$ implies $A\le_iB$, and $A\le_iB$ implies $\pi(B\setminus A)\ge\pi(A\setminus B)$.

\emph{Round 13: leading charges without nestedness.} The next theorem removes every nestedness hypothesis from the leading-charge results, in any number of worlds (\src{S2} \S4, confirmed in \src{S2\_referee} \S\S1--2). It concerns only the edges at $e_i$. Recall that $C_{e_i}=P_i$ and $C_{e_j}=P_j$ (competitiveness at the vertices and strict (G)), so that $E_i(C_{e_i})=\varnothing$ and $E_i(C_{e_j})=R_{ij}$, and that for $p\in e_{ij}$, $E_j(C_p)=R_{ij}\setminus E_i(C_p)$ (Lemma~\ref{lem:Q34}).

\begin{theorem}\label{thm:S21}
Let $n\ge2$, let $K'\subseteq\Di$ be a credal set (for $n=2$, an interval), and let a ranking on $\M_U$ satisfy strict (G), (IB) and (S). Let $\pi$ be a leading charge of world $i$ on $M$. Then either $\pi(M)=0$, or $\pi(D)>0$ for every non-null measurable $D\subseteq M$. In the second case the countably additive part of $\pi$ has a density $\rho>0$ a.e.\ on $M$, $\pi$ endorses every comparison of world $i$ between subsets of $M$, and world $i$ adds up (SP$^{\mathrm{reg}}_i$). No nestedness hypothesis is made.
\end{theorem}
\begin{proof}[Proof sketch]
Fix $j\ne i$ with $\pi(R_{ij})>0$ and put $m(p)\coloneqq\pi(E_iC_p)$ for $p\in e_{ij}$, so $m(e_i)=0$ and $m(e_j)=\pi(R_{ij})$. Every non-null $E\subseteq M$ contains, for every $\epsilon>0$, a non-null $D$ with $\pi(D)<\epsilon$: split $E$ into $k$ disjoint non-null pieces, one of which has $\pi\le\pi(E)/k$. \emph{$m$ is lower semicontinuous.} If $E_iC_{p^*}$ is non-null, take such a $D\subseteq E_iC_{p^*}$ and $A\coloneqq E_iC_{p^*}\setminus D$. Every edge index $q$ with $m(q)<\pi(A)$ has $C_q\preceq_iX^{ij}_A$ by the first property of a leading charge. If such $q$ accumulated at $p^*$, the closed section $S_i(X^{ij}_A)$ would contain $p^*$, against $X^{ij}_A\prec_iC_{p^*}$ (strict (G)). \emph{$m$ is upper semicontinuous.} If $m(q_k)>m(q^*)+\eta$ with $q_k\to q^*$, take a small non-null $D\subseteq R_{ij}\setminus E_iC_{q^*}$ and $A\coloneqq E_iC_{q^*}\cup D$. The second property gives $C_{q_k}\not\preceq_iX^{ij}_A$, so competitiveness on $\supp q_k\ni j$ gives $C_{q_k}\preceq_jX^{ij}_A$, and closedness gives $C_{q^*}\preceq_jX^{ij}_A$, against $X^{ij}_A\prec_jC_{q^*}$. \emph{A generic level.} So $m$ is continuous and attains every value in $[0,\pi(R_{ij})]$. Values of weak local extrema of $m$ are countable, so some value $v\in(0,\pi(R_{ij}))$ is not one of them; at any $p$ with $m(p)=v$, every neighbourhood contains indices with $m<v$ and indices with $m>v$. The two arguments above, with closedness, give $C_p\preceq_iX^{ij}_A$ and $C_p\preceq_jX^{ij}_A$ for every $A\subseteq R_{ij}$ with $\pi(A)=v$. \emph{No $\pi$-null sets.} If $N'\subseteq R_{ij}$ is non-null with $\pi(N')=0$, then either $A\coloneqq E_iC_p\setminus N'$ (if $N'\cap E_iC_p$ is non-null) or $A\coloneqq E_iC_p\cup N'$ (otherwise) has level $v$, and $X^{ij}_A$ is strictly better than $C_p$ at world $i$, respectively $j$; a contradiction. So $\pi$ is strictly positive on $R_{ij}$. \emph{Linking} ($n\ge3$): $R_{ij}\cap R_{ik}$ contains a neighbourhood of $e_i-\eta(e_j+e_k)$, which lies in $U$ because $K'\subseteq\Di$, so positivity spreads to every $R_{ik}$, and $M=\bigcup_{k\ne i}R_{ik}$. Finally, if $\{\rho=0\}$ were non-null, removing a carrier of the purely finitely additive part from it would leave a non-null $\pi$-null set. Endorsement is the second property, and $\pi+\mu|_{N_+}$ is a strictly positive weight endorsing every relevant comparison (whose error sets lie in $M$).
\end{proof}
\status{proved (\src{S2} Thm.~S2.1); confirmed, every step re-derived (\src{S2\_referee} \S1); so two readers. The proof above uses the referee's two simplifications: no Yosida--Hewitt carrier lemma is needed to find small sets, and continuity of $m$ serves only to produce one generic index. Inputs: strict (G) at worlds $i$ and $j$; (IB), through the normal form; competitiveness and closed sections at edge indices, for the edge test models only; for $n\ge3$, the linking step. Not used: (C), (LA) beyond (S), nestedness, completeness, countable additivity of $\pi$, interior indices. This is consistent with the locality barrier (Proposition~\ref{prop:Q31}), which concerns indices in compact subsets of $\Di$; the edges lie outside every such set. \emph{Why no nestedness is needed:} earlier proofs parametrised the edge selection by its chain; here it is parametrised by its $\pi$-level, and at a generic level the selected model is optimal at both edge worlds among all edge test models of that level, so it can neither contain nor avoid a non-null $\pi$-null set. Proposition~\ref{prop:B2}, \src{N26} Thm.~N26.1 and \src{P4} Thm.~B3 are consistent with it (\src{S2\_referee} \S1).}

\begin{corollary}\label{cor:S21}
\begin{enumerate}[leftmargin=*,label=(\alph*)]
\item ($n=2$.) Theorem~\ref{thm:E2}, including its lexicographic clause, holds without (Nest), (K) or the chain: strict (G), (IB) and (S) suffice.
\item ($n=2$.) Let world~1 be $\mathrm{lex}(\pi,\lambda')$, where $\pi$ is a bounded charge vanishing on null sets, nonzero on $M$ and zero on some non-null subset of $M$, and $\lambda'$ is any strictly positive tie-breaker. Then no ranking with strict (G), (IB) and (S) has this world~1.
\item Let world $i$ be a lexicographic world as in Theorem~\ref{thm:LX} (on $\M_U$, every axiom), with weakly signed primary chart vector $c^F_i$. Then either $c^F_i=0$ a.e.\ on $M$, or $c^F_i\ne0$ a.e.\ on $M$.
\end{enumerate}
\end{corollary}
\status{(a) proved (\src{S2} Cor.~S2.1$'$(a)); confirmed (\src{S2\_referee} \S2.1). (b) is Proposition~\ref{prop:B2}(b) restated without chain coordinates, as \src{S2\_referee} \S2.2 requires (CONFIRMED WITH REPAIR): Proposition~\ref{prop:B2} identifies $M$ with $[0,1]$ through the chain, which presupposes nesting; $\pi$ is a leading charge of $\mathrm{lex}(\pi,\lambda')$, so (b) is Theorem~\ref{thm:S21}, and Proposition~\ref{prop:B2}(b) is the case obtained by transporting $\nu^+$ and $\lambda$ to $M$ by a measure isomorphism. (c) proved in a referee report by inspection (\src{S2\_referee} \S2.5), one reader: the primary $|c^F_i|\mu$ restricted to $M$ is a leading charge of world $i$, so the ``zeros allowed'' of Theorem~\ref{thm:LX} occur only in this all-or-nothing form on $M$. Nothing beyond the hypotheses of Theorem~\ref{thm:LX} is used.}

\begin{theorem}\label{thm:Q36}
Let $n\ge2$, let $K'\subseteq\Di$ be any credal set, and let a ranking on $\M_U$ satisfy strict (G), (IB), (LA) and (S). Fix a world $i$.
\begin{enumerate}[leftmargin=*,label=(\alph*)]
\item Every leading charge $\pi$ of world $i$ on $M$ either vanishes on $M$ or has a countably additive density $\rho>0$ a.e.\ on $M$; in the second case it endorses every comparison of world $i$ between subsets of $M$, and SP$^{\mathrm{reg}}_i$ holds.
\item[(a$'$)] ($n=3$.) If world $i$ has positive integrable leading densities on $R_{ij}$ and on $R_{ik}$ (in the sense of part~(f)), then it has a leading charge on $M$ that is nonzero there, so (a) applies.
\item If world $i$ is lexicographic with countably many bounded charge levels $\ge0$ vanishing on null sets (countably additive, purely finitely additive or mixed), SP$^{\mathrm{reg}}_i$ holds.
\item If (Nest$_e$) holds at world $i$ on every edge $e_{ij}$, $j\ne i$, and $\le_i$ is complete, SP$^{\mathrm{reg}}_i$ holds.
\item If (Nest$_e$) holds at world $i$ on $e_{ij}$, then strict escape and Lemma~T hold for $\le_i$ on $R_{ij}$.
\end{enumerate}
\end{theorem}
\begin{proof}[Proof sketch]
$M=\bigcup_{j\ne i}R_{ij}$, and for $j\ne k$ (both $\ne i$) $R_{ij}\cap R_{ik}$ contains a neighbourhood of $e_i-\eta(e_j+e_k)$, which lies in $U$ for small $\eta$ because $K'\subseteq\Di$; so it is non-null (the \emph{linking} step). (a) is Theorem~\ref{thm:S21}. Under (Nest$_e$) on every edge at $e_i$ there is also the earlier proof: on each $R_{ij}$, Proposition~\ref{prop:Q35}(4) with Theorem~\ref{thm:E2} gives $\pi|_{R_{ij}}=0$ or a density $>0$ a.e.\ on $R_{ij}$; if $\pi|_M\ne0$, some $\pi|_{R_{ij}}\ne0$, then $\rho>0$ a.e.\ on $R_{ij}\cap R_{ik}$, so $\pi|_{R_{ik}}\ne0$, and so on for every $k$. (a$'$) is Theorem~\ref{thm:B13}: the glued density $\rho_i$ gives the leading charge $\rho_i\mu|_M$. (b) The first level that is nonzero on $M$ is a leading charge on $M$. (c) For non-null $Z\subseteq M$, Proposition~\ref{prop:Q35}(5) partitions $R_{ij}$, and through the linking sets every $R_{ik}$, into finitely many pieces $\le_iZ$; so $\le_i$ on $M$ is a complete fine qualitative probability, and Savage's theorem~\cite{Sav1954,Fis1986} gives a strictly positive, almost agreeing finitely additive probability. (d) is Proposition~\ref{prop:Q35}(4).
\end{proof}
\status{(a), (b) proved without (Nest$_e$) (\src{S2} Thm.~S2.1 and Cor.~S2.1$'$(b); confirmed, \src{S2\_referee} \S\S1, 2.3; two readers). The earlier version of this theorem assumed (Nest$_e$) at world $i$ on every edge at $e_i$ for all four items; under that hypothesis (a), (b), (d) had been proved before (\src{Q3} Thm.~Q3.6; \src{Q1} Thm.~Q1.9, CHECKED there; confirmed, \src{Q3\_referee} \S4.2), and that version is the special case. (c) CHECKED: the one import is Savage's theorem, as in Theorem~\ref{thm:E3}. (d) proved, as before. (c) and (d) still use (Nest$_e$), through Proposition~\ref{prop:Q35}; by Remark~\ref{rem:straighten} it may be replaced there by any nested edge family with (K) and closed sections. For $n=2$ this is Theorems~\ref{thm:E2} and~\ref{thm:E3} (one edge, $R_{12}=M$); the new content for $n\ge3$ is complementarity on edges and linking. For a general acceptance cone, positivity on edge comparisons does not globalise: an endorsing charge restricted to the edges need not endorse comparisons out of interior selected models. This is why the theorem stops at leading charges and complete worlds. \emph{Round 14:} item (a$'$) is \src{B1} Thm.~B1.3 (Theorem~\ref{thm:B13}), added here as \src{B1\_referee} (iv) recommends; proved, confirmed by that referee, one reader. Its proof uses of $K'$ only that the link class $R_{ij}\cap R_{ik}$ is non-null, which is the linking step above. So a \emph{leading} density on the two edge regions at $e_i$ does propagate to all comparisons of world $i$, while an \emph{endorsing} edge weight need not (Proposition~\ref{prop:B1cex}).}

\begin{corollary}\label{cor:Q37}
Assume (H$_U$). Suppose that, for some selection satisfying (S), every world $i$ is (CP-a), or has a leading charge nonzero on $M$ (in particular, is lexicographic with charge levels), or is complete and satisfies (Nest$_e$) on both edges at $e_i$. Then $T_1$ is not covered.
\end{corollary}
\begin{proof}
(CP-a) worlds add up (Corollary~\ref{cor:CPaU}(b)); worlds with a leading charge nonzero on $M$ have SP$^{\mathrm{reg}}$ by Theorem~\ref{thm:Q36}(a), (b) (that is, Theorem~\ref{thm:S21}), for any selection with (S); complete worlds with (Nest$_e$) at $e_i$ have SP$^{\mathrm{reg}}$ by Theorem~\ref{thm:Q36}(c). Worlds with SP$^{\mathrm{reg}}$ add up (\src{N19} Lemma~R1). So the ranking is in (M0), which Theorem~\ref{thm:M0closed}(b) excludes.
\end{proof}
\status{proved at the status of Theorem~\ref{thm:M0closed}(b) (imports \src{N13} Prop.~B and \src{N19} Prop.~R3(a), each read once), together with Theorem~\ref{thm:fullsupport}, Theorem~\ref{thm:S21} (two readers), \src{P4B\_referee} \S8 (one reader) and \src{Q3\_referee}; CHECKED when a complete world is involved (Theorem~\ref{thm:E3} via Savage). Sources: \src{Q3} Cor.~Q3.7 (confirmed, \src{Q3\_referee} \S5), strengthened by \src{S2} Cor.~S2.1$'$(c) (confirmed, \src{S2\_referee} \S2.4). The earlier version, which required (Nest$_e$) at $e_i$ of leading-charge worlds as well, is the special case. The all-(CP-a) case is Corollary~\ref{cor:CPaU}. \emph{Contrapositive} (replacing the earlier (E-a)/(E-b)): every world $i$ of a cover of $T_1$ that does not add up is not (CP-a) and has no leading charge nonzero on $M$ (unconditionally); and either (E-a$'$) some edge family at $e_i$ is not nested at world $i$, or (E-b$'$) all are nested and world $i$ is incomplete, with strict escape and Lemma~T on every $R_{ij}$ (\src{S2\_referee} \S2.4). \emph{Rounds 14--15} (part~(f)). By Theorem~\ref{thm:Q36}(a$'$) the leading-charge clause covers every world with positive integrable leading densities on both edge regions at $e_i$ (\src{B1\_referee} (iv)). By Lemma~\ref{lem:B31}, ``complete (total) and AR on $M$'' is a special case of ``leading charge nonzero on $M$'', so the complete clause adds nothing beyond the leading-charge clause except as one route to AR, through Theorem~\ref{thm:Q36}(c) (nestedness, then fineness, then Savage's theorem); Corollary~\ref{cor:B32} replaces it by ``total and AR on $M$'' (\src{B3\_referee} \S2). Theorem~\ref{thm:C16} adds glued credal worlds to the list (round 15; confirmed by one referee).}

The next lemma is a set identity, valid for any $n$; it drives the round-13 nestedness results (\src{S2} \S3).

\begin{lemma}\label{lem:PL}
Let $C,F\in\M_U$, let $D\subseteq C\sd F$, and put $X\coloneqq C\sd D$, $X'\coloneqq F\sd D$. Then at every world $k$, $(E_kC\setminus E_kX,\,E_kX\setminus E_kC)=(E_kX'\setminus E_kF,\,E_kF\setminus E_kX')$, so for every ranking with (G) and (IB): $C\preceq_kX$ iff $X'\preceq_kF$, and $F\preceq_kX'$ iff $X\preceq_kC$. If $F\in\C_L$ and $C\preceq_kX$ at every world $k$, then $X'\sim_kF$ and $C\sim_kX$ at every world, and $X'\in\C_L$ by (C); if moreover $C\in\C_L$, then also $X\in\C_L$.
\end{lemma}
\status{proved (\src{S2} Lemma~PL, the ``parallelogram lemma'': $C\to X$ and $X'\to F$ are the same exchange); confirmed with repair (\src{S2\_referee} \S3.1): \src{S2} concluded that both $X$ and $X'$ are coherent from $F\in\C_L$ alone; only $X'$ is, and $X\in\C_L$ needs $C\in\C_L$, which holds in every use. NUMERICS: \src{s2py/pl\_identity.py} (the set identity) and \src{s2ref\_py/pl\_orders.py} (both equivalences, with additive, lexicographic and incomplete unanimity worlds, $n\in\{2,3,4\}$, $20{,}000$ trials, no violation; a control with $D\not\subseteq C\sd F$ fails $13{,}607$ times).}

The hypothesis (Nest$_e$) is genuine: it is not derived from the axioms. The edge inputs of Proposition~\ref{prop:Q35} do not force it (Proposition~\ref{prop:plateau} below), and in two worlds the hypotheses of Proposition~\ref{prop:Q35} without (C) do not force it either (Proposition~\ref{prop:R71}). It holds for edge selections that are Bayes for additive edge systems with null ties. A non-null tie set of the edge worlds cannot be split in crossing ways by an edge selection of a system satisfying the axioms if the two edge worlds are additive with integral penalties and the third world is complete, in particular if all worlds are additive (Proposition~\ref{prop:RS2A}), nor, in two worlds, if either world has a nonzero leading charge (Theorem~\ref{thm:S22}). A crossing tie-split would require the third world to be incomplete, with $C$ and $C\sd(P\cup Q)$ incomparable there for every bounded balanced exchange; whether this can happen is open. For $n=3$ the exclusion is proved only for integral (density) penalties at the edge worlds: with purely finitely additive charge penalties, Lyapunov's theorem and the Bayes-support step below are not available. (Correction, round 13. The earlier text here said that (Nest$_e$) ``can fail at exceptional edge weights where a non-null tie set is split in crossing ways''. That is true of the family of Bayes face models, not of the edge selection of any system satisfying the axioms in the cases just listed; \src{S2} \S2, corrected in scope by \src{S2\_referee} \S4.4.)

\begin{proposition}\label{prop:S20}
Let $\psi$ be an additive ranking on $\M_U$ with strictly positive integrable densities $\rho_k$. Take $n=2$, or one edge $e_{ij}$ with weights $(1-t)e_i+te_j$, and put $r\coloneqq\rho_i/(\rho_i+\rho_j)$ on $R_{ij}$ ($R_{12}=M$ if $n=2$).
\begin{enumerate}[leftmargin=*,label=(\alph*)]
\item A $t$-Bayes model errs at world $i$ where $r<t$, is correct at world $i$ where $r>t$, and is free on the tie set $Z_t\coloneqq\{r=t\}$. For $t<t'$, $\{r\le t\}\subseteq\{r<t'\}$. So Bayes error sets at distinct weights are nested, for every tie-break (strictly nested iff $\{t<r<t'\}$ is non-null), and two crossing Bayes models are Bayes at one common weight and differ only on its tie set.
\item Let $n=2$. If $\psi$, with some selection, satisfies every axiom (in particular (S) and (C)), then $Z_t$ is null at every interior $t$.
\end{enumerate}
\end{proposition}
\begin{proof}[Proof sketch]
(a) Comparing the weighted costs $(1-t)\rho_i$ and $t\rho_j$ of the two possible errors. (b) Every $t$-Bayes model at an interior $t$ is undominated, so by Theorem~\ref{thm:kkm} and (C) it lies in $\C_L$ and is coherent. If $Z_t$ were non-null it would be a cone modulo null sets, as a difference of coherent cones, and a bounded split of it is $t$-Bayes and not a cone.
\end{proof}
\status{(a) proved (\src{S2} S2.0(a)); confirmed with a wording repair (\src{S2\_referee} \S4.1): the result is for one edge or $n=2$, and the sets are nested, not ``strictly nested in any $n$'' (for $n\ge3$, at interior weights, nested pairs have codimension~2, Lemma~\ref{lem:nested}). NUMERICS \src{s2py/additive\_tie.py} (exact rationals, 39 weights, every tie-break). (b) proved (\src{S2} S2.0(b)); confirmed (\src{S2\_referee} \S4.1). NUMERICS: a $5\times5$ rectangle of coherent $t_0$-Bayes models with 100 crossing pairs, and an explicit incoherent split with the same profile at both worlds.}

\begin{proposition}\label{prop:RS2A}
Let $n=3$ and let a ranking on $\M_U$ (any $K'\subseteq\Di$) satisfy strict (G), (IB), (S) and (C). Let worlds $i$ and $j$ be additive with integral penalties, $\psi_l(X)=\int_{E_lX}\rho_l\,d\mu$ with $\rho_l>0$ a.e.\ and integrable ($l=i,j$), let $k$ be the third world, and let $C,F\in\C_L$ be competitive on $\{i,j\}$ (for instance $C=C_p$, $F=C_{p'}$ with $p,p'\in e_{ij}$).
\begin{enumerate}[leftmargin=*,label=(\alph*)]
\item If world $k$ is complete, $C$ and $F$ are nested at world $i$. Hence (Nest$_e$) holds on $e_{ij}$; in particular it holds when all three worlds are additive with integral penalties.
\item In general, suppose $C$ and $F$ cross, and put $A\coloneqq E_iF\setminus E_iC$, $B\coloneqq E_iC\setminus E_iF$. Then $A\cup B\subseteq Z_s$ for one $s\in(0,1)$, and for all non-null bounded $P\subseteq A$, $Q\subseteq B$ with $\int_P\rho_i\,d\mu=\int_Q\rho_i\,d\mu$, the models $C$ and $C\sd(P\cup Q)$ are incomparable at world $k$, and so are $F$ and $F\sd(P\cup Q)$.
\item If moreover world $k$ has a leading charge $\pi_k$ nonzero on $M$, a crossing also needs the countably additive density of $\pi_k$ to be proportional to $\rho_i$ a.e.\ on $A\cup B$, and $C$ to be correct at world $k$ a.e.\ on one of $A$, $B$ and wrong at world $k$ a.e.\ on the other.
\end{enumerate}
\end{proposition}
\begin{proof}[Proof sketch]
Competitiveness on $\{i,j\}$ makes the profile $(\psi_i,\psi_j)$ of $C$ weakly Pareto-minimal in a compact convex attainable set (Lyapunov), so $C$ is $s$-Bayes for worlds $i,j$ at some $s$; by Proposition~\ref{prop:S20}(a) a crossing forces one common $s\in(0,1)$ and $A\cup B\subseteq Z_s$, where $\rho_j=\kappa\rho_i$. Coherent models and $P_i$ are cones modulo null sets, so $A$ and $B$ are cones and meet the unit ball in non-null sets; a bounded non-null $P\cup Q$ is not a cone, so $X\coloneqq C\sd(P\cup Q)$ and $X'\coloneqq F\sd(P\cup Q)$ are incoherent (\emph{bounded splices}). The balance makes $X\sim_iC$ and $X\sim_jC$. If $C\preceq_kX$, Lemma~\ref{lem:PL} gives $X'\preceq_lF$ at all three worlds, so $X'\in\C_L$ by undominatedness of $F$ and (C), which is impossible. If $X\preceq_kC$, then $X$ dominates $C$, or is tied with it everywhere and lies in $\C_L$ by (C); both are impossible. So $C$ and $X$ are incomparable at $k$, and the $F$-side follows by Lemma~\ref{lem:PL}; a complete world has no incomparable pairs, which gives (a). (c) With $P,Q$ off one carrier of the purely finitely additive part of $\pi_k$, the two properties of a leading charge turn $\pi_k(E_kX)\ne\pi_k(E_kC)$ into a strict comparison, which (b) excludes; varying $P$ and $Q$ gives the proportionality, with a nonzero constant because the density is $>0$ a.e.\ (Theorem~\ref{thm:S21}).
\end{proof}
\status{proved in a referee report (\src{S2\_referee} \S4.3, Prop.~R-S2.A), one reader. It repairs and generalises \src{S2} S2.0(c) (all three worlds additive, every edge selection nested), which is the special case ``world $k$ additive''. \src{S2}'s proof of S2.0(c) asserted an incoherent Lyapunov half-splice without proof (GAP: Lyapunov may produce a cone); the referee repairs this with bounded splices (\src{S2\_referee} \S4.2), and the proof above needs neither closed sections, generic levels, the bathtub principle nor $K'$. In two worlds the same argument shows that, if both worlds are additive, $\C_L$ has no crossing pair. NUMERICS \src{s2ref\_py/fan\_search.py}: crossing pairs undominated at both models with no all-world tie, 0 when world $k$ is additive or lexicographic (complete), 50 in 3000 random trials when world $k$ is an incomplete unanimity order. Related (\src{S2} S2.4 item~5, strengthened in \src{S2\_referee} \S5, proved): a crossing between a generic edge-selected $C$ and an edge model $F\in\C_L$ requires $C$ to be wrong at world $k$ a.e.\ on $A$ or a.e.\ on $B$, under strict (G) alone. The ``fans'' of \src{S2} S2.4 items 2--4 are derived only for additive edge worlds, and item~4 is a heuristic (GAP).}

On a natural family of coherent edge models (Nest$_e$) is exactly a comonotonicity condition. For $n=3$, the edge $e_{12}$ and $0\le x\le y\le1$ put $M_{x,y}\coloneqq\co(K'\cup\{(a,1-a,0):a\in[x,y]\})$.

\begin{proposition}\label{prop:comonotone}
(i) $D_{M_{x,y}}$ agrees with $P_1$ off $R_{12}$. (ii) On $R_{12}$ the world-1 error set of $D_{M_{x,y}}$ is
\begin{multline*}
\{h_1>0>h_2:\ \min(P_{K'}(h),\,xh_1+(1-x)h_2)\le0\}\\
\cup\ \{h_1<0<h_2:\ \min(P_{K'}(h),\,yh_1+(1-y)h_2)>0\},
\end{multline*}
whose first part grows strictly (on non-null sets) as $x$ decreases and whose second part grows strictly as $y$ decreases. Hence $D_{M_{x,y}}$ and $D_{M_{x',y'}}$ are nested at world~1 iff $(x,y)$ and $(x',y')$ are comonotone, i.e.\ $(x-x')(y-y')\ge0$.
\end{proposition}
\status{proved in a referee report (\src{Q3\_referee} \S4.4), one reader. NUMERICS: \src{shadow\_family.py}, part (5) (a comonotone pair nested, a crossing pair not, with about $10^4$ sample points in each difference). So on this family (Nest$_e$) is the edge form of CONJECTURE~\ref{conj:Q2N}: ``(S) forces comonotone endpoints''. Edge-selected models touch the edge (\src{N17} N17.4(c), applicable by Lemma~\ref{lem:Q34}). Whether every coherent edge-selected model is of the form $D_{M_{x,y}}$, and whether (S) forces comonotone endpoints, is open. Round 13: the edge inputs alone do not force it (Proposition~\ref{prop:plateau}), nor do the hypotheses of Proposition~\ref{prop:Q35} without (C) (Proposition~\ref{prop:R71}); it is proved when the edge worlds are additive with integral penalties and the third world is complete (Proposition~\ref{prop:RS2A}); and on this family a crossing pair can be straightened pairwise under extra hypotheses (Proposition~\ref{prop:S112}).}

\begin{remark}\label{rem:straighten}
Proposition~\ref{prop:Q35}(1)--(5), Theorem~\ref{thm:Q36}(c), (d) and the complete clause of Corollary~\ref{cor:Q37} use the edge selection only through a family $(c_b)_{b\in[0,1]}$ of subsets of $R_{ij}$ with $c_0=\varnothing$ and $c_1=R_{ij}$ that is a chain, satisfies (K) at every interior $b$, has closed $b$-sections at worlds $i$ and $j$ for subsets of $R_{ij}$, and through strict (G) and cancellation. They never use that the family is selected or coherent, or that $\C_L$ is the path. So any \emph{nested edge family} with these properties can replace the selection. Producing such a family from a given selection is called \emph{straightening}; on edges, straightening rather than (Nest$_e$) itself is the right target. After Theorem~\ref{thm:S21}, which removes (Nest$_e$) for leading charges, what straightening would still buy is Theorem~\ref{thm:Q36}(c), for complete worlds. For Theorem~\ref{thm:LF}, Corollary~\ref{cor:IS} and Proposition~\ref{prop:MP} on the edge to be statements about \emph{relevant} comparisons, each member of the family must be tied at worlds $i$ and $j$ with a member of $\C_L$.
\end{remark}
\status{proved by inspection of the recorded proofs (\src{S1} Remark~S1.8, against \src{Q3} \S4.2 and \src{Q3\_referee} \S4.1); confirmed with repair (\src{S1\_referee} \S6). The referee adds Proposition~\ref{prop:Q35}(5) and Theorem~\ref{thm:Q36}(c) to the list, restricts the relevant-comparison reading of LF, IS and MP as stated, and notes that the ``no missing value'' and continuity steps of Lemma~N0 need only a chain indexed by a compact connected set, with $c=\varnothing$ and $c=R_{ij}$ at two index points. \emph{Round 14:} after Corollary~\ref{cor:B32}, Theorem~\ref{thm:Q36}(c) is the case of that corollary in which almost representability comes from fineness, so straightening is one route to AR for total worlds, not the only one; and Proposition~\ref{prop:B34} shows that on an edge a crossing path can coexist with a world that is not fine (\src{B3\_referee} \S\S2, 4).}

\begin{proposition}\label{prop:plateau}
Let $R\cong[0,2)$ with Lebesgue measure $\mu$, and consider the two-world edge system on subsets of $R$ with $c\le_iA$ iff $\mu(c)\le\mu(A)$ and $d\le_jB$ iff $\mu(d)\le\mu(B)$ (world $j$ compares $R\setminus c$). Let $p\mapsto c_p$, $p\in[0,1]$, be continuous in measure with $\mu(c_p)$ nondecreasing, $c_0=\varnothing$ and $c_1=R$; for instance, grow to $[0,\tfrac12)$, slide to $[1,1\tfrac12)$ at constant mass, then grow to $R$. Then strict (G), (IB), (K) at every index and closed sections at both worlds hold, but for the example path $c_{1/3}=[0,\tfrac12)$ and $c_{2/3}=[1,1\tfrac12)$ are disjoint.
\end{proposition}
\status{proved (\src{S1} Example~S1.9; NUMERICS \src{s1py/plateau\_edge.py}, 300 random test sets on a 4000-cell grid); confirmed (\src{S1\_referee} \S7). It is an edge system, not a ranking: its sets are neither coherent nor undominated, and in three worlds they would be dominated through world~3. A nested family exists here, so it is not a counterexample to straightening. It shows that (Nest$_e$) cannot be derived from the inputs of Proposition~\ref{prop:Q35}; any derivation must use (C), or undominatedness through the third world. Proposition~\ref{prop:R71} below is the sharper statement for a full ranking.}

\paragraph{(c) Two worlds with a nested selection: the exclusion side.} Throughout this part assume (Nest) as in part~(c) of \S\ref{sec:open-ordinal} ($n=2$, strict (G), (IB), (LA), (S), chain $c_b$, $d_b=M\setminus c_b$, slabs $[s,t)\coloneqq c_t\setminus c_s$, all modulo null sets), and (C). Besides Lemma~T we use the two drop lemmas of \src{P4B\_referee} \S8.1 (proved, one reader): \emph{(DL)} if $x\in(0,1]$, $Z\subseteq c_{x-\eta}$ is non-null, $F\subseteq d_x$, and $Z\le_1[b,x)\cup F$ for every $b\in(x-\eta,x)$, then $Z\le_1F$; \emph{(DR)} if $x<z$ and $D\subseteq d_z$ is non-null, there is $y>x$ with $L\cup R\le_1L\cup D$ for every $R\subseteq[x,y)$ and every $L$ disjoint from $[x,y)\cup D$. A \emph{cut pair} at $b\in(0,1)$ is a pair $(P,Q)$ with $P\subseteq c_b$ and $Q\subseteq d_b$. Let $\mathcal V_1$ be the set of probability charges on $M$, vanishing on null sets, with $\nu(P)\le\nu(Q)$ for every cut pair with $P\le_1Q$; it is weak-* compact, convex, nonempty and norm-closed, and SP$^{\mathrm{reg}}_1$ means that some $\nu\in\mathcal V_1$ is strictly positive on non-null sets. $\mathcal V_2$ is defined in the mirror way. Write $\bar\mu\coloneqq\mu/\mu(M)$.

\begin{lemma}\label{lem:cutform}
Let $b\in(0,1)$ and let $X$ be regular with $E_1X\cap M=(c_b\setminus P)\cup Q$ for a cut pair $(P,Q)$ at $b$. Then $C_b\preceq_1X$ iff $P\le_1Q$, and $C_b\preceq_2X$ iff $Q\le_2P$. Competitiveness is (K): every cut pair has $P\le_1Q$ or $Q\le_2P$. Undominatedness with (C): no cut pair other than $(\varnothing,\varnothing)$ has both $Q\le_1P$ and $P\le_2Q$. The relevant comparisons at world~1 are exactly the cut pairs with $P\le_1Q$, and at world~2 those with $Q\le_2P$.
\end{lemma}
\status{proved (\src{Q1} Lemma~Q1.0, a reformulation of \src{P4B\_referee} Lemma~N0 and the cut-pair form of \src{P4} \S5.5, \src{P4B\_referee} \S5); confirmed (\src{Q12\_referee} \S1.1).}

The symmetry $b\mapsto1-b$ with the worlds exchanged preserves (Nest), closed sections at both worlds, (K) and strict (G) at both worlds, so every result of \S\ref{sec:open-ordinal}(c) holds for world~2 with $c$ and $d$ exchanged. In particular \emph{Lemma~T$_2$}: for non-null $W$ and $x\in[0,1]$ there is $\eta>0$ such that no non-null $Q$ inside the slab $(x-\eta,x+\eta)$ has $W\le_2Q$ (\src{Q2} \S1, proved by symmetry from \src{P4B\_referee} Lemma~T; confirmed, \src{Q12\_referee} \S0). Lemma~T$_1$ is the same statement for world~1. For non-null $Z$ let $b_0(Z)\coloneqq\max(1-\eta,\tfrac12)$, with $\eta$ given by Lemma~T$_1$ at $x=1$ for $W=Z$.

The next theorem is \src{Q2}'s \emph{left-fineness}.

\begin{theorem}\label{thm:LF}
Let $a\in(0,1)$ and let $Z\subseteq d_a$ be non-null. There are finitely many slabs $J_1,\dots,J_m$ covering $c_a$ such that every measurable $Q\subseteq c_a\cap J_l$ satisfies $Q\le_1Z$, for every $l$; these are relevant comparisons (cut pairs at $a$), so every $\nu\in\mathcal V_1$ has $\nu(c_a)\le m\,\nu(Z)$. Uniformly: there is $\delta=\delta(Z,a)>0$ such that $Q\le_1Z$ for every $Q\subseteq c_a$ lying in a slab of width $<\delta$.
\end{theorem}
\begin{proof}
For $x\in[0,a]$, Lemma~T$_2$ with $W=Z$ gives $\eta_x$. The intervals $(x-\eta_x,x+\eta_x)$ cover $[0,a]$; take a finite subcover and a Lebesgue number $\delta$. For non-null $Q\subseteq c_a\cap J_l$, $(Q,Z)$ is a cut pair at $a$, so (K) gives $Q\le_1Z$ or $Z\le_2Q$, and Lemma~T$_2$ excludes the second; null $Q$ satisfy $Q\le_1Z$ by monotonicity.
\end{proof}
\status{proved (\src{Q2} Thm.~LF); confirmed (\src{Q12\_referee} \S2.1). It is \src{P4B\_referee}'s Lemma~DR made uniform (the ``uniform DR'' that \src{P4B\_referee} \S8.5 asked for), and the left half of Savage fineness without completeness. It bounds charges on the left by charges on the right, never the reverse, which is why it does not finish the proof.}

\begin{corollary}\label{cor:IS}
For $\nu\in\mathcal V_1$ put $a_0(\nu)\coloneqq\sup\{a\in[0,1]:\nu(c_a)=0\}$ and $a^*\coloneqq\inf_{\nu\in\mathcal V_1}a_0(\nu)$.
\begin{enumerate}[leftmargin=*,label=(\roman*)]
\item Every $\nu\in\mathcal V_1$ is positive on every non-null subset of $d_{a_0(\nu)}$.
\item Some $\bar\nu\in\mathcal V_1$ has $a_0(\bar\nu)=a^*$, so it is positive on every non-null subset of $d_{a^*}$.
\item SP$^{\mathrm{reg}}_1$ holds iff $a^*=0$.
\item If $a^*>0$, every $c_a$ with $a<a^*$ is $\mathcal V_1$-null (null for every member), and every $\mathcal V_1$-null set is contained in $c_{a^*}$ modulo null sets. Hence $\neg$SP$^{\mathrm{reg}}_1$ iff some initial segment $c_a$, $a>0$, is $\mathcal V_1$-null; equivalently, for every $\epsilon>0$ there are cut pairs $P_k\le_1Q_k$ and $\lambda_k\ge0$ with $1_{c_a}+\sum_k\lambda_k(1_{Q_k}-1_{P_k})\le\epsilon$ a.e.
\end{enumerate}
\end{corollary}
\status{proved (\src{Q2} Cor.~IS); confirmed, every step re-derived (\src{Q12\_referee} \S2.2). Precision: (iv) shows $\nu(c_a)=0$ for $a<a^*$, not $\nu(c_{a^*})=0$; a member may carry a finitely additive ``left mass'' at $a^*$, so the $\mathcal V_1$-null ideal may be strictly larger than the union of the $c_a$, $a<a^*$. The certificate form is the separation step of \src{P4} \S5.5. Mirror (world~2, final segments $d_b$): proved by symmetry. So the null sets of a failing world are initial segments: a certificate can only export the mass of an initial segment, whatever its targets.}

A \emph{simple certificate} at world~1 is a non-null $Z$ with $Z\le_1A_j$ for a sequence $(A_j)$ of bounded overlap, $\sum_j1_{A_j}\le r$ (for instance pairwise disjoint); then every charge endorsing these comparisons vanishes on $Z$. The comparisons $Z\le_1A_j$ need not be relevant.

\begin{proposition}\label{prop:MP}
Let $Z$ be non-null and $b_0=b_0(Z)$.
\begin{enumerate}[leftmargin=*,label=(\alph*)]
\item For every $A$ with $Z\le_1A$ and every $t\in[b_0,1)$, $d_t\le_2A$; equivalently $d_t\setminus A\le_2A\cap c_t$. This is a relevant world-2 comparison ($C_t\preceq_2X$ with $E_2X=A$).
\item If $Z\le_1A_j$ for a sequence of bounded overlap, every $\nu\in\mathcal V_2$ vanishes on $d_{b_0}$; so world~2 fails SP$^{\mathrm{reg}}_2$, with the mirror certificate $d_{b_0}\le_2A_j$ through the same sets $A_j$.
\end{enumerate}
\end{proposition}
\begin{proof}
(a) Consider the cut pair $(A\cap c_t,\,d_t\setminus A)$ at $t$. If $A\cap c_t\le_1d_t\setminus A$, cancellation gives $A\le_1d_t$, so $Z\le_1d_t\subseteq d_{1-\eta}$, against Lemma~T$_1$. So (K) gives $d_t\setminus A\le_2A\cap c_t$, i.e.\ $d_t\le_2A$. (b) For $\nu\in\mathcal V_2$, $J\nu(d_{b_0})\le\sum_{j\le J}\nu(A_j)\le r$ for every $J$.
\end{proof}
\status{proved (\src{Q2} Prop.~MP, which assumes bounded overlap throughout); confirmed (\src{Q12\_referee} \S2.3), with the observation that (a) needs no overlap condition: overlap is used only in the final step (b). The sharp condition for (b) is Proposition~\ref{prop:N1}(b). This is also \src{Q1} Thm.~Q1.4 after the referee's repair (\src{Q12\_referee} \S1.4): its summability hypothesis and the passage to a subset $Z'$ are superfluous, it holds with $h_1=\eta(Z)$, and \src{Q1} \S0's statement of the world-2 certificate as ``$d_{1-h}\le_2P_j$ for disjoint $P_j$'' is corrected to $d_{1-h}\le_2A_j$ (equivalently $d_{1-h}\setminus A_j\le_2A_j\cap c_{1-h}$); $d_{1-h}\le_2A_j\cap c_{1-h}$ is stronger and is not proved. After the repair \src{Q1} Thm.~Q1.4 and \src{Q2} Prop.~MP have the same conclusion. Summability alone does not suffice for (b) (a decreasing family with a common free filter carries a charge with $\nu(A_j)=1$).}

\begin{corollary}\label{cor:MPloop}
Let $Z$ be non-null. There is $s_0>0$ such that $c_s\le_1A$ for every $s\in(0,s_0]$ and every $A$ with $Z\le_1A$. So a simple world-1 certificate with any exported set upgrades to one exporting an initial segment, with the same targets; and a simple certificate at either world forces one at the other through the same sets.
\end{corollary}
\status{proved (\src{Q2} Cor.~MP-loop: Proposition~\ref{prop:MP}(a), then its mirror with $W=d_{b_0}$ and Lemma~T$_2$ at $x=0$); confirmed (\src{Q12\_referee} \S2.3). ``Same targets'' means the same sets, which are $E_1$ of one model and $E_2$ of another.}

\begin{corollary}\label{cor:MPprime}
If world~1 has a simple certificate, then neither world is complete, (CP-a), or has a nonzero leading charge.
\end{corollary}
\status{proved (\src{Q2} Cor.~MP$'$); confirmed (\src{Q12\_referee} \S2.3). The ``complete'' clause uses Theorem~\ref{thm:E3} and is CHECKED. The (CP-a) clause is valid because every member of a (CP-a) credal set endorses \emph{all} comparisons, including the non-relevant $Z\le_1A_j$. So the open class of \S\ref{sec:open-ordinal}(c) must be inhabited at both worlds at once.}

\begin{proposition}\label{prop:NS}
Let $Z\le_1A_j$ be a simple certificate at world~1 and $b_0=b_0(Z)$. Let $\pi\ge0$ be a charge vanishing on null sets such that world~1 accepts every cut pair at cuts $t\in[b_0,1)$ that $\pi$ strictly favours: $\pi(P)<\pi(Q)\Rightarrow P\le_1Q$. Then $\pi(d_{b_0})=0$.
\end{proposition}
\status{proved (\src{Q2} Prop.~NS: by the proof of Proposition~\ref{prop:MP}(a), $A_j\cap c_t\not\le_1d_t\setminus A_j$, so $\pi(d_t\setminus A_j)\le\pi(A_j)$ and $\pi(d_t)\le2\pi(A_j)\to0$ along a subsequence); confirmed with repair (\src{Q12\_referee} \S2.3): $b_0$ is $b_0(Z)$ of Proposition~\ref{prop:MP}, which depends on $Z$, and $t=b_0$ is in the range. No world-2 hypothesis is used. Consequence: the architecture of \src{N17} N17.8 (an additive band, i.e.\ unanimity over densities bounded below, plus exotic patches) cannot carry a simple certificate, since its band contains the strict $\mu$-comparisons on cut pairs and $\mu(d_{b_0})>0$. Applied with $\pi=\nu_2\in\mathcal V_2$ it re-derives Proposition~\ref{prop:MP}(b), and applied to the construction of part~(e) it locates exactly the failing closedness of Proposition~\ref{prop:Q13} (\src{Q12\_referee} \S2.3).}

\begin{proposition}\label{prop:Q15}
If world~2 satisfies SP$^{\mathrm{reg}}_2$, world~1 has no simple certificate; this applies when world~2 is (CP-a), has a nonzero leading charge (in particular, is lexicographic with charge levels), is complete, or is additive. If moreover world~2 has an endorsing charge $\nu_2$ that is strictly positive on non-null sets and uniformly absolutely continuous ($\nu_2(A)\to0$ as $\mu(A)\to0$; for instance world~2 additive with integrable density), then (a) no non-null $Z$ has $Z\le_1A_j$ for a sequence with $\mu(A_j)\to0$, and (b) there is no \emph{star certificate}: no non-null $N$ such that for every $m$ there are a finite partition $N=P_1\sqcup\dots\sqcup P_k$ and pairwise disjoint sets $Q_{il}$ ($i\le k$, $l\le m$), disjoint from $N$, with $P_i\le_1Q_{il}$.
\end{proposition}
\status{proved (\src{Q1} Cors.~Q1.5, Q1.6); confirmed (\src{Q12\_referee} \S1.4). The first sentence uses Theorems~\ref{thm:E2}, \ref{thm:E3} (CHECKED) and \ref{thm:fullsupport} at world~2 by symmetry. (a) is Proposition~\ref{prop:MP}(a): $\nu_2(d_{b_0})\le\nu_2(A_j)\to0$. (b) chains the comparisons $1_{Q_{il}}-1_{P_i}$ ($i\le k$; disjoint positive and negative parts) to $N\le_1\bigsqcup_iQ_{il}$, and some of these $m$ disjoint unions has $\bar\mu\le1/m$; it needs only (a). Scope: certificates with a \emph{fixed} source. Consequence: any two-world counterexample of these types (simple, vanishing-target and fixed-source star certificates) needs both worlds non-additive. Round 13: this holds without assuming (Nest), because a world that is additive with integrable density (or has a nonzero leading charge, or is monotonely continuous) forces the selection to be nested (Theorem~\ref{thm:S22}, Corollary~\ref{cor:MC}); the scope is unchanged, and ratio certificates (below) are not excluded next to an additive world (\src{S1\_referee} \S5).}

\begin{proposition}\label{prop:N1}
Let $Z$ be non-null, $b_0=b_0(Z)$, and $\mathcal U_1(Z)\coloneqq\{A:Z\le_1A\}$.
\begin{enumerate}[leftmargin=*,label=(\alph*)]
\item $d_t\le_2A$ for every $A\in\mathcal U_1(Z)$ and $t\in[b_0,1)$ (Proposition~\ref{prop:MP}(a)).
\item For $\mathcal A\subseteq\mathcal U_1(Z)$ the following are equivalent: every charge $\nu\ge0$ has $\inf_{A\in\mathcal A}\nu(A)=0$; and $\inf\{\|f\|_\infty:f\in\co\{1_A:A\in\mathcal A\}\}=0$ (\emph{$\mathcal A$ is uniformly Ces\`aro-thin}). If they hold, every $\nu\in\mathcal V_2$ vanishes on $d_{b_0}$, every $\nu\in\mathcal V_1$ vanishes on some $c_{s_0}$ (Corollary~\ref{cor:MPloop}), and both worlds fail SP$^{\mathrm{reg}}$. Bounded overlap is a special case; summability is not.
\item If world~2 has SP$^{\mathrm{reg}}_2$ with weight $\nu_2$, then $\nu_2(A)\ge\nu_2(d_{b_0(Z)})>0$ for every $A\in\mathcal U_1(Z)$, so $\|f\|_\infty\ge\nu_2(d_{b_0(Z)})$ for every $f\in\co\{1_A:A\in\mathcal U_1(Z)\}$. Symmetrically with the worlds exchanged.
\end{enumerate}
\end{proposition}
\status{proved in a referee report (\src{Q12\_referee} \S4, N1; (b) by Sion's minimax theorem on the weak-* compact set of states), one reader. So simple certificates are exactly ``$Z\le_1$ a uniformly Ces\`aro-thin family''. Next to a world with SP$^{\mathrm{reg}}$, the world-1 upper set of each $Z$ is uniformly non-thin, with modulus $\nu_2(d_{b_0(Z)})$; that modulus depends on $Z$ only through the radius of Lemma~T$_1$, which has no rate. This is where general certificates with thin, varying sources escape. Round 13: Proposition~\ref{prop:MPsharp} sharpens the modulus to $\nu_2(d_{t_Z})$, where $1-t_Z$ is exactly the radius of Lemma~T$_1$ at $1$.}

\begin{proposition}\label{prop:N2}
Let $Z$ be non-null, $\mathcal A\subseteq\mathcal U_1(Z)$, $s_1\coloneqq\max\{s:c_s\le_1A\ \forall A\in\mathcal A\}$ and $t_1\coloneqq\min\{t:d_t\le_2A\ \forall A\in\mathcal A\}$ (both attained, by closed sections; $s_1>0$ and $t_1<1$ by Corollary~\ref{cor:MPloop}). If $s_1<t_1$, then $c_{s_1}\le_1d_{t_1}$ and $d_{t_1}\le_2c_{s_1}$: the cut pair $(c_{s_1},d_{t_1})$ is accepted by both worlds, each in its favourable direction.
\end{proposition}
\status{proved in a referee report (\src{Q12\_referee} \S4, N2, via (K) and the mirror of Lemma~DL), one reader. Undominatedness forbids only the reverse acceptances ($Q\le_1P$ and $P\le_2Q$), and every comparison generated by the mirrored ``ping-pong'' of Corollary~\ref{cor:MPloop} goes in the favourable direction. So undominatedness, (C) and (K) give no contradiction from mirrored pairs; a proof that mirrored pairs cannot occur must use closedness quantitatively. This confirms \src{Q2} \S4.3's diagnosis from a second direction. Round 13: the quantitative form suggested here earlier, a rate in Lemma~T, is not known to bear on mirrored pairs, whose sources are fixed (item (R3) of part~(g)); the ping-pong never raises the exported segment (Lemma~\ref{lem:loop}), and the world-2 half of a mirrored pair meets (K) by itself on the cut pairs that would transport the world-1 targets (Lemma~\ref{lem:block}). Closedness is genuinely needed: without it, at both worlds, mirrored pairs exist (Proposition~\ref{prop:S31}).}

\emph{Shape of a surviving certificate} (\src{Q2} \S4.1; proved necessary conditions, confirmed with a wording repair in \src{Q12\_referee} \S2.4). If world~1 fails SP$^{\mathrm{reg}}_1$ under (Nest): (a) the exported set is an initial segment $c_a$, $a>0$, with a certificate as in Corollary~\ref{cor:IS}(iv), and every set disjoint from $c_{a^*}$ is charged by the single endorsing charge $\bar\nu$; (b) for each fixed source $P$ and point $x$ there is $\eta(P,x)>0$ such that no non-null target $Q$ inside the $\eta$-slab at $x$ has $P\le_1Q$, but $\eta$ has no lower bound uniform in $P$ (\src{Q2}'s wording ``no exported mass can be sent into a fixed small slab'' is a reading: Lemma~T fixes the source, and certificate sources vary); (c) left-fineness (Theorem~\ref{thm:LF}) holds and bounds only mass on the left; (d) strict escape and the drop lemmas DL, DR hold; (e) if the certificate can be taken simple, it is mirrored at world~2 through the same targets (Proposition~\ref{prop:MP}), neither world is complete, (CP-a) or has a nonzero leading charge (Corollary~\ref{cor:MPprime}), and neither world accepts all strict cut comparisons of a charge charging the far end (Proposition~\ref{prop:NS}); (f) integrating the certificate against $\bar\nu$ gives $0\le\lambda_k(\bar\nu(Q_k)-\bar\nu(P_k))\le\epsilon$ for each $k$, so each target is $\bar\nu$-almost as light as its source while the sources cover $c_a$ with weight $\ge1-\epsilon$. Charges concentrated at $1^-$ (for instance a Banach limit of $\mu(\cdot\cap d_b)/\mu(d_b)$ as $b\to1$) always lie in $\mathcal V_1$, so $\mathcal V_1\ne\varnothing$ carries no information; nothing proved excludes $a^*=1$, the extreme form of the open case (\src{Q2} \S4.2(ii), confirmed).

\emph{Ratio certificates} (\src{Q12\_referee} \S4, N4; an observation, proved by inspection). Proposition~\ref{prop:Q15}(b) handles certificates with a fixed exported set and disjoint targets in each round. A certificate built from thin sources $P_i$ (pieces of $c_a$ of measure $a/K$), each with $L$ targets $Q_{il}$ that overlap \emph{across} $i$ with multiplicity at most $\epsilon L$, has average $\bar\mu(Q_{il})\le\epsilon/K$ against $\bar\mu(P_i)=a/K$, a scale-free ratio $\epsilon/a$. Neither Proposition~\ref{prop:Q15}(a) nor Lemma~T applies (the sources shrink, and with them possibly the radius of Lemma~T), and chaining does not apply (the targets overlap). So even next to an additive world the open case is exactly ``ratio certificates with thin sources''. \emph{Round 16:} such certificates occur in a full system next to an additive world (Theorem~\ref{thm:D21}, step~7 of its proof).

\emph{Round 13: rates, thresholds and quantitative closedness} (\src{S4}, \src{S3} \S6; refereed in \src{S34\_referee}). The standing hypotheses of this part are kept. For a world-2 error set $B$ write $\varphi_2(B)\coloneqq\min\{t:d_t\le_2B\}$, so that the world-2 section of the regular model with $E_2=B$ is $[\varphi_2(B),1]$. A sequence of non-null sets $Q_k$ \emph{concentrates at} $x\in[0,1]$ if $Q_k$ lies in the slab $(x-r_k,x+r_k)$ with $r_k\to0$. Consider the following rate conditions, each required for every $m>0$:
\begin{itemize}[leftmargin=*]
\item[\textbf{(SR$_1$)}] $\inf\{\varphi_1(P):\bar\mu(P)\ge m\}>0$ (a uniform rate in strict escape at $b=0$); \textbf{(SR$_2$)} $\inf\{1-\varphi_2(W):\bar\mu(W)\ge m\}>0$;
\item[\textbf{(R3$_1$)}] for every $g>0$ there is $\eta>0$ such that, whenever $P\subseteq c_b$, $\bar\mu(P)\ge m$ and $x\ge b+g$, no non-null $Q$ inside the slab $(x-\eta,x+\eta)$ has $P\le_1Q$; \textbf{(R3$_1^{\mathrm u}$)} one $\eta$ works for every $P$ with $\bar\mu(P)\ge m$ and every $x\in[0,1]$; \textbf{(R3$_2$)}, \textbf{(R3$_2^{\mathrm u}$)} the mirror forms ($W\subseteq d_b$, $x\le b-g$, $W\le_2Q$).
\end{itemize}
A \emph{doubly-light sequence} is a sequence of sets $S_k\subseteq M$ with $\inf_k\bar\mu(S_k)>0$, together with non-null $R_k$ concentrating at $1$ and non-null $V_k$ concentrating at $0$, such that $S_k\le_1R_k$ and $S_k\le_2V_k$ for all $k$: sets of measure bounded below that world~1 ranks below sets concentrating at the right end and world~2 ranks below sets concentrating at the left end.

\begin{theorem}\label{thm:RATE}
The following are equivalent: (SR$_1$), (R3$_1^{\mathrm u}$), (R3$_1$), (SR$_2$), (R3$_2^{\mathrm u}$), (R3$_2$), and \textbf{(NDL)}: there is no doubly-light sequence. If (SR$_1$) fails at $m$, there is a doubly-light sequence with $\bar\mu(S_k)\ge m/3$ and $S_k\subseteq[m/3,1-m/3)$.
\end{theorem}
\begin{proof}[Proof sketch]
\emph{(SR$_1$)$\Rightarrow$(R3$_1^{\mathrm u}$).} If (R3$_1^{\mathrm u}$) fails at $m$, there are $P_k$ with $\bar\mu(P_k)\ge m$ and, after passing to a subsequence, non-null $Q_k$ concentrating at one point $\bar x$ with $P_k\le_1Q_k$. With $s_m>0$ the bound of (SR$_1$) at $m$, the closed section $[0,\varphi_1(P_k)]$ contains $s_m$, so $c_{s_m}\le_1P_k\le_1Q_k$, against Lemma~T$_1$ for the fixed source $c_{s_m}$. \emph{(R3$_2$)$\Rightarrow$(SR$_1$).} If $\bar\mu(P_k)\ge m$ and $\varphi_1(P_k)\to0$, put $g\coloneqq m/3$ and take $t_k\in(\varphi_1(P_k),g)$, $t_k\to0$. (K) at $t_k$ and cancellation give $d_{t_k}\cap P_k\le_2c_{t_k}\setminus P_k\eqqcolon V_k$, so $W_k\coloneqq P_k\cap d_g\le_2V_k$ with $\bar\mu(W_k)\ge2m/3$, and $V_k$ is non-null by strict (G) and concentrates at $0$; this violates (R3$_2$). The mirror arguments give (SR$_2$)$\Rightarrow$(R3$_2^{\mathrm u}$) and (R3$_1$)$\Rightarrow$(SR$_2$), and (R3$_i^{\mathrm u}$)$\Rightarrow$(R3$_i$) is trivial; this closes the cycle. A doubly-light sequence violates (R3$_1^{\mathrm u}$). Conversely, if (SR$_1$) fails, keep $W_k\le_2V_k$. Then $\varphi_2(W_k)\to1$: otherwise $d_{1-s'}\le_2W_k\le_2V_k$ along a subsequence for some $s'>0$, against Lemma~T$_2$ at $0$ for the fixed source $d_{1-s'}$; for $\tau_k\in(1-g,\varphi_2(W_k))$ with $\tau_k\to1$, (K) at $\tau_k$ gives $c_{\tau_k}\cap W_k\le_1d_{\tau_k}\setminus W_k\eqqcolon R_k$, and $S_k\coloneqq W_k\cap c_{1-g}$ is doubly light.
\end{proof}
\status{proved (\src{S4} Thm.~RATE); confirmed, every implication re-derived (\src{S34\_referee} \S6). Inputs: (K) at interior cuts, closed sections at regular models (Lemma~B0 and Lemma~T), strict (G), cancellation and transitivity; no (C), and no (LA) beyond Lemma~N0. So the qualitative rate is a single joint property of the two worlds. In a doubly-light sequence no non-null set lies in infinitely many $S_k$ (Lemma~T$_1$), so the sources must move. The configuration is the exact dual of a mirrored pair: a mirrored pair has fixed sources and thin, spread targets (``thin sets heavy at both worlds''); a doubly-light sequence has moving fat sources and concentrated targets (``fat sets light at both worlds''). The comparisons that the definition and (K) produce are cut pairs in the favourable direction, which undominatedness does not forbid (\src{S4} \S2.4, confirmed as a statement about these comparisons; its rider, that any exclusion must use closedness beyond Lemma~T, is a heuristic, \src{S34\_referee} \S6). Rates are per system: no bound depending only on $m$ and the gap holds across systems (Proposition~\ref{prop:ExA}).}

\begin{conjecture}\label{conj:NDL}
Under the standing hypotheses of this part, with closed sections for every model, there is no doubly-light sequence; equivalently (Theorem~\ref{thm:RATE}), the qualitative rate (R3$_1$) holds.
\end{conjecture}
\status{CONJECTURE (\src{S4} Conj.~S4.1; label confirmed, \src{S34\_referee} \S10.1); no evidence either way. It holds whenever one world has a strictly positive endorsing charge vanishing at the far end of the chain (Corollary~\ref{cor:SR}). A counterexample would be the first system with every axiom and without a uniform strict-escape rate. It is not known to bear on (R1) or (R2). This is the restated form of item (R3) of part~(g); the linear form of the rate is false (Proposition~\ref{prop:ExA}).}

\begin{corollary}\label{cor:SR}
If world~2 has an endorsing charge $\nu_2\in\mathcal V_2$ that is strictly positive on non-null sets and has $\nu_2(c_t)\to0$ as $t\downarrow0$, the equivalent conditions of Theorem~\ref{thm:RATE} hold; likewise with the worlds exchanged and $\nu_1(d_t)\to0$ as $t\uparrow1$. This applies whenever world~2 has a strictly positive, uniformly absolutely continuous endorsing charge, in particular when world~2 is additive with integrable density.
\end{corollary}
\begin{proof}[Proof sketch]
The countably additive part of $\nu_2$ has a density $>0$ a.e.\ (\src{N19} Lemma~R1), so $\kappa(m)\coloneqq\inf\{\nu_2(B):\bar\mu(B)\ge m\}>0$. If (SR$_1$) failed at $m$, the relevant comparison $d_{t_k}\cap P_k\le_2c_{t_k}\setminus P_k$ of the proof of Theorem~\ref{thm:RATE}, endorsed by $\nu_2$, would give $\kappa(m/2)\le\nu_2(c_{t_k})\to0$.
\end{proof}
\status{proved (\src{S4} Cor.~SR); confirmed (\src{S34\_referee} \S7). The hypothesis is implied by uniform absolute continuity, so the corollary covers the whole setting of \src{Q1} Q1.C1; it is not implied by SP$^{\mathrm{reg}}_2$ alone (a purely finitely additive part may sit at $0^+$). \emph{Reading.} In the setting of Q1.C1 the qualitative rate holds outright, while (R2) is open there; so the qualitative rate does not imply (R2). \emph{Round 16:} (R2) has a negative answer in that setting and Q1.C1 is false (Theorem~\ref{thm:D21}).}

For non-null $Z$ put $t_Z\coloneqq\sup\{t\in[0,1):Z\le_1d_t\}$; the set of such $t$ is an initial segment and contains $0$. Let $\eta_1(Z,1)\coloneqq\sup\{\eta>0:\text{no non-null }Q\subseteq d_{1-\eta}\text{ has }Z\le_1Q\}$, the radius of Lemma~T$_1$ at $1$, and $\kappa_1(Z)\coloneqq\inf\{\bar\mu(A):Z\le_1A\}$; for non-null $W$ put $\kappa_2(W)\coloneqq\inf\{\bar\mu(B):W\le_2B\}$ and $\mathcal U_2(W)\coloneqq\{B:W\le_2B\}$.

\begin{proposition}\label{prop:MPsharp}
Let $Z$ be non-null.
\begin{enumerate}[leftmargin=*,label=(\alph*)]
\item $t_Z<1$, and $t_Z>0$ unless $Z=M$.
\item For every $A\in\mathcal U_1(Z)$ and every $t\in[t_Z,1)$, $d_t\le_2A$.
\item For every $t<t_Z$ some $A\in\mathcal U_1(Z)$ has $d_t\not\le_2A$. So $\bigcap_{A\in\mathcal U_1(Z)}\{t:d_t\le_2A\}=[t_Z,1]$.
\item For $\nu\in\mathcal V_2$, $\nu(d_{t_Z})\le\inf_{A\in\mathcal U_1(Z)}\nu(A)\le\lim_{t\uparrow t_Z}\nu(d_t)$; for countably additive $\nu$ both bounds equal $\nu(d_{t_Z})$. In particular, if world~2 is additive with weight $\bar\mu$, then $\kappa_1(Z)=1-t_Z$.
\item $\eta_1(Z,1)=1-t_Z$. Hence, if world~2 is additive with weight $\bar\mu$, $\kappa_1(Z)=\eta_1(Z,1)$.
\end{enumerate}
\end{proposition}
\begin{proof}[Proof sketch]
(b) For $t>t_Z$ consider the cut pair $(A\cap c_t,d_t\setminus A)$. If $A\cap c_t\le_1d_t\setminus A$, cancellation gives $A\le_1d_t$, so $Z\le_1d_t$, against $t>t_Z$; so (K) gives $d_t\setminus A\le_2A\cap c_t$, that is $d_t\le_2A$. The point $t=t_Z$ is included because world-2 sections are closed. (c) Take $A\coloneqq d_{t'}$ with $t<t'<t_Z$ and $Z\le_1d_{t'}$, and use strict (G). (a) $t_Z<1$ by Lemma~T$_1$ at $1$; if $t_Z=0$, (b) with $A=Z$ and closedness give $M=d_0\le_2Z$, so $Z=M$. (d) The lower bound is (b) with endorsement; the upper bound uses $A=d_{t'}$, $t'\uparrow t_Z$. (e) If $Z\le_1Q$ with $Q\subseteq d_{1-\eta}$ non-null, then $Z\le_1d_{1-\eta}$ by monotonicity, and conversely $Q=d_t$ is admissible.
\end{proof}
\status{(a)--(c) proved (\src{S4} Prop.~MP$\sharp$); confirmed (\src{S34\_referee} \S8.1). (d) confirmed with repair: \src{S4}'s upper bound $\nu(d_{t_Z})$ would need $Z\le_1d_{t_Z}$, i.e.\ attainment of the supremum, which is not a section of a selected model and is not known to be closed; the bounds above are the referee's, and the lower bound $\nu(d_{t_Z})$ is stronger than \src{S4}'s. (e) proved in a referee report by monotonicity (\src{S34\_referee} \S8.1), one reader. NUMERICS \src{s4py/rate\_examples.py} (2) and \src{s34ref\_py/s4\_exampleA.py}: with world~1 additive with density $(1-x)^{-1/2}$ and world~2 $=\bar\mu$, the minimum of $\bar\mu(A)$ over $A$ with $c_s\le_1A$ equals $1-t_{c_s}$ to $10^{-5}$. Proposition~\ref{prop:MP}(a) is the special case $t\ge b_0(Z)\ge t_Z$, and (d) sharpens the modulus of Proposition~\ref{prop:N1}(c) to $\nu_2(d_{t_Z})$: if either world has SP$^{\mathrm{reg}}$, there is no mirrored pair, with this modulus (\src{S4} (W-a), confirmed).}

\begin{lemma}\label{lem:loop}
\begin{enumerate}[leftmargin=*,label=(\alph*)]
\item For all $s,t\in(0,1)$: $c_s\le_1d_t$ or $d_t\le_2c_s$.
\item Put $T(s)\coloneqq t_{c_s}$ and $T^*(t)\coloneqq\inf\{s:d_t\le_2c_s\}$. For every $s$ and every $t>T(s)$: $\mathcal U_1(c_s)\subseteq\mathcal U_2(d_t)\subseteq\mathcal U_1(c_{s'})$ for every $s'\le T^*(t)$, and $T^*(t)\le s$.
\end{enumerate}
\end{lemma}
\begin{proof}[Proof sketch]
(a) If $s<t$, use (K) at a cut in $(s,t)$. If $t<s$, cancellation turns $c_s\le_1d_t$ into $c_t\le_1d_s$ and $d_t\le_2c_s$ into $d_s\le_2c_t$, and (K) applies at a cut in $(t,s)$. If $s=t$, use (K) at $s$. (b) The inclusions are Proposition~\ref{prop:MPsharp}(b) and its mirror (at $s'=T^*(t)$ by closedness). Since $t>T(s)$, $c_s\not\le_1d_t$, so (a) gives $d_t\le_2c_s$, i.e.\ $T^*(t)\le s$.
\end{proof}
\status{proved (\src{S4} Loop lemma); confirmed with repair (\src{S34\_referee} \S8.2): \src{S4} applied (K) at a cut in $(s,t)$, which is empty when $t\le s$, and $t\in(T(s),s]$ does occur, even in Proposition~\ref{prop:ExA} for $s>3/4$ (there $T(s)=1-(1-\sqrt{1-s})^2$, so $T(0.8)\approx0.694$); (a) is the referee's repair. \emph{Reading.} Each round of the ping-pong of Corollary~\ref{cor:MPloop} returns a world-1 statement about an initial segment no longer than the one it started from: a mirrored pair is a fixed point of the ping-pong, and (K) is used up in reaching it. This sharpens Proposition~\ref{prop:N2}. (\src{S4}'s word ``idempotent'' means ``does not raise the segment''.)}

\begin{lemma}\label{lem:block}
Let $Z$ be non-null, $A\in\mathcal U_1(Z)$ and $t\in[t_Z,1)$. Then $Q\le_2A\cap c_t$ for every $Q\subseteq d_t\setminus A$. So on every cut pair $(A\cap c_t,Q)$ with $Q\subseteq d_t\setminus A$, (K) is met by world~2 in its favourable direction, and (K) never forces the world-1 packing $A\cap c_t\le_1Q$ at a cut $t\ge t_Z$.
\end{lemma}
\status{proved in a referee report (\src{S34\_referee} \S10.2: Proposition~\ref{prop:MPsharp}(b) and monotonicity on the left), one reader. \emph{The blocking effect.} In a mirrored pair this is the self-protection of the configuration: the world-2 half of the pair satisfies (K) by itself on exactly the cut pairs that would transport the thin world-1 targets towards $1$. Next to world~2 $=\bar\mu$ it is harmless, since every $A\in\mathcal U_1(Z)$ then has $\bar\mu(A)\ge1-t_Z$. A proof that mirrored pairs cannot occur must therefore obtain world-1 packing of a thin $A\cap c_t$ from (K) at cuts below $t_Z$, or from world-2 closedness at sources that are not segments. \src{S4} \S5.2 (proved, world~2 $=\bar\mu$: the transport of Theorem~\ref{thm:LF} turns any subset of $c_a$, however thin, into receivers of total measure $\ge a$) is confirmed as stated, but its reading that closedness is ``blind to the fibre'' is misleading: in that class (K) itself forces $P\le_1Q$ for every cut pair with $\bar\mu(Q)>\bar\mu(P)$, and the obstruction is the blocking effect (\src{S34\_referee} \S10.1). \src{S4} \S5.1 (``rates are irrelevant to (R1)'') is an observation, not a proved statement (GAP, overlabelled): the sources of a mirrored pair are fixed, so Lemma~T already gives their radii, and the known consequences of rates are lower bounds on $\varphi_1,\varphi_2$, while (R1) needs an upper bound.}

Say that \textbf{(QC$_1$)} holds if $\kappa_1(Z)>0$ for every non-null $Z$, and \textbf{(QC$_2$)} if $\kappa_2(W)>0$ for every non-null $W$: no non-null set is ranked by world~1 (world~2) no worse than sets of arbitrarily small measure. Write \textbf{(R1$'$)} for the failure of (QC$_1$).

\begin{proposition}\label{prop:QCJ}
(QC$_1$), (QC$_2$), ``$\kappa_1(c_s)>0$ for every $s\in(0,1)$'' and ``$\kappa_2(d_t)>0$ for every $t\in(0,1)$'' are equivalent. Moreover $\kappa_1(Z)\ge\kappa_1(c_{s_0})$ with $s_0=s_0(Z)$ as in Corollary~\ref{cor:MPloop}, and $\kappa_1(c_s)\ge\kappa_2(d_t)$ for every $t\in[T(s),1)$. If (QC$_1$) holds, no non-null $Z$ is $\le_1$ every member of a uniformly Ces\`aro-thin family; that is, there is no mirrored simple pair. So (R1$'$) is a mirrored configuration: there are $s>0$ and sets $A_k$ with $\bar\mu(A_k)\to0$ and $c_s\le_1A_k$, and then $d_t\le_2A_k$ for every $t\ge T(s)$.
\end{proposition}
\begin{proof}[Proof sketch]
If $Z\le_1A_k$ with $\bar\mu(A_k)\to0$, Proposition~\ref{prop:MP}(a), which needs no overlap condition, gives $d_{b_0(Z)}\le_2A_k$: this is the failure of (QC$_2$); the mirror gives the converse. Corollary~\ref{cor:MPloop} gives $\mathcal U_1(Z)\subseteq\mathcal U_1(c_{s_0})$, hence the first inequality, so the segment form implies (QC$_1$). The second inequality is Proposition~\ref{prop:MPsharp}(b). Finally, if $\mathcal A\subseteq\mathcal U_1(Z)$ is uniformly Ces\`aro-thin and $f=\sum_i\lambda_i1_{A_i}$ has $\|f\|_\infty<\epsilon$, then $\min_i\bar\mu(A_i)\le\bar\mu(f)<\epsilon$, so $\kappa_1(Z)=0$.
\end{proof}
\status{the implication from (QC$_1$) to the absence of mirrored pairs is proved (\src{S3} \S6); confirmed (\src{S34\_referee} \S5.1), which notes that it needs no minimax argument. The equivalences and the two inequalities are proved in a referee report (\src{S34\_referee} \S5.2, Prop.~QC-J), one reader. (QC) is \emph{a priori} stronger than the absence of mirrored pairs: no system is known in which one holds and the other fails, and the two negations differ in form only (vanishing measure versus Ces\`aro-thin). So (R1$'$) is implied by a mirrored pair, is itself mirrored, and is a weaker form of (R1): excluding (R1$'$) excludes (R1).}

\begin{conjecture}\label{conj:QC}
Under the standing hypotheses of this part, with closed sections for every model, (QC$_1$) holds.
\end{conjecture}
\status{CONJECTURE (\src{S3} Conj.~S3.C1; label confirmed, \src{S34\_referee} \S5.1). Evidence: it holds next to a world with a strictly positive, uniformly absolutely continuous endorsing charge (Proposition~\ref{prop:Q15}(a)); it holds in band-and-translates worlds with closed world-1 sections (Proposition~\ref{prop:S32}(d)); and it fails at both worlds in the system of Proposition~\ref{prop:S31}, where only closed sections fail. It is the measure analogue of Lemma~T (sets of vanishing measure instead of sets shrinking to finitely many points of the chain). A proof needs a packing step for spread targets; by Proposition~\ref{prop:S33} no system with closed sections accepts the packing itself, and by Lemma~\ref{lem:block} (K) at cuts $t\ge t_Z$ cannot supply it.}

\begin{lemma}\label{lem:UE}
Let $\mathcal A$ be a family of pairwise disjoint sets with $c_x\le_1A$ for every $A\in\mathcal A$, and suppose there is $\kappa>0$ with $c_{x+\kappa}\le_1c_x\cup A$ for every $A\in\mathcal A$ (\emph{uniform escape at $x$}). Then $c_{x+\kappa}\le_1A\cup A'\cup B$ for any three distinct $A,A',B\in\mathcal A$.
\end{lemma}
\status{proved (\src{S4} Lemma~UE, by cancellation and sum closure of disjoint comparisons); confirmed (\src{S34\_referee} \S10.1; in step~3 of \src{S4} the substitution is cancellation by the disjoint set $B\cap d_{x+\kappa}$, not sum closure; the step is correct). Consequences (proved; inconclusive as a route): disjoint three-fold unions of a disjoint thin family are again uniformly thin, so uniform escape at the $s_1$ of Proposition~\ref{prop:N2} raises $s_1$ by $\kappa$ along these unions, and under iteration the escape radii must be summable. Where uniform escape fails at $x$, (K) at $x+\kappa$ gives $A\cap d_{x+\kappa}\le_2[x,x+\kappa)\setminus A$ for some $A\in\mathcal A$: a thin world-2 source below a narrow slab, about which Lemma~T$_2$ and every rate say nothing. So the route ends with exactly a thin residue.}

\begin{proposition}\label{prop:ExA}
Let $n=2$, with the path $b\mapsto C_b$ and the coordinate $x(h)$ of part~(e). Let $\theta:[0,1)\to(0,\infty)$ be continuous, strictly increasing and integrable, with $\theta(x)\to\infty$ as $x\to1$; for instance $\theta(x)=(1-x)^{-1/2}$. Let world~2 be $\psi_2(X)\coloneqq\bar\mu(E_2X)$, let world~1 be $\psi_1(X)\coloneqq\int_{E_1X}\rho_1\,d\bar\mu$ with $\rho_1\coloneqq\theta\circ x$ on $M$ and $\rho_1\coloneqq1$ on $N_+$, and let $C_p$ be the $p$-Bayes model. Then the system satisfies every axiom, with closed sections for every model and $\C_L$ the selection, and both worlds add up. For $\theta(x)=(1-x)^{-1/2}$ and $P_m\coloneqq[\tfrac12-m,\tfrac12)$, $\eta_1(P_m,1)\le m^2/2$.
\end{proposition}
\begin{proof}[Proof sketch]
At an interior $p$ with $r\coloneqq p_2/p_1$, the Bayes model errs at world~1 on $h\in M$ exactly where $\theta(x(h))<r$; tie sets are level sets of $x$, hence null. So $C_p=C_{\theta^{-1}(r)}$ (or $C_0=P_1$ when $r\le\theta(0)$): the selection is coherent, nested and continuous. (W1)--(W4) of Theorem~\ref{thm:witness} hold (Bayes optimality; uniqueness at interior $p$; at a vertex, $\psi_1(X)=0$ forces $X=P_1$; continuity by dominated convergence), so the system is successful. For the radius, $\int_{P_m}\rho_1\,d\bar\mu\le\theta(\tfrac12)m=\sqrt2\,m$, while $\int_{d_{1-\eta}}\rho_1\,d\bar\mu=2\sqrt\eta$; so $P_m\le_1d_{1-\eta}$ once $\eta\ge m^2/2$.
\end{proof}
\status{proved (\src{S4} Example~A; axioms verified through Theorem~\ref{thm:witness}); confirmed (\src{S34\_referee} \S9.1). NUMERICS: $\eta/m^2\to\tfrac12$. Variants: $\theta=(1-x)^{-\alpha}$ gives $\eta_1$ of order $m^{1/(1-\alpha)}$, so any power $m^q$ with $q>1$ occurs (the referee's correction of \src{S4}'s ``$q\ge2$''); a logarithmic profile gives $e^{-c/m}$; and with a world-2 density $|x-x_0|^{-1/2}$ the same happens at an interior point $x_0$. So no rate linear in $\bar\mu(P)$ follows from the axioms, even when world~2 is $\bar\mu$, and rates are per system.}

\begin{proposition}\label{prop:LIN}
Let world~2 be additive with weight $\bar\mu$, and suppose there are $c>0$ and $m_0>0$ with $\eta_1(P,1)\ge c\,\bar\mu(P)$ whenever $\bar\mu(P)\le m_0$. Let $a\in(0,1)$ and $\epsilon<ca$. Then world~1 has no ratio certificate of the following shape, for any $K$ with $a/K\le m_0$: sources $P_1,\dots,P_K\subseteq c_a$ with $\bar\mu(P_i)=a/K$, and for each $i$ targets $Q_{il}$ with $P_i\le_1Q_{il}$ whose average measure is at most $\epsilon/K$.
\end{proposition}
\begin{proof}[Proof sketch]
$\eta_1(P_i,1)\ge ca/K$, so $t_{P_i}\le1-ca/K$ (Proposition~\ref{prop:MPsharp}(e)), and Proposition~\ref{prop:MPsharp}(b) gives $d_t\le_2Q_{il}$ for $t=1-ca/K$. Endorsement by $\bar\mu$ gives $\bar\mu(Q_{il})\ge ca/K$ for every $l$, against the average.
\end{proof}
\status{proved (\src{S4} Prop.~LIN); confirmed with repair (\src{S34\_referee} \S9.2): \src{S4}'s proof compared $\bar\mu(Q_{il})$ with $\nu_2(Q_{il})$, which match only for $\nu_2=\bar\mu$, so the statement is given here for world~2 $=\bar\mu$; for a general SP weight $\nu_2$ with $\nu_2(d_t)\ge c'(1-t)$ the averaging must be read in $\nu_2$, or $\nu_2\le C\bar\mu$ assumed. \emph{Reading.} Through Proposition~\ref{prop:MP}, only a rate linear in $\bar\mu(P)$ excludes the ratio certificates of N4, and Proposition~\ref{prop:ExA} shows that the linear rate is not a consequence of the axioms. A proof of (R2) by rates would have to derive linearity from the failure of SP$^{\mathrm{reg}}$ itself; no such mechanism is known. \emph{Round 16:} (R2) has a negative answer (Theorem~\ref{thm:D21}); in that system the generator sources $P$ have $\eta_1(P,1)$ at most of order $\bar\mu(P)^2$ (\src{D2\_referee} \S5, \src{D2\_referee2} \S2.7), a superlinear rate, as Proposition~\ref{prop:LIN} requires of any ratio certificate next to $\bar\mu$.}

\paragraph{(d) Two worlds without nestedness.} Let $n=2$, $K'=[\alpha',\beta']\subseteq(0,1)$ in the weight $t$ of world~2, $M^+\coloneqq\{g\in U:g_1>0>g_2\}$ with $\tau(g)\coloneqq g_1/(g_1-g_2)\in(\beta',1)$, and $M^-\coloneqq\{g\in U:g_1<0<g_2\}$ with $\sigma(g)\coloneqq-g_1/(g_2-g_1)\in(0,\alpha')$.

\begin{proposition}\label{prop:lattice}
The coherent members of $\M_U$ are the $D_{[\alpha,\beta]}$ with $\alpha\in[0,\alpha']$, $\beta\in[\beta',1]$, and
\[
E_1(D_{[\alpha,\beta]})=\{g\in M^-:\sigma(g)<\alpha\}\cup\{g\in M^+:\tau(g)\le\beta\},\qquad E_2=M\setminus E_1 .
\]
So coherent world-1 error sets form a two-parameter lattice, increasing in $(\alpha,\beta)$ for the product order, and $\C_L$ is nested iff $(\alpha,\beta)$ is comonotone along it.
\end{proposition}
\status{proved (\src{Q2} \S5.1); confirmed (\src{Q12\_referee} \S2.5). NUMERICS: \src{q12ref\_py/} \src{q12\_lattice.py}, 28{,}749 mixed gambles, no mismatch). This is the family $X_{\alpha,\beta}$ of \src{N17} N17.8.}

For any selection (nested or not), $S_1(X)\coloneqq\{p:C_p\preceq_1X\}$ and $S_2(X)$ are closed, $0\in S_1(X)$, $1\in S_2(X)$ and $S_1(X)\cup S_2(X)=[0,1]$, so $S_1(X)\cap S_2(X)\ne\varnothing$: this is Theorem~\ref{thm:kkm} for $n=2$ (\src{Q2} \S5.2).

\begin{lemma}\label{lem:GSE}
Let $n=2$ and $C_p\in\C_L$, and let $X$ be regular with $E_1X\supsetneq E_1C_p$ (non-null difference). Then some selected $C_q$ has $C_q\preceq_1X$, $C_q\prec_2C_p$ and $C_q\not\preceq_1C_p$.
\end{lemma}
\status{proved (\src{Q2} Lemma~G-SE: take $q\in S_1(X)\cap S_2(X)$; $E_2X\subsetneq E_2C_p$ by regularity; undominatedness of $C_p$); confirmed (\src{Q12\_referee} \S2.5). In the nested case it is Proposition~\ref{prop:strictescape}. Without nestedness nothing beyond it is proved: Lemma~N0, the cut-pair form, DL, DR, Lemma~T, Theorem~\ref{thm:LF}, Corollary~\ref{cor:IS} and Proposition~\ref{prop:MP} all parametrise the selection by one cut, while a non-nested selection has two independent cuts (in $M^-$ and in $M^+$) and sections that need not be intervals. In $n=3$, \src{Q2} \S6 records that these results do not transfer to comparisons out of \emph{interior} selected models (a world-1 failure can be absorbed by either of two other worlds); on edges they do transfer, under (Nest$_e$) (Proposition~\ref{prop:Q35}; \src{Q3\_referee} \S4.3 repairs \src{Q2} \S6 accordingly). Round 13: Corollary~\ref{cor:continuum}(d) gives strict escape without nestedness, and Theorems~\ref{thm:S22} and~\ref{thm:MC} give nestedness in the leading-charge and monotonely continuous classes.}

\emph{Round 13: nestedness in two worlds} (\src{S1}, \src{S2}; refereed in \src{S1\_referee}, \S\S1--7, 10, and \src{S2\_referee} \S3). In this part, ``every axiom'' means strict (G), (IB), (LA), (S) and (C) for $n=2$ on $\M_U$. The first result settles Conjecture~\ref{conj:Q2N} below in the leading-charge class.

\begin{theorem}\label{thm:S22}
Let $n=2$ with every axiom. If some world $i$ has a leading charge $\pi$ nonzero on $M$, then $\{E_1C:C\in\C_L\}$ is a chain, and each member of $\C_L$ is determined a.e.\ by its $\pi$-level: $E_iF=E_iF'$ a.e.\ whenever $F,F'\in\C_L$ have $\pi(E_iF)=\pi(E_iF')$. Hence (Nest) holds, and all of part~(c) and of \S\ref{sec:open-ordinal}(c) applies.
\end{theorem}
\begin{proof}[Proof sketch]
Say $i=1$. By Theorem~\ref{thm:S21}, $\pi$ is strictly positive. Take $m$ and the generic levels of its proof on the one edge $\Delta=e_{12}$. For $n=2$ every regular $X$ is an edge test model, so at a generic level $v$ the selected $C(v)$ is optimal at both worlds among all regular $X$ with $\pi(E_1X)=v$, and every member of $\C_L$ of level $v$ is tied with $C(v)$. A generic $C=C(v)$ is nested with every $F\in\C_L$: if they crossed, the differences $A$ and $B$ would lie in opposite halves $M^\mp$ (Proposition~\ref{prop:lattice}); take bounded non-null $P\subseteq A$, $Q\subseteq B$ off one carrier of the purely finitely additive part, with $\pi(P)=\pi(Q)$; then $X\coloneqq C\sd(P\cup Q)$ has level $v$, so $C\preceq_1X$ and $C\preceq_2X$, and Lemma~\ref{lem:PL} puts $X'\coloneqq F\sd(P\cup Q)$ in $\C_L$; but $X'$ is not a cone, hence incoherent. So the generic members form an increasing family ordered by level. For $F\in\C_L$ of level $v_F$, the union $Q_-$ of the $E_1C(v)$ with generic $v<v_F$ lies in $E_1F$ and has $\pi(Q_-)\ge v_F=\pi(E_1F)$, so $E_1F\setminus Q_-$ is $\pi$-null, hence null.
\end{proof}
\status{proved (\src{S2} Thm.~S2.2); confirmed (\src{S2\_referee} \S3.4), which notes that Theorem~\ref{thm:kkm} is not needed in the optimality step and that the last step uses monotonicity of $\pi$, not countable additivity. \emph{Consequence} (proved, \src{S2\_referee} \S3.4): in a two-world system with every axiom in which either world has a nonzero leading charge, all of part~(c) applies unconditionally, including Theorem~\ref{thm:LF}, Corollary~\ref{cor:IS}, Proposition~\ref{prop:MP}, Propositions~\ref{prop:N1}, \ref{prop:N2} and~\ref{prop:Q15}; for example, if world~2 is additive, world~1 has no simple certificate and no fixed-source star certificate. So the open two-world class ``without a nonzero leading charge'' no longer needs (Nest) as a hypothesis. \src{S2}'s cross-agreement (\S3; confirmed, \src{S2\_referee} \S3.2) is the special case of Lemma~\ref{lem:rigid} below with an incoherent splice. Recorded failures (\src{S2} \S5): a pure additive tie fails (C) (Proposition~\ref{prop:S20}(b)); lexicographic tie-breaking with frontiers collapses by this theorem; a half-vanishing leading charge is excluded by Theorem~\ref{thm:S21}. The ``forced anatomy'' that \src{S2} derived for that last case predates Theorem~\ref{thm:S21}, was not re-checked, and is not cited (\src{S2\_referee} \S6).}

\begin{lemma}\label{lem:level}
Let $n=2$ and let the ranking satisfy strict (G) and (S). If $U,V$ are regular and undominated (for instance $U,V\in\C_L$) and $U\sim_iV$ at one world, then $U\sim_jV$ at the other; moreover some selected $C_q$ is tied with both $U$ and $V$ at both worlds.
\end{lemma}
\begin{proof}[Proof sketch]
Say $U\sim_2V$, so $S_2(U)=S_2(V)\eqqcolon G$, which is closed and contains $1$. If $0\in G$, then $P_1\preceq_2U$, and strict (G) forces $U=P_1=V$. Otherwise $O\coloneqq[0,1]\setminus G$ is open, nonempty and proper. By the one-dimensional KKM cover, $O\subseteq S_1(U)\cap S_1(V)$, which is closed; so a boundary point $q$ of $O$ lies in all four sections, and undominatedness of $U$ and $V$ makes the four relations ties.
\end{proof}
\status{proved (\src{S1} Lemma~S1.1, the ``Level Lemma''); confirmed (\src{S1\_referee} \S1), which adds the selected $C_q$. Inputs: (S) (closed sections, and the KKM cover from competitiveness), strict (G) and undominatedness. Not used: (C), (LA), (IB), or coherence beyond regularity.}

\begin{corollary}\label{cor:continuum}
Let $n=2$ with every axiom.
\begin{enumerate}[leftmargin=*,label=(\alph*)]
\item On $\C_L$, $\preceq_1$ is a total preorder, $U\preceq_1V$ iff $V\preceq_2U$, and the ties at one world are exactly the ties at the other (\emph{double ties}).
\item The set $\mathfrak L$ of double-tie classes of $\C_L$, ordered by $\preceq_1$, is a compact linear continuum with end points $[P_1]$ and $[P_2]$, and $\varpi(p)\coloneqq[C_p]$ is a continuous surjection from the index set onto $\mathfrak L$.
\item For every model $X$, the classes $W$ with $W\preceq_1X$ form an interval $[\min,u_X]$, those with $W\preceq_2X$ form an interval $[l_X,\max]$, and $l_X\preceq u_X$; the classes in $[l_X,u_X]$ weakly dominate $X$.
\item (Strict escape without nestedness.) If $X$ is regular, $W\in\C_L$ and $E_1X\supsetneq E_1W$ with non-null difference, then every class in $[l_X,u_X]$ lies strictly above $[W]$; and mirror for $E_1X\subsetneq E_1W$.
\end{enumerate}
\end{corollary}
\status{proved (\src{S1} Cor.~S1.2); confirmed (\src{S1\_referee} \S2). (a) needs only Lemma~\ref{lem:level} and competitiveness. (b) The preimages of closed rays are the closed sections, and a connected compact linear order is a linear continuum. (d) is Lemma~\ref{lem:GSE} read in $\mathfrak L$. Also proved: \emph{inclusion closedness}, $\sup S\preceq_1X$ whenever $E_1X$ contains the essential union of the world-1 error sets of representatives of $S\subseteq\mathfrak L$, and the mirror. So in two worlds $\C_L$ is always one-dimensional in the order sense; what is open is only whether that order is compatible with inclusion of error sets.}

\begin{proposition}\label{prop:S13}
Let $n=2$ with every axiom. Some selection satisfying (S) for the given ranking has nested world-1 error sets iff one can choose one member of each double-tie class of $\C_L$ so that the chosen error sets are nested. In particular, if two distinct classes are singletons whose members cross, every selection for that ranking is non-nested.
\end{proposition}
\status{proved (\src{S1} Prop.~S1.3); confirmed (\src{S1\_referee} \S2): every selection meets every class (Theorem~\ref{thm:kkm}; $\C_L$ does not depend on the selection), and replacing $C_p$ by a tied representative keeps its sections and its competitiveness. So ``some nested selection with the same $\C_L$'' is equivalent to nested class representatives, and is strictly weaker than (Nest) when some class has two crossing members.}

\emph{Crossings.} Let $C,C'\in\C_L$ \emph{cross}, i.e.\ neither world-1 error set contains the other. By Proposition~\ref{prop:lattice}, after relabelling, $E_1C=m\cup B$ and $E_1C'=m\cup A$ with $m=E_1C\cap E_1C'$, $A=\{\alpha_C\le\sigma<\alpha_{C'}\}\subseteq M^-$ and $B=\{\beta_{C'}<\tau\le\beta_C\}\subseteq M^+$, both non-null. For $a\subseteq A$ and $b\subseteq B$ the \emph{splices} are the regular models $X_{a,b}$ with $E_1X_{a,b}=m\cup a\cup(B\setminus b)$ and $Y_{a,b}$ with $E_1Y_{a,b}=m\cup(A\setminus a)\cup b$; so $X_{\varnothing,\varnothing}=C$ and $X_{A,B}=C'$. $X_{a,b}$ is coherent iff $a$ is an initial $\sigma$-band of $A$ and $b$ a final $\tau$-band of $B$ (a \emph{segment pair}). By cancellation, $C\preceq_1X_{a,b}$ iff $b\le_1a$, and $C\preceq_2X_{a,b}$ iff $a\le_2b$. The \emph{pattern} of $(a,b)$ is the pair of relations (strictly smaller, strictly larger, tied, or incomparable) between $a$ and $b$ at worlds~1 and~2. A crossing is \emph{rigid} if $[C]\ne[C']$ in $\mathfrak L$, and \emph{tied} if $C$ and $C'$ are doubly tied.

\begin{lemma}\label{lem:rigid}
Let $n=2$ with every axiom, and let $C,C'\in\C_L$ cross. For every $a\subseteq A$ and $b\subseteq B$, the pattern of $(a,b)$ is (strictly smaller, strictly smaller), (strictly larger, strictly larger) or (tied, tied), and the last occurs only for $(a,b)\in\{(\varnothing,\varnothing),(A,B)\}$. So every splice other than $C$ and $C'$, including the coherent lattice points between them, is strictly better than $C$ at one world and strictly worse at the other.
\end{lemma}
\status{proved (\src{S1} Lemma~S1.4 and Thm.~S1.5, ``crossings are rigid''; CHECKED \src{s1py/patterns.py}); confirmed and strengthened (\src{S1\_referee} \S3). Competitiveness and undominatedness of $C$ and $C'$ leave five patterns; the two mixed ones (a tie at one world, incomparability at the other) are excluded for \emph{every} pair by the KKM cover at $X_{a,b}$ together with Lemma~\ref{lem:level} or Corollary~\ref{cor:continuum}(a), and (tied, tied) at a non-segment pair is excluded by (C). The restriction of (tied, tied) to the two corners is the referee's strengthening (proved): that pattern puts both $X_{a,b}$ and $Y_{a,b}$ in $\C_L$ by (C), and both are coherent only at the four corners, two of which strict (G) excludes. NUMERICS \src{s1ref\_py/patterns\_ref.py} (patterns computed by set algebra; on $5+5$ cells only the four corners have both splices coherent). In a tied crossing, \src{S1}'s ``every segment lattice point doubly tied with $C$ lies in $\C_L$'' is true but vacuous beyond $C'$.}

Say that world $i$ is \textbf{monotonely continuous on $M$}, \textbf{(MC$_i$)}, if for measurable $Y_k,Y,Z\subseteq M$: $Y_k\uparrow Y$ and $Y_k\le_iZ$ for all $k$ imply $Y\le_iZ$; and $Y_k\downarrow Y$ and $Z\le_iY_k$ for all $k$ imply $Z\le_iY$ (a Villegas-type continuity of the set order). Under (IB) the two clauses are equivalent. Weaker are the two halves of \textbf{(MC$_{\mathrm{sel}}$)}, which concern only sections of members of $\C_L$ along monotone chains of regular test models: \textbf{(MC$_{\mathrm{sel}}$)$_1$}, for $W\in\C_L$, if $E_1X_k\downarrow E_1X$ and $W\preceq_1X_k$ for all $k$, then $W\preceq_1X$; \textbf{(MC$_{\mathrm{sel}}$)$_2$}, for $W\in\C_L$, if $E_1X_k\uparrow E_1X$ and $W\preceq_2X_k$ for all $k$, then $W\preceq_2X$. (MC$_i$) implies (MC$_{\mathrm{sel}}$)$_i$.

\begin{theorem}\label{thm:MC}
Let $n=2$ with every axiom. If (MC$_{\mathrm{sel}}$)$_1$ holds, or (MC$_{\mathrm{sel}}$)$_2$ holds, then $\{E_1C:C\in\C_L\}$ is a chain, so the conclusion of Conjecture~\ref{conj:Q2N} holds. In particular this is so if either world is monotonely continuous on $M$. It suffices that the hypothesis holds for $W\in\{C,C'\}$ along splice chains of each putative crossing pair.
\end{theorem}
\begin{proof}[Proof sketch]
Assume (MC$_{\mathrm{sel}}$)$_1$ and let $C,C'$ cross; write $a\lhd b$ if $a$ is strictly smaller than $b$ at both worlds. By Lemma~\ref{lem:rigid} every non-segment pair is ordered by $\lhd$ one way or the other. (1) The bands $a_k\coloneqq\{s\le\sigma<s+1/k\}$, $\alpha_C<s$, are non-initial in $A$ and decrease to $\varnothing$. If $B\lhd a_k$ for all $k$, then $C\preceq_1X_{a_k,B}$, whose world-1 error set $m\cup a_k$ decreases to $m\subsetneq E_1C$, and (MC$_{\mathrm{sel}}$)$_1$ gives a contradiction with strict (G). Fix $a\coloneqq a_k$ with $a\lhd B$. (2) Put $b(\beta)\coloneqq\{\beta<\tau\le\beta_C\}$; the set of $\beta$ with $b(\beta)\lhd a$ is an up-set containing $\beta_C$ and not $\beta_{C'}$; let $\beta^*$ be its infimum. (3) For $\beta>\beta^*$, $C\preceq_1X_{a,b(\beta)}$, and these error sets decrease to that of $X_{a,b(\beta^*)}$ as $\beta\downarrow\beta^*$, so $b(\beta^*)\le_1a$; if $\beta^*=\beta_{C'}$ this contradicts $a\lhd B$. (4) Otherwise, for $\beta\in[\beta_{C'},\beta^*)$, $a\lhd b(\beta)$ gives $C'\preceq_1Y_{a,b(\beta)}$, whose world-1 error sets decrease as $\beta\uparrow\beta^*$; so $a\le_1b(\beta^*)$. (5) Then $a\sim_1b(\beta^*)$ at a non-segment pair, against Lemma~\ref{lem:rigid}. The case (MC$_{\mathrm{sel}}$)$_2$ is the mirror, with $Y$-splices from above at $C'$ and $X$-splices from below at $C$.
\end{proof}
\status{proved (\src{S1} Thm.~S1.6, under (MC$_1$), (MC$_2$), or both halves of (MC$_{\mathrm{sel}}$)); confirmed with repair (\src{S1\_referee} \S4). \src{S1} called (MC$_{\mathrm{sel}}$) weaker than (MC$_1$), which is false as stated, since (MC$_1$) gives only the world-1 half; the referee proves that each half alone suffices (Thm.~S1.6$'$), using the world-1 half at $C'$ along the $Y$-splices where \src{S1} had used the world-2 half at $C$. The equivalence of the two clauses of (MC$_i$) under (IB) is \src{S1\_referee} Prop.~R4.1 (proved, one reader). \emph{Necessary shape of a crossing} (\src{S1} \S4; confirmed with repair, \src{S1\_referee} \S5): in a rigid or tied crossing, every non-null non-final $b\subseteq B$ has $b\lhd A$, and the relation of $b$ with the initial $\sigma$-bands of $A$ switches at some point with no tie; both halves of (MC$_{\mathrm{sel}}$) fail, one at $C$ and the other at $C'$, depending on whether the switch is attained. \src{S1}'s ``for each non-initial $a$ the relation with $b(\beta)$ switches'' is false for $a$ with $B\lhd a$.}

\begin{corollary}\label{cor:MC}
Let $n=2$ with every axiom.
\begin{enumerate}[leftmargin=*,label=(\alph*)]
\item A world that is additive with an integrable density $>0$ a.e., or (CP-a) with a credal set of countably additive members, is monotonely continuous on $M$. So if either world is of this kind, (Nest) holds and all of part~(c) and of \S\ref{sec:open-ordinal}(c) applies without assuming it.
\item If $\{E_1C:C\in\C_L\}$ is not a chain, then both halves of (MC$_{\mathrm{sel}}$) fail, neither world is monotonely continuous on $M$, and neither world has a leading charge nonzero on $M$.
\end{enumerate}
\end{corollary}
\status{(a) proved (\src{S1} Cor.~S1.7(a), (b)): in both cases $X\preceq_iY$ iff $\int w_i(1_{E_iY}-1_{E_iX})\,dP\ge0$ for every $P$ in the credal set, with $w_i>0$ bounded, and dominated convergence applies member by member; confirmed (\src{S1\_referee} \S\S4.1, 5). Finitely additive members can fail (MC) (for instance a Banach-limit charge). Scope (\src{S1\_referee} \S5): the gloss ``any two-world counterexample needs both worlds non-additive'' of Proposition~\ref{prop:Q15} now holds without (Nest), within the scope of that proposition. (b) proved, from Theorem~\ref{thm:MC} and Theorem~\ref{thm:S22} (\src{S1\_referee} \S10, conditional there on the confirmation of S2.2, since given in \src{S2\_referee} \S3.4). The two criteria are logically independent: a lexicographic world with a countably additive first level has a leading charge but can fail (MC), and an incomplete countably additive (CP-a) world has (MC) but in general no single leading charge. So a non-nested two-world system with every axiom needs \emph{both} worlds to fail monotone continuity \emph{and} to lack a nonzero leading charge: neither world is additive, (CP-a) with countably additive members, or lexicographic with a charge first level. These are the worlds of the open SP$^{\mathrm{reg}}$ class (\src{S2} S2.3, confirmed).}

\begin{proposition}\label{prop:R71}
Let $n=2$ on $\M_U$ with $K'=[\alpha',\beta']\subseteq(0,1)$, let $\psi_1(X)\coloneqq\mu(E_1X)$ and $\psi_2(X)\coloneqq\mu(E_2X)$, and let $p\mapsto C_p=D_{[\alpha(p),\beta(p)]}$ be any path of coherent models with $\mu(E_1C_p)$ continuous and nondecreasing, $(\alpha,\beta)(0)=(0,\beta')$ and $(\alpha,\beta)(1)=(\alpha',1)$. Then strict (G), (IB), (LA) and (S) hold, with closed sections for every model, and every $C_p$ is coherent and undominated. If the path is not comonotone, for instance $(0,\beta')\to(\alpha',\beta')\to(\alpha_1,\beta_1)$ at constant mass $\to(\alpha',\beta_1)\to(\alpha',1)$ with $\alpha_1<\alpha'$, $\beta_1>\beta'$, the selection crosses. Only (C) fails.
\end{proposition}
\begin{proof}[Proof sketch]
On $M$ exactly one world errs at each gamble, and on $N_+$ both err or neither does; so $\psi_1(X)+\psi_2(X)=\mu(M)+2\mu(N_+\setminus X)$ for every $X$, and every regular model minimises $\psi_1+\psi_2$. Hence every regular model is undominated and competitive on $\{1,2\}$, and $P_1$, $P_2$ are competitive at the vertices. The sections $\{p:\mu(E_1C_p)\le\psi_1(X)\}$ and $\{p:\mu(M)-\mu(E_1C_p)\le\psi_2(X)\}$ are closed. (C) fails at any incoherent regular model with the error mass of some $C_p$.
\end{proof}
\status{proved in a referee report (\src{S1\_referee} \S7, Prop.~R7.1), one reader. NUMERICS \src{s1ref\_py/mu\_equal\_density.py}: exact angular masses, Gaussian $\mu$, $K'=[0.35,0.6]$; the selected $D_{[0.35,0.6]}$ and $D_{[0.1,0.828]}$ cross, with differences of mass $0.061$ each; 300 random models, no violation of the identity or of competitiveness. It strengthens Proposition~\ref{prop:plateau}, which is not a ranking. So in two worlds the missing ingredient is exactly (C): (C) is what must force nestedness, and the hypotheses of Proposition~\ref{prop:Q35} do not imply (Nest$_e$), even for a full ranking. For $n=3$, ``undominatedness through the third world'' may also enter (\src{S1\_referee} \S7).}

\begin{conjecture}\label{conj:Q2N}
In $n=2$ on $\M_U$, strict (G), (IB), (LA), (S) and (C) force $\{E_1C:C\in\C_L\}$ to be a chain, i.e.\ $(\alpha_p,\beta_p)$ comonotone along the selection.
\end{conjecture}
\status{CONJECTURE (\src{Q2} CONJECTURE Q2.N; label confirmed, \src{Q12\_referee}). If true, the nested analysis of part~(c) is the whole two-world case. The natural route (a crossing $p<q$ and Lemma~\ref{lem:GSE} applied to the lattice join $E_1C_p\cup E_1C_q$) produces a third selected model but no contradiction; in $\mathfrak L$ this is consistent with rigid crossings (\src{S1} \S3, recorded failure). Its edge form is the question of Proposition~\ref{prop:comonotone}; \src{Q1} states the edge version as CONJECTURE Q1.C2. \emph{Round 13.} It is proved when some world has a leading charge nonzero on $M$ (Theorem~\ref{thm:S22}) and when one half of (MC$_{\mathrm{sel}}$) holds, in particular when one world is monotonely continuous on $M$ (Theorem~\ref{thm:MC}); so it is open only when both worlds lack both (Corollary~\ref{cor:MC}(b)). (C) is necessary for it (Proposition~\ref{prop:R71}). A counterexample needs a rigid or a tied crossing; neither has been constructed, and no claim is made that either can occur. By Proposition~\ref{prop:S13}, if the classes of a rigid crossing are singletons, no nested selection exists at all, so neither ``some nested selection exists'' nor ``every non-nested selection can be straightened'' is then available; ``non-nestedness gives a cover nothing'' is not proved. The reduced form is Conjecture~\ref{conj:S1R}.}

\begin{conjecture}\label{conj:S1R}
In $n=2$ on $\M_U$, every ranking with every axiom satisfies (MC$_{\mathrm{sel}}$)$_1$ or (MC$_{\mathrm{sel}}$)$_2$. It would suffice to show this for $W\in\{C,C'\}$ along the splice chains of a crossing pair.
\end{conjecture}
\status{CONJECTURE (\src{S1} Conj.~S1.R, which asked for both halves of (MC$_{\mathrm{sel}}$); reduced to one half at the crossing pair by \src{S1\_referee} \S9). By Theorem~\ref{thm:MC} it implies Conjecture~\ref{conj:Q2N}. (MC$_{\mathrm{sel}}$) is the model-space counterpart of the closedness of (S) in the index; it is close to (LC$_{T_1}$) of Proposition~\ref{prop:Q38} and to the continuity of Conjecture~\ref{conj:villegas}, and whether (S) forces lower closedness is already open (\src{Q3} \S5). \emph{Where the proof stops} (\src{S1} \S4): in the nested theory Lemma~DL turns a statement about a moving test model into one about moving \emph{selected} models against a fixed test model, and then uses closedness in the index; in a crossing, the moving splices are not comparable with a selected chain in that way. A proof needs a selected chain through the crossing rectangle or a closedness statement for test chains.}

\emph{Edges in three worlds} (round 13; \src{S1} \S5, refereed in \src{S1\_referee} \S\S8--9). Let $n=3$ and consider an edge $e_{ij}$. The \emph{edge family} consists of the edge-selected models together with the edge test models tied with them at worlds $i$ and $j$.

\begin{lemma}\label{lem:weakedge}
Let $p$ lie in the open edge $e_{ij}$ and let $X$ satisfy $X\preceq_iC_p$ and $X\preceq_jC_p$. Then $X\sim_iC_p$ or $X\sim_jC_p$.
\end{lemma}
\status{proved (\src{S1} Lemma~S1.10, from competitiveness on $\{i,j\}$); confirmed (\src{S1\_referee} \S8). So an edge test model can two-world dominate $C_p$ only with a tie at one of the two worlds; undominatedness then says only that $X\not\preceq_kC_p$ at the third world.}

Say that \textbf{(EU$_{ij}$)} holds if no edge test model is weakly better than a member of the edge family at both worlds $i$ and $j$ and strictly better at one. It is not derived from the axioms. It holds when worlds $i$ and $j$ are additive with densities $>0$ on $R_{ij}$ (proved, \src{S1} \S5.3; direct proof in \src{S1\_referee} \S8: a competitive set beaten at one world and tied at the other can be improved at both by adding a small non-null set; NUMERICS \src{s1ref\_py/eu\_zonoid.py}, which also shows that it fails when one density vanishes on a cell).

\begin{proposition}\label{prop:S111}
Assume (EU$_{ij}$). Then on $e_{ij}$: (a) two members of the edge family tied at one of the worlds $i,j$ are tied at both; (b) the double-tie classes of the edge family form a compact linear continuum $\mathfrak L_e$, with dominator intervals as in Corollary~\ref{cor:continuum}; (c) for crossing members $U,V$, with crossing pieces $A,B\subseteq R_{ij}$, every pair $a\subseteq A$, $b\subseteq B$ has pattern (strictly smaller, strictly smaller), (strictly larger, strictly larger) or (tied, tied) at worlds $i,j$.
\end{proposition}
\status{proved, conditional on (EU$_{ij}$) (\src{S1} Prop.~S1.11); confirmed (\src{S1\_referee} \S8), which supplies the unstated step that the selected $C_q$ of Lemma~\ref{lem:level} is itself an edge test model (Lemma~\ref{lem:Q34}), so that (EU$_{ij}$) applies to it. The pattern (tied, tied) is not excluded on edges, because (C) needs ties at all three worlds and edge splices differ at the third. Without (EU$_{ij}$), competitiveness alone leaves nine patterns, including a one-world tie with a strict comparison at the other, and the two-world machinery does not start. Doubly tied edge crossings are the only non-nested possibility not excluded (wording of \src{S1\_referee} \S10, replacing \src{S1}'s ``the expected genuine non-nested possibility''); by Proposition~\ref{prop:RS2A} they cannot occur when worlds $i,j$ are additive and the third world is complete.}

\begin{proposition}\label{prop:S112}
Assume (EU$_{ij}$) and that world $i$ is monotonely continuous on $R_{ij}$. Let $U,V$ be members of the edge family with $[U]\prec[V]$ in $\mathfrak L_e$ that cross, with $E_iU=m\cup B$ and $E_iV=m\cup A$.
\begin{enumerate}[leftmargin=*,label=(\alph*)]
\item Some edge test model $X$ with $E_iX\subseteq E_iV$ is doubly tied (at $i$ and $j$) with $U$.
\item If $U=D_{M_{x_U,y_U}}$ and $V=D_{M_{x_V,y_V}}$ (Proposition~\ref{prop:comonotone}) with $x_U>x_V$ and $y_U<y_V$, then $X$ can be taken to be the coherent edge model $D_{M_{x'',y_V}}$ for some $x''\in[x_V,x_U]$.
\end{enumerate}
\end{proposition}
\status{proved, conditional on (EU$_{ij}$) and monotone continuity at world $i$ (\src{S1} Prop.~S1.12: along a chain $a(s)\uparrow A$, continuous in measure, the switch of $a(s)$ against $B$ is a tie by monotone continuity and completeness on cross pairs, and Proposition~\ref{prop:S111}(c) makes it a double tie); confirmed (\src{S1\_referee} \S8), including coherence of $D_{M_{x'',y_V}}$, which is valid since $x''\le x_U\le y_U<y_V$. The straightened model is nested with $V$ and tied with $U$ (\emph{pairwise straightening}). The mirror, with monotone continuity at world $j$, gives a model tied with $V$ whose error set contains $E_iU$.}

\begin{conjecture}\label{conj:S113}
Assume (EU$_{ij}$) and that world $i$ or world $j$ is monotonely continuous on $R_{ij}$. Then there is a nested edge family with (K) and closed sections in the sense of Remark~\ref{rem:straighten}.
\end{conjecture}
\status{CONJECTURE (\src{S1} Conj.~S1.13, global straightening; label confirmed, \src{S1\_referee} \S9). By Remark~\ref{rem:straighten}, its live payoff, given Theorem~\ref{thm:S21}, is Theorem~\ref{thm:Q36}(c) without (Nest$_e$). \emph{Where the maximal-chain argument stops} (\src{S1} \S5.6): for a maximal inclusion-chain of sets doubly tied with edge-selected models and a missed class $W$, with $c^-$ the essential union of the members below $W$ and $W^-$ their supremum class, one gets $c^-\sim_iW^-$ and $W^-\preceq_jc^-$ (the latter's derivation is completed in \src{S1\_referee} \S9); the reverse inequality $c^-\preceq_jW^-$, a closedness of $\{D:X\preceq_jD\}$ that is not a section, is missing. The two missing lemmas are (EU$_{ij}$) from three-world undominatedness, which Lemma~\ref{lem:weakedge} gives only up to a one-world tie, and that closedness.}

\paragraph{(e) Construction side: every axiom except one closed section, and (round 16) a full system.} \src{Q1} tried to build a world that does not add up inside a full system. Let $n=2$, $K'=[k_1,k_2]\subseteq(0,1)$ (first coordinate of the prior), and let $b\mapsto C_b$, $b\in[0,1]$, be the path of coherent models $D_{[\alpha(t),\beta(t)]}$ with $\alpha(t)=k_1(1-t)$, $\beta(t)=1-t(1-k_2)$, reparametrised so that the world-1 error sets $c_b$ increase strictly with $\bar\mu(c_b)=b$; $C_0=P_1$, $C_1=P_2$, and each interior $C_b$ has a credal set strictly larger than $K'$ on both sides. Selection: $p\mapsto C_{b(p_2)}$ for a continuous increasing bijection $b$. Let $x(h)\coloneqq\inf\{b:h\in c_b\}$ on $M$ (image measure Lebesgue on $[0,1]$) and $x(h)\coloneqq2$ on $N_+$; for a measurable $\delta:U\to\{-1,0,1\}$ put $\tau_\delta(x)\coloneqq\int_{\{x(h)\ge x\}}\delta\,d\bar\mu$. World~2 is additive, $\psi_2(X)\coloneqq\bar\mu(E_2X)$. For $a\in(0,1)$ and a family $\mathcal F$ of non-null measurable subsets of $d_a$ put $g_F\coloneqq1_F-1_{c_a}$, $\mathrm{band}\coloneqq\{\delta:\tau_\delta\ge0\text{ on }[0,2]\}$, $\mathcal G(F)\coloneqq\{\delta:\tau_\delta\ge\tau_{g_F}\text{ on }[0,2]\}$, and let world~1 accept $\delta=1_{E_1Y}-1_{E_1X}$ ($X\preceq_1Y$) iff $\delta\in\mathrm{band}\cup\bigcup_{F\in\mathcal F}\mathcal G(F)$.

\begin{proposition}\label{prop:Q12}
For every family $\mathcal F$: world~1 satisfies strict (G) and (IB); the system satisfies (LA), competitiveness of every $C_b$ on its support, undominatedness of every $C_b$, (C), $\C_L=$ the selection, (LF) at every interior $C_b$, and closed world-2 sections for every model. If $\mathcal F$ contains a pairwise disjoint sequence $A_1,A_2,\dots$, then world~1 does not add up: every charge $\ge0$ endorsing its relevant comparisons vanishes on $c_a$.
\end{proposition}
\status{proved (\src{Q1} Props.~Q1.1, Q1.2); confirmed with repair (\src{Q12\_referee} \S\S1.1--1.2). The key step of Q1.1 (sums of exotic acceptances stay in one $\mathcal G(F)$) needs the nonnegative problems placed to the right, $x=2$ (\src{Q1} Remark~2.2). \emph{Repair.} \src{Q1}'s proof ended ``by Fact~0 every undominated model is tied with a selected one''; Theorem~\ref{thm:kkm} cannot be cited, because it assumes closed sections at every world and Proposition~\ref{prop:Q13} shows that world~1 has a non-closed section. The conclusion survives by a direct argument: for any model $X$, with $A=E_1X\cap M$ and $N$ its unaccepted nonnegative problems, $b_2\coloneqq\min(1,\max(0,\bar\mu(A)-\bar\mu(N)))$ has $C_{b_2}\preceq_2X$ and $C_{b_2}\preceq_1X$ (band), and $C_{b_2}$ has no tie partner; so every model other than $C_{b_2}$ is dominated by it, $\C_L$ is the selection, and (C), (LA) hold (NUMERICS \src{q12ref\_py/q12\_fact0.py}, 400 random models).}

\begin{proposition}\label{prop:Q13}
Let $X$ be regular with $E_1X\cap M=D\subseteq d_a$, and write $G_D(x)\coloneqq\bar\mu(D\cap\{x(h)\ge x\})$. For $F\in\mathcal F$ with $G_D\ge G_F$ on $[a,1]$, $\{b:1_D-1_{c_b}\in\mathcal G(F)\}=[0,\beta_F(D)]$ with $\beta_F(D)=\min\bigl(1,\min_{x\in[a,1]}(x+G_D(x)-G_F(x))\bigr)$; for other $F$ it is empty. The band contributes $[0,\min(\bar\mu(D),\min_{x\ge a}(x+G_D(x)))]$. In particular, with $a=0.3$, the dyadic blocks $J_k=[1-(1-a)2^{-k},1-(1-a)2^{-k-1})$, $A_j$ consisting of one sub-interval of length $2^{-j}|J_k|$ of each $J_k$ (disjoint across $j$), $\mathcal F=\{A_j\}$ and $D=[a,a+d)\cup[1-h,1)$ with $d=h=0.1$ (in the coordinate $x$), the world-1 section of $X$ is
\[
S_1(X)=\{b:C_b\preceq_1X\}=[0,\,a+d+h)=[0,0.5),
\]
which is not closed.
\end{proposition}
\status{proved (\src{Q1} Prop.~Q1.3); confirmed, the formula re-derived including the empty case (\src{Q12\_referee} \S1.3). NUMERICS \src{q12ref\_py/q13\_sections.py} with the actual disjoint dyadic family (exact rationals, 12 blocks): for $j=1,2$ the section from $A_j$ is empty; for $j=3,\dots,8$, $\beta_{A_j}(D)=0.41252,0.45626,0.47813,0.48907,0.49453,0.49727$, formula equal to bisection to $10^{-12}$; the band value is $0.2$; $0.5$ is not in the section (\src{Q1}'s own script used density surrogates). \emph{Diagnosis.} With world~2 additive, (K) forces world~1 to contain the band on cut pairs, and the failing section is the configuration Lemma~T forbids; Proposition~\ref{prop:Q15} turns this into a theorem. If the profiles $\{G_F\}$ form a sup-norm compact set, every section is closed; a spread certificate needs $G_{A_j}\to0$ uniformly, so compactness is exactly what $\neg$SP$^{\mathrm{reg}}$ destroys. So $\neg$SP$^{\mathrm{reg}}$ is compatible with every axiom except closedness of the sections at one world, and is decided by the limit rule (R5) of Theorem~\ref{thm:R5}.}

\emph{Other designs} (\src{Q1} \S\S4--5; failure reports, not theorems except where stated). Fibrewise bands fail (K): the two worlds are not jointly complete on cut pairs (proved). A single deciding functional produces reversals unless it is additive, and a ``ping-pong'' of mirrored certificates was neither realised nor refuted in round 12 (not proved; consistent with Proposition~\ref{prop:N2}); round 13 realises it with every axiom except closed sections at both worlds (Proposition~\ref{prop:S31}). In three worlds, a lexicographic world with a vanishing leading level fails on the edges, as Theorem~\ref{thm:Q36} explains. A related one-parameter family shows why steep weights do not help: for $K>0$ let world~1 accept $\delta$ iff $\int_{\{x(h)\ge y\}}e^{Kx(h)}\delta\,d\bar\mu\ge0$ for every $y$, for some $K$. This union of exponentially tilted first-order dominance cones is a convex cone, contains the nonnegative functions and no nonzero $\delta\le0$, so it is a legitimate cancellative, strictly monotone world; it accepts every cut pair with non-null target, so every endorsing charge lives at $1^-$ ($a^*=1$); but for $Z\subseteq c_{x-\eta}$ the section $S_1(c_x\setminus Z)=[0,x)$ is open, which Lemma~DL excludes (\src{Q12\_referee} \S4, N3; proved, one reader; NUMERICS \src{q12\_tilt.py}). A surviving world needs a weight infinitely heavy on spread thin sets and light on concentrated ones, i.e.\ non-additive in the chain direction. No construction is claimed here. \emph{Round 16:} Theorem~\ref{thm:D21} below is such a construction, and the weight exists there as a $\sigma$-finite measure (Proposition~\ref{prop:D2sigma}).

\emph{Round 13: mirrored pairs from the construction side} (\src{S3}; refereed in \src{S34\_referee} \S\S1--4). Keep the path $b\mapsto C_b$, the selection and the coordinate $x$ of this part. Fix $0<a<b_0<1$ and a family $\mathcal A=\{A_j\}$ of pairwise disjoint non-null subsets of $d_a$ such that $F'_j\coloneqq A_j\cap c_{b_0}$ is non-null (for instance the members of the dyadic family of Proposition~\ref{prop:Q13} that meet $[a,b_0)$). The mirror image of world~1 under $x\mapsto1-x$ is defined as follows: put $x'(h)\coloneqq x(h)$ on $M$ and $x'(h)\coloneqq-1$ on $N_+$, $\eta_\delta(x)\coloneqq\int_{\{x'(h)\le x\}}\delta\,d\bar\mu$ (heads instead of tails), $\mathrm{band}_2\coloneqq\{\delta:\eta_\delta\ge0\text{ on }[-1,1]\}$, $g'_j\coloneqq1_{F'_j}-1_{d_{b_0}}$ and $\mathcal G_2(F'_j)\coloneqq\{\delta:\eta_\delta\ge\eta_{g'_j}\text{ on }[-1,1]\}$.

\begin{proposition}\label{prop:S31}
Let world~1 be the world of Proposition~\ref{prop:Q12} with $\mathcal F=\mathcal A$, and let world~2 accept $\delta=1_{E_2Y}-1_{E_2X}$ ($X\preceq_2Y$) iff $\delta\in\mathrm{band}_2\cup\bigcup_j\mathcal G_2(F'_j)$. Then:
\begin{enumerate}[leftmargin=*,label=(\alph*)]
\item both worlds satisfy strict (G) and (IB); (K) holds for every model; no model other than $C_b$ weakly dominates or ties $C_b$; $\C_L$ is the selection, so the world-1 error sets of $\C_L$ form a chain; and (LA), (C) and (LF) at every interior $C_b$ hold;
\item $c_a\le_1A_j$ and $d_{b_0}\le_2A_j$ for every $j$, both relevant comparisons, through one pairwise disjoint (hence uniformly Ces\`aro-thin) family; so both worlds fail SP$^{\mathrm{reg}}$, in mirror image and with the same targets;
\item closed sections fail at both worlds: the world-1 section of the regular model with $E_1X\cap M=d_a$ is $[0,1)$, and the world-2 section of the regular model with $E_2X\cap M=c_{b_0}$ is $(0,1]$.
\end{enumerate}
\end{proposition}
\begin{proof}[Proof sketch]
(a) World~1 is Proposition~\ref{prop:Q12}'s; world~2 is its image under $x\mapsto1-x$, with $N_+$ on the side that enters every condition, so the key step of \src{Q1} Prop.~Q1.1 transports. (K) comes from the two bands: if band$_1$ rejects $1_{A\cup N}-1_{c_b}$ and band$_2$ rejects $1_{(M\setminus A)\cup N}-1_{d_b}$, then for some $x_1<b<x_2$, $\bar\mu([x_1,b)\setminus A)>\bar\mu(N)+\bar\mu(A\cap d_b)\ge\bar\mu(N)+\bar\mu(A\cap[b,x_2))>2\bar\mu(N)+\bar\mu(c_b\setminus A)\ge\bar\mu([x_1,b)\setminus A)$. For undominatedness, the conditions on $(1,2]$ force $N=\varnothing$, the tail at $b$ forces $A\subseteq c_b$, and then world~2 accepts a nonpositive vector only if it is $0$. The band sections are closed initial and final segments covering $[0,1]$, so they meet at some $b^*$, and every $X\ne C_{b^*}$ is dominated by $C_{b^*}$ (the direct argument of the repair of Proposition~\ref{prop:Q12}). (b) $g_j\in\mathcal G(A_j)$ and $g'_j\in\mathcal G_2(F'_j)$, and $d_{b_0}\le_2F'_j\le_2A_j$ by monotonicity. (c) $c_b\le_1d_a$ via $\mathcal G(A_j)$ iff $b\le1-\bar\mu(A_j)$, and via the band iff $b\le1-a$; since $\bar\mu(A_j)\to0$, the section contains every $b<1$, and not $1$ by strict (G). The world-2 statement is the mirror.
\end{proof}
\status{proved (\src{S3} Prop.~S3.1; NUMERICS \src{s3py/mirror\_pair.py}); confirmed (\src{S34\_referee} \S1), which gives the simpler witness (c) of non-closed sections and a wording repair: \src{S3}'s ``$\C_L$ = the selection (so (Nest))'' should read ``$\C_L$ is the chain'', since (Nest) includes (S). NUMERICS \src{s34ref\_py/s31\_grid.py}, with a non-uniform disjoint family: (K) and undominatedness without failure over 600 models and 95 cuts, and section ends $1-\bar\mu(A_j)$ and $\bar\mu(F'_j)$ as predicted. It realises the ping-pong of \src{Q1} \S4 and Proposition~\ref{prop:N2}, which had been neither realised nor refuted. So every axiom except closed sections, at both worlds, is compatible with a mirrored simple pair; closedness is where any exclusion of (R1) must bite.}

\begin{proposition}\label{prop:S32}
Let world~1 be $\mathrm{band}\cup\bigcup_{F\in\mathcal F}\mathcal G(F)$ as in this part, for any family $\mathcal F$ of non-null subsets of $d_a$, and let world~2 be arbitrary. For $d,h>0$ with $a+d+2h<1$ put $D\coloneqq[a,a+d)\cup[1-h,1)$, let $X_D$ be regular with $E_1X_D\cap M=D$, and put $\mathcal F_D\coloneqq\{F\in\mathcal F:G_F\le G_D\text{ on }[a,1]\}$ and $\kappa_D\coloneqq\inf_{F\in\mathcal F_D}\bar\mu(F)$.
\begin{enumerate}[leftmargin=*,label=(\alph*)]
\item For $F\in\mathcal F_D$, $\beta_F(D)=a+d+h-\bar\mu(F)$ in Proposition~\ref{prop:Q13}; the band contributes $[0,d+h]$.
\item If $\kappa_D<a$, then $S_1(X_D)=[0,a+d+h-\kappa_D]$ if the infimum is attained, and $[0,a+d+h-\kappa_D)$ if it is not.
\item $\mathcal F_D$ contains every $F\in\mathcal F$ with $\bar\mu(F)\le h$.
\item If $\inf_{F\in\mathcal F}\bar\mu(F)=0$, then $S_1(X_D)$ is not closed, for every admissible $(d,h)$. Hence, if every world-1 section is closed, then $\inf_{F\in\mathcal F}\bar\mu(F)>0$, and every $A$ with $Z\le_1A$ satisfies $\bar\mu(A)\ge\min(\bar\mu(Z),\inf_{F\in\mathcal F}\bar\mu(F))$; so (QC$_1$) holds and world~1 has no mirrored simple pair.
\end{enumerate}
\end{proposition}
\status{proved (\src{S3} Prop.~S3.2); confirmed with repair (\src{S34\_referee} \S2). Repairs: the section is also half-open when the infimum is positive but not attained; what (d) proves, and uses, is $\inf_{\mathcal F}\bar\mu(F)>0$, not that a minimum exists; and \src{S3}'s ``closedness $\iff$ attained target measure'' holds for the test models $X_D$ only. In general the condition is \emph{necessary, not sufficient}: with $a=0.3$, $F_0=[0.99,1)$, $F_n=[0.5,0.6+1/n)$ ($n\ge5$) and $D=[0.3,0.8)$, the minimum $\bar\mu(F_0)=0.01$ is attained, yet $S_1(X_D)=[0,0.7)$ is open (the referee's counterexample, proved; NUMERICS \src{s34ref\_py/s3\_sections.py} (3); (1) checks the formula against bisection in twelve random cases). Proposition~\ref{prop:Q13}'s numbers are the instance $\beta_{A_j}=a+d+h-\bar\mu(A_j)$. The section sees $F$ only through $\bar\mu(F)$, so thinness in the transverse (radial) direction does not help; and a thin family never has a sup-norm compact set of profiles ($\|G_F\|_\infty=\bar\mu(F)$, so $0$ is a limit profile), so the compactness criterion of Proposition~\ref{prop:Q13} cannot be met by a thin family; the operative necessary condition is $\inf_{\mathcal F}\bar\mu(F)>0$.}

\begin{proposition}\label{prop:S33}
\begin{enumerate}[leftmargin=*,label=(\alph*)]
\item Under (Nest) with (K), strict (G), (IB) and closed sections at regular models: if $Z$ is non-null and $Z\le_1A$, then $A\not\le_1Q$ for every non-null $Q$ in the Lemma-T$_1$ neighbourhood of $1$ for $Z$, or of any finite set of points of the chain.
\item Assume only (G) and (IB), with a chain $c_b$ as in this part. If $Z\subseteq c_s$ is non-null, $s<1$, and there are sets $A_k$ with $Z\le_1A_k\le_1[1-m_k,1)$ and $m_k\to0$, then $c_b\le_1M\setminus Z$ for every $b<1$; so one application of the limit rule (R5) of Theorem~\ref{thm:R5}, or closedness of this section at $b=1$, gives $Z\le_1\varnothing$, outcome (a) of the exact consistency test.
\item A band accepts $A\le_1[1-\bar\mu(A),1)$ for every $A\subseteq M$ (\emph{right-packing}). So a world containing a band, with closed sections, has no simple certificate whose targets have vanishing measure.
\end{enumerate}
\end{proposition}
\status{proved (\src{S3} Lemma~S3.3, the ``packing lemma''); confirmed (\src{S34\_referee} \S3, which adds the case $b\le s$ in (b)). (a) is transitivity and Lemma~T$_1$. This is the mechanism behind every failure in \src{Q1}, the tilt world, and Propositions~\ref{prop:S31} and~\ref{prop:S32}: any construction of (R1) needs both worlds to refuse packing of their thin targets, towards $1$ at world~1 and towards $0$ at world~2, while staying jointly complete on cut pairs (K).}

\begin{proposition}\label{prop:S34}
Identify $M$ with $[0,1]_x\times[0,1]_y$ by an isomorphism of measure algebras that preserves the chain ($x$ the chain coordinate, $y$ a transverse coordinate). Let world~1 be $\mathrm{fband}_1\cup\bigcup_j(g_j+\mathrm{fband}_1)$, where $\mathrm{fband}_1$ is the set of $\delta$ whose tails in $x$ are $\ge0$ on a.e.\ fibre $\{y=\text{const}\}$, $A_j=[a,1)\times Y_j$ with disjoint non-null $Y_j$, and $g_j=1_{A_j}-1_{c_a}$; let world~2 be the mirror (fibrewise heads, translates by $g'_j$). Then (K) fails: for $b\in(a,b_0)$, disjoint $Y',Y''$ of equal positive measure, and $\epsilon>0$ with $a<b-\epsilon$ and $b+\epsilon<b_0$, the cut pair $P\coloneqq[b-\epsilon,b)\times Y'$, $Q\coloneqq[b,b+\epsilon)\times Y''$ at $b$ is accepted by neither world in its competitive direction.
\end{proposition}
\status{proved for the design stated (\src{S3} Claim~S3.4); confirmed (\src{S34\_referee} \S4; NUMERICS \src{s34ref\_py/s34\_fibre.py}: (K) fails on all 124 cut pairs tested, each of which both \emph{global} bands accept). Fibrewise cones never move mass across fibres, so they never pack (Proposition~\ref{prop:S33} is avoided), but they are not jointly complete on cut pairs. \emph{General form} (proved): in a world $K\cup\bigcup_F(g_F+K)$ with $K$ inside the global band cone, every translate member $\delta$ has $\|\delta^+\|_1\ge\bar\mu(F)$, so only cut pairs with $\bar\mu(Q)<\inf_F\bar\mu(F)$ are out of reach of the translates; if that infimum is positive the cones must be jointly complete on small cut pairs, and if it is $0$ Propositions~\ref{prop:S32} and~\ref{prop:S33} apply. This gives a full proof of the round-12 failure report ``fibrewise bands fail (K)''. Other designs (unanimity over lexicographic charge vectors; directed unions of closed cones) are failure reports, with the CONJECTURE (\src{S3} \S5) that the former satisfy SP$^{\mathrm{reg}}$ under (Nest) with (K).}
\emph{Round 16: a full system in which a world does not add up} (\src{D2}; refereed independently by two referees, \src{D2\_referee} and \src{D2\_referee2}, who re-derived every step; their repairs are incorporated). Keep $n=2$ and $K'=[\alpha',\beta']\subseteq(0,1)$ in the weight of world~2, take the path $s\mapsto C_s$ of Theorem~\ref{thm:C21}, and reparametrise it as in this part, so that $b\mapsto C_b$, $b\in[0,1]$, has world-1 error sets $c_b\coloneqq E_1(C_b)\cap M$ increasing with $\bar\mu(c_b)=b$, $C_0=P_1$, $C_1=P_2$ and $E_2(C_b)\cap M=d_b$. Let $x(h)\coloneqq\inf\{b:h\in c_b\}$ on $M$, so that $c_b=\{x<b\}$ modulo null sets and $x$ carries $\bar\mu|_M$ to Lebesgue measure on $[0,1)$. For bounded measurable $f$ on $U$ and $t\in[0,1]$ put
\[
\tau_f(t)\coloneqq\int_{M\cap\{x\ge t\}}f\,d\bar\mu+\int_{N_+}f\,d\bar\mu,\qquad D(f)\coloneqq\tau_f(1)=\int_{N_+}f\,d\bar\mu
\]
(the tails of this part, with the unanimous problems at the right). Every model $X$ has $E_1X=A\cup N_X$ with $A\coloneqq E_1X\cap M$ and $N_X\coloneqq N_+\setminus X$; $X$ is regular iff $N_X$ is null, and $X$ is determined by $E_1X$.

Fix $a\in(0,1)$ and $h_0=2^{-m_0}$ with $1-h_0>a$. A \emph{design} is a family $\mathcal G$ of pairs $\gamma=(P_\gamma,Q_\gamma)$ of non-null measurable subsets of $M$, with $\gamma_*\coloneqq1_{Q_\gamma}-1_{P_\gamma}$, such that:
\begin{itemize}[leftmargin=*]
\item[\textbf{(d1)}] $P_\gamma\subseteq c_a$ and $Q_\gamma\subseteq\{x\ge1-h_0\}$;
\item[\textbf{(d2)}] $\bar\mu(Q_\gamma\cap Q_{\gamma'}\cap\{x\ge t\})>0$ for all $\gamma,\gamma'\in\mathcal G$ (equal or not) and all $t<1$;
\item[\textbf{(d3)}] for every $\eta>0$ only finitely many $\gamma$ have $\bar\mu(P_\gamma)\ge\eta$;
\item[\textbf{(d4)}] for every $\epsilon>0$ there are finitely many $\gamma\in\mathcal G$ and weights $w_\gamma>0$ with $\sum_\gamma w_\gamma1_{P_\gamma}=1_{c_a}$ and $\sum_\gamma w_\gamma1_{Q_\gamma}\le\epsilon$ everywhere on $M$.
\end{itemize}
Given a design, let $\mathcal K_0$ be the set of bounded measurable $f$ with $f\ge0$ a.e.\ on $N_+$, $\tau_f(t)\ge0$ for all $t\in[0,1]$, and $f\ge0$ a.e.\ on $M\cap\{x\ge b^*\}$ for some $b^*<1$; it is a convex cone, closed under addition. World~1 \emph{accepts} a measurable $\delta:U\to\{-1,0,1\}$, written $\delta\in\mathcal A_1$, by the following rules:
\begin{itemize}[leftmargin=*]
\item[\textbf{(R$^-$)}] if $D(\delta)<0$, $\delta$ is rejected;
\item[\textbf{(R$^+$)}] if $D(\delta)>0$, $\delta$ is accepted iff $\tau_\delta(t)\ge0$ for every $t\in[a,1]$;
\item[\textbf{(R$^0$)}] if $D(\delta)=0$, $\delta$ is accepted iff $\delta\in\mathcal K_0$ or $\delta-\gamma_*\in\mathcal K_0$ for some $\gamma\in\mathcal G$.
\end{itemize}
Put $X\preceq_1Y$ iff $1_{E_1Y}-1_{E_1X}\in\mathcal A_1$, and $X\preceq_2Y$ iff $\bar\mu(E_2X)\le\bar\mu(E_2Y)$. In words: the base cone $\mathcal K_0$ is first-order (tail) dominance in the chain coordinate, except that it never moves mass out of a neighbourhood of $x=1$ (the \emph{protected end}); each generator trades a thin left strip $P_\gamma$ for a thin set $Q_\gamma$ spread over every neighbourhood of $x=1$; and once a model errs on a non-null set of unanimous problems, the initial segment $c_a$ costs world~1 nothing.

\begin{lemma}\label{lem:D22}
Designs exist, for every $a\in(0,1)$ and every admissible $h_0$.
\end{lemma}
\begin{proof}[Proof sketch]
Let $J_m\coloneqq\{1-2^{-m}\le x<1-2^{-m-1}\}$ ($m\ge m_0$), with increasing affine maps $u_m:J_m\to[0,1)$, and let $\iota:[0,1)\to[0,1)^{\mathbb N}$, $v\mapsto(v_1,v_2,\dots)$, read $v_n$ off the binary digits of $v$ in the positions whose index has $2$-adic valuation $n-1$ (terminating expansions); $\iota$ is defined everywhere and measure-preserving. At level $n\ge1$ put $K_n\coloneqq2^n$, choose a prime $q_n\ge4^n$, let $\Pi_n\coloneqq q_n^2+q_n+1$ be the number of points of the projective plane $\mathrm{PG}(2,q_n)$, give the point $\pi$ the cell $[\pi/\Pi_n,(\pi+1)/\Pi_n)$, put $L_n\coloneqq\lfloor\Pi_n/K_n\rfloor$, and assign distinct lines to the index pairs $(i,l)$, $i<K_n$, $l<L_n$. Let $\hat Y_{n,i,l}$ be the set of $v$ whose coordinate $v_n$ lies in a cell of a point of the line of $(i,l)$, $P_{n,i,l}\coloneqq\{ia/K_n\le x<(i+1)a/K_n\}$ and $Q_{n,i,l}\coloneqq\bigcup_{m\ge m_0}(J_m\cap u_m^{-1}(\iota^{-1}(\hat Y_{n,i,l})))$. (d1) holds by construction. (d2): every $\{x\ge t\}$ with $t<1$ contains a block $J_m$, on which the intersection has measure $\bar\mu(J_m)$ times the measure of the intersection of the two sets of digit sequences; at one level two lines meet in exactly one point, and different levels use independent coordinates, so this is positive. (d3): $\bar\mu(P_{n,i,l})=a2^{-n}$, and each level is finite. (d4): weights $1/L_n$ at level $n$ give $\sum1_P/L_n=1_{c_a}$; every $v_n$ lies in one cell, and its point lies on $q_n+1$ lines, so $\sum1_Q/L_n\le(q_n+1)/L_n\le2(q_n+1)2^n/q_n^2\to0$.
\end{proof}
\status{proved (\src{D2} Lemma~D2.2); confirmed by both referees (\src{D2\_referee} \S5, \src{D2\_referee2} \S2.1), the incidence facts and the bound checked numerically for the least primes $q_n\ge4^n$ (\src{d2ref\_py/design\_check.py}, \src{d2ref2\_py/design\_check.py}; 0 failures).}

\begin{lemma}\label{lem:D23}
$\mathcal A_1$ contains every $\delta\ge0$, contains no $\delta\le0$ other than $0$, and contains $\delta_1+\delta_2$ whenever $\delta_1,\delta_2\in\mathcal A_1$ and $\delta_1+\delta_2$ is $\{-1,0,1\}$-valued. Hence (Lemma~\ref{lem:vectorform}) $\preceq_1$ is a preorder on $\M_U$ with strict (G) and (IB). World~2 is additive with a density, so it has strict (G) and (IB) as well.
\end{lemma}
\begin{proof}[Proof sketch]
\emph{Fact (F):} if $D(\delta)=0$ and $\delta\in\mathcal A_1$, then $\tau_\delta\ge0$ on $[a,1]$; if $\delta-\gamma_*\in\mathcal K_0$, this is because $\tau_\delta(t)\ge\tau_{\gamma_*}(t)=\bar\mu(Q_\gamma\cap\{x\ge t\})\ge0$ for $t\ge a$, by (d1). \emph{Near-one fact:} if $\delta-\gamma_*\in\mathcal K_0$, then $\delta\ge1_{Q_\gamma}$ a.e.\ on $M\cap\{x\ge t^*\}$ for some $t^*<1$, since $\gamma_*=1_{Q_\gamma}$ on $\{x\ge1-h_0\}$. For $\delta\ge0$, (R$^+$) or $\delta\in\mathcal K_0$ accepts. For $\delta\le0$, $\delta\ne0$: (R$^-$) rejects if $D(\delta)<0$; if $D(\delta)=0$, then $\delta\in\mathcal K_0$ forces $\int\delta\,d\bar\mu=\tau_\delta(0)\ge0$, so $\delta=0$, and $\delta-\gamma_*\in\mathcal K_0$ is impossible by the near-one fact and (d2) with $\gamma'=\gamma$. For sums, $D$ is additive and $D\ge0$ on $\mathcal A_1$. If $D(\delta_1+\delta_2)>0$, (F) and (R$^+$) give the tail condition for the sum. If it is $0$, then $D(\delta_1)=D(\delta_2)=0$, and $\mathcal K_0+\mathcal K_0\subseteq\mathcal K_0$, $(\gamma_*+\mathcal K_0)+\mathcal K_0\subseteq\gamma_*+\mathcal K_0$; two generator translates never add to a $\{-1,0,1\}$-valued function, because the sum is $\ge2$ on $Q_\gamma\cap Q_{\gamma'}\cap\{x\ge t^*\}$, which is non-null by (d2).
\end{proof}
\status{proved (\src{D2} Lemma~D2.3); confirmed by both referees (\src{D2\_referee} \S\S1--2, \src{D2\_referee2} \S2.2). \src{D2\_referee2} \S1 notes that every pair $\delta_1,\delta_2$ with a $\{-1,0,1\}$-valued sum is realised by three models, so transitivity of $\preceq_1$ is equivalent to the sum closure. Acceptance depends on $\delta$ alone, so non-regular models play no special role in strict (G). NUMERICS: 26\,236 sums of accepted vectors and 2\,945 chains (\src{d2ref\_py}), 4\,758 transitivity triples on an adversarial pool (\src{d2ref2\_py}); 0 failures.}

For a model $X$ and $b\in[0,1]$ let $S_1(X)\coloneqq\{b:C_b\preceq_1X\}$.

\begin{lemma}\label{lem:D24}
Let $X$ be any model.
\begin{enumerate}[leftmargin=*,label=(\roman*)]
\item If $\bar\mu(N_X)>0$, then $S_1(X)=[0,\min\{1,\,a+\bar\mu(N_X)+\bar\mu(A\cap\{x\ge a\})\}]$.
\item If $N_X$ is null, $S_1(X)$ is a closed initial segment of $[0,1]$.
\end{enumerate}
\end{lemma}
\begin{proof}[Proof sketch]
Put $\delta_b\coloneqq1_{A\cup N_X}-1_{c_b}$. (i) $D(\delta_b)=\bar\mu(N_X)>0$, so (R$^+$) applies; for $t\ge b$ the tail is $\ge0$, and for $t<b$ it is $\psi(t)-b$ with $\psi(t)\coloneqq t+\bar\mu(N_X)+\bar\mu(A\cap\{x\ge t\})$, which is continuous and nondecreasing; so the condition on $[a,1]$ holds iff $b\le\psi(a)$. (ii) $D(\delta_b)=0$, so $S_1(X)=B_0\cup\bigcup_\gamma B_\gamma$ with $B_0\coloneqq\{b:\delta_b\in\mathcal K_0\}$ and $B_\gamma\coloneqq\{b:\delta_b-\gamma_*\in\mathcal K_0\}$. By monotonicity of $t\mapsto t+\bar\mu(A\cap\{x\ge t\})$, $B_0=[0,\bar\mu(A)]$ if $\bar\mu(A)<1$ and $B_0=[0,1]$ if $A=M$. For $B_\gamma$: $b=1$ is excluded by (d2); for $b<1$ the near-one clause is a condition on $A$ alone, and the tail clause $(b-t)^+\le u_\gamma(t)\coloneqq\bar\mu(A\cap\{x\ge t\})-\tau_{\gamma_*}(t)$ defines an initial segment. \emph{Bounds} (stated before closedness, as \src{D2\_referee} \S3 requires): the tails at $t=0$ and $t=a$ give $\beta_\gamma\coloneqq\sup B_\gamma\le\min\{\bar\mu(A)+\bar\mu(P_\gamma),\,1-\bar\mu(Q_\gamma)\}<1$. \emph{Closedness:} if $b_k\uparrow\beta$ with $b_k\in B_\gamma$, then $\beta<1$ and the tail clause passes to the limit, so $B_\gamma$ is empty or $[0,\beta_\gamma]$. \emph{Union:} if $A\ne M$ and $s\coloneqq\sup S_1(X)>\bar\mu(A)$, every $\gamma$ with $\beta_\gamma>(s+\bar\mu(A))/2$ has $\bar\mu(P_\gamma)>(s-\bar\mu(A))/2$; by (d3) there are finitely many, so $s$ is attained.
\end{proof}
\status{proved (\src{D2} Lemma~D2.4); confirmed by both referees (\src{D2\_referee} \S3, with the reordering of bounds and closedness incorporated; \src{D2\_referee2} \S2.3). It is (d3), not (d2), that makes the union closed. \src{D2\_referee} \S3 contrasts Proposition~\ref{prop:S32}: on its test model $D=[a,a+d)\cup[1-h,1)$ the generator contributions here \emph{decrease} to the base value, while in Proposition~\ref{prop:S32} they increase to an unattained supremum; thin sources reverse the direction. NUMERICS: sections on cut grids in both referees' cell models, 0 failures.}

\begin{theorem}\label{thm:D21}
Let $\mathcal G$ be a design. The two-world ranking $(\preceq_1,\preceq_2)$ on $\M_U$ with selection $p\mapsto C_{p_2}$ satisfies strict (G), (IB), (LA), (S) (competitiveness on $\supp p$, and closed sections for every model at both worlds), (C), and (LF) at every interior $C_b$, and $\C_L=\{C_b:b\in[0,1]\}$. World~2 is additive, with density $1$ with respect to $\bar\mu$. World~1 does not add up: every bounded (equivalently, real-valued) finitely additive $\nu\ge0$ on $U$ that vanishes on null sets and endorses the relevant comparisons of world~1 vanishes on the non-null set $c_a$. The failure is of type (M2): no finite family of relevant comparisons is an obstruction.
\end{theorem}
\begin{proof}[Proof sketch]
For a model $X$ and a cut $b$ put $P\coloneqq c_b\setminus A$ and $Q\coloneqq A\setminus c_b$. Then $E_2X=(M\setminus A)\cup N_X$, $C_b\preceq_2X$ iff $\bar\mu(Q)\le\bar\mu(P)+\bar\mu(N_X)$, and $C_b\preceq_1X$ iff $1_{Q\cup N_X}-1_P\in\mathcal A_1$. (1) Strict (G) and (IB): Lemma~\ref{lem:D23}. (2) Competitiveness: $C_0=P_1$ and $C_1=P_2$ by (G). For interior $b$ with $C_b\not\preceq_2X$, i.e.\ $\bar\mu(Q)>\bar\mu(P)+\bar\mu(N_X)$: if $\bar\mu(N_X)>0$, (R$^+$) accepts, since every tail below $b$ is $\ge\bar\mu(N_X)+\bar\mu(Q)-\bar\mu(P)>0$; if $N_X$ is null, $1_Q-1_P\in\mathcal K_0$, being $\ge0$ on $\{x\ge b\}$, $0$ on $N_+$, and with tails $\ge\bar\mu(Q)-\bar\mu(P)\ge0$ below $b$. So (K) holds for every model. (3) Closed sections: at world~2, $\{b:\bar\mu(A\cap d_b)\le\bar\mu(c_b\setminus A)+\bar\mu(N_X)\}$, closed by continuity in $b$; at world~1, Lemma~\ref{lem:D24}; and $p\mapsto p_2$ is continuous. (4) No weak domination of $C_b$: suppose $X\preceq_1C_b$ and $X\preceq_2C_b$, and put $\delta'\coloneqq1_P-1_{Q\cup N_X}$. (R$^-$) forces $N_X$ null. For $b=1$, $X\preceq_2C_1=P_2$ forces $E_2X$ null, so $X=C_1$ (the separate treatment both referees require). For $b<1$, $\delta'\le0$ on $\{x\ge b\}$, so $\delta'-\gamma_*\in\mathcal K_0$ is impossible by (d2); $\delta'\in\mathcal K_0$ gives $\tau_{\delta'}(b)=-\bar\mu(Q)\ge0$, so $Q$ is null, and then $X\preceq_2C_b$ gives $\bar\mu(P)\le0$; so $X=C_b$. (5) By (4) each $C_b$ is coherent and undominated, hence in $\C_L$, so (S) and (LA) hold; by Theorem~\ref{thm:kkm} every undominated model is tied at both worlds with some $C_b$, hence equals it; so $\C_L=\{C_b\}$ and (C) holds. (6) (LF) at interior $C_b$: their credal sets contain $K'$ in their relative interior (Theorem~\ref{thm:C21}(4)). (7) For $\gamma\in\mathcal G$ let $X_\gamma$ be the regular model with $E_1X_\gamma=(c_a\setminus P_\gamma)\cup Q_\gamma$. Then $1_{E_1X_\gamma}-1_{c_a}=\gamma_*$ and $\gamma_*-\gamma_*=0\in\mathcal K_0$, so $C_a\preceq_1X_\gamma$ is a relevant comparison. If a bounded $\nu\ge0$ endorses these, $\nu(P_\gamma)\le\nu(Q_\gamma)$ for all $\gamma$, and for $\epsilon>0$, with (d4), $\nu(c_a)=\sum w_\gamma\nu(P_\gamma)\le\sum w_\gamma\nu(Q_\gamma)\le\epsilon\,\nu(M)$; so $\nu(c_a)=0$. (8) Type (M2): if $f=\sum_k\lambda_k\delta_k\le0$ with $\lambda_k>0$ and relevant $\delta_k$, each $\delta_k$ has $D(\delta_k)=0$ and is, near $x=1$, either $\ge0$ or $\ge1_{Q_{\gamma_k}}$; $f\le0$ excludes the generator terms by (d2), so $f\in\mathcal K_0$, $\int f\,d\bar\mu\ge0$ and $f=0$.
\end{proof}
\status{proved (\src{D2} Thm.~D2.1); \textbf{confirmed by two independent referees} (\src{D2\_referee}, \src{D2\_referee2}), each of whom re-derived every step and found no error affecting the theorem. Repairs incorporated: \emph{bounded} weights are stated explicitly (step~7 needs $\nu(M)<\infty$); step~4 treats $b=1$ separately; Lemma~\ref{lem:D24}(ii) gives the bounds before closedness. \emph{Why boundedness matters} (\src{D2\_referee} \S5): the extended-valued weight $\nu(S)\coloneqq\bar\mu(S)$ if $S\cap\{x\ge t\}$ is null for some $t<1$, and $\nu(S)\coloneqq\infty$ otherwise, is finitely additive, zero exactly on null sets, and endorses every relevant comparison of world~1; the corpus convention is bounded charges (Definition~\ref{def:addup} and the remark after it, \src{N19} Lemma~R1, the probability charges of $\mathcal V_1$), so the theorem stands as stated. Step~8 is the direct argument of \src{D2\_referee} \S5 and \src{D2\_referee2} \S3.2; it replaces \src{D2}'s appeal to \src{N17} N17.3. \src{D2\_referee2} \S2.5 also proves step~5 without Theorem~\ref{thm:kkm}, by exhibiting for every $X\ne C_b$ a selected model that strictly dominates it. \emph{Why Proposition~\ref{prop:S32}(d) does not apply} (\src{D2\_referee2} \S2.7): there $\inf_\gamma\bar\mu(Q_\gamma)=0$ would make a section open, but Proposition~\ref{prop:S32} concerns worlds $\mathrm{band}\cup\bigcup\mathcal G(F)$ in which each target $F\subseteq d_a$ is usable against every source; here generator $\gamma$ is usable only by sets that contain its own source $P_\gamma$ (the tail at $t=0$ and the near-one clause tie it to $P_\gamma$), and the bound $\beta_\gamma\le\bar\mu(A)+\bar\mu(P_\gamma)$ of Lemma~\ref{lem:D24} replaces the hypothesis of Proposition~\ref{prop:S32}(d). \emph{Consistency} (both referees, re-derived in this system): $\varphi_1(P)\ge\bar\mu(P)$, (SR$_1$) and (QC$_1$) hold, Lemma~T$_1$ holds with a radius depending only on $\bar\mu(Z)$, there is no simple, vanishing-target or fixed-source star certificate (as Proposition~\ref{prop:Q15} requires next to an additive world), and the certificate is a ratio certificate with thin sources whose targets meet in every neighbourhood of $x=1$; in the notation of Corollary~\ref{cor:IS}, $a^*=a$ exactly, since $\bar\mu|_{d_a}$ endorses every relevant comparison (\src{D2\_referee2} \S2.4). No conflict with any proved statement of this subsection was found. \emph{Not claimed:} \src{D2}'s remarks ``why (R$^+$) is needed'', ``why the protected end'' and its \S0 account of ``what made it work'' are heuristic readings (both referees) and are not used. NUMERICS: \src{d2py/d2\_check.py} (reproduced byte for byte by \src{D2\_referee}) and both referees' independent cell models, 0 failures. In \src{D2}'s script the sum-closure test ran only once, and at cuts inside the last block the finite proxy for the near-one clause is stricter than the clause itself, so (K) fails there in the proxy only (\src{D2\_referee} \S\S4, 7; \src{D2\_referee2} \S4); no verdict rests on the scripts.}

\emph{Consequences} (proved, from Theorem~\ref{thm:D21}; both referees). (1) Conjecture N17.B$'$ (\src{N19}) is false, already for $n=2$ with a nested selection, so also in the form N17.B$'$(2) of \src{P4} \S5.5 and \src{P4B\_referee} \S8.5: every hypothesis holds and every bounded endorsing charge of world~1 vanishes on $c_a\subseteq M$. (2) \src{Q1}'s Conjecture Q1.C1 is false: world~2 has the strictly positive, uniformly absolutely continuous endorsing charge $\bar\mu$, and world~1 fails SP$^{\mathrm{reg}}_1$. (3) Question (R2) of part~(g) has a negative answer: world~2 has no simple certificate (it is additive), world~1 has none (Proposition~\ref{prop:Q15}), and world~1 fails SP$^{\mathrm{reg}}$. (4) The failure is of type (M2). \emph{Scope.} This is a two-world result and not a cover of $T_1$. What it shows for $T_1$ is that the exclusion route ``prove N17.B$'$, then apply Theorem~\ref{thm:M0closed}'' cannot work as stated: a world that does not add up is no longer hypothetical. \src{D2} \S5 reads this as reopening the positive route for $T_1$; that is a reading, not a theorem (\src{D2\_referee} \S6), and the installation of the mechanism in three worlds is not attempted there (\src{D2}'s Conjecture~D2.C: the world of Theorem~\ref{thm:D21} can be installed on one edge region $R_{ij}$ of a three-world full system on $\M_U$). Part~(f) below shows, in note \src{E1} (audited in \src{F2}, audit refereed once), that the protected part of the mechanism is invisible to the argument of Theorem~\ref{thm:M0closed} at ends on which $T_1$ is eventually correct, and it excludes covers in which every world is an end-protected decoration with integrable, jointly end-heavy base densities (Theorem~\ref{thm:E15}, Proposition~\ref{prop:E16}); other installations are not excluded (part~(g)).

\emph{Scoping: the failure is exactly non-integrability at the protected end} (\src{D2\_referee2} \S3). Write $\sigma_i(h)=\operatorname{sign}h_i$ as in the Notation of this subsection.

\begin{proposition}\label{prop:D2sigma}
In the system of Theorem~\ref{thm:D21}, built on the design constructed in the proof of Lemma~\ref{lem:D22}, let $\rho(h)\coloneqq1/(1-x(h))$ on $M$ and $\rho\coloneqq1$ on $N_+$, and put $\nu\coloneqq\rho\bar\mu$. Then $\nu$ is countably additive and $\sigma$-finite, positive exactly on non-null sets, finite on every $c_b$ with $b<1$, and it endorses every comparison $X\preceq_1Y$ between regular models, $\nu(E_1X)\le\nu(E_1Y)$ in $[0,\infty]$; in particular it endorses every relevant comparison of world~1. Moreover, for interior $q$, the set $\{g\in U:q_1\rho(g)\sigma_1(g)+q_2\sigma_2(g)>0\}$ accepts $N_+$ and has world-1 error set $\{h\in M:\rho(h)<q_2/q_1\}$ on $M$ modulo null sets, an initial segment in $x$; so these Bayes sets are members of the path $C_b$.
\end{proposition}
\begin{proof}[Proof sketch]
Let $\delta=1_{E_1Y}-1_{E_1X}$ be accepted, $X,Y$ regular, so $D(\delta)=0$ and $\delta=0$ on $N_+$. If $\delta\in\mathcal K_0$, then $\delta\ge0$ on $\{x\ge b^*\}$, so $\int\delta^-\rho<\infty$, and Tonelli applied to $\delta^+$ and $\delta^-$ gives $\int\delta\rho\,d\bar\mu=\rho(0)\tau_\delta(0)+\int_0^1\rho'(s)\tau_\delta(s)\,ds\ge0$, as $\tau_\delta\ge0$ and $\rho'\ge0$. If $\delta-\gamma_*\in\mathcal K_0$, then $\delta\ge1_{Q_\gamma}$ on $\{x\ge t^*\}$, so $\int\delta^+\rho\ge\int_{Q_\gamma\cap\{x\ge t^*\}}\rho=\infty$ (on $J_m$, $\rho\ge2^m$ and $\bar\mu(Q_\gamma\cap J_m)$ is a fixed fraction of $2^{-m-1}$, so every block contributes at least a fixed amount), while $\int\delta^-\rho<\infty$.
\end{proof}
\status{proved in a referee report (\src{D2\_referee2} \S3.1), one reader; NUMERICS (\src{d2ref2\_py/cell\_model.py}: 5\,672 accepted relevant comparisons endorsed, 0 failures). The proof uses only that $\rho>0$ is nondecreasing in $x$ with $\int_{Q_\gamma}\rho\,d\bar\mu=\infty$ for every $\gamma$; any such $\rho$ will do (for instance $2^m$ on $J_m$), for any design for which one exists. \emph{Reading} (\src{D2\_referee2} \S3): world~1 is a sub-relation of an additive order whose strictly positive density is integrable on every $c_b$, $b<1$, and fails to be integrable only at the protected end $\{x\to1\}$; so the failure of SP$^{\mathrm{reg}}$ in Theorem~\ref{thm:D21} is \emph{exactly} a failure of integrability of the endorsing density at that end. That end consists of two rays (Lemma~\ref{lem:F11}): $\R_+(1,0)$, approached from $M^+$, which lies in the interior of $U$, and the boundary ray $\R_+(-\alpha',1-\alpha')$ of $U$, approached from $M^-$. So $\rho$ is integrable on every $c_b$ with $b<1$, hence on every compact subset of $M$, but not on any compact neighbourhood in $U$ of a point of $\R_+(1,0)$ (correction of 15 October: the reading in round 16 named only the boundary ray; \src{F1}, \src{F12\_referee} \S1). In the phrase of the tilt paragraph above, this world has ``a weight infinitely heavy on spread thin sets'', and here that weight is a genuine $\sigma$-finite measure. The refutations above are unaffected: N17.B$'$, Q1.C1 and (R2) concern bounded charges. The same reasoning gives a second proof of type (M2): every $\delta_k^-$ lies in some $c_{b_k}$ with $b_k<1$, so $\int(\sum c_k\delta_k)\rho=\sum c_k\int\delta_k\rho\ge0$ (\src{D2\_referee2} \S3.2).}

The following nearby statement was proposed in round 16 as the $\sigma$-finite form of N17.B$'$. \textbf{It is false} (Theorem~\ref{thm:F13} below); it is kept, marked as refuted, because the notes refer to it.

\begin{conjecture}\label{conj:sigmaB}
Let $n\ge2$ on $\M_U$, with strict (G), (IB), (LA), (S) and (C). Then every world $i$ has a countably additive measure $\rho_i\mu$ that endorses its relevant comparisons, with $\rho_i>0$ a.e.\ and $\rho_i$ locally integrable on $U$.
\end{conjecture}
\status{\textbf{REFUTED (15 October 2026).} False as written, already for $n=2$: the system of Theorem~\ref{thm:D21} is a counterexample under the total-cost and the cancellation readings of endorsement (Theorem~\ref{thm:F13} below). Observed by an external auditor; proved in note \src{F1} (Theorems~F1.2, F1.3), and by a second route, specific to Theorem~\ref{thm:D21}, in note \src{F2} (Prop.~F2.4); both confirmed in \src{F12\_referee} \S\S1--2, one reader. Originally a CONJECTURE, suggested by \src{D2\_referee2} \S3.3 (as ``N17.B$'_\sigma$''), the $\sigma$-finite form of N17.B$'$. \emph{Correction.} The round-16 status said that Theorem~\ref{thm:D21} does not refute it, because the density of Proposition~\ref{prop:D2sigma} ``blows up only as $x\to1$, that is, near the boundary ray of $U$ on the side of $e_2$, and is bounded on every closed subcone of $U$''. That sentence was false, and \src{D2\_referee2} \S3.3 contains the same slip. The end $\{x\to1\}$ also contains the ray $\R_+(1,0)$, which lies in the interior of $U$, so $1/(1-x)$ is not integrable on any compact neighbourhood in $U$ of a point of that ray (Lemma~\ref{lem:F11}), and by Theorem~\ref{thm:F13} no endorsing measure does better. Whether the conclusion of the conjecture holds is decided by the selection, not by where the non-additive trades sit: it holds in the system $\mathrm{D2}'_{\mathrm B}$ of Theorem~\ref{thm:F16}, whose chain ends only at $\partial U$ (Theorem~\ref{thm:F18}(2)). \emph{Repairs.} Conjecture~\ref{conj:F1Cp} below replaces ``locally integrable on $U$'' by ``locally integrable on $M$''; the weaker a.e.-finite target is stated after it. \emph{For $T_1$} (\src{D2\_referee2} \S3.3; now to be read for Conjecture~\ref{conj:F1Cp}): the closing argument of Theorem~\ref{thm:M0closed}(a) uses no integrability (its status), so for $\sigma$-finite supporting measures the question is whether Proposition~\ref{prop:M0} ($\alpha$)--($\gamma$), which use the Lyapunov range and finite Bayes risks, survive, for instance by localising to models that differ only where $\rho_i$ is integrable, such as a neighbourhood of $\partial D_{T_1}$ inside $U$; not attempted. Since $\partial D_{T_1}\setminus0\subseteq M$ (CHECKED, \src{F1} \S2.4), such a localisation needs only Conjecture~\ref{conj:F1Cp}, subject to the caveat in its status. \src{D2}'s Conjecture~D2.C, carried out with the same design, would again give a world that is $\sigma$-finitely additive near the end of its edge chain, so it would refute the bounded form of N17.B$'$ in three worlds. Theorem~\ref{thm:E11} below bears on the same question from the side of Theorem~\ref{thm:M0closed}.}

\emph{Audit of 15 October: Conjecture~\ref{conj:sigmaB} is false, and nearby systems} (note \src{F1}, with \src{F2} \S2.3; refereed once in \src{F12\_referee}, whose repairs are incorporated). The refutation was observed by an external auditor and proved in note \src{F1}. Keep the setting of Theorem~\ref{thm:D21}: $n=2$, $K'=[\alpha',\beta']$ in the weight of world~2, $M^\pm$, $\tau$ and $\sigma$ as in part~(d), and $\bar\mu$ as in this part. Two paths of coherent models are used. \emph{Host~A} is the path of Theorem~\ref{thm:C21}, $K_s=[\alpha's,\beta'+(1-\beta')s]$, which is the host of Theorem~\ref{thm:D21}. \emph{Host~B} is the two-phase path $K_s=[0,\beta'+2(1-\beta')s]$ for $s\le\frac12$ and $K_s=[\alpha'(2s-1),1]$ for $s\ge\frac12$. Each is reparametrised as in this part, which gives $C_b$, $c_b$, $d_b$ and the uniform chain coordinate $x$.

\emph{Readings of endorsement for measures of infinite mass} (\src{F1} \S2.1). Let $\nu$ be a countably additive measure on $U$ with $\nu\ll\mu$ and $\nu(Z)>0$ for every non-null $Z$, and let $B\preceq_iX$ be a relevant comparison. Under the \emph{total-cost} reading, $\nu$ endorses it if $\nu(E_iB)\le\nu(E_iX)$ in $[0,\infty]$; this is the reading of Definition~\ref{def:addup}(b), of Proposition~\ref{prop:D2sigma} and of the shadow notation. Under the \emph{cancellation} reading, $\nu$ endorses it if $\nu(E_iB\setminus E_iX)\le\nu(E_iX\setminus E_iB)$ in $[0,\infty]$. A reading has \emph{property (R)} if, whenever $\nu(E_iX)<\infty$, endorsement implies $\nu(E_iB)\le\nu(E_iX)$. Both readings have (R). The reading that constrains only comparisons whose two costs are both finite would make Conjecture~\ref{conj:sigmaB} true in every system, through the density $|h|^{-2}$, under which every non-null cone has infinite mass; so it is not the conjecture's reading (proved in \src{F1}, confirmed).

\begin{lemma}\label{lem:F11}
On host~A, $x(h)\to1$ if and only if $h$ approaches the ray $\R_+(1,0)$ from within $M^+$ or the ray $\R_+(-\alpha',1-\alpha')$ from within $M^-$. The first ray lies in the interior of $U$, since $P_{K'}(1,0)=1-\beta'>0$; the second lies on $\partial U$, since $P_{K'}(-\alpha',1-\alpha')=0$. For $0<r_1<r_2$ and small $\eta>0$ the set $K_\eta\coloneqq\{h:r_1\le|h|\le r_2,\ |\arg h|\le\eta\}$ is a compact subset of $U$, and there are $\kappa>0$ and $x_0<1$ with $\int_{K_\eta\cap M}g(x)\,d\bar\mu\ge\kappa\int_{x_0}^1g(s)\,ds$ for every Borel $g\ge0$ on $[0,1)$; the same holds near the second ray within $M^-$. In particular $1/(1-x)$ is not integrable on $K_\eta$. On host~B, $x(h)\to1$ only as $h$ approaches $\R_+(-\alpha',1-\alpha')$, and every compact $K\subset U$ has $K\cap M\subseteq c_{b_0}$ for some $b_0<1$.
\end{lemma}
\begin{proof}[Proof sketch]
For $h\in M^-$, $h\in D_{[l,r]}$ iff $\sigma(h)<l$, and world~1 errs there by accepting; for $h\in M^+$, $h\notin D_{[l,r]}$ iff $\tau(h)\le r$, and world~1 errs there by rejecting. So $E_1(D_{[l,r]})\cap M=\{h\in M^-:\sigma(h)<l\}\cup\{h\in M^+:\tau(h)\le r\}$. On host~A this is $\{x_A\le s\}$ with $x_A\coloneqq(\tau-\beta')/(1-\beta')$ on $M^+$ and $x_A\coloneqq\sigma/\alpha'$ on $M^-$, and $x$ is an increasing $C^1$ function of $x_A$; so $x\to1$ iff $\tau\to1$ or $\sigma\to\alpha'$. Under $\bar\mu$ the angle of $h$ has a density bounded above and below on $M$, independent of $|h|$, and $\tau$, $\sigma$ have derivatives in $[\frac12,1]$ as functions of the angle on the closed arcs; changing variables from the angle to $x$ gives the integral bound. On host~B the points of $M^+$ have $x<\bar\mu(M^+)<1$, and on a compact $K\subset U$, $\sigma\le\sigma_0<\alpha'$.
\end{proof}
\status{proved (\src{F1} Lemma~F1.1); confirmed (\src{F12\_referee} \S1), one reader, by brute force from the parametrisation and analytically: for $\alpha'=0.3$, $\beta'=0.6$, $(1-x)/|\arg h|\to2.3095$ at $\R_+(1,0)$, and that ray carries a share $0.4361$ of the levels $x\to1$. NUMERICS \src{f1py/geometry.py}, \src{f2py/f2\_d2\_geometry.py}, \src{f12ref\_py/ref\_geometry.py}. The targets $Q_\gamma$ of Lemma~\ref{lem:D22} are defined through $x$ alone, so they occupy both components of every block near $x=1$.}

\begin{theorem}\label{thm:F12}
Let $n=2$ on $\M_U$ with strict (G), (IB), (LA), (S) and (C), and let world~2 be additive with an integrable density $w>0$; put $\nu_2\coloneqq w\mu$. Then the selection is nested (Corollary~\ref{cor:MC}), with chain $c_b$, coordinate $x$ and (K) at interior cuts. Let $\nu$ be a countably additive measure on $U$, with $\nu\ll\mu$, positive on non-null sets and $\sigma$-finite on $M$, that endorses every relevant world-1 comparison under a reading with property (R), and put $\rho\coloneqq d\nu/d\nu_2$ on $M$. Then:
\begin{enumerate}[leftmargin=*,label=(\alph*)]
\item $\nu(c_b)<\infty$ for every $b<1$;
\item $\rho=\varphi\circ x$ a.e.\ on $M$ for a nondecreasing $\varphi:[0,1)\to[0,\infty)$ that is finite everywhere and positive on $(0,1)$;
\item if no bounded countably additive measure that is positive on non-null sets endorses the relevant world-1 comparisons, then $\int_{d_t}\rho\,d\nu_2=\infty$ for every $t<1$.
\end{enumerate}
\end{theorem}
\begin{proof}[Proof sketch]
For $P\subseteq c_b$ and $Q\subseteq d_b$, $b\in(0,1)$, let $X$ be the regular model with $E_1X=(c_b\setminus P)\cup Q$. If $\nu_2(Q)>\nu_2(P)$, then $C_b\not\preceq_2X$, so (K) gives $C_b\preceq_1X$, a relevant comparison. (a) For large $n$, take $P\coloneqq c_b\cap\{\rho>n\}$ and $Q\subseteq d_b\cap\{\rho\le n\}$ with $\nu_2(Q)>\nu_2(P)$; then $\nu(E_1X)<\infty$, and (R) gives $\nu(c_b)\le\nu(E_1X)$. (b) If $c_b\cap\{\rho>c_2\}$ and $d_b\cap\{\rho<c_1\}$ were both non-null for some $c_1<c_2$, choose non-null $P$, $Q$ in them with $\nu_2(Q)=(1+\epsilon)\nu_2(P)$ and $c_1(1+\epsilon)<c_2$; by (a), (R) and subtraction of finite quantities, $\nu(P)\le\nu(Q)$, which contradicts $\nu(P)>c_2\nu_2(P)>c_1\nu_2(Q)>\nu(Q)$. So $\operatorname{ess\,sup}_{c_b}\rho\le\operatorname{ess\,inf}_{d_b}\rho$ for every $b$, and $\rho$ is a.e.\ a nondecreasing function of $x$. (c) If $\int_{d_t}\rho\,d\nu_2<\infty$, then $\nu|_M+\nu_2|_{N_+}$ is bounded, and it endorses every relevant comparison, since relevant comparisons see only $M$.
\end{proof}
\status{proved (\src{F1} Thm.~F1.2); confirmed (\src{F12\_referee} \S2.2), one reader. (K) at interior cuts is \src{P4B\_referee} Lemma~N0. Only (K) and the reading are used: not transitivity, the generators, or where the targets lie. The comparisons trade mass across a cut between the two components of a level set of $x$; this carries a blow-up at one ray of the end over to the other.}

\begin{theorem}\label{thm:F13}
In the system of Theorem~\ref{thm:D21}, with any design, let $\nu$ be a countably additive measure on $U$, with $\nu\ll\mu$ and positive on non-null sets, that endorses the relevant world-1 comparisons under a reading with property (R).
\begin{enumerate}[leftmargin=*,label=(\arabic*)]
\item If $\nu$ is $\sigma$-finite on $M$, then $\nu|_M=\varphi(x)\bar\mu$ with $\varphi$ as in Theorem~\ref{thm:F12}(b) and $\int_t^1\varphi=\infty$ for every $t<1$. Hence $\nu(K)=\infty$ for every compact $K\subset U$ that has a point of $\R_+(1,0)$ in its interior, and $\nu(V\cap M^-)=\infty$ for every neighbourhood $V$ of a point of $\R_+(-\alpha',1-\alpha')$.
\item No such $\nu$ is finite on every compact subset of $U$.
\item Such $\sigma$-finite measures exist (Proposition~\ref{prop:D2sigma}), and all of them are finite on every compact subset of $M$.
\item Such measures that are not $\sigma$-finite also exist: for $t\in(0,1)$, $\nu_t\coloneqq\bar\mu|_{c_t}+\infty\cdot\bar\mu|_{d_t}$ on $M$, and $\bar\mu$ on $N_+$.
\end{enumerate}
In particular, Conjecture~\ref{conj:sigmaB} is false for $n=2$, under the total-cost and the cancellation readings.
\end{theorem}
\begin{proof}[Proof sketch]
(1) Theorem~\ref{thm:F12}(b) applies with $w=1/\mu(M)$. Step~7 of the proof of Theorem~\ref{thm:D21} excludes every bounded endorsing weight that is positive on $c_a$, so Theorem~\ref{thm:F12}(c) gives $\int^1\varphi=\infty$, and Lemma~\ref{lem:F11} with $g=\varphi$ gives the infinite masses. (2) A measure that is finite on compact subsets of the $\sigma$-compact set $U$ is $\sigma$-finite, so (1) applies. (3) By Theorem~\ref{thm:F12}(a), $\nu(c_b)<\infty$ for $b<1$, and a compact $K\subset M$ has $x(K)\subseteq[0,b_0]$ for some $b_0<1$. (4) Let $\delta$ be the vector of a relevant comparison. If $\delta^+$ meets $d_t$ in a non-null set, the cost of the right-hand model is $\infty$. Otherwise $\delta^+\subseteq c_t$; a generator branch is impossible, since it would make $\delta\ge1_{Q_\gamma}$ near $x=1$; so $\delta\in\mathcal K_0$, and $\bar\mu(\delta^-)\le\bar\mu(\delta^+)$ on $c_t$, where $\nu_t=\bar\mu$. For the conjecture: the measure $\rho_1\mu$ it asks for at world~1, with $\rho_1$ locally integrable on $U$, is finite on compact subsets of $U$, which (2) excludes.
\end{proof}
\status{proved (\src{F1} Thm.~F1.3; the failure was observed by an external auditor); confirmed (\src{F12\_referee} \S2.3, which re-derives (4)), one reader. A second proof of the refutation, specific to Theorem~\ref{thm:D21} and under the total-cost reading for $\nu=g\bar\mu$ with $0<g<\infty$ a.e., is \src{F2} Prop.~F2.4, also confirmed (\src{F12\_referee} \S2.4); neither proof uses the other. NUMERICS \src{f1py/geometry.py}: the integral of $1/(1-x)$ over $K_{0.05}$ truncated at $1-10^{-k}$ grows like $\log$, and stays bounded on host~B. Next to world~2 $=\bar\mu$ on host~A, an endorsing measure that is finite on one compact neighbourhood of a point of $\R_+(1,0)$ has $\int^1\varphi<\infty$, which gives a bounded endorser; so there the $\sigma$-finite and the bounded forms coincide. The statement of Proposition~\ref{prop:D2sigma} is correct; only its reading needed the second ray.}

\begin{theorem}\label{thm:F14}
In the setting of Theorem~\ref{thm:F12}, let $b_0\in(0,1)$ and let $S\subseteq c_{b_0}$ be non-null. There is $\theta=\theta(S,b_0)>0$ such that, for every finite family of world-1 comparisons $P_k\le_1Q_k$ (that is, $1_{Q_k}-1_{P_k}$ is accepted by world~1 in the vector form of Lemma~\ref{lem:vectorform}) with $P_k\cap Q_k=\varnothing$ and $P_k,Q_k\subseteq c_{b_0}$, and every family of weights $w_k>0$ with $\sum_kw_k1_{P_k}\ge1_S$,
\[
\sum_kw_k\,\nu_2(Q_k)\ge\theta.
\]
So no certificate of the kind used in Theorem~\ref{thm:D21} (with $\sum_kw_k1_{P_k}=1_S$ and $\sum_kw_k1_{Q_k}\le\epsilon$, for every $\epsilon>0$) has all its sources and targets in $c_{b_0}$.
\end{theorem}
\begin{proof}[Proof sketch]
Closure of world~1's acceptances under $\{-1,0,1\}$-valued sums (Lemma~\ref{lem:vectorform}) and (G) allow shrinking sources, chaining, and adding comparisons with disjoint supports; (K) at $b_0$ gives $Q\le_1R$ whenever $Q\subseteq c_{b_0}$, $R\subseteq d_{b_0}$ and $\nu_2(R)>\nu_2(Q)$. Put $\Gamma\coloneqq\sum_kw_k\nu_2(Q_k)$ and cap each $w_k$ at $1$. Keeping each $k$ independently with probability $w_k$ yields a subfamily that covers a part of $S$ of $\bar\mu$-mass at least $(1-e^{-1})\bar\mu(S)/2$ at cost at most $4\Gamma$. Make its sources disjoint and move its targets to disjoint sets $R_i\subseteq d_{b_0}$; this gives $Z\le_1R$ with $\bar\mu(Z)$ bounded below and $\nu_2(R)\le4\Gamma+\eta$. Corollary~\ref{cor:SR} gives $s_m>0$ with $c_{s_m}\le_1Z$, and Proposition~\ref{prop:MPsharp}(a), (b) give $t^*<1$ with $d_{t^*}\le_2R$, so $\nu_2(R)\ge\nu_2(d_{t^*})>0$.
\end{proof}
\status{proved (\src{F1} Thm.~F1.4; inputs \src{P4B\_referee} Lemma~N0, Lemma~\ref{lem:vectorform}, Corollary~\ref{cor:SR}, Proposition~\ref{prop:MPsharp}(a), (b)); confirmed (\src{F12\_referee} \S3, including the rounding constants), one reader. NUMERICS \src{f1py/rounding.py}. It uses Corollary~\ref{cor:SR} at world~2, that is, that world~2 has a strictly positive, uniformly absolutely continuous endorsing charge. With both worlds non-additive the question is open.}

\begin{corollary}\label{cor:F15}
In the setting of Theorem~\ref{thm:F14}, let $K\subset U$ be compact with $h_2\le-c|h|$ on $K\cap M^+$ for some $c>0$. Then $K\cap M\subseteq c_{b_0}$ for some $b_0<1$. So no certificate of the kind used in Theorem~\ref{thm:D21}, with sources in some $c_a$, has all its targets in $K$. A certificate whose targets lie in one compact subset of $U$ must therefore, for every $c>0$, use targets with $h_2>-c|h|$, that is, targets accumulating at the ray $\R_+(1,0)$.
\end{corollary}
\begin{proof}[Proof sketch]
By Proposition~\ref{prop:lattice} the selected models are $D_{[l(b),r(b)]}$ with $l$, $r$ nondecreasing, so $c_b=\{h\in M^-:\sigma<l(b)\}\cup\{h\in M^+:\tau\le r(b)\}$. Since $\bar\mu(d_b)\to0$, $r(b)\to1$ and $l(b)\to\alpha'$. On $K\cap M^+$, $\tau\le1/(1+c)$; on $K\cap M^-$, $\sigma\le\sigma_0<\alpha'$ by compactness. Take $b_0\ge a$ with $r(b_0)\ge1/(1+c)$ and $l(b_0)>\sigma_0$, and apply Theorem~\ref{thm:F14}.
\end{proof}
\status{proved (\src{F1} Cor.~F1.5); confirmed (\src{F12\_referee} \S3), one reader. \emph{The last sentence is a necessary condition only} (\src{F12\_referee} \S3): on host~B, or on any chain with $r(b)=1$ for some $b<1$, every compact $K\subset U$ has $K\cap M$ inside some $c_{b_0}$, and no certificate with targets in a compact subset of $U$ exists at all. On host~A such targets can be realised near $\R_+(1,0)$, with payoffs not bounded away from zero (\src{F1} Remark~F1.9; the finiteness step of its proof cited Theorem~F1.2(a), which needs $\sigma$-finiteness, and \src{F12\_referee} \S4.6 replaces it by a stepwise argument from (K) and (R)). So, next to an additive world, the three requirements once proposed for a modified system (every target in one compact subset of $U$, payoffs bounded away from zero, and a certificate) are inconsistent.}

\emph{The systems $\mathrm{D2}'_{\mathrm A}$ and $\mathrm{D2}'_{\mathrm B}$} (\src{F1} \S4). Fix a host $H\in\{\mathrm A,\mathrm B\}$. On $M^+$ let $z(h)$ be the $\bar\mu$-proportion of $M^+$ with smaller $\tau$, and on $M^-$ let $y(h)$ be the $\bar\mu$-proportion of $M^-$ with smaller $\sigma$; both are uniform and independent of $|h|$, $z\to0$ at the boundary ray of $U$ in $M^+$, and $y\to1$ at $\R_+(-\alpha',1-\alpha')$. Fix $0<r_1<r_2$, put $A_0\coloneqq\{r_1\le|h|\le r_2\}$, and fix $z_0,y_0\in(0,1)$ with $\sup_{\{z<z_0\}}x<\inf_{\{y\ge y_0\}}x$ (on host~A, $z_0$ small and $y_0$ near $1$; on host~B, any values). Put $S\coloneqq M^+\cap A_0\cap\{z<z_0\}$, $W\coloneqq M^-\cap A_0\cap\{y\ge y_0\}$ and $a\coloneqq\sup_Sx$, so that $S\subseteq c_a$ and $W\subseteq d_a$; on the closure of $S\cup W$, $|h_1|,|h_2|\ge c|h|\ge cr_1$ for some $c>0$. The design is that of Lemma~\ref{lem:D22}, with the blocks $W\cap\{1-2^{-m}\le y<1-2^{-m-1}\}$ (coordinate $y$) in place of the blocks $J_m$, and with sources $P_{n,i,l}\coloneqq S\cap\{iz_0/K_n\le z<(i+1)z_0/K_n\}$; so $P_\gamma\subseteq S$ and $Q_\gamma\subseteq W$. World~1 is given by the rules (R$^-$), (R$^+$), (R$^0$) of Theorem~\ref{thm:D21} with $\mathcal K_0$ replaced by $\mathcal K_0^W$, in which the clause ``$f\ge0$ a.e.\ on $M\cap\{x\ge b^*\}$ for some $b^*<1$'' becomes ``$f\ge0$ a.e.\ on $W\cap\{x\ge b^*\}$ for some $b^*<1$''. World~2 is additive with density $1$ with respect to $\bar\mu$, and the selection is $p\mapsto C_{p_2}$. Call this two-world ranking $\mathrm{D2}'_H$. Its protected set is $W$, and every generator source and target lies in $S\cup W$; world~1's base order (tail dominance in $x$) still acts on all of $M$.

\begin{theorem}\label{thm:F16}
For $H\in\{\mathrm A,\mathrm B\}$, the two-world ranking $\mathrm{D2}'_H$ on $\M_U$ satisfies strict (G), (IB), (LA), (S) (competitiveness on $\supp p$, and closed sections for every model at both worlds), (C), and (LF) at every interior $C_b$, and $\C_L=\{C_b:b\in[0,1]\}$. World~2 is additive. World~1 does not add up: every bounded finitely additive $\nu\ge0$ on $U$ that vanishes on null sets and endorses the relevant world-1 comparisons has $\nu(S)=0$; the failure is of type (M2). Every generator source and target, and the protected set $W$, lie in one compact subset of $\R^2\setminus\{0\}$ on which $|h_1|$ and $|h_2|$ are bounded away from $0$, whose closure meets $\partial U$ and neither unanimous ray; and the targets overlap at $\partial U$: $\bar\mu(Q_\gamma\cap Q_{\gamma'}\cap\{x\ge t\})>0$ for all $\gamma,\gamma'$ and all $t<1$.
\end{theorem}
\begin{proof}[Proof sketch]
As for Theorem~\ref{thm:D21}, with one observation used throughout: $P_\gamma\cap W=\varnothing$, so $\gamma_*=1_{Q_\gamma}$ on all of $W$. Strict (G): a base $\delta\le0$ has $\int_M\delta\,d\bar\mu=\tau_\delta(0)\ge0$, so $\delta=0$; a generator-branch $\delta\le0$ would be $\ge1$ on the non-null set $Q_\gamma\cap W\cap\{x\ge t^*\}$. Sums: two generator translates would give a sum $\ge2$ on $Q_\gamma\cap Q_{\gamma'}\cap W\cap\{x\ge t^*\}$. Sections: as in Lemma~\ref{lem:D24}, with $B_0=[0,\bar\mu(A)]$ and the generator clause ``$Q_\gamma\cap\{x\ge t\}\subseteq A$ a.e.\ for some $t<1$'', which depends on $A$ only. (K) for every model: if $N_X$ is null, $1_Q-1_P\ge0$ on $\{x\ge b\}$, hence on $W\cap\{x\ge b\}$. No weak domination: as in step~4 of Theorem~\ref{thm:D21}, since $Q_\gamma\cap\{x\ge t\}\subseteq W$. $\C_L$ by the direct argument of \src{D2\_referee2} \S2.5. Certificate: $C_a\preceq_1X_\gamma$ with $E_1X_\gamma=(c_a\setminus P_\gamma)\cup Q_\gamma$ is relevant, and the weights of (d4), with $S$ in place of $c_a$, give $\nu(S)\le\epsilon\,\nu(M)$ for a bounded $\nu$. (M2): the generator terms are excluded on $W$ near $x=1$ by the overlap of the targets.
\end{proof}
\status{proved (\src{F1} Lemmas~F1.5$'$, F1.6a, F1.6b, Thm.~F1.6); confirmed (\src{F12\_referee} \S\S4.1--4.4), one reader, with a wording repair incorporated: \src{F1} said that every non-additive ingredient lies in one compact subset; world~1's base order is tail dominance in $x$, an incomplete, non-additive order acting on all of $M$, so only the generators (sources and targets) and the protected set are confined. NUMERICS \src{f1py/d2prime\_cell.py} (exact cell model, 0 failures in every listed test; a density heavy on $W$ but light on the non-$W$ cells of the same levels fails three relevant cut moves, as Theorem~\ref{thm:F12} predicts). This is the nearby reachable form of the request to place every non-additive trade at payoffs bounded away from zero, with protection at $\partial U$; by Corollary~\ref{cor:F15} the targets cannot also be confined to one compact subset of $U$ that avoids $\R_+(1,0)$.}

\begin{theorem}\label{thm:F18}
Let $\nu$ be a countably additive measure on $U$, with $\nu\ll\mu$ and positive on non-null sets, that endorses the relevant world-1 comparisons of $\mathrm{D2}'_H$ under a reading with property (R).
\begin{enumerate}[leftmargin=*,label=(\arabic*)]
\item $\nu$ is not bounded.
\item If $\nu$ is $\sigma$-finite on $M$, then $\nu|_M=\varphi(x)\bar\mu$ with $\varphi$ nondecreasing and finite on $[0,1)$ and $\int_t^1\varphi=\infty$ for every $t<1$; so $\nu$ is finite on every $c_b$ with $b<1$, and $\nu(V\cap M^-)=\infty$ for every neighbourhood $V$ of a point of $\R_+(-\alpha',1-\alpha')$. On host~A, moreover, $\nu(K)=\infty$ for every compact $K\subset U$ with a point of $\R_+(1,0)$ in its interior; so no endorsing measure of world~1 of $\mathrm{D2}'_{\mathrm A}$ is finite on compact subsets of $U$, although no generator, target or protected set comes near $\R_+(1,0)$. On host~B, world~1 has an endorsing measure whose density is locally integrable on $U$, and together with density $1$ at world~2 the conclusion of Conjecture~\ref{conj:sigmaB} holds in $\mathrm{D2}'_{\mathrm B}$.
\item $\sigma$-finite endorsing measures exist, of the form $\varphi(x)\bar\mu$ on $M$ and $\bar\mu$ on $N_+$, and they are finite on every compact subset of $M$. Endorsing measures that are not $\sigma$-finite also exist.
\end{enumerate}
\end{theorem}
\begin{proof}[Proof sketch]
(1) Step~7 of the proof of Theorem~\ref{thm:F16}. (2) Theorem~\ref{thm:F12} with (1), and Lemma~\ref{lem:F11} (its $M^-$ half, and on host~A also its $M^+$ half); on host~B every compact $K\subset U$ has $K\cap M\subseteq c_{b_0}$ with $b_0<1$. (3) For nondecreasing $\varphi\ge1$, summation by parts (Tonelli on $\delta^+$ and $\delta^-$ separately) gives $\int\delta\varphi(x)\,d\bar\mu=\varphi(0)\tau_\delta(0)+\int\tau_\delta\,d\varphi\ge0$ for every base vector $\delta$; choosing $\varphi$ blockwise so large that $\int_{Q_\gamma\cap\{x\ge t\}}\varphi\,d\bar\mu=\infty$ for every $\gamma$ and $t<1$ (possible because every target meets every neighbourhood of the end) makes every generator comparison endorsed with infinite cost on the right. The measures $\nu_t$ of Theorem~\ref{thm:F13}(4) carry over.
\end{proof}
\status{proved (\src{F1} Lemma~F1.7, Thm.~F1.8); confirmed (\src{F12\_referee} \S\S4.5--4.6), one reader. So a $\sigma$-finite endorser blows up exactly at the far end of the chain of selected error sets: on host~A that is $\R_+(1,0)$ together with $\R_+(-\alpha',1-\alpha')$, on host~B the boundary ray alone. Where the non-additive trades sit plays no role, and whether Conjecture~\ref{conj:sigmaB} fails depends on the selection. Lemma~F1.7 holds for any host and any countable design whose targets meet every neighbourhood of the end; \src{F1}'s parenthetical remark that transitivity forces such overlap is unproved (\src{F12\_referee} \S4.5) and is not used.}

\begin{conjecture}\label{conj:F1Cp}
Let $n\ge2$ on $\M_U$, with strict (G), (IB), (LA), (S) and (C). Then every world $i$ has a countably additive measure $\rho_i\mu$ that endorses its relevant comparisons, with $\rho_i>0$ a.e.\ and $\rho_i$ locally integrable on $M$, the open set of mixed problems of $U$.
\end{conjecture}
\status{CONJECTURE (\src{F1} \S2.4, Conjecture~F1.C$'$, suggested in the audit as the repair of Conjecture~\ref{conj:sigmaB}). The systems of Theorems~\ref{thm:D21} and~\ref{thm:F16} satisfy it (Theorems~\ref{thm:F13}(3), \ref{thm:F18}(3)). In two worlds next to an additive world with integrable density, it holds at the other world iff that world has a $\sigma$-finite endorsing measure (from Theorem~\ref{thm:F12}; \src{F1} \S5, confirmed). \src{F1} conjectures that such a measure always exists there (Conjecture~F1.C, open); in the architecture of Theorem~\ref{thm:D21} it does (Theorem~\ref{thm:F18}(3)). Since $\partial D_{T_1}\setminus0\subseteq M$, the localisation route in the status of Conjecture~\ref{conj:sigmaB} needs only this conjecture. \emph{Caveat} (\src{F12\_referee} \S4.7). The evidence is two-world only. There, next to an additive world, Theorem~\ref{thm:F12} places every blow-up at the far end of the chain, which is $\R_+(1,0)$ or $\partial U$ and never interior to $M$, so the conjecture is automatic whenever a $\sigma$-finite endorser exists. In three worlds without (Nest$_e$), nothing proved excludes ends that shrink to interior rays of $M$ or to the $T_1$-correct side of $\partial D_{T_1}$, and exactly there the conjecture would fail if endorsers must blow up at ends. So its use for $T_1$ depends on a location result for ends of the kind of Proposition~\ref{prop:F21}(b), that is, on (Nest$_e$).}

\emph{The a.e.-finite target} (\src{F2} \S2.3, CONJECTURE; with \src{F12\_referee} \S5.2). Every world of a full system has a countably additive endorsing measure whose density is positive and finite a.e. The system of Theorem~\ref{thm:D21} satisfies it, with $1/(1-x)$. \emph{Reading.} Under the total-cost reading this target is weak: a density that is only finite a.e.\ may give both models of a comparison infinite cost, and then the comparison is endorsed vacuously. So the target, and any use of it, must read endorsement in the cancellation form, or require the left-hand model $B$ of each comparison $B\preceq_iX$ to have finite cost $\rho_i\mu(E_iB)$. For $T_1$ it would also need versions of Theorems~\ref{thm:E11} and~\ref{thm:E15} for shadows that are non-integrable only at ends; these are not proved (part~(g)).

\paragraph{(f) Rounds 14 to 16: the link class, total worlds, glued credal worlds, general domains, and what Theorem~\ref{thm:M0closed} needs.} Round 14 (\src{BRIEF13}) asked for an exact cover of $T_1$: a construction with every axiom verified. Three attempts were made and none produced a cover: a world whose interior comparisons absorb the corner cycle (\src{B1}), a total world with crossing edge selections (\src{B3}), and an ordinal resolution of exceptional additive ties on a two-sided domain (\src{B2}). Each failure report contains proved results, and each was refereed once (\src{B1\_referee}, \src{B3\_referee}, \src{B2\_referee}); where a referee repaired a note, the referee's version is given. Round 15 (\src{BRIEF14}) followed the escape left open by Corollary~\ref{cor:B14} below: worlds whose edge orders are unanimities over several densities (\src{C1}, \src{C2}). The round-15 notes were refereed once in \src{C12\_referee}, which found nothing wrong and asked for small repairs; its verdicts and repairs are used below. Rounds 14 and 15 found no world that fails to add up in a full system, for any target; round 16 found one in two worlds (Theorem~\ref{thm:D21}, end of part~(e)). Round 16 also contributes two blocks at the end of this part: (LF) excludes the full-atom residual on general decomposable domains (\src{D1}, refereed once in \src{D1\_referee}), and what the argument of Theorem~\ref{thm:M0closed} needs, with the consequence that the mechanism of Theorem~\ref{thm:D21} cannot carry a cover of $T_1$ in which every world is an end-protected decoration with integrable, jointly end-heavy base densities (\src{E1}, proved in a note; audited in \src{F2}, which writes out the hypotheses and narrows the scope; audit refereed once in \src{F12\_referee}). Round 17 follows in part~(f$'$), and round 19 in part~(f$''$). \src{B1} \S1 reads the additive contradiction of Section~\ref{sec:main} as a closed cycle of exchange-rate links (three facet links near $T_1$, three corner identifications through the corner equation, three edge pinnings by the box) and asks at which link a non-additive world could break it (CHECKED, a reading); the results below locate the answer on the edges.

\emph{Sign classes and the link class} (\src{B1} \S0; confirmed, \src{B1\_referee} (i)). Let $n=3$ and $\{i,j,k\}=\{1,2,3\}$. For $l=1,2,3$ let $Q_l$ be the set of $h\in U$ whose $l$-th coordinate has the sign opposite to that of both other coordinates. The $Q_l$ are pairwise disjoint and their union is $M$ up to the null coordinate planes; $R_{ij}=Q_i\cup Q_j$, and $Q_i=R_{ij}\cap R_{ik}$ is the \emph{link class} of world $i$. It is non-null: it contains a neighbourhood of $e_i-\eta(e_j+e_k)$ (the linking step of Theorem~\ref{thm:Q36}), and for $T_1$ also a neighbourhood of a corner ray of $D_{T_1}$ ($r_{23}\in Q_1$, $r_{31}\in Q_2$, $r_{12}\in Q_3$). Every measurable $A\subseteq M$ is the world-$i$ error set of the regular model $P_i\sd A$, and for regular $Y,Z$, $Y\sd Z=E_iY\sd E_iZ\subseteq M$.

Let $S\subseteq M$ be measurable and non-null, and write $\rho(A)\coloneqq\int_A\rho\,d\mu$. A positive integrable $\rho$ on $S$ is a \emph{leading density of world $i$ on $S$} if, for all regular $Y,Z$ with $Y\sd Z\subseteq S$,
\[
\rho(E_iZ\setminus E_iY)>\rho(E_iY\setminus E_iZ)\ \Longrightarrow\ Y\preceq_iZ ;\tag{L1}
\]
it is \emph{additive on $S$} if moreover $Y\preceq_iZ$ iff $\rho(E_iZ\setminus E_iY)\ge\rho(E_iY\setminus E_iZ)$.

\begin{lemma}\label{lem:B10}
Let world $i$ satisfy strict (G), and let $\rho$ satisfy (L1) on $S$. Then, for all regular $Y,Z$ with $Y\sd Z\subseteq S$,
\[
Y\preceq_iZ\ \Longrightarrow\ \rho(E_iZ\setminus E_iY)\ge\rho(E_iY\setminus E_iZ).\tag{L2}
\]
So $\rho\mu|_S$ endorses every comparison of world $i$ between regular models that differ inside $S$.
\end{lemma}
\begin{proof}[Proof sketch]
If $Y\preceq_iZ$ and $g\coloneqq\rho(E_iY\setminus E_iZ)-\rho(E_iZ\setminus E_iY)>0$, then $S\setminus E_iZ$ is non-null. Take a non-null $e\subseteq S\setminus E_iZ$ with $\rho(e)<g/2$ and put $Z'\coloneqq Z\sd e$. Then $Z\prec_iZ'$ by strict (G), and (L1) gives $Z'\preceq_iY$; so $Z\prec_iZ'\preceq_iY\preceq_iZ$, which is impossible in a preorder.
\end{proof}
\status{proved (\src{B1} Lemma~B1.0); confirmed (\src{B1\_referee} (ii)), one reader. The proof uses only that $\rho\mu$ is atomless, not that $\rho>0$; the strengthening of Corollary~\ref{cor:B14} uses this.}

\begin{lemma}\label{lem:B11}
Let world $i$ satisfy strict (G) and have leading densities $\rho^{ij}$ on $R_{ij}$ and $\rho^{ik}$ on $R_{ik}$. Then $\rho^{ik}=c\,\rho^{ij}$ a.e.\ on $Q_i$ for a constant $c>0$.
\end{lemma}
\begin{proof}[Proof sketch]
If $f\coloneqq\rho^{ik}/\rho^{ij}$ is not a.e.\ constant on $Q_i$, choose $a$ and $\delta>0$ with $S_-\coloneqq\{f<a-\delta\}$ and $S_+\coloneqq\{f>a+\delta\}$ non-null, and $\eta>0$ with $(a-\delta)(1+\eta)<a+\delta$. Take non-null $B\subseteq S_+$ and $A\subseteq S_-$ with $\rho^{ij}(A)=(1+\eta)\rho^{ij}(B)$; then $\rho^{ij}(A)>\rho^{ij}(B)$ and $\rho^{ik}(A)<\rho^{ik}(B)$. The regular models $Y\coloneqq P_i\sd B$ and $Z\coloneqq P_i\sd A$ have $Y\sd Z=A\cup B\subseteq Q_i$; (L1) for $\rho^{ij}$ gives $Y\preceq_iZ$, and (L2) for $\rho^{ik}$ (Lemma~\ref{lem:B10}) then gives $\rho^{ik}(A)\ge\rho^{ik}(B)$, a contradiction.
\end{proof}
\status{proved (\src{B1} Lemma~B1.1, ``gluing''); confirmed (\src{B1\_referee} (ii)), one reader, who notes that the intermediate-value property of an atomless measure (Sierpi\'nski) suffices where the note cites Lyapunov. The \emph{glued density} $\rho_i$ is $\rho^{ij}$ on $R_{ij}$ and $\rho^{ik}/c$ on $R_{ik}$; it is well defined, positive and integrable on $R_{ij}\cup R_{ik}=M$, and leading on each edge region.}

\begin{lemma}\label{lem:B12}
Let $\rho$ be a positive integrable density on $M$, let $C,X$ be regular, $D\coloneqq C\sd X$, and suppose $\Delta\coloneqq\rho(E_iX)-\rho(E_iC)\ge0$. Then there is a finite chain of regular models $C=Z_0,Z_1,\dots,Z_N=X$ such that each difference $Z_t\sd Z_{t+1}$ lies in $R_{ij}$ or in $R_{ik}$, and $\rho(E_iZ_{t+1})-\rho(E_iZ_t)=\Delta/N$ for every $t$.
\end{lemma}
\begin{proof}[Proof sketch]
Put $D_l\coloneqq D\cap Q_l$, and split each $D_l\cap E_iC$ and each $D_l\setminus E_iC$ into $m$ pieces of equal $\rho$-mass. \emph{Case 1: $Q_i\setminus D$ is non-null.} It contains a non-null set on which $C$ is correct at world $i$, or one on which $C$ errs; a small subset of it serves as a reserve. Each round has two steps: one toggles a piece of $D_i$ and a piece of $D_j$ (inside $R_{ij}$), the other a piece of $D_k$ (inside $R_{ik}$). The reserve lies in $Q_i$, hence in both regions; it is toggled in the first step of the round and back in the second, and the order of the two steps is chosen according to the sign needed, so that every step changes the $\rho$-error by exactly $\Delta/(2m)$. The reserve needed has mass $|\Delta/(2m)-s_a/m|$, where $s_a$ is the net change from $D_i\cup D_j$, so it fits for large $m$. \emph{Case 2: $Q_i\subseteq D$ a.e.} Split $D=D'\sqcup D''$ so that the six masses $\rho(D_l\cap E_iC)$, $\rho(D_l\setminus E_iC)$ are halved, and go from $C$ to $C\sd D'$ and then to $X$; each leg has room in the link class and net change $\Delta/2$, and one $m$ serves both legs.
\end{proof}
\status{proved (\src{B1} Lemma~B1.2, the ``transfer chain''); confirmed (\src{B1\_referee} (iii)), one reader, who adds the sentence on a common $m$ in Case~2 (the reserve needed in each leg decreases like $1/m$). NUMERICS: \src{b1py/transfer\_chain.py} (4000 exact instances, 2000 of them without room in the link class) and the referee's independent \src{b1ref\_py/ref\_checks.py} \S2 (1683 instances, both cases and $\Delta=0$, with equality of the step changes checked).}

\begin{theorem}\label{thm:B13}
Let $n=3$, and let world $i$ of a ranking on $\M_U$ satisfy strict (G) and have leading densities on $R_{ij}$ and on $R_{ik}$, with glued density $\rho_i$. Then for all regular $C,X$:
\begin{enumerate}[leftmargin=*,label=(\alph*)]
\item $\rho_i(E_iX)>\rho_i(E_iC)$ implies $C\preceq_iX$;
\item $C\preceq_iX$ implies $\rho_i(E_iX)\ge\rho_i(E_iC)$.
\end{enumerate}
Hence $\rho_i\mu|_M$ is a leading charge of world $i$ on $M$; $\nu_i\coloneqq\rho_i\mu|_M+\mu|_{N_+}$ is positive on non-null sets and endorses every comparison of world $i$ between regular models; and world $i$ adds up. If both edge restrictions are additive, world $i$ is additive on regular models, with density $\rho_i$.
\end{theorem}
\begin{proof}[Proof sketch]
(a) Lemma~\ref{lem:B12} gives a chain whose steps lie inside one edge region and increase the $\rho_i$-error strictly, so (L1) on that region gives $Z_t\preceq_iZ_{t+1}$, and transitivity gives $C\preceq_iX$. (b) is Lemma~\ref{lem:B10} with $S=M$, using (a) as (L1) on $M$. In the additive case, the chain for $\rho_i(E_iX)\ge\rho_i(E_iC)$ has steps with change $\ge0$, each accepted by the additive edge representation.
\end{proof}
\status{proved (\src{B1} Thm.~B1.3); confirmed (\src{B1\_referee} (iv)), one reader, who checks that $\nu_i$ meets Definition~\ref{def:addup}(b), (c) and recommends the framing as Theorem~\ref{thm:Q36}(a$'$). Uses: transitivity of $\preceq_i$, strict (G) at world $i$, regular models, Lyapunov's theorem. Not used: (IB), (S), (C), (LA), completeness, any other world, any selection. NUMERICS \src{b1py/t1\_chain.py} (a sampled $T_1$ instance, $\varepsilon=0.3$, $s=0.9$: at each world the chain uses only edge steps; the strict comparisons that the script prints rest on Lemma~\ref{lem:B10}, not on the chain alone, \src{B1\_referee} (vi)). It is consistent with Theorem~\ref{thm:S21} and with the status of Theorem~\ref{thm:Q36}: an \emph{endorsing} edge weight need not propagate to interior comparisons, but a \emph{leading} edge density does. Without leading densities the edge comparisons of a world do not determine its other comparisons:}

\begin{proposition}\label{prop:B1cex}
Let $n=3$, and let $a\subseteq Q_j$ and $b\subseteq Q_k$ be non-null. For measurable $A,B\subseteq M$ put $A\le B$ iff $A\subseteq B$ a.e., or $A\setminus B\subseteq a$ and $B\setminus A\supseteq b$. Then $\le$ is a strictly monotone preorder satisfying de Finetti cancellation, so, as an order on world-$i$ error sets, it is a world with strict (G) and (IB). For $A,B$ with $A\sd B\subseteq R_{ij}$ or $A\sd B\subseteq R_{ik}$, $A\le B$ iff $A\subseteq B$ a.e.; but $a\le b$, although $a\not\subseteq b$.
\end{proposition}
\status{proved in a referee report (\src{B1\_referee} \S0, counterexample), one reader; NUMERICS: brute force on six atoms (\src{b1ref\_py/ref\_checks.py} \S3). It concerns a single world, not a full ranking. A comparison of the second kind that is not an inclusion has $A\setminus B$ meeting $Q_j$ and $B\setminus A$ meeting $Q_k$, so it lies in neither edge region. So the comparisons of a world inside its two edge regions determine its comparisons on $M$ \emph{when} they admit leading densities (Theorem~\ref{thm:B13}), and not in general.}

\begin{corollary}\label{cor:B14}
Assume (H$_U$). Then some world $i$ and some $j\ne i$ are such that no positive integrable density on $R_{ij}$ is a leading density of world $i$ on $R_{ij}$. In particular some world is not additive on comparisons between regular models that differ inside one of its edge regions. More precisely, every world of the ranking that does not add up has, on one of its two edge regions, no leading density.
\end{corollary}
\begin{proof}
A world with leading densities on both of its edge regions adds up (Theorem~\ref{thm:B13}). So a world that does not add up lacks one; and if every world had both, the ranking would be in (M0), which Theorem~\ref{thm:M0closed}(b) excludes.
\end{proof}
\status{proved at the status of Theorem~\ref{thm:M0closed}(b) (\src{B1} Cor.~B1.4); confirmed (\src{B1\_referee} (v)), one reader; equivalently through Theorem~\ref{thm:Q36}(a$'$) and Corollary~\ref{cor:Q37}. \emph{Strengthening} (CHECKED, by inspection of the proof of Theorem~\ref{thm:S21}; \src{B1\_referee} (v)): ``positive'' may be weakened to ``integrable, $\ge0$ and not a.e.\ zero on $R_{ij}$''. That proof uses the leading property only for comparisons between $C_q$, $q\in e_{ij}$, and edge test models $X^{ij}_A$, whose differences lie in $R_{ij}$, together with strict (G) at worlds $i$ and $j$ and competitiveness and closed sections at edge indices. It needs both properties of a leading charge; for a density $\rho\ge0$ satisfying (L1), the second property (L2) is supplied by Lemma~\ref{lem:B10}, which does not use positivity. So such a $\rho$ is strictly positive on $R_{ij}$, and Lemma~\ref{lem:B11} applies. \emph{Scope.} The corollary concerns integrable, that is countably additive, densities only. Lemma~\ref{lem:B11} uses the Radon--Nikodym theorem and Lemma~\ref{lem:B12} uses Lyapunov's theorem, and neither covers finitely additive edge charges or non-integrable densities. Gluing of finitely additive leading charges is open: the natural statement, that leading charges of world $i$ on $R_{ij}$ and on $R_{ik}$ that are positive on $Q_i$ are proportional there, is not proved (\src{B1\_referee} (v), scope remark).}

\begin{corollary}\label{cor:B15}
Let $n=3$ and let world $i$ satisfy strict (G).
\begin{enumerate}[leftmargin=*,label=(\alph*)]
\item Let $\psi_i(X)\coloneqq\int_{E_iX}\rho_i\,d\mu$ with $\rho_i>0$ a.e.\ and integrable on $M$, and suppose that world $i$ accepts every $\psi_i$-strict comparison between regular models that differ inside $R_{ij}$ or inside $R_{ik}$. Then world $i$ contains every $\psi_i$-strict comparison between regular models, and only comparisons that $\psi_i$ endorses. So adding comparisons to world $i$ adds only $\psi_i$-ties, and a reversed comparison $C\preceq_iX$ with $\psi_i(C)>\psi_i(X)$ contradicts Theorem~\ref{thm:B13}(b): by Lemma~\ref{lem:B10} it would produce a strict cycle through a strict-(G) perturbation.
\item Let $\psi^1,\dots,\psi^r$ be additive orders at world $i$ with positive integrable densities $\rho^s$ on $M$, and suppose that, on each of $R_{ij}$ and $R_{ik}$, all $\psi^s$ induce the same comparisons between regular models that differ inside that region. Then $\rho^s=a_s\rho^1$ a.e.\ on $M$ for constants $a_s>0$, so $\psi^s=a_s\psi^1$ on regular models, and the minimum, the maximum, or any other monotone aggregate of the $\psi^s$ induces the order of $\psi^1$.
\end{enumerate}
\end{corollary}
\begin{proof}[Proof sketch]
(a) is Theorem~\ref{thm:B13}, with $\rho_i$ leading on both regions. (b) On $R_{ij}$, $\rho^s$ and $\rho^1$ are both leading densities of the order of $\psi^1$; the argument of Lemma~\ref{lem:B11}, with $R_{ij}$ in place of $Q_i$, gives $\rho^s=a_s\rho^1$ there, and likewise $\rho^s=b_s\rho^1$ on $R_{ik}$; $a_s=b_s$ because $Q_i$ is non-null.
\end{proof}
\status{proved (\src{B1} Cor.~B1.5) with the referee's corrections (\src{B1\_referee} (vi)), one reader. In (a), \src{B1} derived the reversal's strict cycle from the first claim of Theorem~\ref{thm:B13}, which gives only a weak comparison; the contradiction is with the second claim, as stated here, and the hypothesis is stated precisely. (b) is the referee's ordinal replacement for \src{B1}'s version, which read ``agree on the edge regions'' as equality of measures and is then trivial. These were the two mechanisms proposed in round 14 for absorbing the corner cycle into interior comparisons (an additive base with extra or reversed interior comparisons, and the minimum or maximum of several additive scores); both collapse. By Corollary~\ref{cor:B14}, the cycle can be absorbed only at a world that has no leading density on some edge region; the rest of this part concerns such worlds.}

\emph{Total worlds} (round 14; \src{B3}, refereed in \src{B3\_referee}, one reader). In this block $n\ge2$ and $K'\subseteq\Di$ are arbitrary unless stated. World $i$ is \emph{total} if $\le_i$ is a total preorder on the measurable subsets of $M$; this is the sense of ``complete'' in Theorems~\ref{thm:E3} and~\ref{thm:Q36}(c), not that of \S\ref{sec:CPa-def}, where a complete (CP-a) world has a singleton credal set. World $i$ is \emph{almost representable on $M$} (\emph{AR}) if some finitely additive $\pi\ge0$ on $M$, vanishing on null sets, with $\pi(M)>0$, satisfies $A\le_iB\Rightarrow\pi(A)\le\pi(B)$ for all measurable $A,B\subseteq M$. A \emph{top level} of world $i$ on the subsets of $R_{ij}$ (or of $M$) is a total preorder $\trianglelefteq$ there with $A\le_iB\Rightarrow A\trianglelefteq B$; it is \emph{leading} if moreover $A\vartriangleleft B\Rightarrow A\le_iB$, and \emph{cancellative} if $A\trianglelefteq B$ iff $A\setminus B\trianglelefteq B\setminus A$. For a total world every top level is leading. A set $D$ is \emph{$\trianglelefteq$-null} if $D$ is $\trianglelefteq$-equivalent to $\varnothing$. For $p\in e_{ij}$ write $c_p\coloneqq E_i(C_p)\subseteq R_{ij}$ (Lemma~\ref{lem:Q34}).

\begin{lemma}\label{lem:B31}
Let world $i$ be total. A charge $\pi\ge0$ on $M$ that vanishes on null sets and has $\pi(M)>0$ is a leading charge of world $i$ on $M$ iff it almost agrees with $\le_i$, i.e.\ $A\le_iB\Rightarrow\pi(A)\le\pi(B)$.
\end{lemma}
\status{proved (\src{B3} Lemma~B3.1: almost agreement gives the second property of a leading charge, and with totality the first); confirmed (\src{B3\_referee} \S1), one reader. Only totality on the subsets of $M$ is used.}

\begin{corollary}\label{cor:B32}
Let a ranking on $\M_U$ ($n\ge2$, any $K'\subseteq\Di$) satisfy strict (G), (IB) and (S). Every total world that is AR on $M$ has a strictly positive almost-agreeing charge, which endorses every comparison of the world between subsets of $M$, and the world adds up; no nestedness is used. Hence the clause ``complete and (Nest$_e$) on both edges at $e_i$'' of Corollary~\ref{cor:Q37} may be replaced by ``total and AR on $M$''. Under (H$_U$), a total world of a cover of $T_1$ that does not add up is not AR on $M$, and some edge family at $e_i$ is not nested at world $i$.
\end{corollary}
\status{proved at the status of Theorem~\ref{thm:S21} and Corollary~\ref{cor:Q37} (\src{B3} Cor.~B3.2: Lemma~\ref{lem:B31} with Theorem~\ref{thm:S21}); confirmed with a wording repair (\src{B3\_referee} \S2), one reader. \src{B3} also wrote that a complete world in a cover must be non-AR; that is false as worded, since a cover may contain total worlds that add up next to another world that does not, and only the form stated here is claimed. The last clause is the contrapositive of Theorem~\ref{thm:Q36}(c) (CHECKED, Savage). By Lemma~\ref{lem:B31}, ``total and AR on $M$'' is a special case of ``leading charge nonzero on $M$''.}

\begin{lemma}\label{lem:B33}
Let a ranking on $\M_U$ ($n\ge2$, any $K'\subseteq\Di$) satisfy strict (G), (IB) and (S). Fix an edge $e_{ij}$, a point $p^*\in e_{ij}$ and a measurable $A\subseteq R_{ij}$.
\begin{enumerate}[leftmargin=*,label=(E\arabic*)]
\item If $c_{p^*}\le_iA$ fails, then $c_q\le_iA$ fails for all $q\in e_{ij}$ near $p^*$.
\item If $A\supsetneq c_{p^*}$ with non-null difference, then $c_q\le_iA$ for all $q\in e_{ij}$ near $p^*$.
\end{enumerate}
\end{lemma}
\begin{proof}[Proof sketch]
For $p\in e_{ij}$, $C_p\preceq_iX^{ij}_A$ iff $c_p\le_iA$ (Lemma~\ref{lem:Q34} and the normal form). (E1): the section $S_i(X^{ij}_A)$ is closed. (E2): strict (G) at world $j$ gives $X^{ij}_A\prec_jC_{p^*}$, so by closedness of $S_j(X^{ij}_A)$ some neighbourhood $N$ of $p^*$ has $C_q\not\preceq_jX^{ij}_A$ for $q\in N$; competitiveness on $\supp q\subseteq\{i,j\}$ then gives $C_q\preceq_iX^{ij}_A$ ($q=e_j$ is impossible, and $q=e_i$ is trivial).
\end{proof}
\status{proved (\src{B3} Lemma~B3.3, ``order continuity on edges''); confirmed (\src{B3\_referee} \S3), one reader. Uses: closed sections at the edge test models, strict (G) at world $j$, competitiveness at edge indices, and (IB) through the normal form. Not used: (C), (LA), the third world. It is the semicontinuity step of Theorem~\ref{thm:S21} with the charge removed. For a total world and non-null $D\subseteq c_{p^*}$, $D'\subseteq R_{ij}\setminus c_{p^*}$ it gives, in $\le_i$ (not as sets), $c_{p^*}\setminus D<_ic_q\le_ic_{p^*}\cup D'$ for $q$ near $p^*$.}

\begin{proposition}\label{prop:B34}
Let $R=[0,2)$ with Lebesgue measure $\lambda$ (every atomless edge region is measure-isomorphic to it), let $\mathcal U$ be an ultrafilter of the measure algebra of $R$ containing every interval $(1,1+h)$, $h>0$, let $0<\theta<1$, and put $\pi\coloneqq\lambda+\theta\,\delta_{\mathcal U}$. Consider the two-world edge system in which world $i$ ranks its error sets by $\pi$ ($c\le_iA$ iff $\pi(c)\le\pi(A)$) and world $j$ compares complements by $\pi$ ($R\setminus c\le_jR\setminus A$ iff $\pi(c)\ge\pi(A)$), with the path $c_p=[0,3p\theta)$ for $p\le\tfrac13$, $c_p=(1,1+\tfrac32(p-\tfrac13))$ for $\tfrac13<p\le\tfrac23$, and, for $p>\tfrac23$, $c_p$ growing continuously in measure from $(1,1\tfrac12)$ to $R$ while keeping $(1,1\tfrac12)$. Then:
\begin{enumerate}[leftmargin=*,label=(\alph*)]
\item both worlds satisfy strict (G) and (IB);
\item $m(p)\coloneqq\pi(c_p)$ is continuous, with $m(0)=0$ and $m(1)=2+\theta$; so (K) holds at every index (trivially, since both worlds rank by one total order) and the sections at both worlds are closed;
\item $c_{1/3}=[0,\theta)$ and $c_{1/2}$ are disjoint and non-null, so the path crosses; it jumps in measure at $p=\tfrac13$, but not in $\pi$;
\item $\le_i$ is not fine, and no nested edge family with (K) and closed sections (Remark~\ref{rem:straighten}) exists for this pair of edge worlds;
\item $\pi$ is a strictly positive leading charge, so world $i$ is AR and adds up.
\end{enumerate}
\end{proposition}
\status{proved (\src{B3} Prop.~B3.4); confirmed with repairs (\src{B3\_referee} \S4), one reader. The ultrafilter exists by the ultrafilter lemma (BPI, a weak form of the axiom of choice), extending the filter base formed by the classes of the intervals $(1,1+h)$; $\delta_{\mathcal U}$ is a $\{0,1\}$-valued charge that vanishes on null sets and is not countably additive. In (d), the referee's argument replaces \src{B3}'s appeal to Remark~\ref{rem:straighten} and Proposition~\ref{prop:Q35}(5) (Savage): along any chain from $\varnothing$ to $R$, membership in $\mathcal U$ switches once, $\pi$ jumps by at least $\theta$ there, and since the values $\pi(A)$ fill $[0,2+\theta]$, some section is not closed; this is proved and uses neither Savage's theorem nor (K). Non-fineness: every finite partition has a piece in $\mathcal U$, of value $\ge\theta>\pi(Z)$ for $Z\notin\mathcal U$ with $\lambda(Z)<\theta$. In (e), leading charges in Theorem~\ref{thm:S21} may be finitely additive (\src{B3\_referee} \S4.4). NUMERICS \src{b3py/prop4\_edge.py} (its (K) test is a tautology). Failure of (Nest$_e$) was known (Propositions~\ref{prop:plateau}, \ref{prop:R71}); the new content is non-fineness together with crossing. \emph{Scope} (corrected, \src{B3\_referee} \S4.5): in a two-world ranking with every axiom, \emph{including (C)}, this world is impossible, since Theorem~\ref{thm:S22} (which needs (C)) forces nesting and Theorem~\ref{thm:E3} then fineness. Without (C) it is realised in a two-world ranking on $\M_U$ in which both worlds rank error sets $E$ by $\mathrm{lex}(\mu(E\cap N_+),\pi(E\cap M))$ for a charge $\pi=\mu+\theta\delta_{\mathcal U}$, with $\mathcal U$ an ultrafilter of the measure algebra of $M^+$ and $\theta>0$ small: strict (G), (IB), (LA) and (S) hold, with closed sections for every model, the coherent selection crosses, world~1 is total and not fine, and only (C) fails (\src{B3\_referee} Prop.~R1, proved in that report, one reader; NUMERICS \src{b3ref\_py/two\_world\_ultra.py}). Whether a third world can realise it in a full three-world ranking is not settled. For $T_1$ it does not matter, by (e): here crossing defeats fineness but yields only an AR world, which Corollary~\ref{cor:B32} already covers.}

\begin{theorem}\label{thm:B35}
Let a ranking on $\M_U$ ($n\ge2$, any $K'\subseteq\Di$) satisfy strict (G), (IB) and (S). Let $\trianglelefteq$ be a cancellative leading top level of world $i$ on the subsets of $R_{ij}$, and suppose every non-null $W\subseteq R_{ij}$ contains a non-null $\trianglelefteq$-null set. Then all subsets of $R_{ij}$ are $\trianglelefteq$-equivalent. If this holds on every edge at $e_i$ and $\trianglelefteq$ is a cancellative leading top level on $M$, then $\trianglelefteq$ is trivial on $M$.
\end{theorem}
\begin{proof}[Proof sketch]
Let $f(p)$ be the class of $c_p$. If $c_{p^*}$ is non-null, take a non-null $\trianglelefteq$-null $D\subseteq c_{p^*}$. Strict (G) gives $c_{p^*}\setminus D<_ic_{p^*}$, so by (E1) $c_q\le_ic_{p^*}\setminus D$ fails near $p^*$; as the top level is leading, $c_{p^*}\setminus D\trianglelefteq c_q$, and by cancellation $c_{p^*}\setminus D$ is equivalent to $c_{p^*}$. Symmetrically, (E2) with a $\trianglelefteq$-null $D'\subseteq R_{ij}\setminus c_{p^*}$ gives $c_q\trianglelefteq c_{p^*}$. (The null and co-null cases are direct.) So $f$ is locally constant on the connected edge, $\varnothing=c_{e_i}$ is equivalent to $c_{e_j}=R_{ij}$, and by monotonicity every subset of $R_{ij}$ is $\trianglelefteq$-null. Every $A\subseteq M$ is a finite disjoint union of pieces inside the $R_{ik}$, $k\ne i$, and adding $\trianglelefteq$-null pieces does not change the class.
\end{proof}
\status{proved (\src{B3} Thm.~B3.5, for total worlds); confirmed (\src{B3\_referee} \S5.1), one reader. The extension to leading top levels of any world, total or not, is \src{B3\_referee} Remark~R2 (proved by inspection: totality was used once, to pass from ``$c_q\le_ic_{p^*}\setminus D$ fails'' to $c_{p^*}\setminus D\trianglelefteq c_q$, and the leading property gives this directly), one reader. Uses only Lemma~\ref{lem:B33}, strict (G) and the leading property. Not used: (C), (LA), a third world, nestedness, coherence.}

\begin{corollary}\label{cor:B36}
Let a ranking on $\M_U$ ($n\ge2$, any $K'\subseteq\Di$) satisfy strict (G), (IB) and (S). Let $\mathcal U_1,\mathcal U_2,\dots$ be countably many ultrafilters of the measure algebra of $M$, put $S(A)\coloneqq\{j:A\in\mathcal U_j\}$, and let $\le_0$ be a cancellative comparative probability on sets of indices. Then no world has a nontrivial leading top level of the form $A\trianglelefteq B\iff S(A)\le_0S(B)$. In particular the worlds of \emph{Independent Rankings Without Additive Structure} (a comparative probability over finitely many ultrafilter signatures that is not almost representable, as the first level), both in the two-tier form that their order takes on $\M_U$ and in the one-tier transport to error sets, occur in no such ranking, for any $n$ and any target; nor do $\{0,1\}$-valued or purely finitely additive leading first levels.
\end{corollary}
\status{proved (\src{B3} Cor.~B3.6): $S(A\setminus B)=S(A)\setminus S(B)$ gives cancellation; every non-null set contains a non-null set lying in no $\mathcal U_j$ (countably many removals of small total measure, by Lemma~4 of that paper), which is $\trianglelefteq$-null; Theorem~\ref{thm:B35} applies; for pure charges the carrier lemma (\src{N14} Lemma~Y) gives dense null sets, which recovers Corollary~\ref{cor:S21}(b) for pure charges. Confirmed (\src{B3\_referee} \S5.2), with the description of that paper's order repaired (\S5.3): on $\M_U$ it is lexicographic over two signature tiers (the ultrafilters concentrated on $\{h_i<0\}$, then those on $\{h_i>0\}$), still cancellative and signature-monotone, so the corollary covers it; it is non-AR only if the non-representable comparative probability is not split between the tiers. Extended to leading top levels by Remark~R2; one reader.}

\begin{proposition}\label{prop:B37}
Let $\mathcal U_1,\dots,\mathcal U_m$ be distinct ultrafilters of the measure algebra of $M$ with signature map $S$, let $\nu\ll\mu$ be a non-atomic finite measure on $M$, and let $\kappa$ be a real function on the subsets of $\{1,\dots,m\}$ with $\kappa(\varnothing)=0$ and $4\max|\kappa|<\nu(M)$. Suppose $\Phi(A)\coloneqq\nu(A)+\kappa(S(A))$ defines a leading top level of world $i$ on $M$: $A\le_iB\Rightarrow\Phi(A)\le\Phi(B)$, and $\Phi(A)<\Phi(B)\Rightarrow A\le_iB$. Then $\kappa$ is additive, $\kappa_j\coloneqq\kappa(\{j\})\ge0$, and $\pi\coloneqq\nu+\sum_j\kappa_j\delta_{\mathcal U_j}$ is a leading charge of world $i$ on $M$. If the ranking satisfies strict (G), (IB) and (S) on $\M_U$, world $i$ adds up (Theorem~\ref{thm:S21}).
\end{proposition}
\begin{proof}[Proof sketch]
Disjoint small sets $G_j\in\mathcal U_j$ and sets inside their complement, which lies in no $\mathcal U_j$, realise any two signatures $s,t$ together with any difference of $\nu$-mass in an interval of length almost $2\nu(M)$. Cancellation of $\le_i$ and the leading property give $\Phi(A)<\Phi(B)\Rightarrow\Phi(A\setminus B)\le\Phi(B\setminus A)$ and the same with $A$ and $B$ exchanged; with the realisable differences this forces $\kappa(s)-\kappa(t)=\kappa(s\setminus t)-\kappa(t\setminus s)$ for all $s,t$, hence additivity, and monotonicity against small members of $\mathcal U_j$ gives $\kappa_j\ge0$. Then $\Phi=\pi$.
\end{proof}
\status{proved (\src{B3} Prop.~B3.7, for total worlds); confirmed (\src{B3\_referee} \S6.1) and extended to leading top levels (Remark~R2), one reader. NUMERICS \src{b3py/prop7\_kappa.py} (the algebraic core only). It uses neither (S) nor edges: it concerns the form of a top level. \src{B3}'s reading that escaping Theorem~\ref{thm:B35} forces additivity is HEURISTIC: additivity is proved only for value-form top levels with $4\max|\kappa|<\nu(M)$, mixed regimes are not analysed, and no result shows that crossing is the only way to defeat fineness. \emph{Consequence} (\src{B3\_referee} \S6.2): a world of a cover of $T_1$ that does not add up, total or not, has no leading charge nonzero on $M$ (Corollary~\ref{cor:Q37}), no nontrivial cancellative leading top level with dense null sets (Theorem~\ref{thm:B35}), and no leading top level $\nu+\kappa\circ S$ as above.}

\begin{conjecture}\label{conj:B3C}
Let a ranking on $\M_U$ ($n\ge2$, any $K'\subseteq\Di$) satisfy strict (G), (IB) and (S), and let world $i$ be total. Then world $i$ is AR on $M$.
\end{conjecture}
\status{CONJECTURE (\src{B3} Conj.~B3.C); plausible and open (\src{B3\_referee} \S7), with the implications corrected there. It \emph{would follow from} the edge statement that no total, cancellative, strictly monotone, non-AR order on the measure algebra of $M$ admits, on each edge at $e_i$, a path with (K) and closed sections at both edge worlds; \src{B3}'s ``equivalently'' is wrong, since a ranking supplies such edge systems (Lemma~\ref{lem:Q34}) but an edge system need not extend to a ranking. \emph{If true}, every total world of such a ranking adds up (Corollary~\ref{cor:B32}); so in a cover of $T_1$ every world that does not add up is not total, and no ranking all of whose worlds are total covers $T_1$. Total worlds may still occur in a cover, next to a non-total world that does not add up; \src{B3}'s ``complete worlds are excluded from covers outright'' is wrong as worded. \emph{If false}, a counterexample is a total non-AR order whose top level has no dense family of null sets and that is not exchangeable against a non-atomic measure in the sense of Proposition~\ref{prop:B37}; none is known. Neither (C) nor the third world is used in Lemma~\ref{lem:B33}--Proposition~\ref{prop:B37}, so a proof that (C) and the third world force (Nest$_e$) would not settle the total case, and Proposition~\ref{prop:B34} shows that straightening cannot be the general route to the conjecture.}

\begin{proposition}\label{prop:B3R3}
Let $n=2$ on $\M_U$, let $\le$ be a total, cancellative, strictly monotone preorder on the measure algebra of $M$ with chain of classes $\mathbb L$, and let both worlds rank error sets $E$ by $\mathrm{lex}(\mu(E\cap N_+),\,E\cap M\text{ under }\le)$. Then:
\begin{enumerate}[leftmargin=*,label=(\alph*)]
\item strict (G), (IB) and (LA) hold, and every regular model is undominated and competitive on $\{1,2\}$;
\item a selection of coherent models with $C_{e_1}=P_1$ and $C_{e_2}=P_2$ satisfies (S) iff $p\mapsto[E_1C_p\cap M]$ is continuous from $\Delta$ into $\mathbb L$ with its order topology;
\item if (S) holds, $\mathbb L$ is order-isomorphic to $[0,1]$, so $\le$ has a real-valued representation;
\item (C) fails.
\end{enumerate}
\end{proposition}
\status{proved in a referee report (\src{B3\_referee} \S7.2, Prop.~R3, with Lemma~R3.0: for a cancellative order, $A\le B$ iff $M\setminus B\le M\setminus A$), one reader. Proposition~\ref{prop:R71} is the case where $\le$ is the $\mu$-order, and the two-world realisation in the status of Proposition~\ref{prop:B34} the case of the $\pi$-order. So for $n=2$, and without (C), what Conjecture~\ref{conj:B3C} asserts about this ranking is the following order-theoretic statement.}

\begin{conjecture}\label{conj:B3Cp}
Every total, cancellative, strictly monotone order on the measure algebra of $M$ whose chain of classes is the continuous image of $[0,1]$ under $p\mapsto[c_p]$, for some path of coherent error sets $c_p$ running between $\varnothing$ and $M$, is AR.
\end{conjecture}
\status{CONJECTURE (\src{B3\_referee} \S7.2, B3.C$'$). Conjecture~\ref{conj:B3C} implies it (Proposition~\ref{prop:B3R3}), so a counterexample to it refutes Conjecture~\ref{conj:B3C}; a counterexample would also give a two-world ranking without (C) whose total world does not add up. The referee recommends it as the first test. It is consistent with what is known: the orders of Corollary~\ref{cor:B36} have gaps between signature blocks, so their chain of classes is not connected, and the $\pi$-order of Proposition~\ref{prop:B34} has chain of classes $[0,2+\theta]$ and is AR. A next lemma suggested by the referee, Proposition~\ref{prop:B37} with ``$4\max|\kappa|<\nu(M)$'' replaced by ``the chain of classes is connected'', is a CONJECTURE; the referee's sketch for orders that factor through a vector measure with a full difference range is HEURISTIC and is not claimed.}

\emph{Glued credal worlds} (round 15; \src{C1} \S\S1--6, refereed in \src{C12\_referee} \S3, one reader). Let $n=3$, and let world $i$ satisfy strict (G) and (IB), considered on regular models. Let $\Gamma_{ij}$ be a nonempty, norm-compact, convex set of integrable densities on $R_{ij}$, each $>0$ a.e.\ there, and $\Gamma_{ik}$ the same on $R_{ik}$. \textbf{(F)} means that both are finitely generated, $\Gamma_{ij}=\co\{\rho_1,\dots,\rho_a\}$ and $\Gamma_{ik}=\co\{\sigma_1,\dots,\sigma_b\}$. For regular $Y,Z$ with $Y\sd Z\subseteq R_{ij}$, write $Y\sqsubseteq_{ij}Z$ if $\gamma(E_iY\setminus E_iZ)\le\gamma(E_iZ\setminus E_iY)$ for every $\gamma\in\Gamma_{ij}$ (\emph{edge unanimity}), and $Y\sqsubset_{ij}Z$ if all these inequalities are strict (\emph{strict edge unanimity}); define $\sqsubseteq_{ik}$, $\sqsubset_{ik}$ likewise. The \emph{glued closure} $\preceq^{\mathrm{gl}}$ is the smallest relation on regular models that is reflexive, transitive and (IB)-closed (whether $Y\preceq Z$ holds depends only on $(E_iY\setminus E_iZ,\,E_iZ\setminus E_iY)$) and contains $\sqsubseteq_{ij}\cup\sqsubseteq_{ik}$. World $i$ is \emph{$\Gamma$-leading} if it contains $\sqsubset_{ij}\cup\sqsubset_{ik}$; for singletons this means leading densities on both edge regions. The \emph{link cones} are $\mathcal K_{ij}\coloneqq\operatorname{cone}\{\gamma|_{Q_i}:\gamma\in\Gamma_{ij}\}$ and $\mathcal K_{ik}\coloneqq\operatorname{cone}\{\gamma|_{Q_i}:\gamma\in\Gamma_{ik}\}$, closed convex cones in $L^1(Q_i)$. The \emph{glued credal set} is
\[
\Gamma^{\mathrm{gl}}\coloneqq\{\nu\text{ a finite measure on }M,\ \nu\ll\mu:\ \nu|_{R_{ij}}\in\operatorname{cone}\Gamma_{ij},\ \nu|_{R_{ik}}\in\operatorname{cone}\Gamma_{ik}\},
\]
a fibred product (not an intersection of credal sets) whose restriction to $Q_i$ is $\mathcal K_{ij}\cap\mathcal K_{ik}$; $\Gamma^{\mathrm{gl}}_1\coloneqq\{\nu\in\Gamma^{\mathrm{gl}}:\nu(M)=1\}$. Write $Y\preceq_\Gamma Z$ if $\nu(E_iY)\le\nu(E_iZ)$ for every $\nu\in\Gamma^{\mathrm{gl}}$, and call the pair $(Y,Z)$ \emph{$\Gamma^{\mathrm{gl}}$-strict} if the inequality is strict for every $\nu\in\Gamma^{\mathrm{gl}}\setminus\{0\}$. \textbf{(UI)} (under (F)): every nonzero function in the span of the restrictions $\rho_l|_{Q_i}$, $\sigma_m|_{Q_i}$ is nonzero a.e.\ on $Q_i$ (a generic condition on the design).

\begin{theorem}\label{thm:C11}
\begin{enumerate}[leftmargin=*,label=(\alph*)]
\item If $\mathcal K_{ij}\cap\mathcal K_{ik}=\{0\}$, there are regular $Y,Y',Z$, with all differences inside $Q_i$, such that $E_iY\subsetneq E_iY'$ with non-null difference, $Y'\sqsubset_{ij}Z$ and $Z\sqsubset_{ik}Y$. So no preorder with strict (G) contains $\sqsubset_{ij}\cup\sqsubset_{ik}$; in particular $\preceq^{\mathrm{gl}}$ violates strict (G).
\item If $\mathcal K_{ij}\cap\mathcal K_{ik}\ne\{0\}$, then $\Gamma^{\mathrm{gl}}\ne\{0\}$, every $\nu\in\Gamma^{\mathrm{gl}}\setminus\{0\}$ has a density $>0$ a.e.\ on $M$, and $\preceq^{\mathrm{gl}}\subseteq\preceq_\Gamma$; $\preceq^{\mathrm{gl}}$ satisfies strict (G) and (IB) and has no strict cycle, and every $\nu\in\Gamma^{\mathrm{gl}}\setminus\{0\}$ endorses every comparison of $\preceq^{\mathrm{gl}}$.
\end{enumerate}
\end{theorem}
\begin{proof}[Proof sketch]
(a) Hahn--Banach, applied to the compact convex bases of the two cones, gives a bounded $h$ on $Q_i$ with $\int h\gamma\,d\mu\ge c>0$ for the restrictions of all $\gamma\in\Gamma_{ij}$ and $\int h\sigma\,d\mu\le-c$ for those of all $\sigma\in\Gamma_{ik}$. Lyapunov's theorem for a finite net of $\Gamma_{ij}\cup\Gamma_{ik}$ realises $h$ by sets $A\subseteq\{h>0\}$ and $B\subseteq\{h<0\}$, so that $\gamma(A)-\gamma(B)\ge c/2$ and $\sigma(B)-\sigma(A)\ge c/2$. With $Y\coloneqq P_i\sd B$, $Z\coloneqq P_i\sd A$ and $Y'\coloneqq P_i\sd(B\cup e)$ for a small non-null $e\subseteq Q_i\setminus(A\cup B)$ (norm-compact sets of densities are uniformly absolutely continuous), $Y\prec Y'$, $Y'\sqsubset_{ij}Z$ and $Z\sqsubset_{ik}Y$ form a strict cycle. (b) A nonzero $\lambda\gamma_0|_{Q_i}=\lambda'\gamma'_0|_{Q_i}$ in the intersection glues to a member of $\Gamma^{\mathrm{gl}}$; a nonzero member has both scalars positive, hence a density $>0$ a.e.; $\preceq_\Gamma$ is an (IB)-closed preorder containing both edge unanimities, so it contains $\preceq^{\mathrm{gl}}$; for strict (G), an added error set splits into its parts in $R_{ij}$ and in $Q_k$, each an edge step.
\end{proof}
\status{proved (\src{C1} Thm.~C1.1); confirmed (\src{C12\_referee} \S3), one reader, with a cosmetic repair in (a) (the three models agree off $Q_i$), and an independent check (\src{c12ref\_py/ref\_c1.py}: in 400 random cell instances exactly one alternative holds). NUMERICS (\src{c1py/glue\_dichotomy.py}, Part~A, exact rationals): of 300 random cell instances, 176 had $\mathcal K_{ij}\cap\mathcal K_{ik}=\{0\}$, each with an explicit strict cycle, and 124 a glued density positive on every cell. For singletons, (a) is the contradiction in the proof of Lemma~\ref{lem:B11} and (b) is the glued density.}

\begin{theorem}\label{thm:C14}
Assume (F), (UI) and $\mathcal K_{ij}\cap\mathcal K_{ik}\ne\{0\}$. If the pair $(C,X)$ of regular models is $\Gamma^{\mathrm{gl}}$-strict, there is a finite chain of regular models from $C$ to $X$ whose differences lie inside $R_{ij}$ or inside $R_{ik}$ and each of whose steps is a strict unanimity for the generators of its region. Hence, for regular $C,X$,
\[
(C,X)\ \Gamma^{\mathrm{gl}}\text{-strict}\ \Longrightarrow\ C\preceq^{\mathrm{gl}}X\ \Longrightarrow\ C\preceq_\Gamma X ,
\]
so the glued closure is the (CP-a) order with credal set $\Gamma^{\mathrm{gl}}$, up to comparisons at which some nonzero member of $\Gamma^{\mathrm{gl}}$ gives a tie.
\end{theorem}
\begin{proof}[Proof sketch]
Farkas' lemma splits the error difference into a part with a uniform positive margin against every generator of $\Gamma_{ij}$ and a part with such a margin against every generator of $\Gamma_{ik}$, the two parts exchanging a transfer on $Q_i$; this uses that $\Gamma^{\mathrm{gl}}$-strictness has a positive margin on the compact base $\Gamma^{\mathrm{gl}}_1$. (UI) lets the transfer be realised by sets in the room left in $Q_i$, and the rounds of Lemma~\ref{lem:B12} (``raise first, lower later'', or the mirror) then give a chain of strict edge unanimities; if $Q_i\subseteq C\sd X$, split into two legs as in Case~2 there. The second implication is Theorem~\ref{thm:C11}(b).
\end{proof}
\status{proved (\src{C1} Lemma~C1.3, Thm.~C1.4); confirmed (\src{C12\_referee} \S3), one reader, every step re-derived (the Farkas step also checked numerically, 790 of 790 cases). Whether all comparisons at which some member of $\Gamma^{\mathrm{gl}}$ ties lie in $\preceq^{\mathrm{gl}}$ is not settled; it is not used below. NUMERICS Part~B: the chain built in 92 $\Gamma^{\mathrm{gl}}$-strict instances (3 to 524\,289 steps); the script checks only that each step is \emph{weakly} accepted by the generators of its region, while the lemma claims strict acceptance (\src{C12\_referee} \S3). \emph{Gluing refines the edges} (\src{C1} Lemma~C1.2, Prop.~C1.5, under (F), (UI); confirmed, \src{C12\_referee} \S3, which adds the line that the restricted cone is pointed, so the separating functional can be made strictly positive on it): the restriction of $\Gamma^{\mathrm{gl}}$ to $R_{ij}$ is $\{\gamma\in\operatorname{cone}\Gamma_{ij}:\gamma|_{Q_i}\in\mathcal K_{ik}\}$, so the closure's own edge order on $R_{ij}$ is $\Gamma_{ij}$-unanimity only if $\mathcal K_{ij}\subseteq\mathcal K_{ik}$, and both edge orders survive gluing exactly iff $\mathcal K_{ij}=\mathcal K_{ik}$; NUMERICS Part~C (14 of 14 instances). With Theorem~\ref{thm:C11} this is the \emph{glue dichotomy}: the glued closure of two edge unanimities either collapses (a strict 2-cycle in $Q_i$) or is unanimity over $\Gamma^{\mathrm{gl}}$, whose members are global positive densities; it is never both non-collapsing and free of a global density. \emph{Scope} (\src{C12\_referee} \S3): this concerns the glued closure and worlds whose comparisons are $\Gamma^{\mathrm{gl}}$-weak; \src{C1}'s headline that the design is dead with no third case holds in that scope only. $\Gamma$-leading worlds with further comparisons (the band of Theorem~\ref{thm:C17}) remain open (Question~C1.Q).}

\begin{theorem}\label{thm:C16}
Assume (H$_U$), and suppose every world of the ranking is of one of the following kinds: (i) a \emph{glued credal world}, i.e.\ for some $\Gamma_{ij},\Gamma_{ik}$ as above ((F) and (UI) not required) the world contains $\sqsubset_{ij}\cup\sqsubset_{ik}$ and all its comparisons between regular models lie in $\preceq_\Gamma$ (for instance, its comparisons are exactly those of $\preceq^{\mathrm{gl}}$); (ii) (CP-a); (iii) a world with a leading charge nonzero on $M$; (iv) total and AR on $M$. Then $T_1$ is not covered.
\end{theorem}
\begin{proof}
For a world of kind (i), $\mathcal K_{ij}\cap\mathcal K_{ik}=\{0\}$ contradicts strict (G) by Theorem~\ref{thm:C11}(a); otherwise any $\nu\in\Gamma^{\mathrm{gl}}\setminus\{0\}$ has a density $>0$ a.e.\ on $M$ (Theorem~\ref{thm:C11}(b)), and $\nu+\mu|_{N_+}$ is positive on non-null sets of $U$ and endorses every comparison between regular models, so the world adds up. Kinds (ii)--(iv) add up by Corollary~\ref{cor:CPaU}(b), Theorem~\ref{thm:S21} and Corollary~\ref{cor:B32}. So the ranking is in (M0), which Theorem~\ref{thm:M0closed}(b) excludes.
\end{proof}
\status{proved (\src{C1} Thm.~C1.6); confirmed at the status of Theorem~\ref{thm:M0closed}(b) (\src{C12\_referee} \S3), one reader. Kind (i) assumes that every comparison is $\Gamma^{\mathrm{gl}}$-weak, so it is nearly (CP-a) by hypothesis; the residual of the design is the non-total band of Theorem~\ref{thm:C17} (Question~C1.Q). It uses Theorem~\ref{thm:C11} only, with no (F), (UI), selection or third world. A glued credal world supplies a whole credal set of supporting assessments, each a global positive density, and any one of them carries the potential that the corner equation and the box contradict.}

\begin{theorem}\label{thm:C17}
Assume (F) and (UI), and let world $i$ be $\Gamma$-leading. Then $\mathcal K_{ij}\cap\mathcal K_{ik}\ne\{0\}$, and for all regular $C,X$: if $(C,X)$ is $\Gamma^{\mathrm{gl}}$-strict then $C\preceq_iX$; and if $C\preceq_iX$ then some $\nu\in\Gamma^{\mathrm{gl}}\setminus\{0\}$ has $\nu(E_iX)\ge\nu(E_iC)$, i.e.\ no comparison of world $i$ is $\Gamma^{\mathrm{gl}}$-strictly reversed.
\end{theorem}
\status{proved (\src{C1} Thm.~C1.7); confirmed (\src{C12\_referee} \S3), one reader. The first claim is Theorem~\ref{thm:C11}(a), the second the chain of Theorem~\ref{thm:C14}, whose steps are strict unanimities, and the third the perturbation argument of Lemma~\ref{lem:B10} with a small set obtained by Lyapunov equipartition against the generators. So world $i$ lies in the \emph{$\Gamma^{\mathrm{gl}}$-band}: it contains strict $\Gamma^{\mathrm{gl}}$-unanimity and is contained in the relation ``not $\Gamma^{\mathrm{gl}}$-strictly reversed''. For singletons the band has zero width and this is Theorem~\ref{thm:B13}.}

\begin{corollary}\label{cor:C18}
Assume (F) and (UI). A total $\Gamma$-leading world is fine; hence it is AR on $M$, with a strictly positive almost-agreeing probability, and in a ranking with strict (G), (IB) and (S) on $\M_U$ it adds up (Lemma~\ref{lem:B31}, Corollary~\ref{cor:B32}). So under (H$_U$) a world of a cover of $T_1$ that does not add up is not a total $\Gamma$-leading world.
\end{corollary}
\status{CHECKED (\src{C1} Cor.~C1.8); confirmed at CHECKED (\src{C12\_referee} \S3), one reader, with a repair; the one import is Savage's theorem, as in Theorem~\ref{thm:Q36}(c). Fineness: take $S'\subsetneq S$ with $S\setminus S'$ non-null; then $\min_{\Gamma^{\mathrm{gl}}_1}\nu(S')>0$, and a Lyapunov equipartition of $M$ into $N$ pieces, each carrying $1/N$ of every generator integral, makes every piece $\Gamma^{\mathrm{gl}}$-strictly below $S'$ for large $N$; so every piece is $\le_iS'$ by Theorem~\ref{thm:C17}, and strictly below $S$ by strict (G). (\src{C1} had only ``$\le_iS$'', which is not the strict fineness Savage's theorem needs.)}

\begin{proposition}\label{prop:C19}
Let world $i$ satisfy strict (G) and (IB), and let $\Gamma_{ij},\Gamma_{ik}$ be as above with $\Gamma^{\mathrm{gl}}\ne\{0\}$ and $\Gamma^{\mathrm{gl}}_1$ compact. If every finite family of comparisons of world $i$ between regular models is endorsed by a common $\nu\in\Gamma^{\mathrm{gl}}\setminus\{0\}$, then a single $\nu^*\in\Gamma^{\mathrm{gl}}\setminus\{0\}$ endorses all of them, and world $i$ adds up. Otherwise world $i$ carries a \emph{finite $\Gamma^{\mathrm{gl}}$-obstruction}: comparisons $C_s\preceq_iX_s$ ($s=1,\dots,r$) and weights $c_s\ge0$ with $\sum_sc_s=1$ such that, with $\delta_s\coloneqq1_{E_iX_s}-1_{E_iC_s}$, $\sum_sc_s\,\nu(\delta_s)<0$ for every $\nu\in\Gamma^{\mathrm{gl}}_1$.
\end{proposition}
\status{proved (\src{C1} Prop.~C1.9: the finite intersection property on the compact $\Gamma^{\mathrm{gl}}_1$, positivity of $\nu^*$ from Theorem~\ref{thm:C11}(b), and Sion's minimax theorem for the contrapositive); confirmed (\src{C12\_referee} \S3), one reader, who notes that $\Gamma^{\mathrm{gl}}_1$ is norm-compact whenever $\Gamma_{ij}$, $\Gamma_{ik}$ are, and checks exclusivity and exhaustiveness numerically (\src{c12ref\_py/ref\_c1.py}). In a $\Gamma$-leading world each obstructing comparison is individually $\Gamma^{\mathrm{gl}}$-possible (Theorem~\ref{thm:C17}), while a convex combination is $\Gamma^{\mathrm{gl}}$-strictly negative. \src{C1} adds that the obstructing comparisons cannot be chained and must have crossing supports, as in \src{N17} N17.3; the referee marks this unproved, and it is HEURISTIC here.}

\emph{Open question} (\src{C1} Question~C1.Q). Is there a strictly monotone, cancellative order on the measure algebra of $M$ that is not total (a world with strict (G) and (IB)), that is $\Gamma$-leading for some finitely generated $\Gamma_{ij},\Gamma_{ik}$ with (UI), and that carries a finite $\Gamma^{\mathrm{gl}}$-obstruction? If not, every world that is $\Gamma$-leading for such a pair adds up (Proposition~\ref{prop:C19}, Corollary~\ref{cor:C18}), and the whole family ``credal on the edges, free in the interior'' is excluded from covers of $T_1$ by Theorem~\ref{thm:M0closed}(b). If so, such a world is the only candidate of this kind, and it would still have to satisfy (S), with closed sections for every model, and (C). This is a finite-cancellation question of Kraft--Pratt--Seidenberg type inside a band of fixed finite-dimensional width; it involves no selection, no target and no third world. \src{C12\_referee} \S4 observes that Question~C1.Q and Conjecture~\ref{conj:C2C} below are the same residual seen from two sides (a non-total world with credal edge orders whose interior comparisons are not $\Gamma^{\mathrm{gl}}$-weak), and proposes the single target: under (S) with closed sections for every model, a $\Gamma$-leading world with (F) and (UI) admits no finite $\Gamma^{\mathrm{gl}}$-obstruction.

\emph{Density unanimities on edges, and the statement (LD)} (round 15; \src{C2} \S\S0--4, refereed in \src{C12\_referee} \S\S1--2, one reader). Corollary~\ref{cor:B14} suggests the statement \textbf{(LD)}: in a ranking with every axiom, every world has a leading density on each of its edge regions. With Theorem~\ref{thm:B13} and Theorem~\ref{thm:M0closed}(b) it would exclude every cover of $T_1$ on $\M_U$. (LD) is false, already for $n=2$, where the one edge region is $R_{12}=M$.

\emph{Setting of the next theorem.} Let $n=2$, $K'=[\alpha',\beta']\subseteq(0,1)$ in the weight of world~2, and $M^\pm$, $\sigma$, $\tau$ as in part~(d). Put $x(g)\coloneqq\sigma(g)/\alpha'$ on $M^-$ and $x(g)\coloneqq(\tau(g)-\beta')/(1-\beta')$ on $M^+$, and $C_s\coloneqq D_{[\alpha's,\,\beta'+(1-\beta')s]}$ for $s\in[0,1]$. Then $E_1(C_s)\cap M=\{x<s\}$ modulo null sets, $C_0=P_1$, $C_1=P_2$, and every $C_s$ is coherent, regular and contained in $U$. Fix a positive integrable $\rho_0$ on $U$, and let $\mathcal A$ and $\mathcal B$ be nonempty families of \emph{members}: each $a\in\mathcal A$ is a measure $\nu_a$ equal to $F_a(x)\rho_0\mu$ on $M$ and to $\varphi_a\mu$ on $N_+$, and each $b\in\mathcal B$ a measure $\nu_b$ equal to $G_b(x)\rho_0\mu$ on $M$ and to $\gamma_b\mu$ on $N_+$, where $F_a,G_b:(0,1)\to(0,\infty)$ are measurable, $F_a\rho_0$ and $G_b\rho_0$ are integrable on $M$, and $\varphi_a,\gamma_b>0$ are integrable on $N_+$. Put $X\preceq_1Y$ iff $\nu_a(E_1X)\le\nu_a(E_1Y)$ for every $a\in\mathcal A$, and $X\preceq_2Y$ iff $\nu_b(E_2X)\le\nu_b(E_2Y)$ for every $b\in\mathcal B$; the selection is $p\mapsto C_{p_2}$. Hypotheses: \textbf{(Co)} every ratio $r_{ab}\coloneqq F_a/G_b$ is nondecreasing on $(0,1)$; \textbf{(St)} some $r_{a^*b^*}$ is strictly increasing.

\begin{theorem}\label{thm:C21}
Under (Co) and (St):
\begin{enumerate}[leftmargin=*,label=(\arabic*)]
\item both world orders are preorders on $\M_U$ with strict (G) and (IB);
\item (S) holds: every $C_s$ with $s\in(0,1)$ is competitive on $\{1,2\}$, $C_0$ on $\{1\}$ and $C_1$ on $\{2\}$, and for every $X\in\M_U$ the sections $\{p:C_p\preceq_iX\}$ are closed;
\item every $X\in\M_U$ other than the $C_s$ is strictly dominated by some $C_s$, and each $C_s$ is undominated and tied at both worlds only with itself; hence $\C_L=\{C_s:s\in[0,1]\}$, and (LA) and (C) hold;
\item each $C_s$ with $s\in(0,1)$ is covered, with (LF);
\item world~1 has a leading density on $M$ (even an integrable $\rho\ge0$ that is not a.e.\ zero and satisfies (L1)) iff all $F_a$ are proportional to one another; if two members are not proportional, it has no nonzero leading charge on $M$ either, finitely additive or not; and symmetrically for world~2.
\end{enumerate}
In particular $\mathcal A=\{1+x,\,1+3x^2\}$, $\mathcal B=\{1\}$, $\rho_0\equiv1$ (and all $\varphi_a=\gamma_b=1$) is a full two-world system in which world~1 has no leading density and no nonzero leading charge; with $\mathcal B=\{1,\,2-x\}$ both worlds lack them.
\end{theorem}
\begin{proof}[Proof sketch]
(2) If $C_s\not\preceq_2X$, some $b$ has $\nu_b(A)>\nu_b(c)$, where $A\coloneqq E_1X\cap M$ and $c\coloneqq\{x<s\}$. By (Co) there is $k\ge0$ with $r_{ab}\le k$ on $\{x<s\}$ and $r_{ab}\ge k$ on $\{x\ge s\}$, so $\nu_a(A)-\nu_a(c)\ge k(\nu_b(A)-\nu_b(c))\ge0$ for every $a$, and $C_s\preceq_1X$. The sections are intersections of closed initial (at world~1) or final (at world~2) segments in $s$. (3) The sections of $X$ meet. For $s\in(0,1)$, a model tied with $C_s$ at both worlds equals $C_s$, because $r_{a^*b^*}$ is strictly increasing (here $k=r_{a^*b^*}(s-)<\infty$); for $s=0$ a world-1 tie gives $\nu_a(E_1X)=0$, so $X=P_1=C_0$ by strict (G), and $s=1$ is the mirror at world~2. (5) A leading density vanishing on a non-null set is excluded directly (pairs of sets inside and outside its zero set), and a positive one must be proportional to every member, by the argument of Lemma~\ref{lem:B11}. Two non-proportional members give disjoint error sets $Q$ and $P_e$ that are incomparable at world~1 for a family of small sets $e$; the first property of a leading charge, applied to $(Q,P_e)$ and to $(P_e,Q)$, gives $\pi(P_e)=\pi(Q)$, so $\pi$ vanishes on a non-null set; then $\pi(M)=0$, by Theorem~\ref{thm:S21} or directly, by comparing that set with a piece of $M$ of smaller $\nu_a$-mass and positive $\pi$-mass.
\end{proof}
\status{proved (\src{C2} Thm.~C2.1); confirmed (\src{C12\_referee} \S1), one reader, every axiom re-derived, including closed sections for non-regular models and (LF), with three small repairs incorporated above: the endpoint cases $s\in\{0,1\}$ of the tie argument (G2), leading densities allowed to vanish (G1), and only the first property of a leading charge used (G3); the referee also supplies the proof of $\pi(M)=0$ without Theorem~\ref{thm:S21}. Hypothesis (St) is needed: without it the tie argument, hence (C) and $\C_L=\{C_s\}$, can fail. NUMERICS \src{c2py/ld\_two\_world.py} (800 cells with exact angle masses: the path, competitiveness, domination, closed sections, explicit leading-density violations, and a negative control in which a non-comonotone pair fails competitiveness). The families need not be finite. Both worlds are (CP-a) with countably additive members, so each world adds up, is monotonely continuous (Corollary~\ref{cor:MC}(a)) and has a nested $\C_L$ (Theorem~\ref{thm:MC}); nothing here contradicts the earlier results, and Theorem~\ref{thm:S22} needs a leading charge, which this system lacks. So a world with no leading density and no leading charge, ``incomplete in an essential way'' (unanimity over non-proportional densities, the first alternative of \src{B1} \S3), exists inside a full two-world system; but its worlds add up, so it is no step towards a cover of $T_1$. The ``envy form'' of \src{C2} (C2.2: world~1 the unanimity of $\rho_a$ and $\rho_a+\lambda\rho_2$, i.e.\ of its own loss and of its own loss net of $\lambda$ times world~2's) is a special case.}

Say that \emph{world $i$ is a density unanimity on $R_{ij}$} if there is a family $\mathcal A$ of integrable densities $\rho_a>0$ a.e.\ on $R_{ij}$ such that, for measurable $A,B\subseteq R_{ij}$, $A\le_iB$ iff $\int_A\rho_a\,d\mu\le\int_B\rho_a\,d\mu$ for every $a$; for world $j$ and a family $\mathcal B$ the same, comparing complements in $R_{ij}$. Write $r_{ab}\coloneqq\rho_a/\rho_b$.

\begin{theorem}\label{thm:C23}
Let $n\ge2$ and let a ranking on $\M_U$ (any $K'\subseteq\Di$) satisfy strict (G) at worlds $i$ and $j$, (IB) and (S). Let worlds $i$ and $j$ be density unanimities on $R_{ij}$, with families $\mathcal A$ and $\mathcal B$. Then:
\begin{enumerate}[leftmargin=*,label=(\alph*)]
\item (separation) for every $p\in e_{ij}$ and all $a,b$: $\operatorname{ess\,sup}_{c_p}r_{ab}\le\operatorname{ess\,inf}_{R_{ij}\setminus c_p}r_{ab}$;
\item (sweep) for all $a,b$ and $\lambda_1<\lambda_2$ with $\{r_{ab}=\lambda_1\}$ and $\{r_{ab}=\lambda_2\}$ null, some $p\in e_{ij}$ has $\{r_{ab}<\lambda_1\}\subseteq c_p\subseteq\{r_{ab}\le\lambda_2\}$ modulo null sets;
\item (comonotonicity) for all $a,b,a',b'$, the functions $r_{ab}$ and $r_{a'b'}$ are a.e.\ comonotone on $R_{ij}$;
\item world $i$ has a nonzero leading density on $R_{ij}$ iff all members of $\mathcal A$ are proportional on $R_{ij}$.
\end{enumerate}
\end{theorem}
\status{proved (\src{C2} Thm.~C2.3); confirmed (\src{C12\_referee} \S2), one reader, who notes that the null-level hypothesis is not needed in (b), and that (d) needs the sentence on densities that vanish (as in Theorem~\ref{thm:C21}). (a) If it fails, swapping a small piece of $c_p$ where $r_{ab}$ is high against a piece of its complement where it is low gives an edge test model $X$ with $C_p\not\preceq_iX$ and $C_p\not\preceq_jX$, against competitiveness. (b) The two sets of indices whose $c_p$ contain $\{r_{ab}<\lambda_1\}$, or avoid $\{r_{ab}>\lambda_2\}$, are nonempty and cover the connected edge; if they were disjoint, each would be the trace on the edge of a closed section of an edge test model; so they meet. (c) follows from (a) and (b). Uses: strict (G) at $i$ and $j$, the normal form, competitiveness at open-edge indices, $C_{e_i}=P_i$, $C_{e_j}=P_j$, closed sections at edge test models. Not used: (C), (LA), the third world, nestedness, coherence, totality. So on an edge, competitiveness and closed choice force the edge credal sets to be comonotone, and give a leading member only when the set is a single ray. For $n=2$, Theorem~\ref{thm:C23}(c) is a necessary condition and Theorem~\ref{thm:C21} a sufficient one (\src{C2} Cor.~C2.4, which called this an exact characterisation; the referee notes that the two conditions are stated for different classes, GAP in wording).}

\begin{corollary}\label{cor:C25}
No argument that uses only the restriction of a ranking to one edge (worlds $i$ and $j$, the edge test models, strict (G) and the normal form, competitiveness at edge indices, (K) and closed sections at edge test models) proves (LD), or its case $n=3$. Proposition~\ref{prop:Q35}, Theorem~\ref{thm:S21}, Lemma~\ref{lem:B33} and Theorem~\ref{thm:B35} are arguments of this kind. Any proof of (LD) for $n=3$ must use an input outside this list: the third world, the link class, interior indices, or the shape of coherent edge models in $\R^3$.
\end{corollary}
\status{proved (\src{C2} Cor.~C2.5); confirmed for the listed inputs (\src{C12\_referee} \S2), one reader, whose rewording of the last sentence is used (\src{C2} said ``the link class or the third world'' only; a transported abstract edge system is not a three-world ranking and does not test interior indices or the geometry of coherent edge models). The system of Theorem~\ref{thm:C21} is such an edge system ($R_{12}=M$) with every one of these properties and no leading density, and it can be transported to any $R_{ij}$ by an isomorphism of measure algebras that carries the chain to a chain.}

\begin{proposition}\label{prop:C26}
Let $n=3$ and assume strict (G), (IB) and (S). Let world~1 be a density unanimity on $R_{12}$ and on $R_{13}$ with one family $\mathcal A$, and let world~2 be additive on $R_{12}$ with density $\rho_2$ and world~3 additive on $R_{13}$ with density $\rho_3$. Fix $a\in\mathcal A$ and put $A\coloneqq\rho_2/\rho_a$, $B\coloneqq\rho_3/\rho_a$ (the two edge chain coordinates seen from world~1) and $\psi_b\coloneqq\rho_b/\rho_a$. Let $W\subseteq Q_1$ be non-null, and suppose that the image of $\mu|_W$ under $(A,B)$ is equivalent to Lebesgue measure on a connected open set $D\subseteq(0,\infty)^2$, with projections $I_A$, $I_B$. Then for every $b\in\mathcal A$ there is $\kappa_b>0$ with $\psi_b=\kappa_b$ a.e.\ on $W$, on $\{h\in R_{12}:A(h)\in I_A\}$ and on $\{h\in R_{13}:B(h)\in I_B\}$. In particular, if $I_A$ contains the essential range of $A$ on $R_{12}$ and $I_B$ that of $B$ on $R_{13}$, every $\rho_b$ is proportional to $\rho_a$ on $M$, and $\rho_a$ is a leading density of world~1 on both edge regions.
\end{proposition}
\status{proved (\src{C2} Prop.~C2.6, ``link rigidity''); confirmed (\src{C12\_referee} \S2), one reader. By Theorem~\ref{thm:C23}(c) on $R_{12}$, $\psi_b$ is a function of $A$ on $W$, and by the same on $R_{13}$ a function of $B$; a function of $A$ that equals a function of $B$ a.e.\ on a set where $(A,B)$ has a law equivalent to Lebesgue measure on a connected open set is constant there (Fubini on rectangles); monotonicity extends the constant to the stated regions.}

\begin{corollary}\label{cor:C27}
Under the hypotheses of Proposition~\ref{prop:C26}, suppose that on each sign class the densities $\rho_a$, $\rho_2$, $\rho_3$ are positive constants times $|h_1|$, $|h_2|$, $|h_3|$ respectively (this covers the linear densities and the normalised-box densities of \src{Coverage after Neutral Abstention} on $\M_U$). Then world~1 is additive on $M$ with density $\rho_a$. Hence, under (H$_U$), if worlds~2 and~3 are additive with densities and world~1's edge orders satisfy the hypotheses of Proposition~\ref{prop:C26} with full projections, world~1 adds up (Theorem~\ref{thm:B13}), the ranking is in (M0), and $T_1$ is not covered.
\end{corollary}
\status{proved (\src{C2} Cor.~C2.7); confirmed for the stated densities (\src{C12\_referee} \S2), one reader, the last sentence at the status of Theorem~\ref{thm:M0closed}(b); that the normalised-box densities are of this form was not re-checked by the referee. On $W\coloneqq(-,+,+)\cap U\subseteq Q_1$ the map $(A,B)$ is a submersion onto an open connected set with both projections $(0,\infty)$. The hypotheses say nothing about world~1's comparisons out of interior selected models; Theorem~\ref{thm:B13} determines them. \emph{Why the two-world mechanism stalls in three worlds} (\src{C2} Remark~C2.8; the identity is confirmed, \src{C12\_referee} \S2; its consequences for the edges are a sketch, CHECKED, not proved): members of the form ``own loss minus multiples of the other worlds' losses'' make every interior base-Bayes model Bayes for every member, but on the edges they move the edge weights off the edges unless the link class is priced in an aligned way, and then Theorem~\ref{thm:C23} leaves no competitive edge selection.}

\begin{conjecture}\label{conj:C2C}
Let $n=3$ on $\M_U$ (any $K'\subseteq\Di$), with every axiom and closed sections for every model. If world $i$ is a density unanimity on $R_{ij}$ and on $R_{ik}$, then its members are proportional on $R_{ij}\cup R_{ik}$.
\end{conjecture}
\status{CONJECTURE (\src{C2} Conj.~C2.C). Evidence: Proposition~\ref{prop:C26} and Corollary~\ref{cor:C27} (it holds whenever the link class is priced non-degenerately, in particular for linear and box neighbours) and Remark~C2.8. A counterexample would need aligned link pricing (the pair $(\rho_j/\rho_a,\rho_k/\rho_a)$ degenerate on $Q_i$), interior selected models that are Bayes for every member at once, and coherence of all these Bayes sets. The conjecture says nothing about worlds without a leading density whose edge orders are not density unanimities (band-and-translates, lexicographic with charges, total and non-AR); for those (LD) in three worlds stays open.}

\emph{General decomposable domains: face models at a target weighting} (round 14; \src{B2} \S\S1--6, refereed in \src{B2\_referee}, one reader). In this block $n=3$, $\M$ is a decomposable domain with $D_T\in\M$ for an all-facing triangle $T$ (for instance $T_1$), and (LF) is not assumed. Let $\psi$ be a strictly guiding additive ranking with integral penalties $\ell$ on $\M$, with penalty vectors $c$, $c'$, $d$ as in Section~\ref{sec:setting} (additive rankings and charts). For $x\in\Di$ let $A^*_x(g)$ be the set of advices in $A(g)$ that minimise $x\cdot\ell(g,\cdot)$. A \emph{face model at $x$} is a coherent $D_K\in\M$ whose advice lies in $A^*_x(g)$ for a.e.\ $g$ (the analogue of the $p^*$-face models of \S\ref{sec:M0-closed}). \textbf{($\gamma$)}: for a.e.\ $q\in\Di$ the $q$-Bayes model is a.e.\ unique and coherent. At $x$, a pair whose advice range contains $\{+,0,-\}$ has a \emph{line tie} if exactly two advices are optimal (a $c$-tie, $c\cdot x=0>c'\cdot x$; a $c'$-tie, $c'\cdot x=0>c\cdot x$; or a $d$-tie, $d\cdot x=0<c\cdot x=c'\cdot x$), and an \emph{apex tie} if all three are optimal ($c\cdot x=c'\cdot x=0$ with $c$ and $c'$ not parallel). Let $\mathcal H_T\coloneqq\{m\in\Delta:m\cdot w\ge0$ for the six crossings $w$ of $\partial D_T\}$; for $T_1$ it is $T_1$ together with three caps $\mathrm{cap}_{ij}$ beyond its edges $[v_i,v_j]$ (CHECKED; NUMERICS \src{b2py/geom\_checks.py} G3).

\begin{theorem}\label{thm:B21}
Let $T$ be an all-facing triangle, $D_T\in\M$, $\psi$ as above with ($\gamma$), and let $D_T$ be $\psi$-Bayes at $x\in\Di$.
\begin{enumerate}[leftmargin=*,label=(\alph*)]
\item For each $i$, the model $D^i$ that takes the world-$i$-best advice in $A^*_x(g)$ is coherent, $D^i=D_{K^i}$, and $T\subseteq K^i\subseteq\co(T\cup\{e_i\})$.
\item Every face model $K$ at $x$ has $P_K(w)=0$ at the six crossings $w$ of $\partial D_T$; hence $K\subseteq\mathcal H_T$, and $K$ touches the six tangent lines.
\item Every face model at $x$ contains $T$; so $D_T$ accepts every gamble that some face model accepts, and rejects every gamble that some face model rejects.
\item For a.e.\ $g$, $A^*_x(g)$ is one of $\{+\}$, $\{0\}$, $\{-\}$, $\{+,0\}$, $\{0,-\}$: there are no apex ties and no $d$-ties at $x$.
\end{enumerate}
\end{theorem}
\begin{proof}[Proof sketch]
(a) For a.e.\ direction $\theta$ the Bayes models at $x+s\theta$, $s\downarrow0$, converge to the lexicographic minimiser of $(x\cdot\ell,\theta\cdot\ell)$, which is coherent; $\theta\to e_i-x$ gives $D^i$, and comparison with $D_T$ gives $K^i\subseteq\co(T\cup\{e_i\})$ and a betweenness condition (K2) on gambles on which $T$ gives no advice. (b) Each $K^c$ is pinned at its own crossings (the $c$-th coordinates vanish there) by a two-line argument with no ball lemma. If some face model escapes pinning at a $c$-crossing, then $e_c\in K^c$; this propagates to $e_a\in K^a$ for every $a$; a ball estimate for line ties of one type (the transfer of steps 1--4 of Prop.~F, Lemma~\ref{lem:ball}) collapses each $K^a$ to $\{e_a\}$; and the case $K^a=\{e_a\}$ for all $a$ is impossible, because the arcs of directions $\theta$ with $e_a\in K(\theta)$ would force the no-advice cone of some gamble to meet two different gaps between the vertex directions. Concavity on the facets then gives the arc signs. (c) $T\subseteq K^i$ by pinning and (K2) at $v_i$; if a face model lay strictly on the positive side of a line through $\operatorname{int}T$, some $D^i$ would accept the gamble while $K^i\supseteq T$ straddles it. (d) follows from (c), since $T\subseteq K^i\cap K^j$.
\end{proof}
\status{proved (\src{B2} Thm.~B2.1); confirmed with repairs (\src{B2\_referee} \S1), one reader. Repairs: in Lemma~9 step~3 (the arcs), \src{B2} inferred a lower bound on the complement of an arc from an upper bound; the referee's argument is used: for $g$ with $g_a>0>g_b,g_c$ and generic $\theta$ in the open convex reject cone $R_-(g)$, which contains $\theta_b$, $\theta_c$ and hence the closed short arc between them, $D^\theta$ rejects $g$, so by Fubini the set of $\theta$ with $e_a\in K(\theta)$ misses that arc and a fixed neighbourhood of $\theta_b$ and $\theta_c$; ``one good gamble'' is replaced by ``a.e.\ gamble in an open set''; and the set of such $\theta$ is measurable by Tonelli. In Lemma~7 (the ball estimate), the tied levels are taken as ``good'' levels with essential bounds, and uniformity in the radius follows by chaining over a concentric sub-ball; the statement that no ball centred on a coordinate plane consists of line ties is proved only for line ties of one type, which is all that is used. $D^i\in\M$ follows from Remark~\ref{rem:advicerange}, not from closure in measure. \emph{Scope} (upgraded by the referee, \S2, PROVED): every geometric input (\src{B2}'s G1, G2, G4--G7) holds for every all-facing triangle, and all-facing implies general position; \src{B2} had this only as CHECKED (NUMERICS \src{b2py/random\_triangles.py}, 300 triangles; referee R1, 3000 triangles). Uses only strict guidance, ($\gamma$) and $D_T\in\M$; not (LF), (S), (C), ($\beta$), ($\delta$). It is the two-sided generalisation of Proposition~\ref{prop:crosspin} and Remark~\ref{rem:midarc}. On $\M_U$, \src{B2}'s lemmas with Proposition~\ref{prop:carrying}(c), (d) give an alternative proof of the $p^*\in E$ half of Theorem~\ref{thm:M0closed}(a), after a bridging step through \src{N25} Thm.~4.1(d); it shares the ball lemma and Proposition~\ref{prop:carrying}(c), (d) and is not shorter, contrary to \src{B2} Remark~5.3 (\src{B2\_referee} \S8). What is new on $\M_U$ is an alternative proof of Remark~\ref{rem:midarc} ($K\supseteq T_1$ for every face model).}

\emph{Rankings that add up on a general domain.} \src{B2} \S1 extends Definition~\ref{def:addup} to a decomposable domain: world $i$ \emph{adds up} if some bounded strictly guiding penalty $\lambda_i$ and some finitely additive $\nu_i$, positive on non-null sets of problems whose advice range has more than one element, satisfy $\nu_i(\lambda_i(\cdot,a_B))\le\nu_i(\lambda_i(\cdot,a_X))$ for every relevant comparison $B\preceq_iX$ ($B\in\C_L$, $X$ regular); on $\M_U$ with the error indicator this is Definition~\ref{def:addup}. \emph{Claim B2.0} (\src{B2} Prop.~B2.0): under strict (G), (IB), (LA), (S), (C), $T_1\in\C_L$ and (M0) in this sense, the supporting assessment $\psi_i\coloneqq\int\lambda_i\rho_i\,d\mu$ ($\rho_i$ the density of the countably additive part of $\nu_i$) satisfies ($\alpha^*$) ($D_{T_1}$ is $\psi$-Bayes at some $p^*\in\Di$), ($\beta$), ($\gamma$) and ($\delta$). \emph{Status: not proved.} \src{B2} labels it CHECKED, as a transfer of \src{N13} Prop.~B and \src{N19} Prop.~R3(a); \src{B2\_referee} \S4 finds a GAP (documentation): the general-domain definition is new; the ($\gamma$) step transfers verbatim; ($\beta$) must be written with regular splices for the new definition (the \src{N19} Prop.~R3(a) step); and ($\alpha$) must be re-derived from $T_1\in\C_L$, since \src{N13} cites \src{N11} Cor.~3.2(d), which is for (CP-a). Until Claim~B2.0 is written out, every statement below about rankings on general domains is \emph{conditional on Claim B2.0}. Theorem~\ref{thm:B21} concerns additive fields and does not need it.

\begin{corollary}\label{cor:B22}
Conditional on Claim B2.0: in every ranking on a decomposable domain containing $D_{T_1}$, with strict (G), (IB), (LA), (S), (C) and $T_1\in\C_L$, whose worlds all add up in the sense above, the supporting assessment has at $p^*$ only $c$-ties (orienting each tied pair so that $T_1$ accepts it), $D_{T_1}$ accepts all of them, and every face model at $p^*$ lies between $T_1$ and $\mathcal H_{T_1}$. In particular $p^*$ carries no apex tie, and a cover that breaks apex ties at the target weighting ordinally (the ``fan design'' of \src{B2} \S3) is impossible on every decomposable domain.
\end{corollary}
\status{proved (\src{B2} Cor.~B2.2) from Theorem~\ref{thm:B21}(c), (d), \emph{conditional on Claim B2.0}; confirmed at that status (\src{B2\_referee} \S4), one reader. \src{B2}'s description of the structure a fan would need (\S3) is superseded by Theorem~\ref{thm:B21}(d) and is not claimed.}

\begin{proposition}\label{prop:B213}
Let $\psi$ satisfy the hypotheses of Theorem~\ref{thm:B21} for $T=T_1$ at $x=p^*$, and suppose the ties at $p^*$ are non-null ($p^*\in E$). Then:
\begin{enumerate}[leftmargin=*,label=(\roman*)]
\item the strict-Bayes model at $p^*$ (the unique optimal advice where there is one, no advice elsewhere) is a coherent face model $D_{K^\flat}$, with $T_1\subsetneq K^\flat\subseteq\mathcal H_{T_1}$, so $K^\flat$ enters at least one cap;
\item the ties at $p^*$ are exactly the gambles of $D_{T_1}\setminus D_{K^\flat}$ (oriented), and they form at most three \emph{full corner $c$-atoms}, one for each cap that $K^\flat$ enters: each is an oriented $c$-atom with normal $b$, $b\cdot p^*=0$, with the sign pattern of the adjacent corner ray, carrying the shells over corner arcs of the two adjacent facets, and every face model that accepts it has lower prevision identically $0$ on those arcs, so the atom never straddles $\partial D_{T_1}$;
\item the face models at $p^*$ form a zonotope with at most three generators, whose all-accept point is $T_1$; this point is not extreme iff the normals of the atoms positively span $(p^*)^\perp$, which needs three atoms, one at each corner, and it is then in the interior of the zonotope, but not in general at its centre.
\end{enumerate}
\end{proposition}
\status{proved (\src{B2} Prop.~13); confirmed with repairs (\src{B2\_referee} \S5), one reader. The quantity that vanishes on the arcs is a lower prevision at $p^*$, not the $\pi^+$ of Proposition~\ref{prop:carrying}, which is defined at chord points; ``at its centre'' is corrected to ``in its interior'' (the centre needs $\sum_k\int_{G_k}c\,d\mu=0$); the cap pattern (\src{B2}'s G9) is proved for $T_1$ by the referee; the description in (ii) and (iii) is the referee's. \src{B2}'s remark that this is the only place where ordinal tie-breaking could still matter is heuristic: when $T_1$ is a vertex of the face, the residual is not closed either.}

\emph{The full-atom residual.} Call the configuration of Proposition~\ref{prop:B213} the \emph{full-atom residual} (\src{B2} calls it (R2); it is unrelated to item (R2) of part~(g)). On $\M_U$ it is closed by Proposition~\ref{prop:carrying}(d), which uses that the ties of a carrying atom persist along its whole chord. On a two-sided domain the tie of a gamble $g$ of the atom stops at the apex of $g$ on the chord, where $c'(g)\cdot q=0$, and $d$-atoms on the same chord can add tied levels, so that argument does not transfer (\src{B2} \S5; confirmed, \src{B2\_referee} \S6). It cannot be closed by any argument local at $p^*$:

\begin{proposition}\label{prop:B2ex}
Let $m\in\operatorname{int}\mathrm{cap}_{12}$, $K_0\coloneqq\co(T_1\cup\{m\})$ and $G\coloneqq D_{T_1}\setminus\cl D_{K_0}$; every gamble in $G$ has sign pattern $(-,-,+)$. Let the advice range be $\{+,0\}$ on $G$, $\{0,-\}$ on $-G$, and elsewhere the single advice of $D_{T_1}$ (a decomposable domain, Remark~\ref{rem:advicerange}), and on $G$ put $\ell(g,+)\coloneqq0$ and $\ell(g,0)\coloneqq b\coloneqq(-1,-1,2)$. Then the $q$-Bayes model is $D_{T_1}$ if $q\cdot b>0$ and $D_{K_0}$ if $q\cdot b<0$, with ties on the chord $\{q\cdot b=0\}$, which contains $u$. So ($\gamma$), ($\beta$) and ($\delta$) hold, $D_{T_1}$ is Bayes at $u\in E$, every hypothesis and every conclusion of Theorem~\ref{thm:B21} holds (the face models satisfy $T_1\subseteq K\subseteq K_0\subseteq\mathcal H_{T_1}$), and the configuration is the full-atom residual with one atom. Only (LF) fails (and (C), if $\psi$ is read as a ranking, which is irrelevant for a supporting assessment).
\end{proposition}
\status{proved in a referee report (\src{B2\_referee} \S6), one reader; the pattern claim is also checked numerically (item R4 of the referee's script \src{b2ref\_py/ref\_checks.py}). So the residual can be closed only by an argument that uses (LF) globally, as on $\M_U$, where Lemma~\ref{lem:UVN} supplies vertex neighbourhoods. \src{B2}'s Conjecture~B2.R (``the residual is empty'') is false as stated; the referee restates it as: \emph{under (LF) at $D_{T_1}$, and ($\beta$), ($\delta$), the full-atom residual is empty} (CONJECTURE). A nearby target proposed by the referee (CONJECTURE): Proposition~\ref{prop:carrying}(d) should transfer when, for a.e.\ gamble of the carrying atom, the apex lies outside $\Di$ (automatic if $-$ is not in its advice range) and no other tied gambles have chart lines on the chord, given a two-sided form of frozen support (\src{N25} Thm.~4.1(e)) and Lemma~\ref{lem:UVN} from (LF); this would leave only the genuinely two-sided case. \emph{Round 16:} the restated conjecture holds under (LF) alone, without ($\beta$), ($\delta$) and without Claim~B2.0 (Theorem~\ref{thm:D1}, refereed once).}
\emph{Round 16: (LF) excludes exceptional ties at the target weighting} (\src{D1}; refereed once, \src{D1\_referee}, which found the proof complete after five fixable repairs, all incorporated below). In this block $n=3$, $\M$ is a decomposable domain with $D_{T_1}\in\M$ and (LF) at $D_{T_1}$: there is an open $N\supseteq\partial D_{T_1}\setminus0$ such that for a.e.\ $h\in N$ the advice range contains acceptance of $h$ and silence. Let $\psi$ be a strictly guiding additive ranking on $\M$ with integral penalties and an integrable envelope, with ($\gamma$), and let $E(\psi)\subseteq\Di$ be the set of weightings at which the ties are non-null.

\begin{theorem}\label{thm:D1}
Let $\varepsilon\in(0,1)$, let $\M$ and $\psi$ be as just stated, and let $D_{T_1}$ be $\psi$-Bayes at $p^*\in\Di$. Then $p^*\notin E(\psi)$: the ties at $p^*$ are null, and $D_{T_1}$ is the a.e.\ unique $\psi$-Bayes model at $p^*$. In particular, under (LF) the full-atom residual (Proposition~\ref{prop:B213} with $p^*\in E$) is empty.
\end{theorem}
\begin{proof}[Proof sketch]
Suppose $p^*\in E$. By Proposition~\ref{prop:B213} there is a full corner atom; after relabelling it sits at $r_{12}$, with normal $b$ of sign pattern $(-,-,+)$, $b\cdot p^*=0$, chord $\ell\coloneqq\{b\cdot q=0\}$ cutting off $e_3$, and carrying the shells over arcs $\alpha\subseteq\relint F_1$ and $\alpha'\subseteq\relint F_2$ next to $r_{12}$. (1) \emph{Directional families.} For a direction $\delta$ that is parallel to no atom line and to the $c$-, $c'$- or $d$-line of no non-null set of pairs, and for \emph{every} $q\in\Di$, the lexicographic minimiser of $(q\cdot\ell,\delta\cdot\ell)$ is a coherent model $D_{K^\delta_q}$, the limit of the Bayes models at $q+s\delta$, $s\downarrow0$; $q\mapsto K_q$ is Hausdorff-continuous at every $q\notin E$; the sets $\{q:g\in D_{K^\delta_q}\}$ are convex, so $P^\delta(\cdot,g)\coloneqq P_{K^\delta_\cdot}(g)$ is quasi-concave; (LF) gives vertex neighbourhoods as in Lemma~\ref{lem:UVN} for gambles near $x\in N$ with $x_l>0$ (respectively $x_l<0$); at a chord point $t$ that is neither a fan point nor a point whose pencil is charged by the non-atomic part of the chart directions, the one-sided limits $K^\pm(t)$ exist, and every accepted gamble in $D_{K^+(t)}\sd D_{K^-(t)}$ is a $c$-, $c'$- or $d$-tie on that chord, hence has the strict sign pattern $\pm\sigma'$ of its normal; and $K^\delta_{p^*}$ is a face model, so it contains $T_1$ (Theorem~\ref{thm:B21}(c)). (2) Every Hausdorff cluster value of $K_q$ as $q\to p^*$, $q\notin E$, is a face model, and $K^\pm(t)$ converge to face models along each half-ray of $\ell$. (3) \emph{Two-sided frozen support:} for a.e.\ $t\in\ell$ near $p^*$, $v_1,v_2\in K^+(t)$ and $P_{K^+(t)}\le0$ on $F_1\cup F_2$. All statements ``for a.e.\ $h$ in a compact subarc $\alpha_0$'' are read on the two-dimensional cone $\R_+\alpha_0$, by positive homogeneity. Each shell level over $h\in\alpha_0$ is tied or negated at $t$; by coherence each shell is either entirely accepted or entirely negated by $K^+(t)$, because $K^+(t)$ lies in no line: it is Hausdorff-close to a face model containing $T_1$, by the convergence in (2). A.e.\ shell gamble is tied at $p^*$, so most shells are accepted as $t\to p^*$. On an accepted shell the lower prevision of $K^+(t)$ at $h$ is exactly $0$ (a.e., then by continuity): levels just below $0$ would be $c$-ties or $d$-ties with tie vector $\propto b$, the first impossible at $p^*$ and the second excluded because $K^-(t)$ converges to a face model containing $T_1$. Concavity (Lemma~\ref{lem:caps}(a)) and the support lemma of \src{B2} then give $v_1\in K^+(t)$, and the same on $\alpha'$ gives $v_2$. No (LF) is used here. (4) Hence (TS$_A$) fails on $F_1$ and $F_2$: middle-arc gambles have $P(q,\cdot)\le0$ near $p^*$ off $E$, since otherwise, by (LF) and convexity, the positivity set would reach an interval of the half-ray of $\ell$ on which (3) gives the opposite sign. (5) This gives open pinned sets $O$, with $p^*\in\cl O$ and $v\in K_q$ on $O\setminus E$, at the three crossings $w_{21}=(-1,0,\kappa)$, $w_{31}=(-1,\kappa,0)$ of $F_1$ and $w_{32}=(\kappa,-1,0)$ of $F_2$, from (4), concavity along the facet, and the wedge argument of Lemma~\ref{lem:nopin}. (6) \emph{Walls across chords:} a.e.\ segment from $O$ to $V_l\cap\Di$ contains a point $q\notin E$ with $K_q\subseteq B(v,\eta)$. Off $E$ this is the argument of Lemma~\ref{lem:PL2} (\src{R10b} Lemma~M); at a generic point of an atom line, an exterior collar in $D_{K^o}\setminus D_{K^s}$ would contain a ball centred on a coordinate plane, whose two halves have sign patterns that differ in one coordinate, so at least one of them is not $\pm\sigma'$, against (1). Step~1 of Lemma~M (no $d$-faces of near-facet exterior gambles in the ball) is replaced by the open-chart argument of \src{D1\_referee} \S5: if the open set $R(-h)$ met the ball, it would contain $q\notin E$ with $K_q\subseteq\{m\cdot h<0\}\cap W_v(\mathcal U)\subseteq B(v,\rho)$, against the choice of the ball. (7) \emph{Surround:} by step (iii) of the proof of Theorem~\ref{thm:TSalways} (\src{N28} Lemma~3.4, which needs only a positive distance from $p^*$ and one generic segment per vertex), $p^*$ lies in the convex hull of the sets $\{q\notin E:K_q\subseteq B(v_i,\eta)\}$, $i=1,2$; for $h^3\in\relint F_3$ these lie in the convex set $\{P^\delta(\cdot,h^3)>0\}$, while $P^\delta(p^*,h^3)\le P_{T_1}(h^3)=0$.
\end{proof}
\status{proved (\src{D1} Thm.~D1, Lemmas~D1.1--D1.6); confirmed by one referee (\src{D1\_referee}), one reader. Repairs incorporated: the genericity of chord points excludes both fan points and charged-pencil points, and $c'$-ties are included in the one-sided limits (Lemma~D1.1(d)); the vertex-neighbourhood clauses are stated for $x\in N$; the bookkeeping is done on the cone over $\alpha_0$, not on the arc; in step~(3) the non-collinearity of $K^+(t)$ comes from its convergence to a face model containing $T_1$ (the third item of \src{D1} Lemma~D1.2), not from the diameter bound of the first item, which does not give it; ``a.e.\ shell gamble'' replaces ``a.e.\ gamble of the atom''; the exact zero is proved a.e.\ on the accepted shells and then by continuity; the two geometric facts \src{D1} supported by NUMERICS (GA, GB) have one-line proofs (\src{D1\_referee} \S4); and in step~(6), where \src{D1} invoked \src{R4} Thm.~B off $E$, which says nothing about $d$-faces lying on an atom line, the open-chart argument is used, which also removes the dependence on \src{R4} Thm.~B. The referee confirmed the two steps \src{D1} flagged (the shell dichotomy and exclusion of $d$-ties; the collar ball against the pattern $\pm\sigma'$). \emph{Uses:} strict guidance, an integrable envelope, ($\gamma$), $D_{T_1}$ Bayes at an interior $p^*$, and (LF), only through vertex neighbourhoods at middle-arc gambles of $F_1,F_2$ and at the three crossings, never near the corner rays. \emph{Not used:} Claim~B2.0, ($\beta$), ($\delta$), (LA), (S), (C), Proposition~\ref{prop:carrying}, and the cascade of Theorem~\ref{thm:TSalways} (steps (i), (ii)): the failure of (TS$_A$) on two facets gives walls toward all three vertices directly. So the conjecture restated by \src{B2\_referee} (status of Proposition~\ref{prop:B2ex}) holds under (LF) alone, as a statement about additive fields. Proposition~\ref{prop:B2ex} and Proposition~\ref{prop:D10} show that (LF) cannot be dropped. On $\M_U$ this is an independent proof of the $p^*\in E$ half of Theorem~\ref{thm:M0closed}(a), which does not need frozen support along the whole chord.}

\begin{proposition}\label{prop:D10}
Pick $m_k\in\operatorname{int}\mathrm{cap}_k$ for the three caps of $\mathcal H_{T_1}$, and put $G_k\coloneqq D_{T_1}\setminus\cl D_{\co(T_1\cup\{m_k\})}$. Let the advice range be $\{+,0\}$ on $G_k$, $\{0,-\}$ on $-G_k$, and elsewhere the single advice of $D_{T_1}$, and on $G_k$ put $\ell(g,+)\coloneqq0$ and $\ell(g,0)\coloneqq b_k$, with $b_{12}=(-1,-1,2)$, $b_{23}=(2,-1,-1)$, $b_{31}=(-1,2,-1)$ for the caps beyond the edges $[v_1,v_2]$, $[v_2,v_3]$, $[v_3,v_1]$. Then every Bayes model at every $q\in\Di$ is coherent, $E=\ell_{12}\cup\ell_{23}\cup\ell_{31}$ with $\ell_k\coloneqq\{q\cdot b_k=0\}$, the face models at $u$ form a zonotope with three generators in whose interior $T_1$ lies, each world-$i$-best model rejects exactly the two atoms at the corners adjacent to $v_i$, and the vertex positivity of Lemma~\ref{lem:UVN} fails: the vertex model at $e_2$ contains $v_1$.
\end{proposition}
\status{proved (\src{D1} Example~D1.0 and the $K^i$ structure); confirmed (\src{D1\_referee} \S7), which proves the coherence claim exactly ($D_{K_1}\cap D_{K_2}=D_{\co(K_1\cup K_2)}$) and notes that the $G_k$ are disjoint by their sign patterns; NUMERICS \src{d1py/three\_atom\_example.py}. It is the three-atom version of Proposition~\ref{prop:B2ex}, with the symmetry of $T_1$: (LF) fails, and with it vertex positivity, so every $K_q$ contains $T_1$ and nothing is contradictory.}

\emph{Remark (heuristic).} \src{D1} \S2 notes that at an atom corner the condition (NS) is replaced by a gluing condition (the facet chord images of the two adjacent facets must overlap), and states that this needs $x_k+1/x_{k+1}>\kappa+1$, the inequality of \src{P1}. The note gives no derivation and no script, so this is not claimed; it is not used (\src{D1\_referee} \S7).

\begin{corollary}\label{cor:D17}
Let a ranking on a decomposable domain with (LF) at $D_{T_1}$ satisfy strict (G), (IB), (LA), (S), (C) and $T_1\in\C_L$, with every world adding up in the sense above, and suppose its supporting assessment $\psi$ satisfies ($\alpha^*$) and ($\gamma$) of Claim~B2.0 and has densities $\rho_i>0$ a.e.\ on the problems whose advice range has more than one element. Then $p^*\notin E(\psi)$: $T_1$ is never one of several tied $\psi$-Bayes models at $p^*$, so ordinal tie-breaking at the target weighting has nothing to act on.
\end{corollary}
\status{proved (\src{D1} Cor.~D1.7) from Theorem~\ref{thm:D1}; confirmed (\src{D1\_referee} \S8). It holds at the status of Claim~B2.0 (not proved; a documentation GAP), with a reduced list of imports: of Claim~B2.0 only ($\alpha^*$), ($\gamma$) and the positivity of the densities are used; ($\beta$), ($\delta$) and the regular-splice step of \src{N13} Prop.~B are not. The positivity item is the general-domain form of \src{N19} Lemma~R1, without which $\psi$ is not strictly guiding; the referee adds it to the steps still to be written for Claim~B2.0.}

\emph{The branch $p^*\notin E$} (\src{B2} \S6: an audit of where the routes of Theorem~\ref{thm:gm} use Bayes uniqueness at particular weightings; CHECKED there, corrected by \src{B2\_referee} \S7). The entries for Theorem~\ref{thm:TS}, Propositions~\ref{prop:edgelimits}, \ref{prop:clines} and Lemmas~\ref{lem:forcedlines}, \ref{lem:window} are a reasonable reading: these steps survive a null exceptional set $E$ with $p^*\notin\cl E$ (for forced lines one needs chord points off $E$ converging to the window point, which holds because a line that is not an atom line meets $E$ in a countable set). The corrected conclusion is conditional: \emph{route I1 survives an exceptional set $E$ with $p^*\notin\cl E$, conditionally on (TS), edge confinement and $K_q\to\{e_b\}$ for ($\gamma$)-fields.} Not audited: Theorem~\ref{thm:TSalways} (its walls are defined through $K_q$ at every $q$), Theorems~\ref{thm:cuts} and~\ref{thm:apex}, Proposition~\ref{prop:vertexfree}(i), and the every-$q$ chart form of Proposition~\ref{prop:CF}($\Rightarrow$), which needs Hausdorff continuity of $K_q$ near the edges, and that fails at points of $E$. Moreover the supporting assessment satisfies ($\alpha$)--($\delta$), not (LA), (S), (C), and the audit covers only the uses of Bayes uniqueness. \src{B2}'s summary sentence, that route I1 survives every exceptional set that stays away from $p^*$, is overstated and is not adopted. Quasi-linearity of tie-breaking families (Lemma~\ref{lem:quasilin}), which carried Proposition~\ref{prop:chainchords} on $\M_U$, is not available on two-sided domains (its quasi-convexity half needs half-plane charts), so a two-sided form of Proposition~\ref{prop:chainchords} must first replace it. The case in which $E$ accumulates at $p^*$ is open. \emph{Round 16.} Under (LF), Theorem~\ref{thm:D1} empties the branch $p^*\in E$, so only $p^*\notin E$ remains. \emph{The case $E=\varnothing$} needs an argument too (\src{D1\_referee} \S9): then $\psi$ has Bayes uniqueness at every interior weighting, but Theorem~\ref{thm:gm} is stated for rankings with (LA), (S), (C), so applying it to $\psi$ needs Theorem~\ref{thm:ccibu} (author-checked only) or an audit that its routes use nothing beyond Bayes uniqueness, (TP), (LF), ($\beta$) and ($\delta$) (\src{B2\_referee} \S7 notes that Proposition~\ref{prop:necessary}(b), (c) are used); this case holds at that status. \emph{The case $E\ne\varnothing$:} the chain must survive exceptional chords away from $p^*$, possibly accumulating at $p^*$. The tools of Theorem~\ref{thm:D1} (directional families coherent at every weighting, one-sided limits at chord points, the wall step across chords) are what the audit above lacked. Its step~(6) shows that the wall step crosses chords, given an open convex $G$ with $G\setminus E\subseteq\Gamma_v$ and a pinned set; Theorem~\ref{thm:TSalways} also uses \src{N28} Prop.~3.1 and Lemma~3.2, which need Bayes uniqueness near $p^*$. \src{D1} proposes, as a CONJECTURE, to prove Theorems~\ref{thm:TSalways}, \ref{thm:cuts}, \ref{thm:apex} and Proposition~\ref{prop:vertexfree} for ($\gamma$)-fields with $p^*\notin E$, which would close the case of rankings that add up on every decomposable domain with (LF).

\emph{Consequence for (M0) on general decomposable domains.} Conditional on Claim B2.0, a ranking on a decomposable domain containing $D_{T_1}$ whose worlds all add up covers $T_1$ only if its supporting assessment either has $p^*\in E$ and is in the full-atom residual, which needs (LF) to be closed (Proposition~\ref{prop:B2ex}), or has $p^*\notin E$ with $E$ accumulating at $p^*$ or with one of the unaudited inputs of route I1 failing for ($\gamma$)-fields. \emph{Round 16:} under (LF) the first alternative is empty (Theorem~\ref{thm:D1}, independent of Claim~B2.0; for the ranking, Corollary~\ref{cor:D17}, at the status of Claim~B2.0 with the reduced list of imports), so what remains under (LF) is the branch $p^*\notin E$: $E=\varnothing$ at the status of Theorem~\ref{thm:ccibu}, and $E\ne\varnothing$ conditionally on the unaudited inputs of route~I1.
\emph{Round 16: what Theorem~\ref{thm:M0closed} needs, and the mechanism of Theorem~\ref{thm:D21} at $T_1$} (\src{E1}; \textbf{proved in a note}; audited in note \src{F2}, which writes out every hypothesis and corrects the scope, and the audit was refereed once in \src{F12\_referee}; every result of this block carries that status, and where \src{E1}'s own summaries overstate, the corrected scope is given). Proposition~\ref{prop:F21} is proved in \src{F2}. Assume (H$_U$). For a shadow $\nu=(\rho_i\mu)_i$ (Notation of this subsection) put $c^\nu_i\coloneqq\rho_i\sigma_i$, $B^\nu_t\coloneqq\{g\in U:t\cdot c^\nu(g)>0\}$ (its $t$-Bayes model) and $Z^\nu_t\coloneqq\{g\in U:t\cdot c^\nu(g)=0\}$. Call $\nu$ \emph{admissible} if $D_{T_1}=B^\nu_{p^*}\cup S$ for some $p^*\in\Di$ and some measurable $S\subseteq Z^\nu_{p^*}$. A \emph{dominator comparison} of $\nu$ at $t$ is a comparison $C\preceq_kB^\nu_t$ in which the coherent model $C$ satisfies $C\preceq_kB^\nu_t$ at every world $k$.

\begin{theorem}\label{thm:E11}
Assume (H$_U$), let $\nu$ be an admissible shadow, and put $I_\nu\coloneqq\{t\in\Di:Z^\nu_t\text{ is null and }B^\nu_t\text{ is incoherent}\}$.
\begin{enumerate}[leftmargin=*,label=(\alph*)]
\item $I_\nu$ is not Lebesgue-null.
\item For every $t\in I_\nu$, Theorem~\ref{thm:kkm} gives an undominated selected $C$ with $C\preceq_kB^\nu_t$ at every $k$; and for every coherent $C$ with this property some world $j$ has $\nu_j(E_jC)>\nu_j(E_jB^\nu_t)$, so the dominator comparison $C\preceq_jB^\nu_t$ is not endorsed by $\nu$.
\item If some admissible $\nu$ endorses, at every world, the dominator comparisons supplied by Theorem~\ref{thm:kkm} at a.e.\ $t\in\Di$, then (H$_U$) is contradictory.
\end{enumerate}
\end{theorem}
\begin{proof}[Proof sketch]
(a) For each $g$ the set $\{t:t\cdot c^\nu(g)=0\}$ is a line, so $Z^\nu_t$ is null for a.e.\ $t$ (Tonelli). If $I_\nu$ were null, $B^\nu_t$ would be coherent for a.e.\ $t\in\Di$, which is ($\gamma$) for the strictly signed field $c^\nu$; with admissibility, Theorem~\ref{thm:M0closed}(a) excludes this. (b) A coherent $C$ is regular. Since $Z^\nu_t$ is null, $B^\nu_t$ is the a.e.\ unique pointwise minimiser of $t\cdot\psi^\nu$ over all models, $\psi^\nu_k(X)\coloneqq\nu_k(E_kX)$; if every comparison $C\preceq_kB^\nu_t$ were endorsed, $t\cdot\psi^\nu(C)\le t\cdot\psi^\nu(B^\nu_t)$, so $C=B^\nu_t$ a.e., against incoherence. Theorem~\ref{thm:kkm} uses (LA) and (S). (c) Under that endorsement, (b) makes $I_\nu$ null, against (a).
\end{proof}
\status{proved in a note (\src{E1} Thm.~E1.1, from Theorem~\ref{thm:M0closed}(a) and Theorem~\ref{thm:kkm}); re-derived in the audit \src{F2} \S\S1.2--1.3, confirmed in \src{F12\_referee} \S5.1, one reader. \emph{Hypotheses:} (H$_U$) and an admissible shadow, integrable as in the Notation; (a) uses no integrability, and (b) uses that $\psi^\nu$ is finite. \emph{What it says.} The argument of Theorem~\ref{thm:M0closed} uses only the dominator comparisons $C\preceq_kB^\nu_t$ of a single admissible shadow, at all three worlds, for $t$ in a set of full measure in $\Di$, and only the dominators supplied by Theorem~\ref{thm:kkm}. $T_1$'s own out-comparisons are never used: admissibility is obtained by \emph{choosing} $\nu$, and admissible shadows exist whatever the ranking (for $\nu^L$ of Proposition~\ref{prop:Q31}, $B_t=D_{T_1}$ with null ties on $Q_L$), whereas Proposition~\ref{prop:M0} derives ($\alpha$) from endorsement of $T_1$'s out-comparisons, because there the shadow is extracted from the ranking. If every world adds up with countably additive weights, these weights form an admissible shadow that endorses everything (Proposition~\ref{prop:M0}($\alpha$), via \src{N17} Prop.~2.1), so (c) applies; for purely finitely additive weights Theorem~\ref{thm:M0closed}(b) keeps its own route (\src{N13} Prop.~B), and (c) is a different sufficient condition, not a restatement. $I_\nu$ cannot be confined to a compact subset of $\Di$ in general ($I_{\nu^L}\cap Q_L=\varnothing$, Proposition~\ref{prop:Q32}). If two worlds $j$, $k$ have densities $\rho_j$, $\rho_k$ that endorse all their comparisons (for instance leading densities, Lemma~\ref{lem:B10}), taking $a\rho_j$, $b\rho_k$ as two components of $\nu$ puts every failure in (b) at the third world $i$, \emph{provided} the triple $(\nu_i,a\rho_j,b\rho_k)$ is admissible for some $\nu_i$ and constants $a,b>0$. That needs $a\rho_j/(b\rho_k)$ bounded below on the set where $T_1$ is correct only at $j$, and symmetrically at $k$; raising $\nu_i$ cannot help there, because $T_1$ errs at $i$. Such a completion is not automatic (correction of 15 October; \src{F2} \S1.3, \src{F12\_referee} \S5.1).}

\begin{theorem}\label{thm:E12}
Assume (H$_U$). For every $\eta>0$ there are an admissible shadow $\nu'$, a weight $t\in\Di$ with null ties at which $B^{\nu'}_t$ is incoherent and $\mu(B^{\nu'}_t\sd D_{T_1})<\eta$, a selected undominated $C$ with $C\preceq_kB^{\nu'}_t$ at every world $k$, and a world $j$ at which $C\preceq_jB^{\nu'}_t$ is not endorsed by $\nu'$.
\end{theorem}
\begin{proof}[Proof sketch]
Start from $\nu^L$ ($L>2$) and the cone $C_\eta$ of the proof of Theorem~\ref{thm:R3}, and multiply the weights of the worlds at which $T_1$ is correct by $M\coloneqq2/\min_it_i+1$ off $C_\eta$; call the result $\nu'$. It stays admissible with $p^*=u$: off $C_\eta$, $LMu_i\ge LM/3>2/3$, and on $C_\eta$ some world is correct and $L/3>2/3$, so $B^{\nu'}_u=D_{T_1}$ with null ties. At a generic $t$ near $e_j$ with $w_j<0$, $B^{\nu'}_t=D_{T_1}$ off $C_\eta$, and a non-null part of $C_\eta$ is rejected, so $B^{\nu'}_t$ is incoherent and $\eta$-close to $T_1$, as in Theorem~\ref{thm:R3}. Theorem~\ref{thm:kkm} gives $C$, and Theorem~\ref{thm:E11}(b) gives the unendorsed comparison.
\end{proof}
\status{proved in a note (\src{E1} Thm.~E1.2); re-derived in \src{F2} \S1.3, confirmed in \src{F12\_referee} \S5.1, one reader; it needs no hypothesis on any world beyond (H$_U$); NUMERICS \src{e1py/e1\_bathtub.py} ($\varepsilon=0.3$, $s=0.9$, $L=5$; convexity violated at 1\,798 of 1\,800 removed points tested, the other two not needed). It is the conclusion of case~(ii) of Theorem~\ref{thm:R3} without its case split, with admissible shadows in place of shadows in $\mathcal N$. Theorem~\ref{thm:R3} needed the split only to keep $\nu'$ in $\mathcal N$, and by Theorem~\ref{thm:E11} that is unnecessary; the dominator may be $T_1$ itself. So a world that does not add up breaks the argument of Theorem~\ref{thm:M0closed} only if, for \emph{every} admissible $\nu$, some dominator comparison into the incoherent $B^\nu_t$ is unendorsed on a non-null set of weights, and these incoherent models may be taken arbitrarily close to $T_1$ in measure. Leaving $T_1$'s own out-comparisons unendorsed (case~(i) of Theorem~\ref{thm:R3}) breaks nothing.}

\emph{Ends and protection.} An \emph{end} is a decreasing sequence $S=(S_m)$ of non-null measurable subsets of $M$ with $\mu(\bigcap_mS_m)=0$. A $\{-1,0,1\}$-valued $\delta$ is \emph{$S$-positive} if $\mu(\{\delta=1\}\cap S_m)>0$ for every $m$, and \emph{eventually $\ge0$} (respectively \emph{$=0$}) on $S$ if $\delta\ge0$ (respectively $\delta=0$) a.e.\ on some $S_m$. World $i$ is \emph{end-protected} with a countable family $\mathcal S_i$ of ends if, for every relevant comparison $C\preceq_iX$, the vector $1_{E_iX}-1_{E_iC}$ is $S$-positive for some $S\in\mathcal S_i$ or eventually $\ge0$ on every $S\in\mathcal S_i$. World~1 of Theorem~\ref{thm:D21} is of this form, with the single end $S_m\coloneqq M\cap\{x\ge1-2^{-m}\}$: a relevant comparison has $D(\delta)=0$, and either $\delta\in\mathcal K_0$, which is eventually $\ge0$ by the near-one clause, or $\delta-\gamma_*\in\mathcal K_0$, so that $\delta\ge1_{Q_\gamma}$ near $x=1$, which is $S$-positive by (d2) (\src{E1} \S2.1, from Lemma~\ref{lem:D23}). ``Targets that meet at every scale'' is $S$-positivity at an end.

\begin{lemma}\label{lem:E13}
Assume (H$_U$).
\begin{enumerate}[leftmargin=*,label=(\alph*)]
\item Let worlds $2$ and $3$ have densities $\rho_2,\rho_3$ that endorse all their comparisons between regular models, and let world~1 be end-protected with family $\mathcal S_1$. Then for every $S\in\mathcal S_1$, $\mu(S_m\cap E_1(T_1))=0$ for some $m$.
\item Assume moreover (Nest$_e$) at world~1 on $e_{12}$ and $e_{13}$. Then for every non-null $A\subseteq M$, world~1 has a relevant comparison whose vector is $-1$ on a non-null subset of $A$. So no non-null fixed region has all relevant world-1 comparisons $\ge0$ on it.
\end{enumerate}
\end{lemma}
\begin{proof}[Proof sketch]
(a) Suppose $\mu(S_m\cap E_1(T_1))>0$ for all $m$. Let $R_0\coloneqq(+,-,-)\cap D_{T_1}$, which is non-null (it contains a neighbourhood of a point just inside $D_{T_1}$ near $r_{23}$). Choose depths for the countably many members of $\mathcal S_1$ so that $R$, $R_0$ minus their deep parts, is non-null, and let $p\subseteq\bigcup_{m\ge m_0}(S_m\cap E_1(T_1))$ have non-null pieces at every depth and $\rho_2(p)+\rho_3(p)<\min(\rho_2(R),\rho_3(R))$. Then $X\coloneqq T_1\sd(p\cup R)$ is regular and strictly better than $T_1$ for $\rho_2$ and $\rho_3$, so by endorsement $T_1\not\preceq_2X$ and $T_1\not\preceq_3X$; $T_1$ is competitive on $\{1,2,3\}$ (through Theorem~\ref{thm:kkm} it is tied at every world with a selected model), so $T_1\preceq_1X$, a relevant comparison with vector $1_R-1_p$, which is neither eventually $\ge0$ on $S$ nor $S'$-positive for any $S'\in\mathcal S_1$. (b) $M=R_{12}\cup R_{13}$; say $\mu(A\cap R_{12})>0$. By Proposition~\ref{prop:Q35}(1) the edge error sets $c_b$ increase continuously to $R_{12}$, so some interior $b$ has $\mu(A\cap c_b)>0$ and $d_b$ non-null; trading a non-null piece of $A\cap c_b$ for a set $R\subseteq d_b$ of larger $\rho_2$-mass gives a model strictly $\rho_2$-better than the selected edge model, and competitiveness on $\{1,2\}$ gives the comparison.
\end{proof}
\status{proved in a note (\src{E1} Lemma~E1.3); re-derived in \src{F2} \S1.3, confirmed in \src{F12\_referee} \S5, one reader; NUMERICS (\src{e1py/e1\_heavy.py}: an instance of (a)). \emph{Hypotheses} (\src{F2} \S1.2). Part~(a) uses only the comparisons out of $T_1$ at worlds~2 and~3, so it suffices that $\rho_2\mu$ and $\rho_3\mu$ endorse those. With integrable densities, as in the Notation, nothing more is needed. \src{F2} relaxes integrability to a.e.\ finiteness; that relaxation is valid only if endorsement is read in the cancellation form, or if $\rho_k\mu(E_kT_1)<\infty$ for $k=2,3$, because otherwise $T_1\preceq_kX$ compares two infinite costs and is endorsed vacuously (\src{F12\_referee} \S5.2). The argument of (a) works at any selected $C_p$ competitive on $\{1,2,3\}$ with $C_p\cap(+,-,-)$ non-null, and on $e_{12}$ it uses world~2 alone. \emph{What (a) proves about location} (corrected 15 October). Under its hypotheses, every end of world~1 is eventually $T_1$-correct at world~1; nothing more. The round-16 status went on to say that such an end can sit only at limits of near-unanimous problems (for world~1, at gambles approaching the $e_1$ ray inside $(+,-,-)$, or the face $\{h_2=0\}$ inside $(+,-,+)$), and part~(f) said that protection cannot be moved off the near-unanimous limits. That is unproved, and the proved constraint admits more (\src{F2} \S1.4, \src{F12\_referee} \S5.3). $T_1$ rejects every gamble on a collar of $\partial U$, so world~1 is correct there wherever $h_1<0$ (Proposition~\ref{prop:F21}(a)); $T_1$-correct regions also include the correct side of $\partial D_{T_1}$ and interior regions. Under (Nest$_e$) on both edges, with the hypotheses of Proposition~\ref{prop:F21}(b), the ends lie at the near-unanimous boundary or at $\partial U\cap\{h_1<0\}$. No three-world protected world with an end at $\partial U$ has been constructed; in two worlds, the end $x\to1$ of Theorem~\ref{thm:D21} contains a ray of $\partial U$ (Lemma~\ref{lem:F11}), and $\mathrm{D2}'_{\mathrm B}$ has its whole protected set at $\partial U$ (Theorem~\ref{thm:F16}). Part~(b) excludes fixed protected regions only; the same construction, with pieces at every depth, gives Proposition~\ref{prop:F21}(b).}

\begin{proposition}\label{prop:F21}
Assume (H$_U$).
\begin{enumerate}[leftmargin=*,label=(\alph*)]
\item There is $\eta_0>0$ such that $D_{T_1}\cap\{h:P_{K'}(h)<\eta_0|h|\}$ is null. So on this collar of $\partial U$ the model $T_1$ rejects every gamble, and it is correct at world $i$ a.e.\ on the part of the collar where $h_i<0$. Hence every end of world $i$ that eventually lies in $\{h:P_{K'}(h)<\eta|h|,\ h_i<0\}$ for some $\eta\le\eta_0$ satisfies the hypothesis of Lemma~\ref{lem:E14}, with no hypothesis on any world.
\item Suppose further: (Nest$_e$) at world~1 on $e_{12}$ and on $e_{13}$; densities $\rho_2,\rho_3>0$, finite a.e., such that $\rho_2\mu$ endorses the world-2 comparisons $C_p\preceq_2X$ with $p$ in the open edge $e_{12}$ and $X$ regular, and $\rho_3\mu$ the world-3 comparisons $C_p\preceq_3X$ with $p$ in the open edge $e_{13}$, endorsement being read in the cancellation form or with $\rho_k\mu(E_kC_p)<\infty$; and world~1 end-protected with a countable family $\mathcal S_1$. Then for every end of world~1 and every $\eta>0$ its pieces eventually lie: in $(+,-,-)$, within angle $\eta$ of the $e_1$ ray; in $(+,-,+)$, within angle $\eta$ of the plane $\{h_2=0\}$, and in $(+,+,-)$, within angle $\eta$ of $\{h_3=0\}$; in $(-,+,+)$, $(-,+,-)$ and $(-,-,+)$, within angle $\eta$ of $\partial U$. So the ends lie at the near-unanimous boundary or at $\partial U\cap\{h_1<0\}$, and $T_1$ is eventually correct at world~1 at both.
\end{enumerate}
\end{proposition}
\begin{proof}[Proof sketch]
(a) $\cl D_{T_1}\setminus0\subseteq U$ (\S\ref{sec:aos}), and $P_{K'}$ is continuous and positive on the compact set $\cl D_{T_1}\cap\{|h|=1\}$, so its minimum $\eta_0$ there is positive. (b) First, for every end $S$, every $b<1$ and all large $m$, $S_m\cap c_b$ is null, where $c_b$ are the world-1 error sets of the edge selection on $e_{12}$: otherwise, trading pieces of $S_m\cap c_b$ at every depth for a set $R\subseteq d_b$ off the deep ends with larger $\rho_2$-mass gives a model $X$ strictly $\rho_2$-better than $C_p$, so $C_p\not\preceq_2X$, and competitiveness on $\{1,2\}$ gives $C_p\preceq_1X$ with a vector that violates end-protection. Second, as $p\to e_2$ the sets $c_b$ fill $R_{12}$ (Proposition~\ref{prop:Q35}(1)) and $K_p\to\co(K'\cup\{e_2\})$ in Hausdorff distance, so $d_b\subseteq\{m_2<\eta|\cdot|\}$ for $b$ near $1$, with $m_2(h)\coloneqq\min(P_{K'}(h),h_2)$ if $h_1<0<h_2$ and $m_2(h)\coloneqq-h_2$ if $h_1>0>h_2$. The same holds on $e_{13}$. Combining the two constraints sign class by sign class gives the list; in $(-,+,-)$ and $(-,-,+)$ a small $h_2$ (respectively $h_3$) with the other coordinates negative forces $P_{K'}(h)<0$, so the end must approach $\partial U$.
\end{proof}
\status{proved (\src{F2} Props.~F2.1, F2.2, Cor.~F2.3); confirmed (\src{F12\_referee} \S5.3), one reader, subject to the reading of endorsement stated in (b), which is the referee's repair of \src{F2}'s relaxation to a.e.\ finiteness (\src{F12\_referee} \S5.2). Part~(b) imports the Hausdorff convergence of \src{R1} Prop.~3.1, as used in Lemma~\ref{lem:quasilin}(b): CHECKED, not re-proved. Part~(a) covers only ends that eventually lie in a collar of $\partial U$ inside $\{h_i<0\}$; an end approaching $\partial U\cap\{h_i=0\}$ is not covered. Without (Nest$_e$), only Lemma~\ref{lem:E13}(a) and its variant constrain the ends, and nothing proved excludes ends shrinking to the $T_1$-correct side of $\partial D_{T_1}$, to interior rays, or to mixed directions on the planes $\{h_2=0\}$, $\{h_3=0\}$ in the interior of $U$, such as $(-1,0,b)$ with $b\gtrsim2.15$ (NUMERICS \src{f2py/f2\_ends\_T1.py}, $\varepsilon=0.3$, $s=0.9$; within margin $2\cdot10^{-3}$ of $\partial U$, $T_1$ accepts nothing in any sign class).}

\begin{lemma}\label{lem:E14}
Let $\nu$ be admissible, $D_{T_1}=B^\nu_{p^*}\cup S$ with $S\subseteq Z^\nu_{p^*}$, and for each world $i$ let $\mathcal S_i$ be a countable family of ends with $S_m\cap E_i(T_1)$ null for large $m$. Then there is an admissible $\nu'$ with $\nu'_i=\nu_i(1+f_i)$, $f_i\ge0$ supported in the ends of world $i$, such that for every $t\in\Di$, every $i$ and every $S\in\mathcal S_i$, $B^{\nu'}_t$ is correct at world $i$ a.e.\ on $S_m$ for all $m\ge m(S,t)$. Consequently no vector $1_{E_iB^{\nu'}_t}-1_{E_iC}$, for any model $C$, is $S$-positive for any $S\in\mathcal S_i$.
\end{lemma}
\begin{proof}[Proof sketch]
Enumerate the pairs $(S,l)$ with $S\in\bigcup_i\mathcal S_i$, $l\in\mathbb N$, choose depths $m_l(S)$ beyond which $S\cap E_i(T_1)$ is null and $(\rho_j+\rho_k)(S_{m_l(S)})\le2^{-r(S)}l^{-3}$ ($r(S)$ the index of $S$), and put $f_i\coloneqq\sum_{S\in\mathcal S_i}\sum_ll\,(\rho_j+\rho_k)\rho_i^{-1}1_{S_{m_l(S)}}$; then $\rho_if_i$ is integrable. On $\operatorname{supp}f_i$, $T_1$ is correct at world $i$ a.e., so raising $\rho_i$ there raises the margin at $p^*$ and keeps $\nu'$ admissible. Deep in $S$ the weights of the worlds at which $T_1$ errs are unchanged, and $t\cdot c^{\nu'}(g)\sigma_i(g)\ge(t_il-1)(\rho_j+\rho_k)>0$ once $l>1/t_i$.
\end{proof}
\status{proved in a note (\src{E1} Lemma~E1.4); re-derived in \src{F2} \S1.3, confirmed in \src{F12\_referee} \S5.1, one reader; NUMERICS (\src{e1py/e1\_heavy.py}: on the $(+,-,-)$ end near the $e_1$ ray, unweighted shadows reject 7\,980\,080 sampled points over random $t$, the heavy reweighting none at depth $>1/t_1$, and $B_u=T_1$ is kept; \src{f2py/f2\_ends\_T1.py}~(3): at a $\partial U$ end in $(-,+,+)$, admissibility is kept and there are no world-1 errors deep in the end). \emph{Location-free} (\src{F2} \S\S1.2, 2.1): the lemma uses no hypothesis on the rankings, only that the ends are eventually $T_1$-correct at their world. For ends in a collar of $\partial U$ inside $\{h_i<0\}$ this holds with no hypothesis on any world (Proposition~\ref{prop:F21}(a)). So moving the protection to $\partial U$ does not make it visible: for every admissible $\nu$ the reweighting is admissible, its Bayes models are correct deep in every such end, and no dominator comparison into them is end-positive. At $\partial U\cap\{h_i>0\}$, $T_1$ errs at world~$i$; that location is excluded by Lemma~\ref{lem:E13}(a) under its hypotheses, and by Proposition~\ref{prop:F21}(b).}

\begin{theorem}\label{thm:E15}
Assume (H$_U$). Suppose each world $i$ is end-protected with a countable (possibly empty) family $\mathcal S_i$, and there are integrable densities $\rho_i>0$ such that:
\begin{itemize}[leftmargin=*]
\item[\textbf{(b)}] every relevant comparison of world $i$ whose vector is eventually $=0$ on every $S\in\mathcal S_i$ is endorsed by $\rho_i\mu$ (the \emph{base adds up});
\item[\textbf{(h)}] for every $S\in\mathcal S_i$, $\operatorname{ess\,inf}_{S_m}\rho_i/(\rho_j+\rho_k)\to\infty$ as $m\to\infty$ (the density is \emph{heavy at its ends}).
\end{itemize}
Then there is a contradiction. So no cover of $T_1$ on $\M_U$ has this form.
\end{theorem}
\begin{proof}[Proof sketch]
Let $\psi\coloneqq(\rho_i\mu)_i$. (1) $\psi$ is admissible: if a regular $X$ had $\psi_k(X)<\psi_k(T_1)$ at every $k$, let $Y$ agree with $X$ except deep in every end, where it agrees with $T_1$, with depths so large that $Y$ is still strictly better at every world; competitiveness (through Theorem~\ref{thm:kkm}, as in Lemma~\ref{lem:E13}(a)) gives $T_1\preceq_kY$ for some $k$, a comparison whose vector is eventually $0$ on every end of world $k$, so (b) gives $\psi_k(T_1)\le\psi_k(Y)$, a contradiction. So $\psi(T_1)$ is weakly Pareto in the convex compact (Lyapunov) range, a supporting weight $p^*$ exists, and it is interior, since at a weight on $\partial\Delta$ the non-null set of $(-,-,+)$-gambles that $T_1$ accepts near $r_{12}$ (and its cyclic images) would contradict pointwise minimality. (2) Let $t\in\Di$ have null ties and $C\preceq_iB^\psi_t$ at every $i$. By (h), $B^\psi_t$ is correct at world $i$ deep in every $S\in\mathcal S_i$ (the computation of Lemma~\ref{lem:E14} with $f=0$), so the vector is $\le0$ there and $S$-positive for no $S$; end-protection makes it eventually $\ge0$, hence eventually $=0$; by (b) it is endorsed. (3) Theorem~\ref{thm:E11}(c) gives the contradiction.
\end{proof}
\status{proved in a note (\src{E1} Thm.~E1.5); re-derived in \src{F2} \S1.3 with no gap found, confirmed in \src{F12\_referee} \S5.1, one reader. \emph{Exact hypotheses} (\src{F2} \S1.1): (H$_U$); every world end-protected with a countable, possibly empty, family of ends; and integrable $\rho_i>0$ a.e.\ with (b) and (h) for the triple. Condition~(h) involves the other worlds' densities, so it is not a property of world $i$ alone. Admissibility of $\psi$ is proved in step~(1), not assumed. The theorem uses neither Lemma~\ref{lem:E13} nor any location of the ends; it does use integrability (the Lyapunov range and finite Bayes risks). So it excludes exactly the covers in which every world is an end-protected decoration with integrable, jointly end-heavy base densities. It does not apply when some world's base has no integrable endorsing density that is heavy at its ends relative to integrable endorsers of the other worlds, or when some world fails (b) or end-protection, for instance a second world that fails to add up (\src{F2} \S\S2.2, 4.2). \src{E1}'s ``unconditional general form'' (in any cover, for every admissible $\nu$, the end-heavy $\nu'$ has unendorsed dominator comparisons, none end-positive) holds only when every end is eventually $T_1$-correct at its world, the hypothesis of Lemma~\ref{lem:E14} (\src{F2} \S1.2). \emph{Scope} (\src{E1}): (b) and (h) concern one base density that is heavy at its ends; Lemma~\ref{lem:E14} alone does not reweight a base that adds up only with a density that is not heavy, because (b) is not inherited by $\rho_i(1+f)$. A world with $\mathcal S_i=\varnothing$ satisfies the hypotheses iff an integrable positive density endorses all its relevant comparisons. Additive bases of gauge-mixture type (\src{R10a} Thm.~M) satisfy (h) at the $e_1$-ray end, where the other two coordinates of the penalty vector tend to $0$ while the first stays bounded away from $0$ (CHECKED by inspection). So ``an additive near-cover plus modifications of the type of Theorem~\ref{thm:D21}'' falls under this theorem when its base densities satisfy (b) and (h) jointly. \emph{Round 17} (part~(f$'$)). (h) can be replaced by eventual $T_1$-correctness of every end at its world (Theorem~\ref{thm:H13}), and, for finite families of ends, dropped unless the conflict digraph has a directed cycle (Proposition~\ref{prop:HR1}). Next to two worlds with integrable full endorsers, an a.e.\ finite endorser of one world's base in the cancellation form suffices for the exclusion (Theorem~\ref{thm:H12}).}

\begin{proposition}\label{prop:E16}
In the system of Theorem~\ref{thm:D21}, let $w:[0,1)\to[1,\infty)$ be nondecreasing and integrable, with $w(x)\to\infty$ as $x\to1$, and let $\nu'$ be $w(x)\bar\mu$ on $M$ and $\bar\mu$ on $N_+$ at world~1, and $\bar\mu$ at world~2. Then:
\begin{enumerate}[leftmargin=*,label=(\roman*)]
\item $\nu'$ endorses every relevant world-1 comparison that is not $S$-positive for the end $S$ above;
\item for a.e.\ $t$, $B^{\nu'}_t=C_{x_t}$ with $x_t\coloneqq\inf\{x:t_1w(x)>t_2\}$, a selected model.
\end{enumerate}
So world~1 of Theorem~\ref{thm:D21} satisfies (b) and (h) of Theorem~\ref{thm:E15} in its two-world system: its failure to add up sits entirely in $S$-positive comparisons, which never meet the Bayes models of this shadow.
\end{proposition}
\begin{proof}[Proof sketch]
(i) Such a comparison has $\delta\in\mathcal K_0$ with $D(\delta)=0$; with $\tau(t)\coloneqq\int_{\{x\ge t\}}\delta\,d\bar\mu\ge0$ and $\tau(1)=0$, summation by parts gives $\int w\delta\,d\bar\mu=w(0)\tau(0)+\int\tau\,dw\ge0$, and for unbounded $w$ one passes to the limit through $\min(w,n)$. (ii) On $M$, $\sigma_1=-\sigma_2$, so $t\cdot c(g)=\sigma_1(g)(t_1w(x)-t_2)$; thus $E_1(B^{\nu'}_t)\cap M=\{w<t_2/t_1\}$, an initial segment, and $N_+$ is accepted; ties are non-null only at the countably many values $t_2/t_1$ at which $w$ has a flat piece; a regular model is determined by its world-1 error set.
\end{proof}
\status{proved in a note (\src{E1} Prop.~E1.6); re-derived in \src{F2} \S1.3, confirmed in \src{F12\_referee} \S5.1, one reader. It is a two-world result and says nothing about a three-world installation. It concerns the base alone: by Theorem~\ref{thm:F13}, endorsing base and generators together needs infinite mass near $\R_+(1,0)$, but the base has the integrable end-heavy endorser of (i), whatever the placement of the targets near $x=1$, the $\partial U$ part included; so Theorem~\ref{thm:F13} does not touch Theorem~\ref{thm:E15} (\src{F2} \S2.3). NUMERICS in an exact cell model (script \src{e1\_d2\_invisible.py} in \src{e1py}: no negative $w$-sums over 20\,000 elements of $\mathcal K_0$, every Bayes error set of 3\,000 heavy shadows an initial segment). Compare Proposition~\ref{prop:D2sigma}: there a non-integrable density endorses \emph{every} comparison between regular models; here an integrable one endorses every comparison that is not end-positive.}

\emph{Protection and visibility} (\src{E1} \S0, \S2, corrected after \src{F2} and \src{F12\_referee}). The mechanism of Theorem~\ref{thm:D21} has two parts. \emph{The base} is the selected model out of which the generators compare; in three worlds it can be moved anywhere, $T_1$ included (sources in $E_1(T_1)$, targets in an end at which $T_1$ is correct), which makes every shadow fail on $T_1$'s out-comparisons (case~(i) of Theorem~\ref{thm:R3}); by Theorem~\ref{thm:E11} this is irrelevant. \emph{The protection} is the end. Under the hypotheses of Lemma~\ref{lem:E13}(a) every end is eventually $T_1$-correct at the protected world; for ends in a collar of $\partial U$ inside $\{h_i<0\}$ this holds with no hypothesis (Proposition~\ref{prop:F21}(a)); and under (Nest$_e$) on both edges, with the hypotheses of Proposition~\ref{prop:F21}(b), the ends lie at the near-unanimous boundary or at $\partial U\cap\{h_i<0\}$. The round-16 text said that the protection cannot be moved off the near-unanimous limits; that was unproved, and nothing proved confines it there. At an end on which $T_1$ is eventually correct, the argument of Theorem~\ref{thm:M0closed} may use a shadow that is arbitrarily heavy there (Lemma~\ref{lem:E14}); the Bayes models of that shadow are correct deep in the end, so no dominator comparison against them is end-positive, and none is a protected comparison. Since Theorem~\ref{thm:M0closed} needs only one admissible shadow and its dominator comparisons (Theorem~\ref{thm:E11}), protected failures of adding up at such ends are invisible to it, and Theorem~\ref{thm:E15} turns this into an exclusion of the covers in which every world is an end-protected decoration with integrable, jointly end-heavy base densities. The sentence of \src{D2} \S5, that nothing known excludes covers whose non-additive worlds are of the type of Theorem~\ref{thm:D21}, is superseded only for covers that meet these hypotheses. \src{D2}'s Conjecture~D2.C (installing that world on an edge region $R_{ij}$ of a three-world full system), if true, would give a three-world full system with a world that does not add up. Whether that world could carry a cover of $T_1$ depends on whether its base has an integrable end-heavy endorsing density on all of $M$, compatible with integrable endorsers of the other two worlds; Proposition~\ref{prop:E16} shows this only in the two-world system. (Round 17: heaviness and compatibility are not needed. Next to two worlds with integrable full endorsers, an a.e.\ finite endorser of the installed world's base in the cancellation form already excludes the installation (Theorem~\ref{thm:H12}), and every tuple of base endorsers is admissible (Lemma~\ref{lem:H11}).) So the round-16 statement that Conjecture~D2.C does not bear on $T_1$ was overstated (\src{F2} \S1.2), and so was \src{E1}'s statement that every installation of this type is excluded or inert before a selection is chosen. No selection was built in \src{E1}.

\emph{What a cover still needs} (\src{E1} \S3, as corrected in \src{F2} \S4.3 and \src{F12\_referee} \S5.7). For every admissible shadow $\nu$, the set of weights $t$ at which some world has a dominator comparison into the incoherent $B^\nu_t$ that $\nu$ does not endorse is non-null. Under (H$_U$) this is automatic (Theorem~\ref{thm:E11}(a), (b)), and the world may depend on $\nu$ and $t$; so it is not a restriction on any one world. At an end-protected world whose ends are eventually $T_1$-correct, with $\nu$ replaced by its reweighting $\nu'$ of Lemma~\ref{lem:E14}, any unendorsed comparison at that world is a base comparison, eventually zero at every end, which the protection of Theorem~\ref{thm:D21} does not cover. In two worlds next to an additive world, \src{D2} \S0 states that unprotected ratio certificates can be separated by transport and then summed into a violation of (SR$_1$) or (QC$_1$); that statement is a heuristic reading (\src{D2\_referee} \S7), not a result of the corpus. Whether such failures can be carried in three worlds is the open point; rounds 17 and 19 narrow it (parts~(f$'$), (f$''$)), and it is stated precisely in part~(g).

\paragraph{(f$'$) Round 17: localized failures are invisible, the conflict digraph, and pattern switching.} Round 17 (\src{BRIEF15}) asked again for an exact cover of $T_1$, built on the lessons of rounds 14 to 16. The proposed design moved the conflict between the corner equation and the box into the \emph{compatibility} of the worlds' endorsing densities: the bases of the worlds were to have endorsing densities, but no tuple of them was to be admissible. Three notes were written, and none produced a cover. \src{H1} shows that this design target cannot be met, and that failures of adding up that can be confined to sets of arbitrarily small measure are invisible to the argument of Theorem~\ref{thm:M0closed}. \src{H2} treats two or three worlds that fail to add up, introduces the conflict digraph, and gives a full two-world system in which neither world adds up. \src{J1} tests the dominators of models that agree with a Bayes model except deep in the ends, where they may give any advice, and closes every conflict cycle whose ends can be put in an aligned order. \src{H1} and \src{H2} were refereed once, together, in \src{H12\_referee}. The referee found nothing wrong, asked for repairs of one omitted line, one reading of endorsement, two statements of scope, one deduction and two wordings, and proved a strengthening of Theorem~\ref{thm:H23} (Proposition~\ref{prop:HR1}); the referee's versions are used below. \src{J1} was not refereed in round 17. \emph{Round 19:} it has now been refereed once (\src{J1\_referee}). The referee found nothing wrong, corrected one question and two statements of scope, closed two one-line gaps, and proved a general criterion for aligned orders (Theorem~\ref{thm:JR2}); the statuses of the \src{J1} block below are updated, and the referee's additions follow it. No assumption is added anywhere.

\emph{Densities, readings and localization} (\src{H1} \S1). Assume (H$_U$). In this part a \emph{density} is a measurable $\rho=(\rho_k)_k$ with $0<\rho_k<\infty$ a.e.; it need not be integrable. Put $c^\rho\coloneqq(\rho_k\sigma_k)_k$, and define $B^\rho_t$, $Z^\rho_t$, admissibility and
\[
I_\rho\coloneqq\{t\in\Di:Z^\rho_t\text{ is null and }B^\rho_t\text{ is incoherent}\}
\]
as for shadows in part~(f). Endorsement by $\rho_k\mu$ is read in the total-cost form or in the cancellation form of part~(e). Both readings have property (R), and for integrable $\rho_k$ they coincide. A \emph{localizing family} is a decreasing sequence $\Theta=(\Theta_m)$ of measurable subsets of $M$ with $\mu(\Theta_m)\to0$; the empty family $\Theta_m=\varnothing$ is allowed. Let $C\preceq_kX$ be a world-$k$ comparison with vector $\delta\coloneqq1_{E_kX}-1_{E_kC}$. It is \emph{$\Theta$-localized} if $\delta\le0$ a.e.\ on $\Theta_m$ for some $m$: on $\Theta_m$, $X$ makes no world-$k$ error that $C$ does not make. It is \emph{$(\Theta,T_1)$-localized} if $\delta\le0$ a.e.\ on $\Theta_m\cap\{T_1\text{ is correct at world }k\}$ for some $m$. Suppose world $k$ is end-protected with a countable family $\mathcal S_k$. Choose depths $a_S(m)\uparrow\infty$ with $\sum_{S\in\mathcal S_k}\mu(S_{a_S(m)})\to0$, and put $\Theta_m\coloneqq\bigcup_{S\in\mathcal S_k}S_{a_S(m)}$ (the \emph{deep ends} of world $k$). A $\Theta$-localized relevant comparison is then $S$-positive for no $S\in\mathcal S_k$, so end-protection makes it eventually $\ge0$, hence eventually $=0$, on every end: it is a base comparison. So a density that endorses the base comparisons of world $k$ endorses its $\Theta$-localized relevant comparisons. For $t\in I_\rho$ and a model $Y$ put
\[
\Phi_t(Y)\coloneqq\int_{Y\sd B^\rho_t}|t\cdot c^\rho|\,d\mu\in[0,\infty],\qquad\Phi^*_t\coloneqq\inf\{\Phi_t(C):C\in\M_U\text{ coherent}\}.
\]

\begin{lemma}\label{lem:H10}
Let $\rho$ be a density and $t\in I_\rho$.
\begin{enumerate}[leftmargin=*,label=(\alph*)]
\item $\Phi^*_t>0$, and the infimum is attained ($\Phi^*_t=\infty$ is allowed).
\item If $Y$ is regular and $\rho_k\mu(E_kY)<\infty$ for every $k$, then $\Phi_t(Y)\le\sum_kt_k\,\rho_k\mu(E_kY)<\infty$. If $Y'$ is another such model, then $\Phi_t(Y)-\Phi_t(Y')=\sum_kt_k\bigl[\rho_k\mu(E_kY)-\rho_k\mu(E_kY')\bigr]$, all terms being finite.
\item Let $X$ be regular with $\Phi_t(X)<\Phi^*_t$ and $\rho_k\mu(E_kX)<\infty$ for every $k$, and let $C$ be coherent with $C\preceq_kX$ for every $k$. Then, under any reading of endorsement with property (R), some world $k$ does not endorse $C\preceq_kX$.
\end{enumerate}
\end{lemma}
\begin{proof}[Proof sketch]
(a) The coherent models of $\M_U$ are the $D_K$ with $K\supseteq K'$ (Proposition~\ref{prop:reduction}(a)). These credal sets form a Hausdorff-compact family, and $P_{K_n}\to P_K$ pointwise gives $1_{D_{K_n}}\to1_{D_K}$ off the null set $\{P_K=0\}$, the boundary of a convex cone with interior. Since $\mu$ is finite, the coherent models are compact for convergence in measure (\src{R1} Prop.~3.1). Along a subsequence that attains $\liminf\Phi_t(Y_n)$ and converges a.e., Fatou gives lower semicontinuity of $\Phi_t$, also for $[0,\infty]$-valued $\Phi_t$. So the infimum is attained at a coherent $C^*$. As $B^\rho_t$ is incoherent, $C^*\sd B^\rho_t$ is non-null, and $|t\cdot c^\rho|>0$ a.e.\ because $Z^\rho_t$ is null; so $\Phi_t(C^*)>0$. (b) On $Y\sd B^\rho_t$ the worlds at which $B^\rho_t$ is right are exactly the worlds at which $Y$ errs, so $|t\cdot c^\rho|\le\sum_{k:\,g\in E_kY}t_k\rho_k$ there; integrating gives the bound. For regular $Y$, $1_{E_kY}-1_{E_kY'}=\sigma_k(1_{Y'}-1_Y)$, so $\sum_kt_k\rho_k(1_{E_kY}-1_{E_kY'})=(1_{Y'}-1_Y)\,t\cdot c^\rho=(1_{Y\sd B^\rho_t}-1_{Y'\sd B^\rho_t})|t\cdot c^\rho|$ pointwise; every term is integrable, and integration gives the identity. (c) If every comparison were endorsed, (R) would give $\rho_k\mu(E_kC)\le\rho_k\mu(E_kX)<\infty$ for every $k$. $C$ is regular, so (b) gives $\Phi_t(C)\le\Phi_t(X)<\Phi^*_t\le\Phi_t(C)$.
\end{proof}
\status{proved (\src{H1} Lemma~H1.0); confirmed (\src{H12\_referee} \S\S1.1--1.3), one reader. The first sentence of (b) is the referee's added line: the note used that $\Phi_t(Y)$ is finite without showing it (\src{H12\_referee} \S1.2). NUMERICS (\src{h1py/h1\_localization.py}, $\varepsilon=0.3$, $s=0.9$, shadow $\nu^L$ with $L=5$): at two weights $t$ the model $B_t$ is incoherent; making it correct at world~1 on shrinking collars of $\partial U$ inside $\{h_1<0\}$ costs at most $1.7\cdot10^{-5}$ in $\Phi_t$ at the deepest collar, while the best coherent models found have $\Phi_t\approx5.3\cdot10^{-3}$ and $2.6\cdot10^{-3}$ (upper estimates of $\Phi^*_t$ only). \src{h12ref\_py/h12ref\_checks.py} checks the identity and the bound of (b) pointwise with a non-integrable $\rho_1$. \emph{What it adds.} Theorem~\ref{thm:E11}(b) says that the dominator comparisons into $B^\nu_t$ itself are not all endorsed. Part~(c) says the same for every regular model whose cost is within $\Phi^*_t$ of the Bayes cost. Such a model can be made correct on any prescribed set of small measure, which is where protection lives.}

\begin{lemma}\label{lem:H11}
Assume (H$_U$). Let $\rho$ be a density and $\Theta$ a localizing family such that every $\rho_k$ is integrable on $M\setminus\Theta_m$ for every $m$. Suppose that, for each world $k$, $\rho_k\mu$ endorses in the cancellation form every comparison $T_1\preceq_kY$ with $Y$ regular and $Y=T_1$ a.e.\ on some $\Theta_m$. Then $\rho$ is admissible.
\end{lemma}
\begin{proof}[Proof sketch]
Let $v_k(g)\coloneqq\rho_k(g)$ if $T_1$ is correct at world $k$ at $g$ and $v_k(g)\coloneqq-\rho_k(g)$ otherwise ($g\in M$); toggling the advice of $T_1$ on $A\subseteq M$ changes the world-$k$ error mass by $\int_Av_k$. (1) No $A\subseteq M\setminus\Theta_m$ has $\int_Av_k<0$ for $k=1,2,3$: otherwise $Y\coloneqq T_1\sd A$ is regular and equals $T_1$ on $\Theta_m$; $T_1$ is competitive on $\{1,2,3\}$ (through Theorem~\ref{thm:kkm} it is tied at every world with a selected model), so $T_1\preceq_kY$ for some $k$, and endorsement, with both masses finite, gives $\int_Av_k\ge0$. (2) By Lyapunov's theorem the range $\{\int_Av\,d\mu:A\subseteq M\setminus\Theta_m\}$ is compact and convex, contains $0$, and by (1) misses $-\operatorname{int}\R^3_+$; separation gives $p^{(m)}\in\Delta$ with $p^{(m)}\cdot v\ge0$ a.e.\ on $M\setminus\Theta_m$. (3) A limit point $p^*$ has $p^*\cdot v\ge0$ a.e.\ on $M$, since $\mu(\bigcap_m\Theta_m)=0$. (4) $p^*$ is interior: if $p^*_3=0$, the non-null set of $(-,-,+)$-gambles that $T_1$ accepts near $r_{12}$ has $p^*\cdot v=-p^*_1\rho_1-p^*_2\rho_2<0$; the other cases are cyclic. (5) On $D_{T_1}$, $v=c^\rho$, and off it $v=-c^\rho$, so $B^\rho_{p^*}\subseteq D_{T_1}$ and $D_{T_1}\setminus B^\rho_{p^*}\subseteq Z^\rho_{p^*}$.
\end{proof}
\status{proved (\src{H1} Lemma~H1.1); confirmed (\src{H12\_referee} \S1.3), one reader; re-verified in \src{J1} \S1. Only $T_1$'s own out-comparisons are used, and only those that leave a set of small measure untouched. For integrable $\rho$ and $\Theta_m=\varnothing$ it is Proposition~\ref{prop:M0}($\alpha$) (\src{N17} Prop.~2.1($\alpha$)); for ends and base endorsers it is step~(1) of Theorem~\ref{thm:E15}. The truncation to $M\setminus\Theta_m$ supplies the admissibility step without the Lyapunov range of all of $M$ that \src{F2} \S4.3 named as missing for non-integrable densities. \emph{Consequence for the design target} (proved; confirmed). Suppose every world is end-protected with countably many ends (possibly none), and every world's base has an endorsing density that is integrable, or a.e.\ finite with endorsement in the cancellation form. Take $\Theta_m$ to be the deep ends of all worlds together with $\{\rho_k>m\}\cap M$. Then $T_1\preceq_kY$ with $Y=T_1$ on $\Theta_m$ is a base comparison, so \emph{every} tuple of such base endorsers is admissible. The target of \src{BRIEF15}, base endorsing families with no admissible tuple, is therefore unattainable under (H$_U$) whenever every family is nonempty. Countably many ends are needed for the summable choice of depths.}

\begin{theorem}\label{thm:H12}
Assume (H$_U$). Suppose that worlds~$2$ and~$3$ have integrable densities $\rho_2,\rho_3>0$ a.e.\ that endorse all their relevant comparisons; that world~$1$ has a measurable $\rho_1$ with $0<\rho_1<\infty$ a.e.; that $\Theta$ is a localizing family (possibly empty); and that $\rho_1\mu$ endorses, in the cancellation form, every $\Theta$-localized relevant comparison of world~$1$. Then (H$_U$) is contradicted. Under the total-cost reading the same holds if $\rho_1\mu(E_1T_1)<\infty$.
\end{theorem}
\begin{proof}[Proof sketch]
Put $\Theta'_m\coloneqq(\Theta_m\cup\{\rho_1>m\})\cap M$, a localizing family with $\rho_1\le m$ on $M\setminus\Theta'_m$; a $\Theta'$-localized comparison is $\Theta$-localized. \emph{Admissibility:} Lemma~\ref{lem:H11} with $\Theta'$, since the comparisons $T_1\preceq_1Y$ with $Y=T_1$ on $\Theta'_m$ are $\Theta'$-localized. \emph{A weight in $I_\rho$:} for each $g$ the set $\{t:t\cdot c^\rho(g)=0\}$ is a line, so $Z^\rho_t$ is null for a.e.\ $t$ (Tonelli); if $B^\rho_t$ were coherent for a.e.\ $t$, Theorem~\ref{thm:M0closed}(a) would be contradicted. \emph{Near-Bayes models:} let $X_m\coloneqq B^\rho_t$ on $U\setminus\Theta'_m$, and on $\Theta'_m$ let $X_m$ be correct at world~1. $X_m$ is regular. Where $B^\rho_t$ errs at world~1, $|t\cdot c^\rho|\le t_2\rho_2+t_3\rho_3$, so $\Phi_t(X_m)\le\int_{\Theta'_m}(\rho_2+\rho_3)\,d\mu\to0$; and $\rho_1\mu(E_1X_m)\le m\,\mu(M)$. \emph{Dominators:} Theorem~\ref{thm:kkm} gives a selected, undominated, hence coherent $C_m$ with $C_m\preceq_kX_m$ at every $k$. The world-1 vector is $\le0$ on $\Theta'_m$, where $X_m$ is correct, so that comparison is $\Theta'$-localized and endorsed; worlds~2 and~3 endorse theirs. Once $\Phi_t(X_m)<\Phi^*_t$ this contradicts Lemma~\ref{lem:H10}(c).
\end{proof}
\status{proved (\src{H1} Thm.~H1.2); confirmed (\src{H12\_referee} \S1.4), one reader. The cancellation form is invoked only on comparisons whose positive part has finite $\rho_1$-mass ($T_1\preceq_1T_1\sd A$ with $A\subseteq M\setminus\Theta'_m$, and $C_m\preceq_1X_m$), so the vacuous case $\infty\le\infty$ never enters, and the theorem holds under any reading with (R). \emph{Theorem~\ref{thm:M0closed}(a) without integrability} (\src{H12\_referee} \S1.4): $c^\rho$ and $c^\rho/|c^\rho|$ have the same Bayes sets, tie sets and exceptional set; the second field is strictly signed and bounded, hence integrable; property ($\gamma$) and the form $D_{T_1}=B_{p^*}\cup S$ with $S\subseteq Z_{p^*}$ concern these sets only. So Theorem~\ref{thm:M0closed}(a) for every measurable strictly signed field follows from its integrable case. \emph{Consequences} (\src{H1} Cor.~H1.4; confirmed as scoped, \src{H12\_referee} \S1.5). Suppose two worlds $j$, $k$ have integrable densities that endorse all their relevant comparisons, for instance additive worlds with integral penalties, worlds that add up with countably additive weights, or worlds with leading densities (Lemma~\ref{lem:B10}). Then, under (H$_U$), the third world $i$ is not (a) an end-protected world with countably many ends, anywhere and of any heaviness, whose base has an integrable endorsing density, or an a.e.\ finite one in the cancellation form (base endorsement implies localized endorsement); nor (b) a world that adds up $\sigma$-finitely in the cancellation form ($\Theta=\varnothing$). In two worlds, the density of Proposition~\ref{prop:D2sigma} endorses in the cancellation form: relevant comparisons out of $C_b$ have negative part in $c_b$, which has finite mass, and generator comparisons are endorsed as a finite cost against an infinite one. So a three-world installation of a world that behaves like world~1 of Theorem~\ref{thm:D21}, next to two worlds with integrable full endorsers, is excluded as soon as it keeps an a.e.\ finite endorser in the cancellation form. Within this setting the warning of \src{BRIEF15} (that an ``end-light'' base might collapse) is moot: the system is excluded whether the base is end-heavy or end-light and whatever its core, and no collapse argument is needed (\src{H1} \S5). The nearby target of part~(g) (\src{F2} \S4.3) is thereby proved when only one world's density is non-integrable and the other two worlds have integrable full endorsers.}

\begin{theorem}\label{thm:H13}
Assume (H$_U$). Suppose there are integrable densities $\rho_k>0$ a.e.\ and localizing families $\Theta^k$, $k=1,2,3$ (each possibly empty), such that $\rho_k\mu$ endorses every $(\Theta^k,T_1)$-localized relevant comparison of world $k$. Then (H$_U$) is contradicted. The same holds with ``$\Theta^k$-localized'' in place of ``$(\Theta^k,T_1)$-localized'', provided $\Theta^i_m\cap\Theta^j_m\cap R_{ij}$ is null for all $i\ne j$ and all $m$.
\end{theorem}
\begin{proof}[Proof sketch]
Put $\Theta_m\coloneqq\bigcup_k\Theta^k_m$. Comparisons $T_1\preceq_kY$ with $Y=T_1$ on $\Theta_m$ have $\delta=0$ there, so Lemma~\ref{lem:H11} gives admissibility; take $t\in I_\rho$ as before. First version: let $X_m\coloneqq B^\rho_t$ off $\Theta_m$ and $X_m\coloneqq T_1$ on $\Theta_m$. On $\Theta^k_m\cap\{T_1\text{ correct at }k\}$, $X_m$ is correct at world $k$, so every $C\preceq_kX_m$ is $(\Theta^k,T_1)$-localized; $\Phi_t(X_m)\le\sum_k\int_{\Theta_m}\rho_k\to0$, all costs are finite, and Lemma~\ref{lem:H10}(c) concludes. Second version: on $\Theta^k_m$ let $X_m$ be correct at world $k$. This is consistent, because two worlds disagree only on $R_{ij}$, which the hypothesis excludes; argue as before.
\end{proof}
\status{proved (\src{H1} Thm.~H1.3); confirmed, both versions (\src{H12\_referee} \S1.5), one reader. \emph{Scope.} Endorsement of every $(\Theta^k,T_1)$-localized comparison is \emph{stronger} than base endorsement when world $k$ has ends on which $T_1$ is wrong at world $k$: such a comparison may be positive at such an end. If $\Theta^k_m\subseteq\{T_1\text{ correct at }k\}$, the two notions of localization coincide, and the compatibility condition of the second version holds automatically, since on $R_{ij}$ $T_1$ is correct at exactly one of $i$, $j$. Ends that are eventually $T_1$-correct at their world can be cut to such sets. This is automatic for ends in a collar of $\partial U$ inside $\{h_k<0\}$ (Proposition~\ref{prop:F21}(a)), it holds under the hypotheses of Lemma~\ref{lem:E13}(a), and it holds under (h) of Theorem~\ref{thm:E15} (see below). Integrability is used for $\Phi_t(X_m)\to0$ when $X_m$ errs at a world whose density is not integrable on $\Theta_m$; a version for two or more worlds with non-integrable densities is not proved.}

\emph{What replaces the heaviness hypothesis of Theorem~\ref{thm:E15}} (\src{H12\_referee} \S1.6). \src{H1} states that Theorem~\ref{thm:E15} ``holds without (h)''. Literally that would mean hypothesis (b) alone, with ends on which $T_1$ may be wrong at their world, and that is not what is proved. Theorem~\ref{thm:H12} needs integrable endorsers of \emph{all} relevant comparisons at the two other worlds, and Theorem~\ref{thm:H13} needs every end to be eventually $T_1$-correct at its world. What is proved is this: \emph{Theorem~\ref{thm:E15} holds with (h) replaced by eventual $T_1$-correctness of every end at its world (Theorem~\ref{thm:H13}, with $\Theta^k$ the deep ends of world $k$), a condition that (h) implies; and, for finite families of ends, it holds with (h) dropped altogether unless the conflict digraph defined below has a directed cycle (Theorem~\ref{thm:H23} as strengthened in Proposition~\ref{prop:HR1}).} That (h) implies $T_1$-correctness: step~(1) of Theorem~\ref{thm:E15}, which does not use (h), gives $T_1=B^\psi_{p^*}$ up to ties; deep in an end of world $i$, (h) gives $p^*\cdot c^\psi\sigma_i\ge p^*_i\rho_i-p^*_j\rho_j-p^*_k\rho_k>0$. So Theorem~\ref{thm:E15} is a special case of Theorem~\ref{thm:H13}. Neither result covers ends on which $T_1$ is wrong at their world when the conflict digraph has a cycle. \src{H1} \S6 also lists what survives Theorems~\ref{thm:H12} and~\ref{thm:H13}; the referee's corrected list (\src{H12\_referee} \S1.7) is given at the end of this part.

\emph{Ends refined} (\src{H2} \S1). An end $S$ of world $i$ is \emph{$T_1$-correct} if $\mu(S_m\cap E_i(T_1))=0$ for some $m$, and \emph{$T_1$-wrong} if $\mu(S_m\setminus E_i(T_1))=0$ for some $m$. It lies in an open mixed orthant $O$ if $\mu(S_m\setminus O)=0$ for some $m$; $M$ is the union of the six open mixed orthants up to null sets.

\begin{lemma}\label{lem:H20}
Let world $i$ be end-protected with a countable family $\mathcal S_i$. Replace each $S\in\mathcal S_i$ by its twelve pieces $S\cap O\cap E_i(T_1)$ and $S\cap O\setminus E_i(T_1)$, $O$ ranging over the open mixed orthants, and discard each piece that is null from some depth on. The new family is countable, and finite if $\mathcal S_i$ is; each of its members is an end that lies in one orthant and is $T_1$-correct or $T_1$-wrong at world $i$; world $i$ is end-protected with it; and the base comparisons of world $i$ are the same for both families.
\end{lemma}
\begin{proof}[Proof sketch]
Each kept piece decreases, is non-null at every depth and has null intersection. If $\delta$ is $S$-positive, then at every depth one of the twelve pieces meets $\{\delta=1\}$ non-nully; one piece does so at infinitely many depths, hence at all depths, since the pieces decrease, and that piece is kept. A discarded piece is null from some depth, so ``eventually $\ge0$ (or $=0$) on $S$'' is equivalent to the same statement on every kept piece of $S$.
\end{proof}
\status{proved (\src{H2} Lemma~H2.0); confirmed (\src{H12\_referee} \S2.1), one reader.}

From now on ends are refined in this way. Fix a $T_1$-wrong end $W$ of world $i$ lying in the orthant $O$. On $W$, $T_1$ gives one advice a.e., and it is wrong at world $i$. Let $V(W)$ be the set of worlds $j$ with $\sigma_j\ne\sigma_i$ on $O$, i.e.\ the worlds at which $T_1$ is correct on $W$; it is nonempty because $W\subseteq M$. For each world fix the non-null \emph{improvement region} $R^0_1\coloneqq D_{T_1}\cap(+,-,-)$, $R^0_2\coloneqq D_{T_1}\cap(-,+,-)$, $R^0_3\coloneqq D_{T_1}\cap(-,-,+)$, on which rejecting instead of accepting is bad at world $i$ and good at both other worlds.

\begin{lemma}\label{lem:H21}
Assume (H$_U$). Let world $i$ be end-protected with a countable family $\mathcal S_i$, and let $W\in\mathcal S_i$ be $T_1$-wrong, lying in the orthant $O$.
\begin{enumerate}[leftmargin=*,label=(\alph*)]
\item There are a world $j\in V(W)$ and a non-null $R\subseteq R^0_i$ such that every finitely additive $\nu\ge0$ on $U$ that vanishes on null sets and endorses the relevant world-$j$ comparisons out of $T_1$ has $\nu(R)=0$. In particular world $j$ does not add up. If $V(W)=\{j\}$, then for every non-null $p\subseteq W$ that meets every $W_m$ non-nully, and every non-null $R\subseteq R^0_i$ off the deep parts of the ends of world $i$ (depths fixed in advance), world $j$ accepts $1_p-1_R$, i.e.\ $T_1\preceq_jT_1\sd(p\cup R)$.
\item Suppose moreover that every world $j\in V(W)$ is end-protected with a finite family $\mathcal S_j$ and has a measurable $\rho_j$, $0<\rho_j<\infty$ a.e., such that $\rho_j\mu$ endorses the base comparisons of world $j$ out of $T_1$, endorsement being read in the cancellation form or with $\rho_j\mu(E_jT_1)<\infty$. Then $W$ is \emph{shadowed}: for every $n$ there is $m'$ with
\[
W_{m'}\subseteq\bigcup\{S_n:\ S\text{ a $T_1$-correct end of some }j\in V(W)\}\quad\text{modulo null sets.}
\]
\end{enumerate}
\end{lemma}
\begin{proof}[Proof sketch]
Choose depths for the members of $\mathcal S_i$ so that $R$, a non-null part of $R^0_i$ off their deep parts, exists, and let $p\subseteq W$ be non-null and meet every $W_m$. Put $X\coloneqq T_1\sd(p\cup R)$, a regular model. At a world with the sign of world $i$ on $O$, $X$ is better than $T_1$ on a non-null set and nowhere worse, so strict (G) gives $T_1\not\preceq X$ there. At world $i$ the vector $1_R-1_p$ is $-1$ at every depth of $W$ and positive at no end, so end-protection gives $T_1\not\preceq_iX$. Competitiveness of $T_1$ (through Theorem~\ref{thm:kkm}) then gives $T_1\preceq_jX$ for some $j\in V(W)$, with vector $1_p-1_R$. (a) Split the depth layers of $W$ to obtain pairwise disjoint $p_1,p_2,\dots$, each meeting every $W_m$ non-nully. Some $j$ accepts $1_{p_n}-1_R$ for infinitely many $n$, so an endorsing $\nu$ has $\nu(R)\le\nu(p_n)$ for those $n$, and boundedness gives $\nu(R)=0$. (b) If $W_{m'}\setminus\bigcup_SS_n$ were non-null for every $m'$, one could choose $p$ inside it, meeting every depth, with $\rho_j(p)<\rho_j(R)$ for every $j\in V(W)$, after shrinking $R$ off the sets $S_n$ (finitely many). Then $1_p-1_R$ vanishes on $S_n$ for every end of every $j\in V(W)$, so it is a base comparison out of $T_1$, and endorsement gives $\rho_j(R)\le\rho_j(p)$, a contradiction. $T_1$-wrong ends of the worlds of $V(W)$ are disjoint from $W$ at large depth, so only $T_1$-correct ends are needed in the cover.
\end{proof}
\status{proved (\src{H2} Lemma~H2.1); confirmed (\src{H12\_referee} \S\S2.1--2.2), one reader. \emph{Repair in (b).} The note allows $\rho_j$ finite a.e.\ without fixing the reading of endorsement. Under the total-cost reading with $\rho_j\mu(E_jT_1)=\infty$ the comparison $T_1\preceq_jX$ is endorsed vacuously, and no contradiction follows; so (b) is stated with the cancellation form or with finite cost, as the referee requires (\src{H12\_referee} \S2.2, the same defect as \src{F12\_referee} \S5.2). For integrable $\rho_j$, the case used in Theorem~\ref{thm:H23}, nothing changes. In (b) the depth $n$ may be raised freely, which is what makes the shrinking of $R$ legitimate; it uses finiteness of the families $\mathcal S_j$. \emph{A further consequence} (proved in a referee report, \src{H12\_referee} \S2.1, one reader): if $V(W)=\{j\}$, every a.e.\ finite endorser of world $j$'s comparisons out of $T_1$, in the cancellation form, has $\rho_j\mu(W_m)=\infty$ for every $m$. \emph{Remarks} (\src{H2} \S2). Part~(a) uses only (H$_U$) and end-protection of world $i$. It is the counterpart of Lemma~\ref{lem:E13}(a): that lemma assumes endorsers at the other worlds and concludes that every end is $T_1$-correct, while (a) drops the assumption and concludes that a $T_1$-wrong end needs a partner world that fails to add up. Such ``conflict ends'' can sit wherever $T_1$ is wrong at world $i$, for instance in the collar of $\partial U\cap\{h_i>0\}$, where $T_1$ rejects (Proposition~\ref{prop:F21}(a)), or inside $D_{T_1}\cap\{h_i<0\}$.}

\begin{proposition}\label{prop:H22}
Let $\nu$ be an admissible shadow and, for each world $i$, let $\mathcal S^{\mathrm c}_i$ be a countable family of $T_1$-correct ends of world $i$, with any locations and any overlaps between worlds. The reweighting $\nu'_i=\nu_i(1+f_i)$ of Lemma~\ref{lem:E14} can be done at all worlds at once, and $\nu'$ is admissible and integrable. Moreover, with $\Omega^-(g)$ the set of worlds at which $T_1$ errs at $g$,
\[
\textbf{(h$'$)}\qquad\operatorname*{ess\,inf}_{g\in S_m}\ \rho'_i(g)\Big/\sum_{k\in\Omega^-(g)}\rho'_k(g)\ \longrightarrow\ \infty\quad(m\to\infty)\qquad\text{for every }S\in\mathcal S^{\mathrm c}_i.
\]
\end{proposition}
\begin{proof}[Proof sketch]
At a point of a $T_1$-correct end of world $i$, $T_1$ is right at $i$, and any other world whose $T_1$-correct end contains the point has the same sign; so the worlds whose weights are raised there never conflict. On $S_{m_l(S)}$, $\rho'_i\ge l(\rho_j+\rho_k)\ge l\sum_{k\in\Omega^-(g)}\rho'_k$, because $f_k=0$ for $k\in\Omega^-(g)$.
\end{proof}
\status{proved (\src{H2} Prop.~H2.2); confirmed (\src{H12\_referee} \S2.3), one reader. \emph{Mostly not new.} Lemma~\ref{lem:E14} already reweights all worlds at once; the new content is (h$'$), heaviness relative to the worlds at which $T_1$ errs. Hypothesis (b) of Theorem~\ref{thm:E15} does not transfer to the reweighted shadow, as stated in the scope of that theorem. After Proposition~\ref{prop:HR1}, this proposition is not needed for Theorem~\ref{thm:H23}.}

\begin{definition}\label{def:conflict}
Let every world be end-protected with ends refined as in Lemma~\ref{lem:H20}. The \emph{conflict digraph} $\vec{\mathcal H}$ has vertices $1,2,3$ and an arrow $i\to j$ ($j\ne i$) iff some $T_1$-wrong end $W\in\mathcal S_i$ and some $T_1$-correct end $S\in\mathcal S_j$ satisfy $\mu(W_m\cap S_m)>0$ for every $m$. Then $j\in V(W)$.
\end{definition}

\begin{theorem}\label{thm:H23}
Assume (H$_U$). Suppose every world $i$ is end-protected with a finite, possibly empty, family $\mathcal S_i$, refined as in Lemma~\ref{lem:H20}, and there are $\rho_i\in L^1(\mu)$, $\rho_i>0$ a.e., such that \textbf{(b)} $\rho_i\mu$ endorses every relevant world-$i$ comparison whose vector is eventually $=0$ on every $S\in\mathcal S_i$, and \textbf{(h$'$)} holds for every $T_1$-correct $S\in\mathcal S_i$. If $\vec{\mathcal H}$ has no directed cycle, then (H$_U$) is contradicted.
\end{theorem}
\begin{proof}[Proof sketch]
Let $\psi\coloneqq(\rho_i\mu)_i$. (1) $\psi$ is admissible (Lemma~\ref{lem:H11} with the deep ends of all worlds, or step~(1) of Theorem~\ref{thm:E15}). (2) Lemma~\ref{lem:H21}(b) applies to every $T_1$-wrong end. (3) By (h$'$), for each $t\in\Di$ there is a depth $m(t)$ beyond which $B^\psi_t$ is correct at world $j$ on every $T_1$-correct end of world $j$; by (2), deep in a $T_1$-wrong end of world $i$, $B^\psi_t$ is right at some world of $V(W)$, hence wrong at $i$. So for any model $C$ the vector $\delta_i\coloneqq1_{E_iB^\psi_t}-1_{E_iC}$ is $\le0$ deep in the $T_1$-correct ends of world $i$ and $\ge0$ deep in its $T_1$-wrong ends. (4) Let $t\in I_\psi$ and let $C$ be a dominator of $B^\psi_t$ from Theorem~\ref{thm:kkm}. If $\delta_i$ is not eventually $0$ on some end of world $i$, end-protection makes it $W$-positive for a $T_1$-wrong $W\in\mathcal S_i$. On $\{\delta_i=1\}\cap W$ the model $C$ differs from $B^\psi_t$. This set is covered, at every depth, by finitely many $T_1$-correct ends of worlds of $V(W)$, so by pigeonhole one such end $S\in\mathcal S_j$ meets it at every depth: then $i\to j$, and $\delta_j=-1$ there at every depth, so $\delta_j$ is not eventually $0$ on a $T_1$-correct end of world $j$. Iterating gives an infinite walk in a digraph with three vertices, hence a directed cycle. If $\vec{\mathcal H}$ is acyclic, every dominator comparison is a base comparison and is endorsed by (b); as $B^\psi_t$ is the a.e.\ unique minimiser, $C=B^\psi_t$ a.e., against incoherence (Theorem~\ref{thm:E11}(b)).
\end{proof}
\status{proved (\src{H2} Thm.~H2.3); confirmed (\src{H12\_referee} \S2.4), including the pigeonhole step and the use of finiteness, one reader. The arrow condition $\mu(W_m\cap S_m)>0$ for every $m$ needs the covering of $\{\delta_i=1\}\cap W$ at depth $\max(m,m'(n))$ rather than $m'(n)$, a trivial adjustment (\src{H12\_referee} \S2.4). Finiteness of the families is used exactly in the pigeonhole of step~(4), in the uniform depth $m(t)$, and in Lemma~\ref{lem:H21}(b); with countably many ends the covering in step~(4) could use a different end at each depth, and the argument fails. No location of the ends is used. Theorem~\ref{thm:E15} is the case without $T_1$-wrong ends (no arrows), with (h) replaced by the weaker (h$'$).}

\begin{proposition}\label{prop:HR1}
Theorem~\ref{thm:H23} holds without (h$'$): under (H$_U$), finite refined families $\mathcal S_i$, and $\rho_i\in L^1(\mu)$ with $\rho_i>0$ a.e.\ and (b), there is a contradiction whenever $\vec{\mathcal H}$ is acyclic.
\end{proposition}
\begin{proof}[Proof sketch]
$\psi$ is admissible; take $t\in I_\psi$. Choose depths $a_S(m)\uparrow\infty$ beyond the depth at which each end is $T_1$-correct or $T_1$-wrong, let $\Theta_m$ be the union of the $S_{a_S(m)}$ over all ends of all worlds, and put $X_m\coloneqq T_1$ on $\Theta_m$ and $X_m\coloneqq B^\psi_t$ off $\Theta_m$. $X_m$ is regular, $\Phi_t(X_m)\le\int_{\Theta_m}\sum_k\rho_k\to0$, and all costs are finite. Let $C_m$ be a dominator of $X_m$ from Theorem~\ref{thm:kkm}, with vectors $\delta_k\coloneqq1_{E_kX_m}-1_{E_kC_m}$. Because $X_m=T_1$ deep in the ends, $\delta_k\le0$ deep in the $T_1$-correct ends of world $k$ and $\delta_k\ge0$ deep in its $T_1$-wrong ends; this replaces step~(3) of Theorem~\ref{thm:H23}, the only place where (h$'$) was used. Step~(4) then runs as before with $X_m$ in place of $B^\psi_t$: on $\{\delta_i=1\}\cap W\cap S_n$, with $S$ a $T_1$-correct end of a world $j\in V(W)$ and $n$ large, $X_m=T_1$ is right at $j$ and $C_m\ne X_m$, so $\delta_j=-1$ at every depth. If $\vec{\mathcal H}$ is acyclic, every $\delta_k$ is eventually $0$ on every end, every dominator comparison is a base comparison and is endorsed by (b), and this contradicts Lemma~\ref{lem:H10}(c) once $\Phi_t(X_m)<\Phi^*_t$.
\end{proof}
\status{proved in a referee report (\src{H12\_referee} \S2.4, Prop.~R.1), one reader. The same statement, by the same near-Bayes model, is proved independently in note \src{J1} (Thm.~J1.1; proved in a note, confirmed in \src{J1\_referee} \S2, round 19, one reader). \emph{Consequence.} Heaviness, (h) or (h$'$), plays no role anywhere in the route with finitely many ends and integrable base endorsers, and Proposition~\ref{prop:H22} is not needed for Theorem~\ref{thm:H23}.}

\begin{corollary}\label{cor:H23p}
Under the hypotheses of Proposition~\ref{prop:HR1}, a cover of $T_1$ needs a directed cycle in $\vec{\mathcal H}$: a $2$-cycle $i\rightleftarrows j$ (a $T_1$-wrong end of $i$ meets a $T_1$-correct end of $j$ at every depth, and vice versa) or a $3$-cycle. Hence at least two worlds have $T_1$-wrong ends, and at least two worlds have no integrable density that endorses all their relevant comparisons. If the ends on the cycle have one-element sets $V(W)$, the partner worlds given by Lemma~\ref{lem:H21}(a) are distinct, and at least two worlds fail to add up.
\end{corollary}
\begin{proof}[Proof sketch]
The cycle is Proposition~\ref{prop:HR1}. If two worlds had integrable densities endorsing all their relevant comparisons, the third would have an integrable base endorser, which endorses its $\Theta$-localized comparisons for $\Theta$ its deep ends, and Theorem~\ref{thm:H12} would give a contradiction; no cycle is needed for this. The last sentence follows from Lemma~\ref{lem:H21}(a).
\end{proof}
\status{proved (\src{H2} Cor.~H2.3$'$, corrected in \src{H12\_referee} \S2.5), one reader. \emph{Correction.} The note concluded that ``at least two worlds fail to add up'' because every world on the cycle has a partner that fails to add up. That does not follow when a set $V(W)$ has two elements: in a $2$-cycle $1\rightleftarrows2$ with $W_1$ in $(+,-,-)$ and $W_2$ in $(-,+,-)$, $V(W_1)=\{2,3\}$ and $V(W_2)=\{1,3\}$, and Lemma~\ref{lem:H21}(a) may return world~3 both times. The corrected statement is the referee's: the bound on integrable full endorsers holds in general, through Theorem~\ref{thm:H12}, and the bound on worlds that fail to add up holds when the sets $V(W)$ on the cycle are singletons.}

\emph{The template of a conflict cycle, and the germ exchanges} (\src{H2} \S5, with \src{H12\_referee} \S2.6). \emph{Geometry} (CHECKED). A $2$-cycle between worlds~1 and~3 can sit in the collar of $\partial U$, where $T_1$ rejects (Proposition~\ref{prop:F21}(a)): an end $A$ in the collar part of $(+,+,-)$, $T_1$-wrong at worlds~1 and~2 and $T_1$-correct at world~3; an end $B$ in the collar part of $(-,+,+)$, $T_1$-wrong at worlds~3 and~2 and $T_1$-correct at world~1; world~1 protected at $A$ and $B$, and world~3 protected at $B$ and $A$. If $A$ and $B$ shrink to rays in the relative interiors of two facets of $\partial U$, coherent models can give every combination of advice on $A$ and $B$, so nothing geometric excludes the configuration (the collar geometry is checked in \src{h12ref\_py/h12ref\_checks.py}~(3)). \emph{What a shadow with (b) and (h$'$) forces} (proved, from the proof of Theorem~\ref{thm:H23}). Deep in $A$ and $B$, $B^\psi_t$ rejects, as $T_1$ does. A dominator $C$ of an incoherent $B^\psi_t$ that (b) does not rule out must, along the walk, accept non-null parts of $A$ and of $B$ at every depth. So each of worlds~1 and~3 must accept a comparison that exchanges an error germ at its $T_1$-wrong end for an error germ at its $T_1$-correct end. Two corrections of the referee are incorporated: the note's ``and in fact for every admissible shadow (E1.1)'' is deleted, since for an admissible shadow without (b) the unendorsed dominator comparison may be a base comparison, and then nothing forces an exchange; and the weight of the germs comes from end-protection itself (at the world's own $T_1$-wrong end) and from Lemma~\ref{lem:H21}(a) (germ dominance, at its $T_1$-correct end), not from Lemma~\ref{lem:H21}(b), which is a shadowing statement. \emph{Where the design broke} (\src{H2} \S5.3; the displayed claim proved, confirmed). For every set $G$ of toggles good at world~2, competitiveness of $T_1$ against $T_1\sd(D_A\cup D_B\cup G)$, added to the dominator vector by Lemma~\ref{lem:vectorform}, puts a vector supported off every end into the acceptance set of world~1 or of world~3. A contradiction would follow from (b) if such vectors were realised by relevant comparisons out of selected models; in general they are not, and no world order was found that accepts the exchanges with closed sections while keeping $T_1$ competitive. \emph{Question H2.Q} (\src{H2} \S5.4), as restated by the referee without (h$'$) and with $X_m$ of Proposition~\ref{prop:HR1} in place of $B^\psi_t$ (\src{H12\_referee} \S3): can the dominators of the models $X_m$ make these exchanges, for a non-null set of $t$, with closed sections and $T_1$ competitive? \src{J1} answers it negatively for every aligned configuration, including this template (Corollary~\ref{cor:J13}(a) below).

\begin{theorem}\label{thm:H24}
Let $n=2$, in the setting and notation of Theorem~\ref{thm:D21} (host~A, chain coordinate $x$, $c_b$, $d_b$, $C_b$), and put $\tilde x\coloneqq1-x$ on $M$. Fix a design $\mathcal G$ in the coordinate $x$ (parameters $a$, $h_0$) and a design $\mathcal G'$ in the coordinate $\tilde x$ (parameters $a'$, $h'_0$). Let world~1 accept $X\preceq_1Y$ iff $1_{E_1Y}-1_{E_1X}\in\mathcal A_1$, the acceptance set of world~1 of Theorem~\ref{thm:D21} built on $(x,\mathcal G)$, and let world~2 accept $X\preceq_2Y$ iff $1_{E_2Y}-1_{E_2X}\in\mathcal A'$, the same construction built on $(\tilde x,\mathcal G')$. With the selection $p\mapsto C_{p_2}$ this ranking satisfies strict (G), (IB), (LA), (S) (competitiveness on $\supp p$, and closed sections for every model at both worlds), (C), $\C_L=\{C_b:b\in[0,1]\}$, and (LF) at every interior $C_b$. Neither world adds up: every bounded finitely additive $\nu\ge0$ on $U$ that vanishes on null sets and endorses the relevant comparisons of world~1 vanishes on $c_a$, and every one that endorses those of world~2 vanishes on $d_{1-a'}$.
\end{theorem}
\begin{proof}[Proof sketch]
For a model $X$ with $E_1X\cap M=A$ and unanimous errors $N_X$, put $P\coloneqq c_b\setminus A$ and $Q\coloneqq A\setminus c_b$; then $C_b\preceq_1X$ iff $\delta_1\coloneqq1_Q-1_P+1_{N_X}\in\mathcal A_1$, and $C_b\preceq_2X$ iff $\delta_2\coloneqq1_P-1_Q+1_{N_X}\in\mathcal A'$. (1) Lemmas~\ref{lem:D23} and~\ref{lem:D24} are statements about one world's rule and the chain; $\tilde x$ pushes $\bar\mu|_M$ to Lebesgue measure and $E_2C_b=\{\tilde x<1-b\}$, so they hold at world~2 in the coordinate $\tilde x$, where the sections are closed final segments in $b$. (2) Competitiveness, (K) for every model: if $N_X$ is null and $\bar\mu(Q)\ge\bar\mu(P)$, then $\delta_1\in\mathcal K_0$; if $\bar\mu(Q)<\bar\mu(P)$, the mirrored statement holds for $\delta_2$ in $\mathcal A'$. If $\bar\mu(N_X)>0$, rejection at world~1 forces $\bar\mu(P)>\bar\mu(Q)+\bar\mu(N_X)$ and rejection at world~2 forces $\bar\mu(Q)>\bar\mu(P)+\bar\mu(N_X)$, which are incompatible. (3) No weak domination of $C_b$: (R$^-$) forces $N_X$ null; for $b<1$ the world-1 condition forces $Q$ null, and then strict (G) at world~2 forces $P$ null; $b=1$ is treated separately, as in Theorem~\ref{thm:D21}, since there $Q=A\setminus c_1$ is empty. (4) $\C_L$, (C), (LA) as in steps~(4)--(5) of Theorem~\ref{thm:D21}, and (LF) as in its step~(6). (5) World~1 does not add up by step~(7) of Theorem~\ref{thm:D21}; world~2 by the mirror of that step, with sources in $d_{1-a'}$.
\end{proof}
\status{proved (\src{H2} Thm.~H2.4); confirmed (\src{H12\_referee} \S2.7), every axiom re-derived independently, one reader. (K) was also checked in exact rational arithmetic on random unions of intervals (\src{h12ref\_py/h12ref\_checks.py}~(4): 1\,500 tests, 0 failures; 573 cases accepted only at world~1 and 688 only at world~2). \emph{Wording repair.} In step~(3) the note says that $-\delta_1-\gamma_*\in\mathcal K_0$ is impossible because near $x=1$ it is $\le-1$ on $Q_\gamma$; that is false at $b=1$, where $-\delta_1$ may be $+1$ near $1$. At $b=1$, however, $Q=A\setminus c_1$ is empty, so the conclusion stands; $b=1$ is treated separately, as both referees of \src{D2} required for Theorem~\ref{thm:D21} (\src{H12\_referee} \S2.7). NUMERICS (\src{h2py/h2\_mirrored\_d2.py}, an exact cell model symmetric under $x\mapsto1-x$; 0 failures in 1\,500 competitiveness tests, 1\,500 no-weak-domination tests, 1\,200 strict-(G) tests, 528 transitivity triples and 80 section-shape tests; rerun byte-identically by the referee). \emph{What it shows.} In the full systems of Theorems~\ref{thm:C21}, \ref{thm:D21} and Proposition~\ref{prop:ExA}, the ones \src{H2} examined, some world adds up; here neither does, so two worlds that fail to add up can coexist under every axiom. It is not a conflict cycle: relative to any interior target $C_{b^*}$ both ends are target-correct at their worlds and the bases have integrable endorsers, so in a three-world installation with $T_1$-correct ends Theorem~\ref{thm:H13} applies, with or without heaviness (\src{H12\_referee} \S2.7).}

\emph{Round 17, note \src{J1}: near-Bayes models and pattern switching} (\textbf{proved in a note; round 19: confirmed in \src{J1\_referee}, one reader}; every result of this block carries that status unless stated otherwise). Hypothesis \textbf{(EP)}: (H$_U$) holds; each world $k$ is end-protected with a \emph{finite}, possibly empty, family $\mathcal S_k$, refined as in Lemma~\ref{lem:H20}; and there are $\rho_k\in L^1(\mu)$, $\rho_k>0$ a.e., such that $\rho_k\mu$ endorses every relevant world-$k$ base comparison (hypothesis (b) of Theorem~\ref{thm:H23}). No heaviness is assumed. Put $\psi\coloneqq(\rho_k\mu)_k$, $\mathcal E^k_m\coloneqq\bigcup_{S\in\mathcal S_k}S_m$ and $\mathcal E_m\coloneqq\bigcup_k\mathcal E^k_m$; then $\mu(\mathcal E_m)\to0$.

\begin{lemma}\label{lem:J1C}
Assume (EP). Then $\psi$ is admissible and $I_\psi$ is not null. Let $t\in I_\psi$, and let $X$ be any model with $X=B^\psi_t$ on $U\setminus\mathcal E_m$ and any advice on $\mathcal E_m$. Then $X$ is regular, $\Phi_t(X)\le\int_{\mathcal E_m}\sum_k\rho_k\,d\mu$, and $\psi_k(E_kX)<\infty$ for every $k$. If $m$ is so large that $\int_{\mathcal E_m}\sum_k\rho_k\,d\mu<\Phi^*_t$, Theorem~\ref{thm:kkm} gives a selected, undominated, coherent $C$ with $C\preceq_kX$ at every $k$, all three comparisons are relevant, and at least one of them is not a base comparison.
\end{lemma}
\begin{proof}[Proof sketch]
Admissibility is Lemma~\ref{lem:H11} with $\Theta=(\mathcal E_m)$: a comparison $T_1\preceq_kY$ with $Y=T_1$ on $\mathcal E_m$ is a base comparison. $I_\psi$ is not null by Theorem~\ref{thm:E11}(a). $X$ is regular because $\mathcal E_m\subseteq M$ and $B^\psi_t\supseteq N_+$ a.e. The bound on $\Phi_t$ holds because $X$ differs from $B^\psi_t$ only on $\mathcal E_m$ and $|t\cdot c^\psi|\le\sum_k\rho_k$. By Lemma~\ref{lem:H10}(c) not all three comparisons are endorsed, so by (b) not all are base comparisons.
\end{proof}
\status{proved (\src{J1} \S1, Steps~A--C); confirmed (\src{J1\_referee} \S1), one reader. The referee stresses that the bound on $\Phi_t(X)$ is uniform in the advice pattern on $\mathcal E_m$: this is what allows $m$ to be fixed, after $t$, before the pattern is chosen, and every result of \src{J1} relies on it. \src{J1} re-verifies the parts of Lemmas~\ref{lem:H10} and~\ref{lem:H11} that it uses. This is the \emph{near-Bayes tool}: the tested model need not be the Bayes model, which at a conflict end is wrong at one of the two worlds whatever the shadow, and on the deep ends it may carry any advice pattern.}

For a comparison $C\preceq_kX$ write $\delta_k\coloneqq1_{E_kX}-1_{E_kC}$; on $U$, $\delta_k(g)=0$ iff $C$ and $X$ give the same advice at $g$.

\begin{lemma}\label{lem:J10}
Assume (EP). Let $C\preceq_kX$ be a relevant comparison such that, for every $S\in\mathcal S_k$, $\delta_k\le0$ a.e.\ on $S_{a_S}$ for some depth $a_S$. Then there is a depth $m'$ with $C=X$ a.e.\ on $\mathcal E^k_{m'}$. This holds in particular if $X$ is correct at world $k$ a.e.\ on every $S_{a_S}$.
\end{lemma}
\begin{proof}[Proof sketch]
$\delta_k$ is $S$-positive for no $S\in\mathcal S_k$, so end-protection makes it $\ge0$ a.e.\ on some $S_{b_S}$; hence $\delta_k=0$ a.e.\ on $S_{\max(a_S,b_S)}$. The family is finite, so the maximum over $S$ is a depth.
\end{proof}
\status{proved (\src{J1} Lemma~J1.0); confirmed (\src{J1\_referee} \S2), one reader.}

\emph{Heaviness-free conflict digraph.} \src{J1} Thm.~J1.1 is the statement of Proposition~\ref{prop:HR1}, proved independently with the model $X=T_1$ on $\mathcal E_m$ (proved in \src{J1}; confirmed in \src{J1\_referee} \S2, one reader, with the same depth adjustment as \src{H12\_referee} \S2.4: the arrow needs the covering at a depth $j(M)\ge M$). \emph{Pattern switching.} Fix an order $(k_1,k_2,k_3)$ of the worlds. For $g\in\mathcal E_m$ let $r_m(g)\coloneqq\min\{r:g\in\mathcal E^{k_r}_m\}$. The \emph{switched model} $X_m$ equals $B^\psi_t$ on $U\setminus\mathcal E_m$, and at $g\in\mathcal E_m$ gives the advice that is correct at world $k_{r_m(g)}$. For $s=2,3$ the \emph{conflict set} is $\Xi^s_m\coloneqq\{g\in\mathcal E_m:r_m(g)<s,\ \sigma_{k_{r_m(g)}}(g)\ne\sigma_{k_s}(g)\}$, the points where $X_m$ follows an earlier world whose sign differs from that of world $k_s$. The order is \emph{aligned} if, for infinitely many $m$, the following holds for $s=2,3$ and every $n$:
\[
\textbf{(Al)}\qquad\text{there is a depth $d$ with }\ \mathcal E^{k_s}_d\cap\Xi^s_m\subseteq\textstyle\bigcup_{q<s}\mathcal E^{k_q}_n\quad\text{modulo null sets.}
\]
In words: going deep into a later world's ends, inside a fixed-depth conflicting end of an earlier world, takes one arbitrarily deep into the earlier worlds' ends.

\begin{theorem}\label{thm:J12}
Assume (EP). If some order of the worlds is aligned, (H$_U$) is contradicted.
\end{theorem}
\begin{proof}[Proof sketch]
Take $m$ as in (Al) with $\int_{\mathcal E_m}\sum_k\rho_k<\Phi^*_t$, the switched model $X\coloneqq X_m$, and $C$ as in Lemma~\ref{lem:J1C}. (1) $X$ is correct at world $k_1$ on $\mathcal E^{k_1}_m$, so Lemma~\ref{lem:J10} gives $m'_1$ with $C=X$ a.e.\ on $\mathcal E^{k_1}_{m'_1}$. (2) Take $d$ from (Al) with $s=2$ and $n\coloneqq m'_1$, and let $g\in\mathcal E^{k_2}_d$. If $r_m(g)=2$, or $r_m(g)=1$ and $\sigma_{k_1}(g)=\sigma_{k_2}(g)$, then $X$ is correct at world $k_2$ at $g$, so $\delta_{k_2}(g)\le0$. Otherwise $g\in\Xi^2_m$, so by (Al) $g\in\mathcal E^{k_1}_{m'_1}$ a.e., where $C=X$ and $\delta_{k_2}(g)=0$. So $\delta_{k_2}\le0$ a.e.\ on $\mathcal E^{k_2}_d$, and Lemma~\ref{lem:J10} gives $m'_2$ with $C=X$ a.e.\ on $\mathcal E^{k_2}_{m'_2}$. (3) The same argument with $n\coloneqq\max(m'_1,m'_2)$ and $s=3$ gives $C=X$ a.e.\ on $\mathcal E^{k_3}_{m'_3}$. (4) So every $\delta_k$ is eventually $0$ on every end of world $k$, all three dominator comparisons are base comparisons, and this contradicts Lemma~\ref{lem:J1C}.
\end{proof}
\status{proved (\src{J1} Thm.~J1.2); confirmed (\src{J1\_referee} \S3), one reader, including the order of quantifiers ($t$, then $m$, then $C$, then $m'(C)$, then $d$) and the pairwise sufficient condition below. It uses (b), integrability (Lemma~\ref{lem:J1C}), finiteness of the families (Lemma~\ref{lem:J10} and $\mu(\mathcal E_m)\to0$), end-protection, Theorem~\ref{thm:kkm}, Theorem~\ref{thm:M0closed}(a) and Lemma~\ref{lem:H11}. It uses no heaviness, no location of the ends and no $T_1$-correctness of the ends. \emph{Pairwise sufficient condition} (proved in the note). (Al) holds for $s$ at all large $m$ if, for every $r<s$ and every pair of ends $S\in\mathcal S_{k_r}$, $S'\in\mathcal S_{k_s}$ lying in a common orthant on which $\sigma_{k_r}\ne\sigma_{k_s}$, one of the following holds: (i) $S$ and $S'$ are eventually disjoint; (ii) $S'$ is shadowed by the earlier worlds: for every $n$ there is a depth $d$ with $S'_d\subseteq\bigcup_{q<s}\mathcal E^{k_q}_n$; (iii) for every $m$, $S'\cap S_m$ is shadowed by the earlier worlds: for all $m,n$ there is $d$ with $S'_d\cap S_m\subseteq\bigcup_{q<s}\mathcal E^{k_q}_n$. Lemma~\ref{lem:H21}(b) supplies (ii) for a $T_1$-wrong end $S'$ of world $k_s$ when every world of $V(S')$ that has ends comes earlier in the order.}

\begin{corollary}\label{cor:J13}
Assume (EP). Each of the following configurations contradicts (H$_U$), so none carries a cover of $T_1$.
\begin{enumerate}[leftmargin=*,label=(\alph*)]
\item The template above: worlds~$1$ and~$3$ have the common ends $A$, $B$, and $\mathcal S_2$ is any finite refined family (\src{J1} assumed $\mathcal S_2\subseteq\{A,B\}$; the referee showed that this is unnecessary). Take the order $(1,3,2)$. Concretely, the switched model accepts deep in $A$ and rejects deep in $B$; the dominator must copy it there, and then all three comparisons are base comparisons. So the germ exchanges above cannot be made: the configuration is inconsistent with (EP).
\item Any configuration in which, for any two ends of different worlds lying in a common orthant with conflicting signs, either (i) holds or each end is shadowed by the other. Then every order is aligned, whatever the number of conflict cycles.
\item A $2$-cycle $i\rightleftarrows j$ in which the third world $l$ has $\mathcal S_l=\varnothing$ (for instance an additive world, or one that adds up with an integrable density) and the arrow $i\to j$ is \emph{aligned}: every $T_1$-correct end $S'$ of $j$ and every $T_1$-wrong end $W$ of $i$ lying in a common orthant satisfy: for all $m,n$ there is $d$ with $S'_d\cap W_m\subseteq\mathcal E^i_n$. Here the order $(i,j,l)$ is aligned.
\end{enumerate}
\end{corollary}
\begin{proof}[Proof sketch]
Theorem~\ref{thm:J12} with the pairwise sufficient condition. In (a), with the order $(1,3,2)$: the ends of world~3 are ends of world~1, so (ii) holds for them. For a $T_1$-wrong end $S'$ of world~2, $V(S')\subseteq\{1,3\}$, both earlier, so (ii) holds by Lemma~\ref{lem:H21}(b). For a $T_1$-correct end $S'$ of world~2 the only conflicting pairs are with $A$ (in $(+,+,-)$, against world~3) and with $B$ (in $(-,+,+)$, against world~1); as ends of the earlier conflicting world these are $T_1$-correct, so both members of the pair are $T_1$-correct at worlds of opposite sign, and the pair is eventually disjoint: case (i) (\src{J1\_referee} \S4.1). In (c), check the pairs $S\in\mathcal S_i$, $S'\in\mathcal S_j$ on conflicting orthants: if $S'$ is $T_1$-wrong at $j$, then $V(S')\subseteq\{i,l\}$ and $\mathcal S_l=\varnothing$, so (ii) holds by Lemma~\ref{lem:H21}(b). If $S'$ is $T_1$-correct at $j$ and $S$ is $T_1$-correct at $i$, then beyond the purity depths $T_1$ is right at $i$ on $S_m$ and at $j$ on $S'_m$, with $\sigma_i\ne\sigma_j$ on the orthant, so $S_m\cap S'_m$ is null: case (i). If $S'$ is $T_1$-correct at $j$ and $S$ is $T_1$-wrong at $i$, the hypothesis is (iii). The step for world $l$ is vacuous.
\end{proof}
\status{proved (\src{J1} Cor.~J1.3); confirmed (\src{J1\_referee} \S4), one reader. \emph{Repairs.} (a) is the referee's stronger form, with no condition on $\mathcal S_2$ (proved in the referee report, \S4.1, one reader). (c) had a one-line gap: the note said that a $T_1$-correct $S'$ forces $S$ to be a $T_1$-wrong end of $i$, which is not literally true, since $T_1$'s advice varies within an orthant; the case of a $T_1$-correct $S$ is case (i), as in the proof above (\src{J1\_referee} \S4.3). In (b), ``shadowed'' means: for every $M$ there is $j$ with $S'_j\subseteq S_M$. For every configuration in (a)--(c), Question H2.Q has a negative answer. The note's ``for every $t\in I_\psi$'' adds nothing: Theorem~\ref{thm:J12} contradicts (H$_U$) itself, so one $t$ suffices. In particular the template of \src{H2} \S5.1, the only concrete conflict-cycle design in the corpus, is closed. With $\mathcal S_l=\varnothing$ as in (c) the conflict set is exactly the union of the pair pieces, so ``no aligned order'' is exactly ``both arrows misaligned'' there (\src{J1\_referee} \S4.3).}

\emph{What survives pattern switching.} By Theorem~\ref{thm:J12}, a cover under (EP) needs \emph{every} order of the worlds to violate (Al). In the simplest case, a $2$-cycle $i\rightleftarrows j$ next to a world without ends (Corollary~\ref{cor:J13}(c)), it needs \emph{both} arrows to be misaligned; call such a cycle \emph{doubly misaligned}. The arrow $i\to j$ is misaligned if some $T_1$-wrong end $W$ of $i$ and some $T_1$-correct end $S$ of $j$ in a common orthant, with $W$ shadowed by $S$ (Lemma~\ref{lem:H21}(b)), are such that for some $m$ the set $S\cap W_m$ is not shadowed by the ends of world $i$: at every depth of $S$, a non-null set of points of $W_m$ stays outside some fixed $\mathcal E^i_n$, i.e.\ is shallow in the ends of world $i$. \emph{Misaligned arrows exist as configurations of sets} (\src{J1} Example~J1.4; CHECKED; nothing about rankings is claimed). Let $O$ be a mixed orthant, $\hat r$ a direction in $O$, and $\gamma$ an arc of directions in $O$ with $\hat r$ in its relative interior. Let $W_m$ be the cone over the spherical cap of radius $\theta_m$ around $\hat r$, and $S_m$ the cone over the $\theta_m$-tube around $\gamma$, with $\theta_m\downarrow0$. Both are ends, and $W_m\subseteq S_m$, so $W$ is shadowed by $S$. For $\theta_n<\theta_m$, the set $S_d\cap W_m\setminus W_n$ contains the cone over the directions within $\theta_d$ of $\gamma$ at angular distance in $(\theta_n,\theta_m)$ from $\hat r$, which is non-null for every depth $d$. So if $W$ is the only end of world $i$ in $O$, the arrow is misaligned. With two such pairs, one for each arrow (for instance in the orthants $(+,+,-)$ and $(-,+,+)$ of the template), the configuration escapes Theorem~\ref{thm:J12}. (Round 19, \src{J1\_referee} \S5: confirmed at CHECKED level, with two precisions. The starting radius $\theta_0$ must be small enough that every cap and tube lies in $U\cap O$, and in a region where $T_1$'s advice is constant if the ends are to be $T_1$-wrong at $i$ and $T_1$-correct at $j$; otherwise start at a later depth. And ``escapes Theorem~\ref{thm:J12}'' holds when $\mathcal S_l=\varnothing$ and $W$ is the only end of $i$ in $O$, since then each of the six orders reduces to one of the two arrows. NUMERICS \src{j1ref\_py/j1ref\_geometry.py}, in the $\partial U$ collar of $(+,+,-)$ with $\varepsilon=0.3$, $s=0.9$: the tube of radius $0.04$ lies $18\%$ outside $U$ and $T_1$ accepts on $20\%$ of it; from radius $0.01$ on there are no exceptions; the cap lies inside the tube and the shell is hit at every tested depth $d\le10$.)

\begin{proposition}\label{prop:J15}
Assume (EP) and the situation of Corollary~\ref{cor:J13}(c) without the alignment hypothesis. Take the order $(i,j,l)$, $m$ large, the switched model $X$, and its dominator $C$ from Lemma~\ref{lem:J1C}. Then:
\begin{enumerate}[leftmargin=*,label=(\arabic*)]
\item $C=X$ a.e.\ on $\mathcal E^i_{m'}$ for some $m'$, and world $l$'s comparison is a base comparison;
\item world $j$'s comparison $C\preceq_jX$ is not a base comparison: it is positive at some $T_1$-correct end $S'$ of $j$, and deep in $S'$ its positive set lies in the shells $W_m\setminus W_{m'}$ of the $T_1$-wrong ends $W$ of world $i$ in the orthant of $S'$, where $C$ gives the advice that is correct at world $j$ against the advice of $X$, which is correct at world $i$.
\end{enumerate}
The same holds with $i$ and $j$ interchanged, under the order $(j,i,l)$.
\end{proposition}
\begin{proof}[Proof sketch]
(1) is step~(1) of Theorem~\ref{thm:J12}, and $\mathcal S_l=\varnothing$ makes every world-$l$ comparison a base comparison. If $C\preceq_jX$ were also a base comparison, Lemma~\ref{lem:J1C} would be contradicted; so, by end-protection, $\delta_j$ is $S'$-positive for some $S'\in\mathcal S_j$. If $S'$ were $T_1$-wrong at $j$, then $V(S')\subseteq\{i,l\}$ and Lemma~\ref{lem:H21}(b) would give $S'_n\subseteq\mathcal E^i_{m'}$ for large $n$, where $C=X$; so $S'$ is $T_1$-correct at $j$. Deep in $S'$, $X$ is wrong at $j$ only where it follows world $i$ against the sign of world $j$, i.e.\ in $\mathcal E^i_m$ on the orthant of $S'$. An end of $i$ that contains such a point and is $T_1$-correct at $i$ is eventually disjoint from $S'$ (the argument of Corollary~\ref{cor:J13}(c)), so the relevant ends of $i$ are $T_1$-wrong; points of $\mathcal E^i_{m'}$ are excluded by (1).
\end{proof}
\status{proved (\src{J1} Prop.~J1.5); confirmed (\src{J1\_referee} \S6), one reader, with the same one-line gap as Corollary~\ref{cor:J13}(c), closed in the proof above (the note said that the ends of $i$ there are $T_1$-wrong). Taking $M=m'$, the proposition implies that $i\to j$ is misaligned, consistent with Corollary~\ref{cor:J13}(c). So the exchanges that \src{H2} \S5.2 placed at the ends must happen in shells of world $i$'s $T_1$-wrong ends that stay at bounded depth in $W$ while going to every depth in $S'$. \emph{Why pattern switching stops here} (\src{J1} \S4.2). The depth $m'=m'(C)$ to which the dominator copies $X$ is not controlled in advance; Lemma~\ref{lem:J10} gives existence only. For a misaligned arrow the shell $S'\cap(W_m\setminus W_{m'})$ meets every $S'_n$ non-nully, for every $m'$, so $C$ can take world $j$'s advice there at every depth of $S'$: a protected, $S'$-positive comparison at world $j$, which (b) does not constrain. Reversing the order moves the problem to the other arrow. Three repairs fail at stated points: refining the ends (the pieces $S'\setminus W$ do not decrease, and the positive set is again confined to the same shell); giving each point to the world in whose end it is deeper (world $i$'s own comparison can then be positive near $\hat r$, where both ends are deep, and the first step is lost); and re-choosing $X$ after seeing $C$ (strict (G) gives the comparison in the wrong direction for transitivity at world $j$, and taking the dominator of the new model restarts the regress with a new uncontrolled depth). The referee confirms these as descriptions; no impossibility is claimed, and none follows (\src{J1\_referee} \S7). A small by-product of the third (proved in the referee report, \S7, one reader): let $X^*$ be $X$ made correct at world $j$ on $Q\coloneqq(\mathcal E^i_m\setminus\mathcal E^i_{m'})\cap\{\sigma_i\ne\sigma_j\}$; if $C\preceq_lX^*$ (for instance if $\sigma_l=\sigma_i$ on $Q$), then $C\not\preceq_jX^*$. So the dominators of $X$ and $X^*$ are different models, which is the regress just described, not an obstruction.}

\emph{Question J1.Q} (open; no conjecture is stated in \src{J1}; the wording is corrected as \src{J1\_referee} \S\S8.1, 10 requires). Assume (EP), and let $\vec{\mathcal H}$ be a $2$-cycle $i\rightleftarrows j$ with both arrows misaligned, whose third world $l$ is incident to no arrow. For every tuple $\psi$ of integrable base endorsers (such tuples are admissible, Lemma~\ref{lem:H11}), every $t\in I_\psi$ and every large $m$, can the dominators carry the shell exchanges of Proposition~\ref{prop:J15} under every axiom: strict (G), (IB), transitivity, closed sections for every model, and competitiveness of $T_1$? If not, no cover in class (EP) has this configuration, which is case (R-i) below. \emph{Round 19 corrections.} (1) The note asked the question ``for every integrable admissible $\psi$''. That is too broad: Proposition~\ref{prop:J15} uses (b) for $\psi$, and for an admissible $\psi$ without (b) the unendorsed dominator comparison can be a base comparison, so that nothing forces shells (the defect found in \src{H2} \S5.2). (2) ``$\mathcal S_l=\varnothing$'' is replaced by ``$l$ incident to no arrow'', which Corollary~\ref{cor:JR3}(c) below allows. (3) \src{J1} states the consequence of a negative answer as ``no cover of $T_1$ on $\M_U$ has finitely many end-protected worlds with integrable base endorsers''. That inference is unproved (\src{J1\_referee} \S8.2). A negative answer closes only case (R-i). The configurations in class (EP) with no aligned order that are not of this form ($3$-cycles, and $2$-cycles whose third world has ends) are not analysed in \src{J1}; Theorem~\ref{thm:JR2} and Corollary~\ref{cor:JR3} below narrow them. (4) \src{J1} \S0 also says that no conflict cycle survives if one of its arrows is aligned. This is proved for $2$-cycles whose third world is incident to no arrow, and for pure $3$-cycles (Corollary~\ref{cor:JR3}(b), (c)); it does not follow from Theorem~\ref{thm:JR2} for digraphs with more arrows. If the answer to J1.Q is yes, Example~J1.4 (a cap inside a tube, twice, in two conflicting orthants) is the geometric template. \emph{Scope limits of \src{J1}} (\src{J1} \S6). With countably many ends, Lemma~\ref{lem:J10} takes a maximum over the family and $\mu(\mathcal E_m)\to0$ needs summable depths; the depths $m'_S$ are then not uniform, and (Al) would have to be required end by end, with only finitely many earlier ends meeting each later end deeply. This is not proved. Non-integrable base endorsers are untouched: Lemma~\ref{lem:J1C} needs $\int_{\mathcal E_m}\sum_k\rho_k\to0$. The installation step of \src{BRIEF15} (a three-world ranking with every axiom) was not reached: the template is closed, and the surviving configuration has no world orders yet. No scripts were used in \src{J1}. (Round 19: the countable forms of Proposition~\ref{prop:HR1} and Theorem~\ref{thm:J12} are proved in \src{K1}, with alignment required end by end; Theorems~\ref{thm:K14} and~\ref{thm:K15} in part~(f$''$).)

\emph{Round 19, \src{J1\_referee}: alignment of orders from alignment of arrows} (proved in a referee report, \src{J1\_referee} \S8, one reader). Assume (EP), with the notation of Theorem~\ref{thm:J12}. Call an arrow $a\to b$ of $\vec{\mathcal H}$ \emph{aligned} if it satisfies the hypothesis of Corollary~\ref{cor:J13}(c): every $T_1$-wrong $W\in\mathcal S_a$ and every $T_1$-correct $S'\in\mathcal S_b$ lying in a common orthant satisfy, for all $m,n$, $S'_d\cap W_m\subseteq\mathcal E^a_n$ for some depth $d$; otherwise it is \emph{misaligned}. Let $\vec{\mathcal H}_{\mathrm{mis}}$ be the set of misaligned arrows.

\begin{theorem}\label{thm:JR2}
Assume (EP), and fix an order $(k_1,k_2,k_3)$ of the worlds. Suppose that:
\begin{enumerate}[leftmargin=*,label=(\alph*)]
\item for every $r<s$, every $T_1$-wrong $W\in\mathcal S_{k_r}$ and every $T_1$-correct $S'\in\mathcal S_{k_s}$ lying in a common orthant, and all large $m$: for every $n$ there is a depth $d$ with $S'_d\cap W_m\subseteq\bigcup_{q<s}\mathcal E^{k_q}_n$ (every forward arrow is aligned relative to the earlier worlds);
\item for every $T_1$-wrong $S'\in\mathcal S_{k_2}$, $T_1$-correct $S\in\mathcal S_{k_1}$ and $T_1$-correct $T\in\mathcal S_{k_3}$ lying in a common orthant, and all large $m$: for every $n$ there is a depth $d$ with $S'_d\cap S_m\cap T_n\subseteq\mathcal E^{k_1}_n$.
\end{enumerate}
Then the order is aligned, so (H$_U$) is contradicted. Condition (b) holds vacuously unless $\vec{\mathcal H}$ has both arrows $k_2\to k_1$ and $k_2\to k_3$.
\end{theorem}
\begin{proof}[Proof sketch]
Verify the pairwise sufficient condition of Theorem~\ref{thm:J12}. Take $m$ beyond every orthant-purity and $T_1$-purity depth and beyond every depth from which an eventually disjoint pair is disjoint. Let $S\in\mathcal S_{k_r}$, $S'\in\mathcal S_{k_s}$, $r<s$, lie in a common orthant on which $\sigma_{k_r}\ne\sigma_{k_s}$. (1) If both are $T_1$-correct, or both $T_1$-wrong, at their worlds, then on $S_m\cap S'_m$ the single advice of $T_1$ would be right (respectively wrong) at two worlds of opposite sign; so the intersection is null, case (i). (2) $S$ $T_1$-wrong and $S'$ $T_1$-correct: this is (a), i.e.\ case (iii). (3) $S$ $T_1$-correct, $S'$ $T_1$-wrong, $s=3$: $V(S')\subseteq\{k_1,k_2\}$, both earlier, so (ii) holds by Lemma~\ref{lem:H21}(b). (4) $S$ $T_1$-correct, $S'$ $T_1$-wrong, $s=2$: given $N$, apply Lemma~\ref{lem:H21}(b) at depth $n\coloneqq\max(N,d_0)$; then $S'_d$ lies in the union of the $T_n$ over the $T_1$-correct ends $T$ of the worlds of $V(S')\subseteq\{k_1,k_3\}$. The terms with $T\in\mathcal S_{k_1}$ lie in $\mathcal E^{k_1}_N$. A $T\in\mathcal S_{k_3}$ that does not meet $S'$ at every depth is eventually disjoint from it, so its term is null for $n,d\ge d_0$. For the remaining $T$ (an arrow $k_2\to k_3$), (b) gives $S'_d\cap S_m\cap T_n\subseteq\mathcal E^{k_1}_N$. So $S'_d\cap S_m\subseteq\mathcal E^{k_1}_N$, which is (iii). If $S'$ does not meet $S$ at every depth, $S_m\cap S'_d$ is null, which gives the vacuity statement together with the argument for $T$.
\end{proof}
\status{proved in a referee report (\src{J1\_referee} \S8.3, Thm.~R.2), one reader.}

\begin{corollary}\label{cor:JR3}
Assume (EP). Alignment of an arrow $a\to b$ implies condition (a) of Theorem~\ref{thm:JR2} for that arrow in every order in which $a$ precedes $b$.
\begin{enumerate}[leftmargin=*,label=(\alph*)]
\item If $\vec{\mathcal H}_{\mathrm{mis}}$ has no directed cycle, condition (a) holds in every order in which all misaligned arrows point backward (such orders exist). So (H$_U$) is contradicted as soon as condition (b) holds for one of these orders, and that is automatic if the middle world of the order does not have arrows to both other worlds.
\item Let $\vec{\mathcal H}=\{a\to b,\ b\to c,\ c\to a\}$ be a pure $3$-cycle. In the order $(a,c,b)$ the only forward arrow is $a\to b$, and the middle world $c$ has one outgoing arrow. So if $a\to b$ is aligned, even only relative to $\mathcal E^a\cup\mathcal E^c$ in place of $\mathcal E^a$, (H$_U$) is contradicted. A pure $3$-cycle can carry a cover only if all three arrows are misaligned in this weaker sense.
\item Let $\vec{\mathcal H}=\{i\to j,\ j\to i\}$, with the third world $l$ incident to no arrow ($\mathcal S_l$ arbitrary, possibly nonempty). The orders $(i,j,l)$ and $(l,i,j)$ are aligned iff $i\to j$ is aligned, and the mirror orders behave likewise for $j\to i$. Proposition~\ref{prop:J15} holds verbatim with the order $(l,i,j)$.
\end{enumerate}
\end{corollary}
\begin{proof}[Proof sketch]
(a) and (b) are Theorem~\ref{thm:JR2}. In (c), with $l$ first, world $l$'s step closes; world $i$'s step needs nothing, since no arrow involves $l$; in world $j$'s step a $T_1$-wrong end of $j$ is covered by the $T_1$-correct ends of $i$, since those of $l$ are eventually disjoint from it; and the conflict points in $S'$ led by $l$ are null.
\end{proof}
\status{proved in a referee report (\src{J1\_referee} \S8.3, Cor.~R.3), one reader. NUMERICS (\src{j1ref\_py/j1ref\_orders.py}): all ten digraphs on three vertices with a directed cycle, up to relabelling. If condition (b) is assumed to hold, a set of misaligned arrows escapes Theorem~\ref{thm:JR2} iff it contains a directed cycle (0 mismatches over all subsets): a cycle of misaligned arrows has a forward arrow in every order. If condition (b) is assumed to fail wherever it is needed, there are further survivors with no misaligned cycle, exactly when the unique order that makes $\vec{\mathcal H}_{\mathrm{mis}}$ point backward has a middle world with two outgoing arrows; example $\vec{\mathcal H}=\{1\to2,1\to3,2\to1\}$ with $\vec{\mathcal H}_{\mathrm{mis}}=\{1\to3,2\to1\}$. For the complete digraph, Theorem~\ref{thm:JR2} alone excludes nothing if (b) fails.}

\emph{The precise residual of class (EP)} (\src{J1\_referee} \S8.4). By Theorems~\ref{thm:J12}, \ref{thm:JR2}, Proposition~\ref{prop:HR1} and Corollary~\ref{cor:JR3}, a cover of $T_1$ in class (EP) (the rankings with (EP); see part~(g)) needs one of the following.
\begin{itemize}[leftmargin=*]
\item[\textbf{(R-i)}] $\vec{\mathcal H}$ is a $2$-cycle whose third world is incident to no arrow, and both arrows are misaligned. A negative answer to Question~J1.Q, in the corrected wording above, closes exactly this case.
\item[\textbf{(R-ii)}] $\vec{\mathcal H}$ is a pure $3$-cycle $a\to b\to c\to a$, and each arrow $a\to b$ is misaligned relative to $\mathcal E^a\cup\mathcal E^c$. No analogue of Proposition~\ref{prop:J15} is proved. By the argument of Theorem~\ref{thm:JR2}, in the order $(a,c,b)$ the closure reaches worlds $a$ and $c$, and the positive set of world $b$'s comparison lies in shells of the $T_1$-wrong ends of $a$ that are shallow in $\mathcal E^a\cup\mathcal E^c$ (CHECKED in \src{J1\_referee}, not proved).
\item[\textbf{(R-iii)}] $\vec{\mathcal H}$ has at least three arrows and is not a pure $3$-cycle; and either $\vec{\mathcal H}_{\mathrm{mis}}$ contains a directed cycle, or every order that makes $\vec{\mathcal H}_{\mathrm{mis}}$ point backward has a middle world with arrows to both other worlds at which condition (b) of Theorem~\ref{thm:JR2} fails. Here the third world's comparison is not shown to be a base comparison, and no description by shells is available.
\end{itemize}
\emph{Question J1.Q$'$} (\src{J1\_referee} \S8.4; open, no conjecture stated). Under (EP), if no order of the worlds satisfies (Al), is (H$_U$) contradictory? A positive answer excludes every cover in class (EP); Question~J1.Q is its case (R-i).

\emph{What survives round 17} (\src{H1} \S6, corrected in \src{H12\_referee} \S\S1.7, 3). \src{H1} stated that a cover needs ``exactly one'' of three cases: (A) a non-localizable failure at one world, next to two integrable worlds; (B) mutual sheltering between two exceptional worlds; and endorsers at the other worlds that are only finitely additive. The list is neither exclusive (the last two cases can occur together) nor exhaustive. The corrected statement is: a cover of $T_1$ on $\M_U$ that is not excluded by Theorems~\ref{thm:H12}, \ref{thm:H13} and Proposition~\ref{prop:HR1} has \emph{at least one} of the following properties (the referee's labels).
\begin{itemize}[leftmargin=*]
\item[\textbf{(A$'$)}] \emph{Non-localizable failure.} Some world $i$ is such that, for every localizing family $\Theta$ and every a.e.\ finite $\rho_i>0$, some $\Theta$-localized relevant comparison of world $i$ is not endorsed by $\rho_i\mu$ in the cancellation form. (This is (A) without its clause on the other two worlds. It contains the case of worlds that add up only through purely finitely additive weights, which have no a.e.\ finite endorser. \emph{Round 19:} by Theorem~\ref{thm:K31}(c)--(e), every world that adds up is (Loc$_{\mathrm{int}}$), since the countably additive part of an endorsing weight is positive a.e.\ and endorses every comparison localized off the carrier of its purely finitely additive part; so no world that adds up falls under (A$'$). By Theorem~\ref{thm:K31}, (A$'$) consists of worlds with a certificate and worlds in the class (N) of part~(f$''$).) What is left here is a failure carried by comparisons that leave every set of small measure alone: an (M1) obstruction, or an (M2) family whose sources and targets have measure bounded below on a fixed core. In two worlds next to an additive world no certificate of the type of Theorem~\ref{thm:D21} has its sources and targets in such a core (Theorem~\ref{thm:F14}), which is evidence against (A$'$); the three-world case is open.
\item[\textbf{(B)}] \emph{Sheltering.} The localizations of two worlds cannot be made $T_1$-correct, so the compatibility that Theorem~\ref{thm:H13} needs fails: for instance, some world has an end on which $T_1$ is wrong at that world at every depth, inside $R_{ij}$, sheltered by the protected set of another world. If every world is end-protected with a finite family and has an integrable base endorser, this happens only along a directed cycle of $\vec{\mathcal H}$ (Theorem~\ref{thm:H23}, Proposition~\ref{prop:HR1}); with countably many ends an arrow of $\vec{\mathcal H}$ is not excluded. \emph{Round 19:} with countably many ends and integrable base endorsers, only cyclic configurations with no aligned order remain (Theorems~\ref{thm:K14}, \ref{thm:K15}; \src{K1\_referee} \S9); with finitely many ends, only the cases (R-i)--(R-iii) above. Sheltering in general, without ends, is what the patch theorem of part~(f$''$) leaves (Corollary~\ref{cor:K35p}(c)).
\item[\textbf{(D)}] \emph{Several non-integrable endorsers.} Every world has an a.e.\ finite endorser of its localized comparisons, and the localizations are compatible, but at least two worlds' endorsers are non-integrable, or one world's endorser is non-integrable and the other worlds have only localized, not full, integrable endorsers. This is a genuine obstacle to the method: if $\min(\rho_i,\rho_j)$ is not integrable on $R_{ij}$, every model has infinite cost at world $i$ or at world $j$ there, so no model with finite costs at both worlds is available near $B^\rho_t$.
\end{itemize}
Part~(g) uses this list for the residual.

\paragraph{(f$''$) Round 19: what a world that does not add up can be, finite certificates, and countably many ends.} Round 19 (\src{BRIEF16}) asked for a cover of $T_1$, or a theorem that settles $T_1$ on $\M_U$, or the exact mechanism a cover must use. Three notes were written, and none produced a cover or a settling theorem. \src{K3} classifies the worlds that do not add up by the Hahn--Banach and Yosida--Hewitt theorems and identifies the exact residual; it also proves a patch theorem that unifies Theorems~\ref{thm:H12}, \ref{thm:H13} and pattern switching and extends them to $\sigma$-finite densities at several worlds. \src{K2} treats worlds with a finite certificate (mode (M1)). \src{K1} extends the exclusions for end-protected worlds from finitely many to countably many ends. Each note was refereed once (\src{K1\_referee}, \src{K2\_referee}, \src{K3\_referee}), and \src{J1} was refereed in \src{J1\_referee} (part~(f$'$)). No referee found anything wrong in the mathematics. The repairs concern statements of scope and one sufficient condition; they are incorporated below, and the referees' own results (Proposition~\ref{prop:KR62}, Proposition~\ref{prop:KQA}, the configuration after Theorem~\ref{thm:K15}) are marked as proved in a referee report. Unless stated otherwise every result of this part is \emph{proved in a note and confirmed by one referee}. The blocks are given in the order of the argument of part~(g): \src{K3}, then \src{K2}, then \src{K1}. No assumption is added anywhere.

\emph{Round 19, note \src{K3}: the Yosida--Hewitt classification} (\src{K3}; \src{K3\_referee}). Fix a world $i$ of a ranking on $\M_U$ ($n\ge2$, any credal set $K'\subseteq\Di$) with strict (G) and (IB); the other axioms enter only through $\C_L$. Let $\mathfrak W\coloneqq\mathfrak W_i$ be the set of vectors $\delta=1_{E_iX}-1_{E_iB}$ of the relevant comparisons $B\preceq_iX$ (Definition~\ref{def:addup}). Coherent and regular models are correct on $N_+$ at every world, so every $\delta\in\mathfrak W$ vanishes on $N_+$, and everything below lives on $M$. Let $\mathrm{ba}$ be the space of bounded finitely additive set functions on the measurable subsets of $M$ that vanish on null sets; it is the dual of $L^\infty(M,\mu)$. For $\mathfrak G\subseteq\mathfrak W$ let $\mathfrak N(\mathfrak G)$ be the set of $\nu\in\mathrm{ba}$ with $\nu\ge0$, $\nu(M)=1$ and $\nu(\delta)\ge0$ for every $\delta\in\mathfrak G$ (the bounded endorsers of $\mathfrak G$), and put $\mathfrak N\coloneqq\mathfrak N(\mathfrak W)$. By the Yosida--Hewitt theorem each $\nu\ge0$ in $\mathrm{ba}$ is uniquely $\nu=\nu_c+\nu_p$, with $\nu_c=\rho_\nu\mu$ for some $0\le\rho_\nu\in L^1(M)$ and $\nu_p\ge0$ purely finitely additive. Modes are taken world by world: world $i$ is in mode (M0) if it adds up, in mode (M1) if it carries a finite obstruction (a \emph{certificate}, \S\ref{sec:addup-def}), and in mode (M2) otherwise. A family $\mathfrak G$ satisfies \textbf{(NC$_{\mathfrak G}$)} if no finite nonnegative combination of its members is $\le-1$ a.e.; a world that is not in mode (M1) satisfies (NC$_{\mathfrak W}$), since such a combination would be a certificate. Localizing families $\Theta$ are as in part~(f$'$), the empty family included; write $\Theta\le\Theta'$ if $\Theta_m\subseteq\Theta'_m$ for all $m$, $(\Theta\cup\Theta')_m\coloneqq\Theta_m\cup\Theta'_m$, and $\mathfrak W_\Theta$ for the $\Theta$-localized members of $\mathfrak W$. World $i$ is \textbf{localizable}, \textbf{(Loc)}, if for some localizing $\Theta$ and some measurable $\rho$ with $0<\rho<\infty$ a.e.\ on $M$, $\rho\mu(\delta^-)\le\rho\mu(\delta^+)$ in $[0,\infty]$ for every $\delta\in\mathfrak W_\Theta$ (the cancellation form): this is the hypothesis of Theorem~\ref{thm:H12} on its third world. It is \textbf{(Loc$_{\mathrm{int}}$)} if this holds with $\rho\in L^1(M)$; the condition then reads $\int\rho\delta\,d\mu\ge0$. (\src{K3} calls these classes (S) and (S$_{\mathrm{int}}$); they are renamed because (S) is the axiom of closed competitive choice.) World $i$ is \textbf{(EP$_{\mathrm b}$)} if it is end-protected with a countable family of ends and its base comparisons are endorsed in the cancellation form by an a.e.\ finite $\rho>0$. The residual class is \textbf{(N)}: mode (M2) and not (Loc). Finally, for $\mathfrak G\subseteq\mathfrak W$, a non-null $B'\subseteq M$ and $c>0$, a \emph{$c$-ratio certificate on $B'$ from $\mathfrak G$} is a pair $(\hat\delta,g)$, with $\hat\delta$ a finite nonnegative combination of members of $\mathfrak G$ and $g\in L^\infty$, $g\ge0$, supported in $B'$, $\int g\,d\mu=1$, such that $\hat\delta\le c-g$ a.e.\ on $M$. The combination removes one unit of $\mu$-mass from $B'$, spread as $g$, while its positive excess is at most $c$. The certificates that make $c_a$ invisible to bounded endorsers in Theorem~\ref{thm:D21} are of this kind, with $g$ uniform on $c_a$. The blind sets $\Lambda_\Theta$ below are unrelated to the error indicator $\Lambda_i$ of \S\ref{sec:CPa-def}.

\begin{theorem}\label{thm:K31}
Let world $i$ have strict (G) and (IB).
\begin{enumerate}[leftmargin=*,label=(\alph*)]
\item $\mathfrak N\ne\varnothing$ iff (NC$_{\mathfrak W}$) holds. In particular, if world $i$ is not in mode (M1), then $\mathfrak N(\mathfrak W_\Theta)\supseteq\mathfrak N\ne\varnothing$ for every localizing $\Theta$.
\item Let $\mathfrak G\subseteq\mathfrak W$ and $\nu\in\mathfrak N(\mathfrak G)$. There is a localizing family $\Theta^\nu$ with $\nu_p(M\setminus\Theta^\nu_m)\le2^{-m}$ for every $m$, and $\int\rho_\nu\delta\,d\mu\ge0$ for every $\Theta^\nu$-localized $\delta\in\mathfrak G$.
\end{enumerate}
Assume (NC$_{\mathfrak W}$) for the rest of the theorem.
\begin{enumerate}[leftmargin=*,label=(\alph*),resume]
\item There is $\nu^*\in\mathfrak N$ such that the set $\Lambda_\varnothing\coloneqq M\setminus\{\rho_{\nu^*}>0\}$, unique up to null sets, satisfies $\rho_\nu=0$ a.e.\ on $\Lambda_\varnothing$ for every $\nu\in\mathfrak N$: on $\Lambda_\varnothing$ every bounded endorser is purely finitely additive. World $i$ adds up iff $\Lambda_\varnothing$ is null. In mode (M2), $\Lambda_\varnothing$ is non-null.
\item For a localizing $\Theta$ let $\Lambda_\Theta$ be the set of (c) for $\mathfrak W_\Theta$ in place of $\mathfrak W$; $\Lambda_\varnothing$ is the case of the empty family. Then $\Theta\mapsto\Lambda_\Theta$ is antitone, and $\Lambda_\Theta\subseteq\Lambda_\varnothing$. The \emph{non-localizable blind set} $\Lambda_\infty\coloneqq\operatorname*{ess\,inf}_\Theta\Lambda_\Theta$ equals $\Lambda_{\Theta^\infty}$ for some localizing $\Theta^\infty$. There are a localizing $\hat\Theta\ge\Theta^\infty$ and $\rho_\infty\in L^1(M)$, $\rho_\infty\ge0$, with $\{\rho_\infty>0\}=M\setminus\Lambda_\infty$ a.e.\ and $\Lambda_{\hat\Theta}=\Lambda_\infty$, such that $\int\rho_\infty\delta\,d\mu\ge0$ for every $\delta\in\mathfrak W_{\hat\Theta}$.
\item (Loc$_{\mathrm{int}}$) holds iff $\Lambda_\infty$ is null. Hence (Loc$_{\mathrm{int}}$) implies (Loc), and every world in (N) has a non-null $\Lambda_\infty$.
\item Let $B\subseteq M$ be non-null and $\Theta$ localizing. Then $B\subseteq\Lambda_\Theta$ (mod null) iff for every non-null $B'\subseteq B$ and every $c>0$ there is a $c$-ratio certificate on $B'$ from $\mathfrak W_\Theta$. Hence $B\subseteq\Lambda_\infty$ iff this holds for every localizing $\Theta$.
\end{enumerate}
\end{theorem}
\begin{proof}[Proof sketch]
(a) If (NC$_{\mathfrak W}$) holds, $-1$ is not in the sup-norm closure of the convex cone generated by $\mathfrak W$ and the nonnegative functions (otherwise a combination is $\le-1+\eta$ with $\eta<1$, and rescaling gives one $\le-1$). Hahn--Banach separation gives a functional that is $\ge0$ on that cone and positive at $1$; normalised, it lies in $\mathfrak N$. Conversely $\nu\in\mathfrak N$ gives $\nu(\sum\lambda\delta)\ge0>\nu(-1)$. (b) Put $\pi\coloneqq\nu_p$ and $\lambda(A)\coloneqq\inf_{B\subseteq A}(\pi(B)+\mu(A\setminus B))$, the infimum $\pi\wedge\mu$ in $\mathrm{ba}$. It is $\le\mu$, hence countably additive, hence $0$ because $\pi$ is purely finitely additive. So for every $\varepsilon>0$ there is $E$ with $\mu(E)<\varepsilon$ and $\pi(M\setminus E)<\varepsilon$. Take such $E_k$ for $\varepsilon=2^{-k}$ and $\Theta^\nu_m\coloneqq\bigcup_{k\ge m}E_k$. For $\delta\in\mathfrak G$ with $\delta\le0$ on $\Theta^\nu_m$ and $m'\ge m$, $0\le\nu(\delta)\le\int\rho_\nu\delta\,d\mu+0+2^{-m'}$; let $m'\to\infty$. (c) Countable convex combinations of members of $\mathfrak N$ lie in $\mathfrak N$, and their densities are the combinations of the densities (the criterion of (b) shows that the purely finitely additive parts stay purely finitely additive). So the supports $\{\rho_\nu>0\}$ are closed under countable unions, and a largest one exists because $\mu$ is finite. If $\Lambda_\varnothing$ is null, $\nu^*+\mu|_{N_+}$, normalised, is a weight as in Definition~\ref{def:addup}(b). Conversely a strictly positive endorsing weight has an a.e.\ positive density (\src{N19} Lemma~R1, re-derived in \src{K3}). (d) $\Theta\le\Theta'$ gives $\mathfrak W_{\Theta'}\subseteq\mathfrak W_\Theta$, hence $\Lambda_{\Theta'}\subseteq\Lambda_\Theta$. The family is downward directed ($\Lambda_{\Theta\cup\Theta'}\subseteq\Lambda_\Theta\cap\Lambda_{\Theta'}$), so its essential infimum is attained along a countable subfamily $\Theta^{(k)}$. The diagonal union $\Theta^\infty_m\coloneqq\bigcup_k\Theta^{(k)}_{m+j_k}$, with $\mu(\Theta^{(k)}_{j_k})<2^{-k}$, is localizing, and a $\Theta^\infty$-localized $\delta$ is $\Theta^{(k)}$-localized for every $k$. Apply (c) to $\mathfrak W_{\Theta^\infty}$ and (b) to the resulting endorser, which gives a family $\Theta'$, and put $\hat\Theta\coloneqq\Theta^\infty\cup\Theta'$. (e) An integrable $\rho>0$ endorsing $\mathfrak W_\Theta$ gives a member of $\mathfrak N(\mathfrak W_\Theta)$ with full support, so $\Lambda_\Theta$ and hence $\Lambda_\infty$ are null; conversely use (d). (f) For $\nu\in\mathrm{ba}_+$, $\nu\ge c\mu|_{B'}$ iff $\rho_\nu\ge c$ a.e.\ on $B'$. So $B\not\subseteq\Lambda_\Theta$ iff some $\nu\in\mathfrak N(\mathfrak W_\Theta)$ has $\nu\ge c\mu|_{B'}$ for some non-null $B'\subseteq B$ and $c>0$. Hahn--Banach applied to the cone of the functions $\hat\delta+f-c\bigl(\int_{B'}f\,d\mu\bigr)\cdot1+h$ ($\hat\delta$ a combination from $\mathfrak W_\Theta$, $f,h\ge0$) shows that no such $\nu$ exists iff ratio certificates exist; (NC$_{\mathfrak W}$) excludes the case in which $f$ puts little mass on $B'$.
\end{proof}
\status{proved (\src{K3} Thm.~K3.1); confirmed (\src{K3\_referee} \S1), one reader. \emph{Repairs.} Parts (c)--(f) need $\mathfrak N\ne\varnothing$, i.e.\ (NC$_{\mathfrak W}$), which the note's standing hypotheses omitted; it holds whenever world $i$ is not in mode (M1), the only case in which (c)--(f) are used (alternatively, put $\Lambda_\Theta\coloneqq M$ when $\mathfrak N(\mathfrak W_\Theta)=\varnothing$). Only $\Lambda_\infty$ is canonical (mod null); $\hat\Theta$ and $\rho_\infty$ depend on choices (the note called $\hat\Theta$ canonical). The endorser of (a) is bounded and nonnegative, not strictly positive. \emph{An observation of the referee.} The inclusion of (N) in $\{\Lambda_\infty\text{ non-null}\}$ may be strict: a world may be (Loc) with a non-integrable density while $\Lambda_\infty$ is non-null. Truncating $\rho$ at level $m$ and adding $\{\rho>m\}$ to the localizing family does not give (Loc$_{\mathrm{int}}$), because the level needed depends on the comparison. \emph{Reading.} The route of Hahn--Banach, Yosida--Hewitt and localization gives exactly (b): the purely finitely additive part of an endorser always sits on sets of small measure, and its countably additive part endorses every comparison localized off them. The route stops only where that countably additive part vanishes on a non-null set. For world~1 of Theorem~\ref{thm:D21}, $\Lambda_\varnothing=c_a$ (\src{K3}), but localizing at the end of the chain repairs this: $\Lambda_\infty$ is null, so that world is (Loc$_{\mathrm{int}}$). For the tilt world of Theorem~\ref{thm:K33} the zero set is all of $M$, and no localization repairs it. So every world that does not add up either is in mode (M1), or has no finite certificate and a non-null $\Lambda_\infty$, or has no finite certificate and is (Loc$_{\mathrm{int}}$).}

\begin{proposition}\label{prop:K32}
\begin{enumerate}[leftmargin=*,label=(\alph*)]
\item Every world is end-protected with the empty family of ends, and then every relevant comparison is a base comparison.
\item A world that is (EP$_{\mathrm b}$) is (Loc); it is (Loc$_{\mathrm{int}}$) if its base endorser is integrable.
\item Hence the statement ``every world that does not add up is in mode (M1), or (EP$_{\mathrm b}$), or (Loc)'' is equivalent to (N)$\,=\varnothing$.
\end{enumerate}
\end{proposition}
\begin{proof}[Proof sketch]
(a) holds by the definition: the second clause of end-protection is vacuous. (b) is the deep-ends argument of part~(f$'$): with summable depths, a localized comparison is $S$-positive for no end $S$, so end-protection makes it eventually $\ge0$, hence eventually $=0$, on every end: it is a base comparison. (c) follows from (b).
\end{proof}
\status{proved (\src{K3} Prop.~K3.2); confirmed (\src{K3\_referee} \S2), one reader. So bare end-protection carries no information, and the content of end-protection is its base density. The classification ``(M1), end-protected, or localizable'' that \src{K3} set out to prove is therefore exactly the statement (N)$\,=\varnothing$. (The referee notes that this classification is not stated in \src{BRIEF16} as filed.)}

\begin{theorem}\label{thm:K33}
Let $n=2$, in the setting of part~(e) (the interval $K'=[k_1,k_2]\subseteq(0,1)$, the path $b\mapsto C_b$, the chain $c_b=E_1C_b\cap M$ with $\bar\mu(c_b)=b$, $d_b=M\setminus c_b$, and the chain coordinate $x$, with $x\coloneqq2$ on $N_+$). Let world~2 be additive, $X\preceq_2Y$ iff $\bar\mu(E_2X)\le\bar\mu(E_2Y)$, and let world~1 be the tilt world of part~(e): $X\preceq_1Y$ iff, for some $\kappa>0$, the vector $\delta\coloneqq1_{E_1Y}-1_{E_1X}$ has $\int_{\{x\ge y\}}e^{\kappa x}\delta\,d\bar\mu\ge0$ for every $y\in[0,2]$. Take the selection $p\mapsto C_{p_2}$.
\begin{enumerate}[leftmargin=*,label=(\alph*)]
\item This ranking satisfies strict (G) and (IB) at both worlds, (LA), competitiveness of every $C_b$ on its support (indeed (K) for every model), (C) with $\C_L=\{C_b:b\in[0,1]\}$, (LF) at every interior $C_b$, and closed world-2 sections for every model. The world-1 sections are not closed.
\item World~1 is in mode (M2), and its non-localizable blind set is $\Lambda_\infty=M$.
\item World~1 is not (Loc): for every localizing $\Theta$ and every measurable $\rho$ with $0<\rho<\infty$ a.e., some accepted, relevant, $\Theta$-localized world-1 comparison has $\rho\mu(\delta^-)>\rho\mu(\delta^+)$. So world~1 is in (N).
\end{enumerate}
\end{theorem}
\begin{proof}[Proof sketch]
(a) For $\kappa<\kappa'$, Stieltjes integration by parts with the nondecreasing weight $e^{(\kappa'-\kappa)x}$ shows that the $\kappa$-cone lies in the $\kappa'$-cone; so the accepted set is an increasing union of convex cones. It contains every $\delta\ge0$ and, looking at $y=0$, no nonzero $\delta\le0$; Lemma~\ref{lem:vectorform} gives strict (G) and (IB). For a model $X$ put $A\coloneqq E_1X\cap M$, $N\coloneqq N_+\setminus X$, $P\coloneqq c_b\setminus A$, $Q\coloneqq A\setminus c_b$; the world-1 vector of $C_b\preceq_1X$ is $1_Q-1_P+1_N$ and the world-2 vector is $1_P-1_Q+1_N$. If $Q\cup N$ is non-null, every tail is positive for $\kappa$ large, so $C_b\preceq_1X$; otherwise $C_b\preceq_2X$. This is (K). A weak dominator of $C_b$ must have $Q$ and $N$ null (tail of the reversed world-1 vector at $y=b$), and then $P$ null (world~2). Every $X\ne C_b$ is strictly dominated by $C_{b_2}$, $b_2\coloneqq\min(1,\max(0,\bar\mu(A)-\bar\mu(N)))$ (the repair of \src{Q12\_referee}), directly and without Theorem~\ref{thm:kkm}, which needs the closed sections this system lacks. So $\C_L=\{C_b\}$, and (C) and (LA) hold. The world-1 section of $c_x\setminus Z$, $Z\subseteq c_{x-\eta}$ non-null, is $[0,x)$. (b) The world-1 error sets of $\C_L$ form a chain, so there is no certificate (\src{N17} N17.3, Lemma~3.1). For non-null $Q\subseteq d_b$ the comparison $C_b\preceq_1X$ with $E_1X\cap M=Q$ is accepted, so a bounded endorser has $\nu(c_b)\le\nu(Q)$; disjoint such $Q_n$ give $\nu(c_b)=0$. Choosing the $Q_n$ in $d_b\setminus\Theta_m$ with $\mu(\Theta_m)<\mu(d_b)/2$ gives the same for every $\Theta$, so $\Lambda_\Theta=M$ for every $\Theta$. (c) Given $\Theta$ and $\rho$, take $b\in(0,1)$, non-null $P\subseteq c_b$ with $0<\rho\mu(P)<\infty$, and non-null $Q\subseteq(d_b\setminus\Theta_m)\cap\{\rho\le n\}$ with $\rho\mu(Q)<\rho\mu(P)$. The comparison with $E_1X\cap M=(c_b\setminus P)\cup Q$ is accepted, relevant and $\Theta$-localized.
\end{proof}
\status{proved (\src{K3} Thm.~K3.3); confirmed (\src{K3\_referee} \S3), every axiom re-derived, one reader. \emph{Credit.} World~1 is the tilt world of \src{Q12\_referee} N3 (part~(e)); that report already showed that it accepts every cut pair with non-null target, that every endorsing charge lives at $1^-$, and that its sections are open. Theorem~\ref{thm:K33} completes it: the full axiom check, with the direct domination of \src{Q12\_referee}'s repair; $\Lambda_\Theta=M$ for every localizing $\Theta$; and the failure of (Loc) for every a.e.\ finite $\rho$. NUMERICS (\src{k3py/tilted\_nonS.py}, exact piecewise-constant data in 60-digit arithmetic): 300 of 300 cut pairs with targets of measure $10^{-4}$ to $10^{-2}$ are accepted for some $\kappa$; for 200 random positive step densities and random localizing sets of measure $0.05$, an accepted localized cut pair with $\rho\mu(P)>\rho\mu(Q)$ is found every time; the section of $c_{0.7}\setminus[0.4,0.5)$ contains $b=0.6999$ but not $b=0.7$. \emph{Consequence.} Every axiom except closedness of the world-1 sections holds, and world~1 is in (N). So any proof that mode (M2) implies (Loc) must use closed sections at the failing world. Lemma~DL of part~(e) is exactly what excludes the tilt world, and it is a two-world lemma for nested chains.}

\begin{lemma}\label{lem:K34}
Let $\mathfrak F\subseteq\mathfrak W$ be a family with $\mu(\delta^+\cap\delta'^+)>0$ for all $\delta,\delta'\in\mathfrak F$ and $\inf_{\delta\in\mathfrak F}\mu(\delta^+)=0$.
\begin{enumerate}[leftmargin=*,label=(\alph*)]
\item There is a localizing family $\Theta^{\mathfrak F}$ such that no member of $\mathfrak F$ is $\Theta^{\mathfrak F}$-localized.
\item If moreover some a.e.\ finite $\rho>0$ (respectively integrable $\rho>0$) and some localizing $\Theta$ are such that $\rho\mu$ endorses every $\Theta$-localized member of $\mathfrak W\setminus\mathfrak F$ in the cancellation form, then world $i$ is (Loc) (respectively (Loc$_{\mathrm{int}}$)), with the localizing family $\Theta\cup\Theta^{\mathfrak F}$.
\end{enumerate}
\end{lemma}
\begin{proof}[Proof sketch]
(a) Pick $\delta_k\in\mathfrak F$ with $\mu(\delta_k^+)<2^{-k}$ and put $\Theta^{\mathfrak F}_m\coloneqq\bigcup_{k\ge m}\delta_k^+$. Every $\delta\in\mathfrak F$ meets each $\delta_k^+$ non-nully, so it is not $\le0$ on any $\Theta^{\mathfrak F}_m$. (b) A $(\Theta\cup\Theta^{\mathfrak F})$-localized $\delta$ is $\Theta$-localized and, by (a), not in $\mathfrak F$.
\end{proof}
\status{proved (\src{K3} Lemma~K3.4); confirmed (\src{K3\_referee} \S4.1), one reader. The referee adds (proved in the referee report, one reader) that the lemma holds for countably many such families, by the diagonal union of Theorem~\ref{thm:K31}(d), and, more simply, for any family that is positive at every depth of one localizing family $T$ ($\delta^+\cap T_m$ non-null for all $m$), with no pairwise meeting and no small members. The second form is what the protected worlds of Theorems~\ref{thm:D21} and~\ref{thm:H24} instantiate, with $T$ the protected end; those worlds are (Loc$_{\mathrm{int}}$) by localizing at the end, as the referee checked for world~1 of Theorem~\ref{thm:H24}. \emph{What the lemma says about a world in (N)} (wording repaired as \src{K3\_referee} \S4.2 requires). For every family $\mathfrak F$ as in the lemma, every localizing $\Theta$ and every a.e.\ finite $\rho>0$, some $\Theta$-localized member of $\mathfrak W\setminus\mathfrak F$ is not endorsed by $\rho\mu$ in the cancellation form; this is the contrapositive of (b). The note drew two further readings: that any proof of the classification must use closedness against certificates whose positive parts do not pairwise meet, and that a world in (N) must accept comparisons with many pairwise disjoint small targets. Neither follows: the unendorsed objects are single comparisons, and positive parts that avoid every prescribed small set and do not lie in one intersecting family can still meet pairwise (long thin strips, any two of which cross). They are not used here.}

\emph{$T_1$: several worlds and $\sigma$-finite densities} (\src{K3} \S5). From here on (H$_U$) holds. For measurable $\rho_k\ge0$ on $U$, finite a.e.\ (with $\rho_k\coloneqq1$ on $N_+$, which affects no relevant comparison), put $c\coloneqq(\rho_k\sigma_k)_k$, $B_t\coloneqq\{t\cdot c>0\}$, $Z_t\coloneqq\{t\cdot c=0\}$, and $\Phi_t$, $\Phi^*_t$ as in part~(f$'$); property ($\gamma$) is that of \S\ref{sec:M0-closed}: $B_t$ is coherent (mod null) for a.e.\ $t\in\Di$. Lemma~\ref{lem:H10}(a), (b) hold for $\rho_k\ge0$: (a) uses only that $Z_t$ is null and $B_t$ incoherent, and (b) only the pointwise identity and finiteness (\src{K3\_referee} \S5.1).

\begin{theorem}\label{thm:K35}
Assume (H$_U$). Suppose that for each world $k$ there are a measurable $\rho_k:U\to[0,\infty)$, finite a.e., and a localizing family $\Theta^k$, and that there is a regular model $Y^\circ$, the \emph{patch}, such that:
\begin{itemize}[leftmargin=*]
\item[\textbf{(Y1)}] $\rho_k\mu$ endorses, in the cancellation form, every relevant world-$k$ comparison whose vector is $\le0$ a.e.\ on $\Theta^k_m\setminus E_k(Y^\circ)$ for some $m$;
\item[\textbf{(Y2)}] $\int_{E_k(Y^\circ)}\rho_k\,d\mu<\infty$ for every $k$;
\item[\textbf{(Y3)}] $c\ne0$ a.e.\ on $M$.
\end{itemize}
Then $c$ has property ($\gamma$). If moreover $\rho_k>0$ a.e.\ for every $k$, then (H$_U$) is contradicted.
\end{theorem}
\begin{proof}[Proof sketch]
Put $\Theta_m\coloneqq M\cap\bigcup_k(\Theta^k_m\cup\{\rho_k>m\})$, a localizing family off which every $\rho_k\le m$. By (Y3) and Tonelli, $Z_t$ is null for a.e.\ $t$; so if ($\gamma$) fails, some $t\in\Di$ has $Z_t$ null and $B_t$ incoherent. Let $X_m\coloneqq B_t$ on $U\setminus\Theta_m$ and $X_m\coloneqq Y^\circ$ on $\Theta_m$. It is regular, its costs are finite (at most $m\mu(M)$ off $\Theta_m$, plus (Y2)), and, by the bound of Lemma~\ref{lem:H10}(b), $\Phi_t(X_m)\le\sum_kt_k\int_{\Theta_m\cap E_k(Y^\circ)}\rho_k\,d\mu\to0$. Theorem~\ref{thm:kkm} gives a coherent selected $C_m$ with $C_m\preceq_kX_m$ at every $k$. Its world-$k$ vector is $\le0$ off $E_kX_m$, in particular on $\Theta^k_m\setminus E_k(Y^\circ)$, so (Y1) and property (R) give $\rho_k\mu(E_kC_m)\le\rho_k\mu(E_kX_m)<\infty$. Lemma~\ref{lem:H10}(b) gives $\Phi_t(C_m)\le\Phi_t(X_m)$, and Lemma~\ref{lem:H10}(a) gives $\Phi_t(C_m)\ge\Phi^*_t>0$: a contradiction for large $m$. If every $\rho_k>0$ a.e., Lemma~\ref{lem:H11} applies with $\Theta$ (its comparisons $T_1\preceq_kY$, $Y=T_1$ on $\Theta_m$, are covered by (Y1)), so $T_1$ is admissible at an interior weight; $c$ is strictly signed, and Theorem~\ref{thm:M0closed}(a), which needs no integrability, says that ($\gamma$) fails.
\end{proof}
\status{proved (\src{K3} Thm.~K3.5); confirmed (\src{K3\_referee} \S5.1), every step re-derived with $\rho_k\ge0$, one reader. \src{K3} calls (Y1)--(Y3) (P1)--(P3); they are renamed to avoid the note names. \emph{Reading} (\src{K3}; labelled a reading by the referee). Theorems~\ref{thm:H12}, \ref{thm:H13} and the near-Bayes models of \src{J1} are choices of a patch: $P_1$, $T_1$, and $T_1$ or a switched model deep in the ends.}

\begin{corollary}\label{cor:K35p}
Assume (H$_U$).
\begin{enumerate}[leftmargin=*,label=(\alph*)]
\item (Theorem~\ref{thm:H12}.) With $Y^\circ\coloneqq P_1$ and $\Theta^2=\Theta^3=\varnothing$, the theorem recovers Theorem~\ref{thm:H12}, with $\rho_1$ only a.e.\ finite; $\rho_2$ and $\rho_3$ need only be integrable on $R_{12}$ and $R_{13}$ respectively.
\item (Theorem~\ref{thm:H13}, $\sigma$-finite form.) With $Y^\circ\coloneqq T_1$: if every $\rho_k$ is a.e.\ finite and positive, endorses every $(\Theta^k,T_1)$-localized relevant comparison of world $k$ in the cancellation form, and is integrable on $E_k(T_1)$, then (H$_U$) is contradicted. Theorem~\ref{thm:H13} needed every $\rho_k$ integrable on $M$; non-integrability is now allowed wherever $T_1$ is correct at that world, in particular at every $T_1$-correct end.
\item (Every world localizable, compatible localizations.) Let every world $k$ be (Loc) with $(\rho_k,\Theta^k)$, and put $\rho_\pm\coloneqq\sum_{k:\sigma_k=\pm1}\rho_k$. Suppose that, for some $m_0$, (i) $\mu(\Theta^i_{m_0}\cap\Theta^j_{m_0}\cap R_{ij})=0$ for all $i\ne j$, and (ii) $\sum_k\int_{\Theta^k_{m_0}}\sum_{j:\sigma_j\ne\sigma_k}\rho_j\,d\mu+\int_{M\setminus\bigcup_k\Theta^k_{m_0}}\min(\rho_+,\rho_-)\,d\mu<\infty$. Then (H$_U$) is contradicted.
\end{enumerate}
\end{corollary}
\begin{proof}[Proof sketch]
(a), (b): Theorem~\ref{thm:K35} with the stated patches; $E_1(P_1)=\varnothing$, and ``$(\Theta^k,T_1)$-localized'' is ``$\le0$ on $\Theta^k_m\setminus E_k(T_1)$''. (c) Let $Y^\circ$ be correct at $k$ on $\Theta^k_{m_0}$, which is consistent by (i), since off $R_{ij}$ worlds $i$ and $j$ agree; off $\bigcup_k\Theta^k_{m_0}$ let it accept iff $\rho_-<\rho_+$, and let it accept $N_+$. For $m\ge m_0$, $E_k(Y^\circ)\cap\Theta^k_m$ is null, so (Y1) is the endorsement of (Loc). The sum in (ii) bounds $\sum_k\int_{E_k(Y^\circ)}\rho_k$ (the note says ``is''; it is equal only when the $\Theta^k_{m_0}$ are disjoint).
\end{proof}
\status{proved (\src{K3} Cor.~K3.5$'$); confirmed (\src{K3\_referee} \S5.2), one reader; the remark on $R_{12}$, $R_{13}$ in (a) is the referee's. \emph{The residual of (c).} If every world is (Loc), a cover needs, for every $m_0$, either \emph{sheltering}: $\Theta^i_{m_0}\cap\Theta^j_{m_0}\cap R_{ij}$ non-null for some $i\ne j$, so that two disagreeing worlds protect a common region at every depth (this is case (B) of part~(f$'$)); or \emph{infinite minority mass}: the sum in (ii) is infinite, i.e.\ two disagreeing coalitions both carry infinite mass on a common region off the localizing sets, or some world's density is not integrable on the localizing set of a world that disagrees with it. The referee notes that the second alternative is the literal negation of (ii); the genuine content is Theorem~\ref{thm:K35} itself and Lemma~\ref{lem:K36}. \emph{Sharpening} (proved in the referee report, \src{K3\_referee} \S5.3, one reader). If Question~Q$_{\mathrm{K3}}$ below has answer No, every world that is not in mode (M1) is (Loc$_{\mathrm{int}}$) (worlds that add up included, since a null $\Lambda_\varnothing$ gives a null $\Lambda_\infty$); then all densities in (c) can be taken integrable, (ii) holds automatically, and a cover of $T_1$ needs a world in mode (M1) or sheltering. Infinite minority mass cannot occur.}

\begin{lemma}\label{lem:K36}
Assume (H$_U$), and let every world $j$ be (Loc) with $(\rho_j,\Theta^j)$. Let $O$ be an open mixed orthant, $W\subseteq O\cap E_k(T_1)$ measurable, and $V(W)\coloneqq\{j:\sigma_j\ne\sigma_k\text{ on }O\}$. Then at least one of the following holds:
\begin{itemize}[leftmargin=*]
\item[] \emph{sheltered}: $\rho_k\mu\bigl(W\setminus\bigcup_{j\in V(W)}\Theta^j_m\bigr)<\infty$ for some $m$;
\item[] \emph{shadowed}: some $j\in V(W)$ has $\rho_j\mu(W\setminus\Theta^j_m)=\infty$ for every $m$; on $W$, $T_1$ is correct at $j$, and $j$ disagrees with $k$.
\end{itemize}
\end{lemma}
\begin{proof}[Proof sketch]
Suppose $W$ is not sheltered. Take a non-null $R\subseteq R^0_k\setminus\Theta^k_{m_1}$ on which every $\rho_j$ is bounded, and pairwise disjoint $p_n\subseteq W\setminus\bigcup_{j\in V(W)}\Theta^j_n$ with $\rho_k\mu(R)<\rho_k\mu(p_n)<\infty$ and every $\rho_j\mu(p_n)$ finite. Put $X_n\coloneqq T_1\sd(p_n\cup R)$. At world $k$ the vector $1_R-1_{p_n}$ is $\Theta^k$-localized and would be endorsed if accepted, which is false; at the worlds with the sign of $k$ on $O$, $X_n$ is strictly better; so competitiveness of $T_1$ (Theorem~\ref{thm:kkm}) gives $T_1\preceq_jX_n$ for some $j\in V(W)$. That vector, $1_{p_n}-1_R$, is positive only on $p_n$, which misses $\Theta^j_n$; so it is $\Theta^j$-localized, and endorsement gives $\rho_j\mu(R)\le\rho_j\mu(p_n)$. Some $j$ occurs for infinitely many $n$, and then $\rho_j\mu(W\setminus\Theta^j_m)=\infty$ for every $m$.
\end{proof}
\status{proved (\src{K3} Lemma~K3.6); confirmed (\src{K3\_referee} \S5.4), one reader. It is the $\sigma$-finite analogue of Lemma~\ref{lem:H21}(b), where a $T_1$-wrong end is shadowed by $T_1$-correct ends. So $T_1$-wrong blow-ups of $\sigma$-finite endorsers are always either sheltered or shadowed by a blow-up of a disagreeing, $T_1$-correct world on the same set, and Theorem~\ref{thm:K35} fails exactly when such blow-ups cannot be separated.}

\emph{$T_1$: one world without a finite certificate next to two integrable worlds} (\src{K3} \S6).

\begin{lemma}\label{lem:K37p}
Let $K\supseteq K'$ be a credal set and $i\ne j$. The model $D_K$ agrees with $P_i$ a.e.\ on the mixed problems of $U$ at which worlds $i$ and $j$ agree iff $K\cap[e_i,e_j]\ne\varnothing$ and $K\subseteq\co(K'\cup[e_i,e_j])$.
\end{lemma}
\begin{proof}[Proof sketch]
Take $i=1$, $j=2$; the problems are those of $U\cap[(+,+,-)\cup(-,-,+)]$. If $m^*\in K\cap[e_1,e_2]$, then $m^*\cdot g<0$ on $(-,-,+)$, and $D_K\supseteq D_{\co(K'\cup[e_1,e_2])}=U\cap\{g_1>0,g_2>0\}$. Conversely, $D_K\cap M$ is open; rejection of $(-1,-1,L)\in U$ for large $L$ gives points of $K$ converging to $[e_1,e_2]$; and a point of $K$ outside $\co(K'\cup[e_1,e_2])$ is separated by some $g\in U\cap(+,+,-)$ that $D_K$ rejects on a neighbourhood.
\end{proof}
\status{proved (\src{K3} Lemma~K3.7$'$); confirmed (\src{K3\_referee} \S6.1), one reader. NUMERICS (\src{k3py/between\_lemma.py}): 400 random polytopes; all 114 predicted to agree agree, and 285 of the 286 predicted to disagree disagree on the sample. The note's diagnosis of the remaining case inspected a different random stream; the referee's script \src{k3ref\_py/between\_diag.py} reproduces the main loop: the case has an extra point with third coordinate $9.19\cdot10^{-5}$, and $D_K$ accepts $(-1,-1,L)$ for $L=10^5,10^6,10^8$, as the lemma requires, but not for $L\le10^4$, beyond the sampled range $L\le e^8$. So it is a sampling limit, as the note said.}

\begin{theorem}\label{thm:K37}
Assume (H$_U$). Let worlds~$2$ and~$3$ have integrable $\rho_2,\rho_3>0$ a.e.\ that endorse all their relevant comparisons. Then world~1 is in mode (M1) or in (N). Suppose it is not in mode (M1). Let $\rho_1\coloneqq\rho_\infty$ and $\Lambda\coloneqq\Lambda_\infty$ (Theorem~\ref{thm:K31}(d) at world~1); more generally, let $\rho_1\ge0$ be any integrable density that endorses every $\Theta$-localized relevant world-1 comparison for some localizing $\Theta$, and $\Lambda\coloneqq\{\rho_1=0\}\cap M$. Put $c\coloneqq(\rho_1\sigma_1,\rho_2\sigma_2,\rho_3\sigma_3)$.
\begin{enumerate}[leftmargin=*,label=(\alph*)]
\item $\Lambda$ is non-null, and $c$ has property ($\gamma$).
\item \emph{Splices.} For $\theta\in[0,1]$ let $G_\theta\coloneqq\{\theta\rho_2\sigma_2+(1-\theta)\rho_3\sigma_3>0\}$, and let $S_\theta$ be the model equal to $P_1$ on $M\setminus\Lambda$ and to $G_\theta$ on $\Lambda$, accepting $N_+$. Then $S_\theta$ is coherent for a.e.\ $\theta\in(0,1)$, and $S_1$ ($P_2$ on $\Lambda$) and $S_0$ ($P_3$ on $\Lambda$) are coherent.
\item[(b$'$)] \emph{Rigidity.} There are coherent $C_{12}=D_{K_1}$ and $C_{13}=D_{K_0}$ in $\M_U$ with $K_1\cap[e_1,e_2]\ne\varnothing$, $K_1\subseteq\co(K'\cup[e_1,e_2])$, $K_0\cap[e_1,e_3]\ne\varnothing$ and $K_0\subseteq\co(K'\cup[e_1,e_3])$, such that a.e.\ $\Lambda=(E_1(C_{12})\cap R_{12})\cup(E_1(C_{13})\cap R_{13})$, and $E_1(C_{12})=E_1(C_{13})$ on $R_{12}\cap R_{13}$. In every open mixed orthant $O$, $\Lambda\cap O$ is a.e.\ $D_{K_*}\cap O$ if $g_1<0$ on $O$, and $O\setminus D_{K_*}$ if $g_1>0$ on $O$, for some $K_*\in\{K_1,K_0\}$: a convex cone, or the complement in $O$ of one.
\item[(c)] \emph{Admissibility.} If $\Lambda\not\supseteq E_1(T_1)\cap M$ (mod null), then $D_{T_1}=B_{p^*}\cup S$ with $S\subseteq Z_{p^*}$ for some $p^*\in\Di$.
\end{enumerate}
\end{theorem}
\begin{proof}[Proof sketch]
First sentence: mode (M0) at world~1 is excluded by Theorem~\ref{thm:M0closed}(b), since worlds~2 and~3 add up with the weights $\rho_k\mu$, and (Loc) by Theorem~\ref{thm:H12}. (a) If $\Lambda$ were null, Theorem~\ref{thm:H12} would apply; ($\gamma$) is Theorem~\ref{thm:K35} with $Y^\circ=P_1$, $\Theta^1=\Theta$, $\Theta^2=\Theta^3=\varnothing$ ((Y3) holds since $\rho_2>0$). (b) The map $(s,\theta)\mapsto(1-s,s\theta,s(1-\theta))$ is a $C^1$ bijection of $(0,1)^2$ onto $\Di$ with Jacobian of size $s$; by (a) and Fubini, for a.e.\ $\theta$ the models $B_{t(s,\theta)}$ are coherent for a.e.\ $s$, and the $\theta$-tie set is null. As $s\to0$, $B_{t(s,\theta)}\to S_\theta$ in measure (sign of $\sigma_1$ off $\Lambda$, sign of $\theta\rho_2\sigma_2+(1-\theta)\rho_3\sigma_3$ on $\Lambda$), and coherent models are closed under convergence in measure (proof of Lemma~\ref{lem:H10}(a)); $\theta\to1,0$ gives $S_1$, $S_0$. (b$'$) $S_1=P_1$ where $\sigma_1=\sigma_2$, so Lemma~\ref{lem:K37p} applies to its credal set $K_1$; on $R_{12}$, $S_1=P_1$ off $\Lambda$ and $S_1=U\setminus P_1$ on $\Lambda$, so $\Lambda\cap R_{12}=E_1(S_1)\cap R_{12}$; likewise $S_0$ on $R_{13}$; $M=R_{12}\cup R_{13}$. (c) Steps 1--3 and 5 of Lemma~\ref{lem:H11} use no positivity, and give $p^*\in\Delta$ with $p^*\cdot v\ge0$ a.e. Boundary points other than $e_1$ are excluded as in its step~4, using $\rho_2,\rho_3>0$; and $p^*=e_1$ would force $\rho_1=0$ on $E_1(T_1)\cap M$.
\end{proof}
\status{proved (\src{K3} Thm.~K3.7); confirmed (\src{K3\_referee} \S6.2), one reader. \src{K3} labels the four parts ($\alpha$), ($\beta$), ($\beta'$), ($\gamma$); they are relabelled (a), (b), (b$'$), (c) here, because ($\gamma$) is the name of the property in (a). ``Endorses exactly'' in the note is read as ``endorses every'' (referee). NUMERICS (\src{k3py/splice\_rigidity.py}, rerun identically by the referee): a blind set built from $K_1=\co(K'\cup\{e_1,\tfrac12(e_1+e_2)\})$ and $K_0=\co(K'\cup\{e_1,\tfrac12(e_1+e_3)\})$ covers $11\%$ of the mixed sample, $S_1$ and $S_0$ coincide with $D_{K_1}$ and $D_{K_0}$, and there are no midpoint-convexity violations in $20\,000$ sums; in that example $\Lambda\cap R_{12}\cap R_{13}$ is empty, so the consistency line is not tested there (referee). A slab blind set $\{-3<g_2/g_1<-1\}$ in $(+,-,+)$, not of the rigid form, gives 836 violations in $200\,000$ sums, so $S_1$ is incoherent, as (b$'$) predicts. \emph{Reading.} On each edge region of world~1, $\Lambda$ is the world-1 error set of a coherent model between $P_1$ and $P_j$: the shape of the blind set $c_a=E_1(C_a)$ of Theorem~\ref{thm:D21}, installed on both edges at once. But the blind set of Theorem~\ref{thm:D21} is localizable, and $\Lambda_\infty$ is not.}

\emph{Question Q$_{\mathrm A}$} (\src{K3} \S6; open, no conjecture stated). Is there an integrable field $c=(\rho_1\sigma_1,\rho_2\sigma_2,\rho_3\sigma_3)$ on $U$, with $\rho_2,\rho_3>0$ a.e.\ and $\rho_1\ge0$ vanishing exactly on a non-null set $\Lambda$ of the form of Theorem~\ref{thm:K37}(b$'$), that has property ($\gamma$) and satisfies $D_{T_1}=B_{p^*}\cup S$ with $S\subseteq Z_{p^*}$ for some $p^*\in\Di$? A negative answer is Theorem~\ref{thm:M0closed}(a) for fields with one weakly signed component; with Theorem~\ref{thm:K37} it would give: next to two integrable worlds, the third world of a cover is in mode (M1) or has $\Lambda_\infty\supseteq E_1(T_1)\cap M$. The proof of Theorem~\ref{thm:M0closed}(a) uses strict signs at every world, and \src{K3} did not check which steps survive a zero component.

\begin{proposition}\label{prop:KQA}
Let $c$, $\Lambda$, $K_1$, $K_0$ and $p^*\in\Di$ be as in Question~Q$_{\mathrm A}$, with $\Lambda$ described as in Theorem~\ref{thm:K37}(b$'$), and put $\theta^*\coloneqq p^*_2/(p^*_2+p^*_3)$.
\begin{enumerate}[leftmargin=*,label=(\alph*)]
\item On the doubly disagreeing region, $\Lambda$ lies in the world-1 error set of $T_1$: $\Lambda\cap R_{12}\cap R_{13}\subseteq E_1(T_1)$ a.e. Equivalently, $D_{T_1}\cap(+,-,-)\subseteq D_{K_1}\cap D_{K_0}$ and $D_{K_1}\cap(-,+,+)$, $D_{K_0}\cap(-,+,+)\subseteq D_{T_1}$: on $R_{12}\cap R_{13}$ the world-1 error sets of $C_{12}$ and $C_{13}$ lie inside that of $T_1$.
\item There is a coherent model $C^*$ equal to $P_1$ on $M\setminus\Lambda$ and to $T_1$ on $\Lambda\setminus Z_{p^*}$; so $E_1(C^*)=E_1(T_1)\cap\Lambda$ off $Z_{p^*}$.
\item On $\Lambda$, $t\cdot c=(t_2+t_3)\bigl(\theta\rho_2\sigma_2+(1-\theta)\rho_3\sigma_3\bigr)$ with $\theta=t_2/(t_2+t_3)$: on $\Lambda$ the Bayes sets depend on $t$ only through $\theta$, and every oriented atom contained in $\Lambda$ has its chart line through the vertex $e_1$.
\end{enumerate}
\end{proposition}
\begin{proof}[Proof sketch]
(a) On $\Lambda\cap(+,-,-)$, $c=(0,-\rho_2,-\rho_3)$, so $p^*\cdot c<0$ with no tie and $T_1$ rejects there; on $\Lambda\cap(-,+,+)$, $p^*\cdot c>0$ and $T_1$ accepts. In both cases $T_1$ errs at world~1. (b) On $\Lambda$, $B_{p^*}=G_{\theta^*}$; on the tie set $Z_{p^*}\cap\Lambda$, $\sigma_2\ne\sigma_3$ and the sign of $\theta\rho_2\sigma_2+(1-\theta)\rho_3\sigma_3$ is that of $(\theta-\theta^*)\sigma_2$. So the one-sided limits of the coherent splices $S_\theta$ of Theorem~\ref{thm:K37}(b) as $\theta\to\theta^*$ exist and are coherent. (c) is the computation, with $c_1=0$ on $\Lambda$.
\end{proof}
\status{proved in a referee report (\src{K3\_referee} \S7), one reader. \emph{Where Theorem~\ref{thm:M0closed}(a) breaks} (a reading in the same report). Theorem~\ref{thm:Estructure}(b) (the chart lines miss the vertices and cut off the odd vertex) and Lemma~\ref{lem:ball} (strict signs give the orthant) are exactly the steps that fail for chart lines through $e_1$. So Question~Q$_{\mathrm A}$ is not a routine extension of Theorem~\ref{thm:M0closed}(a); a proof of ``No'' would have to treat chords through $e_1$ as a new configuration beside chords, atoms and fans.}

\emph{Two worlds next to an additive world} (\src{K3} \S4.3).

\begin{proposition}\label{prop:K38}
In the setting of Theorem~\ref{thm:F12} ($n=2$, strict (G), (IB), (LA), (S), (C), world~2 additive with an integrable density $w>0$, nested selection with chain $c_b$ and uniform coordinate $x$), let $\Theta$ be localizing and $\nu\in\mathfrak N(\mathfrak W_\Theta)$ at world~1. Then $d\nu_c/d(w\mu)=\varphi\circ x$ a.e.\ on $M$ for some nondecreasing $\varphi:[0,1)\to[0,\infty)$. Hence every $\Lambda_\Theta$, and $\Lambda_\infty$, is an initial segment $c_a$ of the chain, mod null. Write $\Lambda_\infty=c_{a_\infty}$; world~1 is (Loc$_{\mathrm{int}}$) iff $a_\infty=0$.
\end{proposition}
\begin{proof}[Proof sketch]
By Theorem~\ref{thm:K31}(b), $\rho w\mu$, $\rho\coloneqq d\nu_c/d(w\mu)$, endorses every relevant comparison localized at $\Theta''\coloneqq\Theta\cup\Theta'$ for some $\Theta'$. If $c_b\cap\{\rho>c_2\}$ and $(d_b\setminus\Theta''_m)\cap\{\rho<c_1\}$ were both non-null with $c_1(1+\varepsilon)<c_2$, the cut pair $C_b\preceq_1X$, $E_1X\cap M=(c_b\setminus P)\cup Q$, with $P$, $Q$ in those sets and $w\mu(Q)=(1+\varepsilon)w\mu(P)$, would be accepted by (K), relevant and $\Theta''$-localized, and not endorsed. So $\operatorname{ess\,sup}_{c_b}\rho\le\operatorname{ess\,inf}_{d_b\setminus\Theta''_m}\rho$ for every $m$, and letting $m\to\infty$ the proof of Theorem~\ref{thm:F12}(b) applies with $\rho\ge0$. The supports are final segments of the chain, and so is their maximal union.
\end{proof}
\status{proved (\src{K3} Prop.~K3.8); confirmed (\src{K3\_referee} \S8), one reader. For world~1 of Theorem~\ref{thm:D21}, $a_\infty=0$; for the tilt world (Theorem~\ref{thm:K33}, which lacks closed world-1 sections), $a_\infty=1$.}

\begin{conjecture}\label{conj:K3C}
In the setting of Proposition~\ref{prop:K38}, with closed sections at both worlds, $a_\infty=0$: world~1 is (Loc$_{\mathrm{int}}$).
\end{conjecture}
\status{conjecture (\src{K3} K3.C), correctly labelled (\src{K3\_referee} \S8). \emph{Placement} (referee). It is the localized repair of \src{Q1}'s Conjecture Q1.C1, which Theorem~\ref{thm:D21} refutes as stated although that world has $a_\infty=0$, and the localized form of question (R2) of part~(g); \src{Q12\_referee}'s diagnosis applies, so a proof needs a rate in Lemma~T, i.e.\ (R3). \emph{Evidence, as corrected.} By Theorem~\ref{thm:K31}(f), $a_\infty>0$ means ratio certificates on every non-null $B'\subseteq c_{a_\infty}$ from comparisons localized at the end of the chain. \src{K3} stated that Theorem~\ref{thm:F14} excludes every such certificate that is chain-free (its positive parts disjoint from its negative parts). That does not follow (\src{K3\_referee} \S8): Theorem~\ref{thm:F14} needs a cover $\sum w_k1_{P_k}\ge1_S$ of a fixed set $S$, with a constant depending on $S$ and no rate, while the density $g$ of a ratio certificate is uncontrolled and may concentrate as $c\to0$. Theorem~\ref{thm:F14} excludes only chain-free ratio certificates whose $g$ is at least a fixed $\gamma$ on a set of fixed measure (for instance uniform $g$), which is weaker evidence, since uniform $g$ is not equivalent to Theorem~\ref{thm:K31}(f)(ii); excluding general chain-free ratio certificates needs the same missing rate (R3). Chained certificates, in which targets of some comparisons are partly sources of others, are not addressed.}

\emph{Question Q$_{\mathrm{K3}}$} (\src{K3} \S7; open, the governing question of the note). In a full system (every axiom, closed sections for every model), can some world have a non-null non-localizable blind set $\Lambda_\infty$? Equivalently (Theorem~\ref{thm:K31}(f), under (NC$_{\mathfrak W}$)): is there a non-null $B\subseteq M$ such that, for every localizing $\Theta$ and every $c>0$, every non-null $B'\subseteq B$ carries a $c$-ratio certificate made of $\Theta$-localized relevant comparisons? \emph{No} would give the classification in its strong form: every world that does not add up is in mode (M1) or (Loc$_{\mathrm{int}}$); and then, by the referee's sharpening of Corollary~\ref{cor:K35p}, a cover of $T_1$ needs a world in mode (M1) or sheltering. \emph{Yes}, with a world that is not (Loc), would refute the classification. By Theorem~\ref{thm:K33} no axiom other than closedness of the sections excludes such a world. Aimed at $T_1$ next to two integrable worlds, it would have to be blind, non-localizably, on a set of the rigid shape of Theorem~\ref{thm:K37}(b$'$), and Question~Q$_{\mathrm A}$ would need a positive answer. The two-world form next to an additive world is Conjecture~\ref{conj:K3C}. For $T_1$, settling the case without a finite certificate needs Question~Q$_{\mathrm{K3}}$ (or Conjecture~\ref{conj:K3C} on the edges), and, for a single exceptional world next to two integrable worlds, also Question~Q$_{\mathrm A}$ (\src{K3\_referee} \S9). \src{K3} built no three-world system and chose no selection through $T_1$.

\emph{Round 19, note \src{K2}: finite certificates (mode (M1)) in full systems} (\src{K2}; \src{K2\_referee}). A \emph{full system} is a ranking on $\M_U$ with strict (G), (IB), (LA), (S) (competitiveness on $\supp p$ and closed sections for every model) and (C). A \emph{certificate at world $i$} is a finite family of relevant comparisons $B_k\preceq_iX_k$ with vectors $\delta_k\coloneqq1_{E_iX_k}-1_{E_iB_k}$ and weights $\lambda_k>0$ such that $\sum_k\lambda_k\delta_k\le0$ a.e.\ and $\ne0$. Let $\mathcal C_{\pm}$ be the set of measurable $\{-1,0,1\}$-valued functions on $U$, and $\mathcal A_i\subseteq\mathcal C_\pm$ the acceptance set of world $i$ (Lemma~\ref{lem:vectorform}): it contains every $\delta\ge0$, no nonzero $\delta\le0$, and is closed under sums that stay in $\mathcal C_\pm$. If $\delta_1,\dots,\delta_m\in\mathcal A_i$ have pairwise disjoint positive parts and pairwise disjoint negative parts, then $\sum_k\delta_k\in\mathcal A_i$ (\src{N17} N17.1(b), the chaining lemma). Sign classes $Q_l$ and edge regions $R_{ij}=Q_i\cup Q_j$ are as in part~(f), and leading densities as in (L1) there. A toggle of the world-$i$ error status on a set flips the world-$j$ error status the same way on $\{\sigma_i=\sigma_j\}$ and the opposite way on $R_{ij}$; so for models $X,Y$ and $\delta\coloneqq1_{E_iX}-1_{E_iY}$,
\[
\rho_j(E_jX)-\rho_j(E_jY)=-\int\delta\,\tilde\rho_j\,d\mu,\qquad\tilde\rho_j\coloneqq\rho_j\text{ on }R_{ij},\quad\tilde\rho_j\coloneqq-\rho_j\text{ on }\{\sigma_i=\sigma_j\}.\tag{K2}
\]

\emph{Divisible vectors.} A vector $\delta\in\mathcal A_i$ is \emph{divisible} if there is a finite family $\mathcal M_\delta$ of finite measures on $U$, absolutely continuous with respect to $\mu$ and containing $\mu|_U$, such that every $\delta'\in\mathcal C_\pm$ with $\delta'^+\subseteq\delta^+$, $\delta'^-\subseteq\delta^-$ and, for one $f\in(0,1]$, $m(\delta'^\pm)=f\cdot m(\delta^\pm)$ for every $m\in\mathcal M_\delta$, lies in $\mathcal A_i$. In words: every proportional part of the move is accepted.

\begin{theorem}\label{thm:K21}
Let world $i$ have strict (G) and (IB) (any $n$, any target). If $\delta_1,\dots,\delta_r\in\mathcal A_i$ are divisible and $\lambda_1,\dots,\lambda_r\ge0$ satisfy $\sum_k\lambda_k\delta_k\le0$ a.e., then $\sum_k\lambda_k\delta_k=0$ a.e. In particular no certificate consists of divisible comparisons.
\end{theorem}
\begin{proof}[Proof sketch]
Suppose the sum is negative on a non-null set. (1) On the atoms $c$ of the finite algebra generated by the $\delta_k^\pm$, the cone $\{\lambda\ge0:D\lambda\le0\}$, $D_{ck}\coloneqq\delta_k(c)$, is rational polyhedral; by Minkowski--Weyl some integer $n\ge0$ has $Dn\le0$ and $(Dn)_{c_0}<0$. (2) Let $N_0\coloneqq\max_c\sum_{k:\delta_k(c)=-1}n_k\ge1$. By Lyapunov's theorem each cell splits into $N_0$ pieces of equal mass for every measure in $\bigcup_k\mathcal M_{\delta_k}$. (3) In each cell give each $k$ with $\delta_k(c)=-1$ its own $n_k$ pieces, and each $k$ with $\delta_k(c)=+1$ its own $n_k$ pieces among those already given to negatives (possible since $Dn\le0$); a positive piece of $k$ is never a negative piece of $k$. (4) The parts $\delta'_k$ ($+1$ on $k$'s positive pieces, $-1$ on its negative pieces) are proportional with $f=n_k/N_0$, hence accepted. (5) They have pairwise disjoint positive parts and pairwise disjoint negative parts, so their sum is accepted by the chaining lemma; it is $0$ or $-1$ on each piece, and $-1$ on some piece of $c_0$. This contradicts strict (G).
\end{proof}
\status{proved (\src{K2} Thm.~K2.1); confirmed (\src{K2\_referee} \S1), one reader. CHECKED (\src{k2py/divisibility.py}): the assignment of steps (3)--(5) on $3\,000$ random integer certificates, exact, 0 failures. It uses neither (S) nor (C).}

\emph{Forced comparisons.} A relevant comparison $B\preceq_iX$ is \emph{forced} if there is a selected $C_p$ with $B\sim_lC_p$ at every world $l$, $p_i>0$, and such that every $j\in\supp p\setminus\{i\}$ has a leading density $\rho_j$ on $M$ with $\rho_j(E_jX)<\rho_j(E_jB)$.

\begin{proposition}\label{prop:K22}
A forced comparison holds, $B\preceq_iX$, and its vector $\delta=1_{E_iX}-1_{E_iB}$ is divisible.
\end{proposition}
\begin{proof}[Proof sketch]
For $j\in\supp p\setminus\{i\}$, $B\sim_jC_p$ and Lemma~\ref{lem:B10} applied both ways give $\rho_j(E_jB)=\rho_j(E_jC_p)$, so $X\prec_jC_p$ by (L1). Competitiveness of $C_p$ on $\supp p$ leaves $C_p\preceq_iX$, and $B\sim_iC_p$ gives $B\preceq_iX$. For divisibility take $\mathcal M_\delta\coloneqq\{\mu|_U\}\cup\{\rho_j\mu|_{R_{ij}},\ \rho_j\mu|_{\{\sigma_i=\sigma_j\}\cap M}:j\in\supp p\setminus\{i\}\}$. For a proportional part $\delta'$ with fraction $f$, the regular model $X'$ with $E_iX'=E_iB\sd\supp\delta'$ has vector $\delta'$, and by (K2), $\rho_j(E_jX')-\rho_j(E_jB)=f\bigl(\rho_j(E_jX)-\rho_j(E_jB)\bigr)<0$; so $B\preceq_iX'$ is forced through the same $C_p$.
\end{proof}
\status{proved (\src{K2} Prop.~K2.2); confirmed (\src{K2\_referee} \S2), one reader, including the transport to tied members of $\C_L$.}

\begin{corollary}\label{cor:K22p}
In a full system every certificate at world $i$ contains a comparison that is not divisible, and in particular not forced. If every world other than $i$ has a leading density on $M$, every certificate at world $i$ contains a comparison $B\preceq_iX$ such that, for every selected $C_p$ tied with $B$ at every world with $p_i>0$, some $j\in\supp p\setminus\{i\}$ has $\rho_j(E_jX)\ge\rho_j(E_jB)$.
\end{corollary}
\status{proved (\src{K2} Cor.~K2.2$'$); confirmed (\src{K2\_referee} \S2), one reader. The second sentence is stated with $\rho_j$ as the referee requires: the note said ``$X$ is not strictly better than $B$ at some $j$'', which is stronger, since with a leading density $X\prec_jB$ and $\rho_j(E_jX)=\rho_j(E_jB)$ can coexist. \emph{Remarks} (\src{K2}). The generators $g_*$ of the designs of Theorem~\ref{thm:D21} are not divisible (CHECKED in \src{K2}; plausible but not re-derived by the referee), and neither are the exotic vectors of \src{N17} N17.7 (confirmed by the referee: at $f=\frac12$ a proportional part violates the defining inequalities). So rigidity is necessary for a certificate, not an artefact of those designs. One rigid comparison with a large weight next to divisible ones is not excluded: the proof needs proportional pieces that do not reuse a support.}

\begin{theorem}\label{thm:K23}
Let $n=2$, and let the ranking satisfy every axiom on $\M_U$. If either world has a certificate, then:
\begin{enumerate}[leftmargin=*,label=(\alph*)]
\item $\{E_1C:C\in\C_L\}$ is not a chain, so Conjecture~\ref{conj:Q2N} fails for this ranking;
\item neither world has a leading charge nonzero on $M$, and neither half of (MC$_{\mathrm{sel}}$) holds (Corollary~\ref{cor:MC}(b)); in particular neither world is additive with a density, (CP-a) with countably additive members, lexicographic with a charge first level, or monotonely continuous;
\item the certificate contains two bases whose error sets cross, every non-corner splice pair between them is strictly ordered the same way at both worlds (Lemma~\ref{lem:rigid}), and the set of double-tie classes of $\C_L$ is a compact linear continuum (Corollary~\ref{cor:continuum}(b)).
\end{enumerate}
\end{theorem}
\begin{proof}[Proof sketch]
On regular coherent models $E_2=M\setminus E_1$, so the world-2 error sets of $\C_L$ form a chain iff the world-1 error sets do; then \src{N17} N17.3, Lemma~3.1 (refereed in \src{N19}), excludes certificates at both worlds, and applied to the bases alone it needs two of them to cross. (b) is Corollary~\ref{cor:MC}(b), the contrapositive of Theorems~\ref{thm:S22} and~\ref{thm:MC}.
\end{proof}
\status{proved (\src{K2} Thm.~K2.3); confirmed (\src{K2\_referee} \S3), one reader. It assembles N17.8(c), Lemma~3.1 of N17.3 and Corollary~\ref{cor:MC}(b). In (c) the shape at a crossing is the necessary shape in the status of Theorem~\ref{thm:MC} (the relation of $b$ with the initial bands of $A$ switches with no tie). \emph{Reading.} A two-world system with a certificate is a counterexample to Conjecture~\ref{conj:Q2N} of the hardest kind, with no leading charge and no (MC$_{\mathrm{sel}}$) at either world; so the two-world route to mode (M1) is at least as hard as refuting that conjecture, which has resisted rounds 13 to 18. Without (C) a crossing selection exists (Proposition~\ref{prop:R71}), but both its worlds are additive, so it has no certificate.}

\begin{theorem}\label{thm:K24}
Let $n=3$, a full system, $i$ one world, and $j,k$ the other two, additive with integral penalties $\psi_l(X)=\int_{E_lX}\rho_l\,d\mu$, $\rho_l>0$ a.e.\ integrable ($l=j,k$). For $s\in(0,1)$ let $F_s\coloneqq\{g\in U:s\rho_j\sigma_j+(1-s)\rho_k\sigma_k>0\}$; let $\mathcal T_{jk}$ be the countable set of $s$ for which $\{s\rho_j=(1-s)\rho_k\}\cap R_{jk}$ is non-null; and put $r\coloneqq\rho_k/\rho_j$, $\kappa(s)\coloneqq s/(1-s)$. For $s\notin\mathcal T_{jk}$:
\begin{enumerate}[leftmargin=*,label=(\alph*)]
\item $F_s$ is a.e.\ equal to a selected model, and $F_s\in\C_L$;
\item no regular $X\ne F_s$ has $\psi_j(X)\le\psi_j(F_s)$ and $\psi_k(X)\le\psi_k(F_s)$; so $F_s$ is undominated whatever world $i$'s order is, and no comparison $F_s\preceq_iX$ is forced through a selected index $p$ with $\{j,k\}\subseteq\supp p$;
\item $E_i(F_s)=Q_i\cup(Q_j\cap\{r<\kappa(s)\})\cup(Q_k\cap\{r>\kappa(s)\})$ a.e.;
\item for $s<s'$ outside $\mathcal T_{jk}$, $F_s$ and $F_{s'}$ cross iff $\mu(Q_j\cap\{\kappa(s)\le r<\kappa(s')\})>0$ and $\mu(Q_k\cap\{\kappa(s)<r\le\kappa(s')\})>0$;
\item if a certificate at world $i$ has only bases $F_{s_1},\dots,F_{s_m}$ ($s_l\notin\mathcal T_{jk}$), every comparison based at the least $s$ has its gain in $Q_j$, and every comparison based at the largest $s$ has its gain in $Q_k$; these gains are disjoint.
\end{enumerate}
\end{theorem}
\begin{proof}[Proof sketch]
(a) For $s\notin\mathcal T_{jk}$, $F_s$ is the a.e.\ unique minimiser of $s\psi_j+(1-s)\psi_k$; Theorem~\ref{thm:kkm} gives a selected undominated $C$ with $C\preceq_lF_s$ at every $l$, and additivity at $j$ and $k$ makes $C$ a minimiser, so $C=F_s$ a.e. (b) Such an $X$ would be a second minimiser. (c) Class by class: on $Q_i$ worlds $j$, $k$ agree against $i$; on $Q_j$, $F_s$ sides with $j$ iff $r<\kappa(s)$, which is where it errs at $i$; symmetrically on $Q_k$; it is correct on $N_+$. (d) follows from (c). (e) By Lemma~3.1 of \src{N17} N17.3, the gain of a comparison at a base lies in the error set of another base; for the least $s$ the other bases' error sets meet $Q_k\cap\{r\le\kappa(s)\}$ only in null sets.
\end{proof}
\status{proved (\src{K2} Thm.~K2.4); confirmed (\src{K2\_referee} \S4), one reader. NUMERICS (\src{k2py/frontier\_and\_edges.py}~(A), Brier densities $\rho_l=|g_l|$ for $T_1$ with $\varepsilon=0.3$, $s=0.9$, $60\,000$ samples): formula (c) has 0 mismatches at three weights, and $F_{0.4}$, $F_{0.6}$ cross; the referee's check found 0 of $243\,722$ candidate gain points outside the class predicted by (e). \emph{Repairs of the readings} (\src{K2\_referee} \S4). (1) \emph{Crossing needs interleaving.} The frontier models cross only when the ratio $r$ interleaves on $Q_j$ and $Q_k$, as it does for the Brier densities; if $\operatorname{ess\,sup}_{Q_j}r\le\operatorname{ess\,inf}_{Q_k}r$ the frontier is a chain, and nothing excludes this when the $F_s$ are coherent. So the crossing bases that \src{N17} N17.3 requires of a certificate are available in this class when $r$ interleaves, and not otherwise. (2) The note's ``an (M1) world has no positive endorser at all, so it sits in case (A) of \src{H1}'' does not follow: a certificate excludes bounded positive weights and densities finite on its supports, but not an a.e.\ finite density of infinite mass there. The intended conclusion holds by Theorem~\ref{thm:H12} directly: next to two worlds with integrable full endorsers the third world of a cover is not (Loc), certificate or not. (3) In (e), ``the extreme pair is executable as a pair'' should be read through the scheme below: disjoint gains give $\int(g_k+g_l)\varphi\le\int_U\varphi$, so condition (pair) of Theorem~\ref{thm:K25} fails for that pair.}

\emph{A sufficient construction scheme} (\src{K2} \S5). Let $c=(\rho_1\sigma_1,\rho_2\sigma_2,\rho_3\sigma_3)$ be a strictly signed integrable field on $U=D_{K'}$ whose Bayes sets $B_q\coloneqq\{g\in U:q\cdot c(g)>0\}$ are coherent with null ties for \emph{every} $q\in\Delta$; the Brier field $c(g)=g$ is an example, with $B_q=D_{\co(K'\cup\{q\})}$. Worlds~2 and~3 are additive with penalties $\psi_j(X)=\int_{E_jX}\rho_j\,d\mu$, and the selection is $q\mapsto B_q$. With $\tilde\rho_j$ as in (K2) relative to world~1, put $\varphi_a\coloneqq a_0\rho_1+a_1\tilde\rho_2+a_2\tilde\rho_3$ for $a\in\R^3_{\ge0}$. A \emph{forced vector} is $t=1_{E_1X}-1_{E_1B_q}$ with $q_1>0$, $X$ any model, and $\psi_j(X)<\psi_j(B_q)$ for every $j\in\supp q\setminus\{1\}$. Let $\mathcal Y$ be a set of functions $\varphi_a$ with $\rho_1\in\mathcal Y$, $a_0>0$ for every member, every member $\ge0$ on $U$, and $\int t\varphi\,d\mu\ge0$ for every member $\varphi$ and every forced vector $t$. Put $\mathcal K_{\mathcal Y}\coloneqq\{f\in L^\infty(U):\int f\varphi\,d\mu\ge0\text{ for all }\varphi\in\mathcal Y\}$.

\begin{theorem}\label{thm:K25}
In this setting let $g_1,\dots,g_m\in\mathcal C_\pm$ satisfy:
\begin{itemize}[leftmargin=*]
\item[\textbf{(fit)}] $g_k=1_{E_1X_k}-1_{E_1B_{q_k}}$ for some $q_k\in\Delta$ and regular $X_k$;
\item[\textbf{(endorse)}] for every $k$ some $\varphi\in\mathcal Y$ has $\int g_k\varphi\,d\mu>0$;
\item[\textbf{(pair)}] for every $k,l$ ($k=l$ allowed) some $\varphi\in\mathcal Y$ has $\int(g_k+g_l)\varphi\,d\mu>\int_U\varphi\,d\mu$;
\item[\textbf{(cert)}] some $\lambda\ge0$ has $\sum_k\lambda_kg_k\le0$ and $\ne0$.
\end{itemize}
Let world~1 accept $X\preceq_1Y$ iff $1_{E_1Y}-1_{E_1X}\in\mathcal A_1\coloneqq\mathcal C_\pm\cap\bigl(\mathcal K_{\mathcal Y}\cup\bigcup_k(g_k+\mathcal K_{\mathcal Y})\bigr)$. Then this world, with worlds~2 and~3 and the selection $q\mapsto B_q$, satisfies strict (G), (IB), (LA), (S) with closed sections for every model, and (C) with $\C_L=\{B_q:q\in\Delta\}$; and world~1 has a certificate.
\end{theorem}
\begin{proof}[Proof sketch]
(1) $\mathcal A_1$ contains every $\delta\ge0$; it contains no nonzero $\delta\le0$ ($\rho_1>0$, and the (endorse)-witness excludes $g_k+\mathcal K_{\mathcal Y}$); it is closed under sums in $\mathcal C_\pm$, since a sum from $g_k+\mathcal K_{\mathcal Y}$ and $g_l+\mathcal K_{\mathcal Y}$ leaves $\mathcal C_\pm$ by (pair). Lemma~\ref{lem:vectorform} gives strict (G) and (IB). (2) Competitiveness: at $q_1=0$ by $(2,3)$-Bayes optimality; at $q=e_1$ trivially; at $q_1>0$ either some $j$ has $\psi_j(B_q)\le\psi_j(X)$, or the vector is forced and lies in $\mathcal K_{\mathcal Y}$. (3) If $Y$ weakly dominates $B_q$, Bayes optimality gives $\int w\rho_1\le0$, $w\coloneqq1_{E_1B_q}-1_{E_1Y}$, with equality only for $Y=B_q$, and (K2) gives $\int w\tilde\rho_j\le0$; so $\int w\varphi<0$ for every $\varphi\in\mathcal Y$ when $Y\ne B_q$, and $w\notin\mathcal A_1$. Theorem~\ref{thm:kkm} then gives $\C_L=\{B_q\}$, all coherent, hence (LA) and (C). (4) $q\mapsto B_q$ is continuous in measure, so each world-1 section is a finite union of closed sets; worlds~2 and~3 have closed sections by continuity. (5) $g_k\in\mathcal A_1$ since $0\in\mathcal K_{\mathcal Y}$, and (cert) makes the comparisons $B_{q_k}\preceq_1X_k$ a certificate.
\end{proof}
\status{proved (\src{K2} Thm.~K2.5); confirmed (\src{K2\_referee} \S5), every axiom checked, closed sections and (C) included, one reader. \src{K2} names the conditions (fit), (E), (P), (Cert); they are renamed because (E) and (P) are used elsewhere. \emph{An addition of the referee} (proved in the referee report, one reader): requiring forced vectors to lie in $\mathcal K_{\mathcal Y}$ is no restriction for small trades. If $\varphi_k$ is the (endorse)-witness of $g_k$, with $\gamma_k\coloneqq\int g_k\varphi_k>0$, a forced vector $t$ with $\int_{t^+}\varphi_k<\gamma_k$ for all $k$ is not in any $g_k+\mathcal K_{\mathcal Y}$, so if it is accepted it lies in $\mathcal K_{\mathcal Y}$. \emph{Interior indices} (\src{K2}, confirmed). For a forced vector at an interior $q$, Bayes optimality gives $\int t\varphi_a>0$ for every $a\ge0$; so the constraints on $\mathcal Y$ come from the two world-1 edges.}

\emph{Edge conditions.} Put $r_2\coloneqq\rho_2/\rho_1$. \textbf{(SPREAD$_{12}$)}: there are $\theta>0$, $a\ge0$ and $\eta>0$ such that, for every $\varepsilon\in(0,\theta)$, both $Q_1\cap\{\theta<r_2<\theta+\varepsilon,\ \rho_3/\rho_2>a+\eta\}$ and $Q_1\cap\{\theta-\varepsilon<r_2\le\theta,\ \rho_3/\rho_2<a\}$ are non-null. \textbf{(ADJ$_{12}$)}: there are $\theta>0$ and $a>0$ such that, for every $\varepsilon\in(0,\theta)$, both $Q_1\cap\{\theta<r_2<\theta+\varepsilon,\ \rho_3/\rho_2>a\}$ and $Q_2\cap\{\theta-\varepsilon<r_2\le\theta,\ \rho_3/\rho_2>a\}$ are non-null. (SPREAD$_{13}$) and (ADJ$_{13}$) are the same with $2$ and $3$ exchanged.

\begin{theorem}\label{thm:K26}
In the setting of Theorem~\ref{thm:K25}, suppose (SPREAD$_{12}$) or (ADJ$_{12}$), and (SPREAD$_{13}$) or (ADJ$_{13}$). Then every $a\ge0$ such that $\varphi_a$ endorses all forced vectors has $a_1=a_2=0$. So $\mathcal Y\subseteq\{a_0\rho_1\}$, no $g_1,\dots,g_m$ satisfy both (endorse) and (cert), and the scheme of Theorem~\ref{thm:K25} is empty over such bases.
\end{theorem}
\begin{proof}[Proof sketch]
At the edge index $q=(q_1,q_2,0)$ with $q_1/q_2=\theta$, $E_1(B_q)\cap R_{12}=\{r_2>\theta\}$. For non-null $B\subseteq R_{12}\cap\{r_2>\theta\}$ and $B'\subseteq R_{12}\cap\{r_2\le\theta\}$ with $\rho_2\mu(B')=(1+\eta')\rho_2\mu(B)$, moving the world-1 errors from $B$ to $B'$ lowers $\psi_2$, so $t=1_{B'}-1_B$ is forced and endorsement needs $\int_{B'}\varphi_a\ge\int_B\varphi_a$. On $Q_1$, $\varphi_a=\rho_2(a_0/r_2+a_1+a_2\rho_3/\rho_2)$; on $Q_2$, $\varphi_a=\rho_2(a_0/r_2+a_1-a_2\rho_3/\rho_2)$. Under (SPREAD$_{12}$), with $B$, $B'$ in the two sets at scale $\varepsilon$ and $\varepsilon,\eta'\downarrow0$, this gives $a_2\eta\le0$; under (ADJ$_{12}$) it gives $2a_2a\le0$. So $a_2=0$, and the edge $(1,3)$ gives $a_1=0$. Then (endorse) gives $\int g_k\rho_1>0$ for all $k$, while (cert) gives $\sum_k\lambda_k\int g_k\rho_1<0$.
\end{proof}
\status{proved (\src{K2} Thm.~K2.6); confirmed (\src{K2\_referee} \S6.1), one reader. NUMERICS (\src{k2py/frontier\_and\_edges.py}~(B), $\varepsilon=0.3$, $s=0.9$, $\theta=1$): at the edge $(1,2)$, $a=(1,0,0.05)$ and $(1,0,0.2)$ are violated; the referee's check confirms that $a=(1,0.3,0)$ and $(1,0.05,0)$ are violated at the edge $(1,3)$, which the note's script did not test.}

\begin{proposition}\label{prop:K26a}
The Brier base ($\rho_l=|g_l|$) satisfies (SPREAD$_{12}$) and (SPREAD$_{13}$) for every $K'=u+s(T_1-u)$.
\end{proposition}
\begin{proof}[Proof sketch]
On the open set $Q_1\cap U\cap\{g_1>0\}$ the map $g\mapsto(|g_2/g_1|,|g_3/g_2|)$ is smooth with surjective differential, so preimages of open subsets of its image are non-null; choose a rectangle around $(\theta,a+\eta/2)$ inside the image.
\end{proof}
\status{proved (\src{K2} Prop.~K2.6a); confirmed (\src{K2\_referee} \S6.1), one reader. Only (SPREAD) is proved for the Brier base; (ADJ$_{12}$) holds there numerically (\src{k2ref\_py/k2ref\_checks.py}~(2): at $\theta=1$ both sets are non-empty for $a$ up to $10$). The note's ``both hold for the Brier base, PROVED'' is corrected accordingly.}

\begin{proposition}\label{prop:KR62}
Let the base be the Brier base. Let $\varphi\in L^1(U)$, $\varphi\ge0$, satisfy $\int t\varphi\,d\mu\ge0$ for every two-set forced trade $t$ at the world-1 edges $(1,2)$ and $(1,3)$; the trades may be restricted to arbitrarily small mass. Then $\varphi=c\rho_1$ a.e.\ on $M$ for some $c\ge0$.
\end{proposition}
\begin{proof}[Proof sketch]
(1) For each $\theta>0$ the edge trades give
\[
\operatorname*{ess\,sup}_{R_{12}\cap\{r_2>\theta\}}\varphi/\rho_2\ \le\ \operatorname*{ess\,inf}_{R_{12}\cap\{r_2\le\theta\}}\varphi/\rho_2 .
\]
(2) So $\varphi/\rho_2=H(r_2)$ a.e.\ on $R_{12}$ for a nonincreasing $H$ (the level sets of $r_2$ are null), and likewise $\varphi/\rho_3=H_3(r_3)$ on $R_{13}$. (3) On the link class $Q_1\cap\{g_1<0\}$, with $x=r_2$, $y=r_3$, the two descriptions give $xH(x)=yH_3(y)$ a.e.\ for the law of $(x,y)$, which is equivalent to Lebesgue measure on an open up-closed region whose sections are overlapping rays; Fubini gives $xH(x)\equiv c$. (4) So $\varphi=c\rho_2/r_2=c\rho_1$ on $R_{12}$, likewise on $R_{13}$, and $M=R_{12}\cup R_{13}$.
\end{proof}
\status{proved in a referee report (\src{K2\_referee} \S6.2, Prop.~R6.2), one reader; the region check (3) found 0 violations of up- and down-closure in $2\cdot10^5$ samples. \emph{Consequences} (the same report). (i) In Theorem~\ref{thm:K25}, $\mathcal Y$ may be any set of nonnegative integrable functions containing $\rho_1$; equivalently $\mathcal K_{\mathcal Y}$ may be any weak-$*$ closed convex cone. Over the Brier base every member of $\mathcal Y$ is then $c\rho_1$ on $M$, and since the $g_k$ vanish on $N_+$, (endorse) and (cert) clash as in Theorem~\ref{thm:K26}; by the referee's addition to Theorem~\ref{thm:K25}, this holds even if forced vectors are only required to be accepted. So over the Brier base the whole template of dual-cone decorations is empty, not only its linear form, with no (SPREAD) or (ADJ) hypothesis. (ii) This supplies the endorser half of the evidence for Conjecture~\ref{conj:K26p} in the Brier case: the weak-$*$ closed convex cone generated by the edge-forced trades and the nonnegative functions contains every $\rho_1$-strict improvement in its closure. (iii) For general bases the same proof works under null level sets of $r_2$ on $R_{12}$ and of $r_3$ on $R_{13}$ and a linked-sections condition on the law of $(r_2,r_3)$ on $Q_1$.}

\begin{conjecture}\label{conj:K26p}
Let worlds~2 and~3 be additive with integrable densities, in a full system. Suppose that on the indices with $q_1>0$ the selected world-1 error sets are those of the Bayes family of an additive field with (SPREAD) or (ADJ) at both world-1 edges. Then world~1 has a leading density, so it adds up, whatever its order; in particular it has no certificate.
\end{conjecture}
\status{conjecture (\src{K2} K2.6$'$). \emph{Evidence.} The endorser side holds for the Brier base by Proposition~\ref{prop:KR62}. What is missing is the passage from the weak-$*$ closed convex hull of the forced cone to the acceptance set itself, which is neither convex nor closed (a transfer chain in the manner of Lemma~\ref{lem:B12}). \src{K2\_referee} recommends this step for the Brier--Bayes selection as the nearby target: if it holds, a Bayes-shaped selection over Brier-additive worlds~2 and~3 forces world~1 to add up; if it fails, the failure must sit where pairwise sums leave $\mathcal C_\pm$, which is where a non-decorative world with a certificate would have to live.}

\begin{lemma}\label{lem:K27}
Let world $i$ have strict (G), let (S) hold (any $n$, any target), let $e_ie_j$ be an edge, and let $p_m\to e_j$ along the open edge. There is no non-null $A\subseteq R_{ij}$ with $A\cap E_i(C_{p_m})=\varnothing$ for all $m$.
\end{lemma}
\begin{proof}[Proof sketch]
Such $C_{p_m}$ err at world $j$ on all of $A$. Take non-null $A'\subsetneq A$ with $A\setminus A'$ non-null and $X\coloneqq P_j\sd A'$; then $X\prec_jC_{p_m}$, so competitiveness gives $C_{p_m}\preceq_iX$, and closedness gives $C_{e_j}=P_j\preceq_iX$. But $E_iX=R_{ij}\setminus A'\subsetneq E_iP_j$, so $X\prec_iP_j$.
\end{proof}
\status{proved (\src{K2} Lemma~K2.7); confirmed (\src{K2\_referee} \S7), one reader. Its use: since $M=R_{ij}\cup R_{ik}$, every non-null subset of $M$ meets world-$i$ error sets of edge-selected models with $p_i>0$; next to a world with a leading density those models carry forced comparisons. The rest of the note's ``use'' paragraph is heuristic.}

\emph{What \src{K2} gives for $T_1$, with the overstatements removed} (\src{K2} \S8, \src{K2\_referee} \S8). (1) If a world of a cover has a certificate, the certificate contains a rigid, non-divisible comparison, which next to worlds with leading densities is not forced (Corollary~\ref{cor:K22p}). The dominator comparisons used in parts~(f) and~(f$'$) are not forced at the dominator's index, so \src{K2} does not exclude a certificate made of them. (2) Next to two worlds additive with integrable densities, crossing bases are available through the frontier family when the density ratio interleaves (Theorem~\ref{thm:K24}). (3) Over additive bases whose Bayes models are coherent at every weight and which satisfy (SPREAD) or (ADJ) at both world-1 edges, and over the Brier base for every weak-$*$ closed acceptance cone, the dual-cone decoration scheme is empty (Theorem~\ref{thm:K26}, Proposition~\ref{prop:KR62}). As a route to $T_1$ that scheme is excluded in any case: its $\C_L$ is the Bayes family of an additive field, which never contains $T_1$ (Theorem~\ref{thm:main}). The note's further statements (``protection needs endorsers that concentrate on the overlaps \dots\ what is left is $\rho_i$'' as a general obstruction, and ``a $T_1$ cover with an (M1) world has a non-Bayes selection on an open set of indices, with a forced cone that is not a half-plane family'') hold only for decorations of the type of Theorem~\ref{thm:K25}; in general they would need Conjecture~\ref{conj:K26p}, and ``on an open set of indices'' is not derived. They are not stated here. \emph{Question Q-K2} (\src{K2} \S9; open). Is there a full system on $\M_U$ with two worlds additive with integrable densities whose third world has a certificate? A negative answer would close mode (M1) next to two such worlds. No full system with a world in mode (M1) is known, for any target.

\emph{Round 19, note \src{K1}: countably many ends} (\src{K1}; \src{K1\_referee}). \textbf{Class (EP$_\infty$)}: (H$_U$) holds; every world $k$ is end-protected with a \emph{countable}, possibly empty, family $\mathcal S_k$, refined as in Lemma~\ref{lem:H20}; and there are $\rho_k\in L^1(\mu)$, $\rho_k>0$ a.e., such that $\rho_k\mu$ endorses every relevant base comparison of world $k$. Class (EP) of part~(f$'$) is the case of finite families. A world that adds up with a countably additive weight is in the class with $\mathcal S_k=\varnothing$ (Radon--Nikodym). A world that adds up only through a purely finitely additive weight is not in the class with $\mathcal S_k=\varnothing$; whether it is in the class through a nonempty family of ends is not decided (the referee's correction of the note's ``is not in the class''). Each refined end $S$ has a \emph{purity depth} $n_S$ beyond which it lies in one orthant and is $T_1$-correct or $T_1$-wrong. Put $\psi\coloneqq(\rho_k\mu)_k$. With summable depths, Lemma~\ref{lem:H11} makes $\psi$ admissible, and $I_\psi$ is not null (Theorem~\ref{thm:M0closed}(a)); \src{H12\_referee} \S1.3 records that Lemma~\ref{lem:H11} needs only countably many ends. Fix $t\in I_\psi$, write $B_t\coloneqq B^\psi_t$, and put $\mathrm{Cost}_t(Y)\coloneqq\sum_kt_k\rho_k\mu(E_kY)$. By Lemma~\ref{lem:H10}, $\mathrm{Cost}_t(Y)-\mathrm{Cost}_t(B_t)=\Phi_t(Y)$ for every regular $Y$ with finite costs, so every coherent $C$ has $\mathrm{Cost}_t(C)\ge\underline{\mathrm{Cost}}_t\coloneqq\mathrm{Cost}_t(B_t)+\Phi^*_t$, with equality attained. A model $A$ is \emph{$k$-safe} if it is regular and there are depths $(a_S)_{S\in\mathcal S_k}$ such that $A$ is correct at world $k$ a.e.\ on $S_{a_S}$ for every $S\in\mathcal S_k$; if $\mathcal S_k=\varnothing$, every regular model is $k$-safe. A \emph{depth assignment} on a family of ends is a map $S\mapsto a_S\in\mathbb N$; for an assignment $a$ on all ends of all worlds put $\mathcal E^k(a)\coloneqq\bigcup_{S\in\mathcal S_k}S_{a_S}$ and $\mathcal E(a)\coloneqq\bigcup_k\mathcal E^k(a)$.

\begin{lemma}\label{lem:K10}
Assume (EP$_\infty$). Let $C\preceq_kX$ be a relevant comparison with vector $\delta\coloneqq1_{E_kX}-1_{E_kC}$ such that $\delta\le0$ a.e.\ on $S_{a_S}$ for every $S\in\mathcal S_k$, for some depths $a_S$. Then the comparison is a base comparison, and $C=X$ a.e.\ on $S_{m'_S}$ for some depths $m'_S\ge a_S$. In particular every relevant comparison $C\preceq_kA$ into a $k$-safe model $A$ is a base comparison.
\end{lemma}
\begin{proof}[Proof sketch]
$\delta$ is $S$-positive for no $S$, so end-protection makes it $\ge0$ on some $S_{b_S}$; hence $\delta=0$ a.e.\ on $S_{\max(a_S,b_S)}$, end by end. On $U$, $\delta=0$ iff $C$ and $X$ give the same advice.
\end{proof}
\status{proved (\src{K1} Lemmas~K1.0, K1.0$'$); confirmed (\src{K1\_referee} \S2.1), one reader. It is Lemma~\ref{lem:J10} without the maximum over the family: ``base'' is a statement end by end, so no uniform depth is needed.}

\begin{lemma}\label{lem:K11}
Assume (EP$_\infty$). Let $W\in\mathcal S_i$ be $T_1$-wrong, lying in the orthant $O$, and let $\mathcal T$ be the set of $T_1$-correct ends of the worlds of $V(W)$.
\begin{enumerate}[leftmargin=*,label=(\alph*)]
\item For every depth assignment $M$ on $\mathcal T$ there is $m'$ with $W_{m'}\subseteq\bigcup_{S\in\mathcal T}S_{M_S}$ modulo null sets.
\item If $P\subseteq W$ meets $W_m$ non-nully for every $m$, some $S\in\mathcal T$ has $\mu(P\cap W_m\cap S_n)>0$ for all $m,n$.
\end{enumerate}
\end{lemma}
\begin{proof}[Proof sketch]
(a) Enlarging $M$ only shrinks the right side, so enlarge it beyond the purity depths and extend it to all ends $E$ of all worlds by depths $N_E$ with $\sum_E\mu(E_{N_E})<\mu(R^0_i)/2$. $T_1$-wrong ends of worlds of $V(W)$ are disjoint from $W$ at large depth, since $T_1$ is right at those worlds on $W$. The set $R\coloneqq R^0_i\setminus\bigcup_EE_{N_E}$ is non-null. If the claim failed, choose non-null $p_{m'}\subseteq W_{m'}\setminus\bigcup_{S\in\mathcal T}S_{M_S}$ with $\sum_{m'}\rho_j\mu(p_{m'})<\rho_j\mu(R)$ for every $j\in V(W)$, and $p\coloneqq\bigcup p_{m'}$, and argue as in Lemma~\ref{lem:H21}: with $X\coloneqq T_1\sd(p\cup R)$, competitiveness of $T_1$ gives $T_1\preceq_jX$ for some $j\in V(W)$; its vector $1_p-1_R$ vanishes at depth $M_E$ on every end of $j$, so it is a base comparison, and endorsement gives $\rho_j\mu(R)\le\rho_j\mu(p)$, against the choice of $p$. (b) Enumerate $\mathcal T=\{S^1,S^2,\dots\}$. If each $S^k$ had $(m_k,n_k)$ with $P\cap W_{m_k}\cap S^k_{n_k}$ null, choose $M_k\ge n_k$ with $\mu(S^k_{M_k})\le2^{-k-1}\mu(P\cap W_{m_k})$ and apply (a): it gives $m'$ with $\mu(P\cap W_{m'})\le\frac12\mu(P\cap W_{m'})$, so $P\cap W_{m'}$ is null, a contradiction.
\end{proof}
\status{proved (\src{K1} Lemma~K1.1, Cor.~K1.1$'$); confirmed (\src{K1\_referee} \S\S4.1, 4.2), re-derived step by step, one reader. (a) is a genuine extension of Lemma~\ref{lem:H21}(b) to countable families in the integrable case. (b) replaces the pigeonhole over finitely many ends in Theorem~\ref{thm:H23} and Proposition~\ref{prop:HR1}, and absorbs the depth adjustment of \src{H12\_referee} \S2.4: $\mu(W_m\cap S_n)>0$ for all $m,n$ is the arrow condition of $\vec{\mathcal H}$.}

\begin{theorem}\label{thm:K12}
Assume (EP$_\infty$) and fix $t\in I_\psi$. Call $(A_1,A_2,A_3)$ an \emph{admissible safe triple} if each $A_k$ is $k$-safe and $\sum_kt_k\rho_k\mu(E_kA_k)<\underline{\mathrm{Cost}}_t$. If some $C\in\C_L$ and some admissible safe triple satisfy $C\preceq_kA_k$ for $k=1,2,3$, then (H$_U$) is contradicted.
\end{theorem}
\begin{proof}[Proof sketch]
Each comparison is relevant and, by Lemma~\ref{lem:K10}, a base comparison, so $\rho_k\mu(E_kC)\le\rho_k\mu(E_kA_k)$. Summing, $\mathrm{Cost}_t(C)<\underline{\mathrm{Cost}}_t$, which no coherent $C$ can satisfy.
\end{proof}
\status{proved (\src{K1} Thm.~K1.2); confirmed (\src{K1\_referee} \S2.1), one reader. Lemma~\ref{lem:H10}(c) is the case $A_1=A_2=A_3$. \emph{Remarks of the referee} (\S2.2). Admissibility is a weak constraint: the triple $(P_1,P_2,P_3)$ is always admissible, and for it every $p$ with $|\supp p|\ge2$ and $C_p\notin\{P_k\}$ is anti-safe in the sense of Theorem~\ref{thm:K16}. The binding constraint is $C\preceq_kA_k$, so the useful triples are the worst admissible ones, the canonical triple below. The tool is used only in Corollary~\ref{cor:K13} and Theorem~\ref{thm:K16}; Theorems~\ref{thm:K14} and~\ref{thm:K15} rest on Lemmas~\ref{lem:K10} and~\ref{lem:K11} (the note attributed them to this theorem).}

\emph{The canonical triple.} For a depth assignment $a$ on all ends of all worlds with $a_S\ge n_S$, let $A^a_k$ be $B_t$ off $\mathcal E(a)$, correct at world $k$ on $\mathcal E^k(a)$, and wrong at world $k$ on $\mathcal E(a)\setminus\mathcal E^k(a)$. It is regular and $k$-safe, and $\sum_kt_k\rho_k\mu(E_kA^a_k)\le\mathrm{Cost}_t(B_t)+\int_{\mathcal E(a)}\sum_k\rho_k\,d\mu$; so the triple is admissible as soon as $\int_{\mathcal E(a)}\sum_k\rho_k<\Phi^*_t$, which holds for all large summable assignments. Every $k$-safe model at depths $a$ that equals $B_t$ off $\mathcal E(a)$ is at least as good as $A^a_k$ at world $k$, by (G).

\begin{corollary}\label{cor:K13}
Assume (EP$_\infty$). If some regular $X$ satisfies $X\preceq_kA_k$ for every $k$ for some admissible safe triple, (H$_U$) is contradicted. In particular there is no cover in class (EP$_\infty$) if, at some depth assignment $a_0$, deep ends of different worlds never conflict: $\mu(\mathcal E^i(a_0)\cap\mathcal E^j(a_0)\cap R_{ij})=0$ for all $i\ne j$.
\end{corollary}
\begin{proof}[Proof sketch]
Theorem~\ref{thm:kkm} gives a selected $C_p$ with $C_p\preceq_kX\preceq_kA_k$; apply Theorem~\ref{thm:K12}. For the second statement, conflict-freeness is preserved when the assignment is enlarged, so take $a\ge a_0$ admissible, and let $X$ be $B_t$ off $\mathcal E(a)$ and correct on $\mathcal E(a)$ at every world $k$ with $g\in\mathcal E^k(a)$; this is consistent, and $X\preceq_kA^a_k$ by (G).
\end{proof}
\status{proved (\src{K1} Cor.~K1.3); confirmed (\src{K1\_referee} \S3), one reader, with the hypothesis stated at an assignment $a_0$ and then enlarged (referee's precision). It is the second version of Theorem~\ref{thm:H13} for countable families. The note's claim of a gain ``at some depth instead of every depth'' is withdrawn: the condition is monotone in the depth, and the localizing families of Theorem~\ref{thm:H13} may be re-indexed. The gain is countable families.}

\begin{theorem}\label{thm:K14}
Assume (EP$_\infty$), and let $\vec{\mathcal H}$ be the conflict digraph of Definition~\ref{def:conflict} (an arrow $i\to l$ iff some $T_1$-wrong $W\in\mathcal S_i$ and $T_1$-correct $S\in\mathcal S_l$ have $\mu(W_m\cap S_m)>0$ for every $m$). If $\vec{\mathcal H}$ has no directed cycle, (H$_U$) is contradicted.
\end{theorem}
\begin{proof}[Proof sketch]
Take an assignment $a$ with $a_S\ge n_S$ and $\int_{\mathcal E(a)}\sum_k\rho_k<\Phi^*_t$, the model $X\coloneqq T_1$ on $\mathcal E(a)$ and $B_t$ off it, and a dominator $C$ from Theorem~\ref{thm:kkm}; put $\delta_k\coloneqq1_{E_kX}-1_{E_kC}$. At a $T_1$-correct end of world $k$, $\delta_k\le0$ deep in the end. If $\delta_i$ is not eventually $0$ on some end, end-protection makes it positive on some end, which is then a $T_1$-wrong $W\in\mathcal S_i$; on $P\coloneqq\{\delta_i=1\}\cap W_{n_W}$, $C$ and $X$ give different advice. Lemma~\ref{lem:K11}(b) gives $l\in V(W)$ and a $T_1$-correct $S\in\mathcal S_l$ with $\mu(P\cap W_m\cap S_n)>0$ for all $m,n$, so $i\to l$, and there $\delta_l=-1$ at every depth; so $\delta_l$ is not eventually $0$ on $S$, and the step repeats at $l$. An infinite walk on three vertices contains a directed cycle. If $\vec{\mathcal H}$ is acyclic, all three comparisons are base comparisons and are endorsed, against Lemma~\ref{lem:H10}(c).
\end{proof}
\status{proved (\src{K1} Thm.~K1.4); confirmed (\src{K1\_referee} \S4.3), one reader. This is Proposition~\ref{prop:HR1} (and \src{J1} Thm.~J1.1) for countable families, without heaviness.}

\emph{Switched models with countable families.} Fix an order $(k_1,k_2,k_3)$ and an assignment $a$. For $g\in\mathcal E(a)$ let $r(g)\coloneqq\min\{r:g\in\mathcal E^{k_r}(a)\}$; the switched model $X^a$ is $B_t$ off $\mathcal E(a)$ and gives at $g\in\mathcal E(a)$ the advice correct at world $k_{r(g)}$. Its conflict sets are $\Xi^s(a)\coloneqq\{g\in\mathcal E(a):r(g)<s,\ \sigma_{k_{r(g)}}(g)\ne\sigma_{k_s}(g)\}$, $s=2,3$. Condition \textbf{(Al$_\infty$)} holds at $a$ if, for $s=2,3$ and every assignment $M$ on $\bigcup_{q<s}\mathcal S_{k_q}$, there is an assignment $J\ge a$ on $\mathcal S_{k_s}$ with
\[
\Bigl(\bigcup_{S'\in\mathcal S_{k_s}}S'_{J_{S'}}\Bigr)\cap\Xi^s(a)\subseteq\bigcup_{q<s}\bigcup_{S\in\mathcal S_{k_q}}S_{M_S}\quad\text{modulo null sets.}
\]
The order is \emph{aligned} if for every assignment $a_0$ there is $a\ge a_0$ at which (Al$_\infty$) holds. For finite families and uniform depths, (Al$_\infty$) is (Al) of Theorem~\ref{thm:J12}; with non-uniform depths ``aligned'' is weaker, so the next theorem strengthens Theorem~\ref{thm:J12} slightly even for finite families.

\begin{theorem}\label{thm:K15}
Assume (EP$_\infty$). If some order of the worlds is aligned, (H$_U$) is contradicted.
\end{theorem}
\begin{proof}[Proof sketch]
Choose $a$ with (Al$_\infty$) and $\int_{\mathcal E(a)}\sum_k\rho_k<\Phi^*_t$, the switched model $X\coloneqq X^a$ and its dominator $C$. (1) $X$ is correct at $k_1$ on $\mathcal E^{k_1}(a)$, so by Lemma~\ref{lem:K10} the world-$k_1$ comparison is a base comparison and $C=X$ on $S_{m'_S}$ for $S\in\mathcal S_{k_1}$. (2) Apply (Al$_\infty$) with $s=2$ and $M\coloneqq m'$. At $g\in S'_{J_{S'}}$, either $X$ is correct at $k_2$, so $\delta_{k_2}(g)\le0$, or $g\in\Xi^2(a)$, so $g\in S_{m'_S}$ for some $S\in\mathcal S_{k_1}$, where $C=X$ and $\delta_{k_2}(g)=0$. Lemma~\ref{lem:K10} applies again. (3) The same with $s=3$. All three comparisons are base comparisons, against Lemma~\ref{lem:H10}(c).
\end{proof}
\status{proved (\src{K1} Thm.~K1.5); confirmed (\src{K1\_referee} \S4.4), including the order of quantifiers, one reader. \emph{Sufficient condition, with the referee's repair} (\src{K1\_referee} \S5.1). Let $S'\in\mathcal S_{k_s}$ be $T_1$-wrong, with every world of $V(S')$ that has ends placed before $k_s$; then Lemma~\ref{lem:K11}(a) gives (Al$_\infty$) for the part of $\Xi^s(a)$ inside $S'$. Pairs of ends in a common orthant that are both $T_1$-correct, or both $T_1$-wrong, at worlds of opposite sign are disjoint beyond their purity depths. So the order is aligned if every later end on a conflicting orthant is of one of these kinds, or if there is one assignment $a_1$ such that every remaining later end $S'$ on a conflicting orthant has a depth $j(S')$ with $\mu\bigl(S'_{j(S')}\cap\bigcup_{q<s}\mathcal E^{k_q}(a_1)\bigr)=0$; then every $a\ge a_1$ (and $\ge a_0$) satisfies (Al$_\infty$) for these ends, with $J_{S'}\coloneqq\max(j(S'),a_{S'})$. The note's pairwise reading, ``each such later end is eventually disjoint from each earlier end'', is false with countably many ends, since it can force infinite depths on an earlier end. So Corollary~\ref{cor:J13}(a) extends to countable families, (c) extends with its hypothesis read with assignments, and (b) needs the uniform reading for its disjoint pairs. Theorem~\ref{thm:JR2} and Corollary~\ref{cor:JR3} are not extended as stated; case~(4) of Theorem~\ref{thm:JR2} needs the same uniform reading.}

\emph{Persistent conflict without an arrow} (\src{K1\_referee} \S5.2; CHECKED, a configuration of sets with no ranking). \src{K1} said that conflicts persist at every depth only when the deep ends of two worlds overlap on $R_{ij}$ at every depth, as an arrow does. That is true in one direction only. In a conflicting orthant let world~1 have one $T_1$-wrong end $W$, the cones over the caps of radius $0.2\cdot2^{-m}$ around a ray $r$, and world~2 the $T_1$-correct ends $S'_l$, the cones over the caps of radius $0.01\cdot2^{-l-n}$ around rays at angle $0.15\cdot2^{-l}$ from $r$. Each pair $(W,S'_l)$ is disjoint from depth $l+1$ on, so there is no arrow; but for every depth $a$ of $W$ and every $M>a$ the end $S'_a$ lies, at every depth, in the shell $W_a\setminus W_M$. So $\mathcal E^1(a)\cap\mathcal E^2(a')\cap R_{12}$ is non-null for every pair of assignments, and (Al$_\infty$) fails for every order placing world~1 before world~2 unless other earlier ends cover the shells. (\src{k1ref\_py/countable\_conflict.py}: the angle inequalities hold on a grid with 0 failures, and the non-nullity is confirmed by Monte Carlo.) So an arrow forces a persistent conflict, and with countably many ends a persistent conflict can also come from infinitely many pairwise eventually disjoint ends. With countable families, Corollary~\ref{cor:K13} and Theorem~\ref{thm:K14} are logically independent exclusions.

\begin{theorem}\label{thm:K16}
Assume (EP$_\infty$). Call $p\in\Delta$ \emph{anti-safe} for an admissible safe triple $A$ if $C_p\not\preceq_kA_k$ for every $k\in\supp p$.
\begin{enumerate}[leftmargin=*,label=(\alph*)]
\item In a cover in class (EP$_\infty$), for every tuple $\rho$ of integrable base endorsers, every $t\in I_\psi$ and every admissible safe triple $A$, the selection has an anti-safe point for $A$.
\item Let $p$ be anti-safe for the canonical triple $A^a$ ($a$ admissible). Then: (1) $|\supp p|\ge2$; (2) for every regular $X$ that equals $B_t$ off $\mathcal E(a)$, the worlds $k\in\supp p$ with $C_p\preceq_kX$ (there is at least one) are all worlds at which $X$ errs on a non-null part of $\mathcal E^k(a)$; (3) $\supp p$ contains $i\ne j$ with $\mu(\mathcal E^i(a)\cap\mathcal E^j(a)\cap R_{ij})>0$; (4) if $\supp p=\{i,j\}$, then $C_p\preceq_jY$ for every $i$-safe $Y$ at depths $a$ that equals $B_t$ off $\mathcal E(a)$, and $C_p\preceq_iY$ for every such $j$-safe $Y$, while $C_p\not\preceq_iY$ and $C_p\not\preceq_jY'$ for every such $i$-safe $Y$ and $j$-safe $Y'$.
\item Conversely, if for one $\rho$, one $t\in I_\psi$ and one admissible safe triple no point of $\Delta$ is anti-safe, there is no cover in class (EP$_\infty$) with that ranking.
\end{enumerate}
\end{theorem}
\begin{proof}[Proof sketch]
(a) The sets $F_k\coloneqq\{p:C_p\preceq_kA_k\}$ are closed by (S). If no point were anti-safe, every face $\Delta_S$ would lie in $\bigcup_{k\in S}F_k$, and the KKM lemma would give $p^*\in F_1\cap F_2\cap F_3$, against Theorem~\ref{thm:K12}. (c) is (a) read backwards. (b) (1) At a vertex, $C_{e_k}$ is competitive on $\{k\}$. (2) A model correct at world $k$ on $\mathcal E^k(a)$ is at least as good as $A^a_k$ at world $k$, by (G). (3) Otherwise the model of Corollary~\ref{cor:K13}, built from the worlds of $\supp p$, contradicts (2). (4) follows from (2) and competitiveness on $\{i,j\}$.
\end{proof}
\status{proved (\src{K1} Thm.~K1.6); confirmed (\src{K1\_referee} \S\S6.1, 6.2), one reader. The referee checks that the constraints of (b)(4) are not contradictory in themselves: they place $C_p$ strictly between a $j$-safe model and an $i$-safe one at world $j$. By the configuration above, the conflicting pair of (b)(3) need not give an arrow of $\vec{\mathcal H}$. \emph{What the theorem is and is not} (\src{K1\_referee} \S6.3). \src{K1} poses Question K1.Q (is there a ranking in class (EP$_\infty$) whose selection has an anti-safe point for every $\rho$, $t$ and admissible safe triple?) and proves that a cover in the class exists iff K1.Q has a positive answer. That equivalence is correct but tautological: a ranking as in K1.Q is a cover in the class by definition, and by (a) membership in the class already implies the anti-safe clause. So K1.Q is the class question restated; it is not a reduction, it does not separate the selection from the ranking, and it removes no axiom. What the theorem gives is a necessary condition, (a) and (b), and an exclusion test, (c). \emph{A selection-level test that would be a genuine reduction} (suggested by the referee; not proved, not conjectured, untested): whether any two-world order on $R_{ij}$, with closed sections and competitiveness on the edge, can realise (b)(4) at every large $a$ for edge indices; if none can, every anti-safe index is interior. \emph{Failed repairs} (\src{K1} \S6, confirmed as descriptions). A monotone iteration of $i$-safe switched models has the proved properties stated in \src{K1}, but (S) gives closedness in $p$ for a fixed model, never continuity in the model, so the limit keeps a shell at one of the two worlds; a fixed point through (IB) would give a contradiction through Lemma~\ref{lem:H10}(c), but its sets have a second argument that moves with $p$, and (S) says nothing about them. The note's counting heuristic makes no statement and is not used.}

\emph{Question J1.Q$'_\infty$} (\src{K1\_referee} \S6.3; open, no conjecture stated). Under (EP$_\infty$), if $\vec{\mathcal H}$ has a directed cycle and no order of the worlds satisfies (Al$_\infty$), is (H$_U$) contradictory? By Theorems~\ref{thm:K14} and~\ref{thm:K15} this is what remains of class (EP$_\infty$); it is a strict sub-question of the class question, and its finite form is Question~J1.Q$'$ of part~(f$'$). So untouched case (iii) of \src{BRIEF16} (countably many ends) is closed except along cyclic configurations with no aligned order. The note's ``doubly misaligned cycles are exactly the configurations in which no order is aligned'' holds only for case (R-i) of part~(f$'$) (\src{K1\_referee} \S5.3).

\paragraph{(g) What a cover of $T_1$ still needs, and the precise remaining questions.} \emph{What a cover on $\M_U$ still needs} (rewritten in round 16, again after the audit of 15 October, following \src{F2} \S4.3 as corrected in \src{F12\_referee} \S5.7, again in round 17, following \src{H12\_referee} \S3 and \src{J1}, and again in round 19, following \src{K1}, \src{K2}, \src{K3} and the four referee reports of parts~(f$'$) and~(f$''$)). Assume (H$_U$). In order:
\begin{enumerate}[leftmargin=*,label=(\arabic*)]
\item \emph{A cover needs a world that does not add up} (Corollary~\ref{cor:M0closed}(a)). Such worlds are no longer hypothetical: Theorem~\ref{thm:D21} gives one in a full two-world system on $\M_U$, next to an additive world, endorsed by a $\sigma$-finite density that fails to be integrable only at the protected end of its chain (Proposition~\ref{prop:D2sigma}); that end contains the ray $\R_+(1,0)$, interior to $U$, so no endorsing measure is finite on compact subsets of $U$ (Theorem~\ref{thm:F13}). Theorem~\ref{thm:F16} gives two more, with generators and protected set where both payoffs are bounded away from zero and the protection at $\partial U$. Theorem~\ref{thm:H24} gives a full two-world system in which neither world adds up. None is known in three worlds or in a cover.
\item \emph{Such a world is of one of three kinds} (Theorem~\ref{thm:K31}). (i) It carries a finite certificate (mode (M1)). (ii) It has no finite certificate, and its non-localizable blind set $\Lambda_\infty$ is non-null: for every localizing family, ratio certificates made of localized comparisons remove mass from every part of $\Lambda_\infty$ (Theorem~\ref{thm:K31}(f)). (iii) It has no finite certificate and $\Lambda_\infty$ is null; then it is localizable with an integrable density, (Loc$_{\mathrm{int}}$) (Theorem~\ref{thm:K31}(e)). A world that is localizable only with non-integrable densities is of kind (ii) unless it has a certificate. A world that adds up is (Loc$_{\mathrm{int}}$) as well.
\item \emph{Kind (i): a finite certificate.} Every certificate contains a rigid comparison, one whose proportional parts are not all accepted (Theorem~\ref{thm:K21}); next to worlds with leading densities such a comparison is not forced by competitiveness (Proposition~\ref{prop:K22}, Corollary~\ref{cor:K22p}). Over additive bases of Brier type the natural construction scheme, a dual-cone decoration of a Bayes base (Theorem~\ref{thm:K25}), is empty (Theorem~\ref{thm:K26}, Proposition~\ref{prop:KR62}), and as a route to $T_1$ it is excluded in any case (Theorem~\ref{thm:main}). In two worlds a certificate makes the system a counterexample to Conjecture~\ref{conj:Q2N} (Theorem~\ref{thm:K23}). No full system with a certificate is known, for any target (Question~Q-K2, Conjecture~\ref{conj:K26p}).
\item \emph{Kind (ii): a non-null blind set $\Lambda_\infty$.} Whether this can happen in a full system is open: Question~Q$_{\mathrm{K3}}$ of part~(f$''$). No axiom other than closedness of the sections excludes it (Theorem~\ref{thm:K33}). Next to two worlds with integrable full endorsers, such a world is blind on a set of the rigid shape of Theorem~\ref{thm:K37}(b$'$), and, unless $\Lambda_\infty\supseteq E_1(T_1)\cap M$, it needs a positive answer to Question~Q$_{\mathrm A}$; in two worlds next to an additive world it is excluded if Conjecture~\ref{conj:K3C} holds.
\item \emph{Kind (iii): localizable worlds.} Suppose every world of the cover is localizable with a.e.\ finite densities, (Loc); this holds in particular if no world is of kind (i) or (ii). Then, by the patch theorem (Theorem~\ref{thm:K35}, Corollary~\ref{cor:K35p}(c)), a cover needs \emph{sheltering}, two disagreeing worlds whose localizing sets overlap on $R_{ij}$ at every depth, or \emph{infinite minority mass}. Infinite minority mass needs a world that is not (Loc$_{\mathrm{int}}$), and such a world, if it has no certificate, has a non-null $\Lambda_\infty$. So if Question~Q$_{\mathrm{K3}}$ has answer No, infinite minority mass is excluded, and a cover needs a world with a certificate or sheltering (\src{K3\_referee} \S5.3). $T_1$-wrong blow-ups of $\sigma$-finite endorsers are always sheltered or shadowed (Lemma~\ref{lem:K36}). For sheltering by ends: within class (EP$_\infty$) it survives only along conflict cycles with no aligned order (Theorems~\ref{thm:K14}, \ref{thm:K15}), which is Question~J1.Q$'_\infty$; with finitely many ends, only in the cases (R-i)--(R-iii) of part~(f$'$), which is Question~J1.Q$'$, with Question~J1.Q its case (R-i).
\end{enumerate}
Every world of a cover that does not add up has, in addition, all of the following properties (each with the status of the result cited).
\begin{itemize}[leftmargin=*]
\item It is not (CP-a), and it has no leading charge nonzero on $M$ (Corollary~\ref{cor:CPaU}(b), Theorem~\ref{thm:S21}).
\item On one of its two edge regions $R_{ij}$, $R_{ik}$ it has no leading density (Theorem~\ref{thm:B13}, Corollary~\ref{cor:B14}); indeed no integrable $\rho\ge0$ that is not a.e.\ zero on that region is leading there (CHECKED). This concerns countably additive densities only: a world with finitely additive leading charges on both edge regions is not covered, because gluing such charges through the link class is open.
\item It is not a glued credal world (Theorem~\ref{thm:C16}). If it is $\Gamma$-leading for finitely generated $\Gamma_{ij},\Gamma_{ik}$ with (UI), it is not total (Corollary~\ref{cor:C18}, CHECKED) and it carries a finite $\Gamma^{\mathrm{gl}}$-obstruction (Proposition~\ref{prop:C19}). (Round 15; refereed once.)
\item If it is total, it is not AR on $M$ (Corollary~\ref{cor:B32}), and some edge family at $e_i$ is not nested at world $i$ (Theorem~\ref{thm:Q36}(c), CHECKED). Only this is proved; that such a world cannot be total at all would follow from Conjecture~\ref{conj:B3C}, which is open. Total or not, it has no nontrivial cancellative leading top level with dense null sets (Theorem~\ref{thm:B35}; in particular none of signature type, Corollary~\ref{cor:B36}) and no leading top level $\nu+\kappa\circ S$ with $4\max|\kappa|<\nu(M)$ (Proposition~\ref{prop:B37}).
\item If world $j$ is additive with a density on $R_{ij}$, world $k$ is additive with a density on $R_{ik}$, and world $i$ is a density unanimity on both regions with one family $\mathcal A$, then on no non-null part of $Q_i$ does the pair of edge chain coordinates have a law equivalent to Lebesgue measure on a connected open set whose projections contain their essential ranges (Proposition~\ref{prop:C26}); when the densities of the reference member and of the neighbours are of linear or normalised-box type this case cannot occur (Corollary~\ref{cor:C27}). (Round 15; refereed once.)
\end{itemize}
\emph{The cover as a whole} (round 16, proved in note \src{E1}, audited in note \src{F2}, audit refereed once; rewritten in round 17 after \src{H1} and \src{H2}, refereed once in \src{H12\_referee}, and after \src{J1}, refereed once in round 19; extended in round 19 after \src{K1} and \src{K3}, each refereed once). Every cover shows non-additive behaviour at $T_1$ itself (Theorem~\ref{thm:R3}). For every admissible shadow $\nu$, Theorem~\ref{thm:E11}(a), (b) give a non-null set of weights $t$ at which \emph{some} world $j$, which may depend on $\nu$ and $t$, has a dominator comparison into the incoherent $B^\nu_t$ that $\nu$ does not endorse; for suitable $\nu$ these incoherent models lie within any $\eta$ of $T_1$ (Theorem~\ref{thm:E12}). Under (H$_U$) this is automatic, so it is not a restriction on any one world. Round 17 makes it robust: the same holds for the dominator comparisons into every regular model whose cost is within $\Phi^*_t$ of the Bayes cost (Lemma~\ref{lem:H10}(c)), and such models may agree with the Bayes model off any set of small measure and be correct on it. So the failures that a cover needs cannot be confined to sets of vanishing measure at a world whose density is integrable off them. In particular the following covers are excluded.
\begin{itemize}[leftmargin=*]
\item Covers in which every world is an end-protected decoration with integrable, jointly end-heavy base densities (Theorem~\ref{thm:E15}), now a special case of Theorem~\ref{thm:H13}.
\item Covers with two worlds whose integrable densities endorse all their relevant comparisons and a third world whose $\Theta$-localized comparisons, for some localizing family $\Theta$, are endorsed in the cancellation form by an a.e.\ finite density (Theorem~\ref{thm:H12}). This includes every end-protected third world with countably many ends whose base has such an endorser, end-heavy or end-light, and every third world that adds up $\sigma$-finitely in the cancellation form. So a three-world installation of world~1 of Theorem~\ref{thm:D21} next to two such worlds is excluded as soon as it keeps an a.e.\ finite endorser in the cancellation form.
\item Covers in which every world has an integrable density that endorses its $(\Theta^k,T_1)$-localized comparisons (Theorem~\ref{thm:H13}); in particular covers by end-protected worlds with integrable base endorsers whose ends are all eventually $T_1$-correct at their worlds, with no heaviness.
\item Covers in which every world is end-protected with finitely many ends and has an integrable base endorser, and the conflict digraph is acyclic (Theorem~\ref{thm:H23}, Proposition~\ref{prop:HR1}; the proposition is proved in a referee report).
\item The same class with some aligned order of the worlds, in particular the configurations of Corollary~\ref{cor:J13} and those of Theorem~\ref{thm:JR2} and Corollary~\ref{cor:JR3} (Theorem~\ref{thm:J12}, refereed once in round 19; Theorem~\ref{thm:JR2} and Corollary~\ref{cor:JR3} proved in a referee report).
\item Round 19: covers in which every world is end-protected with \emph{countably} many ends and has an integrable base endorser, and the conflict digraph is acyclic, or some order of the worlds is aligned in the sense (Al$_\infty$), or deep ends of different worlds never conflict at some depth assignment (Theorems~\ref{thm:K14}, \ref{thm:K15}, Corollary~\ref{cor:K13}).
\item Round 19: covers for which a patch exists (Theorem~\ref{thm:K35}): a regular model on whose world-$k$ error set every density $\rho_k$ is integrable, with $\rho_k$ a.e.\ finite and positive and endorsing the world-$k$ comparisons localized off that error set. This contains Theorem~\ref{thm:H12} with $\rho_2$, $\rho_3$ integrable only on $R_{12}$, $R_{13}$; Theorem~\ref{thm:H13} with densities integrable only on the error sets of $T_1$ (Corollary~\ref{cor:K35p}(b)); and covers in which every world is localizable with compatible localizations and finite minority mass (Corollary~\ref{cor:K35p}(c)).
\item Round 19: covers in which some world has a certificate made of forced comparisons only (Corollary~\ref{cor:K22p}).
\end{itemize}

\emph{The precise residual of the end-protected classes} (\src{H12\_referee} \S3, narrowed by \src{J1}, \src{J1\_referee} and \src{K1}). Let \emph{class (EP)} consist of the rankings with (H$_U$) in which every world is end-protected with a finite family of ends, refined as in Lemma~\ref{lem:H20}, and has an integrable density, positive a.e., that endorses its base comparisons (hypothesis (EP) of part~(f$'$)). Within class (EP), a cover of $T_1$ needs all of the following.
\begin{enumerate}[leftmargin=*,label=(\arabic*)]
\item A directed cycle, a $2$-cycle or a $3$-cycle, in the conflict digraph $\vec{\mathcal H}$ (Theorem~\ref{thm:H23}, Proposition~\ref{prop:HR1}). Hence at least two worlds have $T_1$-wrong ends and at least two worlds have no integrable density that endorses all their relevant comparisons; if the ends on the cycle have one-element sets $V(W)$, at least two worlds fail to add up (Corollary~\ref{cor:H23p}).
\item For every admissible $\psi$ with (b), every $t\in I_\psi$ and every large $m$, with $X_m=T_1$ deep in the ends and $X_m=B^\psi_t$ elsewhere: every dominator of $X_m$ from Theorem~\ref{thm:kkm} is positive at a $T_1$-wrong end of a world on the cycle; along the walk it gives the advice opposite to that of $T_1$ on non-null sets at every depth of the $T_1$-wrong ends of the cycle; and it errs at every depth on the $T_1$-correct ends that shelter them (the germ exchanges; \src{H12\_referee} \S3).
\item No order of the three worlds is aligned (Theorem~\ref{thm:J12}; refereed once in round 19). By Theorem~\ref{thm:JR2} and Corollary~\ref{cor:JR3} the cover is then in one of the cases (R-i)--(R-iii) of part~(f$'$). In case (R-i), a $2$-cycle whose third world is incident to no arrow, both arrows are misaligned, so the cycle is doubly misaligned, and the exchanges take place in shells of the $T_1$-wrong ends, at bounded depth in them (Proposition~\ref{prop:J15}); whether this can happen under every axiom is Question~J1.Q. In case (R-ii), a pure $3$-cycle, all three arrows are misaligned in the weaker sense of Corollary~\ref{cor:JR3}(b), and the positive set lies in shells (CHECKED, not proved). In case (R-iii) no description by shells is available. The configurations of Corollary~\ref{cor:J13}, the template of \src{H2} among them, are excluded. The single question that would settle class (EP) is Question~J1.Q$'$.
\end{enumerate}
\emph{Class (EP$_\infty$)}, the same with countable families of ends (part~(f$''$)). Theorems~\ref{thm:K14} and~\ref{thm:K15} give (1) and (3) with (Al$_\infty$) in place of (Al), and Corollary~\ref{cor:K13} excludes conflict-free depth assignments. What is left is Question~J1.Q$'_\infty$. Theorem~\ref{thm:K16} adds a necessary condition at the level of the selection: for every admissible safe triple some selection index is anti-safe, with two worlds of its support in persistent conflict.

\emph{Outside class (EP)} (\src{H12\_referee} \S\S1.7, 3; the cases (A$'$), (B), (D) at the end of part~(f$'$); round 19 additions marked). Nothing proved excludes a cover of $T_1$ on $\M_U$ of one of the following kinds.
\begin{itemize}[leftmargin=*]
\item \emph{Non-integrable endorsers at two or more worlds} (case (D)). Every world has an a.e.\ finite endorser of its localized comparisons, but at least two of these are non-integrable, or one is and the other worlds have only localized integrable endorsers. Theorems~\ref{thm:H12} and~\ref{thm:H13} do not apply, and the near-Bayes models need not have finite costs. The nearby target of round 16 (\src{F2} \S4.3: Theorems~\ref{thm:E11}(c) and~\ref{thm:E15} for shadows that are finite a.e.\ and integrable off every deep end) is proved when only one world's density is non-integrable and the other two worlds have integrable full endorsers (Theorem~\ref{thm:H12}); otherwise it is open, and it needs the referee's added condition that the cost differences on $C\sd B^\nu_t$ be finite (\src{F12\_referee} \S5.7). \emph{Round 19:} the patch theorem proves the exclusion for a.e.\ finite densities at several worlds whenever some regular model has integrable error mass at every world and every density endorses the comparisons localized off that error set (Theorem~\ref{thm:K35}, Corollary~\ref{cor:K35p}(b), (c)). What is left of case (D) is sheltering and infinite minority mass (Corollary~\ref{cor:K35p}(c)); the latter needs a world that is not (Loc$_{\mathrm{int}}$), and is excluded if Question~Q$_{\mathrm{K3}}$ has answer No.
\item \emph{Countably many ends.} Theorem~\ref{thm:H23}, Proposition~\ref{prop:HR1} and Theorem~\ref{thm:J12} need finite families: the pigeonhole of the walk, the uniform depth of Lemma~\ref{lem:J10}, and $\mu(\mathcal E_m)\to0$. Theorem~\ref{thm:H13} still excludes the case in which every end is eventually $T_1$-correct at its world. \emph{Round 19} (\src{K1}, \src{K1\_referee} \S9): with integrable base endorsers (class (EP$_\infty$)), an acyclic conflict digraph and an order aligned in the sense (Al$_\infty$) are excluded (Theorems~\ref{thm:K14}, \ref{thm:K15}), and so are conflict-free depth assignments (Corollary~\ref{cor:K13}); what remains inside (EP$_\infty$) is Question~J1.Q$'_\infty$. With countably many ends a persistent conflict need not come from an arrow (the configuration after Theorem~\ref{thm:K15}).
\item \emph{Failures not of end-protected form.} Protection that is not given by ends, and non-localizable failures (case (A$'$)): a world whose localized comparisons, for every localizing family, have no a.e.\ finite endorser in the cancellation form. This includes (M1) obstructions and (M2) families whose sources and targets have measure bounded below on a fixed core; in two worlds next to an additive world, no certificate of the type of Theorem~\ref{thm:D21} is of the last kind (Theorem~\ref{thm:F14}). \emph{Round 19:} by Proposition~\ref{prop:K32}, end-protection without a base density carries no information, so the question is localizability itself; and case (A$'$) is exactly ``not (Loc)'', which by Theorem~\ref{thm:K31} consists of worlds with a certificate and of worlds in (N), which have a non-null $\Lambda_\infty$. A world that adds up, even only through weights with a nonzero purely finitely additive part, is (Loc$_{\mathrm{int}}$) (Theorem~\ref{thm:K31}(c)--(e); \src{K3\_referee} \S5.3), so it is not in case (A$'$).
\end{itemize}
Known examples: Theorem~\ref{thm:C21} gives a world without a leading density in a full system, and it adds up; Theorems~\ref{thm:D21} and~\ref{thm:F16} give worlds that do not add up in full two-world systems, and for the world of Theorem~\ref{thm:D21} the failure is invisible in the sense above (Proposition~\ref{prop:E16}); Theorem~\ref{thm:H24} gives a full two-world system in which neither world adds up, and it is not a conflict cycle. Theorem~\ref{thm:K33} gives a world in the class (N), with $\Lambda_\infty=M$, in a two-world system with every axiom except closed sections at that world. No world with the properties of the residual is known, and no full system with a world of kind (i) or (ii) is known. What remains is listed below: the two-world questions (R1)--(R4) of rounds 12 and 13, which bear on $T_1$ through the edge systems (Proposition~\ref{prop:Q35}, Remark~\ref{rem:straighten}), with (R2) now answered negatively; the three-world questions of rounds 14 to 19; and the questions off the $U$-domain.

\emph{Two worlds.} Under (Nest), with every axiom and closed sections for every model (\src{Q12\_referee} \S6):
\begin{itemize}[leftmargin=*]
\item[\textbf{(R1)}] \emph{Mirrored simple pairs.} Is there a system with a non-null $Z$ and a uniformly Ces\`aro-thin family $\mathcal A$ with $Z\le_1A$ for all $A\in\mathcal A$? By Proposition~\ref{prop:N1} and Corollary~\ref{cor:MPloop} this is the same as $c_{s_0}\le_1A$ for all $A\in\mathcal A$, for some $s_0\in(0,1)$, and then $d_t\le_2A$ for all $A\in\mathcal A$ and all $t\ge T(s_0)$ (Proposition~\ref{prop:MPsharp}). By Proposition~\ref{prop:N2} and Lemma~\ref{lem:loop} the ping-pong generates only favourable-direction comparisons and never raises the exported segment, so a proof that no such system exists must use closedness quantitatively; without closed sections at both worlds such systems exist (Proposition~\ref{prop:S31}). \emph{Weaker target} (round 13): (R1$'$), the failure of (QC$_1$), is the same question with ``$\bar\mu(A_k)\to0$'' in place of ``Ces\`aro-thin''; it is implied by (R1), is itself mirrored (Proposition~\ref{prop:QCJ}), and excluding it (Conjecture~\ref{conj:QC}) excludes (R1). \emph{Known constraints on a construction:} both worlds fail SP$^{\mathrm{reg}}$, and neither is complete or (CP-a) or has a leading charge nonzero on $M$; the targets do not concentrate (Lemma~T); no world packs its targets towards its far end (Proposition~\ref{prop:S33}); (K) at cuts $t\ge T(s_0)$ cannot transport them (Lemma~\ref{lem:block}); and the system is neither of band-and-translates type with closed sections (Proposition~\ref{prop:S32}) nor of fibrewise band-and-translates type (Proposition~\ref{prop:S34}). \emph{Any exclusion} must obtain world-1 packing of a thin $A\cap c_t$ from (K) at cuts below $T(s_0)$, or from world-2 closedness at sources that are not segments (\src{S34\_referee} \S12).
\item[\textbf{(R2)}] \emph{General certificates.} If neither world has a simple certificate, must both satisfy SP$^{\mathrm{reg}}$? Equivalently (Corollary~\ref{cor:IS}): can $c_a$ be $\mathcal V_1$-null through ratio certificates whose sources are thin pieces of $c_a$, with overlapping targets? This is open even when world~2 is additive (\src{Q1} CONJECTURE Q1.C1). Round 13: there the qualitative rate is available (Corollary~\ref{cor:SR}), so a proof of (R2) must use the failure of SP$^{\mathrm{reg}}$ itself; and the only rate that the route through Proposition~\ref{prop:MP} can use is linear (Proposition~\ref{prop:LIN}), which is not a consequence of the axioms (Proposition~\ref{prop:ExA}). \emph{Round 16: answered negatively.} In the full system of Theorem~\ref{thm:D21}, under (Nest) and (C), world~2 is additive (so it has no simple certificate), world~1 has no simple certificate (Proposition~\ref{prop:Q15}), and world~1 fails SP$^{\mathrm{reg}}$: $c_a$ is $\mathcal V_1$-null through ratio certificates whose sources are thin pieces of $c_a$ and whose targets overlap in every neighbourhood of the end of the chain. So Q1.C1 is false. With $\sigma$-finite endorsing measures in place of charges, world~1 does satisfy the analogue of SP$^{\mathrm{reg}}$ (Proposition~\ref{prop:D2sigma}). The form of that analogue with densities locally integrable on $U$, Conjecture~\ref{conj:sigmaB}, is false (Theorem~\ref{thm:F13}); the form with local integrability on $M$ is Conjecture~\ref{conj:F1Cp}.
\item[\textbf{(R3)}] \emph{A rate in Lemma~T.} For non-null $P\subseteq c_b$ and $x>b$, does the radius $\eta(P,x)$ of Lemma~T$_1$ admit a positive lower bound depending only on $\bar\mu(P)$ and $x-b$? This qualitative rate is equivalent to the same rate at world~2, to a uniform rate in strict escape at $b=0$ at either world, and to the absence of doubly-light sequences (sets of measure bounded below that world~1 ranks below sets concentrating at $1$ and world~2 below sets concentrating at $0$) (\src{S4} Thm.~RATE, Theorem~\ref{thm:RATE}). It holds whenever one world has a strictly positive endorsing charge vanishing at the far end of the chain, in particular next to an additive world with integrable density (\src{S4} Cor.~SR, Corollary~\ref{cor:SR}); so it does not give \src{Q1} Q1.C1, which is open in exactly that setting (round 16: Q1.C1 is false, Theorem~\ref{thm:D21}, where this rate holds), and it does not by itself give a Vitali argument for N17.B$'$. Against the ratio certificates of N4, the route of Proposition~\ref{prop:MP} needs a rate linear in $\bar\mu(P)$, $\eta_1(P,1)\ge c\,\bar\mu(P)$ (\src{S4} Prop.~LIN, Proposition~\ref{prop:LIN}); this fails in a system with every axiom and both worlds additive ($\eta_1(P,1)\le\bar\mu(P)^2/2$, \src{S4} Example~A, Proposition~\ref{prop:ExA}), so it is not a consequence of the axioms and could only come from the failure of SP$^{\mathrm{reg}}$ itself. No rate is known to bear on (R1), whose sources are fixed. \emph{Restated} (round 13): what remains of (R3) is Conjecture~\ref{conj:NDL} (no doubly-light sequences), open only when neither world has an end-vanishing endorsing charge; its linear form is false. (Correction: the earlier text called (R3) ``the missing lemma common to (R1) and (R2)'' and said that it ``would give Q1.C1, and with Theorem~\ref{thm:LF} a Vitali-type argument for N17.B$'$ in two worlds''. By Corollary~\ref{cor:SR} the qualitative rate is a theorem in the setting of Q1.C1, where world~1 has no simple certificate, so each of these claims would amount to an unproved claim of Q1.C1 itself, and in (R1) the sources are fixed; \src{S4} \S4.3, \src{S34\_referee} \S9.3.)
\item[\textbf{(R4)}] \emph{Nestedness from (S).} CONJECTURE~\ref{conj:Q2N} in two worlds, and its edge form (Nest$_e$) in three: on the coherent edge models of Proposition~\ref{prop:comonotone}, ``(S) forces comonotone endpoints''. Round 13: in two worlds it is true when one world has a nonzero leading charge (Theorem~\ref{thm:S22}) or is monotonely continuous (Theorem~\ref{thm:MC}); it is open only when both worlds lack both (Corollary~\ref{cor:MC}(b)); (C) is necessary for it (Proposition~\ref{prop:R71}); the reduced form is Conjecture~\ref{conj:S1R}; and in a non-nested system strict escape still holds (Corollary~\ref{cor:continuum}(d)). On edges, (Nest$_e$) holds in the cases of Proposition~\ref{prop:RS2A} and is not implied by the edge inputs (Proposition~\ref{prop:plateau}); the right target is straightening (Remark~\ref{rem:straighten}), proved pairwise under (EU$_{ij}$) and monotone continuity (Proposition~\ref{prop:S112}) and conjectured globally (Conjecture~\ref{conj:S113}). After Theorem~\ref{thm:S21}, edge nestedness matters only for complete worlds and for general acceptance cones. Round 14: after Corollary~\ref{cor:B32}, for total worlds it matters only as one route to AR on $M$, which Conjecture~\ref{conj:B3C} would make unnecessary; and on an edge a crossing path can coexist with a world that is not fine (Proposition~\ref{prop:B34}; realised without (C) in a two-world ranking, see its status).
\end{itemize}
Any example must fail SP$^{\mathrm{reg}}$ at both worlds if its certificate is simple, be non-additive in the chain direction at both worlds (Propositions~\ref{prop:NS}, \ref{prop:Q15} and the tilt example), and have closed sections; in band-and-translates worlds, closed world-1 sections force $\inf\bar\mu(F)>0$ over the generating family, hence (QC$_1$) and no mirrored pair (\src{S3} Prop.~S3.2, Proposition~\ref{prop:S32}). (Correction: the earlier text said that an example must ``pass the closedness test of Proposition~\ref{prop:Q13}, which requires a compact family of target profiles that a thin family cannot be''; compactness is sufficient there, not necessary, and the necessary condition in that class is $\inf\bar\mu(F)>0$, \src{S34\_referee} \S9.3.)

\emph{How a positive two-world answer would lift to $T_1$.} Proposition~\ref{prop:Q35} shows that, under (Nest$_e$), the restriction of a three-world system to each $R_{ij}$, with the edge selection, satisfies every hypothesis of the nested two-world results; by Remark~\ref{rem:straighten}, any nested edge family with (K) and closed sections will do. Hence:
\begin{itemize}[leftmargin=*]
\item for worlds with a nonzero leading charge (Nest$_e$) has been removed: they are excluded from covers of $T_1$ on $\M_U$ outright (Theorems~\ref{thm:S21}, \ref{thm:Q36}(a), (b), Corollary~\ref{cor:Q37}). Complete worlds still need it: a proof of (R4) in its edge form, or of straightening (Conjecture~\ref{conj:S113}), would remove (Nest$_e$) from Theorem~\ref{thm:Q36}(c) and from the complete clause of Corollary~\ref{cor:Q37}, so that complete worlds would be excluded outright too. Round 14: by Corollary~\ref{cor:B32} a total world needs only to be AR on $M$, and edge nestedness or straightening is one route to that; Conjecture~\ref{conj:B3C} would exclude total worlds that do not add up with no nestedness at all;
\item a negative answer to (R1) together with a positive answer to (R2) (for instance through (R3)), i.e.\ N17.B$'$ in two nested worlds, applies verbatim on each $R_{ij}$: every world of a cover of $T_1$ with nested edges would have a strictly positive charge endorsing its edge comparisons on each $R_{ij}$. For leading charges (now without (Nest$_e$), Theorem~\ref{thm:S21}) and Savage measures, positivity globalises through the linking sets $R_{ij}\cap R_{ik}$ (Theorem~\ref{thm:Q36}); for a general acceptance cone it does not, and one further step is needed: that no $\neg$SP$^{\mathrm{reg}}$ certificate is carried only by comparisons out of interior selected models with crossing error sets (\src{Q1} Q1.C5; by \src{N17} N17.3 such a certificate, if of type (M1), uses at least four cyclic comparisons with overlapping gains). Round 14: positivity also globalises from leading densities on the two edge regions (Theorem~\ref{thm:B13}), but not from the edge comparisons alone (Proposition~\ref{prop:B1cex}). Round 15: for a world that is $\Gamma$-leading for finitely generated $\Gamma_{ij},\Gamma_{ik}$ with (UI), the further step is exactly Question~C1.Q of part~(f) (Theorem~\ref{thm:C17}, Proposition~\ref{prop:C19}). Round 16: (R2) has a negative answer and N17.B$'$ fails in two nested worlds (Theorem~\ref{thm:D21}), so this route cannot be completed as stated; its analogue with densities locally integrable on $U$ (Conjecture~\ref{conj:sigmaB}) is false too (Theorem~\ref{thm:F13}), its analogue on $M$ (Conjecture~\ref{conj:F1Cp}) is open, and the mechanism of Theorem~\ref{thm:D21} cannot carry a cover of $T_1$ in which every world is an end-protected decoration with integrable, jointly end-heavy base densities (Theorem~\ref{thm:E15}; audited in \src{F2}, audit refereed once); round 17 removes the heaviness and excludes more (Theorems~\ref{thm:H12}, \ref{thm:H13}, Proposition~\ref{prop:HR1}; part~(f$'$));
\item independently, the first case of Theorem~\ref{thm:R3} should be excluded: that at some world the comparisons out of $T_1$ admit no strictly positive endorsing shadow. This is a statement about $T_1$'s out-comparisons at one world only; Lemma~IC (\src{R11}) and the bathtub structure of $\mathcal N_i$ are the natural tools. Excluding it puts every cover in case (ii), so unendorsed relevant comparisons against incoherent $X_k\to T_1$ are always present, and Proposition~\ref{prop:Q38} then makes far dominators cost a far tie partner or a failure of (LC$_{T_1}$). Round 16: Theorem~\ref{thm:E12} (proved in a note, audited in \src{F2}, audit refereed once) gives unendorsed comparisons against incoherent models within $\eta$ of $T_1$ without excluding case~(i), with admissible shadows in place of shadows in $\mathcal N$ (the dominator may then be $T_1$ itself), and by Theorem~\ref{thm:E11} the argument of Theorem~\ref{thm:M0closed} never uses $T_1$'s out-comparisons; so excluding case~(i) is no longer needed for that purpose.
\end{itemize}
\emph{Three worlds: questions from rounds 14 to 19.}
\begin{itemize}[leftmargin=*]
\item \emph{Finitely additive gluing.} Are leading charges of world $i$ on $R_{ij}$ and on $R_{ik}$ that are positive on $Q_i$ proportional there? A positive answer would extend Corollary~\ref{cor:B14} to finitely additive edge charges (scope remark in its status).
\item \emph{Total worlds.} Conjecture~\ref{conj:B3C}: every total world of a ranking with strict (G), (IB) and (S) on $\M_U$ is AR. Its first test is the order-theoretic Conjecture~\ref{conj:B3Cp}. If it holds, every world of a cover of $T_1$ that does not add up is not total.
\item \emph{Credal edges.} Question~C1.Q of part~(f): can a non-total world that is $\Gamma$-leading for finitely generated $\Gamma_{ij},\Gamma_{ik}$ with (UI) carry a finite $\Gamma^{\mathrm{gl}}$-obstruction? A negative answer excludes every such world from covers of $T_1$.
\item \emph{Density-unanimity edges.} Conjecture~\ref{conj:C2C}: in three worlds with every axiom, a world that is a density unanimity on both of its edge regions has proportional members, hence a leading density on both, hence adds up. A counterexample needs aligned pricing of the link class and interior selected models that are Bayes for every member at once.
\item \emph{Other edge orders.} Worlds without a leading density on some edge region whose edge orders are neither density unanimities nor $\Gamma$-leading for finitely generated pairs with (UI) (for instance non-Archimedean edge orders, or $\Gamma$-leading only for pairs without (UI)) are not covered by any result above. By Corollary~\ref{cor:C25}, an exclusion of such worlds, like any proof of (LD) for $n=3$, must use an input beyond the restriction of the ranking to one edge: the third world, the link class, interior indices, or the shape of coherent edge models in $\R^3$.
\item \emph{At $T_1$ itself.} Case (i) of Theorem~\ref{thm:R3} (the third item of the lifting list above); after Theorem~\ref{thm:E12} its exclusion is no longer needed to obtain unendorsed comparisons near $T_1$.
\item \emph{$\sigma$-finite endorsement} (round 16; audit of 15 October). Conjecture~\ref{conj:sigmaB} (a countably additive endorsing measure with an a.e.\ positive density, locally integrable on $U$) is false (Theorem~\ref{thm:F13}); Theorem~\ref{thm:D21} already shows that the bounded form fails. Its repair, Conjecture~\ref{conj:F1Cp} (local integrability on $M$), is open, and its use for $T_1$ depends on (Nest$_e$) (caveat in its status). The weaker a.e.-finite target of part~(e) is open, and must be read with endorsement in the cancellation form or with finite cost of the left-hand model. In two worlds next to an additive world, whether a $\sigma$-finite endorser always exists is \src{F1}'s Conjecture~F1.C. For $T_1$ the question is whether Proposition~\ref{prop:M0} survives for such measures (status of Conjecture~\ref{conj:sigmaB}), and whether Theorems~\ref{thm:E11} and~\ref{thm:E15} extend to shadows that are non-integrable only at ends. Round 17 proves this when only one world's density is non-integrable and the other two worlds have integrable full endorsers (Theorem~\ref{thm:H12}); it is open for two or more such worlds (case (D) of part~(f$'$)). Round 19 proves it for several such worlds whenever a patch exists (Theorem~\ref{thm:K35}, Corollary~\ref{cor:K35p}); what is left is sheltering and infinite minority mass.
\item \emph{Conflict cycles, and the cases outside class (EP)} (rounds 17 and 19; \src{H1}, \src{H2}, \src{H12\_referee}, \src{J1}, \src{J1\_referee}, \src{K1}, \src{K1\_referee}). Question~J1.Q: can a doubly misaligned $2$-cycle whose third world is incident to no arrow carry the shell exchanges of Proposition~\ref{prop:J15} under every axiom, for every tuple of integrable base endorsers? Question~J1.Q$'$: can a ranking in class (EP) in which no order of the worlds is aligned cover $T_1$? Question~J1.Q$'_\infty$: the same for class (EP$_\infty$) with (Al$_\infty$). \src{J1} states that a negative answer to J1.Q would exclude every cover in class (EP); that inference is unproved (\src{J1\_referee} \S8.2): a negative answer closes only case (R-i), and cases (R-ii) and (R-iii) remain. No conjecture is stated. Outside class (EP): case (A$'$), case (B) with countably many ends, case (D), and protection not of end type. \src{D2}'s Conjecture~D2.C (installing the world of Theorem~\ref{thm:D21} on an edge region) is open; next to two worlds with integrable full endorsers, an installation that keeps an a.e.\ finite endorser of its base in the cancellation form is excluded (Theorem~\ref{thm:H12}), with no condition of heaviness or compatibility.
\item \emph{Worlds without a localizable endorser, and finite certificates} (round 19; \src{K2}, \src{K3}, \src{K2\_referee}, \src{K3\_referee}). Question~Q$_{\mathrm{K3}}$: can a world of a full system have a non-null non-localizable blind set $\Lambda_\infty$? Its two-world form next to an additive world is Conjecture~\ref{conj:K3C}, whose evidence is weaker than \src{K3} stated (status there). Question~Q$_{\mathrm A}$: can an integrable field with one component vanishing on a set of the rigid shape of Theorem~\ref{thm:K37}(b$'$) have property ($\gamma$) and $T_1$ as a Bayes model at an interior weight? Proposition~\ref{prop:KQA} gives necessary conditions and locates where the proof of Theorem~\ref{thm:M0closed}(a) breaks. Question~Q-K2: is there a full system with two worlds additive with integrable densities whose third world has a certificate? Conjecture~\ref{conj:K26p}: a Bayes-shaped selection over such worlds forces the third world to add up; for the Brier base, by Proposition~\ref{prop:KR62}, it reduces to the passage from the closed convex hull of the forced cone to the acceptance set.
\end{itemize}
\emph{Off the $U$-domain.}
\begin{itemize}[leftmargin=*]
\item Additive rankings with integral penalties are excluded on every decomposable domain with (LF) (Theorem~\ref{thm:gm}); with bounded charge penalties, given interior Bayes minimisers (Corollary~\ref{cor:gmcharges}).
\item Rankings whose worlds all add up: on general decomposable domains Claim~B2.0 (the transfer of Proposition~\ref{prop:M0}) must first be written out. Conditional on it, apex ties at the target weighting are excluded on every decomposable domain (Corollary~\ref{cor:B22}). Round 16: under (LF) the full-atom residual is empty, for every strictly guiding integral ($\gamma$)-field, independently of Claim~B2.0 (Theorem~\ref{thm:D1}, refereed once, which proves the restated conjecture in the status of Proposition~\ref{prop:B2ex}); without (LF) it occurs (Propositions~\ref{prop:B2ex}, \ref{prop:D10}). For rankings this holds at the status of Claim~B2.0 with the reduced list of imports of Corollary~\ref{cor:D17}. What remains under (LF) is the branch $p^*\notin E$: the case $E=\varnothing$, at the status of Theorem~\ref{thm:ccibu} or after an audit of the routes of Theorem~\ref{thm:gm} for fields with Bayes uniqueness everywhere; and the case $E\ne\varnothing$, where route~I1 survives only conditionally on the unaudited inputs listed in part~(f), or is open if $E$ accumulates at $p^*$ (\src{D1}'s proposed next step, a CONJECTURE, is in part~(f)).
\item (CP-a) rankings beyond $\M_U$: Conjecture~\ref{conj:U} and Corollary~\ref{cor:residual} remain the live statements (scope of Corollary~\ref{cor:CPaU}). Rankings with a world that does not add up have not been examined off $\M_U$. The world-level results Theorem~\ref{thm:S21}, Lemma~\ref{lem:B33}, Theorems~\ref{thm:B35} and~\ref{thm:C23} are stated on $\M_U$ for any $K'\subseteq\Di$, not on other domains.
\end{itemize}
\emph{Status:} (R1) open, with (R1$'$) the weaker target of (R1); (R2) answered negatively in round 16 (Theorem~\ref{thm:D21}); (R3) restated as Conjecture~\ref{conj:NDL}, which holds next to an end-vanishing endorsing charge, and its linear form is false; (R4) proved in two worlds when one world has a nonzero leading charge or is monotonely continuous, open otherwise, and in its edge form proved in the cases of Proposition~\ref{prop:RS2A} and open otherwise (straightening, Conjecture~\ref{conj:S113}), and for total worlds now only one route to AR (Corollary~\ref{cor:B32}). The lifting statements are as stated in Proposition~\ref{prop:Q35} and Theorems~\ref{thm:Q36} and~\ref{thm:B13} (proved; complete worlds CHECKED), and the further interior step is open in general, and is Question~C1.Q for $\Gamma$-leading worlds. Finitely additive gluing, Conjectures~\ref{conj:B3C}, \ref{conj:B3Cp}, \ref{conj:C2C}, Question~C1.Q and Claim~B2.0 are open. Round 16: (R2) is answered negatively, and N17.B$'$ (two worlds, also nested) and Q1.C1 are false (Theorem~\ref{thm:D21}, confirmed by two referees); the restated full-atom conjecture is proved under (LF) (Theorem~\ref{thm:D1}, refereed once); Theorems~\ref{thm:E11}, \ref{thm:E12}, \ref{thm:E15}, Lemmas~\ref{lem:E13}, \ref{lem:E14} and Proposition~\ref{prop:E16} are proved in a note, audited in \src{F2} (which states their exact hypotheses), and the audit was refereed once. Audit of 15 October: Conjecture~\ref{conj:sigmaB} is false (Theorem~\ref{thm:F13}, refereed once); Lemma~\ref{lem:F11}, Theorems~\ref{thm:F12}, \ref{thm:F14}, \ref{thm:F16}, \ref{thm:F18}, Corollary~\ref{cor:F15} and Proposition~\ref{prop:F21} are proved and refereed once; Conjecture~\ref{conj:F1Cp} and the a.e.-finite target are open. Round 17: Lemmas~\ref{lem:H10}, \ref{lem:H11}, \ref{lem:H20}, \ref{lem:H21}, Proposition~\ref{prop:H22}, Theorems~\ref{thm:H12}, \ref{thm:H13}, \ref{thm:H23}, \ref{thm:H24} and Corollary~\ref{cor:H23p} are proved and refereed once (\src{H12\_referee}), with the referee's repairs; Proposition~\ref{prop:HR1} is proved in that referee report, one reader, and independently in \src{J1}; Lemmas~\ref{lem:J1C}, \ref{lem:J10}, Theorem~\ref{thm:J12}, Corollary~\ref{cor:J13} and Proposition~\ref{prop:J15} are proved in a note and, in round 19, confirmed once (\src{J1\_referee}), with the referee's repairs. Round 19: Theorem~\ref{thm:JR2} and Corollary~\ref{cor:JR3} are proved in a referee report (\src{J1\_referee}), one reader; Theorems~\ref{thm:K31}, \ref{thm:K33}, \ref{thm:K35}, \ref{thm:K37}, Propositions~\ref{prop:K32}, \ref{prop:K38}, Lemmas~\ref{lem:K34}, \ref{lem:K36}, \ref{lem:K37p} and Corollary~\ref{cor:K35p} (\src{K3}), Theorems~\ref{thm:K21}, \ref{thm:K23}, \ref{thm:K24}, \ref{thm:K25}, \ref{thm:K26}, Propositions~\ref{prop:K22}, \ref{prop:K26a}, Corollary~\ref{cor:K22p} and Lemma~\ref{lem:K27} (\src{K2}), and Lemmas~\ref{lem:K10}, \ref{lem:K11}, Theorems~\ref{thm:K12}, \ref{thm:K14}, \ref{thm:K15}, \ref{thm:K16} and Corollary~\ref{cor:K13} (\src{K1}) are proved in notes and confirmed once, with the referees' repairs; Propositions~\ref{prop:KR62} and~\ref{prop:KQA} are proved in referee reports, one reader. Questions~J1.Q, J1.Q$'$, J1.Q$'_\infty$, Q$_{\mathrm{K3}}$, Q$_{\mathrm A}$ and Q-K2 and Conjectures~\ref{conj:K3C} and~\ref{conj:K26p} are open, and so are the cases (R-i)--(R-iii) of class (EP).

\subsection{Other targets and more worlds}\label{sec:open-other}
\begin{itemize}[leftmargin=*]
\item \emph{Other triangles.} Theorem~\ref{thm:main} covers every all-facing triangle in general position. For the crossing classes $(1,2,3)$ and non-facing $(2,2,2)$ the analogue is \src{R15} Conjecture R15, supported by numerics on about 1340 random triangles ($\min\Lambda_{\max}=1.0045$ for $(1,2,3)$, $1.2187$ for $(2,2,2)$; corrected values in \src{R18} \S7). The first-order arguments of \S\S4.1--4.5 apply facet by facet wherever a facet has a crossing with a shared positive coordinate; the pinching algebra is specific to the all-facing class, and for $(1,2,3)$ a nondegeneracy at the one-crossing facet is needed (\src{R18} \S5). \emph{Status:} open.
\item \emph{Segments and boundary targets} ($T_3$, $T_4$): open. Several lemmas need $\dim K=2$ (\src{R12} \S1 repair; \src{R1} Lemma 3.5).
\item \emph{$n\ge4$.} Not addressed. Vertex confinement, Lemma UVN, Lemma M, \src{R10b} Thm.~PL and Cor.~PL2 hold in every $n$; Theorem~\ref{thm:ccibu}, the box/corner system and the pinching are three-world arguments. On all-model domains, v3's additive obstruction excludes full additive coverage of targets with three vertices of $\Delta$ (\src{Full Coverage \dots v3}, App.~C).
\end{itemize}

\section{Positive results and approximation}\label{sec:positive}
\begin{sloppypar}
\begin{itemize}[leftmargin=*]
\item \emph{Gauge mixtures.} For compact convex $M_i\subseteq\R^3_{\ge0}\setminus0$ with $a_ie_i\in M_i$ and $\nu(M_i)\subseteq\co(\nu(M_j)\cup\{e_i\})$ ($\nu$ = normalisation), the penalties built from $f_i(h)=\min_{m\in M_i}m\cdot h$ give an additive system on a decomposable domain whose Bayes model at every $p\in\Delta$ is unique and equals $D_{K_p}$, $K_p=\nu(\sum_ip_iM_i)$; every interior $K_{p^*}$ is covered with (LF) (\src{R10a} Thm.~M; refereed once, \src{R12}). Examples include a regular hexagon around $u$ and 9-gons converging to $T_1$ in the Hausdorff metric (\src{R10a} Prop.~M$'$; refereed once, \src{R12}). No gauge mixture covers a triangle in general position (\src{R10a} Thm.~R, ray lemma; refereed once).
\item \emph{Other covered classes.} Precise models; normalised boxes and their conditional closures (\src{Conditional Shunning of Error.tex}); every coherent target when $n=2$ (\src{Full Coverage \dots v3}); the vacuous model under its accepted exception.
\item \emph{Approximation.} For an interior triangle $L$ and its parallel bodies, $\lp_{L+rB}(g)=\lp_L(g)-r\|P_Hg\|$, so $D_{L+rB}\uparrow D_L$ as $r\downarrow0$; parallel bodies are covered (CPM manuscript v19, cone domains). This is a genuine approximation result, but the representable class is not closed in the target topology: flat facets appear only at $r=0$ (\src{CONTAMINATION\_NOGO\_SCOPE\_AND\_GEOMETRY.md} \S\S4--5).
\item \emph{Earlier no-go for continuous maps.} For continuous, positively homogeneous coherence-preserving maps with weak, generically strict coordinate signs, the three-state contamination triangle at any interior baseline has no realiser at any interior prior, by the same facet maps, corner equations and box as in Proposition~\ref{prop:T1} (\src{CONTAMINATION\_NOGO\_SCOPE\_AND\_GEOMETRY.md} \S2; earlier work, not re-refereed here). In the ranking setting, on domains non-adjacent near the target, Theorem~\ref{thm:main} needs no continuity assumption; on general decomposable domains with (LF) the answer is also negative, again with no continuity assumption (Theorem~\ref{thm:gm}).
\end{itemize}
\end{sloppypar}

\begin{thebibliography}{9}
\bibitem{HU1977} F.~Hiai and H.~Umegaki, Integrals, conditional expectations, and martingales of multivalued functions, \emph{Journal of Multivariate Analysis} 7 (1977) 149--182.
\bibitem{W1991} P.~Walley, \emph{Statistical Reasoning with Imprecise Probabilities}, Chapman and Hall, London, 1991.
\bibitem{W2000} P.~Walley, Towards a unified theory of imprecise probability, \emph{International Journal of Approximate Reasoning} 24 (2000) 125--148.
\bibitem{Sav1954} L.~J.~Savage, \emph{The Foundations of Statistics}, Wiley, New York, 1954.
\bibitem{Fis1986} P.~C.~Fishburn, The axioms of subjective probability, \emph{Statistical Science} 1 (1986) 335--358.
\bibitem{V1964} C.~Villegas, On qualitative probability $\sigma$-algebras, \emph{Annals of Mathematical Statistics} 35 (1964) 1787--1796.
\bibitem{MLS1987} J.-E.~Mart\'inez-Legaz and I.~Singer, Lexicographical separation in $\R^n$, \emph{Linear Algebra and its Applications} 90 (1987) 147--163.
\end{thebibliography}

\section*{Sources}\label{sec:sources}
\addcontentsline{toc}{section}{Sources}
Project notes, in the working folder unless stated. Short codes in the text refer to these files.
\begin{description}[leftmargin=*,font=\normalfont\ttfamily,style=nextline]
\item[BRIEF.md, BRIEF2.md, BRIEF3.md--BRIEF8.md] Setting, axioms, (LF), decomposable domains, the realisation problem, weakened guidance, (CP-a).
\item[BRIEF9.md] Round 10: the organised attack on CF$(T_1)$ (construction or impossibility, no added assumptions).
\item[BRIEF10.md] Round 11: the positive ordinal attempt (build non-additive rankings covering $T_1$ exactly, or find a genuine obstruction).
\item[BRIEF11.md] Round 12: does a world that does not add up exist inside a full system satisfying every axiom (construction or proof, any target, any $n\ge2$)?
\item[BRIEF12.md] Round 13: edge nestedness (R4) and mirrored simple pairs (R1), from both sides.
\item[BRIEF13.md] Round 14: ambitious attempts at an exact cover of $T_1$ (construction with every axiom, or one honest failure report).
\item[BRIEF14.md] Round 15: the escape left by Corollary~B1.4; glued credal worlds (C1) and the leading-density statement (LD) (C2).
\item[BRIEF15.md] Round 17: an ambitious positive attempt for $T_1$ built on the lessons of rounds 14--16 (move the corner conflict into the compatibility of the worlds' endorsing densities).
\item[BRIEF16.md] Round 19: settle $T_1$ on the $U$-domain, or find the mechanism a cover must use (a cover, a settling theorem, or a single proved obstruction or precise open question).
\item[Coverage after Neutral Abstention.tex] The exact consistency test (Thm.~closure, \S3) and normalised boxes (Generalized dominance paper folder).
\item[Vacuous Targets and Transfinite Construction.tex] \S5: unions, the reserved-neighbourhood invariant, the extension lemma (same folder).
\item[Three Complete Rankings Cannot Cover.tex, Complete Rankings for Any Number of Worlds.tex] No complete ranking with (G), (IB), (S) covers the vacuous model, $n=3$ and $n\ge3$ (same folder; cited at the level of their statements).
\item[Loss Dominance Generalized v2.tex] Rankings as preorders, competitiveness, (E), (LA), (A), (C), (S), KKM dominance theorem (Generalized dominance paper folder).
\item[Full Coverage by Homothetic Credal Families v3.tex] Binary problems, four-rank guidance order, (G), (IB); two-world coverage; all-model obstruction (same folder).
\item[Consolidated Statements - Decomposable Domains 2026-10-12.md] Status of R1--R9 (R1 table of charts for each advice range; chart theorem, continuity; R2 Lemma U; R4 Thms A, B, Prop.~G1).
\item[R1\_notes.md, R2\_notes.md, R4\_notes.md, R6\_referee.md, R8\_referee.md] Additive rankings on decomposable domains (charts, Bayes uniqueness, adjacency, ray equations); referee reports.
\item[R10a\_notes.md, R10b\_notes.md, R12\_referee.md] Gauge mixtures, ray lemma; vertex confinement, Lemma UVN, Lemma M, Thm.~PL, Cor.~PL2; referee report.
\item[R11\_notes.md, R13\_referee.md] Lemma IC, Lemma K, Thm.~W; referee report.
\item[R15\_notes.md, R18\_referee.md] Corner calculus, all-facing triangles, box and Thm.~6.2; referee report.
\item[CCIBU.tex] Closed choice from interior Bayes uniqueness in three worlds.
\item[N12\_notes.md, N14\_referee.md] Reduction to a credal family, chords, edge box, (PC)+(NS) $\Rightarrow$ contradiction, two-slope projectivity; referee report with Thm.~C2 and the non-adjacent reduction.
\item[N15\_notes.md, N18\_referee.md, N20\_referee.md] Straight levels, pinning, pivot chords, no sliding edges, slope--rotation inequalities, two-stage blow-up; two independent referee reports.
\item[N21\_notes.md, N23\_referee.md] General decomposable domains; referee report with Lemma R23.1 and Cors.~R23.2--R23.4.
\item[N11\_notes.md, N13\_referee.md, N22\_notes.md, N24\_referee.md] Coherent almost-preference rankings; referee reports with Props A, B, R1, R2.
\item[N1--N7, N8--N10, N16, N17\_notes.md, N19\_referee.md] Weakened guidance, non-additive constructions, adversarial attempts, rankings that do not add up.
\item[N25\_notes.md, N25\_referee.md] The exceptional-tie case M0: structure of $E(\psi)$, coherent envelopes, atoms, genericity, consequences of (E$_\psi$), the target at an exceptional weighting; referee report with Prop.~F (no fan points) and corrected residuals R-a$'$, R-b$'$.
\item[N26\_notes.md, N26\_referee.md] Resolving exceptional ties lexicographically: the two-world construction, the class $\Ltwo$ in three worlds; referee report with Prop.~R1 and the sharpened Conjecture N26.R.
\item[N27\_notes.md, N27\_referee.md] The compact-collar repair of $N_0$, ray transfer of charts, the main theorem under permitted ties; referee report.
\item[N28\_notes.md, N28\_referee.md] Competing acceptances on general (LF) decomposable domains: (TS) in every cover, cut apexes, vertex-free interiors, CF$(T)$, diagonal symmetries; referee report with Lemma R1, Prop.~R3 and the sector-cap example C1.
\item[TS\_second\_referee.md] Second, independent referee report on Thm.~TSalways and its inputs, Cor.~TSalways, Thm.~cuts, Thm.~apex (wording repair), Lemma R1, Prop.~vertexfree and Prop.~CF($\Rightarrow$) in its every-$q$ form (repairs); the two corrections to \S\ref{sec:general-open}.
\item[I1\_notes.md, I1\_referee.md] Route I1: edge maps, exact edge limits, pinned chart lines, forced supporting lines, the middle window and the three-line determinant; no cover of $T_1$; referee report (CONFIRMED, repairs R1--R7).
\item[I2\_notes.md, I2\_referee.md] Route I2: no straddling on the positive side of a chord, positive state near the vertices, the precise collar, transport of one weight, transition points and the blow-up contradiction; referee report with repair R1 and Prop.~D$^*$.
\item[C2\_notes.md] Round 15 (top of the file): (LD) is false (Thm.~C2.1), edge comonotonicity (C2.3), edge-only arguments (C2.5), link rigidity (C2.6, C2.7), Conjecture C2.C; refereed in C12\_referee.md. Archive (round 10): Route C2: discretisations DCF and DCF$^+$ (NUMERICS), no straddling, chord structure, forced caps, gluing, pinning of the cap weights, two caps cannot coexist.
\item[C1\_notes.md] Round 15 (top of the file): glued credal worlds, the dichotomy (C1.1), the closure as unanimity over the glued credal set (C1.3, C1.4), the design excluded (C1.6), the band (C1.7--C1.9), Question C1.Q; refereed in C12\_referee.md. Archive (round 10): construction side: gauge mixtures are split, no injective precise open region (Thm.~A), no locally definable solution (Thm.~B), wild-solution analysis.
\item[C1C2\_referee.md] Referee report on C2 and C1: Prop.~P, hand elimination for all $x>1$ (empty exceptional set), repaired Thm.~A, Theorem~R with two routes; adjudication of the chord claim.
\item[C12\_referee.md] Referee report on the round-15 notes C1 and C2: everything confirmed; repairs G1--G3 of C2.1, the wording of C2.4, C2.5, C2.8, the scope of C1.6, strict fineness in C1.8, the closing remark of C1.9 marked unproved; the single target for C1.Q and C2.C.
\item[PropD\_notes.md] Theorem D (precise regions are constant or diagonal-projective, under the d-line condition only), Theorem T (transport), the empty-interior case; not yet refereed.
\item[P1\_notes.md] Round 11, case (M0) from the construction side: straight levels for $K^\flat$, one-sided corner inequalities plus box (Lemma P1.1), corner atoms under (HC), split facets (NUMERICS).
\item[P2\_notes.md] Round 11: vector form of the test, holonomy (NUMERICS), flexible worlds, Theorem LX, rigid acceptances, nested criterion, design attempts.
\item[P3\_notes.md] Round 11: the extension lemma; exchange transport, the limit rule and the exact extension theorem, compactness, two counterexamples, the witnessed extension theorem, the Villegas-type conjecture.
\item[P4\_notes.md] Round 11: Part A, the chord-tolerant chain (quasi-linearity, ball lemma, PL2 across chords), $p^*\in E$, carrying atoms, the short argument; Part B, N17.B$'$ in two worlds with a nested selection (strict escape, Thm.~B3, Prop.~B2, removal lemma).
\item[M0\_referee.md] Referee report on P4 Part A and P1; Prop.~R.1 (crossing pinning), Prop.~R.2, Theorem R (no ranking that adds up covers $T_1$ on $\M_U$).
\item[M0\_second\_referee.md] Second, independent referee report on Prop.~R.1 (strengthened), Prop.~R.2 and Theorem R; what singles out $T_1$; completed dependency list.
\item[P2P3\_referee.md] Referee report on P3 and P2: confirmations, the narrowing of the counterexamples, the replacement of P2.3, the class repair of Theorem LX.
\item[P4B\_referee.md] Referee report on P4 Part B: Lemma N0, repaired Prop.~B2, Lemma T, Theorems E2 and E3, the precise open class in two worlds.
\item[CPaU\_referee.md] Referee report on Remark CPaU (now Corollary~\ref{cor:CPaU}) and the last sentence of Remark~\ref{rem:EPchords}: the field $c^P$, the translation into Theorem~\ref{thm:M0closed}(a), the logic of (U), scope, Route B (CP-a worlds add up), the corrupted \texttt{\textbackslash ref}.
\item[Q1\_notes.md] Round 12, construction side: cut form, band plus exotic cones (every axiom except closed sections at one world), exact sections, mirroring (Q1.4), spread and star certificates, edge transfer (Q1.7--Q1.10), failure reports, conjectures Q1.C1--C5.
\item[Q2\_notes.md] Round 12, exclusion side in two worlds: Lemma T$_2$, left-fineness, initial segments, mirror pairs, MP$'$, no additive scaffold, necessary conditions, the non-nested lattice, G-SE, Conjecture Q2.N.
\item[Q3\_notes.md] Round 12, exclusion side for $T_1$: locality barrier, shadows, edge complementarity and edge reduction, three-world transfer, Corollary Q3.7, far dominators and tie toggles.
\item[Q3\_referee.md] Referee report on Q3, Q1.7--Q1.10 and Q2 \S6: confirmations and repairs; Lemma R1, Prop.~R2, Theorem R3; (Nest$_e$) as comonotonicity; corrected Conjecture Q2.E.
\item[Q12\_referee.md] Referee report on the two-world parts of Q1 and Q2: confirmations, the repair of Q1.2 (Fact 0 not available) and Q1.4, the exact section in the dyadic family; new results N1--N4; the precise remaining question (R1)--(R4).
\item[S1\_notes.md] Round 13, (R4) proof side: Level Lemma, the continuum $\C_L/{\sim}$, nested representatives, rigid crossings, monotone continuity forces a chain, the edge results under (EU$_{ij}$), pairwise straightening, Conjectures S1.13 and S1.R.
\item[S1\_referee.md] Referee report on S1: confirmations, the repair of S1.6 (each half of (MC$_{\mathrm{sel}}$) suffices), the corner-pair strengthening of S1.5, the repaired necessary shape, Remark S1.8 with repair, Prop.~R7.1 (only (C) fails), the combined two-world statement.
\item[S2\_notes.md] Round 13, (R4) construction side: the tie remark made precise (S2.0), the parallelogram lemma, Theorem S2.1 (leading charges without nestedness), Theorem S2.2, three-world anatomy (S2.4), failure reports.
\item[S2\_referee.md] Referee report on S2: S2.1 and S2.2 confirmed; repairs of PL, S2.0(a), (c), B2(b) and S2.4; Prop.~R-S2.A; the LX-class corollary; the corrected tie remark.
\item[S3\_notes.md] Round 13, (R1) construction side: mirrored pair without closed sections (S3.1), exact sections of band-and-translates worlds (S3.2), the packing lemma (S3.3), the fibrewise route (S3.4), Conjecture QC.
\item[S4\_notes.md] Round 13, (R1)--(R3) exclusion side: Theorem RATE, Corollary SR, MP$\sharp$, the loop lemma, Example~A, Prop.~LIN, Lemma UE, the (R1) failure report, Conjecture S4.1.
\item[S34\_referee.md] Referee report on S3 and S4: confirmations and repairs (S3.2, MP$\sharp$(d), the loop lemma, LIN); Prop.~QC-J; the blocking lemma; the corrected (R3) item of \S\ref{sec:open-core}(g); the precise remaining form of (R1)--(R3).
\item[B1\_notes.md] Round 14: sign classes and the link class, Lemmas B1.0--B1.2 (gluing, transfer chain), Theorem B1.3, Corollaries B1.4, B1.5; failure report on absorbing the corner cycle.
\item[B1\_referee.md] Referee report on B1: all core results confirmed; the strengthening of B1.4 via S2.1 (CHECKED); corrections of B1.5(a), (b); the counterexample to the headline of B1 \S0; scope remark on finitely additive gluing.
\item[B2\_notes.md] Round 14: tie anatomy on two-sided domains, Theorem B2.1 (face models at a target weighting), Prop.~B2.0, Cor.~B2.2, Prop.~13, the residual (R2), the audit of the branch $p^*\notin E$.
\item[B2\_referee.md] Referee report on B2: Theorem B2.1 confirmed with repairs (Lemmas 7, 9), upgraded to every all-facing triangle; Prop.~B2.0 a documentation GAP; tie anatomy at $p^*$; a field realising (R2) without (LF); corrected audit.
\item[B3\_notes.md] Round 14: total worlds, Lemma B3.1, Cor.~B3.2, order continuity B3.3, the ultrafilter example B3.4, Theorems and Props.~B3.5--B3.7, Conjecture B3.C.
\item[B3\_referee.md] Referee report on B3: confirmations and scope repairs; Prop.~R1 (two-world realisation without (C)); Remark R2 (leading top levels); Prop.~R3 and Conjecture B3.C$'$.
\item[D1\_notes.md] Round 16: corner atoms through $p^*$ against (LF); directional families on two-sided domains, two-sided frozen support, walls across chords, Theorem~D1, Example~D1.0, Corollary~D1.7.
\item[D1\_referee.md] Referee report on D1: Theorem~D1 confirmed after five repairs (genericity, measure bookkeeping, non-collinearity, the open-chart replacement of Step~1 of Lemma~M); Corollary~D1.7 with the reduced import list; the case $E=\varnothing$.
\item[D2\_notes.md] Round 16: a full two-world system on $\M_U$ in which a world does not add up (designs from projective planes, the protected end, Theorem~D2.1, Lemmas~D2.2--D2.4); Conjecture~D2.C.
\item[D2\_referee.md] Referee report on D2: everything confirmed; bounded weights, the case $b=1$, the order of bounds in D2.4; heuristic remarks relabelled; direct proof of type (M2).
\item[D2\_referee2.md] Second, independent referee report on D2: everything confirmed; the $\sigma$-finite endorsement (failure exactly at the protected end), the suggested $\sigma$-finite conjecture (false as written: F1, F2), a direct proof of $\C_L=\{C_b\}$, why Prop.~S3.2(d) does not apply.
\item[E1\_notes.md] Round 16 (audited in F2; the audit refereed once in F12 referee): what Theorem M0closed needs (one admissible shadow's dominator comparisons), Theorem~R3 without its case split, ends and end-protection, end-heavy shadows, the exclusion of end-protected decorations, the invisibility of the failure in D2.
\item[F1\_notes.md] Audit of 15 October: the geometry of the end of the chain of D2 (two rays, one interior to $U$); readings of endorsement and property (R); chain monotonicity of $\sigma$-finite endorsers next to an additive world (F1.2); the refutation of Conjecture sigmaB (F1.3); no compact-target certificate next to an additive world (F1.4, F1.5); the systems D2$'_{\mathrm A}$, D2$'_{\mathrm B}$ (F1.5$'$--F1.8); Conjectures F1.C and F1.C$'$.
\item[F2\_notes.md] Audit of 15 October: the exact hypotheses of E1.1--E1.6; where ends can sit (F2.1--F2.3); E1.4 and E1.5 at $\partial U$; a second refutation of Conjecture sigmaB (F2.4); B3 on total worlds; corrected summary statements and the open route.
\item[F12\_referee.md] Referee report on F1 and F2: nothing wrong at theorem level; repairs (the last sentence of F1.5 is a necessary condition only; the wording of F1.6; the finiteness step of F1.9; the caveat on F1.C$'$; the reading of endorsement under a.e.\ finiteness; the ends sentence and the B3 sentence; item~5 of the open route restated).
\item[H1\_notes.md] Round 17: the stability lemma (H1.0), automatic admissibility (H1.1), one exceptional world with a localized endorser (H1.2), several worlds with $T_1$-localized endorsers (H1.3), Cor.~H1.4; failure report on moving the corner conflict into the compatibility of endorsing densities.
\item[H2\_notes.md] Round 17: refined ends (H2.0), conflict ends and shadowing (H2.1), joint reweighting and (h$'$) (H2.2), the conflict digraph (H2.3, H2.3$'$), the template of a conflict cycle and Question H2.Q, the mirrored two-world system in which neither world adds up (H2.4).
\item[H12\_referee.md] Referee report on H1 and H2: nothing wrong; the omitted bound in H1.0(b); Theorem M0closed(a) without integrability; the narrowed replacement of (h); the corrected list of what survives; the reading in H2.1(b); Prop.~R.1 ((h$'$) unnecessary); the corrected H2.3$'$; the corrections of H2 \S5.2; every axiom of H2.4 re-derived; the precise residual.
\item[J1\_notes.md] Round 17 (refereed once in round 19, J1 referee): near-Bayes models, the closure lemma, H2.3 without (h$'$) (J1.1), pattern switching under alignment (J1.2), closed configurations (J1.3), misaligned arrows (J1.4), the shell exchanges (J1.5), Question J1.Q, scope limits.
\item[J1\_referee.md] Round 19, referee report on J1: nothing wrong; J1.3(a) strengthened; one-line gaps in J1.3(c) and J1.5 closed; the radii in J1.4; the wording of J1.Q; the closure claim unproved (a negative answer closes only a $2$-cycle with an isolated third world); Theorem~R.2 and Corollary~R.3 (alignment of orders from alignment of arrows); the residual (R-i)--(R-iii); Question J1.Q$'$.
\item[K1\_notes.md] Round 19: class (EP$_\infty$) (countably many ends); safe closure (K1.0), shadowing for countable families (K1.1, K1.1$'$), per-world safe tests (K1.2), the KKM form (K1.3), acyclic digraphs (K1.4), aligned orders (K1.5), anti-safe selection indices and Question K1.Q (K1.6).
\item[K1\_referee.md] Referee report on K1: nothing wrong in K1.0--K1.6; the equivalence with K1.Q tautological; the sufficient condition after K1.5 repaired (uniform reading); a persistent conflict without an arrow (configuration, CHECKED); the residual Question J1.Q$'_\infty$.
\item[K2\_notes.md] Round 19: finite certificates in full systems; the divisibility theorem (K2.1), forced comparisons (K2.2, K2.2$'$), two worlds (K2.3), the frontier family (K2.4), the decoration scheme (K2.5), the edge obstruction (K2.6, K2.6a), the edge sweep (K2.7), Conjecture K2.6$'$, Question Q-K2.
\item[K2\_referee.md] Referee report on K2: nothing wrong; K2.2$'$ in terms of $\rho_j$; crossing needs interleaving; only (SPREAD) proved for Brier; Proposition~R6.2 (every endorser of the edge trades over Brier is $c\rho_1$ on $M$); overstatements in the readings removed.
\item[K3\_notes.md] Round 19: the Yosida--Hewitt classification (K3.1), end-protection adds nothing (K3.2), the tilt world completed (K3.3), pairwise-meeting protection (K3.4), the patch theorem (K3.5, K3.5$'$, K3.6), one world next to two integrable worlds (K3.7$'$, K3.7), two worlds next to an additive world (K3.8), Conjecture K3.C, Questions Q$_{\mathrm A}$ and Q$_{\mathrm{K3}}$.
\item[K3\_referee.md] Referee report on K3: nothing wrong; (NC) made explicit; $\hat\Theta$ and $\rho_\infty$ not canonical; K3.3 credited to Q12 referee N3; the readings of K3.4 repaired; the sharpening under Q$_{\mathrm{K3}}$ = No; necessary conditions for Q$_{\mathrm A}$; the evidence for K3.C corrected.
\item[Independent Rankings Without Additive Structure.tex] Complete rankings whose first level is a non-representable comparative probability over ultrafilter signatures (Generalized dominance paper folder).
\item[START HERE -- Coverage Status.md] Glossary (terminology of \S\ref{sec:addup-def}).
\item[Auditor's note (14 Oct 2026)] Corrections adopted by the author: relevant comparisons, the content of the additive reduction, non-adjacency, ``supporting additive assessment'', and the (M1)/(M2) wording.
\item[w1991.txt] Extracts of Walley (1991) \S\S3.7--3.8 (lines 8200--8600).
\item[CONTAMINATION\_NOGO\_SCOPE\_AND\_GEOMETRY.md] Continuous CPM no-go for contamination triangles (CPMs folder; earlier work).
\end{description}

\end{document}
